\nwfilename{nw/many-sorted-term.nw}\nwbegindocs{0}% -*- mode: poly-noweb; noweb-code-mode: sml-mode; -*-% ===> this file was generated automatically by noweave --- better not edit it
%%
%% many-sorted-term.nw
%% 
%% Made by Alex Nelson <pqnelson@gmail.com>
%% Login   <alex@lisp>
%% 
%% Started on  2026-04-04T18:09:28-0700
%% Last update 2026-04-04T18:09:28-0700
%% 

\section{Syntax for terms}

\begin{node}
Since a proof assistant resembles an interpreter in many ways, it
should not be surprising that we will first focus on implement
abstract syntax trees for terms and formulas for three-sorted
first-order logic. 
\end{node}

\begin{node}
The grammar for a term in $FS_{0}$'s three-sorted logic is, more or
less, the same as in single-sorted logic. Since we are working with
\emph{three} sorts, we isolate it in is own module. We would have the
BNF for such things be:
\begin{center}
\begin{tabular}{rcl}
$\langle$\textit{sort}$\rangle$ & $::=$ & \texttt{Ind} $\mid$ \texttt{Fun} $\mid$ \texttt{Class}\\
$\langle$\textit{term}$\rangle$ & $::=$ & $\langle$\textit{variable}$\rangle$ \\
& $\mid$ & $\langle$\textit{function}$\rangle$
\end{tabular}
\end{center}
Remember, a ``function'' here should be thought of as a ``term constructor''.

We've abstracted away the ``concrete syntax'' for the sorted
first-order logic, so we can treat variables as an ordered pair
consisting of a lexeme (string representation of its identifier) and
its sort. Similarly, we can treat a function as an ordered triple
consisting of a lexeme, its sort, and a list of terms for its
arguments.

Before getting too far ahead of ourselves, let us implement a simple
encoding for the sort. We will just make everything public, because
there's not much to hide. There are three data constructors (one for
each sort), a comparison function, and a function to get the string
representation for the sort.

Also we should bear in mind the user may have a ``mismatch'' between
the sorts expected for arguments of predicates and functions, and the
arguments supplied for them. This situation should raise an
exception. We will create a {\Tt{}Sort.Mismatch\nwendquote} exception, with a string
message for what went wrong.
\end{node}

\nwenddocs{}\nwbegincode{1}\sublabel{NW1rTmOJ-3ZqggC-1}\nwmargintag{{\nwtagstyle{}\subpageref{NW1rTmOJ-3ZqggC-1}}}\moddef{Sort.sml~{\nwtagstyle{}\subpageref{NW1rTmOJ-3ZqggC-1}}}\endmoddef\nwstartdeflinemarkup\nwenddeflinemarkup
structure Sort = struct
exception Mismatch of string;
datatype t = IND | FUN | CLASS;

(* compare : \nwlinkedidentc{Sort.t}{NW1rTmOJ-3ZqggC-1} * \nwlinkedidentc{Sort.t}{NW1rTmOJ-3ZqggC-1} -> order *)
fun compare(IND,IND) = EQUAL
  | compare(IND, _) = LESS
  | compare(FUN,IND) = GREATER
  | compare(FUN,FUN) = EQUAL
  | compare(FUN, _) = LESS
  | compare(CLASS, CLASS) = EQUAL
  | compare(CLASS, _) = GREATER;

(* to_string : \nwlinkedidentc{Sort.t}{NW1rTmOJ-3ZqggC-1} -> string *)
fun to_string (IND) = "IND"
  | to_string (FUN) = "FUN"
  | to_string (CLASS) = "CLASS";
end;
\nwindexdefn{\nwixident{Sort.t}}{Sort.t}{NW1rTmOJ-3ZqggC-1}\nwindexdefn{\nwixident{Sort.compare}}{Sort.compare}{NW1rTmOJ-3ZqggC-1}\nwindexdefn{\nwixident{Sort.to{\_}string}}{Sort.to:unstring}{NW1rTmOJ-3ZqggC-1}\nwindexdefn{\nwixident{Sort.IND}}{Sort.IND}{NW1rTmOJ-3ZqggC-1}\nwindexdefn{\nwixident{Sort.FUN}}{Sort.FUN}{NW1rTmOJ-3ZqggC-1}\nwindexdefn{\nwixident{Sort.CLASS}}{Sort.CLASS}{NW1rTmOJ-3ZqggC-1}\eatline
\nwnotused{Sort.sml}\nwidentdefs{\\{{\nwixident{Sort.CLASS}}{Sort.CLASS}}\\{{\nwixident{Sort.compare}}{Sort.compare}}\\{{\nwixident{Sort.FUN}}{Sort.FUN}}\\{{\nwixident{Sort.IND}}{Sort.IND}}\\{{\nwixident{Sort.t}}{Sort.t}}\\{{\nwixident{Sort.to{\_}string}}{Sort.to:unstring}}}\nwendcode{}\nwbegindocs{2}\nwdocspar
\begin{node}
Now, a three-sorted term is either a variable (which is a string and a
sort) or a function (which is a string, a sort, and a list of arguments).
We do not ``hardcode'' any of the primitive notions from $FS_{0}$ in
this representation. Instead, they will be ``initialized'' in the
state of the kernel.
\end{node}

\nwenddocs{}\nwbegincode{3}\sublabel{NW1rTmOJ-4LsRWC-1}\nwmargintag{{\nwtagstyle{}\subpageref{NW1rTmOJ-4LsRWC-1}}}\moddef{Term.sml~{\nwtagstyle{}\subpageref{NW1rTmOJ-4LsRWC-1}}}\endmoddef\nwstartdeflinemarkup\nwenddeflinemarkup
structure Term = struct
datatype t = Var of string * \nwlinkedidentc{Sort.t}{NW1rTmOJ-3ZqggC-1}
           | Fun of string * \nwlinkedidentc{Sort.t}{NW1rTmOJ-3ZqggC-1} * t list;

\LA{}Test if two terms have the same sort~{\nwtagstyle{}\subpageref{NW1rTmOJ-1WQNps-1}}\RA{}

\LA{}Predicates of terms about their sorts~{\nwtagstyle{}\subpageref{NW1rTmOJ-3FfXTO-1}}\RA{}

\LA{}Free variables in a term~{\nwtagstyle{}\subpageref{NW1rTmOJ-2nRtky-1}}\RA{}

\LA{}Substitution of a term for a variable in a term~{\nwtagstyle{}\subpageref{NW1rTmOJ-3OIyM5-1}}\RA{}

\LA{}Comparing two terms~{\nwtagstyle{}\subpageref{NW1rTmOJ-2yJYsC-1}}\RA{}

\LA{}Equality of two terms~{\nwtagstyle{}\subpageref{NW1rTmOJ-3810sY-1}}\RA{}

\LA{}Test if term contains subterm~{\nwtagstyle{}\subpageref{NW1rTmOJ-Re21r-1}}\RA{}

\LA{}Determine a fresh variable given some free variables~{\nwtagstyle{}\subpageref{NW1rTmOJ-4Ykq9M-1}}\RA{}

\LA{}Serialize a term~{\nwtagstyle{}\subpageref{NW1rTmOJ-40h2U3-1}}\RA{}
end;
\nwindexdefn{\nwixident{Term.t}}{Term.t}{NW1rTmOJ-4LsRWC-1}\eatline
\nwnotused{Term.sml}\nwidentdefs{\\{{\nwixident{Term.t}}{Term.t}}}\nwidentuses{\\{{\nwixident{Sort.t}}{Sort.t}}}\nwindexuse{\nwixident{Sort.t}}{Sort.t}{NW1rTmOJ-4LsRWC-1}\nwendcode{}\nwbegindocs{4}\nwdocspar
\begin{node}
We provide a minimal abstraction for terms. As it stands now, without
any ``helper'' or ``uniform'' constructors, it is about 44 lines of
code. If we ``hardcoded'' the $FS_{0}$ primitives, that would require
about 412 lines of code --- about $9.36$ times as much code.
\end{node}

\begin{node}
We provide a uniform way to access the sort of a term, and test if two
terms have the same sort.
\end{node}

\nwenddocs{}\nwbegincode{5}\sublabel{NW1rTmOJ-1WQNps-1}\nwmargintag{{\nwtagstyle{}\subpageref{NW1rTmOJ-1WQNps-1}}}\moddef{Test if two terms have the same sort~{\nwtagstyle{}\subpageref{NW1rTmOJ-1WQNps-1}}}\endmoddef\nwstartdeflinemarkup\nwusesondefline{\\{NW1rTmOJ-4LsRWC-1}}\nwenddeflinemarkup
(* sort : \nwlinkedidentc{Term.t}{NW1rTmOJ-4LsRWC-1} -> \nwlinkedidentc{Sort.t}{NW1rTmOJ-3ZqggC-1} *)
fun sort (Var (_, s)) = s
  | sort (Fun (_, s, _)) = s;

(* have_same_sort : \nwlinkedidentc{Term.t}{NW1rTmOJ-4LsRWC-1} -> \nwlinkedidentc{Term.t}{NW1rTmOJ-4LsRWC-1} -> bool *)
fun have_same_sort (lhs : t) (rhs : t) =
  (sort lhs) = (sort rhs);
\nwindexdefn{\nwixident{Term.sort}}{Term.sort}{NW1rTmOJ-1WQNps-1}\nwindexdefn{\nwixident{Term.have{\_}same{\_}sort}}{Term.have:unsame:unsort}{NW1rTmOJ-1WQNps-1}\eatline
\nwused{\\{NW1rTmOJ-4LsRWC-1}}\nwidentdefs{\\{{\nwixident{Term.have{\_}same{\_}sort}}{Term.have:unsame:unsort}}\\{{\nwixident{Term.sort}}{Term.sort}}}\nwidentuses{\\{{\nwixident{Sort.t}}{Sort.t}}\\{{\nwixident{Term.t}}{Term.t}}}\nwindexuse{\nwixident{Sort.t}}{Sort.t}{NW1rTmOJ-1WQNps-1}\nwindexuse{\nwixident{Term.t}}{Term.t}{NW1rTmOJ-1WQNps-1}\nwendcode{}\nwbegindocs{6}\nwdocspar
\begin{node}
We have some helper functions for common predicates.
\end{node}

\nwenddocs{}\nwbegincode{7}\sublabel{NW1rTmOJ-3FfXTO-1}\nwmargintag{{\nwtagstyle{}\subpageref{NW1rTmOJ-3FfXTO-1}}}\moddef{Predicates of terms about their sorts~{\nwtagstyle{}\subpageref{NW1rTmOJ-3FfXTO-1}}}\endmoddef\nwstartdeflinemarkup\nwusesondefline{\\{NW1rTmOJ-4LsRWC-1}}\nwenddeflinemarkup
(* is_var : \nwlinkedidentc{Term.t}{NW1rTmOJ-4LsRWC-1} -> bool *)
fun is_var (Var _) = true
  | is_var _ = false;

(* is_ind : \nwlinkedidentc{Term.t}{NW1rTmOJ-4LsRWC-1} -> bool *)
fun is_ind (tm : t) = (\nwlinkedidentc{Sort.IND}{NW1rTmOJ-3ZqggC-1} = sort tm);

(* is_fun : \nwlinkedidentc{Term.t}{NW1rTmOJ-4LsRWC-1} -> bool *)
fun is_fun (tm : t) = (\nwlinkedidentc{Sort.FUN}{NW1rTmOJ-3ZqggC-1} = sort tm);

(* is_class : \nwlinkedidentc{Term.t}{NW1rTmOJ-4LsRWC-1} -> bool *)
fun is_class (tm : t) = (\nwlinkedidentc{Sort.CLASS}{NW1rTmOJ-3ZqggC-1} = sort tm);
\nwindexdefn{\nwixident{Term.is{\_}ind}}{Term.is:unind}{NW1rTmOJ-3FfXTO-1}\nwindexdefn{\nwixident{Term.is{\_}fun}}{Term.is:unfun}{NW1rTmOJ-3FfXTO-1}\nwindexdefn{\nwixident{Term.is{\_}class}}{Term.is:unclass}{NW1rTmOJ-3FfXTO-1}\eatline
\nwused{\\{NW1rTmOJ-4LsRWC-1}}\nwidentdefs{\\{{\nwixident{Term.is{\_}class}}{Term.is:unclass}}\\{{\nwixident{Term.is{\_}fun}}{Term.is:unfun}}\\{{\nwixident{Term.is{\_}ind}}{Term.is:unind}}}\nwidentuses{\\{{\nwixident{Sort.CLASS}}{Sort.CLASS}}\\{{\nwixident{Sort.FUN}}{Sort.FUN}}\\{{\nwixident{Sort.IND}}{Sort.IND}}\\{{\nwixident{Term.t}}{Term.t}}}\nwindexuse{\nwixident{Sort.CLASS}}{Sort.CLASS}{NW1rTmOJ-3FfXTO-1}\nwindexuse{\nwixident{Sort.FUN}}{Sort.FUN}{NW1rTmOJ-3FfXTO-1}\nwindexuse{\nwixident{Sort.IND}}{Sort.IND}{NW1rTmOJ-3FfXTO-1}\nwindexuse{\nwixident{Term.t}}{Term.t}{NW1rTmOJ-3FfXTO-1}\nwendcode{}\nwbegindocs{8}\nwdocspar
\begin{node}
The free variables of a term are just \emph{all} the variables
appearing in the term. Later, formulas will remove the bound variables.
\end{node}

\nwenddocs{}\nwbegincode{9}\sublabel{NW1rTmOJ-2nRtky-1}\nwmargintag{{\nwtagstyle{}\subpageref{NW1rTmOJ-2nRtky-1}}}\moddef{Free variables in a term~{\nwtagstyle{}\subpageref{NW1rTmOJ-2nRtky-1}}}\endmoddef\nwstartdeflinemarkup\nwusesondefline{\\{NW1rTmOJ-4LsRWC-1}}\nwenddeflinemarkup
(* fv : \nwlinkedidentc{Term.t}{NW1rTmOJ-4LsRWC-1} -> \nwlinkedidentc{Term.t}{NW1rTmOJ-4LsRWC-1} list *)
fun fv (x as (Var _)) = [x]
  | fv (Fun (_, _, args)) = List.foldr (fn (arg, fvs) => (fv arg) @ fvs)
                                       []
                                       args;
\nwindexdefn{\nwixident{Term.fv}}{Term.fv}{NW1rTmOJ-2nRtky-1}\eatline
\nwused{\\{NW1rTmOJ-4LsRWC-1}}\nwidentdefs{\\{{\nwixident{Term.fv}}{Term.fv}}}\nwidentuses{\\{{\nwixident{Term.t}}{Term.t}}}\nwindexuse{\nwixident{Term.t}}{Term.t}{NW1rTmOJ-2nRtky-1}\nwendcode{}\nwbegindocs{10}\nwdocspar
\begin{node}
Substitution amounts to replacing a variable \emph{of a given sort}
with a term \emph{of the same sort} in a term recursively by:
\begin{enumerate}
\item $x[t/x]=t$ 
\item $y[t/x]=y$ if $y\neq x$
\item $(f(t_{1},\dots,t_{n}))[t/x]=(f(t_{1}[t/x],\dots,t_{n}[t/x]))$
\end{enumerate}
Although we \emph{should} restrict our focus to the situation where
$x$ is a variable, I suspect there might be situations where we want
to replace \emph{entire subterms}.
\end{node}

\nwenddocs{}\nwbegincode{11}\sublabel{NW1rTmOJ-3OIyM5-1}\nwmargintag{{\nwtagstyle{}\subpageref{NW1rTmOJ-3OIyM5-1}}}\moddef{Substitution of a term for a variable in a term~{\nwtagstyle{}\subpageref{NW1rTmOJ-3OIyM5-1}}}\endmoddef\nwstartdeflinemarkup\nwusesondefline{\\{NW1rTmOJ-4LsRWC-1}}\nwenddeflinemarkup
(* subst : \nwlinkedidentc{Term.t}{NW1rTmOJ-4LsRWC-1} -> \nwlinkedidentc{Term.t}{NW1rTmOJ-4LsRWC-1} -> \nwlinkedidentc{Term.t}{NW1rTmOJ-4LsRWC-1} -> \nwlinkedidentc{Term.t}{NW1rTmOJ-4LsRWC-1} *)
fun subst (x : t) (replacement : t) (tm : t) =
  if not (have_same_sort x replacement)
  then raise Sort.Mismatch "\nwlinkedidentc{Term.subst}{NW1rTmOJ-3OIyM5-1}: variable and replacement have different sorts"
  else if x = tm then replacement
  else (case tm of
         (Var _) => if x = tm then replacement else tm
       | (Fun(f,s,args)) => Fun(f,s,map (subst x replacement) args));
\nwindexdefn{\nwixident{Term.subst}}{Term.subst}{NW1rTmOJ-3OIyM5-1}\eatline
\nwused{\\{NW1rTmOJ-4LsRWC-1}}\nwidentdefs{\\{{\nwixident{Term.subst}}{Term.subst}}}\nwidentuses{\\{{\nwixident{Term.t}}{Term.t}}}\nwindexuse{\nwixident{Term.t}}{Term.t}{NW1rTmOJ-3OIyM5-1}\nwendcode{}\nwbegindocs{12}\nwdocspar
\begin{node}
Comparing two terms follows the following rules:
\begin{enumerate}
\item First compare the sort of both sides; if they differ, then this
  is the result
\item Variables are less than or equal to everything
\item Functions are greater than or equal to everything
\item Variables are compared by sort, then by name
\item Functions are compared by sort, then by name, then
  lexicographically on the arguments
\end{enumerate}
\end{node}


\nwenddocs{}\nwbegincode{13}\sublabel{NW1rTmOJ-2yJYsC-1}\nwmargintag{{\nwtagstyle{}\subpageref{NW1rTmOJ-2yJYsC-1}}}\moddef{Comparing two terms~{\nwtagstyle{}\subpageref{NW1rTmOJ-2yJYsC-1}}}\endmoddef\nwstartdeflinemarkup\nwusesondefline{\\{NW1rTmOJ-4LsRWC-1}}\nwenddeflinemarkup
(* compare : \nwlinkedidentc{Term.t}{NW1rTmOJ-4LsRWC-1} * \nwlinkedidentc{Term.t}{NW1rTmOJ-4LsRWC-1} -> order *)
fun compare(Var(x1,s1), Var(x2,s2)) =
    (case \nwlinkedidentc{Sort.compare}{NW1rTmOJ-3ZqggC-1}(s1, s2) of
      EQUAL => String.compare(x1, x2)
    | r => r)
  | compare(Var _, _) = LESS
  | compare(Fun _, Var _) = GREATER
  | compare(Fun(f1,s1,args1), Fun(f2,s2,args2)) =
    (case \nwlinkedidentc{Sort.compare}{NW1rTmOJ-3ZqggC-1}(s1,s2) of
      EQUAL => (case String.compare(f1, f2) of
                 EQUAL => compare_lists args1 args2
               | r => r)
    | r => r)
(* compare_lists : \nwlinkedidentc{Term.t}{NW1rTmOJ-4LsRWC-1} list -> \nwlinkedidentc{Term.t}{NW1rTmOJ-4LsRWC-1} list -> order *)
and compare_lists [] [] = EQUAL
  | compare_lists [] ys = LESS
  | compare_lists xs [] = GREATER
  | compare_lists (x::xs) (y::ys) =
     (case compare(x,y) of
       EQUAL => compare_lists xs ys
     | r => r);
\nwindexdefn{\nwixident{Term.compare}}{Term.compare}{NW1rTmOJ-2yJYsC-1}\nwindexdefn{\nwixident{Term.compare{\_}lists}}{Term.compare:unlists}{NW1rTmOJ-2yJYsC-1}\eatline
\nwused{\\{NW1rTmOJ-4LsRWC-1}}\nwidentdefs{\\{{\nwixident{Term.compare}}{Term.compare}}\\{{\nwixident{Term.compare{\_}lists}}{Term.compare:unlists}}}\nwidentuses{\\{{\nwixident{Sort.compare}}{Sort.compare}}\\{{\nwixident{Term.t}}{Term.t}}}\nwindexuse{\nwixident{Sort.compare}}{Sort.compare}{NW1rTmOJ-2yJYsC-1}\nwindexuse{\nwixident{Term.t}}{Term.t}{NW1rTmOJ-2yJYsC-1}\nwendcode{}\nwbegindocs{14}\nwdocspar
\begin{node}
The \emph{syntactic} equality of terms can be determined in a more
performant manner. Two variables are equal if their payload are equal
(i.e., they have the same sort and the same lexeme). Two functions are
equal if they have the same name, same sort, and their arguments are
syntactically equal (determined by recursively invoking {\Tt{}\nwlinkedidentq{Term.eq}{NW1rTmOJ-3810sY-1}\nwendquote}
pairwise on the arguments).
\end{node}

\nwenddocs{}\nwbegincode{15}\sublabel{NW1rTmOJ-3810sY-1}\nwmargintag{{\nwtagstyle{}\subpageref{NW1rTmOJ-3810sY-1}}}\moddef{Equality of two terms~{\nwtagstyle{}\subpageref{NW1rTmOJ-3810sY-1}}}\endmoddef\nwstartdeflinemarkup\nwusesondefline{\\{NW1rTmOJ-4LsRWC-1}}\nwenddeflinemarkup
(* eq : \nwlinkedidentc{Term.t}{NW1rTmOJ-4LsRWC-1} -> \nwlinkedidentc{Term.t}{NW1rTmOJ-4LsRWC-1} -> bool *)
fun eq (Var(x1,s1)) (Var(x2,s2)) = (s1 = s2) andalso (x1 = x2)
  | eq (Fun(f1,s1,args1)) (Fun(f2,s2,args2)) =
      (s1 = s2) andalso (f1 = f2) andalso (eq_lists args1 args2)
  | eq _ _ = false
(* eq_lists : \nwlinkedidentc{Term.t}{NW1rTmOJ-4LsRWC-1} list -> \nwlinkedidentc{Term.t}{NW1rTmOJ-4LsRWC-1} list -> bool *)
and eq_lists [] [] = true
  | eq_lists (x::xs) (y::ys) = (eq x y) andalso (eq_lists xs ys)
  | eq_lists _ _ = false;
\nwindexdefn{\nwixident{Term.eq}}{Term.eq}{NW1rTmOJ-3810sY-1}\nwindexdefn{\nwixident{Term.eq{\_}lists}}{Term.eq:unlists}{NW1rTmOJ-3810sY-1}\eatline
\nwused{\\{NW1rTmOJ-4LsRWC-1}}\nwidentdefs{\\{{\nwixident{Term.eq}}{Term.eq}}\\{{\nwixident{Term.eq{\_}lists}}{Term.eq:unlists}}}\nwidentuses{\\{{\nwixident{Term.t}}{Term.t}}}\nwindexuse{\nwixident{Term.t}}{Term.t}{NW1rTmOJ-3810sY-1}\nwendcode{}\nwbegindocs{16}\nwdocspar
\begin{node}
We also want to check if a term contains a subterm. This is not a
``strict'' comparison: any term occurs in itself.
Our intuition should be thinking of $x\leq y$ rather than $x<y$. 
\end{node}

\nwenddocs{}\nwbegincode{17}\sublabel{NW1rTmOJ-Re21r-1}\nwmargintag{{\nwtagstyle{}\subpageref{NW1rTmOJ-Re21r-1}}}\moddef{Test if term contains subterm~{\nwtagstyle{}\subpageref{NW1rTmOJ-Re21r-1}}}\endmoddef\nwstartdeflinemarkup\nwusesondefline{\\{NW1rTmOJ-4LsRWC-1}}\nwenddeflinemarkup
(* occurs_in : \nwlinkedidentc{Term.t}{NW1rTmOJ-4LsRWC-1} -> \nwlinkedidentc{Term.t}{NW1rTmOJ-4LsRWC-1} -> bool *)
fun occurs_in (subtm : t) (tm : t) : bool =
  (eq subtm tm) orelse
  (case tm of
    (Var _) => false
  | (Fun(_,_,args)) => List.exists (occurs_in subtm) args);
\nwindexdefn{\nwixident{Term.occurs{\_}in}}{Term.occurs:unin}{NW1rTmOJ-Re21r-1}\eatline
\nwused{\\{NW1rTmOJ-4LsRWC-1}}\nwidentdefs{\\{{\nwixident{Term.occurs{\_}in}}{Term.occurs:unin}}}\nwidentuses{\\{{\nwixident{Term.t}}{Term.t}}}\nwindexuse{\nwixident{Term.t}}{Term.t}{NW1rTmOJ-Re21r-1}\nwendcode{}\nwbegindocs{18}\nwdocspar
\begin{node}
We will want to find a fresh variable which does not appear in a list
of variables.
\end{node}

\nwenddocs{}\nwbegincode{19}\sublabel{NW1rTmOJ-4Ykq9M-1}\nwmargintag{{\nwtagstyle{}\subpageref{NW1rTmOJ-4Ykq9M-1}}}\moddef{Determine a fresh variable given some free variables~{\nwtagstyle{}\subpageref{NW1rTmOJ-4Ykq9M-1}}}\endmoddef\nwstartdeflinemarkup\nwusesondefline{\\{NW1rTmOJ-4LsRWC-1}}\nwenddeflinemarkup
local
  fun fresh_var_iter (x : string) (s : \nwlinkedidentc{Sort.t}{NW1rTmOJ-3ZqggC-1}) (vars : t list) n = 
    let val var = Var(x ^ (Int.toString n), s)
    in if List.exists (eq var) vars
       then fresh_var_iter x s vars (n + 1)
       else var
    end;
in

fun fresh_var (x : string) (s : \nwlinkedidentc{Sort.t}{NW1rTmOJ-3ZqggC-1}) (vars : t list) = 
  let val var = Var(x, s)
  in if List.exists (eq var) vars
     then fresh_var_iter x s vars 0
     else var
  end;

end;
\nwindexdefn{\nwixident{Term.fresh{\_}var}}{Term.fresh:unvar}{NW1rTmOJ-4Ykq9M-1}\eatline
\nwused{\\{NW1rTmOJ-4LsRWC-1}}\nwidentdefs{\\{{\nwixident{Term.fresh{\_}var}}{Term.fresh:unvar}}}\nwidentuses{\\{{\nwixident{Sort.t}}{Sort.t}}}\nwindexuse{\nwixident{Sort.t}}{Sort.t}{NW1rTmOJ-4Ykq9M-1}\nwendcode{}\nwbegindocs{20}\nwdocspar
\begin{node}
For debugging purposes, it will be useful to ``serialize'' a term (to
produce a string representation of the code needed to construct it).
\end{node}

\nwenddocs{}\nwbegincode{21}\sublabel{NW1rTmOJ-40h2U3-1}\nwmargintag{{\nwtagstyle{}\subpageref{NW1rTmOJ-40h2U3-1}}}\moddef{Serialize a term~{\nwtagstyle{}\subpageref{NW1rTmOJ-40h2U3-1}}}\endmoddef\nwstartdeflinemarkup\nwusesondefline{\\{NW1rTmOJ-4LsRWC-1}}\nwenddeflinemarkup
fun serialize (Var(x,s)) = "Var("^x^", "^(\nwlinkedidentc{Sort.to_string}{NW1rTmOJ-3ZqggC-1} s)^")"
  | serialize (Fun(f, s, args)) = "Fun("^f^
                                  ", "^
                                  (\nwlinkedidentc{Sort.to_string}{NW1rTmOJ-3ZqggC-1} s)^
                                  ", ["^
                                  (String.concatWith ", " (map serialize args))^
                                  "])";
\nwindexdefn{\nwixident{Term.serialize}}{Term.serialize}{NW1rTmOJ-40h2U3-1}\eatline

\nwused{\\{NW1rTmOJ-4LsRWC-1}}\nwidentdefs{\\{{\nwixident{Term.serialize}}{Term.serialize}}}\nwidentuses{\\{{\nwixident{Sort.to{\_}string}}{Sort.to:unstring}}}\nwindexuse{\nwixident{Sort.to{\_}string}}{Sort.to:unstring}{NW1rTmOJ-40h2U3-1}\nwendcode{}\nwfilename{nw/many-sorted-formula.nw}\nwbegindocs{0}% -*- mode: poly-noweb; noweb-code-mode: sml-mode; -*-
%%
%% many-sorted-formula.nw
%% 
%% Made by Alex Nelson <pqnelson@gmail.com>
%% Login   <alex@lisp>
%% 
%% Started on  2026-04-04T18:49:02-0700
%% Last update 2026-04-04T18:49:02-0700
%%

\section{Syntax for formulas}

\begin{node}
We need to use opaque types for formulas, because we want to use smart
constructors to ensure quantifiers are given variable {\Tt{}\nwlinkedidentq{Term.t}{NW1rTmOJ-4LsRWC-1}\nwendquote} instances.
Otherwise it'd be too easy to accidentally change the sort of a bound
variable, which will produce bugs later on (bugs that will be
difficult to trace).

Consequently, we will want to provide \define{Destructuring} functions
to ``break apart'' the abstract syntax tree for a formula, so we can
recover its constituent parts. We have one destructuring function for
each data constructor. If the user mistakenly applies an incorrect
destructuring formula to a different formula (e.g., {\Tt{}dest{\_}forall\nwendquote}
applied to an implication), then an exception should be raised. This
is precisely the {\Tt{}\nwlinkedidentq{Formula.Dest}{NW3w0iRO-2jnj9W-1}\nwendquote} exception's role.
\end{node}

\nwenddocs{}\nwbegincode{1}\sublabel{NW3w0iRO-5BH7S-1}\nwmargintag{{\nwtagstyle{}\subpageref{NW3w0iRO-5BH7S-1}}}\moddef{Formula.sig~{\nwtagstyle{}\subpageref{NW3w0iRO-5BH7S-1}}}\endmoddef\nwstartdeflinemarkup\nwenddeflinemarkup
signature Formula = sig
  exception Dest of string;
  type t;
  val is_atom : t -> bool;
  \LA{}Specification for formulas~{\nwtagstyle{}\subpageref{NW3w0iRO-2HPkc8-1}}\RA{}
end;
\nwnotused{Formula.sig}\nwendcode{}\nwbegindocs{2}\nwdocspar

\nwenddocs{}\nwbegincode{3}\sublabel{NW3w0iRO-2jnj9W-1}\nwmargintag{{\nwtagstyle{}\subpageref{NW3w0iRO-2jnj9W-1}}}\moddef{Formula.sml~{\nwtagstyle{}\subpageref{NW3w0iRO-2jnj9W-1}}}\endmoddef\nwstartdeflinemarkup\nwenddeflinemarkup
structure Formula :> Formula = struct
exception Dest of string;

datatype t = \nwlinkedidentc{False}{NW3w0iRO-2jnj9W-1}
           | \nwlinkedidentc{Pred}{NW3w0iRO-2jnj9W-1} of string * \nwlinkedidentc{Term.t}{NW1rTmOJ-4LsRWC-1} list
           | \nwlinkedidentc{Not}{NW3w0iRO-2jnj9W-1} of t
           | \nwlinkedidentc{And}{NW3w0iRO-2jnj9W-1} of t * t
           | \nwlinkedidentc{Or}{NW3w0iRO-2jnj9W-1} of t * t
           | \nwlinkedidentc{Implies}{NW3w0iRO-2jnj9W-1} of t * t
           | \nwlinkedidentc{Iff}{NW3w0iRO-2jnj9W-1} of t * t
           | \nwlinkedidentc{Forall}{NW3w0iRO-2jnj9W-1} of \nwlinkedidentc{Term.t}{NW1rTmOJ-4LsRWC-1} * t
           | \nwlinkedidentc{Exists}{NW3w0iRO-2jnj9W-1} of \nwlinkedidentc{Term.t}{NW1rTmOJ-4LsRWC-1} * t;

fun is_atom (\nwlinkedidentc{False}{NW3w0iRO-2jnj9W-1}) = true
  | is_atom (\nwlinkedidentc{Pred}{NW3w0iRO-2jnj9W-1} _) = true
  | is_atom _ = false;

\LA{}Constructors for formulas~{\nwtagstyle{}\subpageref{NW3w0iRO-3SlqND-1}}\RA{}

\LA{}Destructors for formulas~{\nwtagstyle{}\subpageref{NW3w0iRO-1Yt1Jl-1}}\RA{}

\LA{}Predicates for formulas~{\nwtagstyle{}\subpageref{NW3w0iRO-3NyJLv-1}}\RA{}

\LA{}Substitution in a formula~{\nwtagstyle{}\subpageref{NW3w0iRO-1jtQRV-1}}\RA{}

\LA{}Folding formulas~{\nwtagstyle{}\subpageref{NW3w0iRO-1UOPDY-1}}\RA{}

\LA{}Free and bound variables in a formula~{\nwtagstyle{}\subpageref{NW3w0iRO-jRgYZ-1}}\RA{}

\LA{}Check if a formula contains a variable~{\nwtagstyle{}\subpageref{NW3w0iRO-4DOLFK-1}}\RA{}

\LA{}Precedence for connectives~{\nwtagstyle{}\subpageref{NW3w0iRO-3yI6te-1}}\RA{}

\LA{}Comparing formulas~{\nwtagstyle{}\subpageref{NW3w0iRO-35s47m-1}}\RA{}

\LA{}Test if two formulas are equal~{\nwtagstyle{}\subpageref{NW3w0iRO-5gQzE-1}}\RA{}

\LA{}Test a term is free for a variable in a formula~{\nwtagstyle{}\subpageref{NW3w0iRO-ZZefY-1}}\RA{}

\LA{}Test for $E^{+}$-formula~{\nwtagstyle{}\subpageref{NW3w0iRO-1O1TnJ-1}}\RA{}

\LA{}Serialize a formula~{\nwtagstyle{}\subpageref{NW3w0iRO-3H9tUA-1}}\RA{}

\LA{}Take the union of applying a function to atoms~{\nwtagstyle{}\subpageref{NW3w0iRO-28e6MC-1}}\RA{}
end;
\nwindexdefn{\nwixident{Formula.Dest}}{Formula.Dest}{NW3w0iRO-2jnj9W-1}\nwindexdefn{\nwixident{Formula.t}}{Formula.t}{NW3w0iRO-2jnj9W-1}\nwindexdefn{\nwixident{False}}{False}{NW3w0iRO-2jnj9W-1}\nwindexdefn{\nwixident{Pred}}{Pred}{NW3w0iRO-2jnj9W-1}\nwindexdefn{\nwixident{Not}}{Not}{NW3w0iRO-2jnj9W-1}\nwindexdefn{\nwixident{And}}{And}{NW3w0iRO-2jnj9W-1}\nwindexdefn{\nwixident{Or}}{Or}{NW3w0iRO-2jnj9W-1}\nwindexdefn{\nwixident{Implies}}{Implies}{NW3w0iRO-2jnj9W-1}\nwindexdefn{\nwixident{Iff}}{Iff}{NW3w0iRO-2jnj9W-1}\nwindexdefn{\nwixident{Forall}}{Forall}{NW3w0iRO-2jnj9W-1}\nwindexdefn{\nwixident{Exists}}{Exists}{NW3w0iRO-2jnj9W-1}\eatline
\nwnotused{Formula.sml}\nwidentdefs{\\{{\nwixident{And}}{And}}\\{{\nwixident{Exists}}{Exists}}\\{{\nwixident{False}}{False}}\\{{\nwixident{Forall}}{Forall}}\\{{\nwixident{Formula.Dest}}{Formula.Dest}}\\{{\nwixident{Formula.t}}{Formula.t}}\\{{\nwixident{Iff}}{Iff}}\\{{\nwixident{Implies}}{Implies}}\\{{\nwixident{Not}}{Not}}\\{{\nwixident{Or}}{Or}}\\{{\nwixident{Pred}}{Pred}}}\nwidentuses{\\{{\nwixident{Term.t}}{Term.t}}}\nwindexuse{\nwixident{Term.t}}{Term.t}{NW3w0iRO-2jnj9W-1}\nwendcode{}\nwbegindocs{4}\nwdocspar
\begin{node}
The canonical contradiction $\bot$ can be constructed, and we can test
for it. But it does not make sense to ``destructure'' it, so we omit
its destructor.

The canonical ``tautology'' \textit{verum} $\top$ is an abbreviation
for $\neg\bot$. Like the canonical contradiction, there's no
destructuring function for it (it is ``atomic'').
\end{node}

\nwenddocs{}\nwbegincode{5}\sublabel{NW3w0iRO-2HPkc8-1}\nwmargintag{{\nwtagstyle{}\subpageref{NW3w0iRO-2HPkc8-1}}}\moddef{Specification for formulas~{\nwtagstyle{}\subpageref{NW3w0iRO-2HPkc8-1}}}\endmoddef\nwstartdeflinemarkup\nwusesondefline{\\{NW3w0iRO-5BH7S-1}}\nwprevnextdefs{\relax}{NW3w0iRO-2HPkc8-2}\nwenddeflinemarkup
val mk_verum : unit -> t;
val mk_falsum : unit -> t;
val mk_pred : string * \nwlinkedidentc{Term.t}{NW1rTmOJ-4LsRWC-1} list -> t;
\nwalsodefined{\\{NW3w0iRO-2HPkc8-2}\\{NW3w0iRO-2HPkc8-3}\\{NW3w0iRO-2HPkc8-4}\\{NW3w0iRO-2HPkc8-5}\\{NW3w0iRO-2HPkc8-6}\\{NW3w0iRO-2HPkc8-7}\\{NW3w0iRO-2HPkc8-8}\\{NW3w0iRO-2HPkc8-9}\\{NW3w0iRO-2HPkc8-A}\\{NW3w0iRO-2HPkc8-B}\\{NW3w0iRO-2HPkc8-C}\\{NW3w0iRO-2HPkc8-D}\\{NW3w0iRO-2HPkc8-E}\\{NW3w0iRO-2HPkc8-F}\\{NW3w0iRO-2HPkc8-G}\\{NW3w0iRO-2HPkc8-H}}\nwused{\\{NW3w0iRO-5BH7S-1}}\nwidentuses{\\{{\nwixident{Term.t}}{Term.t}}}\nwindexuse{\nwixident{Term.t}}{Term.t}{NW3w0iRO-2HPkc8-1}\nwendcode{}\nwbegindocs{6}\nwdocspar
\nwenddocs{}\nwbegincode{7}\sublabel{NW3w0iRO-3SlqND-1}\nwmargintag{{\nwtagstyle{}\subpageref{NW3w0iRO-3SlqND-1}}}\moddef{Constructors for formulas~{\nwtagstyle{}\subpageref{NW3w0iRO-3SlqND-1}}}\endmoddef\nwstartdeflinemarkup\nwusesondefline{\\{NW3w0iRO-2jnj9W-1}}\nwprevnextdefs{\relax}{NW3w0iRO-3SlqND-2}\nwenddeflinemarkup
fun mk_falsum () = \nwlinkedidentc{False}{NW3w0iRO-2jnj9W-1};

fun mk_verum () = \nwlinkedidentc{Not}{NW3w0iRO-2jnj9W-1}(\nwlinkedidentc{False}{NW3w0iRO-2jnj9W-1});

fun mk_pred args = \nwlinkedidentc{Pred}{NW3w0iRO-2jnj9W-1} args;
\nwindexdefn{\nwixident{Formula.mk{\_}falsum}}{Formula.mk:unfalsum}{NW3w0iRO-3SlqND-1}\nwindexdefn{\nwixident{Formula.mk{\_}pred}}{Formula.mk:unpred}{NW3w0iRO-3SlqND-1}\nwindexdefn{\nwixident{Formula.mk{\_}verum}}{Formula.mk:unverum}{NW3w0iRO-3SlqND-1}\eatline
\nwalsodefined{\\{NW3w0iRO-3SlqND-2}\\{NW3w0iRO-3SlqND-3}\\{NW3w0iRO-3SlqND-4}}\nwused{\\{NW3w0iRO-2jnj9W-1}}\nwidentdefs{\\{{\nwixident{Formula.mk{\_}falsum}}{Formula.mk:unfalsum}}\\{{\nwixident{Formula.mk{\_}pred}}{Formula.mk:unpred}}\\{{\nwixident{Formula.mk{\_}verum}}{Formula.mk:unverum}}}\nwidentuses{\\{{\nwixident{False}}{False}}\\{{\nwixident{Not}}{Not}}\\{{\nwixident{Pred}}{Pred}}}\nwindexuse{\nwixident{False}}{False}{NW3w0iRO-3SlqND-1}\nwindexuse{\nwixident{Not}}{Not}{NW3w0iRO-3SlqND-1}\nwindexuse{\nwixident{Pred}}{Pred}{NW3w0iRO-3SlqND-1}\nwendcode{}\nwbegindocs{8}\nwdocspar

\nwenddocs{}\nwbegincode{9}\sublabel{NW3w0iRO-2HPkc8-2}\nwmargintag{{\nwtagstyle{}\subpageref{NW3w0iRO-2HPkc8-2}}}\moddef{Specification for formulas~{\nwtagstyle{}\subpageref{NW3w0iRO-2HPkc8-1}}}\plusendmoddef\nwstartdeflinemarkup\nwusesondefline{\\{NW3w0iRO-5BH7S-1}}\nwprevnextdefs{NW3w0iRO-2HPkc8-1}{NW3w0iRO-2HPkc8-3}\nwenddeflinemarkup
val dest_pred : t -> string * \nwlinkedidentc{Term.t}{NW1rTmOJ-4LsRWC-1} list;
\nwused{\\{NW3w0iRO-5BH7S-1}}\nwidentuses{\\{{\nwixident{Term.t}}{Term.t}}}\nwindexuse{\nwixident{Term.t}}{Term.t}{NW3w0iRO-2HPkc8-2}\nwendcode{}\nwbegindocs{10}\nwdocspar
\nwenddocs{}\nwbegincode{11}\sublabel{NW3w0iRO-1Yt1Jl-1}\nwmargintag{{\nwtagstyle{}\subpageref{NW3w0iRO-1Yt1Jl-1}}}\moddef{Destructors for formulas~{\nwtagstyle{}\subpageref{NW3w0iRO-1Yt1Jl-1}}}\endmoddef\nwstartdeflinemarkup\nwusesondefline{\\{NW3w0iRO-2jnj9W-1}}\nwprevnextdefs{\relax}{NW3w0iRO-1Yt1Jl-2}\nwenddeflinemarkup
fun dest_pred (\nwlinkedidentc{Pred}{NW3w0iRO-2jnj9W-1} args) = args
  | dest_pred _ = raise Dest "dest_pred";
\nwindexdefn{\nwixident{Formula.dest{\_}pred}}{Formula.dest:unpred}{NW3w0iRO-1Yt1Jl-1}\eatline
\nwalsodefined{\\{NW3w0iRO-1Yt1Jl-2}\\{NW3w0iRO-1Yt1Jl-3}\\{NW3w0iRO-1Yt1Jl-4}}\nwused{\\{NW3w0iRO-2jnj9W-1}}\nwidentdefs{\\{{\nwixident{Formula.dest{\_}pred}}{Formula.dest:unpred}}}\nwidentuses{\\{{\nwixident{Pred}}{Pred}}}\nwindexuse{\nwixident{Pred}}{Pred}{NW3w0iRO-1Yt1Jl-1}\nwendcode{}\nwbegindocs{12}\nwdocspar
\nwenddocs{}\nwbegincode{13}\sublabel{NW3w0iRO-2HPkc8-3}\nwmargintag{{\nwtagstyle{}\subpageref{NW3w0iRO-2HPkc8-3}}}\moddef{Specification for formulas~{\nwtagstyle{}\subpageref{NW3w0iRO-2HPkc8-1}}}\plusendmoddef\nwstartdeflinemarkup\nwusesondefline{\\{NW3w0iRO-5BH7S-1}}\nwprevnextdefs{NW3w0iRO-2HPkc8-2}{NW3w0iRO-2HPkc8-4}\nwenddeflinemarkup
val is_verum : t -> bool;
val is_falsum : t -> bool;
val is_pred : t -> bool;
\nwused{\\{NW3w0iRO-5BH7S-1}}\nwendcode{}\nwbegindocs{14}\nwdocspar
\nwenddocs{}\nwbegincode{15}\sublabel{NW3w0iRO-3NyJLv-1}\nwmargintag{{\nwtagstyle{}\subpageref{NW3w0iRO-3NyJLv-1}}}\moddef{Predicates for formulas~{\nwtagstyle{}\subpageref{NW3w0iRO-3NyJLv-1}}}\endmoddef\nwstartdeflinemarkup\nwusesondefline{\\{NW3w0iRO-2jnj9W-1}}\nwprevnextdefs{\relax}{NW3w0iRO-3NyJLv-2}\nwenddeflinemarkup
fun is_falsum (\nwlinkedidentc{False}{NW3w0iRO-2jnj9W-1}) = true
  | is_falsum _ = false;

fun is_verum (\nwlinkedidentc{Not}{NW3w0iRO-2jnj9W-1}(\nwlinkedidentc{False}{NW3w0iRO-2jnj9W-1})) = true
  | is_verum _ = false;

fun is_pred (\nwlinkedidentc{Pred}{NW3w0iRO-2jnj9W-1} _) = true
  | is_pred _ = false;
\nwindexdefn{\nwixident{Formula.is{\_}falsum}}{Formula.is:unfalsum}{NW3w0iRO-3NyJLv-1}\nwindexdefn{\nwixident{Formula.is{\_}verum}}{Formula.is:unverum}{NW3w0iRO-3NyJLv-1}\nwindexdefn{\nwixident{Formula.is{\_}pred}}{Formula.is:unpred}{NW3w0iRO-3NyJLv-1}\eatline
\nwalsodefined{\\{NW3w0iRO-3NyJLv-2}\\{NW3w0iRO-3NyJLv-3}\\{NW3w0iRO-3NyJLv-4}}\nwused{\\{NW3w0iRO-2jnj9W-1}}\nwidentdefs{\\{{\nwixident{Formula.is{\_}falsum}}{Formula.is:unfalsum}}\\{{\nwixident{Formula.is{\_}pred}}{Formula.is:unpred}}\\{{\nwixident{Formula.is{\_}verum}}{Formula.is:unverum}}}\nwidentuses{\\{{\nwixident{False}}{False}}\\{{\nwixident{Not}}{Not}}\\{{\nwixident{Pred}}{Pred}}}\nwindexuse{\nwixident{False}}{False}{NW3w0iRO-3NyJLv-1}\nwindexuse{\nwixident{Not}}{Not}{NW3w0iRO-3NyJLv-1}\nwindexuse{\nwixident{Pred}}{Pred}{NW3w0iRO-3NyJLv-1}\nwendcode{}\nwbegindocs{16}\nwdocspar
\begin{node}
Negation is the only unary logical connective. Intuitionistic logic
can abbreviate negation as ``implies $\bot$''. The advantage to such
an abbreviation is that negation-introduction and elimination
inference rules do not need to be present in the kernel. We can
implement them as derived inference rules, i.e., as functions outside
the kernel using kernel primitives.
\end{node}

\nwenddocs{}\nwbegincode{17}\sublabel{NW3w0iRO-2HPkc8-4}\nwmargintag{{\nwtagstyle{}\subpageref{NW3w0iRO-2HPkc8-4}}}\moddef{Specification for formulas~{\nwtagstyle{}\subpageref{NW3w0iRO-2HPkc8-1}}}\plusendmoddef\nwstartdeflinemarkup\nwusesondefline{\\{NW3w0iRO-5BH7S-1}}\nwprevnextdefs{NW3w0iRO-2HPkc8-3}{NW3w0iRO-2HPkc8-5}\nwenddeflinemarkup
val mk_not : t -> t;
val dest_not : t -> t;
val is_not : t -> bool;
\nwused{\\{NW3w0iRO-5BH7S-1}}\nwendcode{}\nwbegindocs{18}\nwdocspar

\nwenddocs{}\nwbegincode{19}\sublabel{NW3w0iRO-3SlqND-2}\nwmargintag{{\nwtagstyle{}\subpageref{NW3w0iRO-3SlqND-2}}}\moddef{Constructors for formulas~{\nwtagstyle{}\subpageref{NW3w0iRO-3SlqND-1}}}\plusendmoddef\nwstartdeflinemarkup\nwusesondefline{\\{NW3w0iRO-2jnj9W-1}}\nwprevnextdefs{NW3w0iRO-3SlqND-1}{NW3w0iRO-3SlqND-3}\nwenddeflinemarkup
fun mk_not fm = \nwlinkedidentc{Not}{NW3w0iRO-2jnj9W-1} fm;
\nwindexdefn{\nwixident{Formula.mk{\_}not}}{Formula.mk:unnot}{NW3w0iRO-3SlqND-2}\eatline
\nwused{\\{NW3w0iRO-2jnj9W-1}}\nwidentdefs{\\{{\nwixident{Formula.mk{\_}not}}{Formula.mk:unnot}}}\nwidentuses{\\{{\nwixident{Not}}{Not}}}\nwindexuse{\nwixident{Not}}{Not}{NW3w0iRO-3SlqND-2}\nwendcode{}\nwbegindocs{20}\nwdocspar
\nwenddocs{}\nwbegincode{21}\sublabel{NW3w0iRO-1Yt1Jl-2}\nwmargintag{{\nwtagstyle{}\subpageref{NW3w0iRO-1Yt1Jl-2}}}\moddef{Destructors for formulas~{\nwtagstyle{}\subpageref{NW3w0iRO-1Yt1Jl-1}}}\plusendmoddef\nwstartdeflinemarkup\nwusesondefline{\\{NW3w0iRO-2jnj9W-1}}\nwprevnextdefs{NW3w0iRO-1Yt1Jl-1}{NW3w0iRO-1Yt1Jl-3}\nwenddeflinemarkup
fun dest_not (\nwlinkedidentc{Not}{NW3w0iRO-2jnj9W-1} fm) = fm
  | dest_not _ = raise Dest "dest_not";
\nwindexdefn{\nwixident{Formula.dest{\_}not}}{Formula.dest:unnot}{NW3w0iRO-1Yt1Jl-2}\eatline
\nwused{\\{NW3w0iRO-2jnj9W-1}}\nwidentdefs{\\{{\nwixident{Formula.dest{\_}not}}{Formula.dest:unnot}}}\nwidentuses{\\{{\nwixident{Not}}{Not}}}\nwindexuse{\nwixident{Not}}{Not}{NW3w0iRO-1Yt1Jl-2}\nwendcode{}\nwbegindocs{22}\nwdocspar
\nwenddocs{}\nwbegincode{23}\sublabel{NW3w0iRO-3NyJLv-2}\nwmargintag{{\nwtagstyle{}\subpageref{NW3w0iRO-3NyJLv-2}}}\moddef{Predicates for formulas~{\nwtagstyle{}\subpageref{NW3w0iRO-3NyJLv-1}}}\plusendmoddef\nwstartdeflinemarkup\nwusesondefline{\\{NW3w0iRO-2jnj9W-1}}\nwprevnextdefs{NW3w0iRO-3NyJLv-1}{NW3w0iRO-3NyJLv-3}\nwenddeflinemarkup
fun is_not (\nwlinkedidentc{Not}{NW3w0iRO-2jnj9W-1} _) = true
  | is_not _ = false;
\nwindexdefn{\nwixident{Formula.is{\_}not}}{Formula.is:unnot}{NW3w0iRO-3NyJLv-2}\eatline
\nwused{\\{NW3w0iRO-2jnj9W-1}}\nwidentdefs{\\{{\nwixident{Formula.is{\_}not}}{Formula.is:unnot}}}\nwidentuses{\\{{\nwixident{Not}}{Not}}}\nwindexuse{\nwixident{Not}}{Not}{NW3w0iRO-3NyJLv-2}\nwendcode{}\nwbegindocs{24}\nwdocspar
\begin{node}
The binary connectives all have similar constructors, destructors, and
predicates. A more extensible way to implement formulas would be to
have a single data constructor for binary connectives with a string
holding the connective. This is the approach
Paulson~\cite{paulson1990designing} takes.
\end{node}

\nwenddocs{}\nwbegincode{25}\sublabel{NW3w0iRO-2HPkc8-5}\nwmargintag{{\nwtagstyle{}\subpageref{NW3w0iRO-2HPkc8-5}}}\moddef{Specification for formulas~{\nwtagstyle{}\subpageref{NW3w0iRO-2HPkc8-1}}}\plusendmoddef\nwstartdeflinemarkup\nwusesondefline{\\{NW3w0iRO-5BH7S-1}}\nwprevnextdefs{NW3w0iRO-2HPkc8-4}{NW3w0iRO-2HPkc8-6}\nwenddeflinemarkup
val mk_and : t * t -> t;
val mk_or : t * t -> t;
val mk_imp : t * t -> t;
val mk_iff : t * t -> t;

val dest_and : t -> t * t;
val dest_or : t -> t * t;
val dest_imp : t -> t * t;
val dest_iff : t -> t * t;

val is_and : t -> bool;
val is_or : t -> bool;
val is_imp : t -> bool;
val is_iff : t -> bool;
\nwused{\\{NW3w0iRO-5BH7S-1}}\nwendcode{}\nwbegindocs{26}\nwdocspar

\nwenddocs{}\nwbegincode{27}\sublabel{NW3w0iRO-3SlqND-3}\nwmargintag{{\nwtagstyle{}\subpageref{NW3w0iRO-3SlqND-3}}}\moddef{Constructors for formulas~{\nwtagstyle{}\subpageref{NW3w0iRO-3SlqND-1}}}\plusendmoddef\nwstartdeflinemarkup\nwusesondefline{\\{NW3w0iRO-2jnj9W-1}}\nwprevnextdefs{NW3w0iRO-3SlqND-2}{NW3w0iRO-3SlqND-4}\nwenddeflinemarkup
fun mk_and args = \nwlinkedidentc{And}{NW3w0iRO-2jnj9W-1} args;

fun mk_or args = \nwlinkedidentc{Or}{NW3w0iRO-2jnj9W-1} args;

fun mk_imp args = \nwlinkedidentc{Implies}{NW3w0iRO-2jnj9W-1} args;

fun mk_iff args = \nwlinkedidentc{Iff}{NW3w0iRO-2jnj9W-1} args;
\nwindexdefn{\nwixident{Formula.mk{\_}and}}{Formula.mk:unand}{NW3w0iRO-3SlqND-3}\nwindexdefn{\nwixident{Formula.mk{\_}or}}{Formula.mk:unor}{NW3w0iRO-3SlqND-3}\nwindexdefn{\nwixident{Formula.mk{\_}imp}}{Formula.mk:unimp}{NW3w0iRO-3SlqND-3}\nwindexdefn{\nwixident{Formula.mk{\_}iff}}{Formula.mk:uniff}{NW3w0iRO-3SlqND-3}\eatline
\nwused{\\{NW3w0iRO-2jnj9W-1}}\nwidentdefs{\\{{\nwixident{Formula.mk{\_}and}}{Formula.mk:unand}}\\{{\nwixident{Formula.mk{\_}iff}}{Formula.mk:uniff}}\\{{\nwixident{Formula.mk{\_}imp}}{Formula.mk:unimp}}\\{{\nwixident{Formula.mk{\_}or}}{Formula.mk:unor}}}\nwidentuses{\\{{\nwixident{And}}{And}}\\{{\nwixident{Iff}}{Iff}}\\{{\nwixident{Implies}}{Implies}}\\{{\nwixident{Or}}{Or}}}\nwindexuse{\nwixident{And}}{And}{NW3w0iRO-3SlqND-3}\nwindexuse{\nwixident{Iff}}{Iff}{NW3w0iRO-3SlqND-3}\nwindexuse{\nwixident{Implies}}{Implies}{NW3w0iRO-3SlqND-3}\nwindexuse{\nwixident{Or}}{Or}{NW3w0iRO-3SlqND-3}\nwendcode{}\nwbegindocs{28}\nwdocspar
\nwenddocs{}\nwbegincode{29}\sublabel{NW3w0iRO-1Yt1Jl-3}\nwmargintag{{\nwtagstyle{}\subpageref{NW3w0iRO-1Yt1Jl-3}}}\moddef{Destructors for formulas~{\nwtagstyle{}\subpageref{NW3w0iRO-1Yt1Jl-1}}}\plusendmoddef\nwstartdeflinemarkup\nwusesondefline{\\{NW3w0iRO-2jnj9W-1}}\nwprevnextdefs{NW3w0iRO-1Yt1Jl-2}{NW3w0iRO-1Yt1Jl-4}\nwenddeflinemarkup
fun dest_and (\nwlinkedidentc{And}{NW3w0iRO-2jnj9W-1} args) = args
  | dest_and _ = raise Dest "dest_and";

fun dest_or (\nwlinkedidentc{Or}{NW3w0iRO-2jnj9W-1} args) = args
  | dest_or _ = raise Dest "dest_or";

fun dest_imp (\nwlinkedidentc{Implies}{NW3w0iRO-2jnj9W-1} args) = args
  | dest_imp _ = raise Dest "dest_imp";

fun dest_iff (\nwlinkedidentc{Iff}{NW3w0iRO-2jnj9W-1} args) = args
  | dest_iff _ = raise Dest "dest_iff";
\nwindexdefn{\nwixident{Formula.dest{\_}and}}{Formula.dest:unand}{NW3w0iRO-1Yt1Jl-3}\nwindexdefn{\nwixident{Formula.dest{\_}or}}{Formula.dest:unor}{NW3w0iRO-1Yt1Jl-3}\nwindexdefn{\nwixident{Formula.dest{\_}imp}}{Formula.dest:unimp}{NW3w0iRO-1Yt1Jl-3}\nwindexdefn{\nwixident{Formula.dest{\_}iff}}{Formula.dest:uniff}{NW3w0iRO-1Yt1Jl-3}\eatline
\nwused{\\{NW3w0iRO-2jnj9W-1}}\nwidentdefs{\\{{\nwixident{Formula.dest{\_}and}}{Formula.dest:unand}}\\{{\nwixident{Formula.dest{\_}iff}}{Formula.dest:uniff}}\\{{\nwixident{Formula.dest{\_}imp}}{Formula.dest:unimp}}\\{{\nwixident{Formula.dest{\_}or}}{Formula.dest:unor}}}\nwidentuses{\\{{\nwixident{And}}{And}}\\{{\nwixident{Iff}}{Iff}}\\{{\nwixident{Implies}}{Implies}}\\{{\nwixident{Or}}{Or}}}\nwindexuse{\nwixident{And}}{And}{NW3w0iRO-1Yt1Jl-3}\nwindexuse{\nwixident{Iff}}{Iff}{NW3w0iRO-1Yt1Jl-3}\nwindexuse{\nwixident{Implies}}{Implies}{NW3w0iRO-1Yt1Jl-3}\nwindexuse{\nwixident{Or}}{Or}{NW3w0iRO-1Yt1Jl-3}\nwendcode{}\nwbegindocs{30}\nwdocspar
\nwenddocs{}\nwbegincode{31}\sublabel{NW3w0iRO-3NyJLv-3}\nwmargintag{{\nwtagstyle{}\subpageref{NW3w0iRO-3NyJLv-3}}}\moddef{Predicates for formulas~{\nwtagstyle{}\subpageref{NW3w0iRO-3NyJLv-1}}}\plusendmoddef\nwstartdeflinemarkup\nwusesondefline{\\{NW3w0iRO-2jnj9W-1}}\nwprevnextdefs{NW3w0iRO-3NyJLv-2}{NW3w0iRO-3NyJLv-4}\nwenddeflinemarkup
fun is_and (\nwlinkedidentc{And}{NW3w0iRO-2jnj9W-1} _) = true
  | is_and _ = false;

fun is_or (\nwlinkedidentc{Or}{NW3w0iRO-2jnj9W-1} _) = true
  | is_or _ = false;

fun is_imp (\nwlinkedidentc{Implies}{NW3w0iRO-2jnj9W-1} _) = true
  | is_imp _ = false;

fun is_iff (\nwlinkedidentc{Iff}{NW3w0iRO-2jnj9W-1} _) = true
  | is_iff _ = false;
\nwindexdefn{\nwixident{Formula.is{\_}and}}{Formula.is:unand}{NW3w0iRO-3NyJLv-3}\nwindexdefn{\nwixident{Formula.is{\_}or}}{Formula.is:unor}{NW3w0iRO-3NyJLv-3}\nwindexdefn{\nwixident{Formula.is{\_}imp}}{Formula.is:unimp}{NW3w0iRO-3NyJLv-3}\nwindexdefn{\nwixident{Formula.is{\_}iff}}{Formula.is:uniff}{NW3w0iRO-3NyJLv-3}\eatline
\nwused{\\{NW3w0iRO-2jnj9W-1}}\nwidentdefs{\\{{\nwixident{Formula.is{\_}and}}{Formula.is:unand}}\\{{\nwixident{Formula.is{\_}iff}}{Formula.is:uniff}}\\{{\nwixident{Formula.is{\_}imp}}{Formula.is:unimp}}\\{{\nwixident{Formula.is{\_}or}}{Formula.is:unor}}}\nwidentuses{\\{{\nwixident{And}}{And}}\\{{\nwixident{Iff}}{Iff}}\\{{\nwixident{Implies}}{Implies}}\\{{\nwixident{Or}}{Or}}}\nwindexuse{\nwixident{And}}{And}{NW3w0iRO-3NyJLv-3}\nwindexuse{\nwixident{Iff}}{Iff}{NW3w0iRO-3NyJLv-3}\nwindexuse{\nwixident{Implies}}{Implies}{NW3w0iRO-3NyJLv-3}\nwindexuse{\nwixident{Or}}{Or}{NW3w0iRO-3NyJLv-3}\nwendcode{}\nwbegindocs{32}\nwdocspar
\begin{node}
Quantifiers are essentially ordered pairs of a {\Tt{}\nwlinkedidentq{Term.t}{NW1rTmOJ-4LsRWC-1}\nwendquote} variable and
a formula.
\end{node}

\nwenddocs{}\nwbegincode{33}\sublabel{NW3w0iRO-2HPkc8-6}\nwmargintag{{\nwtagstyle{}\subpageref{NW3w0iRO-2HPkc8-6}}}\moddef{Specification for formulas~{\nwtagstyle{}\subpageref{NW3w0iRO-2HPkc8-1}}}\plusendmoddef\nwstartdeflinemarkup\nwusesondefline{\\{NW3w0iRO-5BH7S-1}}\nwprevnextdefs{NW3w0iRO-2HPkc8-5}{NW3w0iRO-2HPkc8-7}\nwenddeflinemarkup
  val mk_forall : \nwlinkedidentc{Term.t}{NW1rTmOJ-4LsRWC-1} * t -> t;
  val mk_exists : \nwlinkedidentc{Term.t}{NW1rTmOJ-4LsRWC-1} * t -> t;
  
  val dest_forall : t -> \nwlinkedidentc{Term.t}{NW1rTmOJ-4LsRWC-1} * t;
  val dest_exists : t -> \nwlinkedidentc{Term.t}{NW1rTmOJ-4LsRWC-1} * t;

  val is_forall : t -> bool;
  val is_exists : t -> bool;
\nwused{\\{NW3w0iRO-5BH7S-1}}\nwidentuses{\\{{\nwixident{Term.t}}{Term.t}}}\nwindexuse{\nwixident{Term.t}}{Term.t}{NW3w0iRO-2HPkc8-6}\nwendcode{}\nwbegindocs{34}\nwdocspar

\nwenddocs{}\nwbegincode{35}\sublabel{NW3w0iRO-3SlqND-4}\nwmargintag{{\nwtagstyle{}\subpageref{NW3w0iRO-3SlqND-4}}}\moddef{Constructors for formulas~{\nwtagstyle{}\subpageref{NW3w0iRO-3SlqND-1}}}\plusendmoddef\nwstartdeflinemarkup\nwusesondefline{\\{NW3w0iRO-2jnj9W-1}}\nwprevnextdefs{NW3w0iRO-3SlqND-3}{\relax}\nwenddeflinemarkup
fun mk_forall (x, fm) =
  if Term.is_var x then \nwlinkedidentc{Forall}{NW3w0iRO-2jnj9W-1}(x, fm)
  else raise Fail "mk_forall: Term must be a variable";

fun mk_exists (x, fm) =
  if Term.is_var x then \nwlinkedidentc{Exists}{NW3w0iRO-2jnj9W-1}(x, fm)
  else raise Fail "mk_exists: Term must be a variable";
\nwindexdefn{\nwixident{Formula.mk{\_}forall}}{Formula.mk:unforall}{NW3w0iRO-3SlqND-4}\nwindexdefn{\nwixident{Formula.mk{\_}exists}}{Formula.mk:unexists}{NW3w0iRO-3SlqND-4}\eatline
\nwused{\\{NW3w0iRO-2jnj9W-1}}\nwidentdefs{\\{{\nwixident{Formula.mk{\_}exists}}{Formula.mk:unexists}}\\{{\nwixident{Formula.mk{\_}forall}}{Formula.mk:unforall}}}\nwidentuses{\\{{\nwixident{Exists}}{Exists}}\\{{\nwixident{Forall}}{Forall}}}\nwindexuse{\nwixident{Exists}}{Exists}{NW3w0iRO-3SlqND-4}\nwindexuse{\nwixident{Forall}}{Forall}{NW3w0iRO-3SlqND-4}\nwendcode{}\nwbegindocs{36}\nwdocspar
\nwenddocs{}\nwbegincode{37}\sublabel{NW3w0iRO-1Yt1Jl-4}\nwmargintag{{\nwtagstyle{}\subpageref{NW3w0iRO-1Yt1Jl-4}}}\moddef{Destructors for formulas~{\nwtagstyle{}\subpageref{NW3w0iRO-1Yt1Jl-1}}}\plusendmoddef\nwstartdeflinemarkup\nwusesondefline{\\{NW3w0iRO-2jnj9W-1}}\nwprevnextdefs{NW3w0iRO-1Yt1Jl-3}{\relax}\nwenddeflinemarkup
fun dest_forall (\nwlinkedidentc{Forall}{NW3w0iRO-2jnj9W-1} args) = args
  | dest_forall _ = raise Dest "dest_forall";

fun dest_exists (\nwlinkedidentc{Exists}{NW3w0iRO-2jnj9W-1} args) = args
  | dest_exists _ = raise Dest "dest_exists";
\nwused{\\{NW3w0iRO-2jnj9W-1}}\nwidentuses{\\{{\nwixident{Exists}}{Exists}}\\{{\nwixident{Forall}}{Forall}}}\nwindexuse{\nwixident{Exists}}{Exists}{NW3w0iRO-1Yt1Jl-4}\nwindexuse{\nwixident{Forall}}{Forall}{NW3w0iRO-1Yt1Jl-4}\nwendcode{}\nwbegindocs{38}\nwdocspar

\nwenddocs{}\nwbegincode{39}\sublabel{NW3w0iRO-3NyJLv-4}\nwmargintag{{\nwtagstyle{}\subpageref{NW3w0iRO-3NyJLv-4}}}\moddef{Predicates for formulas~{\nwtagstyle{}\subpageref{NW3w0iRO-3NyJLv-1}}}\plusendmoddef\nwstartdeflinemarkup\nwusesondefline{\\{NW3w0iRO-2jnj9W-1}}\nwprevnextdefs{NW3w0iRO-3NyJLv-3}{\relax}\nwenddeflinemarkup
fun is_forall (\nwlinkedidentc{Forall}{NW3w0iRO-2jnj9W-1} _) = true
  | is_forall _ = false;

fun is_exists (\nwlinkedidentc{Exists}{NW3w0iRO-2jnj9W-1} _) = true
  | is_exists _ = false;
\nwused{\\{NW3w0iRO-2jnj9W-1}}\nwidentuses{\\{{\nwixident{Exists}}{Exists}}\\{{\nwixident{Forall}}{Forall}}}\nwindexuse{\nwixident{Exists}}{Exists}{NW3w0iRO-3NyJLv-4}\nwindexuse{\nwixident{Forall}}{Forall}{NW3w0iRO-3NyJLv-4}\nwendcode{}\nwbegindocs{40}\nwdocspar

\begin{node}
Substitution is defined by the rules:
\begin{enumerate}
\item $\bot[t/x]=\bot$
\item $(P[t_{1},\dots,t_{n}])[t/x] = P[t_{1}[t/x], \dots, t_{n}[t/x]]$
\item $(\neg\phi)[t/x]=\neg(\phi[t/x])$
\item $(\phi\land\psi)[t/x]=(\phi[t/x])\land(\psi[t/x])$
\item $(\phi\lor\psi)[t/x]=(\phi[t/x])\lor(\psi[t/x])$
\item $(\phi\implies\psi)[t/x]=(\phi[t/x])\implies(\psi[t/x])$
\item $(\phi\iff\psi)[t/x]=(\phi[t/x])\iff(\psi[t/x])$
\item $(\forall x\ldotp\phi)[t/x]=\forall x\ldotp\phi$
\item $(\forall y\ldotp\phi)[t/x]=\forall y\ldotp(\phi[t/x])$ if
  $y\neq x$
\item $(\exists x\ldotp\phi)[t/x]=\exists x\ldotp\phi$
\item $(\exists y\ldotp\phi)[t/x]=\exists y\ldotp(\phi[t/x])$ if
  $y\neq x$
\end{enumerate}
I suppose an error should be raised if $x$ is not a variable, but I
don't feel strongly enough about it.

We only need to check once that the variable $x$ and its replacement
$t$ have the same sort. The code reflects this, and then iterates
through the subformulas, making all replacements as needed.
\end{node}

\nwenddocs{}\nwbegincode{41}\sublabel{NW3w0iRO-2HPkc8-7}\nwmargintag{{\nwtagstyle{}\subpageref{NW3w0iRO-2HPkc8-7}}}\moddef{Specification for formulas~{\nwtagstyle{}\subpageref{NW3w0iRO-2HPkc8-1}}}\plusendmoddef\nwstartdeflinemarkup\nwusesondefline{\\{NW3w0iRO-5BH7S-1}}\nwprevnextdefs{NW3w0iRO-2HPkc8-6}{NW3w0iRO-2HPkc8-8}\nwenddeflinemarkup
val subst : \nwlinkedidentc{Term.t}{NW1rTmOJ-4LsRWC-1} -> \nwlinkedidentc{Term.t}{NW1rTmOJ-4LsRWC-1} -> t -> t;
\nwused{\\{NW3w0iRO-5BH7S-1}}\nwidentuses{\\{{\nwixident{Term.t}}{Term.t}}}\nwindexuse{\nwixident{Term.t}}{Term.t}{NW3w0iRO-2HPkc8-7}\nwendcode{}\nwbegindocs{42}\nwdocspar

\nwenddocs{}\nwbegincode{43}\sublabel{NW3w0iRO-1jtQRV-1}\nwmargintag{{\nwtagstyle{}\subpageref{NW3w0iRO-1jtQRV-1}}}\moddef{Substitution in a formula~{\nwtagstyle{}\subpageref{NW3w0iRO-1jtQRV-1}}}\endmoddef\nwstartdeflinemarkup\nwusesondefline{\\{NW3w0iRO-2jnj9W-1}}\nwenddeflinemarkup
(* subst : \nwlinkedidentc{Term.t}{NW1rTmOJ-4LsRWC-1} -> \nwlinkedidentc{Term.t}{NW1rTmOJ-4LsRWC-1} -> \nwlinkedidentc{Formula.t}{NW3w0iRO-2jnj9W-1} -> \nwlinkedidentc{Formula.t}{NW3w0iRO-2jnj9W-1} *)
local
fun iter x t fm =
 (case fm of
   (\nwlinkedidentc{False}{NW3w0iRO-2jnj9W-1}) => fm
 | (\nwlinkedidentc{Pred}{NW3w0iRO-2jnj9W-1}(P,args)) => \nwlinkedidentc{Pred}{NW3w0iRO-2jnj9W-1}(P, map (\nwlinkedidentc{Term.subst}{NW1rTmOJ-3OIyM5-1} x t) args)
 | (\nwlinkedidentc{Not}{NW3w0iRO-2jnj9W-1} A) => \nwlinkedidentc{Not}{NW3w0iRO-2jnj9W-1}(iter x t A)
 | (\nwlinkedidentc{And}{NW3w0iRO-2jnj9W-1}(A,B)) => \nwlinkedidentc{And}{NW3w0iRO-2jnj9W-1}(iter x t A, iter x t B)
 | (\nwlinkedidentc{Or}{NW3w0iRO-2jnj9W-1}(A,B)) => \nwlinkedidentc{Or}{NW3w0iRO-2jnj9W-1}(iter x t A, iter x t B)
 | (\nwlinkedidentc{Implies}{NW3w0iRO-2jnj9W-1}(A,B)) => \nwlinkedidentc{Implies}{NW3w0iRO-2jnj9W-1}(iter x t A, iter x t B)
 | (\nwlinkedidentc{Iff}{NW3w0iRO-2jnj9W-1}(A,B)) => \nwlinkedidentc{Iff}{NW3w0iRO-2jnj9W-1}(iter x t A, iter x t B)
 | (\nwlinkedidentc{Forall}{NW3w0iRO-2jnj9W-1}(y,A)) => if \nwlinkedidentc{Term.eq}{NW1rTmOJ-3810sY-1} x y then fm
                    else \nwlinkedidentc{Forall}{NW3w0iRO-2jnj9W-1}(y, iter x t A)
 | (\nwlinkedidentc{Exists}{NW3w0iRO-2jnj9W-1}(y,A)) => if \nwlinkedidentc{Term.eq}{NW1rTmOJ-3810sY-1} x y then fm
                    else \nwlinkedidentc{Exists}{NW3w0iRO-2jnj9W-1}(y, iter x t A));
in
fun subst x t fm =
  if not(\nwlinkedidentc{Term.have_same_sort}{NW1rTmOJ-1WQNps-1} x t)
  then raise Sort.Mismatch "\nwlinkedidentc{Formula.subst}{NW3w0iRO-1jtQRV-1}: variable and replacement have different sorts"
  else iter x t fm
end;
\nwindexdefn{\nwixident{Formula.subst}}{Formula.subst}{NW3w0iRO-1jtQRV-1}\eatline
\nwused{\\{NW3w0iRO-2jnj9W-1}}\nwidentdefs{\\{{\nwixident{Formula.subst}}{Formula.subst}}}\nwidentuses{\\{{\nwixident{And}}{And}}\\{{\nwixident{Exists}}{Exists}}\\{{\nwixident{False}}{False}}\\{{\nwixident{Forall}}{Forall}}\\{{\nwixident{Formula.t}}{Formula.t}}\\{{\nwixident{Iff}}{Iff}}\\{{\nwixident{Implies}}{Implies}}\\{{\nwixident{Not}}{Not}}\\{{\nwixident{Or}}{Or}}\\{{\nwixident{Pred}}{Pred}}\\{{\nwixident{Term.eq}}{Term.eq}}\\{{\nwixident{Term.have{\_}same{\_}sort}}{Term.have:unsame:unsort}}\\{{\nwixident{Term.subst}}{Term.subst}}\\{{\nwixident{Term.t}}{Term.t}}}\nwindexuse{\nwixident{And}}{And}{NW3w0iRO-1jtQRV-1}\nwindexuse{\nwixident{Exists}}{Exists}{NW3w0iRO-1jtQRV-1}\nwindexuse{\nwixident{False}}{False}{NW3w0iRO-1jtQRV-1}\nwindexuse{\nwixident{Forall}}{Forall}{NW3w0iRO-1jtQRV-1}\nwindexuse{\nwixident{Formula.t}}{Formula.t}{NW3w0iRO-1jtQRV-1}\nwindexuse{\nwixident{Iff}}{Iff}{NW3w0iRO-1jtQRV-1}\nwindexuse{\nwixident{Implies}}{Implies}{NW3w0iRO-1jtQRV-1}\nwindexuse{\nwixident{Not}}{Not}{NW3w0iRO-1jtQRV-1}\nwindexuse{\nwixident{Or}}{Or}{NW3w0iRO-1jtQRV-1}\nwindexuse{\nwixident{Pred}}{Pred}{NW3w0iRO-1jtQRV-1}\nwindexuse{\nwixident{Term.eq}}{Term.eq}{NW3w0iRO-1jtQRV-1}\nwindexuse{\nwixident{Term.have{\_}same{\_}sort}}{Term.have:unsame:unsort}{NW3w0iRO-1jtQRV-1}\nwindexuse{\nwixident{Term.subst}}{Term.subst}{NW3w0iRO-1jtQRV-1}\nwindexuse{\nwixident{Term.t}}{Term.t}{NW3w0iRO-1jtQRV-1}\nwendcode{}\nwbegindocs{44}\nwdocspar
\begin{node}
Free variables in  a formula is defined recursively on subformulas by
the rules:
\begin{enumerate}
\item $FV(\bot)=\{\}$
\item $FV(P[t_{1},\dots,t_{n}])=FV(t_{1})\cup(\cdots)\cup FV(t_{n})$
\item $FV(\neg A)=FV(A)$
\item $FV(A\land B)=FV(A)\cup FV(B)$
\item $FV(A\lor B)=FV(A)\cup FV(B)$
\item $FV(A\implies B)=FV(A)\cup FV(B)$
\item $FV(A\iff B)=FV(A)\cup FV(B)$
\item $FV(\forall x\ldotp A)=FV(A)\setminus\{x\}$
\item $FV(\exists x\ldotp A)=FV(A)\setminus\{x\}$
\end{enumerate}
\end{node}


\nwenddocs{}\nwbegincode{45}\sublabel{NW3w0iRO-2HPkc8-8}\nwmargintag{{\nwtagstyle{}\subpageref{NW3w0iRO-2HPkc8-8}}}\moddef{Specification for formulas~{\nwtagstyle{}\subpageref{NW3w0iRO-2HPkc8-1}}}\plusendmoddef\nwstartdeflinemarkup\nwusesondefline{\\{NW3w0iRO-5BH7S-1}}\nwprevnextdefs{NW3w0iRO-2HPkc8-7}{NW3w0iRO-2HPkc8-9}\nwenddeflinemarkup
val fv : t -> \nwlinkedidentc{Term.t}{NW1rTmOJ-4LsRWC-1} list;
\nwused{\\{NW3w0iRO-5BH7S-1}}\nwidentuses{\\{{\nwixident{Term.t}}{Term.t}}}\nwindexuse{\nwixident{Term.t}}{Term.t}{NW3w0iRO-2HPkc8-8}\nwendcode{}\nwbegindocs{46}\nwdocspar

\nwenddocs{}\nwbegincode{47}\sublabel{NW3w0iRO-jRgYZ-1}\nwmargintag{{\nwtagstyle{}\subpageref{NW3w0iRO-jRgYZ-1}}}\moddef{Free and bound variables in a formula~{\nwtagstyle{}\subpageref{NW3w0iRO-jRgYZ-1}}}\endmoddef\nwstartdeflinemarkup\nwusesondefline{\\{NW3w0iRO-2jnj9W-1}}\nwprevnextdefs{\relax}{NW3w0iRO-jRgYZ-2}\nwenddeflinemarkup
(* fv : \nwlinkedidentc{Formula.t}{NW3w0iRO-2jnj9W-1} -> \nwlinkedidentc{Term.t}{NW1rTmOJ-4LsRWC-1} list *)
fun fv (\nwlinkedidentc{False}{NW3w0iRO-2jnj9W-1}) = []
  | fv (\nwlinkedidentc{Pred}{NW3w0iRO-2jnj9W-1} (_, args)) = List.foldl (fn (arg, fvs') => (\nwlinkedidentc{Term.fv}{NW1rTmOJ-2nRtky-1} arg) @ fvs')
                                     []
                                     args
  | fv (\nwlinkedidentc{Not}{NW3w0iRO-2jnj9W-1} A) = fv A
  | fv (\nwlinkedidentc{And}{NW3w0iRO-2jnj9W-1}(A, B)) = (fv A) @ (fv B)
  | fv (\nwlinkedidentc{Or}{NW3w0iRO-2jnj9W-1}(A, B)) = (fv A) @ (fv B)
  | fv (\nwlinkedidentc{Implies}{NW3w0iRO-2jnj9W-1}(A, B)) = (fv A) @ (fv B)
  | fv (\nwlinkedidentc{Iff}{NW3w0iRO-2jnj9W-1}(A, B)) = (fv A) @ (fv B)
  | fv (\nwlinkedidentc{Forall}{NW3w0iRO-2jnj9W-1}(x, A)) = List.filter (not o (\nwlinkedidentc{Term.eq}{NW1rTmOJ-3810sY-1} x)) (fv A)
  | fv (\nwlinkedidentc{Exists}{NW3w0iRO-2jnj9W-1}(x, A)) = List.filter (not o (\nwlinkedidentc{Term.eq}{NW1rTmOJ-3810sY-1} x)) (fv A);
\nwalsodefined{\\{NW3w0iRO-jRgYZ-2}}\nwused{\\{NW3w0iRO-2jnj9W-1}}\nwidentuses{\\{{\nwixident{And}}{And}}\\{{\nwixident{Exists}}{Exists}}\\{{\nwixident{False}}{False}}\\{{\nwixident{Forall}}{Forall}}\\{{\nwixident{Formula.t}}{Formula.t}}\\{{\nwixident{Iff}}{Iff}}\\{{\nwixident{Implies}}{Implies}}\\{{\nwixident{Not}}{Not}}\\{{\nwixident{Or}}{Or}}\\{{\nwixident{Pred}}{Pred}}\\{{\nwixident{Term.eq}}{Term.eq}}\\{{\nwixident{Term.fv}}{Term.fv}}\\{{\nwixident{Term.t}}{Term.t}}}\nwindexuse{\nwixident{And}}{And}{NW3w0iRO-jRgYZ-1}\nwindexuse{\nwixident{Exists}}{Exists}{NW3w0iRO-jRgYZ-1}\nwindexuse{\nwixident{False}}{False}{NW3w0iRO-jRgYZ-1}\nwindexuse{\nwixident{Forall}}{Forall}{NW3w0iRO-jRgYZ-1}\nwindexuse{\nwixident{Formula.t}}{Formula.t}{NW3w0iRO-jRgYZ-1}\nwindexuse{\nwixident{Iff}}{Iff}{NW3w0iRO-jRgYZ-1}\nwindexuse{\nwixident{Implies}}{Implies}{NW3w0iRO-jRgYZ-1}\nwindexuse{\nwixident{Not}}{Not}{NW3w0iRO-jRgYZ-1}\nwindexuse{\nwixident{Or}}{Or}{NW3w0iRO-jRgYZ-1}\nwindexuse{\nwixident{Pred}}{Pred}{NW3w0iRO-jRgYZ-1}\nwindexuse{\nwixident{Term.eq}}{Term.eq}{NW3w0iRO-jRgYZ-1}\nwindexuse{\nwixident{Term.fv}}{Term.fv}{NW3w0iRO-jRgYZ-1}\nwindexuse{\nwixident{Term.t}}{Term.t}{NW3w0iRO-jRgYZ-1}\nwendcode{}\nwbegindocs{48}\nwdocspar

\begin{node}
Bound variables in  a formula is defined recursively on subformulas by
the rules:
\begin{enumerate}
\item $BV(\bot)=\{\}$
\item $BV(P[t_{1},\dots,t_{n}])=\{\}$
\item $BV(\neg A)=BV(A)$
\item $BV(A\land B)=BV(A)\cup BV(B)$
\item $BV(A\lor B)=BV(A)\cup BV(B)$
\item $BV(A\implies B)=BV(A)\cup BV(B)$
\item $BV(A\iff B)=BV(A)\cup BV(B)$
\item $BV(\forall x\ldotp A)=BV(A)\cup\{x\}$
\item $BV(\exists x\ldotp A)=BV(A)\cup\{x\}$
\end{enumerate}
\end{node}
\nwenddocs{}\nwbegincode{49}\sublabel{NW3w0iRO-2HPkc8-9}\nwmargintag{{\nwtagstyle{}\subpageref{NW3w0iRO-2HPkc8-9}}}\moddef{Specification for formulas~{\nwtagstyle{}\subpageref{NW3w0iRO-2HPkc8-1}}}\plusendmoddef\nwstartdeflinemarkup\nwusesondefline{\\{NW3w0iRO-5BH7S-1}}\nwprevnextdefs{NW3w0iRO-2HPkc8-8}{NW3w0iRO-2HPkc8-A}\nwenddeflinemarkup
val bv : t -> \nwlinkedidentc{Term.t}{NW1rTmOJ-4LsRWC-1} list;
\nwused{\\{NW3w0iRO-5BH7S-1}}\nwidentuses{\\{{\nwixident{Term.t}}{Term.t}}}\nwindexuse{\nwixident{Term.t}}{Term.t}{NW3w0iRO-2HPkc8-9}\nwendcode{}\nwbegindocs{50}\nwdocspar

\nwenddocs{}\nwbegincode{51}\sublabel{NW3w0iRO-jRgYZ-2}\nwmargintag{{\nwtagstyle{}\subpageref{NW3w0iRO-jRgYZ-2}}}\moddef{Free and bound variables in a formula~{\nwtagstyle{}\subpageref{NW3w0iRO-jRgYZ-1}}}\plusendmoddef\nwstartdeflinemarkup\nwusesondefline{\\{NW3w0iRO-2jnj9W-1}}\nwprevnextdefs{NW3w0iRO-jRgYZ-1}{\relax}\nwenddeflinemarkup
(* bv : \nwlinkedidentc{Formula.t}{NW3w0iRO-2jnj9W-1} -> \nwlinkedidentc{Term.t}{NW1rTmOJ-4LsRWC-1} list *)
fun bv (\nwlinkedidentc{False}{NW3w0iRO-2jnj9W-1}) = []
  | bv (\nwlinkedidentc{Pred}{NW3w0iRO-2jnj9W-1} _) = []
  | bv (\nwlinkedidentc{Not}{NW3w0iRO-2jnj9W-1} A) = bv A
  | bv (\nwlinkedidentc{And}{NW3w0iRO-2jnj9W-1}(A, B)) = (bv A) @ (bv B)
  | bv (\nwlinkedidentc{Or}{NW3w0iRO-2jnj9W-1}(A, B)) = (bv A) @ (bv B)
  | bv (\nwlinkedidentc{Implies}{NW3w0iRO-2jnj9W-1}(A, B)) = (bv A) @ (bv B)
  | bv (\nwlinkedidentc{Iff}{NW3w0iRO-2jnj9W-1}(A, B)) = (bv A) @ (bv B)
  | bv (\nwlinkedidentc{Forall}{NW3w0iRO-2jnj9W-1}(x, A)) = x::(bv A)
  | bv (\nwlinkedidentc{Exists}{NW3w0iRO-2jnj9W-1}(x, A)) = x::(bv A);
\nwindexdefn{\nwixident{Formula.fv}}{Formula.fv}{NW3w0iRO-jRgYZ-2}\nwindexdefn{\nwixident{Formula.bv}}{Formula.bv}{NW3w0iRO-jRgYZ-2}\eatline
\nwused{\\{NW3w0iRO-2jnj9W-1}}\nwidentdefs{\\{{\nwixident{Formula.bv}}{Formula.bv}}\\{{\nwixident{Formula.fv}}{Formula.fv}}}\nwidentuses{\\{{\nwixident{And}}{And}}\\{{\nwixident{Exists}}{Exists}}\\{{\nwixident{False}}{False}}\\{{\nwixident{Forall}}{Forall}}\\{{\nwixident{Formula.t}}{Formula.t}}\\{{\nwixident{Iff}}{Iff}}\\{{\nwixident{Implies}}{Implies}}\\{{\nwixident{Not}}{Not}}\\{{\nwixident{Or}}{Or}}\\{{\nwixident{Pred}}{Pred}}\\{{\nwixident{Term.t}}{Term.t}}}\nwindexuse{\nwixident{And}}{And}{NW3w0iRO-jRgYZ-2}\nwindexuse{\nwixident{Exists}}{Exists}{NW3w0iRO-jRgYZ-2}\nwindexuse{\nwixident{False}}{False}{NW3w0iRO-jRgYZ-2}\nwindexuse{\nwixident{Forall}}{Forall}{NW3w0iRO-jRgYZ-2}\nwindexuse{\nwixident{Formula.t}}{Formula.t}{NW3w0iRO-jRgYZ-2}\nwindexuse{\nwixident{Iff}}{Iff}{NW3w0iRO-jRgYZ-2}\nwindexuse{\nwixident{Implies}}{Implies}{NW3w0iRO-jRgYZ-2}\nwindexuse{\nwixident{Not}}{Not}{NW3w0iRO-jRgYZ-2}\nwindexuse{\nwixident{Or}}{Or}{NW3w0iRO-jRgYZ-2}\nwindexuse{\nwixident{Pred}}{Pred}{NW3w0iRO-jRgYZ-2}\nwindexuse{\nwixident{Term.t}}{Term.t}{NW3w0iRO-jRgYZ-2}\nwendcode{}\nwbegindocs{52}\nwdocspar
\begin{node}
The calculations for bound variables and free variables follows the
same basic template: tell me what happens on atomic formulas, and what
to do with quantifiers, then tell me how to ``recursively combine
results from subformulas''. Computer scientists call this sort of a
function a ``fold''.
\end{node}

\nwenddocs{}\nwbegincode{53}\sublabel{NW3w0iRO-1UOPDY-1}\nwmargintag{{\nwtagstyle{}\subpageref{NW3w0iRO-1UOPDY-1}}}\moddef{Folding formulas~{\nwtagstyle{}\subpageref{NW3w0iRO-1UOPDY-1}}}\endmoddef\nwstartdeflinemarkup\nwusesondefline{\\{NW3w0iRO-2jnj9W-1}}\nwenddeflinemarkup
(* foldl : (\nwlinkedidentc{Formula.t}{NW3w0iRO-2jnj9W-1} * 'a -> 'a) -> (\nwlinkedidentc{Term.t}{NW1rTmOJ-4LsRWC-1} * 'a -> 'a) -> 'a * \nwlinkedidentc{Formula.t}{NW3w0iRO-2jnj9W-1} -> 'a *)
fun foldl (f : t * 'a -> 'a) (q : \nwlinkedidentc{Term.t}{NW1rTmOJ-4LsRWC-1} * 'a -> 'a) (init : 'a) (fm : t) : 'a =
  (case fm of
    \nwlinkedidentc{False}{NW3w0iRO-2jnj9W-1} => f(fm, init)
  | (\nwlinkedidentc{Pred}{NW3w0iRO-2jnj9W-1} _) => f(fm, init)
  | (\nwlinkedidentc{Not}{NW3w0iRO-2jnj9W-1} A) => foldl f q init A 
  | (\nwlinkedidentc{And}{NW3w0iRO-2jnj9W-1}(A,B)) => foldl f q (foldl f q init A) B
  | (\nwlinkedidentc{Or}{NW3w0iRO-2jnj9W-1}(A,B)) => foldl f q (foldl f q init A) B
  | (\nwlinkedidentc{Implies}{NW3w0iRO-2jnj9W-1}(A,B)) => foldl f q (foldl f q init A) B
  | (\nwlinkedidentc{Iff}{NW3w0iRO-2jnj9W-1}(A,B)) => foldl f q (foldl f q init A) B
  | (\nwlinkedidentc{Forall}{NW3w0iRO-2jnj9W-1}(x,A)) => foldl f q (q(x, init)) A
  | (\nwlinkedidentc{Exists}{NW3w0iRO-2jnj9W-1}(x,A)) => foldl f q (q(x, init)) A);

(* foldr : (\nwlinkedidentc{Formula.t}{NW3w0iRO-2jnj9W-1} * 'a -> 'a) -> (\nwlinkedidentc{Term.t}{NW1rTmOJ-4LsRWC-1} * 'a -> 'a) -> 'a * \nwlinkedidentc{Formula.t}{NW3w0iRO-2jnj9W-1} -> 'a *)
fun foldr f q init fm =
  (case fm of
    \nwlinkedidentc{False}{NW3w0iRO-2jnj9W-1} => f(fm, init)
  | (\nwlinkedidentc{Pred}{NW3w0iRO-2jnj9W-1} _) => f(fm, init)
  | (\nwlinkedidentc{Not}{NW3w0iRO-2jnj9W-1} A) => foldr f q init A 
  | (\nwlinkedidentc{And}{NW3w0iRO-2jnj9W-1}(A,B)) => foldr f q (foldr f q init B) A
  | (\nwlinkedidentc{Or}{NW3w0iRO-2jnj9W-1}(A,B)) => foldr f q (foldr f q init B) A
  | (\nwlinkedidentc{Implies}{NW3w0iRO-2jnj9W-1}(A,B)) => foldr f q (foldr f q init B) A
  | (\nwlinkedidentc{Iff}{NW3w0iRO-2jnj9W-1}(A,B)) => foldr f  q(foldr f q init B) A
  | (\nwlinkedidentc{Forall}{NW3w0iRO-2jnj9W-1}(x,A)) => foldr f q (q(x, init)) A
  | (\nwlinkedidentc{Exists}{NW3w0iRO-2jnj9W-1}(x,A)) => foldr f q (q(x, init)) A);
\nwindexdefn{\nwixident{Formula.foldl}}{Formula.foldl}{NW3w0iRO-1UOPDY-1}\nwindexdefn{\nwixident{Formula.foldr}}{Formula.foldr}{NW3w0iRO-1UOPDY-1}\eatline
\nwused{\\{NW3w0iRO-2jnj9W-1}}\nwidentdefs{\\{{\nwixident{Formula.foldl}}{Formula.foldl}}\\{{\nwixident{Formula.foldr}}{Formula.foldr}}}\nwidentuses{\\{{\nwixident{And}}{And}}\\{{\nwixident{Exists}}{Exists}}\\{{\nwixident{False}}{False}}\\{{\nwixident{Forall}}{Forall}}\\{{\nwixident{Formula.t}}{Formula.t}}\\{{\nwixident{Iff}}{Iff}}\\{{\nwixident{Implies}}{Implies}}\\{{\nwixident{Not}}{Not}}\\{{\nwixident{Or}}{Or}}\\{{\nwixident{Pred}}{Pred}}\\{{\nwixident{Term.t}}{Term.t}}}\nwindexuse{\nwixident{And}}{And}{NW3w0iRO-1UOPDY-1}\nwindexuse{\nwixident{Exists}}{Exists}{NW3w0iRO-1UOPDY-1}\nwindexuse{\nwixident{False}}{False}{NW3w0iRO-1UOPDY-1}\nwindexuse{\nwixident{Forall}}{Forall}{NW3w0iRO-1UOPDY-1}\nwindexuse{\nwixident{Formula.t}}{Formula.t}{NW3w0iRO-1UOPDY-1}\nwindexuse{\nwixident{Iff}}{Iff}{NW3w0iRO-1UOPDY-1}\nwindexuse{\nwixident{Implies}}{Implies}{NW3w0iRO-1UOPDY-1}\nwindexuse{\nwixident{Not}}{Not}{NW3w0iRO-1UOPDY-1}\nwindexuse{\nwixident{Or}}{Or}{NW3w0iRO-1UOPDY-1}\nwindexuse{\nwixident{Pred}}{Pred}{NW3w0iRO-1UOPDY-1}\nwindexuse{\nwixident{Term.t}}{Term.t}{NW3w0iRO-1UOPDY-1}\nwendcode{}\nwbegindocs{54}\nwdocspar

\nwenddocs{}\nwbegincode{55}\sublabel{NW3w0iRO-2HPkc8-A}\nwmargintag{{\nwtagstyle{}\subpageref{NW3w0iRO-2HPkc8-A}}}\moddef{Specification for formulas~{\nwtagstyle{}\subpageref{NW3w0iRO-2HPkc8-1}}}\plusendmoddef\nwstartdeflinemarkup\nwusesondefline{\\{NW3w0iRO-5BH7S-1}}\nwprevnextdefs{NW3w0iRO-2HPkc8-9}{NW3w0iRO-2HPkc8-B}\nwenddeflinemarkup
val contains_var : \nwlinkedidentc{Term.t}{NW1rTmOJ-4LsRWC-1} -> t -> bool;
\nwused{\\{NW3w0iRO-5BH7S-1}}\nwidentuses{\\{{\nwixident{Term.t}}{Term.t}}}\nwindexuse{\nwixident{Term.t}}{Term.t}{NW3w0iRO-2HPkc8-A}\nwendcode{}\nwbegindocs{56}\nwdocspar
\nwenddocs{}\nwbegincode{57}\sublabel{NW3w0iRO-4DOLFK-1}\nwmargintag{{\nwtagstyle{}\subpageref{NW3w0iRO-4DOLFK-1}}}\moddef{Check if a formula contains a variable~{\nwtagstyle{}\subpageref{NW3w0iRO-4DOLFK-1}}}\endmoddef\nwstartdeflinemarkup\nwusesondefline{\\{NW3w0iRO-2jnj9W-1}}\nwenddeflinemarkup
fun contains_var x =
 (fn (\nwlinkedidentc{False}{NW3w0iRO-2jnj9W-1}) => false
   | (\nwlinkedidentc{Pred}{NW3w0iRO-2jnj9W-1} (_, args)) => List.exists (\nwlinkedidentc{Term.occurs_in}{NW1rTmOJ-Re21r-1} x) args
   | (\nwlinkedidentc{Not}{NW3w0iRO-2jnj9W-1} A) => contains_var x A
   | (\nwlinkedidentc{And}{NW3w0iRO-2jnj9W-1}(A,B)) => (contains_var x A) orelse (contains_var x B)
   | (\nwlinkedidentc{Or}{NW3w0iRO-2jnj9W-1}(A,B)) => (contains_var x A) orelse (contains_var x B)
   | (\nwlinkedidentc{Implies}{NW3w0iRO-2jnj9W-1}(A,B)) => (contains_var x A) orelse (contains_var x B)
   | (\nwlinkedidentc{Iff}{NW3w0iRO-2jnj9W-1}(A,B)) => (contains_var x A) orelse (contains_var x B)
   | (\nwlinkedidentc{Forall}{NW3w0iRO-2jnj9W-1}(y, A)) => (\nwlinkedidentc{Term.eq}{NW1rTmOJ-3810sY-1} x y) orelse (contains_var x A)
   | (\nwlinkedidentc{Exists}{NW3w0iRO-2jnj9W-1}(y, A)) => (\nwlinkedidentc{Term.eq}{NW1rTmOJ-3810sY-1} x y) orelse (contains_var x A));
\nwindexdefn{\nwixident{Formula.contains{\_}var}}{Formula.contains:unvar}{NW3w0iRO-4DOLFK-1}\eatline
\nwused{\\{NW3w0iRO-2jnj9W-1}}\nwidentdefs{\\{{\nwixident{Formula.contains{\_}var}}{Formula.contains:unvar}}}\nwidentuses{\\{{\nwixident{And}}{And}}\\{{\nwixident{Exists}}{Exists}}\\{{\nwixident{False}}{False}}\\{{\nwixident{Forall}}{Forall}}\\{{\nwixident{Iff}}{Iff}}\\{{\nwixident{Implies}}{Implies}}\\{{\nwixident{Not}}{Not}}\\{{\nwixident{Or}}{Or}}\\{{\nwixident{Pred}}{Pred}}\\{{\nwixident{Term.eq}}{Term.eq}}\\{{\nwixident{Term.occurs{\_}in}}{Term.occurs:unin}}}\nwindexuse{\nwixident{And}}{And}{NW3w0iRO-4DOLFK-1}\nwindexuse{\nwixident{Exists}}{Exists}{NW3w0iRO-4DOLFK-1}\nwindexuse{\nwixident{False}}{False}{NW3w0iRO-4DOLFK-1}\nwindexuse{\nwixident{Forall}}{Forall}{NW3w0iRO-4DOLFK-1}\nwindexuse{\nwixident{Iff}}{Iff}{NW3w0iRO-4DOLFK-1}\nwindexuse{\nwixident{Implies}}{Implies}{NW3w0iRO-4DOLFK-1}\nwindexuse{\nwixident{Not}}{Not}{NW3w0iRO-4DOLFK-1}\nwindexuse{\nwixident{Or}}{Or}{NW3w0iRO-4DOLFK-1}\nwindexuse{\nwixident{Pred}}{Pred}{NW3w0iRO-4DOLFK-1}\nwindexuse{\nwixident{Term.eq}}{Term.eq}{NW3w0iRO-4DOLFK-1}\nwindexuse{\nwixident{Term.occurs{\_}in}}{Term.occurs:unin}{NW3w0iRO-4DOLFK-1}\nwendcode{}\nwbegindocs{58}\nwdocspar
\begin{node}
The precedence is $\bot > \mbox{predicates} > \neg > \land > \lor > \implies > \iff > \forall > \exists$.
\end{node}

\nwenddocs{}\nwbegincode{59}\sublabel{NW3w0iRO-2HPkc8-B}\nwmargintag{{\nwtagstyle{}\subpageref{NW3w0iRO-2HPkc8-B}}}\moddef{Specification for formulas~{\nwtagstyle{}\subpageref{NW3w0iRO-2HPkc8-1}}}\plusendmoddef\nwstartdeflinemarkup\nwusesondefline{\\{NW3w0iRO-5BH7S-1}}\nwprevnextdefs{NW3w0iRO-2HPkc8-A}{NW3w0iRO-2HPkc8-C}\nwenddeflinemarkup
val precedence : t -> t -> order;
\nwused{\\{NW3w0iRO-5BH7S-1}}\nwendcode{}\nwbegindocs{60}\nwdocspar

\nwenddocs{}\nwbegincode{61}\sublabel{NW3w0iRO-3yI6te-1}\nwmargintag{{\nwtagstyle{}\subpageref{NW3w0iRO-3yI6te-1}}}\moddef{Precedence for connectives~{\nwtagstyle{}\subpageref{NW3w0iRO-3yI6te-1}}}\endmoddef\nwstartdeflinemarkup\nwusesondefline{\\{NW3w0iRO-2jnj9W-1}}\nwenddeflinemarkup
fun precedence (\nwlinkedidentc{False}{NW3w0iRO-2jnj9W-1}) (\nwlinkedidentc{False}{NW3w0iRO-2jnj9W-1}) = EQUAL
  | precedence (\nwlinkedidentc{False}{NW3w0iRO-2jnj9W-1}) _ = GREATER
  | precedence (\nwlinkedidentc{Pred}{NW3w0iRO-2jnj9W-1} _) (\nwlinkedidentc{False}{NW3w0iRO-2jnj9W-1}) = LESS
  | precedence (\nwlinkedidentc{Pred}{NW3w0iRO-2jnj9W-1} _) (\nwlinkedidentc{Pred}{NW3w0iRO-2jnj9W-1} _) = EQUAL
  | precedence (\nwlinkedidentc{Pred}{NW3w0iRO-2jnj9W-1} _) _ = GREATER
  | precedence (\nwlinkedidentc{Not}{NW3w0iRO-2jnj9W-1} _) (\nwlinkedidentc{False}{NW3w0iRO-2jnj9W-1}) = LESS
  | precedence (\nwlinkedidentc{Not}{NW3w0iRO-2jnj9W-1} _) (\nwlinkedidentc{Pred}{NW3w0iRO-2jnj9W-1} _) = LESS
  | precedence (\nwlinkedidentc{Not}{NW3w0iRO-2jnj9W-1} _) _ = GREATER
  | precedence (\nwlinkedidentc{And}{NW3w0iRO-2jnj9W-1} _) (\nwlinkedidentc{False}{NW3w0iRO-2jnj9W-1}) = LESS
  | precedence (\nwlinkedidentc{And}{NW3w0iRO-2jnj9W-1} _) (\nwlinkedidentc{Pred}{NW3w0iRO-2jnj9W-1} _) = LESS
  | precedence (\nwlinkedidentc{And}{NW3w0iRO-2jnj9W-1} _) (\nwlinkedidentc{Not}{NW3w0iRO-2jnj9W-1} _) = LESS
  | precedence (\nwlinkedidentc{And}{NW3w0iRO-2jnj9W-1} _) (\nwlinkedidentc{And}{NW3w0iRO-2jnj9W-1} _) = EQUAL
  | precedence (\nwlinkedidentc{And}{NW3w0iRO-2jnj9W-1} _) _ = GREATER
  | precedence (\nwlinkedidentc{Or}{NW3w0iRO-2jnj9W-1} _) (\nwlinkedidentc{False}{NW3w0iRO-2jnj9W-1}) = LESS
  | precedence (\nwlinkedidentc{Or}{NW3w0iRO-2jnj9W-1} _) (\nwlinkedidentc{Pred}{NW3w0iRO-2jnj9W-1} _) = LESS
  | precedence (\nwlinkedidentc{Or}{NW3w0iRO-2jnj9W-1} _) (\nwlinkedidentc{Not}{NW3w0iRO-2jnj9W-1} _) = LESS
  | precedence (\nwlinkedidentc{Or}{NW3w0iRO-2jnj9W-1} _) (\nwlinkedidentc{And}{NW3w0iRO-2jnj9W-1} _) = LESS
  | precedence (\nwlinkedidentc{Or}{NW3w0iRO-2jnj9W-1} _) (\nwlinkedidentc{Or}{NW3w0iRO-2jnj9W-1} _) = EQUAL
  | precedence (\nwlinkedidentc{Or}{NW3w0iRO-2jnj9W-1} _) _ = GREATER
  | precedence (\nwlinkedidentc{Implies}{NW3w0iRO-2jnj9W-1} _) (\nwlinkedidentc{False}{NW3w0iRO-2jnj9W-1}) = LESS
  | precedence (\nwlinkedidentc{Implies}{NW3w0iRO-2jnj9W-1} _) (\nwlinkedidentc{Pred}{NW3w0iRO-2jnj9W-1} _) = LESS
  | precedence (\nwlinkedidentc{Implies}{NW3w0iRO-2jnj9W-1} _) (\nwlinkedidentc{Not}{NW3w0iRO-2jnj9W-1} _) = LESS
  | precedence (\nwlinkedidentc{Implies}{NW3w0iRO-2jnj9W-1} _) (\nwlinkedidentc{And}{NW3w0iRO-2jnj9W-1} _) = LESS
  | precedence (\nwlinkedidentc{Implies}{NW3w0iRO-2jnj9W-1} _) (\nwlinkedidentc{Or}{NW3w0iRO-2jnj9W-1} _) = LESS
  | precedence (\nwlinkedidentc{Implies}{NW3w0iRO-2jnj9W-1} _) (\nwlinkedidentc{Implies}{NW3w0iRO-2jnj9W-1} _) = EQUAL
  | precedence (\nwlinkedidentc{Implies}{NW3w0iRO-2jnj9W-1} _) _ = GREATER
  | precedence (\nwlinkedidentc{Iff}{NW3w0iRO-2jnj9W-1} _) (\nwlinkedidentc{Forall}{NW3w0iRO-2jnj9W-1} _) = GREATER
  | precedence (\nwlinkedidentc{Iff}{NW3w0iRO-2jnj9W-1} _) (\nwlinkedidentc{Exists}{NW3w0iRO-2jnj9W-1} _) = GREATER
  | precedence (\nwlinkedidentc{Iff}{NW3w0iRO-2jnj9W-1} _) (\nwlinkedidentc{Iff}{NW3w0iRO-2jnj9W-1} _) = EQUAL
  | precedence (\nwlinkedidentc{Iff}{NW3w0iRO-2jnj9W-1} _) _ = LESS
  | precedence (\nwlinkedidentc{Forall}{NW3w0iRO-2jnj9W-1} _) (\nwlinkedidentc{Exists}{NW3w0iRO-2jnj9W-1} _) = GREATER
  | precedence (\nwlinkedidentc{Forall}{NW3w0iRO-2jnj9W-1} _) (\nwlinkedidentc{Forall}{NW3w0iRO-2jnj9W-1} _) = EQUAL
  | precedence (\nwlinkedidentc{Forall}{NW3w0iRO-2jnj9W-1} _) _ = LESS
  | precedence (\nwlinkedidentc{Exists}{NW3w0iRO-2jnj9W-1} _) (\nwlinkedidentc{Exists}{NW3w0iRO-2jnj9W-1} _) = EQUAL
  | precedence (\nwlinkedidentc{Exists}{NW3w0iRO-2jnj9W-1} _) _ = LESS
\nwindexdefn{\nwixident{Formula.precedence}}{Formula.precedence}{NW3w0iRO-3yI6te-1}\eatline
\nwused{\\{NW3w0iRO-2jnj9W-1}}\nwidentdefs{\\{{\nwixident{Formula.precedence}}{Formula.precedence}}}\nwidentuses{\\{{\nwixident{And}}{And}}\\{{\nwixident{Exists}}{Exists}}\\{{\nwixident{False}}{False}}\\{{\nwixident{Forall}}{Forall}}\\{{\nwixident{Iff}}{Iff}}\\{{\nwixident{Implies}}{Implies}}\\{{\nwixident{Not}}{Not}}\\{{\nwixident{Or}}{Or}}\\{{\nwixident{Pred}}{Pred}}}\nwindexuse{\nwixident{And}}{And}{NW3w0iRO-3yI6te-1}\nwindexuse{\nwixident{Exists}}{Exists}{NW3w0iRO-3yI6te-1}\nwindexuse{\nwixident{False}}{False}{NW3w0iRO-3yI6te-1}\nwindexuse{\nwixident{Forall}}{Forall}{NW3w0iRO-3yI6te-1}\nwindexuse{\nwixident{Iff}}{Iff}{NW3w0iRO-3yI6te-1}\nwindexuse{\nwixident{Implies}}{Implies}{NW3w0iRO-3yI6te-1}\nwindexuse{\nwixident{Not}}{Not}{NW3w0iRO-3yI6te-1}\nwindexuse{\nwixident{Or}}{Or}{NW3w0iRO-3yI6te-1}\nwindexuse{\nwixident{Pred}}{Pred}{NW3w0iRO-3yI6te-1}\nwendcode{}\nwbegindocs{62}\nwdocspar
\begin{node}
The expectation is that: if {\Tt{}precedence\ lhs\ rhs\ <>\ EQUAL\nwendquote},
then {\Tt{}compare(lhs,rhs)\ =\ precedence\ lhs\ rhs\nwendquote}. Otherwise, we just
recursively work our way down lexicographically comparing subformulas.
\end{node}

\nwenddocs{}\nwbegincode{63}\sublabel{NW3w0iRO-2HPkc8-C}\nwmargintag{{\nwtagstyle{}\subpageref{NW3w0iRO-2HPkc8-C}}}\moddef{Specification for formulas~{\nwtagstyle{}\subpageref{NW3w0iRO-2HPkc8-1}}}\plusendmoddef\nwstartdeflinemarkup\nwusesondefline{\\{NW3w0iRO-5BH7S-1}}\nwprevnextdefs{NW3w0iRO-2HPkc8-B}{NW3w0iRO-2HPkc8-D}\nwenddeflinemarkup
val compare : t * t -> order;
\nwused{\\{NW3w0iRO-5BH7S-1}}\nwendcode{}\nwbegindocs{64}\nwdocspar
\nwenddocs{}\nwbegincode{65}\sublabel{NW3w0iRO-35s47m-1}\nwmargintag{{\nwtagstyle{}\subpageref{NW3w0iRO-35s47m-1}}}\moddef{Comparing formulas~{\nwtagstyle{}\subpageref{NW3w0iRO-35s47m-1}}}\endmoddef\nwstartdeflinemarkup\nwusesondefline{\\{NW3w0iRO-2jnj9W-1}}\nwenddeflinemarkup
fun compare(lhs, rhs) =
 (case precedence lhs rhs of
   EQUAL => (case (lhs, rhs) of
              (\nwlinkedidentc{False}{NW3w0iRO-2jnj9W-1}, \nwlinkedidentc{False}{NW3w0iRO-2jnj9W-1}) => EQUAL
            | (\nwlinkedidentc{Pred}{NW3w0iRO-2jnj9W-1}(P1,args1),\nwlinkedidentc{Pred}{NW3w0iRO-2jnj9W-1}(P2,args2)) =>
              (case String.compare(P1, P2) of
                EQUAL => \nwlinkedidentc{Term.compare_lists}{NW1rTmOJ-2yJYsC-1} args1 args2
              | r => r)
            | (\nwlinkedidentc{Not}{NW3w0iRO-2jnj9W-1} A, \nwlinkedidentc{Not}{NW3w0iRO-2jnj9W-1} B) => compare(A, B)
            | (\nwlinkedidentc{And}{NW3w0iRO-2jnj9W-1}(A1,B1), \nwlinkedidentc{And}{NW3w0iRO-2jnj9W-1}(A2,B2)) => 
                (case compare(A1, A2) of
                  EQUAL => compare(B1,B2)
                | r => r)
            | (\nwlinkedidentc{Or}{NW3w0iRO-2jnj9W-1}(A1,B1), \nwlinkedidentc{Or}{NW3w0iRO-2jnj9W-1}(A2,B2)) => 
                (case compare(A1, A2) of
                  EQUAL => compare(B1,B2)
                | r => r)
            | (\nwlinkedidentc{Implies}{NW3w0iRO-2jnj9W-1}(A1,B1), \nwlinkedidentc{Implies}{NW3w0iRO-2jnj9W-1}(A2,B2)) => 
                (case compare(A1, A2) of
                  EQUAL => compare(B1,B2)
                | r => r)
            | (\nwlinkedidentc{Iff}{NW3w0iRO-2jnj9W-1}(A1,B1), \nwlinkedidentc{Iff}{NW3w0iRO-2jnj9W-1}(A2,B2)) => 
                (case compare(A1, A2) of
                  EQUAL => compare(B1,B2)
                | r => r)
            | (\nwlinkedidentc{Forall}{NW3w0iRO-2jnj9W-1}(x, A), \nwlinkedidentc{Forall}{NW3w0iRO-2jnj9W-1}(y, B)) =>
                (case \nwlinkedidentc{Term.compare}{NW1rTmOJ-2yJYsC-1}(x, y) of
                  EQUAL => compare(A, B)
                | r => r)
            | (\nwlinkedidentc{Exists}{NW3w0iRO-2jnj9W-1}(x, A), \nwlinkedidentc{Exists}{NW3w0iRO-2jnj9W-1}(y, B)) =>
                (case \nwlinkedidentc{Term.compare}{NW1rTmOJ-2yJYsC-1}(x, y) of
                  EQUAL => compare(A, B)
                | r => r)
            | _ => raise Fail "\nwlinkedidentc{Formula.compare}{NW3w0iRO-35s47m-1}: I shouldn't be here...")
   | r => r);
\nwindexdefn{\nwixident{Formula.compare}}{Formula.compare}{NW3w0iRO-35s47m-1}\eatline
\nwused{\\{NW3w0iRO-2jnj9W-1}}\nwidentdefs{\\{{\nwixident{Formula.compare}}{Formula.compare}}}\nwidentuses{\\{{\nwixident{And}}{And}}\\{{\nwixident{Exists}}{Exists}}\\{{\nwixident{False}}{False}}\\{{\nwixident{Forall}}{Forall}}\\{{\nwixident{Iff}}{Iff}}\\{{\nwixident{Implies}}{Implies}}\\{{\nwixident{Not}}{Not}}\\{{\nwixident{Or}}{Or}}\\{{\nwixident{Pred}}{Pred}}\\{{\nwixident{Term.compare}}{Term.compare}}\\{{\nwixident{Term.compare{\_}lists}}{Term.compare:unlists}}}\nwindexuse{\nwixident{And}}{And}{NW3w0iRO-35s47m-1}\nwindexuse{\nwixident{Exists}}{Exists}{NW3w0iRO-35s47m-1}\nwindexuse{\nwixident{False}}{False}{NW3w0iRO-35s47m-1}\nwindexuse{\nwixident{Forall}}{Forall}{NW3w0iRO-35s47m-1}\nwindexuse{\nwixident{Iff}}{Iff}{NW3w0iRO-35s47m-1}\nwindexuse{\nwixident{Implies}}{Implies}{NW3w0iRO-35s47m-1}\nwindexuse{\nwixident{Not}}{Not}{NW3w0iRO-35s47m-1}\nwindexuse{\nwixident{Or}}{Or}{NW3w0iRO-35s47m-1}\nwindexuse{\nwixident{Pred}}{Pred}{NW3w0iRO-35s47m-1}\nwindexuse{\nwixident{Term.compare}}{Term.compare}{NW3w0iRO-35s47m-1}\nwindexuse{\nwixident{Term.compare{\_}lists}}{Term.compare:unlists}{NW3w0iRO-35s47m-1}\nwendcode{}\nwbegindocs{66}\nwdocspar
\begin{node}
We can test if two formulas are equal to each other in a more
optimized manner, by directly checking each connective matches, and
each subformula matches, and each subterm matches.
\end{node}

\nwenddocs{}\nwbegincode{67}\sublabel{NW3w0iRO-2HPkc8-D}\nwmargintag{{\nwtagstyle{}\subpageref{NW3w0iRO-2HPkc8-D}}}\moddef{Specification for formulas~{\nwtagstyle{}\subpageref{NW3w0iRO-2HPkc8-1}}}\plusendmoddef\nwstartdeflinemarkup\nwusesondefline{\\{NW3w0iRO-5BH7S-1}}\nwprevnextdefs{NW3w0iRO-2HPkc8-C}{NW3w0iRO-2HPkc8-E}\nwenddeflinemarkup
val eq : t -> t -> bool;
\nwused{\\{NW3w0iRO-5BH7S-1}}\nwendcode{}\nwbegindocs{68}\nwdocspar

\nwenddocs{}\nwbegincode{69}\sublabel{NW3w0iRO-5gQzE-1}\nwmargintag{{\nwtagstyle{}\subpageref{NW3w0iRO-5gQzE-1}}}\moddef{Test if two formulas are equal~{\nwtagstyle{}\subpageref{NW3w0iRO-5gQzE-1}}}\endmoddef\nwstartdeflinemarkup\nwusesondefline{\\{NW3w0iRO-2jnj9W-1}}\nwenddeflinemarkup
fun eq (\nwlinkedidentc{False}{NW3w0iRO-2jnj9W-1}) (\nwlinkedidentc{False}{NW3w0iRO-2jnj9W-1}) = true
  | eq (\nwlinkedidentc{Pred}{NW3w0iRO-2jnj9W-1}(P,args1)) (\nwlinkedidentc{Pred}{NW3w0iRO-2jnj9W-1}(Q,args2)) = (P=Q) andalso
                                         (\nwlinkedidentc{Term.eq_lists}{NW1rTmOJ-3810sY-1} args1 args2)
  | eq (\nwlinkedidentc{Not}{NW3w0iRO-2jnj9W-1} A) (\nwlinkedidentc{Not}{NW3w0iRO-2jnj9W-1} B) = eq A B
  | eq (\nwlinkedidentc{And}{NW3w0iRO-2jnj9W-1}(A1,A2)) (\nwlinkedidentc{And}{NW3w0iRO-2jnj9W-1}(B1,B2)) = (eq A1 B1) andalso (eq A2 B2)
  | eq (\nwlinkedidentc{Or}{NW3w0iRO-2jnj9W-1}(A1,A2)) (\nwlinkedidentc{Or}{NW3w0iRO-2jnj9W-1}(B1,B2)) = (eq A1 B1) andalso (eq A2 B2)
  | eq (\nwlinkedidentc{Implies}{NW3w0iRO-2jnj9W-1}(A1,A2)) (\nwlinkedidentc{Implies}{NW3w0iRO-2jnj9W-1}(B1,B2)) = (eq A1 B1) andalso (eq A2 B2)
  | eq (\nwlinkedidentc{Iff}{NW3w0iRO-2jnj9W-1}(A1,A2)) (\nwlinkedidentc{Iff}{NW3w0iRO-2jnj9W-1}(B1,B2)) = (eq A1 B1) andalso (eq A2 B2)
  | eq (\nwlinkedidentc{Forall}{NW3w0iRO-2jnj9W-1}(x,A)) (\nwlinkedidentc{Forall}{NW3w0iRO-2jnj9W-1}(y,B)) = (\nwlinkedidentc{Term.eq}{NW1rTmOJ-3810sY-1} x y) andalso (eq A B)
  | eq (\nwlinkedidentc{Exists}{NW3w0iRO-2jnj9W-1}(x,A)) (\nwlinkedidentc{Exists}{NW3w0iRO-2jnj9W-1}(y,B)) = (\nwlinkedidentc{Term.eq}{NW1rTmOJ-3810sY-1} x y) andalso (eq A B)
  | eq _ _ = false;
\nwindexdefn{\nwixident{Formula.eq}}{Formula.eq}{NW3w0iRO-5gQzE-1}\eatline
\nwused{\\{NW3w0iRO-2jnj9W-1}}\nwidentdefs{\\{{\nwixident{Formula.eq}}{Formula.eq}}}\nwidentuses{\\{{\nwixident{And}}{And}}\\{{\nwixident{Exists}}{Exists}}\\{{\nwixident{False}}{False}}\\{{\nwixident{Forall}}{Forall}}\\{{\nwixident{Iff}}{Iff}}\\{{\nwixident{Implies}}{Implies}}\\{{\nwixident{Not}}{Not}}\\{{\nwixident{Or}}{Or}}\\{{\nwixident{Pred}}{Pred}}\\{{\nwixident{Term.eq}}{Term.eq}}\\{{\nwixident{Term.eq{\_}lists}}{Term.eq:unlists}}}\nwindexuse{\nwixident{And}}{And}{NW3w0iRO-5gQzE-1}\nwindexuse{\nwixident{Exists}}{Exists}{NW3w0iRO-5gQzE-1}\nwindexuse{\nwixident{False}}{False}{NW3w0iRO-5gQzE-1}\nwindexuse{\nwixident{Forall}}{Forall}{NW3w0iRO-5gQzE-1}\nwindexuse{\nwixident{Iff}}{Iff}{NW3w0iRO-5gQzE-1}\nwindexuse{\nwixident{Implies}}{Implies}{NW3w0iRO-5gQzE-1}\nwindexuse{\nwixident{Not}}{Not}{NW3w0iRO-5gQzE-1}\nwindexuse{\nwixident{Or}}{Or}{NW3w0iRO-5gQzE-1}\nwindexuse{\nwixident{Pred}}{Pred}{NW3w0iRO-5gQzE-1}\nwindexuse{\nwixident{Term.eq}}{Term.eq}{NW3w0iRO-5gQzE-1}\nwindexuse{\nwixident{Term.eq{\_}lists}}{Term.eq:unlists}{NW3w0iRO-5gQzE-1}\nwendcode{}\nwbegindocs{70}\nwdocspar
\begin{node}
Last, we will want to test if a term $t$ is free for $x$ in a formula
$\phi$. This means that no free variable of $t$ becomes bound when we
substitute $\phi[t/x]$.
\end{node}

\begin{node}
Let $\phi$ be a formula, $x$ be a variable, and $t$ be a term. Suppose
$x$ and $t$ have the same sort. We say ``$t$ is free for $x$ in $\phi$''
by recursively defining it on subformulas:
\begin{enumerate}
\item If $\phi=\bot$, then $t$ is free for $x$ in $\phi$
\item If $\phi=P[t_{1},\dots,t_{n}]$, then $t$ is free for $x$ in $\phi$
\item If $\phi=\neg A$, then we need to check $t$ is free for $x$ in $A$
\item If $\phi$ is $A\land B$ or $A\lor B$ or $A\implies B$ or $A\iff B$,
  then we need to check $t$ is free for $x$ in both $A$ and $B$
\item If $\phi$ is $\exists y\ldotp A$ or $\forall y\ldotp A$ where
  $y\neq x$, if
  $y\in FV(t)$ and $x$ appears in $A$, then $t$ \textbf{is not} free
  for $x$ in $\phi$ (because $y$ will become bound in $t$ after substitution);
On the other hand, if $y\notin FV(t)$, then we recursively check on
the subformula that $t$ is free for $x$ in $A$.
\end{enumerate}
Alternatively, we can perform the following algorithm: compute $V:=FV(t)$.
\begin{enumerate}
\item If $\phi=\bot$, then return true
\item If $\phi=P[t_{1},\dots,t_{n}]$, then return true
\item If $\phi=\neg A$, then perform the algorithm on $A$
\item If $\phi=A\land B$ or $\phi=A\lor B$ or $\phi=A\implies B$ or
  $\phi=A\iff B$, then perform the algorithm on $A$. If the result is
  true, then perform the algorithm on $B$ and return this as the
  result. On the other hand, if the result of the algorithm on $A$ is
  false, then return false as the result.
\item If $\phi=\exists y\ldotp A$ or $\phi=\forall y\ldotp A$ where
  $y=x$, then return true.
\item If $\phi=\exists y\ldotp A$ or $\phi=\forall y\ldotp A$ where
  $y\neq x$,
  \begin{enumerate}
  \item If $y\in V$ and $x\in FV(A)$, then return false
  \item If $x\notin FV(A)$, then return true
  \item If $y\notin V$, then iterate on the subformula $A$ instead of $\phi$
  \end{enumerate}
\end{enumerate}
\end{node}


\nwenddocs{}\nwbegincode{71}\sublabel{NW3w0iRO-2HPkc8-E}\nwmargintag{{\nwtagstyle{}\subpageref{NW3w0iRO-2HPkc8-E}}}\moddef{Specification for formulas~{\nwtagstyle{}\subpageref{NW3w0iRO-2HPkc8-1}}}\plusendmoddef\nwstartdeflinemarkup\nwusesondefline{\\{NW3w0iRO-5BH7S-1}}\nwprevnextdefs{NW3w0iRO-2HPkc8-D}{NW3w0iRO-2HPkc8-F}\nwenddeflinemarkup
val free_for : \nwlinkedidentc{Term.t}{NW1rTmOJ-4LsRWC-1} -> \nwlinkedidentc{Term.t}{NW1rTmOJ-4LsRWC-1} -> t -> bool;
\nwused{\\{NW3w0iRO-5BH7S-1}}\nwidentuses{\\{{\nwixident{Term.t}}{Term.t}}}\nwindexuse{\nwixident{Term.t}}{Term.t}{NW3w0iRO-2HPkc8-E}\nwendcode{}\nwbegindocs{72}\nwdocspar

\nwenddocs{}\nwbegincode{73}\sublabel{NW3w0iRO-ZZefY-1}\nwmargintag{{\nwtagstyle{}\subpageref{NW3w0iRO-ZZefY-1}}}\moddef{Test a term is free for a variable in a formula~{\nwtagstyle{}\subpageref{NW3w0iRO-ZZefY-1}}}\endmoddef\nwstartdeflinemarkup\nwusesondefline{\\{NW3w0iRO-2jnj9W-1}}\nwenddeflinemarkup
fun member x [] = false
  | member x (y::ys) = (x = y) orelse (member x ys);

fun free_for_iter vars x (fm : t) =
  if is_falsum fm then true
  else if is_pred fm then true
  else if is_not fm
  then let val A = dest_not fm
       in free_for_iter vars x A end
  else if is_and fm
  then let val (A,B) = dest_and fm
       in (free_for_iter vars x A) andalso (free_for_iter vars x B) end
  else if is_or fm
  then let val (A,B) = dest_or fm
       in (free_for_iter vars x A) andalso (free_for_iter vars x B) end
  else if is_imp fm
  then let val (A,B) = dest_imp fm
       in (free_for_iter vars x A) andalso (free_for_iter vars x B) end
  else if is_iff fm
  then let val (A,B) = dest_iff fm
       in (free_for_iter vars x A) andalso (free_for_iter vars x B) end
  else if is_forall fm
  then let val (y,A) = dest_forall fm
       in (\nwlinkedidentc{Term.eq}{NW1rTmOJ-3810sY-1} x y) orelse
          if member y vars
          then not(member x (fv A))
          else free_for_iter vars x A
       end
  else let val (y,A) = dest_exists fm
       in (\nwlinkedidentc{Term.eq}{NW1rTmOJ-3810sY-1} x y) orelse
          if member y vars
          then not(member x (fv A))
          else free_for_iter vars x A
       end;

fun free_for (tm : \nwlinkedidentc{Term.t}{NW1rTmOJ-4LsRWC-1}) (x : \nwlinkedidentc{Term.t}{NW1rTmOJ-4LsRWC-1}) fm =
  free_for_iter (\nwlinkedidentc{Term.fv}{NW1rTmOJ-2nRtky-1} tm) x fm;
\nwindexdefn{\nwixident{Formula.free{\_}for}}{Formula.free:unfor}{NW3w0iRO-ZZefY-1}\eatline
\nwused{\\{NW3w0iRO-2jnj9W-1}}\nwidentdefs{\\{{\nwixident{Formula.free{\_}for}}{Formula.free:unfor}}}\nwidentuses{\\{{\nwixident{Term.eq}}{Term.eq}}\\{{\nwixident{Term.fv}}{Term.fv}}\\{{\nwixident{Term.t}}{Term.t}}}\nwindexuse{\nwixident{Term.eq}}{Term.eq}{NW3w0iRO-ZZefY-1}\nwindexuse{\nwixident{Term.fv}}{Term.fv}{NW3w0iRO-ZZefY-1}\nwindexuse{\nwixident{Term.t}}{Term.t}{NW3w0iRO-ZZefY-1}\nwendcode{}\nwbegindocs{74}\nwdocspar
\begin{node}
Serializing a formula is useful for debugging purposes.
\end{node}

\nwenddocs{}\nwbegincode{75}\sublabel{NW3w0iRO-2HPkc8-F}\nwmargintag{{\nwtagstyle{}\subpageref{NW3w0iRO-2HPkc8-F}}}\moddef{Specification for formulas~{\nwtagstyle{}\subpageref{NW3w0iRO-2HPkc8-1}}}\plusendmoddef\nwstartdeflinemarkup\nwusesondefline{\\{NW3w0iRO-5BH7S-1}}\nwprevnextdefs{NW3w0iRO-2HPkc8-E}{NW3w0iRO-2HPkc8-G}\nwenddeflinemarkup
val serialize : t -> string;
\nwused{\\{NW3w0iRO-5BH7S-1}}\nwendcode{}\nwbegindocs{76}\nwdocspar

\nwenddocs{}\nwbegincode{77}\sublabel{NW3w0iRO-3H9tUA-1}\nwmargintag{{\nwtagstyle{}\subpageref{NW3w0iRO-3H9tUA-1}}}\moddef{Serialize a formula~{\nwtagstyle{}\subpageref{NW3w0iRO-3H9tUA-1}}}\endmoddef\nwstartdeflinemarkup\nwusesondefline{\\{NW3w0iRO-2jnj9W-1}}\nwenddeflinemarkup
fun serialize (\nwlinkedidentc{False}{NW3w0iRO-2jnj9W-1}) = "\nwlinkedidentc{False}{NW3w0iRO-2jnj9W-1}"
  | serialize (\nwlinkedidentc{Pred}{NW3w0iRO-2jnj9W-1}(P,args)) = "\nwlinkedidentc{Pred}{NW3w0iRO-2jnj9W-1}("^
                               P^
                               "["^
                               (String.concatWith ", " (map \nwlinkedidentc{Term.serialize}{NW1rTmOJ-40h2U3-1}
                                                            args))^
                               "])"
  | serialize (\nwlinkedidentc{Not}{NW3w0iRO-2jnj9W-1}(A)) = "\nwlinkedidentc{Not}{NW3w0iRO-2jnj9W-1}("^(serialize A)^")"
  | serialize (\nwlinkedidentc{Or}{NW3w0iRO-2jnj9W-1}(A, B)) = "\nwlinkedidentc{Or}{NW3w0iRO-2jnj9W-1}("^(serialize A)^", "^(serialize B)^")"
  | serialize (\nwlinkedidentc{And}{NW3w0iRO-2jnj9W-1}(A, B)) = "\nwlinkedidentc{And}{NW3w0iRO-2jnj9W-1}("^(serialize A)^", "^(serialize B)^")"
  | serialize (\nwlinkedidentc{Implies}{NW3w0iRO-2jnj9W-1}(A, B)) = "\nwlinkedidentc{Implies}{NW3w0iRO-2jnj9W-1}("^(serialize A)^", "^(serialize B)^")"
  | serialize (\nwlinkedidentc{Iff}{NW3w0iRO-2jnj9W-1}(A, B)) = "\nwlinkedidentc{Iff}{NW3w0iRO-2jnj9W-1}("^(serialize A)^", "^(serialize B)^")"
  | serialize (\nwlinkedidentc{Forall}{NW3w0iRO-2jnj9W-1}(x, A)) = "\nwlinkedidentc{Forall}{NW3w0iRO-2jnj9W-1}("^(\nwlinkedidentc{Term.serialize}{NW1rTmOJ-40h2U3-1} x)^", "^(serialize A)^")"
  | serialize (\nwlinkedidentc{Exists}{NW3w0iRO-2jnj9W-1}(x, A)) = "\nwlinkedidentc{Exists}{NW3w0iRO-2jnj9W-1}("^(\nwlinkedidentc{Term.serialize}{NW1rTmOJ-40h2U3-1} x)^", "^(serialize A)^")";
\nwindexdefn{\nwixident{Formula.serialize}}{Formula.serialize}{NW3w0iRO-3H9tUA-1}\eatline
\nwused{\\{NW3w0iRO-2jnj9W-1}}\nwidentdefs{\\{{\nwixident{Formula.serialize}}{Formula.serialize}}}\nwidentuses{\\{{\nwixident{And}}{And}}\\{{\nwixident{Exists}}{Exists}}\\{{\nwixident{False}}{False}}\\{{\nwixident{Forall}}{Forall}}\\{{\nwixident{Iff}}{Iff}}\\{{\nwixident{Implies}}{Implies}}\\{{\nwixident{Not}}{Not}}\\{{\nwixident{Or}}{Or}}\\{{\nwixident{Pred}}{Pred}}\\{{\nwixident{Term.serialize}}{Term.serialize}}}\nwindexuse{\nwixident{And}}{And}{NW3w0iRO-3H9tUA-1}\nwindexuse{\nwixident{Exists}}{Exists}{NW3w0iRO-3H9tUA-1}\nwindexuse{\nwixident{False}}{False}{NW3w0iRO-3H9tUA-1}\nwindexuse{\nwixident{Forall}}{Forall}{NW3w0iRO-3H9tUA-1}\nwindexuse{\nwixident{Iff}}{Iff}{NW3w0iRO-3H9tUA-1}\nwindexuse{\nwixident{Implies}}{Implies}{NW3w0iRO-3H9tUA-1}\nwindexuse{\nwixident{Not}}{Not}{NW3w0iRO-3H9tUA-1}\nwindexuse{\nwixident{Or}}{Or}{NW3w0iRO-3H9tUA-1}\nwindexuse{\nwixident{Pred}}{Pred}{NW3w0iRO-3H9tUA-1}\nwindexuse{\nwixident{Term.serialize}}{Term.serialize}{NW3w0iRO-3H9tUA-1}\nwendcode{}\nwbegindocs{78}\nwdocspar
\begin{node}
Feferman defines $E^{+}$-formulas as those whose atomic formulas are
$t_{1}=t_{2}$, $t_{1}\neq t_{2}$, and $t\in C$, and the only
connectives allowed are disjunction $\lor$ and conjunction $\land$,
and the only possible quantifier is existentially quantified
individual-sorted variables $\exists x\ldotp\phi$ provided $\phi$ is
another $E^{+}$-formula. We will extend this to support $\top$ as an
$E^{+}$-formula.

These are particularly important in $FS_{0}$ because set comprehension
only makes sense for expressions of the form $\{x\mid\phi[x]\}$ where
$\phi$ is an $E^{+}$-formula.

Why are we sticking this function into the {\Tt{}Formula\nwendquote} module? Well,
the naming convention tells us the predicate fully: {\Tt{}\nwlinkedidentq{Formula.is{\_}e{\_}plus}{NW3w0iRO-1O1TnJ-1}\nwendquote}
tests if a formula is an $E^{+}$-formula or not.
\end{node}

\nwenddocs{}\nwbegincode{79}\sublabel{NW3w0iRO-2HPkc8-G}\nwmargintag{{\nwtagstyle{}\subpageref{NW3w0iRO-2HPkc8-G}}}\moddef{Specification for formulas~{\nwtagstyle{}\subpageref{NW3w0iRO-2HPkc8-1}}}\plusendmoddef\nwstartdeflinemarkup\nwusesondefline{\\{NW3w0iRO-5BH7S-1}}\nwprevnextdefs{NW3w0iRO-2HPkc8-F}{NW3w0iRO-2HPkc8-H}\nwenddeflinemarkup
val is_e_plus : t -> bool;
\nwused{\\{NW3w0iRO-5BH7S-1}}\nwendcode{}\nwbegindocs{80}\nwdocspar

\nwenddocs{}\nwbegincode{81}\sublabel{NW3w0iRO-1O1TnJ-1}\nwmargintag{{\nwtagstyle{}\subpageref{NW3w0iRO-1O1TnJ-1}}}\moddef{Test for $E^{+}$-formula~{\nwtagstyle{}\subpageref{NW3w0iRO-1O1TnJ-1}}}\endmoddef\nwstartdeflinemarkup\nwusesondefline{\\{NW3w0iRO-2jnj9W-1}}\nwenddeflinemarkup
fun is_e_plus (\nwlinkedidentc{Not}{NW3w0iRO-2jnj9W-1}(\nwlinkedidentc{False}{NW3w0iRO-2jnj9W-1})) = true
  | is_e_plus (\nwlinkedidentc{Not}{NW3w0iRO-2jnj9W-1}(\nwlinkedidentc{Pred}{NW3w0iRO-2jnj9W-1}("=",[lhs,rhs]))) = \nwlinkedidentc{Term.is_ind}{NW1rTmOJ-3FfXTO-1} lhs andalso
                                           \nwlinkedidentc{Term.is_ind}{NW1rTmOJ-3FfXTO-1} rhs
  | is_e_plus (\nwlinkedidentc{Pred}{NW3w0iRO-2jnj9W-1}("=",[lhs,rhs])) = \nwlinkedidentc{Term.is_ind}{NW1rTmOJ-3FfXTO-1} lhs andalso
                                      \nwlinkedidentc{Term.is_ind}{NW1rTmOJ-3FfXTO-1} rhs
  | is_e_plus (\nwlinkedidentc{Pred}{NW3w0iRO-2jnj9W-1}("in",[lhs,rhs])) = \nwlinkedidentc{Term.is_ind}{NW1rTmOJ-3FfXTO-1} lhs andalso
                                       \nwlinkedidentc{Term.is_class}{NW1rTmOJ-3FfXTO-1} rhs
  | is_e_plus (\nwlinkedidentc{And}{NW3w0iRO-2jnj9W-1}(A,B)) = is_e_plus A andalso is_e_plus B
  | is_e_plus (\nwlinkedidentc{Or}{NW3w0iRO-2jnj9W-1}(A,B)) = is_e_plus A andalso is_e_plus B
  | is_e_plus (\nwlinkedidentc{Exists}{NW3w0iRO-2jnj9W-1}(x,A)) = \nwlinkedidentc{Term.is_ind}{NW1rTmOJ-3FfXTO-1} x andalso is_e_plus A
  | is_e_plus _ = false; 
\nwindexdefn{\nwixident{Formula.is{\_}e{\_}plus}}{Formula.is:une:unplus}{NW3w0iRO-1O1TnJ-1}\eatline
\nwused{\\{NW3w0iRO-2jnj9W-1}}\nwidentdefs{\\{{\nwixident{Formula.is{\_}e{\_}plus}}{Formula.is:une:unplus}}}\nwidentuses{\\{{\nwixident{And}}{And}}\\{{\nwixident{Exists}}{Exists}}\\{{\nwixident{False}}{False}}\\{{\nwixident{Not}}{Not}}\\{{\nwixident{Or}}{Or}}\\{{\nwixident{Pred}}{Pred}}\\{{\nwixident{Term.is{\_}class}}{Term.is:unclass}}\\{{\nwixident{Term.is{\_}ind}}{Term.is:unind}}}\nwindexuse{\nwixident{And}}{And}{NW3w0iRO-1O1TnJ-1}\nwindexuse{\nwixident{Exists}}{Exists}{NW3w0iRO-1O1TnJ-1}\nwindexuse{\nwixident{False}}{False}{NW3w0iRO-1O1TnJ-1}\nwindexuse{\nwixident{Not}}{Not}{NW3w0iRO-1O1TnJ-1}\nwindexuse{\nwixident{Or}}{Or}{NW3w0iRO-1O1TnJ-1}\nwindexuse{\nwixident{Pred}}{Pred}{NW3w0iRO-1O1TnJ-1}\nwindexuse{\nwixident{Term.is{\_}class}}{Term.is:unclass}{NW3w0iRO-1O1TnJ-1}\nwindexuse{\nwixident{Term.is{\_}ind}}{Term.is:unind}{NW3w0iRO-1O1TnJ-1}\nwendcode{}\nwbegindocs{82}\nwdocspar
\begin{node}
We will also want to apply functions of the form $f\colon
P[\vec{t}]\to \alpha~\mathrm{list}$ to predicates, and to just
recursively work our way down subformulas, concatenating the results together.
\end{node}

\nwenddocs{}\nwbegincode{83}\sublabel{NW3w0iRO-2HPkc8-H}\nwmargintag{{\nwtagstyle{}\subpageref{NW3w0iRO-2HPkc8-H}}}\moddef{Specification for formulas~{\nwtagstyle{}\subpageref{NW3w0iRO-2HPkc8-1}}}\plusendmoddef\nwstartdeflinemarkup\nwusesondefline{\\{NW3w0iRO-5BH7S-1}}\nwprevnextdefs{NW3w0iRO-2HPkc8-G}{\relax}\nwenddeflinemarkup
val atom_union : (t -> 'a list) -> t -> 'a list;
\nwused{\\{NW3w0iRO-5BH7S-1}}\nwendcode{}\nwbegindocs{84}\nwdocspar

\nwenddocs{}\nwbegincode{85}\sublabel{NW3w0iRO-28e6MC-1}\nwmargintag{{\nwtagstyle{}\subpageref{NW3w0iRO-28e6MC-1}}}\moddef{Take the union of applying a function to atoms~{\nwtagstyle{}\subpageref{NW3w0iRO-28e6MC-1}}}\endmoddef\nwstartdeflinemarkup\nwusesondefline{\\{NW3w0iRO-2jnj9W-1}}\nwenddeflinemarkup
(* atom_union : (t -> 'a list) -> 'a list *)
fun atom_union f =
  (fn (\nwlinkedidentc{False}{NW3w0iRO-2jnj9W-1}) => []
    | (fm as (\nwlinkedidentc{Pred}{NW3w0iRO-2jnj9W-1} _)) => f fm
    | (\nwlinkedidentc{Not}{NW3w0iRO-2jnj9W-1} A) => atom_union f A
    | (\nwlinkedidentc{And}{NW3w0iRO-2jnj9W-1}(A,B)) => (atom_union f A) @ (atom_union f B)
    | (\nwlinkedidentc{Or}{NW3w0iRO-2jnj9W-1}(A,B)) => (atom_union f A) @ (atom_union f B)
    | (\nwlinkedidentc{Implies}{NW3w0iRO-2jnj9W-1}(A,B)) => (atom_union f A) @ (atom_union f B)
    | (\nwlinkedidentc{Iff}{NW3w0iRO-2jnj9W-1}(A,B)) => (atom_union f A) @ (atom_union f B)
    | (\nwlinkedidentc{Forall}{NW3w0iRO-2jnj9W-1}(_, A)) => atom_union f A
    | (\nwlinkedidentc{Exists}{NW3w0iRO-2jnj9W-1}(_, A)) => atom_union f A);
\nwindexdefn{\nwixident{Formula.atom{\_}union}}{Formula.atom:ununion}{NW3w0iRO-28e6MC-1}\eatline
\nwused{\\{NW3w0iRO-2jnj9W-1}}\nwidentdefs{\\{{\nwixident{Formula.atom{\_}union}}{Formula.atom:ununion}}}\nwidentuses{\\{{\nwixident{And}}{And}}\\{{\nwixident{Exists}}{Exists}}\\{{\nwixident{False}}{False}}\\{{\nwixident{Forall}}{Forall}}\\{{\nwixident{Iff}}{Iff}}\\{{\nwixident{Implies}}{Implies}}\\{{\nwixident{Not}}{Not}}\\{{\nwixident{Or}}{Or}}\\{{\nwixident{Pred}}{Pred}}}\nwindexuse{\nwixident{And}}{And}{NW3w0iRO-28e6MC-1}\nwindexuse{\nwixident{Exists}}{Exists}{NW3w0iRO-28e6MC-1}\nwindexuse{\nwixident{False}}{False}{NW3w0iRO-28e6MC-1}\nwindexuse{\nwixident{Forall}}{Forall}{NW3w0iRO-28e6MC-1}\nwindexuse{\nwixident{Iff}}{Iff}{NW3w0iRO-28e6MC-1}\nwindexuse{\nwixident{Implies}}{Implies}{NW3w0iRO-28e6MC-1}\nwindexuse{\nwixident{Not}}{Not}{NW3w0iRO-28e6MC-1}\nwindexuse{\nwixident{Or}}{Or}{NW3w0iRO-28e6MC-1}\nwindexuse{\nwixident{Pred}}{Pred}{NW3w0iRO-28e6MC-1}\nwendcode{}\nwfilename{nw/kernel.nw}\nwbegindocs{0}% -*- mode: poly-noweb; noweb-code-mode: sml-mode; -*-
%%
%% kernel.nw
%% 
%% Made by Alex Nelson <pqnelson@gmail.com>
%% Login   <alex@lisp>
%% 
%% Started on  2026-04-02T16:30:01-0700
%% Last update 2026-04-02T16:30:01-0700
%% 

\section{Kernel}

\begin{node}
Alright, I will use a natural deduction proof calculus for
three-sorted Intuitionistic First-Order Logic. In an earlier
experiment for single-sorted Intuitionistic First-Order Logic, I used
a Hilbert proof calculus with a natural deduction judgement (I
``hardcoded'' the Deduction theorem into it). But this also required
some tedious work to eliminate (and introduce) logical connectives,
which could be streamlined with natural deduction.
\end{node}

\nwenddocs{}\nwbegincode{1}\sublabel{NW4ICsJQ-2cY0B1-1}\nwmargintag{{\nwtagstyle{}\subpageref{NW4ICsJQ-2cY0B1-1}}}\moddef{Thm.sig~{\nwtagstyle{}\subpageref{NW4ICsJQ-2cY0B1-1}}}\endmoddef\nwstartdeflinemarkup\nwenddeflinemarkup
signature Thm = sig
  exception Fail of string;
  type t;
  val hyps : t -> \nwlinkedidentc{Formula.t}{NW3w0iRO-2jnj9W-1} list;
  val concl : t -> \nwlinkedidentc{Formula.t}{NW3w0iRO-2jnj9W-1};
  val mk_axiom : \nwlinkedidentc{Formula.t}{NW3w0iRO-2jnj9W-1} -> t; (* danger! *)

  \LA{}Specification for inference rules~{\nwtagstyle{}\subpageref{NW4ICsJQ-1XQlFy-1}}\RA{}

  \LA{}Specification for $FS_{0}$ axioms~{\nwtagstyle{}\subpageref{NW4ICsJQ-hYCIS-1}}\RA{}
  (* Database of axioms
  structure State : sig
              type t;
              val add_axiom : t -> string -> \nwlinkedidentc{Formula.t}{NW3w0iRO-2jnj9W-1} -> t;
            end;x
   *)
end
\nwnotused{Thm.sig}\nwidentuses{\\{{\nwixident{Formula.t}}{Formula.t}}}\nwindexuse{\nwixident{Formula.t}}{Formula.t}{NW4ICsJQ-2cY0B1-1}\nwendcode{}\nwbegindocs{2}\nwdocspar

\begin{node}
We start off with an infixed operator approximating the turnstile.
\end{node}

\nwenddocs{}\nwbegincode{3}\sublabel{NW4ICsJQ-YIPBT-1}\nwmargintag{{\nwtagstyle{}\subpageref{NW4ICsJQ-YIPBT-1}}}\moddef{Thm.sml~{\nwtagstyle{}\subpageref{NW4ICsJQ-YIPBT-1}}}\endmoddef\nwstartdeflinemarkup\nwenddeflinemarkup
structure Thm :> Thm = struct
exception Fail of string;

infix |-;
datatype t = |- of (\nwlinkedidentc{Formula.t}{NW3w0iRO-2jnj9W-1} list) * \nwlinkedidentc{Formula.t}{NW3w0iRO-2jnj9W-1};

fun hyps (Gamma |- _) = Gamma;

fun concl (_ |- phi) = phi;

fun mk_axiom phi = ([] |- phi); 

\LA{}Inference rules~{\nwtagstyle{}\subpageref{NW4ICsJQ-2zKkPU-1}}\RA{}

\LA{}Axioms for $FS_{0}$~{\nwtagstyle{}\subpageref{NW4ICsJQ-2aDdLu-1}}\RA{}
end;
\nwindexdefn{\nwixident{Thm.t}}{Thm.t}{NW4ICsJQ-YIPBT-1}\nwindexdefn{\nwixident{Thm.hyps}}{Thm.hyps}{NW4ICsJQ-YIPBT-1}\nwindexdefn{\nwixident{Thm.concl}}{Thm.concl}{NW4ICsJQ-YIPBT-1}\nwindexdefn{\nwixident{Thm.mk{\_}axiom}}{Thm.mk:unaxiom}{NW4ICsJQ-YIPBT-1}\eatline
\nwnotused{Thm.sml}\nwidentdefs{\\{{\nwixident{Thm.concl}}{Thm.concl}}\\{{\nwixident{Thm.hyps}}{Thm.hyps}}\\{{\nwixident{Thm.mk{\_}axiom}}{Thm.mk:unaxiom}}\\{{\nwixident{Thm.t}}{Thm.t}}}\nwidentuses{\\{{\nwixident{Formula.t}}{Formula.t}}}\nwindexuse{\nwixident{Formula.t}}{Formula.t}{NW4ICsJQ-YIPBT-1}\nwendcode{}\nwbegindocs{4}\nwdocspar
\begin{node}
We need some helper functions for taking the union of hypotheses
$\Gamma_{1}\cup\Gamma_{2}$ (implemented as ``{\Tt{}Gamma1\ \nwlinkedidentq{@@}{NW4ICsJQ-2zKkPU-1}\ Gamma2\nwendquote}'')
and removing hypotheses $\Gamma\setminus\{\phi\}$ (implemented as``{\Tt{}\nwlinkedidentq{remove}{NW4ICsJQ-2zKkPU-1}\ phi\ Gamma\nwendquote}'').
I think a slightly more optimized version would sort the hypotheses,
to remove duplicate entries faster. But this would add
an additional $\mathcal{O}(n\log n)$ runtime to each invocation.
\end{node}

\nwenddocs{}\nwbegincode{5}\sublabel{NW4ICsJQ-2zKkPU-1}\nwmargintag{{\nwtagstyle{}\subpageref{NW4ICsJQ-2zKkPU-1}}}\moddef{Inference rules~{\nwtagstyle{}\subpageref{NW4ICsJQ-2zKkPU-1}}}\endmoddef\nwstartdeflinemarkup\nwusesondefline{\\{NW4ICsJQ-YIPBT-1}}\nwprevnextdefs{\relax}{NW4ICsJQ-2zKkPU-2}\nwenddeflinemarkup
infix \nwlinkedidentc{@@}{NW4ICsJQ-2zKkPU-1}
fun [] \nwlinkedidentc{@@}{NW4ICsJQ-2zKkPU-1} ys = ys
  | xs \nwlinkedidentc{@@}{NW4ICsJQ-2zKkPU-1} [] = xs
  | xs \nwlinkedidentc{@@}{NW4ICsJQ-2zKkPU-1} ys = xs @ ys;

fun \nwlinkedidentc{remove}{NW4ICsJQ-2zKkPU-1} phi Gamma =
  List.filter (not o (\nwlinkedidentc{Formula.eq}{NW3w0iRO-5gQzE-1} phi)) Gamma;
\nwindexdefn{\nwixident{@@}}{@@}{NW4ICsJQ-2zKkPU-1}\nwindexdefn{\nwixident{remove}}{remove}{NW4ICsJQ-2zKkPU-1}\eatline
\nwalsodefined{\\{NW4ICsJQ-2zKkPU-2}\\{NW4ICsJQ-2zKkPU-3}\\{NW4ICsJQ-2zKkPU-4}\\{NW4ICsJQ-2zKkPU-5}\\{NW4ICsJQ-2zKkPU-6}\\{NW4ICsJQ-2zKkPU-7}\\{NW4ICsJQ-2zKkPU-8}\\{NW4ICsJQ-2zKkPU-9}\\{NW4ICsJQ-2zKkPU-A}\\{NW4ICsJQ-2zKkPU-B}}\nwused{\\{NW4ICsJQ-YIPBT-1}}\nwidentdefs{\\{{\nwixident{@@}}{@@}}\\{{\nwixident{remove}}{remove}}}\nwidentuses{\\{{\nwixident{Formula.eq}}{Formula.eq}}}\nwindexuse{\nwixident{Formula.eq}}{Formula.eq}{NW4ICsJQ-2zKkPU-1}\nwendcode{}\nwbegindocs{6}\nwdocspar
\begin{node}
We have the inference rules related to assumptions:
\begin{equation}
\vcenter{\infer[\texttt{assume}(\phi,\Gamma)]{\phi,\Gamma\vdash\phi}{}}
\end{equation}
\begin{equation}
\vcenter{\infer[\texttt{disch}(\phi)]{\Gamma\setminus\{\phi\}\vdash\phi\implies\psi}{\Gamma\vdash\psi,&\phi\in\Gamma}}
\end{equation}
\begin{equation}
\vcenter{\infer[\texttt{undisch}]{\phi,\Gamma\vdash\psi}{\Gamma\vdash\phi\implies\psi}}
\end{equation}
\begin{equation}
\vcenter{\infer[\texttt{modus\_ponens}]{\Gamma_{1}\cup\Gamma_{2}\vdash\psi}{\Gamma_{1}\vdash\phi\implies\psi,&\Gamma_{2}\vdash\phi}}
\end{equation}
The other inference rules for connectives are straightforward
``introduction'' and ``elimination'' rules.
\end{node}

\nwenddocs{}\nwbegincode{7}\sublabel{NW4ICsJQ-1XQlFy-1}\nwmargintag{{\nwtagstyle{}\subpageref{NW4ICsJQ-1XQlFy-1}}}\moddef{Specification for inference rules~{\nwtagstyle{}\subpageref{NW4ICsJQ-1XQlFy-1}}}\endmoddef\nwstartdeflinemarkup\nwusesondefline{\\{NW4ICsJQ-2cY0B1-1}}\nwprevnextdefs{\relax}{NW4ICsJQ-1XQlFy-2}\nwenddeflinemarkup
val assume : \nwlinkedidentc{Formula.t}{NW3w0iRO-2jnj9W-1} -> \nwlinkedidentc{Formula.t}{NW3w0iRO-2jnj9W-1} list -> t;
val disch : \nwlinkedidentc{Formula.t}{NW3w0iRO-2jnj9W-1} -> t -> t;
val undisch : t -> t;
val modus_ponens : t -> t -> t;
\nwalsodefined{\\{NW4ICsJQ-1XQlFy-2}\\{NW4ICsJQ-1XQlFy-3}\\{NW4ICsJQ-1XQlFy-4}\\{NW4ICsJQ-1XQlFy-5}\\{NW4ICsJQ-1XQlFy-6}\\{NW4ICsJQ-1XQlFy-7}\\{NW4ICsJQ-1XQlFy-8}\\{NW4ICsJQ-1XQlFy-9}\\{NW4ICsJQ-1XQlFy-A}}\nwused{\\{NW4ICsJQ-2cY0B1-1}}\nwidentuses{\\{{\nwixident{Formula.t}}{Formula.t}}}\nwindexuse{\nwixident{Formula.t}}{Formula.t}{NW4ICsJQ-1XQlFy-1}\nwendcode{}\nwbegindocs{8}\nwdocspar
\nwenddocs{}\nwbegincode{9}\sublabel{NW4ICsJQ-2zKkPU-2}\nwmargintag{{\nwtagstyle{}\subpageref{NW4ICsJQ-2zKkPU-2}}}\moddef{Inference rules~{\nwtagstyle{}\subpageref{NW4ICsJQ-2zKkPU-1}}}\plusendmoddef\nwstartdeflinemarkup\nwusesondefline{\\{NW4ICsJQ-YIPBT-1}}\nwprevnextdefs{NW4ICsJQ-2zKkPU-1}{NW4ICsJQ-2zKkPU-3}\nwenddeflinemarkup
fun assume phi Gamma = (phi::Gamma) |- phi;

fun disch phi (Gamma |- psi) = (\nwlinkedidentc{remove}{NW4ICsJQ-2zKkPU-1} phi Gamma) |- \nwlinkedidentc{Formula.mk_imp}{NW3w0iRO-3SlqND-3}(phi,psi);

fun undisch (Gamma |- A) =
  if not(\nwlinkedidentc{Formula.is_imp}{NW3w0iRO-3NyJLv-3} A)
  then raise Fail "\nwlinkedidentc{Thm.undisch}{NW4ICsJQ-2zKkPU-2}: conclusion must be an implication"
  else let val (phi,psi) = \nwlinkedidentc{Formula.dest_imp}{NW3w0iRO-1Yt1Jl-3} A
       in (phi::Gamma) |- psi end;

fun modus_ponens (Gamma1 |- imp) (Gamma2 |- phi) =
  if not(\nwlinkedidentc{Formula.is_imp}{NW3w0iRO-3NyJLv-3} imp)
  then raise Fail "modus_ponens: major premise expected implication"
  else let
    val (A, B) = \nwlinkedidentc{Formula.dest_imp}{NW3w0iRO-1Yt1Jl-3} imp
  in
    if not(\nwlinkedidentc{Formula.eq}{NW3w0iRO-5gQzE-1} A phi)
    then raise Fail "modus_ponens: premise does not match argument"
    else (Gamma1 \nwlinkedidentc{@@}{NW4ICsJQ-2zKkPU-1} Gamma2) |- B
  end;
\nwindexdefn{\nwixident{Thm.assume}}{Thm.assume}{NW4ICsJQ-2zKkPU-2}\nwindexdefn{\nwixident{Thm.disch}}{Thm.disch}{NW4ICsJQ-2zKkPU-2}\nwindexdefn{\nwixident{Thm.undisch}}{Thm.undisch}{NW4ICsJQ-2zKkPU-2}\nwindexdefn{\nwixident{Thm.modus{\_}ponens}}{Thm.modus:unponens}{NW4ICsJQ-2zKkPU-2}\eatline
\nwused{\\{NW4ICsJQ-YIPBT-1}}\nwidentdefs{\\{{\nwixident{Thm.assume}}{Thm.assume}}\\{{\nwixident{Thm.disch}}{Thm.disch}}\\{{\nwixident{Thm.modus{\_}ponens}}{Thm.modus:unponens}}\\{{\nwixident{Thm.undisch}}{Thm.undisch}}}\nwidentuses{\\{{\nwixident{@@}}{@@}}\\{{\nwixident{Formula.dest{\_}imp}}{Formula.dest:unimp}}\\{{\nwixident{Formula.eq}}{Formula.eq}}\\{{\nwixident{Formula.is{\_}imp}}{Formula.is:unimp}}\\{{\nwixident{Formula.mk{\_}imp}}{Formula.mk:unimp}}\\{{\nwixident{remove}}{remove}}}\nwindexuse{\nwixident{@@}}{@@}{NW4ICsJQ-2zKkPU-2}\nwindexuse{\nwixident{Formula.dest{\_}imp}}{Formula.dest:unimp}{NW4ICsJQ-2zKkPU-2}\nwindexuse{\nwixident{Formula.eq}}{Formula.eq}{NW4ICsJQ-2zKkPU-2}\nwindexuse{\nwixident{Formula.is{\_}imp}}{Formula.is:unimp}{NW4ICsJQ-2zKkPU-2}\nwindexuse{\nwixident{Formula.mk{\_}imp}}{Formula.mk:unimp}{NW4ICsJQ-2zKkPU-2}\nwindexuse{\nwixident{remove}}{remove}{NW4ICsJQ-2zKkPU-2}\nwendcode{}\nwbegindocs{10}\nwdocspar
\begin{node}
The rules for conjunction
\begin{equation}
\vcenter{\infer[\land\mbox{-intro}]{\Gamma_{1}\cup\Gamma_{2}\vdash\phi\land\psi}{\Gamma_{1}\vdash\phi,&\Gamma_{2}\vdash\psi}}
\end{equation}
\begin{equation}
\vcenter{\infer[\land\mbox{-elim-l}]{\Gamma\vdash\phi}{\Gamma\vdash\phi\land\psi}}
\end{equation}
\begin{equation}
\vcenter{\infer[\land\mbox{-elim-r}]{\Gamma\vdash\psi}{\Gamma\vdash\phi\land\psi}}
\end{equation}
\end{node}

\nwenddocs{}\nwbegincode{11}\sublabel{NW4ICsJQ-1XQlFy-2}\nwmargintag{{\nwtagstyle{}\subpageref{NW4ICsJQ-1XQlFy-2}}}\moddef{Specification for inference rules~{\nwtagstyle{}\subpageref{NW4ICsJQ-1XQlFy-1}}}\plusendmoddef\nwstartdeflinemarkup\nwusesondefline{\\{NW4ICsJQ-2cY0B1-1}}\nwprevnextdefs{NW4ICsJQ-1XQlFy-1}{NW4ICsJQ-1XQlFy-3}\nwenddeflinemarkup
val and_intro : t -> t -> t;
val and_elim_l : t -> t;
val and_elim_r : t -> t;
\nwused{\\{NW4ICsJQ-2cY0B1-1}}\nwendcode{}\nwbegindocs{12}\nwdocspar

\nwenddocs{}\nwbegincode{13}\sublabel{NW4ICsJQ-2zKkPU-3}\nwmargintag{{\nwtagstyle{}\subpageref{NW4ICsJQ-2zKkPU-3}}}\moddef{Inference rules~{\nwtagstyle{}\subpageref{NW4ICsJQ-2zKkPU-1}}}\plusendmoddef\nwstartdeflinemarkup\nwusesondefline{\\{NW4ICsJQ-YIPBT-1}}\nwprevnextdefs{NW4ICsJQ-2zKkPU-2}{NW4ICsJQ-2zKkPU-4}\nwenddeflinemarkup
fun and_intro (Gamma1 |- phi) (Gamma2 |- psi) =
  (Gamma1 \nwlinkedidentc{@@}{NW4ICsJQ-2zKkPU-1} Gamma2) |- \nwlinkedidentc{Formula.mk_and}{NW3w0iRO-3SlqND-3}(phi, psi);

fun and_elim_l (Gamma |- phi) =
  if not(\nwlinkedidentc{Formula.is_and}{NW3w0iRO-3NyJLv-3} phi)
  then raise Fail "and_elim_l: expects conjunction"
  else
    let val (A,B) = \nwlinkedidentc{Formula.dest_and}{NW3w0iRO-1Yt1Jl-3} phi
    in Gamma |- A end;

fun and_elim_r (Gamma |- phi) =
  if not(\nwlinkedidentc{Formula.is_and}{NW3w0iRO-3NyJLv-3} phi)
  then raise Fail "and_elim_r: expects conjunction"
  else
    let val (A,B) = \nwlinkedidentc{Formula.dest_and}{NW3w0iRO-1Yt1Jl-3} phi
    in Gamma |- B end;
\nwindexdefn{\nwixident{Formula.and{\_}intro}}{Formula.and:unintro}{NW4ICsJQ-2zKkPU-3}\nwindexdefn{\nwixident{Formula.and{\_}elim{\_}l}}{Formula.and:unelim:unl}{NW4ICsJQ-2zKkPU-3}\nwindexdefn{\nwixident{Formula.and{\_}elim{\_}r}}{Formula.and:unelim:unr}{NW4ICsJQ-2zKkPU-3}\eatline
\nwused{\\{NW4ICsJQ-YIPBT-1}}\nwidentdefs{\\{{\nwixident{Formula.and{\_}elim{\_}l}}{Formula.and:unelim:unl}}\\{{\nwixident{Formula.and{\_}elim{\_}r}}{Formula.and:unelim:unr}}\\{{\nwixident{Formula.and{\_}intro}}{Formula.and:unintro}}}\nwidentuses{\\{{\nwixident{@@}}{@@}}\\{{\nwixident{Formula.dest{\_}and}}{Formula.dest:unand}}\\{{\nwixident{Formula.is{\_}and}}{Formula.is:unand}}\\{{\nwixident{Formula.mk{\_}and}}{Formula.mk:unand}}}\nwindexuse{\nwixident{@@}}{@@}{NW4ICsJQ-2zKkPU-3}\nwindexuse{\nwixident{Formula.dest{\_}and}}{Formula.dest:unand}{NW4ICsJQ-2zKkPU-3}\nwindexuse{\nwixident{Formula.is{\_}and}}{Formula.is:unand}{NW4ICsJQ-2zKkPU-3}\nwindexuse{\nwixident{Formula.mk{\_}and}}{Formula.mk:unand}{NW4ICsJQ-2zKkPU-3}\nwendcode{}\nwbegindocs{14}\nwdocspar
\begin{node}
Disjunction introduction is straightforward:
\begin{equation}
\vcenter{\infer[\lor\mbox{-intro-l}(\phi)]{\Gamma\vdash\phi\lor\psi}{\Gamma\vdash\psi}}
\end{equation}
\begin{equation}
\vcenter{\infer[\lor\mbox{-intro-r}(\psi)]{\Gamma\vdash\phi\lor\psi}{\Gamma\vdash\phi}}
\end{equation}
Disjunction elimination requires some explanation. If we have proven
$\Gamma\vdash A\lor B$ and we want to prove $\Gamma\vdash C$, then we
just need to establish $\Gamma\vdash A\implies C$ and $\Gamma\vdash B\implies C$.
In Mizar, these correspond to the two \texttt{suppose} blocks. The
rule is 
\begin{equation}
\vcenter{\infer[\lor\mbox{-elim}]{\Gamma_{1}\cup\Gamma_{2}\vdash(A\lor B)\implies C}{\Gamma_{1}\vdash A\implies C,&\Gamma_{2}\vdash B\implies C}}
\end{equation}
If any of the premises mismatch each other (e.g., we have
$\Gamma_{2}\vdash A'\implies C$ for some $A\neq A'$), then an
exception is raised.
\end{node}

\nwenddocs{}\nwbegincode{15}\sublabel{NW4ICsJQ-1XQlFy-3}\nwmargintag{{\nwtagstyle{}\subpageref{NW4ICsJQ-1XQlFy-3}}}\moddef{Specification for inference rules~{\nwtagstyle{}\subpageref{NW4ICsJQ-1XQlFy-1}}}\plusendmoddef\nwstartdeflinemarkup\nwusesondefline{\\{NW4ICsJQ-2cY0B1-1}}\nwprevnextdefs{NW4ICsJQ-1XQlFy-2}{NW4ICsJQ-1XQlFy-4}\nwenddeflinemarkup
val or_intro_l : \nwlinkedidentc{Formula.t}{NW3w0iRO-2jnj9W-1} -> t -> t;
val or_intro_r : \nwlinkedidentc{Formula.t}{NW3w0iRO-2jnj9W-1} -> t -> t;
val or_elim : t -> t -> t;
\nwused{\\{NW4ICsJQ-2cY0B1-1}}\nwidentuses{\\{{\nwixident{Formula.t}}{Formula.t}}}\nwindexuse{\nwixident{Formula.t}}{Formula.t}{NW4ICsJQ-1XQlFy-3}\nwendcode{}\nwbegindocs{16}\nwdocspar

\nwenddocs{}\nwbegincode{17}\sublabel{NW4ICsJQ-2zKkPU-4}\nwmargintag{{\nwtagstyle{}\subpageref{NW4ICsJQ-2zKkPU-4}}}\moddef{Inference rules~{\nwtagstyle{}\subpageref{NW4ICsJQ-2zKkPU-1}}}\plusendmoddef\nwstartdeflinemarkup\nwusesondefline{\\{NW4ICsJQ-YIPBT-1}}\nwprevnextdefs{NW4ICsJQ-2zKkPU-3}{NW4ICsJQ-2zKkPU-5}\nwenddeflinemarkup
fun or_intro_l A (Gamma |- B) = (Gamma |- \nwlinkedidentc{Formula.mk_or}{NW3w0iRO-3SlqND-3}(A, B))

fun or_intro_r B (Gamma |- A) = (Gamma |- \nwlinkedidentc{Formula.mk_or}{NW3w0iRO-3SlqND-3}(A, B))

fun or_elim (Gamma1 |- A_imp_C) (Gamma2 |- B_imp_C) =
  if not(\nwlinkedidentc{Formula.is_imp}{NW3w0iRO-3NyJLv-3} A_imp_C)
  then raise Fail "or_elim: first theorem must be an implication"
  else if not(\nwlinkedidentc{Formula.is_imp}{NW3w0iRO-3NyJLv-3} B_imp_C)
  then raise Fail "or_elim: second theorem must be an implication"
  else
    let
      val (A,C1) = \nwlinkedidentc{Formula.dest_imp}{NW3w0iRO-1Yt1Jl-3} A_imp_C;
      val (B,C2) = \nwlinkedidentc{Formula.dest_imp}{NW3w0iRO-1Yt1Jl-3} B_imp_C;
      val A_or_B_imp_C = \nwlinkedidentc{Formula.mk_imp}{NW3w0iRO-3SlqND-3}(\nwlinkedidentc{Formula.mk_or}{NW3w0iRO-3SlqND-3}(A, B), C1);
    in
      if not(\nwlinkedidentc{Formula.eq}{NW3w0iRO-5gQzE-1} C1 C2)
      then raise Fail "or_elim: conclusions of theorems disagree"
      else (Gamma1 \nwlinkedidentc{@@}{NW4ICsJQ-2zKkPU-1} Gamma2) |- A_or_B_imp_C
    end;
\nwindexdefn{\nwixident{Thm.or{\_}elim}}{Thm.or:unelim}{NW4ICsJQ-2zKkPU-4}\nwindexdefn{\nwixident{Thm.or{\_}intro{\_}l}}{Thm.or:unintro:unl}{NW4ICsJQ-2zKkPU-4}\nwindexdefn{\nwixident{Thm.or{\_}intro{\_}r}}{Thm.or:unintro:unr}{NW4ICsJQ-2zKkPU-4}\eatline
\nwused{\\{NW4ICsJQ-YIPBT-1}}\nwidentdefs{\\{{\nwixident{Thm.or{\_}elim}}{Thm.or:unelim}}\\{{\nwixident{Thm.or{\_}intro{\_}l}}{Thm.or:unintro:unl}}\\{{\nwixident{Thm.or{\_}intro{\_}r}}{Thm.or:unintro:unr}}}\nwidentuses{\\{{\nwixident{@@}}{@@}}\\{{\nwixident{Formula.dest{\_}imp}}{Formula.dest:unimp}}\\{{\nwixident{Formula.eq}}{Formula.eq}}\\{{\nwixident{Formula.is{\_}imp}}{Formula.is:unimp}}\\{{\nwixident{Formula.mk{\_}imp}}{Formula.mk:unimp}}\\{{\nwixident{Formula.mk{\_}or}}{Formula.mk:unor}}}\nwindexuse{\nwixident{@@}}{@@}{NW4ICsJQ-2zKkPU-4}\nwindexuse{\nwixident{Formula.dest{\_}imp}}{Formula.dest:unimp}{NW4ICsJQ-2zKkPU-4}\nwindexuse{\nwixident{Formula.eq}}{Formula.eq}{NW4ICsJQ-2zKkPU-4}\nwindexuse{\nwixident{Formula.is{\_}imp}}{Formula.is:unimp}{NW4ICsJQ-2zKkPU-4}\nwindexuse{\nwixident{Formula.mk{\_}imp}}{Formula.mk:unimp}{NW4ICsJQ-2zKkPU-4}\nwindexuse{\nwixident{Formula.mk{\_}or}}{Formula.mk:unor}{NW4ICsJQ-2zKkPU-4}\nwendcode{}\nwbegindocs{18}\nwdocspar
\begin{node}
Negation-related rules
\begin{equation}
\vcenter{\infer[\neg\mbox{-intro}]{\Gamma\vdash\neg\phi}{\Gamma\vdash\phi\implies\bot}}
\end{equation}
\begin{equation}
\vcenter{\infer[\neg\mbox{-elim}]{\Gamma\vdash\phi\implies\bot}{\Gamma\vdash\neg\phi}}
\end{equation}
\end{node}

\nwenddocs{}\nwbegincode{19}\sublabel{NW4ICsJQ-1XQlFy-4}\nwmargintag{{\nwtagstyle{}\subpageref{NW4ICsJQ-1XQlFy-4}}}\moddef{Specification for inference rules~{\nwtagstyle{}\subpageref{NW4ICsJQ-1XQlFy-1}}}\plusendmoddef\nwstartdeflinemarkup\nwusesondefline{\\{NW4ICsJQ-2cY0B1-1}}\nwprevnextdefs{NW4ICsJQ-1XQlFy-3}{NW4ICsJQ-1XQlFy-5}\nwenddeflinemarkup
val not_intro : t -> t;
val not_elim : t -> t;
\nwused{\\{NW4ICsJQ-2cY0B1-1}}\nwendcode{}\nwbegindocs{20}\nwdocspar

\nwenddocs{}\nwbegincode{21}\sublabel{NW4ICsJQ-2zKkPU-5}\nwmargintag{{\nwtagstyle{}\subpageref{NW4ICsJQ-2zKkPU-5}}}\moddef{Inference rules~{\nwtagstyle{}\subpageref{NW4ICsJQ-2zKkPU-1}}}\plusendmoddef\nwstartdeflinemarkup\nwusesondefline{\\{NW4ICsJQ-YIPBT-1}}\nwprevnextdefs{NW4ICsJQ-2zKkPU-4}{NW4ICsJQ-2zKkPU-6}\nwenddeflinemarkup
fun not_intro (Gamma |- phi_imp_false) =
  if not(\nwlinkedidentc{Formula.is_imp}{NW3w0iRO-3NyJLv-3} phi_imp_false)
  then raise Fail "not_intro: theorem is not an implication"
  else let
    val (phi, F) = \nwlinkedidentc{Formula.dest_imp}{NW3w0iRO-1Yt1Jl-3} phi_imp_false;
  in
    if not(\nwlinkedidentc{Formula.is_falsum}{NW3w0iRO-3NyJLv-1} F)
    then raise Fail "not_intro: theorem is not 'implies falsum'"
    else Gamma |- \nwlinkedidentc{Formula.mk_not}{NW3w0iRO-3SlqND-2}(phi)
  end;

fun not_elim (Gamma |- not_phi) =
  if not(\nwlinkedidentc{Formula.is_not}{NW3w0iRO-3NyJLv-2} not_phi)
  then raise Fail "not_elim: theorem is not a negation"
  else let
    val (phi) = \nwlinkedidentc{Formula.dest_not}{NW3w0iRO-1Yt1Jl-2} not_phi;
  in
    Gamma |- \nwlinkedidentc{Formula.mk_imp}{NW3w0iRO-3SlqND-3}(phi, \nwlinkedidentc{Formula.mk_falsum}{NW3w0iRO-3SlqND-1}())
  end;
\nwindexdefn{\nwixident{Thm.not{\_}intro}}{Thm.not:unintro}{NW4ICsJQ-2zKkPU-5}\nwindexdefn{\nwixident{Thm.not{\_}elim}}{Thm.not:unelim}{NW4ICsJQ-2zKkPU-5}\eatline
\nwused{\\{NW4ICsJQ-YIPBT-1}}\nwidentdefs{\\{{\nwixident{Thm.not{\_}elim}}{Thm.not:unelim}}\\{{\nwixident{Thm.not{\_}intro}}{Thm.not:unintro}}}\nwidentuses{\\{{\nwixident{Formula.dest{\_}imp}}{Formula.dest:unimp}}\\{{\nwixident{Formula.dest{\_}not}}{Formula.dest:unnot}}\\{{\nwixident{Formula.is{\_}falsum}}{Formula.is:unfalsum}}\\{{\nwixident{Formula.is{\_}imp}}{Formula.is:unimp}}\\{{\nwixident{Formula.is{\_}not}}{Formula.is:unnot}}\\{{\nwixident{Formula.mk{\_}falsum}}{Formula.mk:unfalsum}}\\{{\nwixident{Formula.mk{\_}imp}}{Formula.mk:unimp}}\\{{\nwixident{Formula.mk{\_}not}}{Formula.mk:unnot}}}\nwindexuse{\nwixident{Formula.dest{\_}imp}}{Formula.dest:unimp}{NW4ICsJQ-2zKkPU-5}\nwindexuse{\nwixident{Formula.dest{\_}not}}{Formula.dest:unnot}{NW4ICsJQ-2zKkPU-5}\nwindexuse{\nwixident{Formula.is{\_}falsum}}{Formula.is:unfalsum}{NW4ICsJQ-2zKkPU-5}\nwindexuse{\nwixident{Formula.is{\_}imp}}{Formula.is:unimp}{NW4ICsJQ-2zKkPU-5}\nwindexuse{\nwixident{Formula.is{\_}not}}{Formula.is:unnot}{NW4ICsJQ-2zKkPU-5}\nwindexuse{\nwixident{Formula.mk{\_}falsum}}{Formula.mk:unfalsum}{NW4ICsJQ-2zKkPU-5}\nwindexuse{\nwixident{Formula.mk{\_}imp}}{Formula.mk:unimp}{NW4ICsJQ-2zKkPU-5}\nwindexuse{\nwixident{Formula.mk{\_}not}}{Formula.mk:unnot}{NW4ICsJQ-2zKkPU-5}\nwendcode{}\nwbegindocs{22}\nwdocspar
\begin{node}
The Intuitionistic rule for eliminating contradictions
\begin{equation}
\vcenter{\infer[\mbox{contr}(\phi)]{\Gamma\vdash\phi}{\Gamma\vdash\bot}}
\end{equation}
\end{node}

\nwenddocs{}\nwbegincode{23}\sublabel{NW4ICsJQ-1XQlFy-5}\nwmargintag{{\nwtagstyle{}\subpageref{NW4ICsJQ-1XQlFy-5}}}\moddef{Specification for inference rules~{\nwtagstyle{}\subpageref{NW4ICsJQ-1XQlFy-1}}}\plusendmoddef\nwstartdeflinemarkup\nwusesondefline{\\{NW4ICsJQ-2cY0B1-1}}\nwprevnextdefs{NW4ICsJQ-1XQlFy-4}{NW4ICsJQ-1XQlFy-6}\nwenddeflinemarkup
val contr : \nwlinkedidentc{Formula.t}{NW3w0iRO-2jnj9W-1} -> t -> t;
\nwused{\\{NW4ICsJQ-2cY0B1-1}}\nwidentuses{\\{{\nwixident{Formula.t}}{Formula.t}}}\nwindexuse{\nwixident{Formula.t}}{Formula.t}{NW4ICsJQ-1XQlFy-5}\nwendcode{}\nwbegindocs{24}\nwdocspar

\nwenddocs{}\nwbegincode{25}\sublabel{NW4ICsJQ-2zKkPU-6}\nwmargintag{{\nwtagstyle{}\subpageref{NW4ICsJQ-2zKkPU-6}}}\moddef{Inference rules~{\nwtagstyle{}\subpageref{NW4ICsJQ-2zKkPU-1}}}\plusendmoddef\nwstartdeflinemarkup\nwusesondefline{\\{NW4ICsJQ-YIPBT-1}}\nwprevnextdefs{NW4ICsJQ-2zKkPU-5}{NW4ICsJQ-2zKkPU-7}\nwenddeflinemarkup
fun contr phi (Gamma |- falsum) =
  if not(\nwlinkedidentc{Formula.is_falsum}{NW3w0iRO-3NyJLv-1} falsum)
  then raise Fail "contr: theorem is not falsum"
  else Gamma |- phi;
\nwindexdefn{\nwixident{Thm.contr}}{Thm.contr}{NW4ICsJQ-2zKkPU-6}\eatline
\nwused{\\{NW4ICsJQ-YIPBT-1}}\nwidentdefs{\\{{\nwixident{Thm.contr}}{Thm.contr}}}\nwidentuses{\\{{\nwixident{Formula.is{\_}falsum}}{Formula.is:unfalsum}}}\nwindexuse{\nwixident{Formula.is{\_}falsum}}{Formula.is:unfalsum}{NW4ICsJQ-2zKkPU-6}\nwendcode{}\nwbegindocs{26}\nwdocspar
\begin{node}
Biconditionals $A\iff B$ may be viewed as an abbreviation for
$(A\implies B)\land(B\implies A)$. The introduction and elimination
rules correspond to the derived rules for the abbreviation.
\begin{equation}
\vcenter{\infer[\Leftrightarrow\mbox{-intro}]{\Gamma_{1}\cup\Gamma_{2}\vdash\phi\iff\psi}{\Gamma_{1}\vdash\phi\implies\psi,&\Gamma_{2}\vdash\psi\implies\phi}}
\end{equation}
\begin{equation}
\vcenter{\infer[\Leftrightarrow\mbox{-elim-l}]{\Gamma\vdash\phi\implies\psi}{\Gamma\vdash\phi\iff\psi}}
\end{equation}
\begin{equation}
\vcenter{\infer[\Leftrightarrow\mbox{-elim-r}]{\Gamma\vdash\psi\implies\phi}{\Gamma\vdash\phi\iff\psi}}
\end{equation}
\end{node}

\nwenddocs{}\nwbegincode{27}\sublabel{NW4ICsJQ-1XQlFy-6}\nwmargintag{{\nwtagstyle{}\subpageref{NW4ICsJQ-1XQlFy-6}}}\moddef{Specification for inference rules~{\nwtagstyle{}\subpageref{NW4ICsJQ-1XQlFy-1}}}\plusendmoddef\nwstartdeflinemarkup\nwusesondefline{\\{NW4ICsJQ-2cY0B1-1}}\nwprevnextdefs{NW4ICsJQ-1XQlFy-5}{NW4ICsJQ-1XQlFy-7}\nwenddeflinemarkup
val iff_intro : t -> t -> t;
val iff_elim_l : t -> t;
val iff_elim_r : t -> t;
\nwused{\\{NW4ICsJQ-2cY0B1-1}}\nwendcode{}\nwbegindocs{28}\nwdocspar
\nwenddocs{}\nwbegincode{29}\sublabel{NW4ICsJQ-2zKkPU-7}\nwmargintag{{\nwtagstyle{}\subpageref{NW4ICsJQ-2zKkPU-7}}}\moddef{Inference rules~{\nwtagstyle{}\subpageref{NW4ICsJQ-2zKkPU-1}}}\plusendmoddef\nwstartdeflinemarkup\nwusesondefline{\\{NW4ICsJQ-YIPBT-1}}\nwprevnextdefs{NW4ICsJQ-2zKkPU-6}{NW4ICsJQ-2zKkPU-8}\nwenddeflinemarkup
fun iff_intro (Gamma1 |- A_imp_B) (Gamma2 |- B_imp_A) =
  if not(\nwlinkedidentc{Formula.is_imp}{NW3w0iRO-3NyJLv-3} A_imp_B)
  then raise Fail "iff_intro: first theorem is not an implication"
  else if not(\nwlinkedidentc{Formula.is_imp}{NW3w0iRO-3NyJLv-3} B_imp_A)
  then raise Fail "iff_intro: second theorem is not an implication"
  else let
    val (A1,B1) = \nwlinkedidentc{Formula.dest_imp}{NW3w0iRO-1Yt1Jl-3}(A_imp_B);
    val (B2,A2) = \nwlinkedidentc{Formula.dest_imp}{NW3w0iRO-1Yt1Jl-3}(B_imp_A);
  in
    if not(\nwlinkedidentc{Formula.eq}{NW3w0iRO-5gQzE-1} A1 A2)
    then raise Fail("iff_intro: premise of first theorem ("^
                    (\nwlinkedidentc{Formula.serialize}{NW3w0iRO-3H9tUA-1} A1)^
                    ") is not conclusion of second theorem ("^
                    (\nwlinkedidentc{Formula.serialize}{NW3w0iRO-3H9tUA-1} A2)^
                    ")")
    else if not(\nwlinkedidentc{Formula.eq}{NW3w0iRO-5gQzE-1} B1 B2)
    then raise Fail("iff_intro: premise of second theorem ("^
                    (\nwlinkedidentc{Formula.serialize}{NW3w0iRO-3H9tUA-1} B2)
                    ^") is not conclusion of first theorem ("^
                    (\nwlinkedidentc{Formula.serialize}{NW3w0iRO-3H9tUA-1} B1)
                    ^")")
    else (Gamma1 \nwlinkedidentc{@@}{NW4ICsJQ-2zKkPU-1} Gamma2) |- \nwlinkedidentc{Formula.mk_iff}{NW3w0iRO-3SlqND-3}(A1,B2)
  end;

fun iff_elim_l (Gamma |- A_iff_B) =
  if not(\nwlinkedidentc{Formula.is_iff}{NW3w0iRO-3NyJLv-3} A_iff_B)
  then raise Fail "iff_elim_l: theorem is not an iff"
  else
    let val (A,B) = \nwlinkedidentc{Formula.dest_iff}{NW3w0iRO-1Yt1Jl-3} A_iff_B
    in Gamma |- \nwlinkedidentc{Formula.mk_imp}{NW3w0iRO-3SlqND-3}(A,B) end;

fun iff_elim_r (Gamma |- A_iff_B) =
  if not(\nwlinkedidentc{Formula.is_iff}{NW3w0iRO-3NyJLv-3} A_iff_B)
  then raise Fail "iff_elim_r: theorem is not an iff"
  else
    let val (A,B) = \nwlinkedidentc{Formula.dest_iff}{NW3w0iRO-1Yt1Jl-3} A_iff_B
    in Gamma |- \nwlinkedidentc{Formula.mk_imp}{NW3w0iRO-3SlqND-3}(B,A) end;
\nwindexdefn{\nwixident{Thm.iff{\_}intro}}{Thm.iff:unintro}{NW4ICsJQ-2zKkPU-7}\nwindexdefn{\nwixident{Thm.iff{\_}elim{\_}l}}{Thm.iff:unelim:unl}{NW4ICsJQ-2zKkPU-7}\nwindexdefn{\nwixident{Thm.iff{\_}elim{\_}r}}{Thm.iff:unelim:unr}{NW4ICsJQ-2zKkPU-7}\eatline
\nwused{\\{NW4ICsJQ-YIPBT-1}}\nwidentdefs{\\{{\nwixident{Thm.iff{\_}elim{\_}l}}{Thm.iff:unelim:unl}}\\{{\nwixident{Thm.iff{\_}elim{\_}r}}{Thm.iff:unelim:unr}}\\{{\nwixident{Thm.iff{\_}intro}}{Thm.iff:unintro}}}\nwidentuses{\\{{\nwixident{@@}}{@@}}\\{{\nwixident{Formula.dest{\_}iff}}{Formula.dest:uniff}}\\{{\nwixident{Formula.dest{\_}imp}}{Formula.dest:unimp}}\\{{\nwixident{Formula.eq}}{Formula.eq}}\\{{\nwixident{Formula.is{\_}iff}}{Formula.is:uniff}}\\{{\nwixident{Formula.is{\_}imp}}{Formula.is:unimp}}\\{{\nwixident{Formula.mk{\_}iff}}{Formula.mk:uniff}}\\{{\nwixident{Formula.mk{\_}imp}}{Formula.mk:unimp}}\\{{\nwixident{Formula.serialize}}{Formula.serialize}}}\nwindexuse{\nwixident{@@}}{@@}{NW4ICsJQ-2zKkPU-7}\nwindexuse{\nwixident{Formula.dest{\_}iff}}{Formula.dest:uniff}{NW4ICsJQ-2zKkPU-7}\nwindexuse{\nwixident{Formula.dest{\_}imp}}{Formula.dest:unimp}{NW4ICsJQ-2zKkPU-7}\nwindexuse{\nwixident{Formula.eq}}{Formula.eq}{NW4ICsJQ-2zKkPU-7}\nwindexuse{\nwixident{Formula.is{\_}iff}}{Formula.is:uniff}{NW4ICsJQ-2zKkPU-7}\nwindexuse{\nwixident{Formula.is{\_}imp}}{Formula.is:unimp}{NW4ICsJQ-2zKkPU-7}\nwindexuse{\nwixident{Formula.mk{\_}iff}}{Formula.mk:uniff}{NW4ICsJQ-2zKkPU-7}\nwindexuse{\nwixident{Formula.mk{\_}imp}}{Formula.mk:unimp}{NW4ICsJQ-2zKkPU-7}\nwindexuse{\nwixident{Formula.serialize}}{Formula.serialize}{NW4ICsJQ-2zKkPU-7}\nwendcode{}\nwbegindocs{30}\nwdocspar
\begin{node}[Quantifier inference rules]
We have the inference rule for generalization
\begin{equation*}
\vcenter{\infer[\forall\mbox{-intro}(x)]{\Gamma\vdash\forall x\ldotp\phi}{\Gamma\vdash\phi,&x\notin FV(\Gamma)}}
\end{equation*}
\end{node}

\nwenddocs{}\nwbegincode{31}\sublabel{NW4ICsJQ-1XQlFy-7}\nwmargintag{{\nwtagstyle{}\subpageref{NW4ICsJQ-1XQlFy-7}}}\moddef{Specification for inference rules~{\nwtagstyle{}\subpageref{NW4ICsJQ-1XQlFy-1}}}\plusendmoddef\nwstartdeflinemarkup\nwusesondefline{\\{NW4ICsJQ-2cY0B1-1}}\nwprevnextdefs{NW4ICsJQ-1XQlFy-6}{NW4ICsJQ-1XQlFy-8}\nwenddeflinemarkup
val forall_intro : \nwlinkedidentc{Term.t}{NW1rTmOJ-4LsRWC-1} -> t -> t;
\nwused{\\{NW4ICsJQ-2cY0B1-1}}\nwidentuses{\\{{\nwixident{Term.t}}{Term.t}}}\nwindexuse{\nwixident{Term.t}}{Term.t}{NW4ICsJQ-1XQlFy-7}\nwendcode{}\nwbegindocs{32}\nwdocspar

\nwenddocs{}\nwbegincode{33}\sublabel{NW4ICsJQ-2zKkPU-8}\nwmargintag{{\nwtagstyle{}\subpageref{NW4ICsJQ-2zKkPU-8}}}\moddef{Inference rules~{\nwtagstyle{}\subpageref{NW4ICsJQ-2zKkPU-1}}}\plusendmoddef\nwstartdeflinemarkup\nwusesondefline{\\{NW4ICsJQ-YIPBT-1}}\nwprevnextdefs{NW4ICsJQ-2zKkPU-7}{NW4ICsJQ-2zKkPU-9}\nwenddeflinemarkup
fun forall_intro (x : \nwlinkedidentc{Term.t}{NW1rTmOJ-4LsRWC-1}) (Gamma |- phi) =
  if not(Term.is_var x)
  then raise Fail "forall_intro: x is not a variable"
  else let fun free_in fm = List.exists (\nwlinkedidentc{Term.eq}{NW1rTmOJ-3810sY-1} x) (\nwlinkedidentc{Formula.fv}{NW3w0iRO-jRgYZ-2} fm);
       in if List.exists free_in Gamma
          then raise Fail "forall_intro: x is free in hypotheses"
          else Gamma |- (\nwlinkedidentc{Formula.mk_forall}{NW3w0iRO-3SlqND-4}(x, phi))
       end;
\nwindexdefn{\nwixident{Thm.forall{\_}intro}}{Thm.forall:unintro}{NW4ICsJQ-2zKkPU-8}\eatline
\nwused{\\{NW4ICsJQ-YIPBT-1}}\nwidentdefs{\\{{\nwixident{Thm.forall{\_}intro}}{Thm.forall:unintro}}}\nwidentuses{\\{{\nwixident{Formula.fv}}{Formula.fv}}\\{{\nwixident{Formula.mk{\_}forall}}{Formula.mk:unforall}}\\{{\nwixident{Term.eq}}{Term.eq}}\\{{\nwixident{Term.t}}{Term.t}}}\nwindexuse{\nwixident{Formula.fv}}{Formula.fv}{NW4ICsJQ-2zKkPU-8}\nwindexuse{\nwixident{Formula.mk{\_}forall}}{Formula.mk:unforall}{NW4ICsJQ-2zKkPU-8}\nwindexuse{\nwixident{Term.eq}}{Term.eq}{NW4ICsJQ-2zKkPU-8}\nwindexuse{\nwixident{Term.t}}{Term.t}{NW4ICsJQ-2zKkPU-8}\nwendcode{}\nwbegindocs{34}\nwdocspar
\begin{node}
Universal quantifier elimination is quite simple provided $t$ is free
for $x$ in $\phi[-]$ (i.e., no free variable will become bound in $\phi[t]$):
\begin{equation}
\vcenter{\infer[\forall\mbox{-elim}]{\Gamma\vdash\phi[t]}{\Gamma\vdash\forall x\ldotp\phi[x]}}
\end{equation}
\end{node}

\nwenddocs{}\nwbegincode{35}\sublabel{NW4ICsJQ-1XQlFy-8}\nwmargintag{{\nwtagstyle{}\subpageref{NW4ICsJQ-1XQlFy-8}}}\moddef{Specification for inference rules~{\nwtagstyle{}\subpageref{NW4ICsJQ-1XQlFy-1}}}\plusendmoddef\nwstartdeflinemarkup\nwusesondefline{\\{NW4ICsJQ-2cY0B1-1}}\nwprevnextdefs{NW4ICsJQ-1XQlFy-7}{NW4ICsJQ-1XQlFy-9}\nwenddeflinemarkup
val forall_elim : \nwlinkedidentc{Term.t}{NW1rTmOJ-4LsRWC-1} -> t -> t;
\nwused{\\{NW4ICsJQ-2cY0B1-1}}\nwidentuses{\\{{\nwixident{Term.t}}{Term.t}}}\nwindexuse{\nwixident{Term.t}}{Term.t}{NW4ICsJQ-1XQlFy-8}\nwendcode{}\nwbegindocs{36}\nwdocspar

\nwenddocs{}\nwbegincode{37}\sublabel{NW4ICsJQ-2zKkPU-9}\nwmargintag{{\nwtagstyle{}\subpageref{NW4ICsJQ-2zKkPU-9}}}\moddef{Inference rules~{\nwtagstyle{}\subpageref{NW4ICsJQ-2zKkPU-1}}}\plusendmoddef\nwstartdeflinemarkup\nwusesondefline{\\{NW4ICsJQ-YIPBT-1}}\nwprevnextdefs{NW4ICsJQ-2zKkPU-8}{NW4ICsJQ-2zKkPU-A}\nwenddeflinemarkup
fun forall_elim (t : \nwlinkedidentc{Term.t}{NW1rTmOJ-4LsRWC-1}) (Gamma |- forall_phi) =
  if not(Formula.is_forall forall_phi)
  then raise Fail "forall_elim: theorem is not universally quantified"
  else let val (x,phi) = Formula.dest_forall forall_phi
       in if not(\nwlinkedidentc{Formula.free_for}{NW3w0iRO-ZZefY-1} t x phi)
          then raise Fail "forall_elim: term is not free for conclusion"
          else Gamma |- (\nwlinkedidentc{Formula.subst}{NW3w0iRO-1jtQRV-1} x t phi)
       end;
\nwindexdefn{\nwixident{Thm.forall{\_}elim}}{Thm.forall:unelim}{NW4ICsJQ-2zKkPU-9}\eatline
\nwused{\\{NW4ICsJQ-YIPBT-1}}\nwidentdefs{\\{{\nwixident{Thm.forall{\_}elim}}{Thm.forall:unelim}}}\nwidentuses{\\{{\nwixident{Formula.free{\_}for}}{Formula.free:unfor}}\\{{\nwixident{Formula.subst}}{Formula.subst}}\\{{\nwixident{Term.t}}{Term.t}}}\nwindexuse{\nwixident{Formula.free{\_}for}}{Formula.free:unfor}{NW4ICsJQ-2zKkPU-9}\nwindexuse{\nwixident{Formula.subst}}{Formula.subst}{NW4ICsJQ-2zKkPU-9}\nwindexuse{\nwixident{Term.t}}{Term.t}{NW4ICsJQ-2zKkPU-9}\nwendcode{}\nwbegindocs{38}\nwdocspar
\begin{node}\label{node:kernel.nw:existential-introduction}
Existential introduction has the same proviso as universal quantifier
elimination:
\begin{equation}
\vcenter{\infer[\exists\mbox{-intro}]{\Gamma\vdash\exists x\ldotp\phi[x]}{\Gamma\vdash\phi[t]}}
\end{equation}
Well, there is some subtlety with existential introduction. For
example, $\phi[t]:=(\forall x\ldotp t=x)$ where $x\notin FV(t)$. Then
$t$ is free for $x$ in $\phi[-]$, but $\exists x\ldotp\phi[x]$ should
be flagged as problematic.

But this situation should never occur because $\phi[x]:=\forall x\ldotp(x=x)$
has $\phi[t]=\forall x\ldotp(x=x)$ which is syntactically different
from $(\forall x\ldotp t=x)$. So an error would be raised, catching
this situation. We add some useful explanatory text in this situation.

A more helpful thing to do would be to rename the bound variable, so
as to avoid this collision.
\end{node}

\nwenddocs{}\nwbegincode{39}\sublabel{NW4ICsJQ-1XQlFy-9}\nwmargintag{{\nwtagstyle{}\subpageref{NW4ICsJQ-1XQlFy-9}}}\moddef{Specification for inference rules~{\nwtagstyle{}\subpageref{NW4ICsJQ-1XQlFy-1}}}\plusendmoddef\nwstartdeflinemarkup\nwusesondefline{\\{NW4ICsJQ-2cY0B1-1}}\nwprevnextdefs{NW4ICsJQ-1XQlFy-8}{NW4ICsJQ-1XQlFy-A}\nwenddeflinemarkup
val exists_intro : \nwlinkedidentc{Term.t}{NW1rTmOJ-4LsRWC-1} -> \nwlinkedidentc{Term.t}{NW1rTmOJ-4LsRWC-1} -> \nwlinkedidentc{Formula.t}{NW3w0iRO-2jnj9W-1} -> t -> t;
\nwused{\\{NW4ICsJQ-2cY0B1-1}}\nwidentuses{\\{{\nwixident{Formula.t}}{Formula.t}}\\{{\nwixident{Term.t}}{Term.t}}}\nwindexuse{\nwixident{Formula.t}}{Formula.t}{NW4ICsJQ-1XQlFy-9}\nwindexuse{\nwixident{Term.t}}{Term.t}{NW4ICsJQ-1XQlFy-9}\nwendcode{}\nwbegindocs{40}\nwdocspar

\nwenddocs{}\nwbegincode{41}\sublabel{NW4ICsJQ-2zKkPU-A}\nwmargintag{{\nwtagstyle{}\subpageref{NW4ICsJQ-2zKkPU-A}}}\moddef{Inference rules~{\nwtagstyle{}\subpageref{NW4ICsJQ-2zKkPU-1}}}\plusendmoddef\nwstartdeflinemarkup\nwusesondefline{\\{NW4ICsJQ-YIPBT-1}}\nwprevnextdefs{NW4ICsJQ-2zKkPU-9}{NW4ICsJQ-2zKkPU-B}\nwenddeflinemarkup
fun exists_intro (x : \nwlinkedidentc{Term.t}{NW1rTmOJ-4LsRWC-1}) (t : \nwlinkedidentc{Term.t}{NW1rTmOJ-4LsRWC-1}) (phi : \nwlinkedidentc{Formula.t}{NW3w0iRO-2jnj9W-1}) (Gamma |- phi_t) =
  if not(Term.is_var(x))
  then raise Fail "exists_intro: x is not a variable"
  else if not(\nwlinkedidentc{Term.have_same_sort}{NW1rTmOJ-1WQNps-1} x t)
  then raise Fail "exists_intro: x and t of different sorts"
  else if not(\nwlinkedidentc{Formula.eq}{NW3w0iRO-5gQzE-1} (\nwlinkedidentc{Formula.subst}{NW3w0iRO-1jtQRV-1} x t phi) phi_t)
  then (if List.exists (\nwlinkedidentc{Term.eq}{NW1rTmOJ-3810sY-1} x) (\nwlinkedidentc{Formula.bv}{NW3w0iRO-jRgYZ-2} phi)
        then raise Fail (
            concat["exists_intro: conclusion mismatch supplied formula ",
                   "(probably x is a bound variable with 't' in its scope)"])
        else raise Fail("exists_intro: conclusion "^
                        (\nwlinkedidentc{Formula.serialize}{NW3w0iRO-3H9tUA-1} phi_t)^
                        " mismatch supplied formula "^
                        (\nwlinkedidentc{Formula.serialize}{NW3w0iRO-3H9tUA-1} (\nwlinkedidentc{Formula.subst}{NW3w0iRO-1jtQRV-1} x t phi))))
  else if not(\nwlinkedidentc{Formula.free_for}{NW3w0iRO-ZZefY-1} t x phi)
  then raise Fail "exists_intro: term is not free for var in formula"
  else Gamma |- (\nwlinkedidentc{Formula.mk_exists}{NW3w0iRO-3SlqND-4}(x, phi))
\nwindexdefn{\nwixident{Thm.exists{\_}intro}}{Thm.exists:unintro}{NW4ICsJQ-2zKkPU-A}\eatline
\nwused{\\{NW4ICsJQ-YIPBT-1}}\nwidentdefs{\\{{\nwixident{Thm.exists{\_}intro}}{Thm.exists:unintro}}}\nwidentuses{\\{{\nwixident{Formula.bv}}{Formula.bv}}\\{{\nwixident{Formula.eq}}{Formula.eq}}\\{{\nwixident{Formula.free{\_}for}}{Formula.free:unfor}}\\{{\nwixident{Formula.mk{\_}exists}}{Formula.mk:unexists}}\\{{\nwixident{Formula.serialize}}{Formula.serialize}}\\{{\nwixident{Formula.subst}}{Formula.subst}}\\{{\nwixident{Formula.t}}{Formula.t}}\\{{\nwixident{Term.eq}}{Term.eq}}\\{{\nwixident{Term.have{\_}same{\_}sort}}{Term.have:unsame:unsort}}\\{{\nwixident{Term.t}}{Term.t}}}\nwindexuse{\nwixident{Formula.bv}}{Formula.bv}{NW4ICsJQ-2zKkPU-A}\nwindexuse{\nwixident{Formula.eq}}{Formula.eq}{NW4ICsJQ-2zKkPU-A}\nwindexuse{\nwixident{Formula.free{\_}for}}{Formula.free:unfor}{NW4ICsJQ-2zKkPU-A}\nwindexuse{\nwixident{Formula.mk{\_}exists}}{Formula.mk:unexists}{NW4ICsJQ-2zKkPU-A}\nwindexuse{\nwixident{Formula.serialize}}{Formula.serialize}{NW4ICsJQ-2zKkPU-A}\nwindexuse{\nwixident{Formula.subst}}{Formula.subst}{NW4ICsJQ-2zKkPU-A}\nwindexuse{\nwixident{Formula.t}}{Formula.t}{NW4ICsJQ-2zKkPU-A}\nwindexuse{\nwixident{Term.eq}}{Term.eq}{NW4ICsJQ-2zKkPU-A}\nwindexuse{\nwixident{Term.have{\_}same{\_}sort}}{Term.have:unsame:unsort}{NW4ICsJQ-2zKkPU-A}\nwindexuse{\nwixident{Term.t}}{Term.t}{NW4ICsJQ-2zKkPU-A}\nwendcode{}\nwbegindocs{42}\nwdocspar
\begin{node}\label{node:kernel:existential-elimination}
Existential elimination takes the form
\begin{equation*}
\vcenter{\infer[\exists\mbox{-elim}(x)]{\Gamma\vdash(\exists x\ldotp\phi)\implies\psi}{\Gamma\vdash\phi\implies\psi,&x\notin FV(\psi),&x\notin FV(\Gamma)}}
\end{equation*}
\end{node}

\nwenddocs{}\nwbegincode{43}\sublabel{NW4ICsJQ-1XQlFy-A}\nwmargintag{{\nwtagstyle{}\subpageref{NW4ICsJQ-1XQlFy-A}}}\moddef{Specification for inference rules~{\nwtagstyle{}\subpageref{NW4ICsJQ-1XQlFy-1}}}\plusendmoddef\nwstartdeflinemarkup\nwusesondefline{\\{NW4ICsJQ-2cY0B1-1}}\nwprevnextdefs{NW4ICsJQ-1XQlFy-9}{\relax}\nwenddeflinemarkup
val exists_elim : \nwlinkedidentc{Term.t}{NW1rTmOJ-4LsRWC-1} -> t -> t;
\nwused{\\{NW4ICsJQ-2cY0B1-1}}\nwidentuses{\\{{\nwixident{Term.t}}{Term.t}}}\nwindexuse{\nwixident{Term.t}}{Term.t}{NW4ICsJQ-1XQlFy-A}\nwendcode{}\nwbegindocs{44}\nwdocspar
\nwenddocs{}\nwbegincode{45}\sublabel{NW4ICsJQ-2zKkPU-B}\nwmargintag{{\nwtagstyle{}\subpageref{NW4ICsJQ-2zKkPU-B}}}\moddef{Inference rules~{\nwtagstyle{}\subpageref{NW4ICsJQ-2zKkPU-1}}}\plusendmoddef\nwstartdeflinemarkup\nwusesondefline{\\{NW4ICsJQ-YIPBT-1}}\nwprevnextdefs{NW4ICsJQ-2zKkPU-A}{\relax}\nwenddeflinemarkup
fun exists_elim (x : \nwlinkedidentc{Term.t}{NW1rTmOJ-4LsRWC-1}) (Gamma |- A_imp_B) =
  if not(Term.is_var x)
  then raise Fail "exists_elim: x is not a variable"
  else if not(\nwlinkedidentc{Formula.is_imp}{NW3w0iRO-3NyJLv-3} A_imp_B)
  then raise Fail "exists_elim: theorem is not an implication"
  else let val (A,B) = \nwlinkedidentc{Formula.dest_imp}{NW3w0iRO-1Yt1Jl-3} A_imp_B;
           fun free_in fm = List.exists (fn y => \nwlinkedidentc{Term.eq}{NW1rTmOJ-3810sY-1} x y) (\nwlinkedidentc{Formula.fv}{NW3w0iRO-jRgYZ-2} fm);
       in if free_in B
          then raise Fail "exists_elim: x is free in conclusion of theorem"
          else if List.exists free_in Gamma
          then raise Fail "exists_elim: x is free in hypotheses"
          else Gamma |- (\nwlinkedidentc{Formula.mk_imp}{NW3w0iRO-3SlqND-3}(\nwlinkedidentc{Formula.mk_exists}{NW3w0iRO-3SlqND-4}(x, A), B))
       end;
\nwindexdefn{\nwixident{Thm.exists{\_}elim}}{Thm.exists:unelim}{NW4ICsJQ-2zKkPU-B}\eatline
\nwused{\\{NW4ICsJQ-YIPBT-1}}\nwidentdefs{\\{{\nwixident{Thm.exists{\_}elim}}{Thm.exists:unelim}}}\nwidentuses{\\{{\nwixident{Formula.dest{\_}imp}}{Formula.dest:unimp}}\\{{\nwixident{Formula.fv}}{Formula.fv}}\\{{\nwixident{Formula.is{\_}imp}}{Formula.is:unimp}}\\{{\nwixident{Formula.mk{\_}exists}}{Formula.mk:unexists}}\\{{\nwixident{Formula.mk{\_}imp}}{Formula.mk:unimp}}\\{{\nwixident{Term.eq}}{Term.eq}}\\{{\nwixident{Term.t}}{Term.t}}}\nwindexuse{\nwixident{Formula.dest{\_}imp}}{Formula.dest:unimp}{NW4ICsJQ-2zKkPU-B}\nwindexuse{\nwixident{Formula.fv}}{Formula.fv}{NW4ICsJQ-2zKkPU-B}\nwindexuse{\nwixident{Formula.is{\_}imp}}{Formula.is:unimp}{NW4ICsJQ-2zKkPU-B}\nwindexuse{\nwixident{Formula.mk{\_}exists}}{Formula.mk:unexists}{NW4ICsJQ-2zKkPU-B}\nwindexuse{\nwixident{Formula.mk{\_}imp}}{Formula.mk:unimp}{NW4ICsJQ-2zKkPU-B}\nwindexuse{\nwixident{Term.eq}}{Term.eq}{NW4ICsJQ-2zKkPU-B}\nwindexuse{\nwixident{Term.t}}{Term.t}{NW4ICsJQ-2zKkPU-B}\nwendcode{}\nwbegindocs{46}\nwdocspar
\begin{xca}
Prove the inference rules provided by the kernel are \textbf{complete}, or find
a counter-example (a formula $\phi$ such that $\models\phi$ but $\nvdash\phi$).
\end{xca}

\begin{xca}
Prove the inference rules provided by the kernel are \textbf{sound}, or find
a counter-example (a formula $\phi$ such that $\nvDash\phi$ but $\vdash\phi$).
\end{xca}

\section*{Axioms}

\begin{node}
  The initial axioms (written in ``curried'' form with implication
  taken to be right associative, so $A\implies B\implies C$ means
  $A\implies(B\implies C)$) taken directly from
  Feferman's paper~\cite[\S16]{feferman1989finitary}.
\end{node}

\begin{node}
We want to avoid the tediousness of invoking $\forall$-elimination,
which means we want to provide axiom schemes. If the user
\emph{really} needs these universally quantified, they can use the
$\forall$-introduction rule.
\begin{enumerate}
  \item $\forall x_{1},x_{2}:\textit{Ind}\ldotp (x_{1},x_{2})\neq0$
  \item $\forall x_{1},x_{2}:\textit{Ind}\ldotp P_{1}(x_{1},x_{2})=x_{1}$
  \item $\forall x_{1},x_{2}:\textit{Ind}\ldotp P_{2}(x_{1},x_{2})=x_{2}$
\end{enumerate}
Care must be taken because {\Tt{}proj1\nwendquote} and {\Tt{}proj2\nwendquote} are ``eager''
constructors, so we need to use {\Tt{}mk{\_}proj1\nwendquote} and {\Tt{}mk{\_}proj2\nwendquote} to state
the intended axioms.
\end{node}

\nwenddocs{}\nwbegincode{47}\sublabel{NW4ICsJQ-hYCIS-1}\nwmargintag{{\nwtagstyle{}\subpageref{NW4ICsJQ-hYCIS-1}}}\moddef{Specification for $FS_{0}$ axioms~{\nwtagstyle{}\subpageref{NW4ICsJQ-hYCIS-1}}}\endmoddef\nwstartdeflinemarkup\nwusesondefline{\\{NW4ICsJQ-2cY0B1-1}}\nwprevnextdefs{\relax}{NW4ICsJQ-hYCIS-2}\nwenddeflinemarkup
val axiom_pair : \nwlinkedidentc{Term.t}{NW1rTmOJ-4LsRWC-1} -> \nwlinkedidentc{Term.t}{NW1rTmOJ-4LsRWC-1} -> t;
val axiom_proj1 : \nwlinkedidentc{Term.t}{NW1rTmOJ-4LsRWC-1} -> \nwlinkedidentc{Term.t}{NW1rTmOJ-4LsRWC-1} -> t;
val axiom_proj2 : \nwlinkedidentc{Term.t}{NW1rTmOJ-4LsRWC-1} -> \nwlinkedidentc{Term.t}{NW1rTmOJ-4LsRWC-1} -> t;
\nwalsodefined{\\{NW4ICsJQ-hYCIS-2}\\{NW4ICsJQ-hYCIS-3}\\{NW4ICsJQ-hYCIS-4}\\{NW4ICsJQ-hYCIS-5}\\{NW4ICsJQ-hYCIS-6}\\{NW4ICsJQ-hYCIS-7}\\{NW4ICsJQ-hYCIS-8}\\{NW4ICsJQ-hYCIS-9}\\{NW4ICsJQ-hYCIS-A}\\{NW4ICsJQ-hYCIS-B}\\{NW4ICsJQ-hYCIS-C}\\{NW4ICsJQ-hYCIS-D}}\nwused{\\{NW4ICsJQ-2cY0B1-1}}\nwidentuses{\\{{\nwixident{Term.t}}{Term.t}}}\nwindexuse{\nwixident{Term.t}}{Term.t}{NW4ICsJQ-hYCIS-1}\nwendcode{}\nwbegindocs{48}\nwdocspar
\nwenddocs{}\nwbegincode{49}\sublabel{NW4ICsJQ-2aDdLu-1}\nwmargintag{{\nwtagstyle{}\subpageref{NW4ICsJQ-2aDdLu-1}}}\moddef{Axioms for $FS_{0}$~{\nwtagstyle{}\subpageref{NW4ICsJQ-2aDdLu-1}}}\endmoddef\nwstartdeflinemarkup\nwusesondefline{\\{NW4ICsJQ-YIPBT-1}}\nwprevnextdefs{\relax}{NW4ICsJQ-2aDdLu-2}\nwenddeflinemarkup
fun axiom_pair x1 x2 =
  mk_axiom(\nwlinkedidentc{FS0.not_equal}{NW2DqMzY-4OjF2N-1} (\nwlinkedidentc{FS0.pair}{NW2DqMzY-4SkDla-1} x1 x2) \nwlinkedidentc{FS0.zero}{NW2DqMzY-1NUSQs-1});

fun axiom_proj1 x1 x2 =
  mk_axiom(\nwlinkedidentc{FS0.equal}{NW2DqMzY-ocj48-1} (\nwlinkedidentc{FS0.mk_proj1}{NW2DqMzY-ilITB-1} (\nwlinkedidentc{FS0.pair}{NW2DqMzY-4SkDla-1} x1 x2)) x1);

fun axiom_proj2 x1 x2 =
  mk_axiom(\nwlinkedidentc{FS0.equal}{NW2DqMzY-ocj48-1} (\nwlinkedidentc{FS0.mk_proj2}{NW2DqMzY-HQiIh-1} (\nwlinkedidentc{FS0.pair}{NW2DqMzY-4SkDla-1} x1 x2)) x2);
\nwindexdefn{\nwixident{Thm.axiom{\_}pair}}{Thm.axiom:unpair}{NW4ICsJQ-2aDdLu-1}\nwindexdefn{\nwixident{Thm.axiom{\_}proj1}}{Thm.axiom:unproj1}{NW4ICsJQ-2aDdLu-1}\nwindexdefn{\nwixident{Thm.axiom{\_}proj2}}{Thm.axiom:unproj2}{NW4ICsJQ-2aDdLu-1}\eatline
\nwalsodefined{\\{NW4ICsJQ-2aDdLu-2}\\{NW4ICsJQ-2aDdLu-3}\\{NW4ICsJQ-2aDdLu-4}\\{NW4ICsJQ-2aDdLu-5}\\{NW4ICsJQ-2aDdLu-6}\\{NW4ICsJQ-2aDdLu-7}\\{NW4ICsJQ-2aDdLu-8}\\{NW4ICsJQ-2aDdLu-9}\\{NW4ICsJQ-2aDdLu-A}\\{NW4ICsJQ-2aDdLu-B}\\{NW4ICsJQ-2aDdLu-C}\\{NW4ICsJQ-2aDdLu-D}}\nwused{\\{NW4ICsJQ-YIPBT-1}}\nwidentdefs{\\{{\nwixident{Thm.axiom{\_}pair}}{Thm.axiom:unpair}}\\{{\nwixident{Thm.axiom{\_}proj1}}{Thm.axiom:unproj1}}\\{{\nwixident{Thm.axiom{\_}proj2}}{Thm.axiom:unproj2}}}\nwidentuses{\\{{\nwixident{FS0.equal}}{FS0.equal}}\\{{\nwixident{FS0.mk{\_}proj1}}{FS0.mk:unproj1}}\\{{\nwixident{FS0.mk{\_}proj2}}{FS0.mk:unproj2}}\\{{\nwixident{FS0.not{\_}equal}}{FS0.not:unequal}}\\{{\nwixident{FS0.pair}}{FS0.pair}}\\{{\nwixident{FS0.zero}}{FS0.zero}}}\nwindexuse{\nwixident{FS0.equal}}{FS0.equal}{NW4ICsJQ-2aDdLu-1}\nwindexuse{\nwixident{FS0.mk{\_}proj1}}{FS0.mk:unproj1}{NW4ICsJQ-2aDdLu-1}\nwindexuse{\nwixident{FS0.mk{\_}proj2}}{FS0.mk:unproj2}{NW4ICsJQ-2aDdLu-1}\nwindexuse{\nwixident{FS0.not{\_}equal}}{FS0.not:unequal}{NW4ICsJQ-2aDdLu-1}\nwindexuse{\nwixident{FS0.pair}}{FS0.pair}{NW4ICsJQ-2aDdLu-1}\nwindexuse{\nwixident{FS0.zero}}{FS0.zero}{NW4ICsJQ-2aDdLu-1}\nwendcode{}\nwbegindocs{50}\nwdocspar
\begin{node}
The axioms for functions:
\begin{enumerate}
\item $Ix=x$
\item $K(a)x=a$
\item $x_{1}=x_{2}\implies D(x_{1},x_{2},y_{1},y_{2})=y_{1}$
\item $x_{1}\neq x_{2}\implies D(x_{1},x_{2},y_{1},y_{2})=y_{2}$
\item $\neg(\exists x_{1},x_{2},y_{1},y_{2}\ldotp u=(x_{1},x_{2},y_{1},y_{2}))\implies Du=0$
\end{enumerate}
Funny enough, Feferman originally intended $D$ to be ``definition by
cases'', but later~\cite[\S5.3]{feferman1982inductively} its name appeared to change to ``decide''.
\end{node}


\nwenddocs{}\nwbegincode{51}\sublabel{NW4ICsJQ-hYCIS-2}\nwmargintag{{\nwtagstyle{}\subpageref{NW4ICsJQ-hYCIS-2}}}\moddef{Specification for $FS_{0}$ axioms~{\nwtagstyle{}\subpageref{NW4ICsJQ-hYCIS-1}}}\plusendmoddef\nwstartdeflinemarkup\nwusesondefline{\\{NW4ICsJQ-2cY0B1-1}}\nwprevnextdefs{NW4ICsJQ-hYCIS-1}{NW4ICsJQ-hYCIS-3}\nwenddeflinemarkup
val axiom_id : \nwlinkedidentc{Term.t}{NW1rTmOJ-4LsRWC-1} -> t;
val axiom_const : \nwlinkedidentc{Term.t}{NW1rTmOJ-4LsRWC-1} -> \nwlinkedidentc{Term.t}{NW1rTmOJ-4LsRWC-1} -> t;
val axiom_cond_true : \nwlinkedidentc{Term.t}{NW1rTmOJ-4LsRWC-1} -> \nwlinkedidentc{Term.t}{NW1rTmOJ-4LsRWC-1} -> \nwlinkedidentc{Term.t}{NW1rTmOJ-4LsRWC-1} -> \nwlinkedidentc{Term.t}{NW1rTmOJ-4LsRWC-1} -> t;
val axiom_cond_false : \nwlinkedidentc{Term.t}{NW1rTmOJ-4LsRWC-1} -> \nwlinkedidentc{Term.t}{NW1rTmOJ-4LsRWC-1} -> \nwlinkedidentc{Term.t}{NW1rTmOJ-4LsRWC-1} -> \nwlinkedidentc{Term.t}{NW1rTmOJ-4LsRWC-1} -> t;
val axiom_cond_err : \nwlinkedidentc{Term.t}{NW1rTmOJ-4LsRWC-1} -> \nwlinkedidentc{Term.t}{NW1rTmOJ-4LsRWC-1} -> \nwlinkedidentc{Term.t}{NW1rTmOJ-4LsRWC-1} -> \nwlinkedidentc{Term.t}{NW1rTmOJ-4LsRWC-1} -> t;
\nwused{\\{NW4ICsJQ-2cY0B1-1}}\nwidentuses{\\{{\nwixident{Term.t}}{Term.t}}}\nwindexuse{\nwixident{Term.t}}{Term.t}{NW4ICsJQ-hYCIS-2}\nwendcode{}\nwbegindocs{52}\nwdocspar
\nwenddocs{}\nwbegincode{53}\sublabel{NW4ICsJQ-2aDdLu-2}\nwmargintag{{\nwtagstyle{}\subpageref{NW4ICsJQ-2aDdLu-2}}}\moddef{Axioms for $FS_{0}$~{\nwtagstyle{}\subpageref{NW4ICsJQ-2aDdLu-1}}}\plusendmoddef\nwstartdeflinemarkup\nwusesondefline{\\{NW4ICsJQ-YIPBT-1}}\nwprevnextdefs{NW4ICsJQ-2aDdLu-1}{NW4ICsJQ-2aDdLu-3}\nwenddeflinemarkup
fun axiom_id x = mk_axiom(\nwlinkedidentc{FS0.equal}{NW2DqMzY-ocj48-1} (\nwlinkedidentc{FS0.mk_id}{NW2DqMzY-D9caK-1} x) x);

fun axiom_const a x = mk_axiom(\nwlinkedidentc{FS0.equal}{NW2DqMzY-ocj48-1} (\nwlinkedidentc{FS0.app}{NW2DqMzY-2oGoBx-1} (\nwlinkedidentc{FS0.const}{NW2DqMzY-22MFug-1} a) x) a);

local
  fun quad x1 x2 y1 y2 = \nwlinkedidentc{FS0.pair}{NW2DqMzY-4SkDla-1} x1 (\nwlinkedidentc{FS0.pair}{NW2DqMzY-4SkDla-1} x2 (\nwlinkedidentc{FS0.pair}{NW2DqMzY-4SkDla-1} y1 y2));
in
fun axiom_cond_true x1 x2 y1 y2 =
  mk_axiom(\nwlinkedidentc{Formula.mk_imp}{NW3w0iRO-3SlqND-3}(\nwlinkedidentc{FS0.equal}{NW2DqMzY-ocj48-1} x1 x2,
                              \nwlinkedidentc{FS0.equal}{NW2DqMzY-ocj48-1} (\nwlinkedidentc{FS0.mk_if_eq}{NW2DqMzY-2vOgWm-1} (quad x1 x2 y1 y2))
                                        y1));
fun axiom_cond_false x1 x2 y1 y2 =
  mk_axiom(\nwlinkedidentc{Formula.mk_imp}{NW3w0iRO-3SlqND-3}(\nwlinkedidentc{FS0.not_equal}{NW2DqMzY-4OjF2N-1} x1 x2,
                              \nwlinkedidentc{FS0.equal}{NW2DqMzY-ocj48-1} (\nwlinkedidentc{FS0.mk_if_eq}{NW2DqMzY-2vOgWm-1} (quad x1 x2 y1 y2))
                                        y2));

fun axiom_cond_err x1 x2 y1 y2 =
  let
    val u = \nwlinkedidentc{Term.fresh_var}{NW1rTmOJ-4Ykq9M-1} "u" \nwlinkedidentc{Sort.IND}{NW1rTmOJ-3ZqggC-1} ((\nwlinkedidentc{Term.fv}{NW1rTmOJ-2nRtky-1} x1) @
                                         (\nwlinkedidentc{Term.fv}{NW1rTmOJ-2nRtky-1} x2) @
                                         (\nwlinkedidentc{Term.fv}{NW1rTmOJ-2nRtky-1} y1) @
                                         (\nwlinkedidentc{Term.fv}{NW1rTmOJ-2nRtky-1} y2));
    val ex = List.foldr \nwlinkedidentc{Formula.mk_exists}{NW3w0iRO-3SlqND-4}
                        (\nwlinkedidentc{FS0.equal}{NW2DqMzY-ocj48-1} u (quad x1 x2 y1 y2))
                        [x1,x2,y1,y2];
  in
    mk_axiom(\nwlinkedidentc{Formula.mk_imp}{NW3w0iRO-3SlqND-3}(\nwlinkedidentc{Formula.mk_not}{NW3w0iRO-3SlqND-2}(ex),
                            \nwlinkedidentc{FS0.equal}{NW2DqMzY-ocj48-1} (\nwlinkedidentc{FS0.mk_if_eq}{NW2DqMzY-2vOgWm-1} u)
                                      \nwlinkedidentc{FS0.zero}{NW2DqMzY-1NUSQs-1}))
  end;
end;
\nwindexdefn{\nwixident{Thm.axiom{\_}id}}{Thm.axiom:unid}{NW4ICsJQ-2aDdLu-2}\nwindexdefn{\nwixident{Thm.axiom{\_}const}}{Thm.axiom:unconst}{NW4ICsJQ-2aDdLu-2}\nwindexdefn{\nwixident{Thm.axiom{\_}cond{\_}true}}{Thm.axiom:uncond:untrue}{NW4ICsJQ-2aDdLu-2}\nwindexdefn{\nwixident{Thm.axiom{\_}cond{\_}false}}{Thm.axiom:uncond:unfalse}{NW4ICsJQ-2aDdLu-2}\nwindexdefn{\nwixident{Thm.axiom{\_}cond{\_}err}}{Thm.axiom:uncond:unerr}{NW4ICsJQ-2aDdLu-2}\eatline
\nwused{\\{NW4ICsJQ-YIPBT-1}}\nwidentdefs{\\{{\nwixident{Thm.axiom{\_}cond{\_}err}}{Thm.axiom:uncond:unerr}}\\{{\nwixident{Thm.axiom{\_}cond{\_}false}}{Thm.axiom:uncond:unfalse}}\\{{\nwixident{Thm.axiom{\_}cond{\_}true}}{Thm.axiom:uncond:untrue}}\\{{\nwixident{Thm.axiom{\_}const}}{Thm.axiom:unconst}}\\{{\nwixident{Thm.axiom{\_}id}}{Thm.axiom:unid}}}\nwidentuses{\\{{\nwixident{Formula.mk{\_}exists}}{Formula.mk:unexists}}\\{{\nwixident{Formula.mk{\_}imp}}{Formula.mk:unimp}}\\{{\nwixident{Formula.mk{\_}not}}{Formula.mk:unnot}}\\{{\nwixident{FS0.app}}{FS0.app}}\\{{\nwixident{FS0.const}}{FS0.const}}\\{{\nwixident{FS0.equal}}{FS0.equal}}\\{{\nwixident{FS0.mk{\_}id}}{FS0.mk:unid}}\\{{\nwixident{FS0.mk{\_}if{\_}eq}}{FS0.mk:unif:uneq}}\\{{\nwixident{FS0.not{\_}equal}}{FS0.not:unequal}}\\{{\nwixident{FS0.pair}}{FS0.pair}}\\{{\nwixident{FS0.zero}}{FS0.zero}}\\{{\nwixident{Sort.IND}}{Sort.IND}}\\{{\nwixident{Term.fresh{\_}var}}{Term.fresh:unvar}}\\{{\nwixident{Term.fv}}{Term.fv}}}\nwindexuse{\nwixident{Formula.mk{\_}exists}}{Formula.mk:unexists}{NW4ICsJQ-2aDdLu-2}\nwindexuse{\nwixident{Formula.mk{\_}imp}}{Formula.mk:unimp}{NW4ICsJQ-2aDdLu-2}\nwindexuse{\nwixident{Formula.mk{\_}not}}{Formula.mk:unnot}{NW4ICsJQ-2aDdLu-2}\nwindexuse{\nwixident{FS0.app}}{FS0.app}{NW4ICsJQ-2aDdLu-2}\nwindexuse{\nwixident{FS0.const}}{FS0.const}{NW4ICsJQ-2aDdLu-2}\nwindexuse{\nwixident{FS0.equal}}{FS0.equal}{NW4ICsJQ-2aDdLu-2}\nwindexuse{\nwixident{FS0.mk{\_}id}}{FS0.mk:unid}{NW4ICsJQ-2aDdLu-2}\nwindexuse{\nwixident{FS0.mk{\_}if{\_}eq}}{FS0.mk:unif:uneq}{NW4ICsJQ-2aDdLu-2}\nwindexuse{\nwixident{FS0.not{\_}equal}}{FS0.not:unequal}{NW4ICsJQ-2aDdLu-2}\nwindexuse{\nwixident{FS0.pair}}{FS0.pair}{NW4ICsJQ-2aDdLu-2}\nwindexuse{\nwixident{FS0.zero}}{FS0.zero}{NW4ICsJQ-2aDdLu-2}\nwindexuse{\nwixident{Sort.IND}}{Sort.IND}{NW4ICsJQ-2aDdLu-2}\nwindexuse{\nwixident{Term.fresh{\_}var}}{Term.fresh:unvar}{NW4ICsJQ-2aDdLu-2}\nwindexuse{\nwixident{Term.fv}}{Term.fv}{NW4ICsJQ-2aDdLu-2}\nwendcode{}\nwbegindocs{54}\nwdocspar
\begin{node}
The compound function axioms.
\begin{enumerate}
\item $\forall f,g,h:\textit{Fun}\ldotp h=\mathcal{P}(f,g)\implies(\forall x:\textit{Ind}\ldotp hx=(fx,gx))$
\item $\forall f,g,h:\textit{Fun}\ldotp h=\mathcal{C}(f,g)\implies(\forall x:\textit{Ind}\ldotp hx=f(gx))$
\end{enumerate}
\end{node}

\nwenddocs{}\nwbegincode{55}\sublabel{NW4ICsJQ-hYCIS-3}\nwmargintag{{\nwtagstyle{}\subpageref{NW4ICsJQ-hYCIS-3}}}\moddef{Specification for $FS_{0}$ axioms~{\nwtagstyle{}\subpageref{NW4ICsJQ-hYCIS-1}}}\plusendmoddef\nwstartdeflinemarkup\nwusesondefline{\\{NW4ICsJQ-2cY0B1-1}}\nwprevnextdefs{NW4ICsJQ-hYCIS-2}{NW4ICsJQ-hYCIS-4}\nwenddeflinemarkup
val axiom_juxt : \nwlinkedidentc{Term.t}{NW1rTmOJ-4LsRWC-1} -> \nwlinkedidentc{Term.t}{NW1rTmOJ-4LsRWC-1} -> t;
val axiom_comp : \nwlinkedidentc{Term.t}{NW1rTmOJ-4LsRWC-1} -> \nwlinkedidentc{Term.t}{NW1rTmOJ-4LsRWC-1} -> t;
\nwused{\\{NW4ICsJQ-2cY0B1-1}}\nwidentuses{\\{{\nwixident{Term.t}}{Term.t}}}\nwindexuse{\nwixident{Term.t}}{Term.t}{NW4ICsJQ-hYCIS-3}\nwendcode{}\nwbegindocs{56}\nwdocspar
\nwenddocs{}\nwbegincode{57}\sublabel{NW4ICsJQ-2aDdLu-3}\nwmargintag{{\nwtagstyle{}\subpageref{NW4ICsJQ-2aDdLu-3}}}\moddef{Axioms for $FS_{0}$~{\nwtagstyle{}\subpageref{NW4ICsJQ-2aDdLu-1}}}\plusendmoddef\nwstartdeflinemarkup\nwusesondefline{\\{NW4ICsJQ-YIPBT-1}}\nwprevnextdefs{NW4ICsJQ-2aDdLu-2}{NW4ICsJQ-2aDdLu-4}\nwenddeflinemarkup
fun axiom_juxt f g =
  let val x = \nwlinkedidentc{Term.fresh_var}{NW1rTmOJ-4Ykq9M-1} "x" \nwlinkedidentc{Sort.IND}{NW1rTmOJ-3ZqggC-1} ((\nwlinkedidentc{Term.fv}{NW1rTmOJ-2nRtky-1} f) @ (\nwlinkedidentc{Term.fv}{NW1rTmOJ-2nRtky-1} g))
  in mk_axiom(\nwlinkedidentc{Formula.mk_forall}{NW3w0iRO-3SlqND-4}(x, \nwlinkedidentc{FS0.equal}{NW2DqMzY-ocj48-1} (\nwlinkedidentc{FS0.app}{NW2DqMzY-2oGoBx-1} (\nwlinkedidentc{FS0.juxt}{NW2DqMzY-1lMRH2-1} f g) x)
                                             (\nwlinkedidentc{FS0.pair}{NW2DqMzY-4SkDla-1} (\nwlinkedidentc{FS0.app}{NW2DqMzY-2oGoBx-1} f x)
                                                       (\nwlinkedidentc{FS0.app}{NW2DqMzY-2oGoBx-1} g x))))
  end;

fun axiom_comp f g =
  let val x = \nwlinkedidentc{Term.fresh_var}{NW1rTmOJ-4Ykq9M-1} "x" \nwlinkedidentc{Sort.IND}{NW1rTmOJ-3ZqggC-1} ((\nwlinkedidentc{Term.fv}{NW1rTmOJ-2nRtky-1} f) @ (\nwlinkedidentc{Term.fv}{NW1rTmOJ-2nRtky-1} g))
  in mk_axiom(\nwlinkedidentc{Formula.mk_forall}{NW3w0iRO-3SlqND-4}(x, \nwlinkedidentc{FS0.equal}{NW2DqMzY-ocj48-1} (\nwlinkedidentc{FS0.app}{NW2DqMzY-2oGoBx-1} (\nwlinkedidentc{FS0.comp}{NW2DqMzY-1dWhhL-1} f g) x)
                                             (\nwlinkedidentc{FS0.app}{NW2DqMzY-2oGoBx-1} f (\nwlinkedidentc{FS0.app}{NW2DqMzY-2oGoBx-1} g x))))
  end;
\nwindexdefn{\nwixident{Thm.axiom{\_}juxt}}{Thm.axiom:unjuxt}{NW4ICsJQ-2aDdLu-3}\nwindexdefn{\nwixident{Thm.axiom{\_}comp}}{Thm.axiom:uncomp}{NW4ICsJQ-2aDdLu-3}\eatline
\nwused{\\{NW4ICsJQ-YIPBT-1}}\nwidentdefs{\\{{\nwixident{Thm.axiom{\_}comp}}{Thm.axiom:uncomp}}\\{{\nwixident{Thm.axiom{\_}juxt}}{Thm.axiom:unjuxt}}}\nwidentuses{\\{{\nwixident{Formula.mk{\_}forall}}{Formula.mk:unforall}}\\{{\nwixident{FS0.app}}{FS0.app}}\\{{\nwixident{FS0.comp}}{FS0.comp}}\\{{\nwixident{FS0.equal}}{FS0.equal}}\\{{\nwixident{FS0.juxt}}{FS0.juxt}}\\{{\nwixident{FS0.pair}}{FS0.pair}}\\{{\nwixident{Sort.IND}}{Sort.IND}}\\{{\nwixident{Term.fresh{\_}var}}{Term.fresh:unvar}}\\{{\nwixident{Term.fv}}{Term.fv}}}\nwindexuse{\nwixident{Formula.mk{\_}forall}}{Formula.mk:unforall}{NW4ICsJQ-2aDdLu-3}\nwindexuse{\nwixident{FS0.app}}{FS0.app}{NW4ICsJQ-2aDdLu-3}\nwindexuse{\nwixident{FS0.comp}}{FS0.comp}{NW4ICsJQ-2aDdLu-3}\nwindexuse{\nwixident{FS0.equal}}{FS0.equal}{NW4ICsJQ-2aDdLu-3}\nwindexuse{\nwixident{FS0.juxt}}{FS0.juxt}{NW4ICsJQ-2aDdLu-3}\nwindexuse{\nwixident{FS0.pair}}{FS0.pair}{NW4ICsJQ-2aDdLu-3}\nwindexuse{\nwixident{Sort.IND}}{Sort.IND}{NW4ICsJQ-2aDdLu-3}\nwindexuse{\nwixident{Term.fresh{\_}var}}{Term.fresh:unvar}{NW4ICsJQ-2aDdLu-3}\nwindexuse{\nwixident{Term.fv}}{Term.fv}{NW4ICsJQ-2aDdLu-3}\nwendcode{}\nwbegindocs{58}\nwdocspar
\begin{node}
Primitive recursion axioms.
\begin{enumerate}
  \item $\forall f,g:\textit{Fun}\ldotp \mathcal{R}(f,g)0=0$
  \item $\forall f,g:\textit{Fun}\ldotp \forall x:\textit{Ind}\ldotp \mathcal{R}(f,g)(x,0)=fx$
  \item $\forall f,g:\textit{Fun}\ldotp\forall x,y,z:\textit{Ind}\ldotp \mathcal{R}(f,g)(x,(y,z))=g(x,y,z,\mathcal{R}(f,g)(x,y),\mathcal{R}(f,g)(x,z)))$
\end{enumerate}
More generally, Feferman~\cite[\S5.2]{feferman1982inductively} shows
us how to extend $FS_{0}$ to support a [nonempty] set of
Urelements $U$, then we need to change the first axiom to something
like $\mathcal{R}(f_{0},f_{1},g)(x,u)=f_{1}(x,u)=f_{0}u$ and
$\mathcal{R}(f_{0},f_{1},g)(x,u)=f_{1}(x,u)$ for all $u\in U$, and
$\mathcal{R}(f_{0},f_{1},g)(x,(y,z))=g(x,y,z,\mathcal{R}(f_{0},f_{1},g)(x,y),\mathcal{R}(f_{0},f_{1},g)(x,z))$. In
other words, we change the $\mathcal{R}(f,g)0=0$ clause to a more
general clause.
\end{node}

\nwenddocs{}\nwbegincode{59}\sublabel{NW4ICsJQ-hYCIS-4}\nwmargintag{{\nwtagstyle{}\subpageref{NW4ICsJQ-hYCIS-4}}}\moddef{Specification for $FS_{0}$ axioms~{\nwtagstyle{}\subpageref{NW4ICsJQ-hYCIS-1}}}\plusendmoddef\nwstartdeflinemarkup\nwusesondefline{\\{NW4ICsJQ-2cY0B1-1}}\nwprevnextdefs{NW4ICsJQ-hYCIS-3}{NW4ICsJQ-hYCIS-5}\nwenddeflinemarkup
val axiom_recur_zero : \nwlinkedidentc{Term.t}{NW1rTmOJ-4LsRWC-1} -> \nwlinkedidentc{Term.t}{NW1rTmOJ-4LsRWC-1} -> t;
val axiom_recur_base : \nwlinkedidentc{Term.t}{NW1rTmOJ-4LsRWC-1} -> \nwlinkedidentc{Term.t}{NW1rTmOJ-4LsRWC-1} -> t;
val axiom_recur_rec : \nwlinkedidentc{Term.t}{NW1rTmOJ-4LsRWC-1} -> \nwlinkedidentc{Term.t}{NW1rTmOJ-4LsRWC-1} -> t;
\nwused{\\{NW4ICsJQ-2cY0B1-1}}\nwidentuses{\\{{\nwixident{Term.t}}{Term.t}}}\nwindexuse{\nwixident{Term.t}}{Term.t}{NW4ICsJQ-hYCIS-4}\nwendcode{}\nwbegindocs{60}\nwdocspar
\nwenddocs{}\nwbegincode{61}\sublabel{NW4ICsJQ-2aDdLu-4}\nwmargintag{{\nwtagstyle{}\subpageref{NW4ICsJQ-2aDdLu-4}}}\moddef{Axioms for $FS_{0}$~{\nwtagstyle{}\subpageref{NW4ICsJQ-2aDdLu-1}}}\plusendmoddef\nwstartdeflinemarkup\nwusesondefline{\\{NW4ICsJQ-YIPBT-1}}\nwprevnextdefs{NW4ICsJQ-2aDdLu-3}{NW4ICsJQ-2aDdLu-5}\nwenddeflinemarkup
fun axiom_recur_zero f g =
  mk_axiom(\nwlinkedidentc{FS0.equal}{NW2DqMzY-ocj48-1} (\nwlinkedidentc{FS0.app}{NW2DqMzY-2oGoBx-1} (\nwlinkedidentc{FS0.recur}{NW2DqMzY-2fxlqE-1} f g) \nwlinkedidentc{FS0.zero}{NW2DqMzY-1NUSQs-1})
                     \nwlinkedidentc{FS0.zero}{NW2DqMzY-1NUSQs-1});

fun axiom_recur_base f g =
  let val x = \nwlinkedidentc{Term.fresh_var}{NW1rTmOJ-4Ykq9M-1} "x" \nwlinkedidentc{Sort.IND}{NW1rTmOJ-3ZqggC-1} ((\nwlinkedidentc{Term.fv}{NW1rTmOJ-2nRtky-1} f) @ (\nwlinkedidentc{Term.fv}{NW1rTmOJ-2nRtky-1} g))
  in mk_axiom(\nwlinkedidentc{Formula.mk_forall}{NW3w0iRO-3SlqND-4}(x,
                                \nwlinkedidentc{FS0.equal}{NW2DqMzY-ocj48-1} (\nwlinkedidentc{FS0.app}{NW2DqMzY-2oGoBx-1} (\nwlinkedidentc{FS0.recur}{NW2DqMzY-2fxlqE-1} f g)
                                                   (\nwlinkedidentc{FS0.pair}{NW2DqMzY-4SkDla-1} x \nwlinkedidentc{FS0.zero}{NW2DqMzY-1NUSQs-1}))
                                          (\nwlinkedidentc{FS0.app}{NW2DqMzY-2oGoBx-1} f x)))
  end;

fun axiom_recur_rec f g =
  let
    val vars = (\nwlinkedidentc{Term.fv}{NW1rTmOJ-2nRtky-1} f) @ (\nwlinkedidentc{Term.fv}{NW1rTmOJ-2nRtky-1} g);
    fun var x = \nwlinkedidentc{Term.fresh_var}{NW1rTmOJ-4Ykq9M-1} x \nwlinkedidentc{Sort.IND}{NW1rTmOJ-3ZqggC-1} vars;
    val x = var "x";
    val y = var "y";
    val z = var "z";
    fun h x1 x2 = \nwlinkedidentc{FS0.app}{NW2DqMzY-2oGoBx-1} (\nwlinkedidentc{FS0.recur}{NW2DqMzY-2fxlqE-1} f g) (\nwlinkedidentc{FS0.pair}{NW2DqMzY-4SkDla-1} x1 x2);
    val lhs = h x (\nwlinkedidentc{FS0.pair}{NW2DqMzY-4SkDla-1} y z);
    val args = List.foldr (fn (a,b) => \nwlinkedidentc{FS0.pair}{NW2DqMzY-4SkDla-1} a b)
                          (h x z)
                          [x, y, z, (h x y)];
    val rhs = \nwlinkedidentc{FS0.app}{NW2DqMzY-2oGoBx-1} g args;
    val eq = \nwlinkedidentc{FS0.equal}{NW2DqMzY-ocj48-1} lhs rhs;
  in
    mk_axiom(List.foldr \nwlinkedidentc{Formula.mk_forall}{NW3w0iRO-3SlqND-4} eq [x,y,z])
  end;
\nwindexdefn{\nwixident{Thm.axiom{\_}recur{\_}zero}}{Thm.axiom:unrecur:unzero}{NW4ICsJQ-2aDdLu-4}\nwindexdefn{\nwixident{Thm.axiom{\_}recur{\_}base}}{Thm.axiom:unrecur:unbase}{NW4ICsJQ-2aDdLu-4}\nwindexdefn{\nwixident{Thm.axiom{\_}recur{\_}rec}}{Thm.axiom:unrecur:unrec}{NW4ICsJQ-2aDdLu-4}\eatline
\nwused{\\{NW4ICsJQ-YIPBT-1}}\nwidentdefs{\\{{\nwixident{Thm.axiom{\_}recur{\_}base}}{Thm.axiom:unrecur:unbase}}\\{{\nwixident{Thm.axiom{\_}recur{\_}rec}}{Thm.axiom:unrecur:unrec}}\\{{\nwixident{Thm.axiom{\_}recur{\_}zero}}{Thm.axiom:unrecur:unzero}}}\nwidentuses{\\{{\nwixident{Formula.mk{\_}forall}}{Formula.mk:unforall}}\\{{\nwixident{FS0.app}}{FS0.app}}\\{{\nwixident{FS0.equal}}{FS0.equal}}\\{{\nwixident{FS0.pair}}{FS0.pair}}\\{{\nwixident{FS0.recur}}{FS0.recur}}\\{{\nwixident{FS0.zero}}{FS0.zero}}\\{{\nwixident{Sort.IND}}{Sort.IND}}\\{{\nwixident{Term.fresh{\_}var}}{Term.fresh:unvar}}\\{{\nwixident{Term.fv}}{Term.fv}}}\nwindexuse{\nwixident{Formula.mk{\_}forall}}{Formula.mk:unforall}{NW4ICsJQ-2aDdLu-4}\nwindexuse{\nwixident{FS0.app}}{FS0.app}{NW4ICsJQ-2aDdLu-4}\nwindexuse{\nwixident{FS0.equal}}{FS0.equal}{NW4ICsJQ-2aDdLu-4}\nwindexuse{\nwixident{FS0.pair}}{FS0.pair}{NW4ICsJQ-2aDdLu-4}\nwindexuse{\nwixident{FS0.recur}}{FS0.recur}{NW4ICsJQ-2aDdLu-4}\nwindexuse{\nwixident{FS0.zero}}{FS0.zero}{NW4ICsJQ-2aDdLu-4}\nwindexuse{\nwixident{Sort.IND}}{Sort.IND}{NW4ICsJQ-2aDdLu-4}\nwindexuse{\nwixident{Term.fresh{\_}var}}{Term.fresh:unvar}{NW4ICsJQ-2aDdLu-4}\nwindexuse{\nwixident{Term.fv}}{Term.fv}{NW4ICsJQ-2aDdLu-4}\nwendcode{}\nwbegindocs{62}\nwdocspar
\begin{node}
The class axioms.
\begin{enumerate}
  \item $\forall x:\textit{Ind}\ldotp x\in\{0\}\iff x=0$ --- although
    this could be ``weakened'' to allow any singleton any we'd have
    $\forall x,y:\textit{Ind}\ldotp x\in\{y\}\iff x=y$
  \item $\forall A:\textit{Class}\ldotp\forall f:\textit{Fun}\ldotp\forall x:\textit{Ind}\ldotp x\in f^{-1}A\iff fx\in A$
  \item $\forall A,B : \textit{Class}\ldotp\forall x:\textit{Ind}\ldotp x\in A\cap B\iff x\in A\land x\in B$
  \item $\forall A,B : \textit{Class}\ldotp\forall x:\textit{Ind}\ldotp x\in A\cup B\iff x\in A\lor x\in B$
\end{enumerate}
\end{node}

\nwenddocs{}\nwbegincode{63}\sublabel{NW4ICsJQ-hYCIS-5}\nwmargintag{{\nwtagstyle{}\subpageref{NW4ICsJQ-hYCIS-5}}}\moddef{Specification for $FS_{0}$ axioms~{\nwtagstyle{}\subpageref{NW4ICsJQ-hYCIS-1}}}\plusendmoddef\nwstartdeflinemarkup\nwusesondefline{\\{NW4ICsJQ-2cY0B1-1}}\nwprevnextdefs{NW4ICsJQ-hYCIS-4}{NW4ICsJQ-hYCIS-6}\nwenddeflinemarkup
val axiom_singleton : \nwlinkedidentc{Term.t}{NW1rTmOJ-4LsRWC-1} -> \nwlinkedidentc{Term.t}{NW1rTmOJ-4LsRWC-1} -> t;
val axiom_preimage : \nwlinkedidentc{Term.t}{NW1rTmOJ-4LsRWC-1} -> \nwlinkedidentc{Term.t}{NW1rTmOJ-4LsRWC-1} -> t;
val axiom_union : \nwlinkedidentc{Term.t}{NW1rTmOJ-4LsRWC-1} -> \nwlinkedidentc{Term.t}{NW1rTmOJ-4LsRWC-1} -> t;
val axiom_intersection : \nwlinkedidentc{Term.t}{NW1rTmOJ-4LsRWC-1} -> \nwlinkedidentc{Term.t}{NW1rTmOJ-4LsRWC-1} -> t;
\nwused{\\{NW4ICsJQ-2cY0B1-1}}\nwidentuses{\\{{\nwixident{Term.t}}{Term.t}}}\nwindexuse{\nwixident{Term.t}}{Term.t}{NW4ICsJQ-hYCIS-5}\nwendcode{}\nwbegindocs{64}\nwdocspar
\nwenddocs{}\nwbegincode{65}\sublabel{NW4ICsJQ-2aDdLu-5}\nwmargintag{{\nwtagstyle{}\subpageref{NW4ICsJQ-2aDdLu-5}}}\moddef{Axioms for $FS_{0}$~{\nwtagstyle{}\subpageref{NW4ICsJQ-2aDdLu-1}}}\plusendmoddef\nwstartdeflinemarkup\nwusesondefline{\\{NW4ICsJQ-YIPBT-1}}\nwprevnextdefs{NW4ICsJQ-2aDdLu-4}{NW4ICsJQ-2aDdLu-6}\nwenddeflinemarkup
fun axiom_singleton obj elt =
  mk_axiom(\nwlinkedidentc{Formula.mk_iff}{NW3w0iRO-3SlqND-3}(\nwlinkedidentc{FS0.In}{NW2DqMzY-eTij6-1} obj (\nwlinkedidentc{FS0.singleton}{NW2DqMzY-14kolG-1} elt),
                          \nwlinkedidentc{FS0.equal}{NW2DqMzY-ocj48-1} obj elt));

fun axiom_preimage f C =
  let val x = \nwlinkedidentc{Term.fresh_var}{NW1rTmOJ-4Ykq9M-1} "x" \nwlinkedidentc{Sort.IND}{NW1rTmOJ-3ZqggC-1} ((\nwlinkedidentc{Term.fv}{NW1rTmOJ-2nRtky-1} f) @ (\nwlinkedidentc{Term.fv}{NW1rTmOJ-2nRtky-1} C))
  in mk_axiom(\nwlinkedidentc{Formula.mk_iff}{NW3w0iRO-3SlqND-3}(\nwlinkedidentc{FS0.In}{NW2DqMzY-eTij6-1} x (\nwlinkedidentc{FS0.preimage}{NW2DqMzY-31WKtJ-1} f C),
                             \nwlinkedidentc{FS0.In}{NW2DqMzY-eTij6-1} (\nwlinkedidentc{FS0.app}{NW2DqMzY-2oGoBx-1} f x) C))
  end;

fun axiom_union A B =
  let val x = \nwlinkedidentc{Term.fresh_var}{NW1rTmOJ-4Ykq9M-1} "x" \nwlinkedidentc{Sort.IND}{NW1rTmOJ-3ZqggC-1} ((\nwlinkedidentc{Term.fv}{NW1rTmOJ-2nRtky-1} A) @ (\nwlinkedidentc{Term.fv}{NW1rTmOJ-2nRtky-1} B))
  in mk_axiom(\nwlinkedidentc{Formula.mk_iff}{NW3w0iRO-3SlqND-3}(\nwlinkedidentc{FS0.In}{NW2DqMzY-eTij6-1} x (\nwlinkedidentc{FS0.union}{NW2DqMzY-22W7o7-1} A B),
                             \nwlinkedidentc{Formula.mk_or}{NW3w0iRO-3SlqND-3}(\nwlinkedidentc{FS0.In}{NW2DqMzY-eTij6-1} x A,
                                           \nwlinkedidentc{FS0.In}{NW2DqMzY-eTij6-1} x B)))
  end;

fun axiom_intersection A B =
  let val x = \nwlinkedidentc{Term.fresh_var}{NW1rTmOJ-4Ykq9M-1} "x" \nwlinkedidentc{Sort.IND}{NW1rTmOJ-3ZqggC-1} ((\nwlinkedidentc{Term.fv}{NW1rTmOJ-2nRtky-1} A) @ (\nwlinkedidentc{Term.fv}{NW1rTmOJ-2nRtky-1} B))
  in mk_axiom(\nwlinkedidentc{Formula.mk_iff}{NW3w0iRO-3SlqND-3}(\nwlinkedidentc{FS0.In}{NW2DqMzY-eTij6-1} x (\nwlinkedidentc{FS0.intersection}{NW2DqMzY-19bkPi-1} A B),
                             \nwlinkedidentc{Formula.mk_and}{NW3w0iRO-3SlqND-3}(\nwlinkedidentc{FS0.In}{NW2DqMzY-eTij6-1} x A,
                                            \nwlinkedidentc{FS0.In}{NW2DqMzY-eTij6-1} x B)))
  end;
\nwindexdefn{\nwixident{Thm.axiom{\_}singleton}}{Thm.axiom:unsingleton}{NW4ICsJQ-2aDdLu-5}\nwindexdefn{\nwixident{Thm.axiom{\_}preimage}}{Thm.axiom:unpreimage}{NW4ICsJQ-2aDdLu-5}\nwindexdefn{\nwixident{Thm.axiom{\_}union}}{Thm.axiom:ununion}{NW4ICsJQ-2aDdLu-5}\nwindexdefn{\nwixident{Thm.axiom{\_}intersection}}{Thm.axiom:unintersection}{NW4ICsJQ-2aDdLu-5}\eatline
\nwused{\\{NW4ICsJQ-YIPBT-1}}\nwidentdefs{\\{{\nwixident{Thm.axiom{\_}intersection}}{Thm.axiom:unintersection}}\\{{\nwixident{Thm.axiom{\_}preimage}}{Thm.axiom:unpreimage}}\\{{\nwixident{Thm.axiom{\_}singleton}}{Thm.axiom:unsingleton}}\\{{\nwixident{Thm.axiom{\_}union}}{Thm.axiom:ununion}}}\nwidentuses{\\{{\nwixident{Formula.mk{\_}and}}{Formula.mk:unand}}\\{{\nwixident{Formula.mk{\_}iff}}{Formula.mk:uniff}}\\{{\nwixident{Formula.mk{\_}or}}{Formula.mk:unor}}\\{{\nwixident{FS0.app}}{FS0.app}}\\{{\nwixident{FS0.equal}}{FS0.equal}}\\{{\nwixident{FS0.In}}{FS0.In}}\\{{\nwixident{FS0.intersection}}{FS0.intersection}}\\{{\nwixident{FS0.preimage}}{FS0.preimage}}\\{{\nwixident{FS0.singleton}}{FS0.singleton}}\\{{\nwixident{FS0.union}}{FS0.union}}\\{{\nwixident{Sort.IND}}{Sort.IND}}\\{{\nwixident{Term.fresh{\_}var}}{Term.fresh:unvar}}\\{{\nwixident{Term.fv}}{Term.fv}}}\nwindexuse{\nwixident{Formula.mk{\_}and}}{Formula.mk:unand}{NW4ICsJQ-2aDdLu-5}\nwindexuse{\nwixident{Formula.mk{\_}iff}}{Formula.mk:uniff}{NW4ICsJQ-2aDdLu-5}\nwindexuse{\nwixident{Formula.mk{\_}or}}{Formula.mk:unor}{NW4ICsJQ-2aDdLu-5}\nwindexuse{\nwixident{FS0.app}}{FS0.app}{NW4ICsJQ-2aDdLu-5}\nwindexuse{\nwixident{FS0.equal}}{FS0.equal}{NW4ICsJQ-2aDdLu-5}\nwindexuse{\nwixident{FS0.In}}{FS0.In}{NW4ICsJQ-2aDdLu-5}\nwindexuse{\nwixident{FS0.intersection}}{FS0.intersection}{NW4ICsJQ-2aDdLu-5}\nwindexuse{\nwixident{FS0.preimage}}{FS0.preimage}{NW4ICsJQ-2aDdLu-5}\nwindexuse{\nwixident{FS0.singleton}}{FS0.singleton}{NW4ICsJQ-2aDdLu-5}\nwindexuse{\nwixident{FS0.union}}{FS0.union}{NW4ICsJQ-2aDdLu-5}\nwindexuse{\nwixident{Sort.IND}}{Sort.IND}{NW4ICsJQ-2aDdLu-5}\nwindexuse{\nwixident{Term.fresh{\_}var}}{Term.fresh:unvar}{NW4ICsJQ-2aDdLu-5}\nwindexuse{\nwixident{Term.fv}}{Term.fv}{NW4ICsJQ-2aDdLu-5}\nwendcode{}\nwbegindocs{66}\nwdocspar
\begin{node}
Inductive generation axioms.
\begin{enumerate}
  \item $\forall A,B : \textit{Class}\ldotp A\subset \mathcal{I}_{2}(A,B)$
  \item $\forall A,B : \textit{Class}\ldotp\forall
    x,y_{1},y_{2}:\textit{Ind}\ldotp (y_{1}\in \mathcal{I}_{2}(A,B)\implies y_{2}\in
    \mathcal{I}_{2}(A,B)\implies(x,y_{1},y_{2})\in B\implies x\in \mathcal{I}_{2}(A,B))$
  \item $\forall A,B,C,X : \textit{Class}\ldotp
    C=\mathcal{I}_{2}(A,B)\implies A\subset X\land(\forall
    x,y_{1},y_{2}:\textit{Ind}\ldotp (y_{1}\in X\implies y_{2}\in
    X\implies(x,y_{1},y_{2})\in B\implies x\in X))\implies C\subset X$
\end{enumerate}
\end{node}
\nwenddocs{}\nwbegincode{67}\sublabel{NW4ICsJQ-hYCIS-6}\nwmargintag{{\nwtagstyle{}\subpageref{NW4ICsJQ-hYCIS-6}}}\moddef{Specification for $FS_{0}$ axioms~{\nwtagstyle{}\subpageref{NW4ICsJQ-hYCIS-1}}}\plusendmoddef\nwstartdeflinemarkup\nwusesondefline{\\{NW4ICsJQ-2cY0B1-1}}\nwprevnextdefs{NW4ICsJQ-hYCIS-5}{NW4ICsJQ-hYCIS-7}\nwenddeflinemarkup
val axiom_induct_base : \nwlinkedidentc{Term.t}{NW1rTmOJ-4LsRWC-1} -> \nwlinkedidentc{Term.t}{NW1rTmOJ-4LsRWC-1} -> t;
val axiom_induct_gen : \nwlinkedidentc{Term.t}{NW1rTmOJ-4LsRWC-1} -> \nwlinkedidentc{Term.t}{NW1rTmOJ-4LsRWC-1} -> t;
val axiom_induct_min : \nwlinkedidentc{Term.t}{NW1rTmOJ-4LsRWC-1} -> \nwlinkedidentc{Term.t}{NW1rTmOJ-4LsRWC-1} -> \nwlinkedidentc{Term.t}{NW1rTmOJ-4LsRWC-1} -> t;
\nwused{\\{NW4ICsJQ-2cY0B1-1}}\nwidentuses{\\{{\nwixident{Term.t}}{Term.t}}}\nwindexuse{\nwixident{Term.t}}{Term.t}{NW4ICsJQ-hYCIS-6}\nwendcode{}\nwbegindocs{68}\nwdocspar
\nwenddocs{}\nwbegincode{69}\sublabel{NW4ICsJQ-2aDdLu-6}\nwmargintag{{\nwtagstyle{}\subpageref{NW4ICsJQ-2aDdLu-6}}}\moddef{Axioms for $FS_{0}$~{\nwtagstyle{}\subpageref{NW4ICsJQ-2aDdLu-1}}}\plusendmoddef\nwstartdeflinemarkup\nwusesondefline{\\{NW4ICsJQ-YIPBT-1}}\nwprevnextdefs{NW4ICsJQ-2aDdLu-5}{NW4ICsJQ-2aDdLu-7}\nwenddeflinemarkup
fun axiom_induct_base A B =
  let val x = \nwlinkedidentc{Term.fresh_var}{NW1rTmOJ-4Ykq9M-1} "x" \nwlinkedidentc{Sort.IND}{NW1rTmOJ-3ZqggC-1} ((\nwlinkedidentc{Term.fv}{NW1rTmOJ-2nRtky-1} A) @ (\nwlinkedidentc{Term.fv}{NW1rTmOJ-2nRtky-1} B))
  in mk_axiom(\nwlinkedidentc{Formula.mk_forall}{NW3w0iRO-3SlqND-4}(x, \nwlinkedidentc{Formula.mk_imp}{NW3w0iRO-3SlqND-3}(\nwlinkedidentc{FS0.In}{NW2DqMzY-eTij6-1} x A,
                                                  \nwlinkedidentc{FS0.In}{NW2DqMzY-eTij6-1} x (\nwlinkedidentc{FS0.induct}{NW2DqMzY-1mMvkv-1} A B))))
  end;

fun axiom_induct_gen A B =
  let
    val vars = ((\nwlinkedidentc{Term.fv}{NW1rTmOJ-2nRtky-1} A) @ (\nwlinkedidentc{Term.fv}{NW1rTmOJ-2nRtky-1} B));
    fun var x = \nwlinkedidentc{Term.fresh_var}{NW1rTmOJ-4Ykq9M-1} x \nwlinkedidentc{Sort.IND}{NW1rTmOJ-3ZqggC-1} vars;
    val x = var "x";
    val y = var "y";
    val z = var "z";
    val C = \nwlinkedidentc{FS0.induct}{NW2DqMzY-1mMvkv-1} A B;
    val triple = \nwlinkedidentc{FS0.pair}{NW2DqMzY-4SkDla-1} x (\nwlinkedidentc{FS0.pair}{NW2DqMzY-4SkDla-1} y z);
    val p1 = \nwlinkedidentc{FS0.In}{NW2DqMzY-eTij6-1} y C;
    val p2 = \nwlinkedidentc{FS0.In}{NW2DqMzY-eTij6-1} z C;
    val p3 = \nwlinkedidentc{FS0.In}{NW2DqMzY-eTij6-1} triple B;
    val p4 = \nwlinkedidentc{FS0.In}{NW2DqMzY-eTij6-1} x C;
    val premises = List.foldr \nwlinkedidentc{Formula.mk_and}{NW3w0iRO-3SlqND-3} p3 [p1, p2];
    val claim = \nwlinkedidentc{Formula.mk_imp}{NW3w0iRO-3SlqND-3}(premises, p4);
  in
    mk_axiom(List.foldr \nwlinkedidentc{Formula.mk_forall}{NW3w0iRO-3SlqND-4} claim [x,y,z])
  end;

fun axiom_induct_min A B X =
  let
    val vars = ((\nwlinkedidentc{Term.fv}{NW1rTmOJ-2nRtky-1} A) @ (\nwlinkedidentc{Term.fv}{NW1rTmOJ-2nRtky-1} B));
    fun var x = \nwlinkedidentc{Term.fresh_var}{NW1rTmOJ-4Ykq9M-1} x \nwlinkedidentc{Sort.IND}{NW1rTmOJ-3ZqggC-1} vars;
    val x = var "x";
    val y = var "y";
    val z = var "z";
    val p0 = \nwlinkedidentc{Formula.mk_forall}{NW3w0iRO-3SlqND-4}(x, \nwlinkedidentc{Formula.mk_imp}{NW3w0iRO-3SlqND-3}(\nwlinkedidentc{FS0.In}{NW2DqMzY-eTij6-1} x A,
                                                 \nwlinkedidentc{FS0.In}{NW2DqMzY-eTij6-1} x X));
    val C = \nwlinkedidentc{FS0.induct}{NW2DqMzY-1mMvkv-1} A B;
    val triple = \nwlinkedidentc{FS0.pair}{NW2DqMzY-4SkDla-1} x (\nwlinkedidentc{FS0.pair}{NW2DqMzY-4SkDla-1} y z);
    val p1 = \nwlinkedidentc{FS0.In}{NW2DqMzY-eTij6-1} y X;
    val p2 = \nwlinkedidentc{FS0.In}{NW2DqMzY-eTij6-1} z X;
    val p3 = \nwlinkedidentc{FS0.In}{NW2DqMzY-eTij6-1} triple B;
    val p4 = \nwlinkedidentc{FS0.In}{NW2DqMzY-eTij6-1} x X;
    val premises = List.foldr \nwlinkedidentc{Formula.mk_and}{NW3w0iRO-3SlqND-3} p3 [p1, p2];
    val claim = \nwlinkedidentc{Formula.mk_imp}{NW3w0iRO-3SlqND-3}(premises, p4);
    val ind_case = List.foldr \nwlinkedidentc{Formula.mk_forall}{NW3w0iRO-3SlqND-4} claim [x,y,z];
    val conc = \nwlinkedidentc{Formula.mk_forall}{NW3w0iRO-3SlqND-4}(x, \nwlinkedidentc{Formula.mk_imp}{NW3w0iRO-3SlqND-3}(\nwlinkedidentc{FS0.In}{NW2DqMzY-eTij6-1} x C,
                                                   \nwlinkedidentc{FS0.In}{NW2DqMzY-eTij6-1} x X));
  in
    mk_axiom(\nwlinkedidentc{Formula.mk_imp}{NW3w0iRO-3SlqND-3}(\nwlinkedidentc{Formula.mk_and}{NW3w0iRO-3SlqND-3}(p0,ind_case),
                            conc))
  end;
\nwindexdefn{\nwixident{Thm.axiom{\_}induct{\_}base}}{Thm.axiom:uninduct:unbase}{NW4ICsJQ-2aDdLu-6}\nwindexdefn{\nwixident{Thm.axiom{\_}induct{\_}gen}}{Thm.axiom:uninduct:ungen}{NW4ICsJQ-2aDdLu-6}\nwindexdefn{\nwixident{Thm.axiom{\_}induct{\_}min}}{Thm.axiom:uninduct:unmin}{NW4ICsJQ-2aDdLu-6}\eatline
\nwused{\\{NW4ICsJQ-YIPBT-1}}\nwidentdefs{\\{{\nwixident{Thm.axiom{\_}induct{\_}base}}{Thm.axiom:uninduct:unbase}}\\{{\nwixident{Thm.axiom{\_}induct{\_}gen}}{Thm.axiom:uninduct:ungen}}\\{{\nwixident{Thm.axiom{\_}induct{\_}min}}{Thm.axiom:uninduct:unmin}}}\nwidentuses{\\{{\nwixident{Formula.mk{\_}and}}{Formula.mk:unand}}\\{{\nwixident{Formula.mk{\_}forall}}{Formula.mk:unforall}}\\{{\nwixident{Formula.mk{\_}imp}}{Formula.mk:unimp}}\\{{\nwixident{FS0.In}}{FS0.In}}\\{{\nwixident{FS0.induct}}{FS0.induct}}\\{{\nwixident{FS0.pair}}{FS0.pair}}\\{{\nwixident{Sort.IND}}{Sort.IND}}\\{{\nwixident{Term.fresh{\_}var}}{Term.fresh:unvar}}\\{{\nwixident{Term.fv}}{Term.fv}}}\nwindexuse{\nwixident{Formula.mk{\_}and}}{Formula.mk:unand}{NW4ICsJQ-2aDdLu-6}\nwindexuse{\nwixident{Formula.mk{\_}forall}}{Formula.mk:unforall}{NW4ICsJQ-2aDdLu-6}\nwindexuse{\nwixident{Formula.mk{\_}imp}}{Formula.mk:unimp}{NW4ICsJQ-2aDdLu-6}\nwindexuse{\nwixident{FS0.In}}{FS0.In}{NW4ICsJQ-2aDdLu-6}\nwindexuse{\nwixident{FS0.induct}}{FS0.induct}{NW4ICsJQ-2aDdLu-6}\nwindexuse{\nwixident{FS0.pair}}{FS0.pair}{NW4ICsJQ-2aDdLu-6}\nwindexuse{\nwixident{Sort.IND}}{Sort.IND}{NW4ICsJQ-2aDdLu-6}\nwindexuse{\nwixident{Term.fresh{\_}var}}{Term.fresh:unvar}{NW4ICsJQ-2aDdLu-6}\nwindexuse{\nwixident{Term.fv}}{Term.fv}{NW4ICsJQ-2aDdLu-6}\nwendcode{}\nwbegindocs{70}\nwdocspar
\begin{node}
We also have induction on the universe.
\begin{enumerate}
  \item $\forall X:\textit{Class}\ldotp 0\in X\land(\forall
    x,y:\textit{Ind}\ldotp(x\in X\implies y\in X\implies(x,y)\in
    X))\implies(\forall x:\textit{Ind}\ldotp x\in X)$.
\end{enumerate}
More generally, as Feferman observed earlier~\cite[\S2.1]{feferman1982inductively}, if we have a [nonempty] set of Urelements $U$, then we
need to change the premise $0\in X$ to $(\forall u\ldotp u\in U\implies u\in X)$.
\end{node}

\nwenddocs{}\nwbegincode{71}\sublabel{NW4ICsJQ-hYCIS-7}\nwmargintag{{\nwtagstyle{}\subpageref{NW4ICsJQ-hYCIS-7}}}\moddef{Specification for $FS_{0}$ axioms~{\nwtagstyle{}\subpageref{NW4ICsJQ-hYCIS-1}}}\plusendmoddef\nwstartdeflinemarkup\nwusesondefline{\\{NW4ICsJQ-2cY0B1-1}}\nwprevnextdefs{NW4ICsJQ-hYCIS-6}{NW4ICsJQ-hYCIS-8}\nwenddeflinemarkup
val axiom_universe_ind : \nwlinkedidentc{Term.t}{NW1rTmOJ-4LsRWC-1} -> t;
\nwused{\\{NW4ICsJQ-2cY0B1-1}}\nwidentuses{\\{{\nwixident{Term.t}}{Term.t}}}\nwindexuse{\nwixident{Term.t}}{Term.t}{NW4ICsJQ-hYCIS-7}\nwendcode{}\nwbegindocs{72}\nwdocspar

\nwenddocs{}\nwbegincode{73}\sublabel{NW4ICsJQ-2aDdLu-7}\nwmargintag{{\nwtagstyle{}\subpageref{NW4ICsJQ-2aDdLu-7}}}\moddef{Axioms for $FS_{0}$~{\nwtagstyle{}\subpageref{NW4ICsJQ-2aDdLu-1}}}\plusendmoddef\nwstartdeflinemarkup\nwusesondefline{\\{NW4ICsJQ-YIPBT-1}}\nwprevnextdefs{NW4ICsJQ-2aDdLu-6}{NW4ICsJQ-2aDdLu-8}\nwenddeflinemarkup
fun axiom_universe_ind X =
  let
    val vars = (\nwlinkedidentc{Term.fv}{NW1rTmOJ-2nRtky-1} X);
    fun var x = \nwlinkedidentc{Term.fresh_var}{NW1rTmOJ-4Ykq9M-1} x \nwlinkedidentc{Sort.IND}{NW1rTmOJ-3ZqggC-1} vars;
    val x = var "x";
    val y = var "y";
    val p0 = \nwlinkedidentc{FS0.In}{NW2DqMzY-eTij6-1} (\nwlinkedidentc{FS0.zero}{NW2DqMzY-1NUSQs-1}) X;
    val claim = \nwlinkedidentc{Formula.mk_imp}{NW3w0iRO-3SlqND-3}(\nwlinkedidentc{Formula.mk_and}{NW3w0iRO-3SlqND-3}(\nwlinkedidentc{FS0.In}{NW2DqMzY-eTij6-1} x X,
                                              \nwlinkedidentc{FS0.In}{NW2DqMzY-eTij6-1} y X),
                               \nwlinkedidentc{FS0.In}{NW2DqMzY-eTij6-1} (\nwlinkedidentc{FS0.pair}{NW2DqMzY-4SkDla-1} x y) X);
    val p1 = List.foldr \nwlinkedidentc{Formula.mk_forall}{NW3w0iRO-3SlqND-4} claim [x,y];
    val conc = \nwlinkedidentc{Formula.mk_forall}{NW3w0iRO-3SlqND-4}(x, \nwlinkedidentc{FS0.In}{NW2DqMzY-eTij6-1} x X);
  in
    mk_axiom(\nwlinkedidentc{Formula.mk_imp}{NW3w0iRO-3SlqND-3}(\nwlinkedidentc{Formula.mk_and}{NW3w0iRO-3SlqND-3}(p0, p1),
                            conc))
  end;
\nwindexdefn{\nwixident{Thm.axiom{\_}universe{\_}ind}}{Thm.axiom:ununiverse:unind}{NW4ICsJQ-2aDdLu-7}\eatline
\nwused{\\{NW4ICsJQ-YIPBT-1}}\nwidentdefs{\\{{\nwixident{Thm.axiom{\_}universe{\_}ind}}{Thm.axiom:ununiverse:unind}}}\nwidentuses{\\{{\nwixident{Formula.mk{\_}and}}{Formula.mk:unand}}\\{{\nwixident{Formula.mk{\_}forall}}{Formula.mk:unforall}}\\{{\nwixident{Formula.mk{\_}imp}}{Formula.mk:unimp}}\\{{\nwixident{FS0.In}}{FS0.In}}\\{{\nwixident{FS0.pair}}{FS0.pair}}\\{{\nwixident{FS0.zero}}{FS0.zero}}\\{{\nwixident{Sort.IND}}{Sort.IND}}\\{{\nwixident{Term.fresh{\_}var}}{Term.fresh:unvar}}\\{{\nwixident{Term.fv}}{Term.fv}}}\nwindexuse{\nwixident{Formula.mk{\_}and}}{Formula.mk:unand}{NW4ICsJQ-2aDdLu-7}\nwindexuse{\nwixident{Formula.mk{\_}forall}}{Formula.mk:unforall}{NW4ICsJQ-2aDdLu-7}\nwindexuse{\nwixident{Formula.mk{\_}imp}}{Formula.mk:unimp}{NW4ICsJQ-2aDdLu-7}\nwindexuse{\nwixident{FS0.In}}{FS0.In}{NW4ICsJQ-2aDdLu-7}\nwindexuse{\nwixident{FS0.pair}}{FS0.pair}{NW4ICsJQ-2aDdLu-7}\nwindexuse{\nwixident{FS0.zero}}{FS0.zero}{NW4ICsJQ-2aDdLu-7}\nwindexuse{\nwixident{Sort.IND}}{Sort.IND}{NW4ICsJQ-2aDdLu-7}\nwindexuse{\nwixident{Term.fresh{\_}var}}{Term.fresh:unvar}{NW4ICsJQ-2aDdLu-7}\nwindexuse{\nwixident{Term.fv}}{Term.fv}{NW4ICsJQ-2aDdLu-7}\nwendcode{}\nwbegindocs{74}\nwdocspar
\begin{node}[Structural induction axioms]
Feferman implicitly has induction axioms for the three sorts. For
individuals, this looks like:
\begin{quote}
Let $\mathcal{P}[-]$ be a unary predicate of individual-sorted
terms. Then: $\mathcal{P}[0]\land(\forall x,y\ldotp\mathcal{P}[x]\land\mathcal{P}[y]\implies\mathcal{P}[(x,y)])\land(\forall f,\forall x\ldotp\mathcal{P}[x]\implies\mathcal{P}[f(x)])\implies\forall x\ldotp\mathcal{P}[x]$ 
where $x$ and $y$ are [fresh] individual-sorted variables and $f$ is a [fresh]
function-sorted individual.
\end{quote}
There are similar induction principles for function-sorted terms and
class-sorted terms. Formally, the structural induction principle for individuals:
\begin{equation}
\begin{split}
\mathcal{P}[0]&\land(\forall x,y:Ind\ldotp\mathcal{P}[x]\land\mathcal{P}[y]\implies\mathcal{P}[(x,y)])\land\\
&\land(\forall f:Fun,\forall x:Ind\ldotp\mathcal{P}[x]\implies\mathcal{P}[f(x)])\implies\\
&\quad\implies(\forall x:Ind\ldotp\mathcal{P}[x])
\end{split}
\end{equation}
First, we will construct a helper function to ensure that we have
fresh variables used for binding.
\end{node}

\nwenddocs{}\nwbegincode{75}\sublabel{NW4ICsJQ-hYCIS-8}\nwmargintag{{\nwtagstyle{}\subpageref{NW4ICsJQ-hYCIS-8}}}\moddef{Specification for $FS_{0}$ axioms~{\nwtagstyle{}\subpageref{NW4ICsJQ-hYCIS-1}}}\plusendmoddef\nwstartdeflinemarkup\nwusesondefline{\\{NW4ICsJQ-2cY0B1-1}}\nwprevnextdefs{NW4ICsJQ-hYCIS-7}{NW4ICsJQ-hYCIS-9}\nwenddeflinemarkup
val Ind_ind : (\nwlinkedidentc{Term.t}{NW1rTmOJ-4LsRWC-1} -> \nwlinkedidentc{Formula.t}{NW3w0iRO-2jnj9W-1}) -> t;
\nwused{\\{NW4ICsJQ-2cY0B1-1}}\nwidentuses{\\{{\nwixident{Formula.t}}{Formula.t}}\\{{\nwixident{Term.t}}{Term.t}}}\nwindexuse{\nwixident{Formula.t}}{Formula.t}{NW4ICsJQ-hYCIS-8}\nwindexuse{\nwixident{Term.t}}{Term.t}{NW4ICsJQ-hYCIS-8}\nwendcode{}\nwbegindocs{76}\nwdocspar

\nwenddocs{}\nwbegincode{77}\sublabel{NW4ICsJQ-2aDdLu-8}\nwmargintag{{\nwtagstyle{}\subpageref{NW4ICsJQ-2aDdLu-8}}}\moddef{Axioms for $FS_{0}$~{\nwtagstyle{}\subpageref{NW4ICsJQ-2aDdLu-1}}}\plusendmoddef\nwstartdeflinemarkup\nwusesondefline{\\{NW4ICsJQ-YIPBT-1}}\nwprevnextdefs{NW4ICsJQ-2aDdLu-7}{NW4ICsJQ-2aDdLu-9}\nwenddeflinemarkup
\LA{}Fresh variable for a formula~{\nwtagstyle{}\subpageref{NW4ICsJQ-2auZja-1}}\RA{}

(* Ind_ind : (\nwlinkedidentc{Term.t}{NW1rTmOJ-4LsRWC-1} -> \nwlinkedidentc{Formula.t}{NW3w0iRO-2jnj9W-1}) -> \nwlinkedidentc{Thm.t}{NW4ICsJQ-YIPBT-1} *)
fun Ind_ind P = 
  let
    val tmp = Term.Var("_", \nwlinkedidentc{Sort.IND}{NW1rTmOJ-3ZqggC-1});
    val x = fresh_var_in "x" \nwlinkedidentc{Sort.IND}{NW1rTmOJ-3ZqggC-1} (P tmp);
    val y = fresh_var_in "y" \nwlinkedidentc{Sort.IND}{NW1rTmOJ-3ZqggC-1} (P x);
    val f = fresh_var_in "f" \nwlinkedidentc{Sort.FUN}{NW1rTmOJ-3ZqggC-1} (P x);
    val base = P (\nwlinkedidentc{FS0.zero}{NW2DqMzY-1NUSQs-1});
    val pair = \nwlinkedidentc{Formula.mk_forall}{NW3w0iRO-3SlqND-4}(x, \nwlinkedidentc{Formula.mk_forall}{NW3w0iRO-3SlqND-4}(y, \nwlinkedidentc{Formula.mk_imp}{NW3w0iRO-3SlqND-3}(\nwlinkedidentc{Formula.mk_and}{NW3w0iRO-3SlqND-3}(P x, P y),
                                                                        P (\nwlinkedidentc{FS0.pair}{NW2DqMzY-4SkDla-1} x y))));
    val app = \nwlinkedidentc{Formula.mk_forall}{NW3w0iRO-3SlqND-4}(f, \nwlinkedidentc{Formula.mk_forall}{NW3w0iRO-3SlqND-4}(x,
                                                     \nwlinkedidentc{Formula.mk_imp}{NW3w0iRO-3SlqND-3}(P x,
                                                                    P(\nwlinkedidentc{FS0.app}{NW2DqMzY-2oGoBx-1} f x))));
    val hypo = \nwlinkedidentc{Formula.mk_and}{NW3w0iRO-3SlqND-3}(base, \nwlinkedidentc{Formula.mk_and}{NW3w0iRO-3SlqND-3}(pair,
                                                   app));
    val conc = \nwlinkedidentc{Formula.mk_forall}{NW3w0iRO-3SlqND-4}(x, P x);
  in
    mk_axiom(\nwlinkedidentc{Formula.mk_imp}{NW3w0iRO-3SlqND-3}(hypo,
                            conc))
  end;
\nwindexdefn{\nwixident{Thm.Ind{\_}ind}}{Thm.Ind:unind}{NW4ICsJQ-2aDdLu-8}\eatline
\nwused{\\{NW4ICsJQ-YIPBT-1}}\nwidentdefs{\\{{\nwixident{Thm.Ind{\_}ind}}{Thm.Ind:unind}}}\nwidentuses{\\{{\nwixident{Formula.mk{\_}and}}{Formula.mk:unand}}\\{{\nwixident{Formula.mk{\_}forall}}{Formula.mk:unforall}}\\{{\nwixident{Formula.mk{\_}imp}}{Formula.mk:unimp}}\\{{\nwixident{Formula.t}}{Formula.t}}\\{{\nwixident{FS0.app}}{FS0.app}}\\{{\nwixident{FS0.pair}}{FS0.pair}}\\{{\nwixident{FS0.zero}}{FS0.zero}}\\{{\nwixident{Sort.FUN}}{Sort.FUN}}\\{{\nwixident{Sort.IND}}{Sort.IND}}\\{{\nwixident{Term.t}}{Term.t}}\\{{\nwixident{Thm.t}}{Thm.t}}}\nwindexuse{\nwixident{Formula.mk{\_}and}}{Formula.mk:unand}{NW4ICsJQ-2aDdLu-8}\nwindexuse{\nwixident{Formula.mk{\_}forall}}{Formula.mk:unforall}{NW4ICsJQ-2aDdLu-8}\nwindexuse{\nwixident{Formula.mk{\_}imp}}{Formula.mk:unimp}{NW4ICsJQ-2aDdLu-8}\nwindexuse{\nwixident{Formula.t}}{Formula.t}{NW4ICsJQ-2aDdLu-8}\nwindexuse{\nwixident{FS0.app}}{FS0.app}{NW4ICsJQ-2aDdLu-8}\nwindexuse{\nwixident{FS0.pair}}{FS0.pair}{NW4ICsJQ-2aDdLu-8}\nwindexuse{\nwixident{FS0.zero}}{FS0.zero}{NW4ICsJQ-2aDdLu-8}\nwindexuse{\nwixident{Sort.FUN}}{Sort.FUN}{NW4ICsJQ-2aDdLu-8}\nwindexuse{\nwixident{Sort.IND}}{Sort.IND}{NW4ICsJQ-2aDdLu-8}\nwindexuse{\nwixident{Term.t}}{Term.t}{NW4ICsJQ-2aDdLu-8}\nwindexuse{\nwixident{Thm.t}}{Thm.t}{NW4ICsJQ-2aDdLu-8}\nwendcode{}\nwbegindocs{78}\nwdocspar
\begin{node}
Structural induction for function-sorted terms looks like,
\begin{equation*}
\begin{split}
\mathcal{P}[\id]&\land\mathcal{P}[\pi_{1}]\land\mathcal{P}[\pi_{2}]\land\mathcal{P}[D]\land(\forall a:Ind\ldotp\mathcal{P}[K_{a}])\land\\
&\land(\forall f,g:Fun\ldotp\mathcal{P}[f]\land\mathcal{P}[g]\implies\mathcal{P}[f\circ g])\land\\
&\land(\forall f,g:Fun\ldotp\mathcal{P}[f]\land\mathcal{P}[g]\implies\mathcal{P}[(f,g)])\land\\
&\land(\forall f,g:Fun\ldotp\mathcal{P}[f]\land\mathcal{P}[g]\implies\mathcal{P}[\mathcal{R}(f,g)])\implies\\
&\quad\implies(\forall f:Fun\ldotp\mathcal{P}[f])
\end{split}
\end{equation*}
\end{node}

\nwenddocs{}\nwbegincode{79}\sublabel{NW4ICsJQ-hYCIS-9}\nwmargintag{{\nwtagstyle{}\subpageref{NW4ICsJQ-hYCIS-9}}}\moddef{Specification for $FS_{0}$ axioms~{\nwtagstyle{}\subpageref{NW4ICsJQ-hYCIS-1}}}\plusendmoddef\nwstartdeflinemarkup\nwusesondefline{\\{NW4ICsJQ-2cY0B1-1}}\nwprevnextdefs{NW4ICsJQ-hYCIS-8}{NW4ICsJQ-hYCIS-A}\nwenddeflinemarkup
val Fun_ind : (\nwlinkedidentc{Term.t}{NW1rTmOJ-4LsRWC-1} -> \nwlinkedidentc{Formula.t}{NW3w0iRO-2jnj9W-1}) -> t;
\nwused{\\{NW4ICsJQ-2cY0B1-1}}\nwidentuses{\\{{\nwixident{Formula.t}}{Formula.t}}\\{{\nwixident{Term.t}}{Term.t}}}\nwindexuse{\nwixident{Formula.t}}{Formula.t}{NW4ICsJQ-hYCIS-9}\nwindexuse{\nwixident{Term.t}}{Term.t}{NW4ICsJQ-hYCIS-9}\nwendcode{}\nwbegindocs{80}\nwdocspar

\nwenddocs{}\nwbegincode{81}\sublabel{NW4ICsJQ-2aDdLu-9}\nwmargintag{{\nwtagstyle{}\subpageref{NW4ICsJQ-2aDdLu-9}}}\moddef{Axioms for $FS_{0}$~{\nwtagstyle{}\subpageref{NW4ICsJQ-2aDdLu-1}}}\plusendmoddef\nwstartdeflinemarkup\nwusesondefline{\\{NW4ICsJQ-YIPBT-1}}\nwprevnextdefs{NW4ICsJQ-2aDdLu-8}{NW4ICsJQ-2aDdLu-A}\nwenddeflinemarkup
(* Fun_ind : (\nwlinkedidentc{Term.t}{NW1rTmOJ-4LsRWC-1} -> \nwlinkedidentc{Formula.t}{NW3w0iRO-2jnj9W-1}) -> t *)
fun Fun_ind P =
  let
    val tmp = Term.Var("_", \nwlinkedidentc{Sort.FUN}{NW1rTmOJ-3ZqggC-1});
    val a = fresh_var_in "a" \nwlinkedidentc{Sort.IND}{NW1rTmOJ-3ZqggC-1} (P tmp);
    val f = fresh_var_in "f" \nwlinkedidentc{Sort.FUN}{NW1rTmOJ-3ZqggC-1} (P tmp);
    val g = fresh_var_in "g" \nwlinkedidentc{Sort.FUN}{NW1rTmOJ-3ZqggC-1} (P tmp);
    fun ind_case (con : \nwlinkedidentc{Term.t}{NW1rTmOJ-4LsRWC-1} -> \nwlinkedidentc{Term.t}{NW1rTmOJ-4LsRWC-1} -> \nwlinkedidentc{Term.t}{NW1rTmOJ-4LsRWC-1}) =
      \nwlinkedidentc{Formula.mk_forall}{NW3w0iRO-3SlqND-4}(f, \nwlinkedidentc{Formula.mk_forall}{NW3w0iRO-3SlqND-4}(g, \nwlinkedidentc{Formula.mk_imp}{NW3w0iRO-3SlqND-3}(\nwlinkedidentc{Formula.mk_and}{NW3w0iRO-3SlqND-3}(P f,
                                                                              P g),
                                                               P(con f g))));
    val premises =
      List.foldr \nwlinkedidentc{Formula.mk_and}{NW3w0iRO-3SlqND-3}
                 (ind_case \nwlinkedidentc{FS0.recur}{NW2DqMzY-2fxlqE-1})
                 [P (FS0.id_const),
                  P (FS0.proj1_const),
                  P (FS0.proj2_const),
                  \nwlinkedidentc{Formula.mk_forall}{NW3w0iRO-3SlqND-4}(a, P (\nwlinkedidentc{FS0.const}{NW2DqMzY-22MFug-1} a)),
                  ind_case (\nwlinkedidentc{FS0.comp}{NW2DqMzY-1dWhhL-1}),
                  ind_case (\nwlinkedidentc{FS0.juxt}{NW2DqMzY-1lMRH2-1})];
    val conc = \nwlinkedidentc{Formula.mk_forall}{NW3w0iRO-3SlqND-4}(f, P f);
  in
    mk_axiom(\nwlinkedidentc{Formula.mk_imp}{NW3w0iRO-3SlqND-3}(premises,
                            conc))
  end;  
\nwindexdefn{\nwixident{Thm.Fun{\_}ind}}{Thm.Fun:unind}{NW4ICsJQ-2aDdLu-9}\eatline
\nwused{\\{NW4ICsJQ-YIPBT-1}}\nwidentdefs{\\{{\nwixident{Thm.Fun{\_}ind}}{Thm.Fun:unind}}}\nwidentuses{\\{{\nwixident{Formula.mk{\_}and}}{Formula.mk:unand}}\\{{\nwixident{Formula.mk{\_}forall}}{Formula.mk:unforall}}\\{{\nwixident{Formula.mk{\_}imp}}{Formula.mk:unimp}}\\{{\nwixident{Formula.t}}{Formula.t}}\\{{\nwixident{FS0.comp}}{FS0.comp}}\\{{\nwixident{FS0.const}}{FS0.const}}\\{{\nwixident{FS0.juxt}}{FS0.juxt}}\\{{\nwixident{FS0.recur}}{FS0.recur}}\\{{\nwixident{Sort.FUN}}{Sort.FUN}}\\{{\nwixident{Sort.IND}}{Sort.IND}}\\{{\nwixident{Term.t}}{Term.t}}}\nwindexuse{\nwixident{Formula.mk{\_}and}}{Formula.mk:unand}{NW4ICsJQ-2aDdLu-9}\nwindexuse{\nwixident{Formula.mk{\_}forall}}{Formula.mk:unforall}{NW4ICsJQ-2aDdLu-9}\nwindexuse{\nwixident{Formula.mk{\_}imp}}{Formula.mk:unimp}{NW4ICsJQ-2aDdLu-9}\nwindexuse{\nwixident{Formula.t}}{Formula.t}{NW4ICsJQ-2aDdLu-9}\nwindexuse{\nwixident{FS0.comp}}{FS0.comp}{NW4ICsJQ-2aDdLu-9}\nwindexuse{\nwixident{FS0.const}}{FS0.const}{NW4ICsJQ-2aDdLu-9}\nwindexuse{\nwixident{FS0.juxt}}{FS0.juxt}{NW4ICsJQ-2aDdLu-9}\nwindexuse{\nwixident{FS0.recur}}{FS0.recur}{NW4ICsJQ-2aDdLu-9}\nwindexuse{\nwixident{Sort.FUN}}{Sort.FUN}{NW4ICsJQ-2aDdLu-9}\nwindexuse{\nwixident{Sort.IND}}{Sort.IND}{NW4ICsJQ-2aDdLu-9}\nwindexuse{\nwixident{Term.t}}{Term.t}{NW4ICsJQ-2aDdLu-9}\nwendcode{}\nwbegindocs{82}\nwdocspar
\begin{node}
Structural induction principle for classes:
\begin{equation*}
\begin{split}
\mathcal{P}[\{0\}]&\land(\forall f:Fun\forall A:Cl\ldotp\mathcal{P}[A]\implies\mathcal{P}[f^{-1}A])\land\\
&\land(\forall A,B:Cl\ldotp\mathcal{P}[A]\land\mathcal{P}[B]\implies\mathcal{P}[A\cup B])\land\\
&\land(\forall A,B:Cl\ldotp\mathcal{P}[A]\land\mathcal{P}[B]\implies\mathcal{P}[A\cap B])\land\\
&\land(\forall A,B:Cl\ldotp\mathcal{P}[A]\land\mathcal{P}[B]\implies\mathcal{P}[I_{2}(A,B)])\implies\\
&\quad\implies(\forall A:Cl\ldotp\mathcal{P}[A])
\end{split}
\end{equation*}
\end{node}

\nwenddocs{}\nwbegincode{83}\sublabel{NW4ICsJQ-hYCIS-A}\nwmargintag{{\nwtagstyle{}\subpageref{NW4ICsJQ-hYCIS-A}}}\moddef{Specification for $FS_{0}$ axioms~{\nwtagstyle{}\subpageref{NW4ICsJQ-hYCIS-1}}}\plusendmoddef\nwstartdeflinemarkup\nwusesondefline{\\{NW4ICsJQ-2cY0B1-1}}\nwprevnextdefs{NW4ICsJQ-hYCIS-9}{NW4ICsJQ-hYCIS-B}\nwenddeflinemarkup
val Class_ind : (\nwlinkedidentc{Term.t}{NW1rTmOJ-4LsRWC-1} -> \nwlinkedidentc{Formula.t}{NW3w0iRO-2jnj9W-1}) -> t;
\nwused{\\{NW4ICsJQ-2cY0B1-1}}\nwidentuses{\\{{\nwixident{Formula.t}}{Formula.t}}\\{{\nwixident{Term.t}}{Term.t}}}\nwindexuse{\nwixident{Formula.t}}{Formula.t}{NW4ICsJQ-hYCIS-A}\nwindexuse{\nwixident{Term.t}}{Term.t}{NW4ICsJQ-hYCIS-A}\nwendcode{}\nwbegindocs{84}\nwdocspar

\nwenddocs{}\nwbegincode{85}\sublabel{NW4ICsJQ-2aDdLu-A}\nwmargintag{{\nwtagstyle{}\subpageref{NW4ICsJQ-2aDdLu-A}}}\moddef{Axioms for $FS_{0}$~{\nwtagstyle{}\subpageref{NW4ICsJQ-2aDdLu-1}}}\plusendmoddef\nwstartdeflinemarkup\nwusesondefline{\\{NW4ICsJQ-YIPBT-1}}\nwprevnextdefs{NW4ICsJQ-2aDdLu-9}{NW4ICsJQ-2aDdLu-B}\nwenddeflinemarkup
(* Class_ind : (\nwlinkedidentc{Term.t}{NW1rTmOJ-4LsRWC-1} -> \nwlinkedidentc{Formula.t}{NW3w0iRO-2jnj9W-1}) -> \nwlinkedidentc{Thm.t}{NW4ICsJQ-YIPBT-1} *)
fun Class_ind P =
  let
    val tmp = Term.Var("_", \nwlinkedidentc{Sort.CLASS}{NW1rTmOJ-3ZqggC-1});
    val a = fresh_var_in "a" \nwlinkedidentc{Sort.IND}{NW1rTmOJ-3ZqggC-1} (P tmp);
    val f = fresh_var_in "f" \nwlinkedidentc{Sort.FUN}{NW1rTmOJ-3ZqggC-1} (P tmp);
    val A = fresh_var_in "A" \nwlinkedidentc{Sort.CLASS}{NW1rTmOJ-3ZqggC-1} (P tmp);
    val B = fresh_var_in "B" \nwlinkedidentc{Sort.CLASS}{NW1rTmOJ-3ZqggC-1} (P tmp);
    fun ind_case (con : \nwlinkedidentc{Term.t}{NW1rTmOJ-4LsRWC-1} -> \nwlinkedidentc{Term.t}{NW1rTmOJ-4LsRWC-1} -> \nwlinkedidentc{Term.t}{NW1rTmOJ-4LsRWC-1}) =
      \nwlinkedidentc{Formula.mk_forall}{NW3w0iRO-3SlqND-4}(A, \nwlinkedidentc{Formula.mk_forall}{NW3w0iRO-3SlqND-4}(B, \nwlinkedidentc{Formula.mk_imp}{NW3w0iRO-3SlqND-3}(\nwlinkedidentc{Formula.mk_and}{NW3w0iRO-3SlqND-3}(P A,
                                                                              P B),
                                                               P(con A B))));
    val premises =
      List.foldr \nwlinkedidentc{Formula.mk_and}{NW3w0iRO-3SlqND-3}
                 (ind_case (\nwlinkedidentc{FS0.induct}{NW2DqMzY-1mMvkv-1}))
                 [P (\nwlinkedidentc{FS0.singleton}{NW2DqMzY-14kolG-1} (\nwlinkedidentc{FS0.zero}{NW2DqMzY-1NUSQs-1})),
                  \nwlinkedidentc{Formula.mk_forall}{NW3w0iRO-3SlqND-4}(f, \nwlinkedidentc{Formula.mk_forall}{NW3w0iRO-3SlqND-4}(A, \nwlinkedidentc{Formula.mk_imp}{NW3w0iRO-3SlqND-3}(P A,
                                                                           P (\nwlinkedidentc{FS0.preimage}{NW2DqMzY-31WKtJ-1} f A)))),
                  ind_case (\nwlinkedidentc{FS0.union}{NW2DqMzY-22W7o7-1}),
                  ind_case (\nwlinkedidentc{FS0.intersection}{NW2DqMzY-19bkPi-1})];
    val conc = \nwlinkedidentc{Formula.mk_forall}{NW3w0iRO-3SlqND-4}(A, P A);
  in
    mk_axiom (\nwlinkedidentc{Formula.mk_imp}{NW3w0iRO-3SlqND-3}(premises, conc))
  end;
\nwindexdefn{\nwixident{Thm.Class{\_}ind}}{Thm.Class:unind}{NW4ICsJQ-2aDdLu-A}\eatline
\nwused{\\{NW4ICsJQ-YIPBT-1}}\nwidentdefs{\\{{\nwixident{Thm.Class{\_}ind}}{Thm.Class:unind}}}\nwidentuses{\\{{\nwixident{Formula.mk{\_}and}}{Formula.mk:unand}}\\{{\nwixident{Formula.mk{\_}forall}}{Formula.mk:unforall}}\\{{\nwixident{Formula.mk{\_}imp}}{Formula.mk:unimp}}\\{{\nwixident{Formula.t}}{Formula.t}}\\{{\nwixident{FS0.induct}}{FS0.induct}}\\{{\nwixident{FS0.intersection}}{FS0.intersection}}\\{{\nwixident{FS0.preimage}}{FS0.preimage}}\\{{\nwixident{FS0.singleton}}{FS0.singleton}}\\{{\nwixident{FS0.union}}{FS0.union}}\\{{\nwixident{FS0.zero}}{FS0.zero}}\\{{\nwixident{Sort.CLASS}}{Sort.CLASS}}\\{{\nwixident{Sort.FUN}}{Sort.FUN}}\\{{\nwixident{Sort.IND}}{Sort.IND}}\\{{\nwixident{Term.t}}{Term.t}}\\{{\nwixident{Thm.t}}{Thm.t}}}\nwindexuse{\nwixident{Formula.mk{\_}and}}{Formula.mk:unand}{NW4ICsJQ-2aDdLu-A}\nwindexuse{\nwixident{Formula.mk{\_}forall}}{Formula.mk:unforall}{NW4ICsJQ-2aDdLu-A}\nwindexuse{\nwixident{Formula.mk{\_}imp}}{Formula.mk:unimp}{NW4ICsJQ-2aDdLu-A}\nwindexuse{\nwixident{Formula.t}}{Formula.t}{NW4ICsJQ-2aDdLu-A}\nwindexuse{\nwixident{FS0.induct}}{FS0.induct}{NW4ICsJQ-2aDdLu-A}\nwindexuse{\nwixident{FS0.intersection}}{FS0.intersection}{NW4ICsJQ-2aDdLu-A}\nwindexuse{\nwixident{FS0.preimage}}{FS0.preimage}{NW4ICsJQ-2aDdLu-A}\nwindexuse{\nwixident{FS0.singleton}}{FS0.singleton}{NW4ICsJQ-2aDdLu-A}\nwindexuse{\nwixident{FS0.union}}{FS0.union}{NW4ICsJQ-2aDdLu-A}\nwindexuse{\nwixident{FS0.zero}}{FS0.zero}{NW4ICsJQ-2aDdLu-A}\nwindexuse{\nwixident{Sort.CLASS}}{Sort.CLASS}{NW4ICsJQ-2aDdLu-A}\nwindexuse{\nwixident{Sort.FUN}}{Sort.FUN}{NW4ICsJQ-2aDdLu-A}\nwindexuse{\nwixident{Sort.IND}}{Sort.IND}{NW4ICsJQ-2aDdLu-A}\nwindexuse{\nwixident{Term.t}}{Term.t}{NW4ICsJQ-2aDdLu-A}\nwindexuse{\nwixident{Thm.t}}{Thm.t}{NW4ICsJQ-2aDdLu-A}\nwendcode{}\nwbegindocs{86}\nwdocspar
\begin{node}
We will want to ensure we are working with a fresh identifier. Towards
that end, we will just accumulate the identifiers for all predicates,
functions, and terms appearing in a formula, and make sure that the
proposed variable does not appear in them.
\end{node}

\nwenddocs{}\nwbegincode{87}\sublabel{NW4ICsJQ-2auZja-1}\nwmargintag{{\nwtagstyle{}\subpageref{NW4ICsJQ-2auZja-1}}}\moddef{Fresh variable for a formula~{\nwtagstyle{}\subpageref{NW4ICsJQ-2auZja-1}}}\endmoddef\nwstartdeflinemarkup\nwusesondefline{\\{NW4ICsJQ-2aDdLu-8}}\nwenddeflinemarkup
fun fresh_var_in (x : string) (s : \nwlinkedidentc{Sort.t}{NW1rTmOJ-3ZqggC-1}) (fm : \nwlinkedidentc{Formula.t}{NW3w0iRO-2jnj9W-1}) =
  let
    fun tm_ident (tm as (Term.Var _)) = [tm]
      | tm_ident (Term.Fun(f,s0,args)) =
          (Term.Var(f,s0))::(List.concat(map tm_ident args));
    fun fm_ident (P : \nwlinkedidentc{Formula.t}{NW3w0iRO-2jnj9W-1}) =
      if \nwlinkedidentc{Formula.is_pred}{NW3w0iRO-3NyJLv-1} P
      then let val (P0,args) = \nwlinkedidentc{Formula.dest_pred}{NW3w0iRO-1Yt1Jl-1} P
           in (Term.Var(P0,\nwlinkedidentc{Sort.IND}{NW1rTmOJ-3ZqggC-1}))::(List.concat(map tm_ident args)) end
      else [];
    val vars = fm_ident fm;
  in
    \nwlinkedidentc{Term.fresh_var}{NW1rTmOJ-4Ykq9M-1} x s vars
  end;
\nwused{\\{NW4ICsJQ-2aDdLu-8}}\nwidentuses{\\{{\nwixident{Formula.dest{\_}pred}}{Formula.dest:unpred}}\\{{\nwixident{Formula.is{\_}pred}}{Formula.is:unpred}}\\{{\nwixident{Formula.t}}{Formula.t}}\\{{\nwixident{Sort.IND}}{Sort.IND}}\\{{\nwixident{Sort.t}}{Sort.t}}\\{{\nwixident{Term.fresh{\_}var}}{Term.fresh:unvar}}}\nwindexuse{\nwixident{Formula.dest{\_}pred}}{Formula.dest:unpred}{NW4ICsJQ-2auZja-1}\nwindexuse{\nwixident{Formula.is{\_}pred}}{Formula.is:unpred}{NW4ICsJQ-2auZja-1}\nwindexuse{\nwixident{Formula.t}}{Formula.t}{NW4ICsJQ-2auZja-1}\nwindexuse{\nwixident{Sort.IND}}{Sort.IND}{NW4ICsJQ-2auZja-1}\nwindexuse{\nwixident{Sort.t}}{Sort.t}{NW4ICsJQ-2auZja-1}\nwindexuse{\nwixident{Term.fresh{\_}var}}{Term.fresh:unvar}{NW4ICsJQ-2auZja-1}\nwendcode{}\nwbegindocs{88}\nwdocspar

\begin{node}
We also have set comprehension available for $FS_{0}$, but first-order
logic is too weak to prove its existence (we cannot quantify over
formulas, or express as a term parametrized by a formula). Using
any $\exists^{+}$-formula $\phi$, we can form a set $S_{\phi}$ such
that if $FV(\phi)\subset\{x_{1},\dots,x_{n}\}$ contain the free
variables of the formula,
\begin{equation}
\vdash(x_{1},\dots,x_{n})\in S_{\phi}\iff\phi.
\end{equation}
Normally, we write $S_{\phi}$ using set builder notation $\{(x_{1},\dots,x_{n})\mid\phi\}$.
We could, ostensibly, prove this in a fragment of second-order logic,
but that is not available to us. Instead, we need to add this as an
axiom schema for our proof assistant.

Care must be taken so that the list of variables given $[x_{1},\dots,x_{n}]$
is transformed into a tuple with the same order $(x_{1},\dots,x_{n})$.

What should happen when $n=0$ variables are given? This would look
like ``$\{\mid\phi\}$'' where $\phi$ has no free variables (or should not
have any). We can do two things here: we can raise an exception, or we
can return the empty set. I'm not sure which should be preferable.
\end{node}

\nwenddocs{}\nwbegincode{89}\sublabel{NW4ICsJQ-hYCIS-B}\nwmargintag{{\nwtagstyle{}\subpageref{NW4ICsJQ-hYCIS-B}}}\moddef{Specification for $FS_{0}$ axioms~{\nwtagstyle{}\subpageref{NW4ICsJQ-hYCIS-1}}}\plusendmoddef\nwstartdeflinemarkup\nwusesondefline{\\{NW4ICsJQ-2cY0B1-1}}\nwprevnextdefs{NW4ICsJQ-hYCIS-A}{NW4ICsJQ-hYCIS-C}\nwenddeflinemarkup
val axiom_class_comp : \nwlinkedidentc{Formula.t}{NW3w0iRO-2jnj9W-1} -> \nwlinkedidentc{Term.t}{NW1rTmOJ-4LsRWC-1} list -> t;
\nwused{\\{NW4ICsJQ-2cY0B1-1}}\nwidentuses{\\{{\nwixident{Formula.t}}{Formula.t}}\\{{\nwixident{Term.t}}{Term.t}}}\nwindexuse{\nwixident{Formula.t}}{Formula.t}{NW4ICsJQ-hYCIS-B}\nwindexuse{\nwixident{Term.t}}{Term.t}{NW4ICsJQ-hYCIS-B}\nwendcode{}\nwbegindocs{90}\nwdocspar
\nwenddocs{}\nwbegincode{91}\sublabel{NW4ICsJQ-2aDdLu-B}\nwmargintag{{\nwtagstyle{}\subpageref{NW4ICsJQ-2aDdLu-B}}}\moddef{Axioms for $FS_{0}$~{\nwtagstyle{}\subpageref{NW4ICsJQ-2aDdLu-1}}}\plusendmoddef\nwstartdeflinemarkup\nwusesondefline{\\{NW4ICsJQ-YIPBT-1}}\nwprevnextdefs{NW4ICsJQ-2aDdLu-A}{NW4ICsJQ-2aDdLu-C}\nwenddeflinemarkup
(* axiom_class_comp : \nwlinkedidentc{Formula.t}{NW3w0iRO-2jnj9W-1} -> \nwlinkedidentc{Term.t}{NW1rTmOJ-4LsRWC-1} list -> \nwlinkedidentc{Thm.t}{NW4ICsJQ-YIPBT-1} *)
fun axiom_class_comp (P : \nwlinkedidentc{Formula.t}{NW3w0iRO-2jnj9W-1}) ([] : \nwlinkedidentc{Term.t}{NW1rTmOJ-4LsRWC-1} list) =
      raise Fail "axiom_class_comp: variable list is empty"
  | axiom_class_comp P ((vars as (h::hs)) : \nwlinkedidentc{Term.t}{NW1rTmOJ-4LsRWC-1} list) =
      let
        fun not_in_vars y = List.all (not o (\nwlinkedidentc{Term.eq}{NW1rTmOJ-3810sY-1} y)) vars;
        val xs : \nwlinkedidentc{Term.t}{NW1rTmOJ-4LsRWC-1} = List.foldl (fn (x,acc) => \nwlinkedidentc{FS0.pair}{NW2DqMzY-4SkDla-1} acc x)
                                     h
                                     hs;
        val S = Term.Fun("S_\{("^
                         (String.concatWith "," (map (fn (Term.Var(x,_)) => x)
                                                     vars))
                         ^")|"^(\nwlinkedidentc{Formula.serialize}{NW3w0iRO-3H9tUA-1} P)^"\}",
                         \nwlinkedidentc{Sort.CLASS}{NW1rTmOJ-3ZqggC-1},
                         []);
      in
        if not(\nwlinkedidentc{Formula.is_e_plus}{NW3w0iRO-1O1TnJ-1} P)
        then raise Fail("axiom_class_comp: formula is not E+")
        else if List.exists not_in_vars (\nwlinkedidentc{Formula.fv}{NW3w0iRO-jRgYZ-2} P)
        then raise Fail("axiom_class_comp: not all free variables of predicate "^
                        "contained in given list of variables")
        else if List.exists (not o Term.is_var) vars
        then raise Fail "axiom_class_comp: given list of variables contains a function"
        else if List.exists (not o \nwlinkedidentc{Term.is_ind}{NW1rTmOJ-3FfXTO-1}) vars
        then raise Fail "axiom_class_comp: given list of variables not all individual-sorted"
        else [] |- (List.foldr \nwlinkedidentc{Formula.mk_forall}{NW3w0iRO-3SlqND-4}
                               (\nwlinkedidentc{Formula.mk_iff}{NW3w0iRO-3SlqND-3} (\nwlinkedidentc{FS0.In}{NW2DqMzY-eTij6-1} xs S,
                                                P))
                               vars)
      end;
\nwindexdefn{\nwixident{Thm.axiom{\_}class{\_}comp}}{Thm.axiom:unclass:uncomp}{NW4ICsJQ-2aDdLu-B}\eatline
\nwused{\\{NW4ICsJQ-YIPBT-1}}\nwidentdefs{\\{{\nwixident{Thm.axiom{\_}class{\_}comp}}{Thm.axiom:unclass:uncomp}}}\nwidentuses{\\{{\nwixident{Formula.fv}}{Formula.fv}}\\{{\nwixident{Formula.is{\_}e{\_}plus}}{Formula.is:une:unplus}}\\{{\nwixident{Formula.mk{\_}forall}}{Formula.mk:unforall}}\\{{\nwixident{Formula.mk{\_}iff}}{Formula.mk:uniff}}\\{{\nwixident{Formula.serialize}}{Formula.serialize}}\\{{\nwixident{Formula.t}}{Formula.t}}\\{{\nwixident{FS0.In}}{FS0.In}}\\{{\nwixident{FS0.pair}}{FS0.pair}}\\{{\nwixident{Sort.CLASS}}{Sort.CLASS}}\\{{\nwixident{Term.eq}}{Term.eq}}\\{{\nwixident{Term.is{\_}ind}}{Term.is:unind}}\\{{\nwixident{Term.t}}{Term.t}}\\{{\nwixident{Thm.t}}{Thm.t}}}\nwindexuse{\nwixident{Formula.fv}}{Formula.fv}{NW4ICsJQ-2aDdLu-B}\nwindexuse{\nwixident{Formula.is{\_}e{\_}plus}}{Formula.is:une:unplus}{NW4ICsJQ-2aDdLu-B}\nwindexuse{\nwixident{Formula.mk{\_}forall}}{Formula.mk:unforall}{NW4ICsJQ-2aDdLu-B}\nwindexuse{\nwixident{Formula.mk{\_}iff}}{Formula.mk:uniff}{NW4ICsJQ-2aDdLu-B}\nwindexuse{\nwixident{Formula.serialize}}{Formula.serialize}{NW4ICsJQ-2aDdLu-B}\nwindexuse{\nwixident{Formula.t}}{Formula.t}{NW4ICsJQ-2aDdLu-B}\nwindexuse{\nwixident{FS0.In}}{FS0.In}{NW4ICsJQ-2aDdLu-B}\nwindexuse{\nwixident{FS0.pair}}{FS0.pair}{NW4ICsJQ-2aDdLu-B}\nwindexuse{\nwixident{Sort.CLASS}}{Sort.CLASS}{NW4ICsJQ-2aDdLu-B}\nwindexuse{\nwixident{Term.eq}}{Term.eq}{NW4ICsJQ-2aDdLu-B}\nwindexuse{\nwixident{Term.is{\_}ind}}{Term.is:unind}{NW4ICsJQ-2aDdLu-B}\nwindexuse{\nwixident{Term.t}}{Term.t}{NW4ICsJQ-2aDdLu-B}\nwindexuse{\nwixident{Thm.t}}{Thm.t}{NW4ICsJQ-2aDdLu-B}\nwendcode{}\nwbegindocs{92}\nwdocspar
\begin{node}
We have extensionality axioms telling us how to prove two functions
are equal to each other, and how two classes are equal to each other.
Functions $f$ and $g$ are equal if and only if for any
individual-sorted term $x$ we have $fx=gx$.
\end{node}
\nwenddocs{}\nwbegincode{93}\sublabel{NW4ICsJQ-hYCIS-C}\nwmargintag{{\nwtagstyle{}\subpageref{NW4ICsJQ-hYCIS-C}}}\moddef{Specification for $FS_{0}$ axioms~{\nwtagstyle{}\subpageref{NW4ICsJQ-hYCIS-1}}}\plusendmoddef\nwstartdeflinemarkup\nwusesondefline{\\{NW4ICsJQ-2cY0B1-1}}\nwprevnextdefs{NW4ICsJQ-hYCIS-B}{NW4ICsJQ-hYCIS-D}\nwenddeflinemarkup
val axiom_fun_eq : \nwlinkedidentc{Term.t}{NW1rTmOJ-4LsRWC-1} -> \nwlinkedidentc{Term.t}{NW1rTmOJ-4LsRWC-1} -> t;
\nwused{\\{NW4ICsJQ-2cY0B1-1}}\nwidentuses{\\{{\nwixident{Term.t}}{Term.t}}}\nwindexuse{\nwixident{Term.t}}{Term.t}{NW4ICsJQ-hYCIS-C}\nwendcode{}\nwbegindocs{94}\nwdocspar
\nwenddocs{}\nwbegincode{95}\sublabel{NW4ICsJQ-2aDdLu-C}\nwmargintag{{\nwtagstyle{}\subpageref{NW4ICsJQ-2aDdLu-C}}}\moddef{Axioms for $FS_{0}$~{\nwtagstyle{}\subpageref{NW4ICsJQ-2aDdLu-1}}}\plusendmoddef\nwstartdeflinemarkup\nwusesondefline{\\{NW4ICsJQ-YIPBT-1}}\nwprevnextdefs{NW4ICsJQ-2aDdLu-B}{NW4ICsJQ-2aDdLu-D}\nwenddeflinemarkup
(* axiom_fun_eq : \nwlinkedidentc{Term.t}{NW1rTmOJ-4LsRWC-1} -> \nwlinkedidentc{Term.t}{NW1rTmOJ-4LsRWC-1} -> \nwlinkedidentc{Thm.t}{NW4ICsJQ-YIPBT-1} *)
fun axiom_fun_eq lhs rhs =
  if not(\nwlinkedidentc{Term.is_fun}{NW1rTmOJ-3FfXTO-1} lhs)
  then raise Fail "axiom_fun_eq: left-hand side is not a function"
  else if not(\nwlinkedidentc{Term.is_fun}{NW1rTmOJ-3FfXTO-1} rhs)
  then raise Fail "axiom_fun_eq: right-hand side is not a function"
  else let val vars = (\nwlinkedidentc{Term.fv}{NW1rTmOJ-2nRtky-1} lhs) @ (\nwlinkedidentc{Term.fv}{NW1rTmOJ-2nRtky-1} rhs);
           val x = \nwlinkedidentc{Term.fresh_var}{NW1rTmOJ-4Ykq9M-1} "x" \nwlinkedidentc{Sort.IND}{NW1rTmOJ-3ZqggC-1} vars;
           val iff = \nwlinkedidentc{Formula.mk_iff}{NW3w0iRO-3SlqND-3};
           val forall = \nwlinkedidentc{Formula.mk_forall}{NW3w0iRO-3SlqND-4};
       in [] |- iff(forall(x, \nwlinkedidentc{FS0.equal}{NW2DqMzY-ocj48-1} (\nwlinkedidentc{FS0.app}{NW2DqMzY-2oGoBx-1} lhs x)
                                        (\nwlinkedidentc{FS0.app}{NW2DqMzY-2oGoBx-1} rhs x)),
                    \nwlinkedidentc{FS0.equal}{NW2DqMzY-ocj48-1} lhs rhs)
       end;
\nwindexdefn{\nwixident{Thm.axiom{\_}fun{\_}eq}}{Thm.axiom:unfun:uneq}{NW4ICsJQ-2aDdLu-C}\eatline
\nwused{\\{NW4ICsJQ-YIPBT-1}}\nwidentdefs{\\{{\nwixident{Thm.axiom{\_}fun{\_}eq}}{Thm.axiom:unfun:uneq}}}\nwidentuses{\\{{\nwixident{Formula.mk{\_}forall}}{Formula.mk:unforall}}\\{{\nwixident{Formula.mk{\_}iff}}{Formula.mk:uniff}}\\{{\nwixident{FS0.app}}{FS0.app}}\\{{\nwixident{FS0.equal}}{FS0.equal}}\\{{\nwixident{Sort.IND}}{Sort.IND}}\\{{\nwixident{Term.fresh{\_}var}}{Term.fresh:unvar}}\\{{\nwixident{Term.fv}}{Term.fv}}\\{{\nwixident{Term.is{\_}fun}}{Term.is:unfun}}\\{{\nwixident{Term.t}}{Term.t}}\\{{\nwixident{Thm.t}}{Thm.t}}}\nwindexuse{\nwixident{Formula.mk{\_}forall}}{Formula.mk:unforall}{NW4ICsJQ-2aDdLu-C}\nwindexuse{\nwixident{Formula.mk{\_}iff}}{Formula.mk:uniff}{NW4ICsJQ-2aDdLu-C}\nwindexuse{\nwixident{FS0.app}}{FS0.app}{NW4ICsJQ-2aDdLu-C}\nwindexuse{\nwixident{FS0.equal}}{FS0.equal}{NW4ICsJQ-2aDdLu-C}\nwindexuse{\nwixident{Sort.IND}}{Sort.IND}{NW4ICsJQ-2aDdLu-C}\nwindexuse{\nwixident{Term.fresh{\_}var}}{Term.fresh:unvar}{NW4ICsJQ-2aDdLu-C}\nwindexuse{\nwixident{Term.fv}}{Term.fv}{NW4ICsJQ-2aDdLu-C}\nwindexuse{\nwixident{Term.is{\_}fun}}{Term.is:unfun}{NW4ICsJQ-2aDdLu-C}\nwindexuse{\nwixident{Term.t}}{Term.t}{NW4ICsJQ-2aDdLu-C}\nwindexuse{\nwixident{Thm.t}}{Thm.t}{NW4ICsJQ-2aDdLu-C}\nwendcode{}\nwbegindocs{96}\begin{node}
Class extensionality should be far more familiar to Mathematicians:
two classes are equal if and only if they contain the same elements.
\end{node}
\nwenddocs{}\nwbegincode{97}\sublabel{NW4ICsJQ-hYCIS-D}\nwmargintag{{\nwtagstyle{}\subpageref{NW4ICsJQ-hYCIS-D}}}\moddef{Specification for $FS_{0}$ axioms~{\nwtagstyle{}\subpageref{NW4ICsJQ-hYCIS-1}}}\plusendmoddef\nwstartdeflinemarkup\nwusesondefline{\\{NW4ICsJQ-2cY0B1-1}}\nwprevnextdefs{NW4ICsJQ-hYCIS-C}{\relax}\nwenddeflinemarkup
val axiom_class_eq : \nwlinkedidentc{Term.t}{NW1rTmOJ-4LsRWC-1} -> \nwlinkedidentc{Term.t}{NW1rTmOJ-4LsRWC-1} -> t;
\nwused{\\{NW4ICsJQ-2cY0B1-1}}\nwidentuses{\\{{\nwixident{Term.t}}{Term.t}}}\nwindexuse{\nwixident{Term.t}}{Term.t}{NW4ICsJQ-hYCIS-D}\nwendcode{}\nwbegindocs{98}\nwdocspar
\nwenddocs{}\nwbegincode{99}\sublabel{NW4ICsJQ-2aDdLu-D}\nwmargintag{{\nwtagstyle{}\subpageref{NW4ICsJQ-2aDdLu-D}}}\moddef{Axioms for $FS_{0}$~{\nwtagstyle{}\subpageref{NW4ICsJQ-2aDdLu-1}}}\plusendmoddef\nwstartdeflinemarkup\nwusesondefline{\\{NW4ICsJQ-YIPBT-1}}\nwprevnextdefs{NW4ICsJQ-2aDdLu-C}{\relax}\nwenddeflinemarkup
(* axiom_class_eq : \nwlinkedidentc{Term.t}{NW1rTmOJ-4LsRWC-1} -> \nwlinkedidentc{Term.t}{NW1rTmOJ-4LsRWC-1} -> \nwlinkedidentc{Thm.t}{NW4ICsJQ-YIPBT-1} *)
fun axiom_class_eq lhs rhs =
  if not(\nwlinkedidentc{Term.is_class}{NW1rTmOJ-3FfXTO-1} lhs)
  then raise Fail "axiom_class_eq: left-hand side is not a class"
  else if not(\nwlinkedidentc{Term.is_class}{NW1rTmOJ-3FfXTO-1} rhs)
  then raise Fail "axiom_class_eq: right-hand side is not a class"
  else let val vars = (\nwlinkedidentc{Term.fv}{NW1rTmOJ-2nRtky-1} lhs) @ (\nwlinkedidentc{Term.fv}{NW1rTmOJ-2nRtky-1} rhs);
           val x = \nwlinkedidentc{Term.fresh_var}{NW1rTmOJ-4Ykq9M-1} "x" \nwlinkedidentc{Sort.IND}{NW1rTmOJ-3ZqggC-1} vars;
           val iff = \nwlinkedidentc{Formula.mk_iff}{NW3w0iRO-3SlqND-3};
           val forall = \nwlinkedidentc{Formula.mk_forall}{NW3w0iRO-3SlqND-4};
       in [] |- iff(forall(x, \nwlinkedidentc{Formula.mk_iff}{NW3w0iRO-3SlqND-3} (\nwlinkedidentc{FS0.In}{NW2DqMzY-eTij6-1} x lhs,
                                              \nwlinkedidentc{FS0.In}{NW2DqMzY-eTij6-1} x rhs)),
                    (\nwlinkedidentc{FS0.equal}{NW2DqMzY-ocj48-1} lhs rhs))
       end;
\nwindexdefn{\nwixident{Thm.axiom{\_}class{\_}eq}}{Thm.axiom:unclass:uneq}{NW4ICsJQ-2aDdLu-D}\eatline
\nwused{\\{NW4ICsJQ-YIPBT-1}}\nwidentdefs{\\{{\nwixident{Thm.axiom{\_}class{\_}eq}}{Thm.axiom:unclass:uneq}}}\nwidentuses{\\{{\nwixident{Formula.mk{\_}forall}}{Formula.mk:unforall}}\\{{\nwixident{Formula.mk{\_}iff}}{Formula.mk:uniff}}\\{{\nwixident{FS0.equal}}{FS0.equal}}\\{{\nwixident{FS0.In}}{FS0.In}}\\{{\nwixident{Sort.IND}}{Sort.IND}}\\{{\nwixident{Term.fresh{\_}var}}{Term.fresh:unvar}}\\{{\nwixident{Term.fv}}{Term.fv}}\\{{\nwixident{Term.is{\_}class}}{Term.is:unclass}}\\{{\nwixident{Term.t}}{Term.t}}\\{{\nwixident{Thm.t}}{Thm.t}}}\nwindexuse{\nwixident{Formula.mk{\_}forall}}{Formula.mk:unforall}{NW4ICsJQ-2aDdLu-D}\nwindexuse{\nwixident{Formula.mk{\_}iff}}{Formula.mk:uniff}{NW4ICsJQ-2aDdLu-D}\nwindexuse{\nwixident{FS0.equal}}{FS0.equal}{NW4ICsJQ-2aDdLu-D}\nwindexuse{\nwixident{FS0.In}}{FS0.In}{NW4ICsJQ-2aDdLu-D}\nwindexuse{\nwixident{Sort.IND}}{Sort.IND}{NW4ICsJQ-2aDdLu-D}\nwindexuse{\nwixident{Term.fresh{\_}var}}{Term.fresh:unvar}{NW4ICsJQ-2aDdLu-D}\nwindexuse{\nwixident{Term.fv}}{Term.fv}{NW4ICsJQ-2aDdLu-D}\nwindexuse{\nwixident{Term.is{\_}class}}{Term.is:unclass}{NW4ICsJQ-2aDdLu-D}\nwindexuse{\nwixident{Term.t}}{Term.t}{NW4ICsJQ-2aDdLu-D}\nwindexuse{\nwixident{Thm.t}}{Thm.t}{NW4ICsJQ-2aDdLu-D}\nwendcode{}\nwbegindocs{100}\nwdocspar
\begin{problem}[Severe]
The sentence ``$\forall x\ldotp0=x\lor\bigl(\exists y,z\ldotp x=(y,z)\bigr)$''
is independent of the axioms of $FS_{0}$.
\end{problem}

This is more critical than I think most people appreciate, because the
entire axiomatization of $FS_{0}$ is predicated on the premise that
every individual-sorted term is either $0$ or else it's an ordered
pair of individual terms. With that implicit premise removed, you can
end up in very unpleasant situations with inductive proofs. But you'd
only realize this if you are foolish enough to use $FS_{0}$ as the
basis for a proof assistant\dots

Of course, this only matters if you are trying to prove a
claim of the form $\forall x\ldotp P[x]$ where $P[-]$ is a
predicate. In practice, this does not really matter since we end up
proving claims of the form,
\begin{equation*}
\forall x\ldotp x\in C\implies P[x].
\end{equation*}
This makes the axiom of the universe useful for such purposes.

\begin{problem}[Minor]
The projection functions are ambiguous on $0$, i.e., $P_{1}0$ and
$P_{2}0$ are not specified but should be (and a good choice would be
$P_{1}0=P_{2}0=0$). 
\end{problem}

\nwenddocs{}\nwfilename{nw/fs0-syntax.nw}\nwbegindocs{0}% -*- mode: poly-noweb; noweb-code-mode: sml-mode; -*-
%%
%% fs0-syntax.nw
%% 
%% Made by Alex Nelson <pqnelson@gmail.com>
%% Login   <alex@lisp>
%% 
%% Started on  2026-04-06T18:37:15-0700
%% Last update 2026-04-06T18:37:15-0700
%% 

\section{Syntactic constructors for $FS_{0}$ primitives}

\begin{node}
I have decided to collect the constructors for all the primitive
notions for $FS_{0}$ in one location. The situation we wish to avoid
is to change conventions later on (e.g., instead of use ``P1'' for the
name of the projection function we might want to use ``car'', ``fst'',
or something else entirely, so we should have one line of code which we
need to change as opposed to dozens).
\end{node}

\nwenddocs{}\nwbegincode{1}\sublabel{NW2DqMzY-3R3yGC-1}\nwmargintag{{\nwtagstyle{}\subpageref{NW2DqMzY-3R3yGC-1}}}\moddef{FS0.sig~{\nwtagstyle{}\subpageref{NW2DqMzY-3R3yGC-1}}}\endmoddef\nwstartdeflinemarkup\nwenddeflinemarkup
signature FS0 = sig
  exception Dest of string;
  (* individuals *)
  \LA{}Specification for $FS_{0}$ individuals~{\nwtagstyle{}\subpageref{NW2DqMzY-24mg92-1}}\RA{}

  (* functions *)
  \LA{}Specification for $FS_{0}$ functions~{\nwtagstyle{}\subpageref{NW2DqMzY-1uF0T3-1}}\RA{}

  (* classes *)
  \LA{}Specification for $FS_{0}$ classes~{\nwtagstyle{}\subpageref{NW2DqMzY-33zuOn-1}}\RA{}

  (* predicates *)
  \LA{}Specification for $FS_{0}$ predicates~{\nwtagstyle{}\subpageref{NW2DqMzY-194t0N-1}}\RA{}
end;
\nwnotused{FS0.sig}\nwendcode{}\nwbegindocs{2}\nwdocspar

\nwenddocs{}\nwbegincode{3}\sublabel{NW2DqMzY-mVmU8-1}\nwmargintag{{\nwtagstyle{}\subpageref{NW2DqMzY-mVmU8-1}}}\moddef{FS0.sml~{\nwtagstyle{}\subpageref{NW2DqMzY-mVmU8-1}}}\endmoddef\nwstartdeflinemarkup\nwenddeflinemarkup
structure FS0 : FS0 = struct
exception Dest of string;

\LA{}Constructors for $FS_{0}$ individuals~{\nwtagstyle{}\subpageref{NW2DqMzY-1NUSQs-1}}\RA{}

\LA{}Constructors for $FS_{0}$ functions~{\nwtagstyle{}\subpageref{NW2DqMzY-3aj0iq-1}}\RA{}

\LA{}Constructors for $FS_{0}$ classes~{\nwtagstyle{}\subpageref{NW2DqMzY-TSvER-1}}\RA{}

\LA{}Constructors for $FS_{0}$ predicates~{\nwtagstyle{}\subpageref{NW2DqMzY-4OjF2N-1}}\RA{}
end;
\nwnotused{FS0.sml}\nwendcode{}\nwbegindocs{4}\nwdocspar

\section*{Constructors for individuals}

\begin{node}
The syntax for individual-sorted terms in $FS_{0}$ is the smallest
inductively-defined set such that
\begin{enumerate}
\item Individual-sorted variables are individual-sorted terms;
\item The constant term $0$ is individual-sorted;
\item If $f$ is a function-sorted term and $t$ is an individual-sorted
  term, then the application $ft$ is an individual-sorted term;
\item If $t_{1}$ and $t_{2}$ are individual-sorted terms, then their
  ordered pair $(t_{1},t_{2})$ is an individual-sorted term;
\item And nothing else.
\end{enumerate}
We provide {\Tt{}zero\nwendquote} as a constant term, as well as constructors and
destructors for ordered {\Tt{}pair\nwendquote}s and function {\Tt{}app\nwendquote}lications.
\end{node}

\nwenddocs{}\nwbegincode{5}\sublabel{NW2DqMzY-24mg92-1}\nwmargintag{{\nwtagstyle{}\subpageref{NW2DqMzY-24mg92-1}}}\moddef{Specification for $FS_{0}$ individuals~{\nwtagstyle{}\subpageref{NW2DqMzY-24mg92-1}}}\endmoddef\nwstartdeflinemarkup\nwusesondefline{\\{NW2DqMzY-3R3yGC-1}}\nwenddeflinemarkup
val zero : \nwlinkedidentc{Term.t}{NW1rTmOJ-4LsRWC-1};
val pair : \nwlinkedidentc{Term.t}{NW1rTmOJ-4LsRWC-1} -> \nwlinkedidentc{Term.t}{NW1rTmOJ-4LsRWC-1} -> \nwlinkedidentc{Term.t}{NW1rTmOJ-4LsRWC-1};
val app : \nwlinkedidentc{Term.t}{NW1rTmOJ-4LsRWC-1} -> \nwlinkedidentc{Term.t}{NW1rTmOJ-4LsRWC-1} -> \nwlinkedidentc{Term.t}{NW1rTmOJ-4LsRWC-1};

val dest_app : \nwlinkedidentc{Term.t}{NW1rTmOJ-4LsRWC-1} -> \nwlinkedidentc{Term.t}{NW1rTmOJ-4LsRWC-1} * \nwlinkedidentc{Term.t}{NW1rTmOJ-4LsRWC-1};
val dest_pair : \nwlinkedidentc{Term.t}{NW1rTmOJ-4LsRWC-1} -> \nwlinkedidentc{Term.t}{NW1rTmOJ-4LsRWC-1} * \nwlinkedidentc{Term.t}{NW1rTmOJ-4LsRWC-1};
\nwused{\\{NW2DqMzY-3R3yGC-1}}\nwidentuses{\\{{\nwixident{Term.t}}{Term.t}}}\nwindexuse{\nwixident{Term.t}}{Term.t}{NW2DqMzY-24mg92-1}\nwendcode{}\nwbegindocs{6}\nwdocspar

\nwenddocs{}\nwbegincode{7}\sublabel{NW2DqMzY-1NUSQs-1}\nwmargintag{{\nwtagstyle{}\subpageref{NW2DqMzY-1NUSQs-1}}}\moddef{Constructors for $FS_{0}$ individuals~{\nwtagstyle{}\subpageref{NW2DqMzY-1NUSQs-1}}}\endmoddef\nwstartdeflinemarkup\nwusesondefline{\\{NW2DqMzY-mVmU8-1}}\nwenddeflinemarkup
val zero = Term.Fun("0", \nwlinkedidentc{Sort.IND}{NW1rTmOJ-3ZqggC-1}, []);

\LA{}Construct an ordered pair from $FS_{0}$ individuals~{\nwtagstyle{}\subpageref{NW2DqMzY-4SkDla-1}}\RA{}

\LA{}Destructuring \code{}cons\edoc{}~{\nwtagstyle{}\subpageref{NW2DqMzY-21ZNyJ-1}}\RA{}

\LA{}Primitive \code{}apply\edoc{} constructor for functions to individuals~{\nwtagstyle{}\subpageref{NW2DqMzY-2oGoBx-1}}\RA{}

\LA{}Destructuring \code{}apply\edoc{}~{\nwtagstyle{}\subpageref{NW2DqMzY-1VpWgD-1}}\RA{}
\nwindexdefn{\nwixident{FS0.zero}}{FS0.zero}{NW2DqMzY-1NUSQs-1}\eatline
\nwused{\\{NW2DqMzY-mVmU8-1}}\nwidentdefs{\\{{\nwixident{FS0.zero}}{FS0.zero}}}\nwidentuses{\\{{\nwixident{Sort.IND}}{Sort.IND}}}\nwindexuse{\nwixident{Sort.IND}}{Sort.IND}{NW2DqMzY-1NUSQs-1}\nwendcode{}\nwbegindocs{8}\nwdocspar
\begin{node}
We should be aware and cogniscent that function application is a major
source of bad performance for $FS_{0}$ as an axiomatic system. Some
functions can be ``optimized away''. We have done this for {\Tt{}proj1\nwendquote}
and {\Tt{}proj2\nwendquote} and {\Tt{}id\nwendquote}. But {\Tt{}const\nwendquote}, {\Tt{}comp\nwendquote}, {\Tt{}juxt\nwendquote}, {\Tt{}recur\nwendquote}
have not been optimized in any sense. The {\Tt{}\nwlinkedidentq{FS0.app}{NW2DqMzY-2oGoBx-1}\nwendquote} function can be
used to perform these optimizations.

If we want to perform these optimizations, then we should replace the
default branch with a {\Tt{}case\nwendquote} expression which dispatches based on
the function {\Tt{}f\nwendquote} and argument {\Tt{}tm\nwendquote} supplied.
\end{node}

\nwenddocs{}\nwbegincode{9}\sublabel{NW2DqMzY-2oGoBx-1}\nwmargintag{{\nwtagstyle{}\subpageref{NW2DqMzY-2oGoBx-1}}}\moddef{Primitive \code{}apply\edoc{} constructor for functions to individuals~{\nwtagstyle{}\subpageref{NW2DqMzY-2oGoBx-1}}}\endmoddef\nwstartdeflinemarkup\nwusesondefline{\\{NW2DqMzY-1NUSQs-1}}\nwenddeflinemarkup
(* \nwlinkedidentc{FS0.app}{NW2DqMzY-2oGoBx-1} : \nwlinkedidentc{Term.t}{NW1rTmOJ-4LsRWC-1} -> \nwlinkedidentc{Term.t}{NW1rTmOJ-4LsRWC-1} -> \nwlinkedidentc{Term.t}{NW1rTmOJ-4LsRWC-1} *)
fun app f tm =
  if not(\nwlinkedidentc{Term.is_fun}{NW1rTmOJ-3FfXTO-1} f)
  then raise Sort.Mismatch "FS0.apply: first argument is not a function"
  else if not(\nwlinkedidentc{Term.is_ind}{NW1rTmOJ-3FfXTO-1} tm)
  then raise Sort.Mismatch "FS0.apply: second argument is not an individual"
  else Term.Fun("apply", \nwlinkedidentc{Sort.IND}{NW1rTmOJ-3ZqggC-1}, [f, tm]);
\nwindexdefn{\nwixident{FS0.app}}{FS0.app}{NW2DqMzY-2oGoBx-1}\eatline
\nwused{\\{NW2DqMzY-1NUSQs-1}}\nwidentdefs{\\{{\nwixident{FS0.app}}{FS0.app}}}\nwidentuses{\\{{\nwixident{Sort.IND}}{Sort.IND}}\\{{\nwixident{Term.is{\_}fun}}{Term.is:unfun}}\\{{\nwixident{Term.is{\_}ind}}{Term.is:unind}}\\{{\nwixident{Term.t}}{Term.t}}}\nwindexuse{\nwixident{Sort.IND}}{Sort.IND}{NW2DqMzY-2oGoBx-1}\nwindexuse{\nwixident{Term.is{\_}fun}}{Term.is:unfun}{NW2DqMzY-2oGoBx-1}\nwindexuse{\nwixident{Term.is{\_}ind}}{Term.is:unind}{NW2DqMzY-2oGoBx-1}\nwindexuse{\nwixident{Term.t}}{Term.t}{NW2DqMzY-2oGoBx-1}\nwendcode{}\nwbegindocs{10}\nwdocspar
\nwenddocs{}\nwbegincode{11}\sublabel{NW2DqMzY-1VpWgD-1}\nwmargintag{{\nwtagstyle{}\subpageref{NW2DqMzY-1VpWgD-1}}}\moddef{Destructuring \code{}apply\edoc{}~{\nwtagstyle{}\subpageref{NW2DqMzY-1VpWgD-1}}}\endmoddef\nwstartdeflinemarkup\nwusesondefline{\\{NW2DqMzY-1NUSQs-1}}\nwenddeflinemarkup
fun dest_app (Term.Fun("apply", \nwlinkedidentc{Sort.IND}{NW1rTmOJ-3ZqggC-1}, [f, tm])) = (f, tm)
  | dest_app tm = raise Dest "\nwlinkedidentc{FS0.dest_app}{NW2DqMzY-1VpWgD-1}";
\nwindexdefn{\nwixident{FS0.dest{\_}app}}{FS0.dest:unapp}{NW2DqMzY-1VpWgD-1}\eatline
\nwused{\\{NW2DqMzY-1NUSQs-1}}\nwidentdefs{\\{{\nwixident{FS0.dest{\_}app}}{FS0.dest:unapp}}}\nwidentuses{\\{{\nwixident{Sort.IND}}{Sort.IND}}}\nwindexuse{\nwixident{Sort.IND}}{Sort.IND}{NW2DqMzY-1VpWgD-1}\nwendcode{}\nwbegindocs{12}\nwdocspar
\begin{node}
Constructing an ordered pair is a primitive operation for individuals
in $FS_{0}$. We need to raise an exception if the arguments supplied
are not {\Tt{}IND\nwendquote} sorts.
\end{node}

\nwenddocs{}\nwbegincode{13}\sublabel{NW2DqMzY-4SkDla-1}\nwmargintag{{\nwtagstyle{}\subpageref{NW2DqMzY-4SkDla-1}}}\moddef{Construct an ordered pair from $FS_{0}$ individuals~{\nwtagstyle{}\subpageref{NW2DqMzY-4SkDla-1}}}\endmoddef\nwstartdeflinemarkup\nwusesondefline{\\{NW2DqMzY-1NUSQs-1}}\nwenddeflinemarkup
(* pair : \nwlinkedidentc{Term.t}{NW1rTmOJ-4LsRWC-1} -> \nwlinkedidentc{Term.t}{NW1rTmOJ-4LsRWC-1} -> \nwlinkedidentc{Term.t}{NW1rTmOJ-4LsRWC-1} *)
fun pair t1 t2 =
  if not(\nwlinkedidentc{Term.is_ind}{NW1rTmOJ-3FfXTO-1} t1)
  then raise Sort.Mismatch "\nwlinkedidentc{FS0.pair}{NW2DqMzY-4SkDla-1}: first term expected to be an individual"
  else if not(\nwlinkedidentc{Term.is_ind}{NW1rTmOJ-3FfXTO-1} t2)
  then raise Sort.Mismatch "\nwlinkedidentc{FS0.pair}{NW2DqMzY-4SkDla-1}: first term expected to be an individual"
  else Term.Fun("cons", \nwlinkedidentc{Sort.IND}{NW1rTmOJ-3ZqggC-1}, [t1, t2]);
\nwindexdefn{\nwixident{FS0.pair}}{FS0.pair}{NW2DqMzY-4SkDla-1}\eatline
\nwused{\\{NW2DqMzY-1NUSQs-1}}\nwidentdefs{\\{{\nwixident{FS0.pair}}{FS0.pair}}}\nwidentuses{\\{{\nwixident{Sort.IND}}{Sort.IND}}\\{{\nwixident{Term.is{\_}ind}}{Term.is:unind}}\\{{\nwixident{Term.t}}{Term.t}}}\nwindexuse{\nwixident{Sort.IND}}{Sort.IND}{NW2DqMzY-4SkDla-1}\nwindexuse{\nwixident{Term.is{\_}ind}}{Term.is:unind}{NW2DqMzY-4SkDla-1}\nwindexuse{\nwixident{Term.t}}{Term.t}{NW2DqMzY-4SkDla-1}\nwendcode{}\nwbegindocs{14}\nwdocspar
\nwenddocs{}\nwbegincode{15}\sublabel{NW2DqMzY-21ZNyJ-1}\nwmargintag{{\nwtagstyle{}\subpageref{NW2DqMzY-21ZNyJ-1}}}\moddef{Destructuring \code{}cons\edoc{}~{\nwtagstyle{}\subpageref{NW2DqMzY-21ZNyJ-1}}}\endmoddef\nwstartdeflinemarkup\nwusesondefline{\\{NW2DqMzY-1NUSQs-1}}\nwenddeflinemarkup
fun dest_pair (Term.Fun("cons", \nwlinkedidentc{Sort.IND}{NW1rTmOJ-3ZqggC-1}, [t1, t2])) = (t1, t2)
  | dest_pair _ = raise Dest "\nwlinkedidentc{FS0.dest_pair}{NW2DqMzY-21ZNyJ-1}";
\nwindexdefn{\nwixident{FS0.dest{\_}pair}}{FS0.dest:unpair}{NW2DqMzY-21ZNyJ-1}\eatline
\nwused{\\{NW2DqMzY-1NUSQs-1}}\nwidentdefs{\\{{\nwixident{FS0.dest{\_}pair}}{FS0.dest:unpair}}}\nwidentuses{\\{{\nwixident{Sort.IND}}{Sort.IND}}}\nwindexuse{\nwixident{Sort.IND}}{Sort.IND}{NW2DqMzY-21ZNyJ-1}\nwendcode{}\nwbegindocs{16}\nwdocspar
\section*{Constructors for functions}

\begin{node}
The set $V_{0}$ of all individual-sorted terms of $FS_{0}$ constitutes
one of the domains of discourse. The next are primitive recursive functions $f\colon V_{0}\to V_{0}$
which are given by the primitive function-sorted term constructors:
\begin{enumerate}
\item The identity function $\id(x)=x$
\item The constant function $k_{a}(x)=a$
\item Projection functions $P_{1}(x,y)=x$ and $P_{2}(x,y)=y$
\item Deciding if two terms are equal
  \begin{equation*}
D(x_{1},x_{2},y_{1},y_{2})=\begin{cases}y_{1} & \mbox{if }x_{1}=x_{2}\\
y_{2} & \mbox{otherwise}
\end{cases}
\end{equation*}
\item For any two function-sorted terms $f$ and $g$, their composition
  $f\circ g$ is another function-sorted term;
\item For any two function-sorted terms $f$ and $g$, their ``product''
  or ``juxtaposition'' $(f,g)$ which behaves intuitively like
  $(f,g)(x,y)=(fx,gy)$ is another function-sorted term;
\item For any two function-sorted terms $f$ and $g$, we can use them
  to construct a new function by primitive recursion, which is the
  function-sorted term denoted $\mathcal{R}(f,g)$.
\end{enumerate}
\end{node}

\nwenddocs{}\nwbegincode{17}\sublabel{NW2DqMzY-3aj0iq-1}\nwmargintag{{\nwtagstyle{}\subpageref{NW2DqMzY-3aj0iq-1}}}\moddef{Constructors for $FS_{0}$ functions~{\nwtagstyle{}\subpageref{NW2DqMzY-3aj0iq-1}}}\endmoddef\nwstartdeflinemarkup\nwusesondefline{\\{NW2DqMzY-mVmU8-1}}\nwenddeflinemarkup
\LA{}Identity function $FS_{0}$ primitive function constant~{\nwtagstyle{}\subpageref{NW2DqMzY-D9caK-1}}\RA{}

\LA{}Decide if two individuals are equal~{\nwtagstyle{}\subpageref{NW2DqMzY-2vOgWm-1}}\RA{}

\LA{}Project out the first component of an $FS_{0}$ ordered pair~{\nwtagstyle{}\subpageref{NW2DqMzY-ilITB-1}}\RA{}

\LA{}Project out the second component of an $FS_{0}$ ordered pair~{\nwtagstyle{}\subpageref{NW2DqMzY-HQiIh-1}}\RA{}

\LA{}Constant function constructor~{\nwtagstyle{}\subpageref{NW2DqMzY-22MFug-1}}\RA{}

\LA{}Compose two $FS_{0}$ functions together~{\nwtagstyle{}\subpageref{NW2DqMzY-1dWhhL-1}}\RA{}

\LA{}Apply two functions componentwise to an $FS_{0}$ pair~{\nwtagstyle{}\subpageref{NW2DqMzY-1lMRH2-1}}\RA{}

\LA{}Recursor constructor for primitive recursion~{\nwtagstyle{}\subpageref{NW2DqMzY-2fxlqE-1}}\RA{}
\nwused{\\{NW2DqMzY-mVmU8-1}}\nwendcode{}\nwbegindocs{18}\nwdocspar

\begin{node}
The identity function is always a critical primitive notion. We
provide the lazy and eager applications of the identity function to a
term. Exceptions are raised if the argument is not an {\Tt{}IND\nwendquote} sorted term.
\end{node}

\nwenddocs{}\nwbegincode{19}\sublabel{NW2DqMzY-1uF0T3-1}\nwmargintag{{\nwtagstyle{}\subpageref{NW2DqMzY-1uF0T3-1}}}\moddef{Specification for $FS_{0}$ functions~{\nwtagstyle{}\subpageref{NW2DqMzY-1uF0T3-1}}}\endmoddef\nwstartdeflinemarkup\nwusesondefline{\\{NW2DqMzY-3R3yGC-1}}\nwprevnextdefs{\relax}{NW2DqMzY-1uF0T3-2}\nwenddeflinemarkup
val id_const : \nwlinkedidentc{Term.t}{NW1rTmOJ-4LsRWC-1};
val id : \nwlinkedidentc{Term.t}{NW1rTmOJ-4LsRWC-1} -> \nwlinkedidentc{Term.t}{NW1rTmOJ-4LsRWC-1};
val mk_id : \nwlinkedidentc{Term.t}{NW1rTmOJ-4LsRWC-1} -> \nwlinkedidentc{Term.t}{NW1rTmOJ-4LsRWC-1};
\nwalsodefined{\\{NW2DqMzY-1uF0T3-2}\\{NW2DqMzY-1uF0T3-3}\\{NW2DqMzY-1uF0T3-4}\\{NW2DqMzY-1uF0T3-5}\\{NW2DqMzY-1uF0T3-6}\\{NW2DqMzY-1uF0T3-7}\\{NW2DqMzY-1uF0T3-8}}\nwused{\\{NW2DqMzY-3R3yGC-1}}\nwidentuses{\\{{\nwixident{Term.t}}{Term.t}}}\nwindexuse{\nwixident{Term.t}}{Term.t}{NW2DqMzY-1uF0T3-1}\nwendcode{}\nwbegindocs{20}\nwdocspar

\nwenddocs{}\nwbegincode{21}\sublabel{NW2DqMzY-D9caK-1}\nwmargintag{{\nwtagstyle{}\subpageref{NW2DqMzY-D9caK-1}}}\moddef{Identity function $FS_{0}$ primitive function constant~{\nwtagstyle{}\subpageref{NW2DqMzY-D9caK-1}}}\endmoddef\nwstartdeflinemarkup\nwusesondefline{\\{NW2DqMzY-3aj0iq-1}}\nwenddeflinemarkup
(* id : \nwlinkedidentc{Term.t}{NW1rTmOJ-4LsRWC-1} -> \nwlinkedidentc{Term.t}{NW1rTmOJ-4LsRWC-1}

Just returns the argument (i.e., 'eagerly applies' the
identity function and returns the result). *)
fun id tm = tm;

val id_const = (Term.Fun("id", \nwlinkedidentc{Sort.FUN}{NW1rTmOJ-3ZqggC-1}, []));

(* mk_id tm : \nwlinkedidentc{Term.t}{NW1rTmOJ-4LsRWC-1} -> \nwlinkedidentc{Term.t}{NW1rTmOJ-4LsRWC-1}

Applies the `id_const` to the argument, and returns the
result. *)
fun mk_id tm = app id_const tm;
\nwindexdefn{\nwixident{FS0.id}}{FS0.id}{NW2DqMzY-D9caK-1}\nwindexdefn{\nwixident{FS0.mk{\_}id}}{FS0.mk:unid}{NW2DqMzY-D9caK-1}\eatline
\nwused{\\{NW2DqMzY-3aj0iq-1}}\nwidentdefs{\\{{\nwixident{FS0.id}}{FS0.id}}\\{{\nwixident{FS0.mk{\_}id}}{FS0.mk:unid}}}\nwidentuses{\\{{\nwixident{Sort.FUN}}{Sort.FUN}}\\{{\nwixident{Term.t}}{Term.t}}}\nwindexuse{\nwixident{Sort.FUN}}{Sort.FUN}{NW2DqMzY-D9caK-1}\nwindexuse{\nwixident{Term.t}}{Term.t}{NW2DqMzY-D9caK-1}\nwendcode{}\nwbegindocs{22}\nwdocspar
\begin{node}
There is also the ``conditional'' primitive function $D$ which acts on
4-tuples of individuals. If the first two components are equal to each
other, then it returns the third component. Otherwise it returns the
fourth component. And if anything else is given, then it returns the
{\Tt{}nil\nwendquote} or {\Tt{}zero\nwendquote} individual constant.

We have provided two functions for this situation:
\begin{enumerate}
\item {\Tt{}mk{\_}if{\_}eq\nwendquote} will lazily apply $D$ to a term, and raise an
  exception if applied to something which is not an individual;
\item {\Tt{}if{\_}eq\nwendquote} will try to destructure the term to test immediately
  if the two arguments are syntactically identical to each other, and
  if so return the third component. But if they are not syntactically
  identical, they could still turn out to be equal, so we just
  delegate to {\Tt{}mk{\_}if{\_}eq\nwendquote}.
\end{enumerate}
\textsc{Note}: we are assuming, contrary to Feferman, that we are
using Lisp conventions with lists writing
$(a,b,c,d)=(a,(b,(c,d)))$. If the reader wants to follow Feferman,
than just change the destructuring to $(t_{123},t_{4})$,
$(t_{12},t_{3})$, and $(t_{1},t_{2})$ instead.

In his PhD thesis, Sean Matthews (Chapter 3, \S2.1) refers to this $D$
as the ``decide'' function.
\end{node}

\nwenddocs{}\nwbegincode{23}\sublabel{NW2DqMzY-1uF0T3-2}\nwmargintag{{\nwtagstyle{}\subpageref{NW2DqMzY-1uF0T3-2}}}\moddef{Specification for $FS_{0}$ functions~{\nwtagstyle{}\subpageref{NW2DqMzY-1uF0T3-1}}}\plusendmoddef\nwstartdeflinemarkup\nwusesondefline{\\{NW2DqMzY-3R3yGC-1}}\nwprevnextdefs{NW2DqMzY-1uF0T3-1}{NW2DqMzY-1uF0T3-3}\nwenddeflinemarkup
val if_eq_const : \nwlinkedidentc{Term.t}{NW1rTmOJ-4LsRWC-1};
val mk_if_eq : \nwlinkedidentc{Term.t}{NW1rTmOJ-4LsRWC-1} -> \nwlinkedidentc{Term.t}{NW1rTmOJ-4LsRWC-1};
val if_eq : \nwlinkedidentc{Term.t}{NW1rTmOJ-4LsRWC-1} -> \nwlinkedidentc{Term.t}{NW1rTmOJ-4LsRWC-1};
\nwused{\\{NW2DqMzY-3R3yGC-1}}\nwidentuses{\\{{\nwixident{Term.t}}{Term.t}}}\nwindexuse{\nwixident{Term.t}}{Term.t}{NW2DqMzY-1uF0T3-2}\nwendcode{}\nwbegindocs{24}\nwdocspar

\nwenddocs{}\nwbegincode{25}\sublabel{NW2DqMzY-2vOgWm-1}\nwmargintag{{\nwtagstyle{}\subpageref{NW2DqMzY-2vOgWm-1}}}\moddef{Decide if two individuals are equal~{\nwtagstyle{}\subpageref{NW2DqMzY-2vOgWm-1}}}\endmoddef\nwstartdeflinemarkup\nwusesondefline{\\{NW2DqMzY-3aj0iq-1}}\nwenddeflinemarkup
val if_eq_const = (Term.Fun("if-equal", \nwlinkedidentc{Sort.FUN}{NW1rTmOJ-3ZqggC-1}, [])); 

(* mk_if_eq : \nwlinkedidentc{Term.t}{NW1rTmOJ-4LsRWC-1} -> \nwlinkedidentc{Term.t}{NW1rTmOJ-4LsRWC-1} *)
fun mk_if_eq tm = app if_eq_const tm;

(* if_eq : \nwlinkedidentc{Term.t}{NW1rTmOJ-4LsRWC-1} -> \nwlinkedidentc{Term.t}{NW1rTmOJ-4LsRWC-1} *)
local
  fun dest_cons (tm : \nwlinkedidentc{Term.t}{NW1rTmOJ-4LsRWC-1}) =
   (case tm of
     (Term.Fun("cons", \nwlinkedidentc{Sort.IND}{NW1rTmOJ-3ZqggC-1}, [t1,t2])) => (t1,t2)
   | _ => raise Match);
in
fun if_eq tm =
  if not(\nwlinkedidentc{Term.is_ind}{NW1rTmOJ-3FfXTO-1} tm)
  then raise Sort.Mismatch "\nwlinkedidentc{FS0.if_eq}{NW2DqMzY-2vOgWm-1}: expects IND"
  else
    let
      val (t1, t234) = dest_cons tm;
      val (t2, t34) = dest_cons t234;
      val (t3, t4) = dest_cons t34;
    in
      if \nwlinkedidentc{Term.eq}{NW1rTmOJ-3810sY-1} t1 t2 then t3
      else mk_if_eq tm
    end
    handle Match => mk_if_eq tm;
end;
\nwindexdefn{\nwixident{FS0.mk{\_}if{\_}eq}}{FS0.mk:unif:uneq}{NW2DqMzY-2vOgWm-1}\nwindexdefn{\nwixident{FS0.if{\_}eq}}{FS0.if:uneq}{NW2DqMzY-2vOgWm-1}\eatline
\nwused{\\{NW2DqMzY-3aj0iq-1}}\nwidentdefs{\\{{\nwixident{FS0.if{\_}eq}}{FS0.if:uneq}}\\{{\nwixident{FS0.mk{\_}if{\_}eq}}{FS0.mk:unif:uneq}}}\nwidentuses{\\{{\nwixident{Sort.FUN}}{Sort.FUN}}\\{{\nwixident{Sort.IND}}{Sort.IND}}\\{{\nwixident{Term.eq}}{Term.eq}}\\{{\nwixident{Term.is{\_}ind}}{Term.is:unind}}\\{{\nwixident{Term.t}}{Term.t}}}\nwindexuse{\nwixident{Sort.FUN}}{Sort.FUN}{NW2DqMzY-2vOgWm-1}\nwindexuse{\nwixident{Sort.IND}}{Sort.IND}{NW2DqMzY-2vOgWm-1}\nwindexuse{\nwixident{Term.eq}}{Term.eq}{NW2DqMzY-2vOgWm-1}\nwindexuse{\nwixident{Term.is{\_}ind}}{Term.is:unind}{NW2DqMzY-2vOgWm-1}\nwindexuse{\nwixident{Term.t}}{Term.t}{NW2DqMzY-2vOgWm-1}\nwendcode{}\nwbegindocs{26}\nwdocspar

\begin{node}
We have the projection operations be primitive functions in
$FS_{0}$. But we optimize applying {\Tt{}proj1\nwendquote} to a term by returning
the result according to the axioms of $FS_{0}$, namely {\Tt{}P1\ (x,\ y)\ =\ x\nwendquote}.

However, there may be times when the user wants a \emph{lazy}
application of the function, rather than \emph{eagerly} returning the
result. For this situation, we provide the {\Tt{}mk{\_}proj1\nwendquote} to return the
{\Tt{}app\ P1\ x\nwendquote} individual term.
\end{node}

\nwenddocs{}\nwbegincode{27}\sublabel{NW2DqMzY-1uF0T3-3}\nwmargintag{{\nwtagstyle{}\subpageref{NW2DqMzY-1uF0T3-3}}}\moddef{Specification for $FS_{0}$ functions~{\nwtagstyle{}\subpageref{NW2DqMzY-1uF0T3-1}}}\plusendmoddef\nwstartdeflinemarkup\nwusesondefline{\\{NW2DqMzY-3R3yGC-1}}\nwprevnextdefs{NW2DqMzY-1uF0T3-2}{NW2DqMzY-1uF0T3-4}\nwenddeflinemarkup
val proj1_const : \nwlinkedidentc{Term.t}{NW1rTmOJ-4LsRWC-1};
val mk_proj1 : \nwlinkedidentc{Term.t}{NW1rTmOJ-4LsRWC-1} -> \nwlinkedidentc{Term.t}{NW1rTmOJ-4LsRWC-1};
val proj1 : \nwlinkedidentc{Term.t}{NW1rTmOJ-4LsRWC-1} -> \nwlinkedidentc{Term.t}{NW1rTmOJ-4LsRWC-1};
\nwused{\\{NW2DqMzY-3R3yGC-1}}\nwidentuses{\\{{\nwixident{Term.t}}{Term.t}}}\nwindexuse{\nwixident{Term.t}}{Term.t}{NW2DqMzY-1uF0T3-3}\nwendcode{}\nwbegindocs{28}\nwdocspar
\nwenddocs{}\nwbegincode{29}\sublabel{NW2DqMzY-ilITB-1}\nwmargintag{{\nwtagstyle{}\subpageref{NW2DqMzY-ilITB-1}}}\moddef{Project out the first component of an $FS_{0}$ ordered pair~{\nwtagstyle{}\subpageref{NW2DqMzY-ilITB-1}}}\endmoddef\nwstartdeflinemarkup\nwusesondefline{\\{NW2DqMzY-3aj0iq-1}}\nwenddeflinemarkup
(* proj1_const : \nwlinkedidentc{Term.t}{NW1rTmOJ-4LsRWC-1} *)
val proj1_const = Term.Fun("P1", \nwlinkedidentc{Sort.FUN}{NW1rTmOJ-3ZqggC-1}, []);

(* mk_proj1 : \nwlinkedidentc{Term.t}{NW1rTmOJ-4LsRWC-1} -> \nwlinkedidentc{Term.t}{NW1rTmOJ-4LsRWC-1}

Lazily apply proj1 to a term. *)
fun mk_proj1 tm = app proj1_const tm;

(* proj1 : \nwlinkedidentc{Term.t}{NW1rTmOJ-4LsRWC-1} -> \nwlinkedidentc{Term.t}{NW1rTmOJ-4LsRWC-1}

Automatically extracts the first component of a pair,
otherwise just applies the "P1" constant. *)
fun proj1 tm =
 (case tm of
   (Term.Fun("cons", \nwlinkedidentc{Sort.IND}{NW1rTmOJ-3ZqggC-1}, [tm', _])) => tm'
   | _ => mk_proj1 tm);
\nwindexdefn{\nwixident{FS0.mk{\_}proj1}}{FS0.mk:unproj1}{NW2DqMzY-ilITB-1}\nwindexdefn{\nwixident{FS0.proj1}}{FS0.proj1}{NW2DqMzY-ilITB-1}\eatline
\nwused{\\{NW2DqMzY-3aj0iq-1}}\nwidentdefs{\\{{\nwixident{FS0.mk{\_}proj1}}{FS0.mk:unproj1}}\\{{\nwixident{FS0.proj1}}{FS0.proj1}}}\nwidentuses{\\{{\nwixident{Sort.FUN}}{Sort.FUN}}\\{{\nwixident{Sort.IND}}{Sort.IND}}\\{{\nwixident{Term.t}}{Term.t}}}\nwindexuse{\nwixident{Sort.FUN}}{Sort.FUN}{NW2DqMzY-ilITB-1}\nwindexuse{\nwixident{Sort.IND}}{Sort.IND}{NW2DqMzY-ilITB-1}\nwindexuse{\nwixident{Term.t}}{Term.t}{NW2DqMzY-ilITB-1}\nwendcode{}\nwbegindocs{30}\nwdocspar
\begin{node}
Similarly, we have the same concerns for projecting out the second
component of an ordered pair.
\end{node}

\nwenddocs{}\nwbegincode{31}\sublabel{NW2DqMzY-1uF0T3-4}\nwmargintag{{\nwtagstyle{}\subpageref{NW2DqMzY-1uF0T3-4}}}\moddef{Specification for $FS_{0}$ functions~{\nwtagstyle{}\subpageref{NW2DqMzY-1uF0T3-1}}}\plusendmoddef\nwstartdeflinemarkup\nwusesondefline{\\{NW2DqMzY-3R3yGC-1}}\nwprevnextdefs{NW2DqMzY-1uF0T3-3}{NW2DqMzY-1uF0T3-5}\nwenddeflinemarkup
val proj2_const : \nwlinkedidentc{Term.t}{NW1rTmOJ-4LsRWC-1};
val mk_proj2 : \nwlinkedidentc{Term.t}{NW1rTmOJ-4LsRWC-1} -> \nwlinkedidentc{Term.t}{NW1rTmOJ-4LsRWC-1};
val proj2 : \nwlinkedidentc{Term.t}{NW1rTmOJ-4LsRWC-1} -> \nwlinkedidentc{Term.t}{NW1rTmOJ-4LsRWC-1};
\nwused{\\{NW2DqMzY-3R3yGC-1}}\nwidentuses{\\{{\nwixident{Term.t}}{Term.t}}}\nwindexuse{\nwixident{Term.t}}{Term.t}{NW2DqMzY-1uF0T3-4}\nwendcode{}\nwbegindocs{32}\nwdocspar

\nwenddocs{}\nwbegincode{33}\sublabel{NW2DqMzY-HQiIh-1}\nwmargintag{{\nwtagstyle{}\subpageref{NW2DqMzY-HQiIh-1}}}\moddef{Project out the second component of an $FS_{0}$ ordered pair~{\nwtagstyle{}\subpageref{NW2DqMzY-HQiIh-1}}}\endmoddef\nwstartdeflinemarkup\nwusesondefline{\\{NW2DqMzY-3aj0iq-1}}\nwenddeflinemarkup
(* proj2_const : \nwlinkedidentc{Term.t}{NW1rTmOJ-4LsRWC-1} *)
val proj2_const = Term.Fun("P2", \nwlinkedidentc{Sort.FUN}{NW1rTmOJ-3ZqggC-1}, []);

(* mk_proj2 : \nwlinkedidentc{Term.t}{NW1rTmOJ-4LsRWC-1} -> \nwlinkedidentc{Term.t}{NW1rTmOJ-4LsRWC-1}

Lazily apply proj2 to a term. *)
fun mk_proj2 tm = app proj2_const tm;

(* proj2 : \nwlinkedidentc{Term.t}{NW1rTmOJ-4LsRWC-1} -> \nwlinkedidentc{Term.t}{NW1rTmOJ-4LsRWC-1}

Automatically extracts the second component of a pair,
otherwise just applies the "P2" constant. *)
fun proj2 tm =
 (case tm of
   (Term.Fun("cons", \nwlinkedidentc{Sort.IND}{NW1rTmOJ-3ZqggC-1}, [_, tm'])) => tm'
   | _ => mk_proj2 tm);
\nwindexdefn{\nwixident{FS0.proj2}}{FS0.proj2}{NW2DqMzY-HQiIh-1}\nwindexdefn{\nwixident{FS0.mk{\_}proj2}}{FS0.mk:unproj2}{NW2DqMzY-HQiIh-1}\eatline
\nwused{\\{NW2DqMzY-3aj0iq-1}}\nwidentdefs{\\{{\nwixident{FS0.mk{\_}proj2}}{FS0.mk:unproj2}}\\{{\nwixident{FS0.proj2}}{FS0.proj2}}}\nwidentuses{\\{{\nwixident{Sort.FUN}}{Sort.FUN}}\\{{\nwixident{Sort.IND}}{Sort.IND}}\\{{\nwixident{Term.t}}{Term.t}}}\nwindexuse{\nwixident{Sort.FUN}}{Sort.FUN}{NW2DqMzY-HQiIh-1}\nwindexuse{\nwixident{Sort.IND}}{Sort.IND}{NW2DqMzY-HQiIh-1}\nwindexuse{\nwixident{Term.t}}{Term.t}{NW2DqMzY-HQiIh-1}\nwendcode{}\nwbegindocs{34}\nwdocspar
\begin{node}
The constant function $K(a)$ will always return $a$ regardless of what
argument it is given: $\forall x\ldotp (K(a))x=a$. We provide a
``constructor'' for the constant function, in contrast to what we did
for the previous function primitives. This may be a source of
performance penalty, depending on whether it is used heavily.
\end{node}

\nwenddocs{}\nwbegincode{35}\sublabel{NW2DqMzY-1uF0T3-5}\nwmargintag{{\nwtagstyle{}\subpageref{NW2DqMzY-1uF0T3-5}}}\moddef{Specification for $FS_{0}$ functions~{\nwtagstyle{}\subpageref{NW2DqMzY-1uF0T3-1}}}\plusendmoddef\nwstartdeflinemarkup\nwusesondefline{\\{NW2DqMzY-3R3yGC-1}}\nwprevnextdefs{NW2DqMzY-1uF0T3-4}{NW2DqMzY-1uF0T3-6}\nwenddeflinemarkup
val const : \nwlinkedidentc{Term.t}{NW1rTmOJ-4LsRWC-1} -> \nwlinkedidentc{Term.t}{NW1rTmOJ-4LsRWC-1};
val dest_const : \nwlinkedidentc{Term.t}{NW1rTmOJ-4LsRWC-1} -> \nwlinkedidentc{Term.t}{NW1rTmOJ-4LsRWC-1};
\nwused{\\{NW2DqMzY-3R3yGC-1}}\nwidentuses{\\{{\nwixident{Term.t}}{Term.t}}}\nwindexuse{\nwixident{Term.t}}{Term.t}{NW2DqMzY-1uF0T3-5}\nwendcode{}\nwbegindocs{36}\nwdocspar
\nwenddocs{}\nwbegincode{37}\sublabel{NW2DqMzY-22MFug-1}\nwmargintag{{\nwtagstyle{}\subpageref{NW2DqMzY-22MFug-1}}}\moddef{Constant function constructor~{\nwtagstyle{}\subpageref{NW2DqMzY-22MFug-1}}}\endmoddef\nwstartdeflinemarkup\nwusesondefline{\\{NW2DqMzY-3aj0iq-1}}\nwenddeflinemarkup
(* const : \nwlinkedidentc{Term.t}{NW1rTmOJ-4LsRWC-1} -> \nwlinkedidentc{Term.t}{NW1rTmOJ-4LsRWC-1} *)
fun const tm =
  if not(\nwlinkedidentc{Term.is_ind}{NW1rTmOJ-3FfXTO-1} tm)
  then raise Sort.Mismatch "\nwlinkedidentc{FS0.const}{NW2DqMzY-22MFug-1}: expects an individual"
  else Term.Fun("const", \nwlinkedidentc{Sort.FUN}{NW1rTmOJ-3ZqggC-1}, [tm]);

\LA{}Destructure \code{}const\edoc{}~{\nwtagstyle{}\subpageref{NW2DqMzY-1qEma8-1}}\RA{}
\nwindexdefn{\nwixident{FS0.const}}{FS0.const}{NW2DqMzY-22MFug-1}\eatline
\nwused{\\{NW2DqMzY-3aj0iq-1}}\nwidentdefs{\\{{\nwixident{FS0.const}}{FS0.const}}}\nwidentuses{\\{{\nwixident{Sort.FUN}}{Sort.FUN}}\\{{\nwixident{Term.is{\_}ind}}{Term.is:unind}}\\{{\nwixident{Term.t}}{Term.t}}}\nwindexuse{\nwixident{Sort.FUN}}{Sort.FUN}{NW2DqMzY-22MFug-1}\nwindexuse{\nwixident{Term.is{\_}ind}}{Term.is:unind}{NW2DqMzY-22MFug-1}\nwindexuse{\nwixident{Term.t}}{Term.t}{NW2DqMzY-22MFug-1}\nwendcode{}\nwbegindocs{38}\nwdocspar
\nwenddocs{}\nwbegincode{39}\sublabel{NW2DqMzY-1qEma8-1}\nwmargintag{{\nwtagstyle{}\subpageref{NW2DqMzY-1qEma8-1}}}\moddef{Destructure \code{}const\edoc{}~{\nwtagstyle{}\subpageref{NW2DqMzY-1qEma8-1}}}\endmoddef\nwstartdeflinemarkup\nwusesondefline{\\{NW2DqMzY-22MFug-1}}\nwenddeflinemarkup
fun dest_const (Term.Fun("const", \nwlinkedidentc{Sort.FUN}{NW1rTmOJ-3ZqggC-1}, [tm])) = tm
  | dest_const _ = raise Dest "\nwlinkedidentc{FS0.dest_const}{NW2DqMzY-1qEma8-1}";
\nwindexdefn{\nwixident{FS0.dest{\_}const}}{FS0.dest:unconst}{NW2DqMzY-1qEma8-1}\eatline
\nwused{\\{NW2DqMzY-22MFug-1}}\nwidentdefs{\\{{\nwixident{FS0.dest{\_}const}}{FS0.dest:unconst}}}\nwidentuses{\\{{\nwixident{Sort.FUN}}{Sort.FUN}}}\nwindexuse{\nwixident{Sort.FUN}}{Sort.FUN}{NW2DqMzY-1qEma8-1}\nwendcode{}\nwbegindocs{40}\nwdocspar
\begin{node}
We provide the constructor for composing two functions together. We
follow the usual Mathematical convention that $(f\circ g)(x)=f(g(x))$.
Again, this is a constructor, not the ``unfolded'' application.

Although, after thinking about it for some time, it might also be wise
to take advantage of associativity of composition to allow \emph{an
arbitrary number} of functions to be composed together. So
$(f\circ(g\circ h))$ would be stored in a flat list $[f,g,h]$. The
when applied, we just need to invoke {\Tt{}List.foldr\ app\ arg\ funs\nwendquote}.
\end{node}
\nwenddocs{}\nwbegincode{41}\sublabel{NW2DqMzY-1uF0T3-6}\nwmargintag{{\nwtagstyle{}\subpageref{NW2DqMzY-1uF0T3-6}}}\moddef{Specification for $FS_{0}$ functions~{\nwtagstyle{}\subpageref{NW2DqMzY-1uF0T3-1}}}\plusendmoddef\nwstartdeflinemarkup\nwusesondefline{\\{NW2DqMzY-3R3yGC-1}}\nwprevnextdefs{NW2DqMzY-1uF0T3-5}{NW2DqMzY-1uF0T3-7}\nwenddeflinemarkup
val comp : \nwlinkedidentc{Term.t}{NW1rTmOJ-4LsRWC-1} -> \nwlinkedidentc{Term.t}{NW1rTmOJ-4LsRWC-1} -> \nwlinkedidentc{Term.t}{NW1rTmOJ-4LsRWC-1};
val dest_comp : \nwlinkedidentc{Term.t}{NW1rTmOJ-4LsRWC-1} -> \nwlinkedidentc{Term.t}{NW1rTmOJ-4LsRWC-1} * \nwlinkedidentc{Term.t}{NW1rTmOJ-4LsRWC-1};
\nwused{\\{NW2DqMzY-3R3yGC-1}}\nwidentuses{\\{{\nwixident{Term.t}}{Term.t}}}\nwindexuse{\nwixident{Term.t}}{Term.t}{NW2DqMzY-1uF0T3-6}\nwendcode{}\nwbegindocs{42}\nwdocspar

\nwenddocs{}\nwbegincode{43}\sublabel{NW2DqMzY-1dWhhL-1}\nwmargintag{{\nwtagstyle{}\subpageref{NW2DqMzY-1dWhhL-1}}}\moddef{Compose two $FS_{0}$ functions together~{\nwtagstyle{}\subpageref{NW2DqMzY-1dWhhL-1}}}\endmoddef\nwstartdeflinemarkup\nwusesondefline{\\{NW2DqMzY-3aj0iq-1}}\nwenddeflinemarkup
(* comp : \nwlinkedidentc{Term.t}{NW1rTmOJ-4LsRWC-1} -> \nwlinkedidentc{Term.t}{NW1rTmOJ-4LsRWC-1} -> \nwlinkedidentc{Term.t}{NW1rTmOJ-4LsRWC-1}

Compose two functions *)
fun comp f g =
  if not(\nwlinkedidentc{Term.is_fun}{NW1rTmOJ-3FfXTO-1} f)
  then raise Sort.Mismatch "\nwlinkedidentc{FS0.comp}{NW2DqMzY-1dWhhL-1}: first argument is not a fun"
  else if not(\nwlinkedidentc{Term.is_fun}{NW1rTmOJ-3FfXTO-1} g)
  then raise Sort.Mismatch "\nwlinkedidentc{FS0.comp}{NW2DqMzY-1dWhhL-1}: second argument is not a fun"
  else Term.Fun("comp", \nwlinkedidentc{Sort.FUN}{NW1rTmOJ-3ZqggC-1}, [f, g]);

\LA{}Destructuring \code{}comp\edoc{}~{\nwtagstyle{}\subpageref{NW2DqMzY-22XS9g-1}}\RA{}
\nwindexdefn{\nwixident{FS0.comp}}{FS0.comp}{NW2DqMzY-1dWhhL-1}\eatline
\nwused{\\{NW2DqMzY-3aj0iq-1}}\nwidentdefs{\\{{\nwixident{FS0.comp}}{FS0.comp}}}\nwidentuses{\\{{\nwixident{Sort.FUN}}{Sort.FUN}}\\{{\nwixident{Term.is{\_}fun}}{Term.is:unfun}}\\{{\nwixident{Term.t}}{Term.t}}}\nwindexuse{\nwixident{Sort.FUN}}{Sort.FUN}{NW2DqMzY-1dWhhL-1}\nwindexuse{\nwixident{Term.is{\_}fun}}{Term.is:unfun}{NW2DqMzY-1dWhhL-1}\nwindexuse{\nwixident{Term.t}}{Term.t}{NW2DqMzY-1dWhhL-1}\nwendcode{}\nwbegindocs{44}\nwdocspar
\nwenddocs{}\nwbegincode{45}\sublabel{NW2DqMzY-22XS9g-1}\nwmargintag{{\nwtagstyle{}\subpageref{NW2DqMzY-22XS9g-1}}}\moddef{Destructuring \code{}comp\edoc{}~{\nwtagstyle{}\subpageref{NW2DqMzY-22XS9g-1}}}\endmoddef\nwstartdeflinemarkup\nwusesondefline{\\{NW2DqMzY-1dWhhL-1}}\nwenddeflinemarkup
fun dest_comp (Term.Fun("comp", \nwlinkedidentc{Sort.FUN}{NW1rTmOJ-3ZqggC-1}, [f, g])) = (f,g)
  | dest_comp _ = raise Dest "\nwlinkedidentc{FS0.dest_comp}{NW2DqMzY-22XS9g-1}";
\nwindexdefn{\nwixident{FS0.dest{\_}comp}}{FS0.dest:uncomp}{NW2DqMzY-22XS9g-1}\eatline
\nwused{\\{NW2DqMzY-1dWhhL-1}}\nwidentdefs{\\{{\nwixident{FS0.dest{\_}comp}}{FS0.dest:uncomp}}}\nwidentuses{\\{{\nwixident{Sort.FUN}}{Sort.FUN}}}\nwindexuse{\nwixident{Sort.FUN}}{Sort.FUN}{NW2DqMzY-22XS9g-1}\nwendcode{}\nwbegindocs{46}\nwdocspar
\begin{node}
We have the constructor for a {\Tt{}juxt\nwendquote}. Where did the name ``juxt''
come from? Clojure. Again, this is defined to be such that $(f,g)x=(fx,gx)$.
The {\Tt{}\nwlinkedidentq{FS0.app}{NW2DqMzY-2oGoBx-1}\nwendquote} function should expand this out.
\end{node}

\nwenddocs{}\nwbegincode{47}\sublabel{NW2DqMzY-1uF0T3-7}\nwmargintag{{\nwtagstyle{}\subpageref{NW2DqMzY-1uF0T3-7}}}\moddef{Specification for $FS_{0}$ functions~{\nwtagstyle{}\subpageref{NW2DqMzY-1uF0T3-1}}}\plusendmoddef\nwstartdeflinemarkup\nwusesondefline{\\{NW2DqMzY-3R3yGC-1}}\nwprevnextdefs{NW2DqMzY-1uF0T3-6}{NW2DqMzY-1uF0T3-8}\nwenddeflinemarkup
val juxt : \nwlinkedidentc{Term.t}{NW1rTmOJ-4LsRWC-1} -> \nwlinkedidentc{Term.t}{NW1rTmOJ-4LsRWC-1} -> \nwlinkedidentc{Term.t}{NW1rTmOJ-4LsRWC-1};
val dest_juxt : \nwlinkedidentc{Term.t}{NW1rTmOJ-4LsRWC-1} -> \nwlinkedidentc{Term.t}{NW1rTmOJ-4LsRWC-1} * \nwlinkedidentc{Term.t}{NW1rTmOJ-4LsRWC-1};
\nwused{\\{NW2DqMzY-3R3yGC-1}}\nwidentuses{\\{{\nwixident{Term.t}}{Term.t}}}\nwindexuse{\nwixident{Term.t}}{Term.t}{NW2DqMzY-1uF0T3-7}\nwendcode{}\nwbegindocs{48}\nwdocspar

\nwenddocs{}\nwbegincode{49}\sublabel{NW2DqMzY-1lMRH2-1}\nwmargintag{{\nwtagstyle{}\subpageref{NW2DqMzY-1lMRH2-1}}}\moddef{Apply two functions componentwise to an $FS_{0}$ pair~{\nwtagstyle{}\subpageref{NW2DqMzY-1lMRH2-1}}}\endmoddef\nwstartdeflinemarkup\nwusesondefline{\\{NW2DqMzY-3aj0iq-1}}\nwenddeflinemarkup
(* juxt : \nwlinkedidentc{Term.t}{NW1rTmOJ-4LsRWC-1} -> \nwlinkedidentc{Term.t}{NW1rTmOJ-4LsRWC-1} -> \nwlinkedidentc{Term.t}{NW1rTmOJ-4LsRWC-1}

Given two functions, applies them componentwise to a `cons`
pair. *)
fun juxt f g =
  if not(\nwlinkedidentc{Term.is_fun}{NW1rTmOJ-3FfXTO-1} f)
  then raise Sort.Mismatch "\nwlinkedidentc{FS0.juxt}{NW2DqMzY-1lMRH2-1}: first argument is not a fun"
  else if not(\nwlinkedidentc{Term.is_fun}{NW1rTmOJ-3FfXTO-1} g)
  then raise Sort.Mismatch "\nwlinkedidentc{FS0.juxt}{NW2DqMzY-1lMRH2-1}: second argument is not a fun"
  else Term.Fun("juxt", \nwlinkedidentc{Sort.FUN}{NW1rTmOJ-3ZqggC-1}, [f, g]);

\LA{}Destructure \code{}juxt\edoc{}~{\nwtagstyle{}\subpageref{NW2DqMzY-hvApp-1}}\RA{}
\nwindexdefn{\nwixident{FS0.juxt}}{FS0.juxt}{NW2DqMzY-1lMRH2-1}\eatline
\nwused{\\{NW2DqMzY-3aj0iq-1}}\nwidentdefs{\\{{\nwixident{FS0.juxt}}{FS0.juxt}}}\nwidentuses{\\{{\nwixident{Sort.FUN}}{Sort.FUN}}\\{{\nwixident{Term.is{\_}fun}}{Term.is:unfun}}\\{{\nwixident{Term.t}}{Term.t}}}\nwindexuse{\nwixident{Sort.FUN}}{Sort.FUN}{NW2DqMzY-1lMRH2-1}\nwindexuse{\nwixident{Term.is{\_}fun}}{Term.is:unfun}{NW2DqMzY-1lMRH2-1}\nwindexuse{\nwixident{Term.t}}{Term.t}{NW2DqMzY-1lMRH2-1}\nwendcode{}\nwbegindocs{50}\nwdocspar
\nwenddocs{}\nwbegincode{51}\sublabel{NW2DqMzY-hvApp-1}\nwmargintag{{\nwtagstyle{}\subpageref{NW2DqMzY-hvApp-1}}}\moddef{Destructure \code{}juxt\edoc{}~{\nwtagstyle{}\subpageref{NW2DqMzY-hvApp-1}}}\endmoddef\nwstartdeflinemarkup\nwusesondefline{\\{NW2DqMzY-1lMRH2-1}}\nwenddeflinemarkup
fun dest_juxt (Term.Fun("juxt", \nwlinkedidentc{Sort.FUN}{NW1rTmOJ-3ZqggC-1}, [f, g])) = (f,g)
  | dest_juxt _ = raise Dest "dest_juxt";
\nwindexdefn{\nwixident{FS0.dest{\_}juxt}}{FS0.dest:unjuxt}{NW2DqMzY-hvApp-1}\eatline
\nwused{\\{NW2DqMzY-1lMRH2-1}}\nwidentdefs{\\{{\nwixident{FS0.dest{\_}juxt}}{FS0.dest:unjuxt}}}\nwidentuses{\\{{\nwixident{Sort.FUN}}{Sort.FUN}}}\nwindexuse{\nwixident{Sort.FUN}}{Sort.FUN}{NW2DqMzY-hvApp-1}\nwendcode{}\nwbegindocs{52}\nwdocspar
\begin{node}
The recursor constructor $\mathcal{R}(f,g)$ is needed for primitive
recursion in $FS_{0}$. I don't pretend to optimize it in any way,
shape, or form.
\end{node}
\nwenddocs{}\nwbegincode{53}\sublabel{NW2DqMzY-1uF0T3-8}\nwmargintag{{\nwtagstyle{}\subpageref{NW2DqMzY-1uF0T3-8}}}\moddef{Specification for $FS_{0}$ functions~{\nwtagstyle{}\subpageref{NW2DqMzY-1uF0T3-1}}}\plusendmoddef\nwstartdeflinemarkup\nwusesondefline{\\{NW2DqMzY-3R3yGC-1}}\nwprevnextdefs{NW2DqMzY-1uF0T3-7}{\relax}\nwenddeflinemarkup
val recur : \nwlinkedidentc{Term.t}{NW1rTmOJ-4LsRWC-1} -> \nwlinkedidentc{Term.t}{NW1rTmOJ-4LsRWC-1} -> \nwlinkedidentc{Term.t}{NW1rTmOJ-4LsRWC-1};
val dest_recur : \nwlinkedidentc{Term.t}{NW1rTmOJ-4LsRWC-1} -> \nwlinkedidentc{Term.t}{NW1rTmOJ-4LsRWC-1} * \nwlinkedidentc{Term.t}{NW1rTmOJ-4LsRWC-1};
\nwused{\\{NW2DqMzY-3R3yGC-1}}\nwidentuses{\\{{\nwixident{Term.t}}{Term.t}}}\nwindexuse{\nwixident{Term.t}}{Term.t}{NW2DqMzY-1uF0T3-8}\nwendcode{}\nwbegindocs{54}\nwdocspar

\nwenddocs{}\nwbegincode{55}\sublabel{NW2DqMzY-2fxlqE-1}\nwmargintag{{\nwtagstyle{}\subpageref{NW2DqMzY-2fxlqE-1}}}\moddef{Recursor constructor for primitive recursion~{\nwtagstyle{}\subpageref{NW2DqMzY-2fxlqE-1}}}\endmoddef\nwstartdeflinemarkup\nwusesondefline{\\{NW2DqMzY-3aj0iq-1}}\nwenddeflinemarkup
(* recur : \nwlinkedidentc{Term.t}{NW1rTmOJ-4LsRWC-1} -> \nwlinkedidentc{Term.t}{NW1rTmOJ-4LsRWC-1} -> \nwlinkedidentc{Term.t}{NW1rTmOJ-4LsRWC-1}

Constructs the recursor given two functions. *)
fun recur f g =
  if not(\nwlinkedidentc{Term.is_fun}{NW1rTmOJ-3FfXTO-1} f)
  then raise Sort.Mismatch "\nwlinkedidentc{FS0.recur}{NW2DqMzY-2fxlqE-1}: first argument is not a fun"
  else if not(\nwlinkedidentc{Term.is_fun}{NW1rTmOJ-3FfXTO-1} g)
  then raise Sort.Mismatch "\nwlinkedidentc{FS0.recur}{NW2DqMzY-2fxlqE-1}: second argument is not a fun"
  else Term.Fun("recur", \nwlinkedidentc{Sort.FUN}{NW1rTmOJ-3ZqggC-1}, [f, g]);

\LA{}Destructure \code{}recur\edoc{}~{\nwtagstyle{}\subpageref{NW2DqMzY-Ir9cb-1}}\RA{}
\nwindexdefn{\nwixident{FS0.recur}}{FS0.recur}{NW2DqMzY-2fxlqE-1}\eatline
\nwused{\\{NW2DqMzY-3aj0iq-1}}\nwidentdefs{\\{{\nwixident{FS0.recur}}{FS0.recur}}}\nwidentuses{\\{{\nwixident{Sort.FUN}}{Sort.FUN}}\\{{\nwixident{Term.is{\_}fun}}{Term.is:unfun}}\\{{\nwixident{Term.t}}{Term.t}}}\nwindexuse{\nwixident{Sort.FUN}}{Sort.FUN}{NW2DqMzY-2fxlqE-1}\nwindexuse{\nwixident{Term.is{\_}fun}}{Term.is:unfun}{NW2DqMzY-2fxlqE-1}\nwindexuse{\nwixident{Term.t}}{Term.t}{NW2DqMzY-2fxlqE-1}\nwendcode{}\nwbegindocs{56}\nwdocspar
\nwenddocs{}\nwbegincode{57}\sublabel{NW2DqMzY-Ir9cb-1}\nwmargintag{{\nwtagstyle{}\subpageref{NW2DqMzY-Ir9cb-1}}}\moddef{Destructure \code{}recur\edoc{}~{\nwtagstyle{}\subpageref{NW2DqMzY-Ir9cb-1}}}\endmoddef\nwstartdeflinemarkup\nwusesondefline{\\{NW2DqMzY-2fxlqE-1}}\nwenddeflinemarkup
fun dest_recur (Term.Fun("recur", \nwlinkedidentc{Sort.FUN}{NW1rTmOJ-3ZqggC-1}, [f, g])) = (f,g)
  | dest_recur _ = raise Dest "\nwlinkedidentc{FS0.dest_recur}{NW2DqMzY-Ir9cb-1}";
\nwindexdefn{\nwixident{FS0.dest{\_}recur}}{FS0.dest:unrecur}{NW2DqMzY-Ir9cb-1}\eatline
\nwused{\\{NW2DqMzY-2fxlqE-1}}\nwidentdefs{\\{{\nwixident{FS0.dest{\_}recur}}{FS0.dest:unrecur}}}\nwidentuses{\\{{\nwixident{Sort.FUN}}{Sort.FUN}}}\nwindexuse{\nwixident{Sort.FUN}}{Sort.FUN}{NW2DqMzY-Ir9cb-1}\nwendcode{}\nwbegindocs{58}\nwdocspar
\section*{Constructors for classes}

\begin{node}
We have our set $V_{0}$ of individual-sorted terms, and we have
discussed function-sorted terms as primitive recursive functions
$f\colon V_{0}\to V_{0}$. Now we can also discuss subclasses of
$V_{0}$. These give us class-sorted terms. They are:
\begin{enumerate}
\item If $t$ is an individual-sorted term, then the singleton $\{t\}$
  is a class-sorted term;\footnote{Originally this was just $\{0\}$ is
  a class-sorted term, but we will see that $\{t\}$ is equally valid.}
\item If $f$ is a function-sorted term and $C$ is a class-sorted term,
  then the preimage $f^{-1}C$ of $C$ under $f$ is a class-sorted term;
\item If $S$ and $T$ are class-sorted terms, then so are their union
  $S\cup T$ and intersection $S\cap T$;
\item If $S$ and $T$ are class-sorted terms (where $T$ is viewed as a
  ternary relation), then they inductively define a set
  $\mathcal{I}_{2}(S,T)$ which intuitively corresponds to saying
  $S\subset \mathcal{I}_{2}(S,T)$ and for any individual-sorted terms
  $x$, $y$, $z$, if $y\in\mathcal{I}_{2}(S,T)$ and $z\in\mathcal{I}_{2}(S,T)$
  and $(x,y,z)\in T$, then $x\in\mathcal{I}_{2}(S,T)$.
\end{enumerate}
\end{node}

\nwenddocs{}\nwbegincode{59}\sublabel{NW2DqMzY-TSvER-1}\nwmargintag{{\nwtagstyle{}\subpageref{NW2DqMzY-TSvER-1}}}\moddef{Constructors for $FS_{0}$ classes~{\nwtagstyle{}\subpageref{NW2DqMzY-TSvER-1}}}\endmoddef\nwstartdeflinemarkup\nwusesondefline{\\{NW2DqMzY-mVmU8-1}}\nwenddeflinemarkup
\LA{}Singleton $FS_{0}$ class~{\nwtagstyle{}\subpageref{NW2DqMzY-14kolG-1}}\RA{}

\LA{}Preimage of an $FS_{0}$ class under an $FS_{0}$ function~{\nwtagstyle{}\subpageref{NW2DqMzY-31WKtJ-1}}\RA{}

\LA{}Unions of two $FS_{0}$ classes~{\nwtagstyle{}\subpageref{NW2DqMzY-22W7o7-1}}\RA{}

\LA{}Intersection of two $FS_{0}$ classes~{\nwtagstyle{}\subpageref{NW2DqMzY-19bkPi-1}}\RA{}

\LA{}Constructor for inductively defined $FS_{0}$ class~{\nwtagstyle{}\subpageref{NW2DqMzY-1mMvkv-1}}\RA{}
\nwused{\\{NW2DqMzY-mVmU8-1}}\nwendcode{}\nwbegindocs{60}\nwdocspar

\begin{node}
The singleton class $\{x\}$ (for any individual term $x$) could
probably be optimized. Well, we are being a bit more generous than
Feferman, since he's axioms only discuss $\{0\}$.
\end{node}

\nwenddocs{}\nwbegincode{61}\sublabel{NW2DqMzY-33zuOn-1}\nwmargintag{{\nwtagstyle{}\subpageref{NW2DqMzY-33zuOn-1}}}\moddef{Specification for $FS_{0}$ classes~{\nwtagstyle{}\subpageref{NW2DqMzY-33zuOn-1}}}\endmoddef\nwstartdeflinemarkup\nwusesondefline{\\{NW2DqMzY-3R3yGC-1}}\nwprevnextdefs{\relax}{NW2DqMzY-33zuOn-2}\nwenddeflinemarkup
val singleton : \nwlinkedidentc{Term.t}{NW1rTmOJ-4LsRWC-1} -> \nwlinkedidentc{Term.t}{NW1rTmOJ-4LsRWC-1};
val dest_singleton : \nwlinkedidentc{Term.t}{NW1rTmOJ-4LsRWC-1} -> \nwlinkedidentc{Term.t}{NW1rTmOJ-4LsRWC-1};
\nwalsodefined{\\{NW2DqMzY-33zuOn-2}\\{NW2DqMzY-33zuOn-3}\\{NW2DqMzY-33zuOn-4}\\{NW2DqMzY-33zuOn-5}}\nwused{\\{NW2DqMzY-3R3yGC-1}}\nwidentuses{\\{{\nwixident{Term.t}}{Term.t}}}\nwindexuse{\nwixident{Term.t}}{Term.t}{NW2DqMzY-33zuOn-1}\nwendcode{}\nwbegindocs{62}\nwdocspar

\nwenddocs{}\nwbegincode{63}\sublabel{NW2DqMzY-14kolG-1}\nwmargintag{{\nwtagstyle{}\subpageref{NW2DqMzY-14kolG-1}}}\moddef{Singleton $FS_{0}$ class~{\nwtagstyle{}\subpageref{NW2DqMzY-14kolG-1}}}\endmoddef\nwstartdeflinemarkup\nwusesondefline{\\{NW2DqMzY-TSvER-1}}\nwenddeflinemarkup
(* singleton : \nwlinkedidentc{Term.t}{NW1rTmOJ-4LsRWC-1} -> \nwlinkedidentc{Term.t}{NW1rTmOJ-4LsRWC-1} *)
fun singleton x =
  if not(\nwlinkedidentc{Term.is_ind}{NW1rTmOJ-3FfXTO-1} x)
  then raise Sort.Mismatch "\nwlinkedidentc{FS0.singleton}{NW2DqMzY-14kolG-1}: expects an individual"
  else Term.Fun("singleton", \nwlinkedidentc{Sort.CLASS}{NW1rTmOJ-3ZqggC-1}, [x]);

\LA{}Destructure \code{}singleton\edoc{}~{\nwtagstyle{}\subpageref{NW2DqMzY-Qd4qS-1}}\RA{}
\nwindexdefn{\nwixident{FS0.singleton}}{FS0.singleton}{NW2DqMzY-14kolG-1}\eatline
\nwused{\\{NW2DqMzY-TSvER-1}}\nwidentdefs{\\{{\nwixident{FS0.singleton}}{FS0.singleton}}}\nwidentuses{\\{{\nwixident{Sort.CLASS}}{Sort.CLASS}}\\{{\nwixident{Term.is{\_}ind}}{Term.is:unind}}\\{{\nwixident{Term.t}}{Term.t}}}\nwindexuse{\nwixident{Sort.CLASS}}{Sort.CLASS}{NW2DqMzY-14kolG-1}\nwindexuse{\nwixident{Term.is{\_}ind}}{Term.is:unind}{NW2DqMzY-14kolG-1}\nwindexuse{\nwixident{Term.t}}{Term.t}{NW2DqMzY-14kolG-1}\nwendcode{}\nwbegindocs{64}\nwdocspar
\nwenddocs{}\nwbegincode{65}\sublabel{NW2DqMzY-Qd4qS-1}\nwmargintag{{\nwtagstyle{}\subpageref{NW2DqMzY-Qd4qS-1}}}\moddef{Destructure \code{}singleton\edoc{}~{\nwtagstyle{}\subpageref{NW2DqMzY-Qd4qS-1}}}\endmoddef\nwstartdeflinemarkup\nwusesondefline{\\{NW2DqMzY-14kolG-1}}\nwenddeflinemarkup
fun dest_singleton (Term.Fun("singleton", \nwlinkedidentc{Sort.CLASS}{NW1rTmOJ-3ZqggC-1}, [x])) = x
  | dest_singleton _ = raise Dest "\nwlinkedidentc{FS0.dest_singleton}{NW2DqMzY-Qd4qS-1}";
\nwindexdefn{\nwixident{FS0.dest{\_}singleton}}{FS0.dest:unsingleton}{NW2DqMzY-Qd4qS-1}\eatline
\nwused{\\{NW2DqMzY-14kolG-1}}\nwidentdefs{\\{{\nwixident{FS0.dest{\_}singleton}}{FS0.dest:unsingleton}}}\nwidentuses{\\{{\nwixident{Sort.CLASS}}{Sort.CLASS}}}\nwindexuse{\nwixident{Sort.CLASS}}{Sort.CLASS}{NW2DqMzY-Qd4qS-1}\nwendcode{}\nwbegindocs{66}\nwdocspar
\begin{node}
The preimage of a class $C$ under a function $f$ is another class
$f^{-1}C$. We do not optimize the construction of this class. In fact,
the class constructors are just treated as primitives.
\end{node}

\nwenddocs{}\nwbegincode{67}\sublabel{NW2DqMzY-33zuOn-2}\nwmargintag{{\nwtagstyle{}\subpageref{NW2DqMzY-33zuOn-2}}}\moddef{Specification for $FS_{0}$ classes~{\nwtagstyle{}\subpageref{NW2DqMzY-33zuOn-1}}}\plusendmoddef\nwstartdeflinemarkup\nwusesondefline{\\{NW2DqMzY-3R3yGC-1}}\nwprevnextdefs{NW2DqMzY-33zuOn-1}{NW2DqMzY-33zuOn-3}\nwenddeflinemarkup
val preimage : \nwlinkedidentc{Term.t}{NW1rTmOJ-4LsRWC-1} -> \nwlinkedidentc{Term.t}{NW1rTmOJ-4LsRWC-1} -> \nwlinkedidentc{Term.t}{NW1rTmOJ-4LsRWC-1};
val dest_preimage : \nwlinkedidentc{Term.t}{NW1rTmOJ-4LsRWC-1} -> \nwlinkedidentc{Term.t}{NW1rTmOJ-4LsRWC-1} * \nwlinkedidentc{Term.t}{NW1rTmOJ-4LsRWC-1};
\nwused{\\{NW2DqMzY-3R3yGC-1}}\nwidentuses{\\{{\nwixident{Term.t}}{Term.t}}}\nwindexuse{\nwixident{Term.t}}{Term.t}{NW2DqMzY-33zuOn-2}\nwendcode{}\nwbegindocs{68}\nwdocspar

\nwenddocs{}\nwbegincode{69}\sublabel{NW2DqMzY-31WKtJ-1}\nwmargintag{{\nwtagstyle{}\subpageref{NW2DqMzY-31WKtJ-1}}}\moddef{Preimage of an $FS_{0}$ class under an $FS_{0}$ function~{\nwtagstyle{}\subpageref{NW2DqMzY-31WKtJ-1}}}\endmoddef\nwstartdeflinemarkup\nwusesondefline{\\{NW2DqMzY-TSvER-1}}\nwenddeflinemarkup
(* preimage : \nwlinkedidentc{Term.t}{NW1rTmOJ-4LsRWC-1} -> \nwlinkedidentc{Term.t}{NW1rTmOJ-4LsRWC-1} -> \nwlinkedidentc{Term.t}{NW1rTmOJ-4LsRWC-1} *)
fun preimage f S =
  if not(\nwlinkedidentc{Term.is_fun}{NW1rTmOJ-3FfXTO-1} f)
  then raise Sort.Mismatch "\nwlinkedidentc{FS0.preimage}{NW2DqMzY-31WKtJ-1}: first argument is not a fun"
  else if not(\nwlinkedidentc{Term.is_class}{NW1rTmOJ-3FfXTO-1} S)
  then raise Sort.Mismatch "\nwlinkedidentc{FS0.preimage}{NW2DqMzY-31WKtJ-1}: second argument is not a class"
  else Term.Fun("preimage", \nwlinkedidentc{Sort.CLASS}{NW1rTmOJ-3ZqggC-1}, [f, S]);

\LA{}Destructure \code{}preimage\edoc{}~{\nwtagstyle{}\subpageref{NW2DqMzY-LSOi5-1}}\RA{}
\nwindexdefn{\nwixident{FS0.preimage}}{FS0.preimage}{NW2DqMzY-31WKtJ-1}\eatline
\nwused{\\{NW2DqMzY-TSvER-1}}\nwidentdefs{\\{{\nwixident{FS0.preimage}}{FS0.preimage}}}\nwidentuses{\\{{\nwixident{Sort.CLASS}}{Sort.CLASS}}\\{{\nwixident{Term.is{\_}class}}{Term.is:unclass}}\\{{\nwixident{Term.is{\_}fun}}{Term.is:unfun}}\\{{\nwixident{Term.t}}{Term.t}}}\nwindexuse{\nwixident{Sort.CLASS}}{Sort.CLASS}{NW2DqMzY-31WKtJ-1}\nwindexuse{\nwixident{Term.is{\_}class}}{Term.is:unclass}{NW2DqMzY-31WKtJ-1}\nwindexuse{\nwixident{Term.is{\_}fun}}{Term.is:unfun}{NW2DqMzY-31WKtJ-1}\nwindexuse{\nwixident{Term.t}}{Term.t}{NW2DqMzY-31WKtJ-1}\nwendcode{}\nwbegindocs{70}\nwdocspar
\nwenddocs{}\nwbegincode{71}\sublabel{NW2DqMzY-LSOi5-1}\nwmargintag{{\nwtagstyle{}\subpageref{NW2DqMzY-LSOi5-1}}}\moddef{Destructure \code{}preimage\edoc{}~{\nwtagstyle{}\subpageref{NW2DqMzY-LSOi5-1}}}\endmoddef\nwstartdeflinemarkup\nwusesondefline{\\{NW2DqMzY-31WKtJ-1}}\nwenddeflinemarkup
fun dest_preimage (Term.Fun("preimage", \nwlinkedidentc{Sort.CLASS}{NW1rTmOJ-3ZqggC-1}, [f, S])) = (f,S)
  | dest_preimage _ = raise Dest "\nwlinkedidentc{FS0.dest_preimage}{NW2DqMzY-LSOi5-1}";
\nwindexdefn{\nwixident{FS0.dest{\_}preimage}}{FS0.dest:unpreimage}{NW2DqMzY-LSOi5-1}\eatline
\nwused{\\{NW2DqMzY-31WKtJ-1}}\nwidentdefs{\\{{\nwixident{FS0.dest{\_}preimage}}{FS0.dest:unpreimage}}}\nwidentuses{\\{{\nwixident{Sort.CLASS}}{Sort.CLASS}}}\nwindexuse{\nwixident{Sort.CLASS}}{Sort.CLASS}{NW2DqMzY-LSOi5-1}\nwendcode{}\nwbegindocs{72}\nwdocspar
\begin{node}
The union of two $FS_{0}$ classes $S\cup T$ cannot be optimized away,
so we leave it as a constructor.
\end{node}

\nwenddocs{}\nwbegincode{73}\sublabel{NW2DqMzY-33zuOn-3}\nwmargintag{{\nwtagstyle{}\subpageref{NW2DqMzY-33zuOn-3}}}\moddef{Specification for $FS_{0}$ classes~{\nwtagstyle{}\subpageref{NW2DqMzY-33zuOn-1}}}\plusendmoddef\nwstartdeflinemarkup\nwusesondefline{\\{NW2DqMzY-3R3yGC-1}}\nwprevnextdefs{NW2DqMzY-33zuOn-2}{NW2DqMzY-33zuOn-4}\nwenddeflinemarkup
val union : \nwlinkedidentc{Term.t}{NW1rTmOJ-4LsRWC-1} -> \nwlinkedidentc{Term.t}{NW1rTmOJ-4LsRWC-1} -> \nwlinkedidentc{Term.t}{NW1rTmOJ-4LsRWC-1};
val dest_union : \nwlinkedidentc{Term.t}{NW1rTmOJ-4LsRWC-1} -> \nwlinkedidentc{Term.t}{NW1rTmOJ-4LsRWC-1} * \nwlinkedidentc{Term.t}{NW1rTmOJ-4LsRWC-1};
\nwused{\\{NW2DqMzY-3R3yGC-1}}\nwidentuses{\\{{\nwixident{Term.t}}{Term.t}}}\nwindexuse{\nwixident{Term.t}}{Term.t}{NW2DqMzY-33zuOn-3}\nwendcode{}\nwbegindocs{74}\nwdocspar
\nwenddocs{}\nwbegincode{75}\sublabel{NW2DqMzY-22W7o7-1}\nwmargintag{{\nwtagstyle{}\subpageref{NW2DqMzY-22W7o7-1}}}\moddef{Unions of two $FS_{0}$ classes~{\nwtagstyle{}\subpageref{NW2DqMzY-22W7o7-1}}}\endmoddef\nwstartdeflinemarkup\nwusesondefline{\\{NW2DqMzY-TSvER-1}}\nwenddeflinemarkup
(* union : \nwlinkedidentc{Term.t}{NW1rTmOJ-4LsRWC-1} -> \nwlinkedidentc{Term.t}{NW1rTmOJ-4LsRWC-1} -> \nwlinkedidentc{Term.t}{NW1rTmOJ-4LsRWC-1} *)
fun union S T =
  if not(\nwlinkedidentc{Term.is_class}{NW1rTmOJ-3FfXTO-1} S)
  then raise Sort.Mismatch "\nwlinkedidentc{FS0.union}{NW2DqMzY-22W7o7-1}: first argument is not a class"
  else if not(\nwlinkedidentc{Term.is_class}{NW1rTmOJ-3FfXTO-1} T)
  then raise Sort.Mismatch "\nwlinkedidentc{FS0.union}{NW2DqMzY-22W7o7-1}: second argument is not a class"
  else Term.Fun("union", \nwlinkedidentc{Sort.CLASS}{NW1rTmOJ-3ZqggC-1}, [S, T]);

\LA{}Destructure \code{}union\edoc{}~{\nwtagstyle{}\subpageref{NW2DqMzY-27dSmE-1}}\RA{}
\nwindexdefn{\nwixident{FS0.union}}{FS0.union}{NW2DqMzY-22W7o7-1}\eatline
\nwused{\\{NW2DqMzY-TSvER-1}}\nwidentdefs{\\{{\nwixident{FS0.union}}{FS0.union}}}\nwidentuses{\\{{\nwixident{Sort.CLASS}}{Sort.CLASS}}\\{{\nwixident{Term.is{\_}class}}{Term.is:unclass}}\\{{\nwixident{Term.t}}{Term.t}}}\nwindexuse{\nwixident{Sort.CLASS}}{Sort.CLASS}{NW2DqMzY-22W7o7-1}\nwindexuse{\nwixident{Term.is{\_}class}}{Term.is:unclass}{NW2DqMzY-22W7o7-1}\nwindexuse{\nwixident{Term.t}}{Term.t}{NW2DqMzY-22W7o7-1}\nwendcode{}\nwbegindocs{76}\nwdocspar
\nwenddocs{}\nwbegincode{77}\sublabel{NW2DqMzY-27dSmE-1}\nwmargintag{{\nwtagstyle{}\subpageref{NW2DqMzY-27dSmE-1}}}\moddef{Destructure \code{}union\edoc{}~{\nwtagstyle{}\subpageref{NW2DqMzY-27dSmE-1}}}\endmoddef\nwstartdeflinemarkup\nwusesondefline{\\{NW2DqMzY-22W7o7-1}}\nwenddeflinemarkup
fun dest_union (Term.Fun("union", \nwlinkedidentc{Sort.CLASS}{NW1rTmOJ-3ZqggC-1}, [S, T])) = (S,T)
  | dest_union _ = raise Dest "\nwlinkedidentc{FS0.dest_union}{NW2DqMzY-27dSmE-1}";
\nwindexdefn{\nwixident{FS0.dest{\_}union}}{FS0.dest:ununion}{NW2DqMzY-27dSmE-1}\eatline
\nwused{\\{NW2DqMzY-22W7o7-1}}\nwidentdefs{\\{{\nwixident{FS0.dest{\_}union}}{FS0.dest:ununion}}}\nwidentuses{\\{{\nwixident{Sort.CLASS}}{Sort.CLASS}}}\nwindexuse{\nwixident{Sort.CLASS}}{Sort.CLASS}{NW2DqMzY-27dSmE-1}\nwendcode{}\nwbegindocs{78}\nwdocspar
\begin{node}
The intersection of two $FS_{0}$ classes $S\cap T$ cannot be optimized away.
\end{node}

\nwenddocs{}\nwbegincode{79}\sublabel{NW2DqMzY-33zuOn-4}\nwmargintag{{\nwtagstyle{}\subpageref{NW2DqMzY-33zuOn-4}}}\moddef{Specification for $FS_{0}$ classes~{\nwtagstyle{}\subpageref{NW2DqMzY-33zuOn-1}}}\plusendmoddef\nwstartdeflinemarkup\nwusesondefline{\\{NW2DqMzY-3R3yGC-1}}\nwprevnextdefs{NW2DqMzY-33zuOn-3}{NW2DqMzY-33zuOn-5}\nwenddeflinemarkup
val intersection : \nwlinkedidentc{Term.t}{NW1rTmOJ-4LsRWC-1} -> \nwlinkedidentc{Term.t}{NW1rTmOJ-4LsRWC-1} -> \nwlinkedidentc{Term.t}{NW1rTmOJ-4LsRWC-1};
val dest_intersection : \nwlinkedidentc{Term.t}{NW1rTmOJ-4LsRWC-1} -> \nwlinkedidentc{Term.t}{NW1rTmOJ-4LsRWC-1} * \nwlinkedidentc{Term.t}{NW1rTmOJ-4LsRWC-1};
\nwused{\\{NW2DqMzY-3R3yGC-1}}\nwidentuses{\\{{\nwixident{Term.t}}{Term.t}}}\nwindexuse{\nwixident{Term.t}}{Term.t}{NW2DqMzY-33zuOn-4}\nwendcode{}\nwbegindocs{80}\nwdocspar
\nwenddocs{}\nwbegincode{81}\sublabel{NW2DqMzY-19bkPi-1}\nwmargintag{{\nwtagstyle{}\subpageref{NW2DqMzY-19bkPi-1}}}\moddef{Intersection of two $FS_{0}$ classes~{\nwtagstyle{}\subpageref{NW2DqMzY-19bkPi-1}}}\endmoddef\nwstartdeflinemarkup\nwusesondefline{\\{NW2DqMzY-TSvER-1}}\nwenddeflinemarkup
(* intersection : \nwlinkedidentc{Term.t}{NW1rTmOJ-4LsRWC-1} -> \nwlinkedidentc{Term.t}{NW1rTmOJ-4LsRWC-1} -> \nwlinkedidentc{Term.t}{NW1rTmOJ-4LsRWC-1} *)
fun intersection S T =
  if not(\nwlinkedidentc{Term.is_class}{NW1rTmOJ-3FfXTO-1} S)
  then raise Sort.Mismatch "\nwlinkedidentc{FS0.intersection}{NW2DqMzY-19bkPi-1}: first argument is not a class"
  else if not(\nwlinkedidentc{Term.is_class}{NW1rTmOJ-3FfXTO-1} T)
  then raise Sort.Mismatch "\nwlinkedidentc{FS0.intersection}{NW2DqMzY-19bkPi-1}: second argument is not a class"
  else Term.Fun("intersection", \nwlinkedidentc{Sort.CLASS}{NW1rTmOJ-3ZqggC-1}, [S, T]);

\LA{}Destructure \code{}intersection\edoc{}~{\nwtagstyle{}\subpageref{NW2DqMzY-4EzxyM-1}}\RA{}
\nwindexdefn{\nwixident{FS0.intersection}}{FS0.intersection}{NW2DqMzY-19bkPi-1}\eatline
\nwused{\\{NW2DqMzY-TSvER-1}}\nwidentdefs{\\{{\nwixident{FS0.intersection}}{FS0.intersection}}}\nwidentuses{\\{{\nwixident{Sort.CLASS}}{Sort.CLASS}}\\{{\nwixident{Term.is{\_}class}}{Term.is:unclass}}\\{{\nwixident{Term.t}}{Term.t}}}\nwindexuse{\nwixident{Sort.CLASS}}{Sort.CLASS}{NW2DqMzY-19bkPi-1}\nwindexuse{\nwixident{Term.is{\_}class}}{Term.is:unclass}{NW2DqMzY-19bkPi-1}\nwindexuse{\nwixident{Term.t}}{Term.t}{NW2DqMzY-19bkPi-1}\nwendcode{}\nwbegindocs{82}\nwdocspar
\nwenddocs{}\nwbegincode{83}\sublabel{NW2DqMzY-4EzxyM-1}\nwmargintag{{\nwtagstyle{}\subpageref{NW2DqMzY-4EzxyM-1}}}\moddef{Destructure \code{}intersection\edoc{}~{\nwtagstyle{}\subpageref{NW2DqMzY-4EzxyM-1}}}\endmoddef\nwstartdeflinemarkup\nwusesondefline{\\{NW2DqMzY-19bkPi-1}}\nwenddeflinemarkup
fun dest_intersection (Term.Fun("intersection", \nwlinkedidentc{Sort.CLASS}{NW1rTmOJ-3ZqggC-1}, [S, T])) = (S,T)
  | dest_intersection _ = raise Dest "\nwlinkedidentc{FS0.dest_intersection}{NW2DqMzY-4EzxyM-1}";
\nwindexdefn{\nwixident{FS0.dest{\_}intersection}}{FS0.dest:unintersection}{NW2DqMzY-4EzxyM-1}\eatline
\nwused{\\{NW2DqMzY-19bkPi-1}}\nwidentdefs{\\{{\nwixident{FS0.dest{\_}intersection}}{FS0.dest:unintersection}}}\nwidentuses{\\{{\nwixident{Sort.CLASS}}{Sort.CLASS}}}\nwindexuse{\nwixident{Sort.CLASS}}{Sort.CLASS}{NW2DqMzY-4EzxyM-1}\nwendcode{}\nwbegindocs{84}\nwdocspar
\begin{node}
The inductively defined $FS_{0}$ class $\mathcal{I}(S,T)$ from the
``base class'' $S$ and the ``inductive generating class'' $T$ may be
optimized away, but there is no obvious way I can see.
\end{node}

\nwenddocs{}\nwbegincode{85}\sublabel{NW2DqMzY-33zuOn-5}\nwmargintag{{\nwtagstyle{}\subpageref{NW2DqMzY-33zuOn-5}}}\moddef{Specification for $FS_{0}$ classes~{\nwtagstyle{}\subpageref{NW2DqMzY-33zuOn-1}}}\plusendmoddef\nwstartdeflinemarkup\nwusesondefline{\\{NW2DqMzY-3R3yGC-1}}\nwprevnextdefs{NW2DqMzY-33zuOn-4}{\relax}\nwenddeflinemarkup
val induct : \nwlinkedidentc{Term.t}{NW1rTmOJ-4LsRWC-1} -> \nwlinkedidentc{Term.t}{NW1rTmOJ-4LsRWC-1} -> \nwlinkedidentc{Term.t}{NW1rTmOJ-4LsRWC-1};
val dest_induct : \nwlinkedidentc{Term.t}{NW1rTmOJ-4LsRWC-1} -> \nwlinkedidentc{Term.t}{NW1rTmOJ-4LsRWC-1} * \nwlinkedidentc{Term.t}{NW1rTmOJ-4LsRWC-1};
\nwused{\\{NW2DqMzY-3R3yGC-1}}\nwidentuses{\\{{\nwixident{Term.t}}{Term.t}}}\nwindexuse{\nwixident{Term.t}}{Term.t}{NW2DqMzY-33zuOn-5}\nwendcode{}\nwbegindocs{86}\nwdocspar

\nwenddocs{}\nwbegincode{87}\sublabel{NW2DqMzY-1mMvkv-1}\nwmargintag{{\nwtagstyle{}\subpageref{NW2DqMzY-1mMvkv-1}}}\moddef{Constructor for inductively defined $FS_{0}$ class~{\nwtagstyle{}\subpageref{NW2DqMzY-1mMvkv-1}}}\endmoddef\nwstartdeflinemarkup\nwusesondefline{\\{NW2DqMzY-TSvER-1}}\nwenddeflinemarkup
(* induct : \nwlinkedidentc{Term.t}{NW1rTmOJ-4LsRWC-1} -> \nwlinkedidentc{Term.t}{NW1rTmOJ-4LsRWC-1} -> \nwlinkedidentc{Term.t}{NW1rTmOJ-4LsRWC-1} *)
fun induct S T =
  if not(\nwlinkedidentc{Term.is_class}{NW1rTmOJ-3FfXTO-1} S)
  then raise Sort.Mismatch "FS0.induction: first argument is not a class"
  else if not(\nwlinkedidentc{Term.is_class}{NW1rTmOJ-3FfXTO-1} T)
  then raise Sort.Mismatch "FS0.induction: second argument is not a class"
  else Term.Fun("induction", \nwlinkedidentc{Sort.CLASS}{NW1rTmOJ-3ZqggC-1}, [S, T]);

\LA{}Destructure \code{}induct\edoc{}~{\nwtagstyle{}\subpageref{NW2DqMzY-4ToA2q-1}}\RA{}
\nwindexdefn{\nwixident{FS0.induct}}{FS0.induct}{NW2DqMzY-1mMvkv-1}\eatline
\nwused{\\{NW2DqMzY-TSvER-1}}\nwidentdefs{\\{{\nwixident{FS0.induct}}{FS0.induct}}}\nwidentuses{\\{{\nwixident{Sort.CLASS}}{Sort.CLASS}}\\{{\nwixident{Term.is{\_}class}}{Term.is:unclass}}\\{{\nwixident{Term.t}}{Term.t}}}\nwindexuse{\nwixident{Sort.CLASS}}{Sort.CLASS}{NW2DqMzY-1mMvkv-1}\nwindexuse{\nwixident{Term.is{\_}class}}{Term.is:unclass}{NW2DqMzY-1mMvkv-1}\nwindexuse{\nwixident{Term.t}}{Term.t}{NW2DqMzY-1mMvkv-1}\nwendcode{}\nwbegindocs{88}\nwdocspar
\nwenddocs{}\nwbegincode{89}\sublabel{NW2DqMzY-4ToA2q-1}\nwmargintag{{\nwtagstyle{}\subpageref{NW2DqMzY-4ToA2q-1}}}\moddef{Destructure \code{}induct\edoc{}~{\nwtagstyle{}\subpageref{NW2DqMzY-4ToA2q-1}}}\endmoddef\nwstartdeflinemarkup\nwusesondefline{\\{NW2DqMzY-1mMvkv-1}}\nwenddeflinemarkup
fun dest_induct (Term.Fun("induction", \nwlinkedidentc{Sort.CLASS}{NW1rTmOJ-3ZqggC-1}, [S, T])) = (S,T)
  | dest_induct _ = raise Dest "\nwlinkedidentc{FS0.dest_induct}{NW2DqMzY-4ToA2q-1}";
\nwindexdefn{\nwixident{FS0.dest{\_}induct}}{FS0.dest:uninduct}{NW2DqMzY-4ToA2q-1}\eatline
\nwused{\\{NW2DqMzY-1mMvkv-1}}\nwidentdefs{\\{{\nwixident{FS0.dest{\_}induct}}{FS0.dest:uninduct}}}\nwidentuses{\\{{\nwixident{Sort.CLASS}}{Sort.CLASS}}}\nwindexuse{\nwixident{Sort.CLASS}}{Sort.CLASS}{NW2DqMzY-4ToA2q-1}\nwendcode{}\nwbegindocs{90}\nwdocspar
\section*{Constructors for predicates}

\begin{node}
The three predicate constructors for $FS_{0}$ are for: 
\begin{enumerate}
\item set membership $\in$
\item equality of terms $=$
\item inequality of terms $\neq$ (which is just an abbreviation $x\neq y$
for $\neg(x=y)$).
\end{enumerate}
\end{node}

\nwenddocs{}\nwbegincode{91}\sublabel{NW2DqMzY-194t0N-1}\nwmargintag{{\nwtagstyle{}\subpageref{NW2DqMzY-194t0N-1}}}\moddef{Specification for $FS_{0}$ predicates~{\nwtagstyle{}\subpageref{NW2DqMzY-194t0N-1}}}\endmoddef\nwstartdeflinemarkup\nwusesondefline{\\{NW2DqMzY-3R3yGC-1}}\nwprevnextdefs{\relax}{NW2DqMzY-194t0N-2}\nwenddeflinemarkup
val not_equal : \nwlinkedidentc{Term.t}{NW1rTmOJ-4LsRWC-1} -> \nwlinkedidentc{Term.t}{NW1rTmOJ-4LsRWC-1} -> \nwlinkedidentc{Formula.t}{NW3w0iRO-2jnj9W-1};
\nwalsodefined{\\{NW2DqMzY-194t0N-2}\\{NW2DqMzY-194t0N-3}}\nwused{\\{NW2DqMzY-3R3yGC-1}}\nwidentuses{\\{{\nwixident{Formula.t}}{Formula.t}}\\{{\nwixident{Term.t}}{Term.t}}}\nwindexuse{\nwixident{Formula.t}}{Formula.t}{NW2DqMzY-194t0N-1}\nwindexuse{\nwixident{Term.t}}{Term.t}{NW2DqMzY-194t0N-1}\nwendcode{}\nwbegindocs{92}\nwdocspar
\nwenddocs{}\nwbegincode{93}\sublabel{NW2DqMzY-4OjF2N-1}\nwmargintag{{\nwtagstyle{}\subpageref{NW2DqMzY-4OjF2N-1}}}\moddef{Constructors for $FS_{0}$ predicates~{\nwtagstyle{}\subpageref{NW2DqMzY-4OjF2N-1}}}\endmoddef\nwstartdeflinemarkup\nwusesondefline{\\{NW2DqMzY-mVmU8-1}}\nwenddeflinemarkup
\LA{}Set membership constructor for $FS_{0}$~{\nwtagstyle{}\subpageref{NW2DqMzY-eTij6-1}}\RA{}

\LA{}Equality constructor for $FS_{0}$ terms~{\nwtagstyle{}\subpageref{NW2DqMzY-ocj48-1}}\RA{}

(* not_equal : \nwlinkedidentc{Term.t}{NW1rTmOJ-4LsRWC-1} -> \nwlinkedidentc{Term.t}{NW1rTmOJ-4LsRWC-1} -> \nwlinkedidentc{Formula.t}{NW3w0iRO-2jnj9W-1} *)
fun not_equal lhs rhs =
  \nwlinkedidentc{Formula.mk_not}{NW3w0iRO-3SlqND-2}(equal lhs rhs);
\nwindexdefn{\nwixident{FS0.not{\_}equal}}{FS0.not:unequal}{NW2DqMzY-4OjF2N-1}\eatline
\nwused{\\{NW2DqMzY-mVmU8-1}}\nwidentdefs{\\{{\nwixident{FS0.not{\_}equal}}{FS0.not:unequal}}}\nwidentuses{\\{{\nwixident{Formula.mk{\_}not}}{Formula.mk:unnot}}\\{{\nwixident{Formula.t}}{Formula.t}}\\{{\nwixident{Term.t}}{Term.t}}}\nwindexuse{\nwixident{Formula.mk{\_}not}}{Formula.mk:unnot}{NW2DqMzY-4OjF2N-1}\nwindexuse{\nwixident{Formula.t}}{Formula.t}{NW2DqMzY-4OjF2N-1}\nwindexuse{\nwixident{Term.t}}{Term.t}{NW2DqMzY-4OjF2N-1}\nwendcode{}\nwbegindocs{94}\nwdocspar
\begin{node}
The only primitive predicate for $FS_{0}$ is the set membership $\in$
which has the generic form $x\in C$ for any $FS_{0}$ individual-sorted
term $x$ and any class-sorted term $C$. 
\end{node}

\nwenddocs{}\nwbegincode{95}\sublabel{NW2DqMzY-194t0N-2}\nwmargintag{{\nwtagstyle{}\subpageref{NW2DqMzY-194t0N-2}}}\moddef{Specification for $FS_{0}$ predicates~{\nwtagstyle{}\subpageref{NW2DqMzY-194t0N-1}}}\plusendmoddef\nwstartdeflinemarkup\nwusesondefline{\\{NW2DqMzY-3R3yGC-1}}\nwprevnextdefs{NW2DqMzY-194t0N-1}{NW2DqMzY-194t0N-3}\nwenddeflinemarkup
val In : \nwlinkedidentc{Term.t}{NW1rTmOJ-4LsRWC-1} -> \nwlinkedidentc{Term.t}{NW1rTmOJ-4LsRWC-1} -> \nwlinkedidentc{Formula.t}{NW3w0iRO-2jnj9W-1};
val dest_in : \nwlinkedidentc{Formula.t}{NW3w0iRO-2jnj9W-1} -> \nwlinkedidentc{Term.t}{NW1rTmOJ-4LsRWC-1} * \nwlinkedidentc{Term.t}{NW1rTmOJ-4LsRWC-1};
\nwused{\\{NW2DqMzY-3R3yGC-1}}\nwidentuses{\\{{\nwixident{Formula.t}}{Formula.t}}\\{{\nwixident{Term.t}}{Term.t}}}\nwindexuse{\nwixident{Formula.t}}{Formula.t}{NW2DqMzY-194t0N-2}\nwindexuse{\nwixident{Term.t}}{Term.t}{NW2DqMzY-194t0N-2}\nwendcode{}\nwbegindocs{96}\nwdocspar

\nwenddocs{}\nwbegincode{97}\sublabel{NW2DqMzY-eTij6-1}\nwmargintag{{\nwtagstyle{}\subpageref{NW2DqMzY-eTij6-1}}}\moddef{Set membership constructor for $FS_{0}$~{\nwtagstyle{}\subpageref{NW2DqMzY-eTij6-1}}}\endmoddef\nwstartdeflinemarkup\nwusesondefline{\\{NW2DqMzY-4OjF2N-1}}\nwenddeflinemarkup
(* In : \nwlinkedidentc{Term.t}{NW1rTmOJ-4LsRWC-1} -> \nwlinkedidentc{Term.t}{NW1rTmOJ-4LsRWC-1} -> \nwlinkedidentc{Formula.t}{NW3w0iRO-2jnj9W-1} *)
fun In x C =
  if not(\nwlinkedidentc{Term.is_ind}{NW1rTmOJ-3FfXTO-1} x)
  then raise Sort.Mismatch "\nwlinkedidentc{FS0.In}{NW2DqMzY-eTij6-1}: first argument is not an ind"
  else if not(\nwlinkedidentc{Term.is_class}{NW1rTmOJ-3FfXTO-1} C)
  then raise Sort.Mismatch "\nwlinkedidentc{FS0.In}{NW2DqMzY-eTij6-1}: second argument is not a class"
  else \nwlinkedidentc{Formula.mk_pred}{NW3w0iRO-3SlqND-1}("in", [x, C]);

\LA{}Destructuring \code{}in\edoc{}~{\nwtagstyle{}\subpageref{NW2DqMzY-4I0O5b-1}}\RA{}
\nwindexdefn{\nwixident{FS0.In}}{FS0.In}{NW2DqMzY-eTij6-1}\eatline
\nwused{\\{NW2DqMzY-4OjF2N-1}}\nwidentdefs{\\{{\nwixident{FS0.In}}{FS0.In}}}\nwidentuses{\\{{\nwixident{Formula.mk{\_}pred}}{Formula.mk:unpred}}\\{{\nwixident{Formula.t}}{Formula.t}}\\{{\nwixident{Term.is{\_}class}}{Term.is:unclass}}\\{{\nwixident{Term.is{\_}ind}}{Term.is:unind}}\\{{\nwixident{Term.t}}{Term.t}}}\nwindexuse{\nwixident{Formula.mk{\_}pred}}{Formula.mk:unpred}{NW2DqMzY-eTij6-1}\nwindexuse{\nwixident{Formula.t}}{Formula.t}{NW2DqMzY-eTij6-1}\nwindexuse{\nwixident{Term.is{\_}class}}{Term.is:unclass}{NW2DqMzY-eTij6-1}\nwindexuse{\nwixident{Term.is{\_}ind}}{Term.is:unind}{NW2DqMzY-eTij6-1}\nwindexuse{\nwixident{Term.t}}{Term.t}{NW2DqMzY-eTij6-1}\nwendcode{}\nwbegindocs{98}\nwdocspar
\nwenddocs{}\nwbegincode{99}\sublabel{NW2DqMzY-4I0O5b-1}\nwmargintag{{\nwtagstyle{}\subpageref{NW2DqMzY-4I0O5b-1}}}\moddef{Destructuring \code{}in\edoc{}~{\nwtagstyle{}\subpageref{NW2DqMzY-4I0O5b-1}}}\endmoddef\nwstartdeflinemarkup\nwusesondefline{\\{NW2DqMzY-eTij6-1}}\nwenddeflinemarkup
fun dest_in fm =
  if not(\nwlinkedidentc{Formula.is_pred}{NW3w0iRO-3NyJLv-1} fm)
  then raise Dest "\nwlinkedidentc{FS0.dest_in}{NW2DqMzY-4I0O5b-1}"
  else (case \nwlinkedidentc{Formula.dest_pred}{NW3w0iRO-1Yt1Jl-1} fm of
         ("in", [x, C]) => (x, C)
       | _ => raise Dest "\nwlinkedidentc{FS0.dest_in}{NW2DqMzY-4I0O5b-1}");
\nwindexdefn{\nwixident{FS0.dest{\_}in}}{FS0.dest:unin}{NW2DqMzY-4I0O5b-1}\eatline
\nwused{\\{NW2DqMzY-eTij6-1}}\nwidentdefs{\\{{\nwixident{FS0.dest{\_}in}}{FS0.dest:unin}}}\nwidentuses{\\{{\nwixident{Formula.dest{\_}pred}}{Formula.dest:unpred}}\\{{\nwixident{Formula.is{\_}pred}}{Formula.is:unpred}}}\nwindexuse{\nwixident{Formula.dest{\_}pred}}{Formula.dest:unpred}{NW2DqMzY-4I0O5b-1}\nwindexuse{\nwixident{Formula.is{\_}pred}}{Formula.is:unpred}{NW2DqMzY-4I0O5b-1}\nwendcode{}\nwbegindocs{100}\nwdocspar
\begin{node}
There is also equality, but this is primitive for individual-sorted
terms. Equality of functions $f=g$ is an abbreviation for
\begin{equation}
\forall x\ldotp(fx=gx).
\end{equation}
Equality of sets $S=T$ is an abbreviation for
\begin{equation}
\forall x\ldotp(x\in S\iff x\in T).
\end{equation}
We just treat them as primitives, with these two abbreviations being
additional axioms for $FS_{0}$.
\end{node}

\nwenddocs{}\nwbegincode{101}\sublabel{NW2DqMzY-194t0N-3}\nwmargintag{{\nwtagstyle{}\subpageref{NW2DqMzY-194t0N-3}}}\moddef{Specification for $FS_{0}$ predicates~{\nwtagstyle{}\subpageref{NW2DqMzY-194t0N-1}}}\plusendmoddef\nwstartdeflinemarkup\nwusesondefline{\\{NW2DqMzY-3R3yGC-1}}\nwprevnextdefs{NW2DqMzY-194t0N-2}{\relax}\nwenddeflinemarkup
val equal : \nwlinkedidentc{Term.t}{NW1rTmOJ-4LsRWC-1} -> \nwlinkedidentc{Term.t}{NW1rTmOJ-4LsRWC-1} -> \nwlinkedidentc{Formula.t}{NW3w0iRO-2jnj9W-1};
val dest_equal : \nwlinkedidentc{Formula.t}{NW3w0iRO-2jnj9W-1} -> \nwlinkedidentc{Term.t}{NW1rTmOJ-4LsRWC-1} * \nwlinkedidentc{Term.t}{NW1rTmOJ-4LsRWC-1};
\nwused{\\{NW2DqMzY-3R3yGC-1}}\nwidentuses{\\{{\nwixident{Formula.t}}{Formula.t}}\\{{\nwixident{Term.t}}{Term.t}}}\nwindexuse{\nwixident{Formula.t}}{Formula.t}{NW2DqMzY-194t0N-3}\nwindexuse{\nwixident{Term.t}}{Term.t}{NW2DqMzY-194t0N-3}\nwendcode{}\nwbegindocs{102}\nwdocspar
\nwenddocs{}\nwbegincode{103}\sublabel{NW2DqMzY-ocj48-1}\nwmargintag{{\nwtagstyle{}\subpageref{NW2DqMzY-ocj48-1}}}\moddef{Equality constructor for $FS_{0}$ terms~{\nwtagstyle{}\subpageref{NW2DqMzY-ocj48-1}}}\endmoddef\nwstartdeflinemarkup\nwusesondefline{\\{NW2DqMzY-4OjF2N-1}}\nwenddeflinemarkup
(* equal : \nwlinkedidentc{Term.t}{NW1rTmOJ-4LsRWC-1} -> \nwlinkedidentc{Term.t}{NW1rTmOJ-4LsRWC-1} -> \nwlinkedidentc{Formula.t}{NW3w0iRO-2jnj9W-1}

TODO: consider 'expanding' this out for functions
 (f = g iff forall x. f x = g x) and for classes
 (S = T iff forall x.(x in S iff x in T)). *)
fun equal lhs rhs =
  if not(\nwlinkedidentc{Term.have_same_sort}{NW1rTmOJ-1WQNps-1} lhs rhs)
  then raise Sort.Mismatch "\nwlinkedidentc{FS0.equal}{NW2DqMzY-ocj48-1}: expects arguments to have same sort"
  else (\nwlinkedidentc{Formula.mk_pred}{NW3w0iRO-3SlqND-1}("=", [lhs, rhs]));

\LA{}Destructuring \code{}equal\edoc{}~{\nwtagstyle{}\subpageref{NW2DqMzY-2FxXJs-1}}\RA{}
\nwindexdefn{\nwixident{FS0.equal}}{FS0.equal}{NW2DqMzY-ocj48-1}\eatline
\nwused{\\{NW2DqMzY-4OjF2N-1}}\nwidentdefs{\\{{\nwixident{FS0.equal}}{FS0.equal}}}\nwidentuses{\\{{\nwixident{Formula.mk{\_}pred}}{Formula.mk:unpred}}\\{{\nwixident{Formula.t}}{Formula.t}}\\{{\nwixident{Term.have{\_}same{\_}sort}}{Term.have:unsame:unsort}}\\{{\nwixident{Term.t}}{Term.t}}}\nwindexuse{\nwixident{Formula.mk{\_}pred}}{Formula.mk:unpred}{NW2DqMzY-ocj48-1}\nwindexuse{\nwixident{Formula.t}}{Formula.t}{NW2DqMzY-ocj48-1}\nwindexuse{\nwixident{Term.have{\_}same{\_}sort}}{Term.have:unsame:unsort}{NW2DqMzY-ocj48-1}\nwindexuse{\nwixident{Term.t}}{Term.t}{NW2DqMzY-ocj48-1}\nwendcode{}\nwbegindocs{104}\nwdocspar
\nwenddocs{}\nwbegincode{105}\sublabel{NW2DqMzY-2FxXJs-1}\nwmargintag{{\nwtagstyle{}\subpageref{NW2DqMzY-2FxXJs-1}}}\moddef{Destructuring \code{}equal\edoc{}~{\nwtagstyle{}\subpageref{NW2DqMzY-2FxXJs-1}}}\endmoddef\nwstartdeflinemarkup\nwusesondefline{\\{NW2DqMzY-ocj48-1}}\nwenddeflinemarkup
fun dest_equal fm =
  if not(\nwlinkedidentc{Formula.is_pred}{NW3w0iRO-3NyJLv-1} fm)
  then raise Dest "\nwlinkedidentc{FS0.dest_equal}{NW2DqMzY-2FxXJs-1}"
  else (case \nwlinkedidentc{Formula.dest_pred}{NW3w0iRO-1Yt1Jl-1} fm of
         ("equal", [lhs, rhs]) => (lhs, rhs)
       | _ => raise Dest "\nwlinkedidentc{FS0.dest_equal}{NW2DqMzY-2FxXJs-1}");
\nwindexdefn{\nwixident{FS0.dest{\_}equal}}{FS0.dest:unequal}{NW2DqMzY-2FxXJs-1}\eatline
\nwused{\\{NW2DqMzY-ocj48-1}}\nwidentdefs{\\{{\nwixident{FS0.dest{\_}equal}}{FS0.dest:unequal}}}\nwidentuses{\\{{\nwixident{Formula.dest{\_}pred}}{Formula.dest:unpred}}\\{{\nwixident{Formula.is{\_}pred}}{Formula.is:unpred}}}\nwindexuse{\nwixident{Formula.dest{\_}pred}}{Formula.dest:unpred}{NW2DqMzY-2FxXJs-1}\nwindexuse{\nwixident{Formula.is{\_}pred}}{Formula.is:unpred}{NW2DqMzY-2FxXJs-1}\nwendcode{}\nwfilename{nw/state.nw}\nwbegindocs{0}% -*- mode: poly-noweb; noweb-code-mode: sml-mode; -*-
%%
%% state.nw
%% 
%% Made by Alex Nelson <pqnelson@gmail.com>
%% Login   <alex@lisp>
%% 
%% Started on  2026-04-05T11:00:55-0700
%% Last update 2026-04-05T11:00:55-0700
%%

\section{State}

\begin{node}
The ``big picture'' is that we are going to build something which
looks like a Lisp interpreter, but instead of the full power of a
Turing-complete Lisp, we will be using $FS_{0}$. This ``$FS_{0}$ Lisp''
can be used as a REPL, which then acts like a logical framework. In
other words, we can use it like a ``syntax-less'' proof assistant with
a fixed foundations of interest. We just ``extend'' the ``Lisp world''
with new block statements as ``just another command'', and instruct
``$FS_{0}$ Lisp'' what to do when processing them.

Or we can use $FS_{0}$ to prove properties about proof
assistants. That is to say, given a ``Lisp engine'', we can prove
theorems about it.

But just as Lisps have a notion of ``environment'' (for binding
identifiers to their values), we need a comparable notion for
$FS_{0}$. Unlike Lisp, there will be two notions of ``environment'':
one for $FS_{0}$ as the metalanguage, and another defined ``within''
the ``Lisp world'' for describing the object logic's
bindings\footnote{After all, when telling $FS_{0}$ how an object logic
handles a new definition, we add bindings to a data structure
manipulatable within the REPL itself.}.

We'll discuss more about implementing a foundations in $FS_{0}$ as a
logical framework when appropriate.
\end{node}

\begin{node}
We can add axioms to the proof assistant, we can define new terms, or
we can define new predicates. Predicates are always well-defined, so
they do not need any proof they are well-defined. We also have the
``initial'' state, corresponding to a ``fresh'' situation where we
only have loaded the primitive notions and axioms for $FS_{0}$.

However, terms need to be proven to be well-defined. This requires an
existence and uniqueness proof, which are the two theorems
required. There are two ways to define a term:
\begin{enumerate}
\item A new term $t$ equals some expression involving existing terms,
  which is an abbreviation (hence does not require any
  well-definedness proofs)
\item A new term $t$ satisfies some condition $P[t]$. This requires
  well-definedness conditions $\exists t\ldotp P[t]$ and
  $\forall t_{1},t_{2}\ldotp P[t_{1}]\implies(P[t_{2}]\implies t_{1}=t_{2})$.
\end{enumerate}
The only free variables which can appear in the definiens (right-hand side) are those
which appear in the definiendum (left-hand side --- remember:
definiendum := definiens).

The ``definitional theorem'' associated with abbreviations is ``lhs = rhs'',
for definitions it's just the defining condition (universally
quantified over free variables), and for predicates it's
``$\mbox{defiendum}\iff\mbox{definiens}$'' (again, universally
quantified over free variables).
\end{node}

\begin{node}
We will use some sort of symbol table to track the terms and
predicates defined. The keys in the table will be strings (the
``lexemes'' for the terms and predicates), the values will be the
record {\Tt{}{\nwlbrace}ind\ :\ \nwlinkedidentq{Thm.t}{NW4ICsJQ-YIPBT-1}\ option,\ fun\ :\ \nwlinkedidentq{Thm.t}{NW4ICsJQ-YIPBT-1}\ option,\ class\ :\ \nwlinkedidentq{Thm.t}{NW4ICsJQ-YIPBT-1}\ option,\ pred\ :\ \nwlinkedidentq{Thm.t}{NW4ICsJQ-YIPBT-1}\ option{\nwrbrace}\nwendquote}
storing the ``definitional theorem'' for each of these.
\end{node}

\begin{node}
We should think a bit about how to handle theories as
``packages''. For example, the ``provided'' primitives (like
{\Tt{}fs0.in\nwendquote}, {\Tt{}fs0.subset\nwendquote}) are prefixed with {\Tt{}fs0.\nwendquote} (the theory
name, with period as a separator). One way to handle this would be to
have a {\Tt{}string\ ->\ SymTable\nwendquote} dictionary. Another way would be to have
the basic type be:
\begin{verbatim}
structure Table = MkRedBlackTree(String);
datatype t = Theory of t Table.t
           | PrimitiveNotion
           | Entry of { ind   : Thm.t option, 
                        fun   : Thm.t option,
                        class : Thm.t option,
                        pred  : Thm.t option };
\end{verbatim}
Then looking up an entry requires ``splitting'' the string by periods,
and recursively looking through theories as needed.
\end{node}

\nwenddocs{}\nwbegincode{1}\sublabel{NWU4iYe-1UXQzr-1}\nwmargintag{{\nwtagstyle{}\subpageref{NWU4iYe-1UXQzr-1}}}\moddef{State.sig~{\nwtagstyle{}\subpageref{NWU4iYe-1UXQzr-1}}}\endmoddef\nwstartdeflinemarkup\nwenddeflinemarkup
signature State = sig
  type t;

  val initial : unit -> t;

  (* expects first \nwlinkedidentc{Term.t}{NW1rTmOJ-4LsRWC-1} to be Fun(..,..,[Var, Var, ..., Var]) *)
  val abbrev_term : t -> \nwlinkedidentc{Term.t}{NW1rTmOJ-4LsRWC-1} -> \nwlinkedidentc{Term.t}{NW1rTmOJ-4LsRWC-1} -> t;
  val def_term : t -> \nwlinkedidentc{Term.t}{NW1rTmOJ-4LsRWC-1} -> \nwlinkedidentc{Formula.t}{NW3w0iRO-2jnj9W-1} -> \{ existence : \nwlinkedidentc{Thm.t}{NW4ICsJQ-YIPBT-1},
                                              uniqueness : \nwlinkedidentc{Thm.t}{NW4ICsJQ-YIPBT-1}\} -> t;
  val get_term : t -> string -> \nwlinkedidentc{Sort.t}{NW1rTmOJ-3ZqggC-1} -> \nwlinkedidentc{Term.t}{NW1rTmOJ-4LsRWC-1} option * t;
  val get_term_defthm : t -> string -> \nwlinkedidentc{Sort.t}{NW1rTmOJ-3ZqggC-1} -> \nwlinkedidentc{Thm.t}{NW4ICsJQ-YIPBT-1} option * t;
  
  (* expects first \nwlinkedidentc{Formula.t}{NW3w0iRO-2jnj9W-1} to be \nwlinkedidentc{Pred}{NW3w0iRO-2jnj9W-1}(..,..,[Var, Var, ..., Var]) *)
  val def_pred : t -> \nwlinkedidentc{Formula.t}{NW3w0iRO-2jnj9W-1} -> \nwlinkedidentc{Formula.t}{NW3w0iRO-2jnj9W-1} -> t;
  val get_pred : t -> string -> \nwlinkedidentc{Thm.t}{NW4ICsJQ-YIPBT-1} option * t;
end;
\nwnotused{State.sig}\nwidentuses{\\{{\nwixident{Formula.t}}{Formula.t}}\\{{\nwixident{Pred}}{Pred}}\\{{\nwixident{Sort.t}}{Sort.t}}\\{{\nwixident{Term.t}}{Term.t}}\\{{\nwixident{Thm.t}}{Thm.t}}}\nwindexuse{\nwixident{Formula.t}}{Formula.t}{NWU4iYe-1UXQzr-1}\nwindexuse{\nwixident{Pred}}{Pred}{NWU4iYe-1UXQzr-1}\nwindexuse{\nwixident{Sort.t}}{Sort.t}{NWU4iYe-1UXQzr-1}\nwindexuse{\nwixident{Term.t}}{Term.t}{NWU4iYe-1UXQzr-1}\nwindexuse{\nwixident{Thm.t}}{Thm.t}{NWU4iYe-1UXQzr-1}\nwendcode{}\nwbegindocs{2}\nwdocspar

\nwenddocs{}\nwfilename{nw/derived.nw}\nwbegindocs{0}% -*- mode: poly-noweb; noweb-code-mode: sml-mode; -*-
%%
%% derived.nw
%% 
%% Made by Alex Nelson <pqnelson@gmail.com>
%% Login   <alex@lisp>
%% 
%% Started on  2026-04-07T10:00:55-0700
%% Last update 2026-04-07T10:00:55-0700
%% 

\section{Derived inference rules}

\epigraph{This is a dull and grueling hell. Everything is bitter and
  leaden, cold and hot, hard and cruel, indistinct and stark, loud and
  random. There is no story to follow, just a succession of impacts, a
  sustained terror, a catalog of pain and exertions\dots}{Dan Abnett,\\
  ``Fragments (All We Have Left)''}% (2025)}

\begin{node}
While natural deduction is more friendly than Hilbert systems, there
are still quite a few theorems needed for making life easier. Further,
there are a few derived inference rules which are useful (which are
\SML/ functions eating in theorems and producing a theorem).

There's no way to present this material in a friendly way. Logic is
not a place to find friends. We will present a ``table of contents''
for results in Intuitionistic propositional logic (\S\ref{node:fs0:derived.nw:toc:propositional-logic})
and another for Intuitionistic first-order logic (\S\ref{node:fs0:derived.nw:toc:predicate-logic}).
\end{node}

\nwenddocs{}\nwbegincode{1}\sublabel{NW1NjfzY-4dPl0p-1}\nwmargintag{{\nwtagstyle{}\subpageref{NW1NjfzY-4dPl0p-1}}}\moddef{Derived.sml~{\nwtagstyle{}\subpageref{NW1NjfzY-4dPl0p-1}}}\endmoddef\nwstartdeflinemarkup\nwenddeflinemarkup
structure Derived = struct
\LA{}Derived results for Intuitionistic propositional logic~{\nwtagstyle{}\subpageref{NW1NjfzY-4CUt2L-1}}\RA{}

\LA{}Derived results for Intuitionistic predicate logic~{\nwtagstyle{}\subpageref{NW1NjfzY-uZfts-1}}\RA{}
end;
\nwnotused{Derived.sml}\nwendcode{}\nwbegindocs{2}\nwdocspar

\begin{node}\label{node:fs0:derived.nw:toc:propositional-logic}
There are a lot of useful theorems to prove, but I have not done so
yet. So I am keeping a list.
\begin{enumerate}\setcounter{enumi}{\value{node}}
\item Given some list of formulas $\Gamma_{1}$ and a theorem
  $\Gamma_{2}\vdash A$, infer $\Gamma_{1}\cup\Gamma_{2}\vdash A$
\item Given $\Gamma_{1}\vdash A\implies B$ and $\Gamma_{2}\vdash B\implies C$,
infer $(\Gamma_{1}\cup\Gamma_{2})\setminus\{A\}\vdash A\implies C$.
\item \checkbox{1} Weak $\neg$-elimination: $\vdash A\implies(\neg A\implies B)$
\item \checkbox{1} Absurdum: $\vdash(A\implies B)\implies((A\implies\neg B)\implies\neg A)$

\nwenddocs{}\nwbegincode{3}\sublabel{NW1NjfzY-4CUt2L-1}\nwmargintag{{\nwtagstyle{}\subpageref{NW1NjfzY-4CUt2L-1}}}\moddef{Derived results for Intuitionistic propositional logic~{\nwtagstyle{}\subpageref{NW1NjfzY-4CUt2L-1}}}\endmoddef\nwstartdeflinemarkup\nwusesondefline{\\{NW1NjfzY-4dPl0p-1}}\nwprevnextdefs{\relax}{NW1NjfzY-4CUt2L-2}\nwenddeflinemarkup
\LA{}Weakening inference rule~{\nwtagstyle{}\subpageref{NW1NjfzY-2n5HPL-1}}\RA{}

\LA{}Hypothetical syllogism derived inference rule~{\nwtagstyle{}\subpageref{NW1NjfzY-2CxI45-1}}\RA{}

\LA{}Weak negation elimination~{\nwtagstyle{}\subpageref{NW1NjfzY-16XbWD-1}}\RA{}

\LA{}Prove $\vdash(A\implies B)\implies((A\implies\neg B)\implies\neg A)$~{\nwtagstyle{}\subpageref{NW1NjfzY-4ZKGAh-1}}\RA{}
\nwalsodefined{\\{NW1NjfzY-4CUt2L-2}\\{NW1NjfzY-4CUt2L-3}\\{NW1NjfzY-4CUt2L-4}\\{NW1NjfzY-4CUt2L-5}\\{NW1NjfzY-4CUt2L-6}\\{NW1NjfzY-4CUt2L-7}}\nwused{\\{NW1NjfzY-4dPl0p-1}}\nwendcode{}\nwbegindocs{4}\nwdocspar
% OR
\item \checkbox{1} $\vdash(A\lor A)\implies A$
\item \checkbox{1} $\vdash(A\lor B)\implies(B\lor A)$
\item \checkbox{1} Associativity: $\vdash(A\lor(B\lor C))\iff((A\lor B)\lor C)$
\item \checkbox{1} Left congruence: $\vdash(A\implies A')\implies(A\lor B\implies A'\lor B)$
\item \checkbox{1} Right congruence: $\vdash(B\implies B')\implies(A\lor B\implies A\lor B')$
\nwenddocs{}\nwbegincode{5}\sublabel{NW1NjfzY-4CUt2L-2}\nwmargintag{{\nwtagstyle{}\subpageref{NW1NjfzY-4CUt2L-2}}}\moddef{Derived results for Intuitionistic propositional logic~{\nwtagstyle{}\subpageref{NW1NjfzY-4CUt2L-1}}}\plusendmoddef\nwstartdeflinemarkup\nwusesondefline{\\{NW1NjfzY-4dPl0p-1}}\nwprevnextdefs{NW1NjfzY-4CUt2L-1}{NW1NjfzY-4CUt2L-3}\nwenddeflinemarkup
\LA{}$\vdash A\lor A\implies A$~{\nwtagstyle{}\subpageref{NW1NjfzY-3wJoEd-1}}\RA{}

\LA{}Disjunction is commutative~{\nwtagstyle{}\subpageref{NW1NjfzY-3WoMbH-1}}\RA{}

\LA{}Disjunction is associative~{\nwtagstyle{}\subpageref{NW1NjfzY-J2rA5-1}}\RA{}

\LA{}Congruence of left disjunct~{\nwtagstyle{}\subpageref{NW1NjfzY-4VDMqe-1}}\RA{}

\LA{}Congruence of right disjunct~{\nwtagstyle{}\subpageref{NW1NjfzY-lrWvY-1}}\RA{}
\nwused{\\{NW1NjfzY-4dPl0p-1}}\nwendcode{}\nwbegindocs{6}\nwdocspar
%\bigbreak% AND
\item \checkbox{1} $\vdash A\implies(B\implies A\land B)$
\item \checkbox{1} $\vdash(A\land B)\implies(B\land A)$
\item \checkbox{1} Associativity: $\vdash(A\land(B\land C))\iff((A\land B)\land C)$
\item \checkbox{1} Left congruence: $\vdash(A\implies A')\implies(A\land B\implies A'\land B)$
\item \checkbox{1} Right congruence: $\vdash(B\implies B')\implies(A\land B\implies A\land B')$

\nwenddocs{}\nwbegincode{7}\sublabel{NW1NjfzY-4CUt2L-3}\nwmargintag{{\nwtagstyle{}\subpageref{NW1NjfzY-4CUt2L-3}}}\moddef{Derived results for Intuitionistic propositional logic~{\nwtagstyle{}\subpageref{NW1NjfzY-4CUt2L-1}}}\plusendmoddef\nwstartdeflinemarkup\nwusesondefline{\\{NW1NjfzY-4dPl0p-1}}\nwprevnextdefs{NW1NjfzY-4CUt2L-2}{NW1NjfzY-4CUt2L-4}\nwenddeflinemarkup
\LA{}Conjunction tautology $\vdash A\implies(B\implies A\land B)$~{\nwtagstyle{}\subpageref{NW1NjfzY-1I5AY8-1}}\RA{}

\LA{}Conjunction is commutative~{\nwtagstyle{}\subpageref{NW1NjfzY-29ie6N-1}}\RA{}

\LA{}Conjunction is associative~{\nwtagstyle{}\subpageref{NW1NjfzY-328uZt-1}}\RA{}

\LA{}Conjunction left congruence~{\nwtagstyle{}\subpageref{NW1NjfzY-4YXcw4-1}}\RA{}

\LA{}Conjunction right congruence~{\nwtagstyle{}\subpageref{NW1NjfzY-TrI6l-1}}\RA{}
\nwused{\\{NW1NjfzY-4dPl0p-1}}\nwendcode{}\nwbegindocs{8}\nwdocspar
%\bigbreak% DISTRIBUTIVITY
\item \checkbox{1} $\vdash A\lor B\implies((A\implies C)\implies((B\implies D)\implies(C\lor D)))$
\item \checkbox{1} $\vdash A\land(B\lor C)\iff((A\land B)\lor(A\land C))$
\item \checkbox{1} $\vdash (A\lor B)\land C\iff((A\land C)\lor(B\land C))$
\item \checkbox{1} $\vdash A\lor(B\land C)\iff((A\lor B)\land(A\lor C))$
\item \checkbox{1} $\vdash (A\land B)\lor C\iff((A\lor C)\land(B\lor C))$
\nwenddocs{}\nwbegincode{9}\sublabel{NW1NjfzY-4CUt2L-4}\nwmargintag{{\nwtagstyle{}\subpageref{NW1NjfzY-4CUt2L-4}}}\moddef{Derived results for Intuitionistic propositional logic~{\nwtagstyle{}\subpageref{NW1NjfzY-4CUt2L-1}}}\plusendmoddef\nwstartdeflinemarkup\nwusesondefline{\\{NW1NjfzY-4dPl0p-1}}\nwprevnextdefs{NW1NjfzY-4CUt2L-3}{NW1NjfzY-4CUt2L-5}\nwenddeflinemarkup
\LA{}Disjunction two-sided congruence~{\nwtagstyle{}\subpageref{NW1NjfzY-h4O6T-1}}\RA{}

\LA{}$\vdash A\land(B\lor C)\iff((A\land B)\lor(A\land C))$~{\nwtagstyle{}\subpageref{NW1NjfzY-efA9w-1}}\RA{}

\LA{}Right distributivity of conjunction over disjunction~{\nwtagstyle{}\subpageref{NW1NjfzY-8du9Q-1}}\RA{}
\nwused{\\{NW1NjfzY-4dPl0p-1}}\nwendcode{}\nwbegindocs{10}\nwdocspar
%\bigbreak% IMPLIES
\item \checkbox{1} Reflexivity: $\vdash A\implies A$
\item \checkbox{1} Frege: $\vdash(A\implies(B\implies C))\implies((A\implies B)\implies(A\implies C))$
\item \checkbox{1} Currying: $\vdash(A\land B\implies C)\implies(A\implies(B\implies C))$
\item \checkbox{1} Uncurrying: $\vdash(A\implies(B\implies C))\implies(A\land B\implies C)$
\item \checkbox{1} Transitivity: $\vdash(A\implies B)\land(B\implies C)\implies(A\implies C)$
\item \checkbox{1} Hypothetical syllogism tautology: $\vdash(A\implies B)\implies((B\implies C)\implies(A\implies C))$
\item \checkbox{1} Contrapositive $\vdash(A\implies B)\implies(\neg B\implies\neg A)$
\item \checkbox{1} Contrapositive 2 $\vdash(A\implies\neg B)\implies(B\implies\neg A)$
\nwenddocs{}\nwbegincode{11}\sublabel{NW1NjfzY-4CUt2L-5}\nwmargintag{{\nwtagstyle{}\subpageref{NW1NjfzY-4CUt2L-5}}}\moddef{Derived results for Intuitionistic propositional logic~{\nwtagstyle{}\subpageref{NW1NjfzY-4CUt2L-1}}}\plusendmoddef\nwstartdeflinemarkup\nwusesondefline{\\{NW1NjfzY-4dPl0p-1}}\nwprevnextdefs{NW1NjfzY-4CUt2L-4}{NW1NjfzY-4CUt2L-6}\nwenddeflinemarkup
\LA{}Theorem: implications is reflexive~{\nwtagstyle{}\subpageref{NW1NjfzY-2PXiqP-1}}\RA{}

\LA{}Frege's tautology (propositional logic)~{\nwtagstyle{}\subpageref{NW1NjfzY-bcGD8-1}}\RA{}

\LA{}Curry implication~{\nwtagstyle{}\subpageref{NW1NjfzY-1iptma-1}}\RA{}

\LA{}Uncurrying implication~{\nwtagstyle{}\subpageref{NW1NjfzY-1JZnQr-1}}\RA{}

\LA{}Hypothetical syllogism~{\nwtagstyle{}\subpageref{NW1NjfzY-1KT69U-1}}\RA{}

\LA{}Transitivity of implies~{\nwtagstyle{}\subpageref{NW1NjfzY-1sm2GN-1}}\RA{}

\LA{}Contrapositive derived rules~{\nwtagstyle{}\subpageref{NW1NjfzY-4KFbFg-1}}\RA{}
\nwused{\\{NW1NjfzY-4dPl0p-1}}\nwendcode{}\nwbegindocs{12}\nwdocspar
\item \checkbox{1} $\vdash\neg(A\lor B)\iff((\neg A)\land(\neg B))$
\item \checkbox{1} $\vdash((\neg A)\lor(\neg B))\implies\neg(A\land B)$
\item \checkbox{1} $\vdash((\neg A)\lor B)\implies(A\implies B)$
\item \checkbox{1} $\vdash(A\implies B)\implies\neg(A\land\neg B)$
\nwenddocs{}\nwbegincode{13}\sublabel{NW1NjfzY-4CUt2L-6}\nwmargintag{{\nwtagstyle{}\subpageref{NW1NjfzY-4CUt2L-6}}}\moddef{Derived results for Intuitionistic propositional logic~{\nwtagstyle{}\subpageref{NW1NjfzY-4CUt2L-1}}}\plusendmoddef\nwstartdeflinemarkup\nwusesondefline{\\{NW1NjfzY-4dPl0p-1}}\nwprevnextdefs{NW1NjfzY-4CUt2L-5}{NW1NjfzY-4CUt2L-7}\nwenddeflinemarkup
\LA{}Negation of disjunction~{\nwtagstyle{}\subpageref{NW1NjfzY-4UPWei-1}}\RA{}

\LA{}Negation of conjunction~{\nwtagstyle{}\subpageref{NW1NjfzY-wYirn-1}}\RA{}

\LA{}$(\neg A)\lor B$ sufficient to prove $A\implies B$~{\nwtagstyle{}\subpageref{NW1NjfzY-3NUP7w-1}}\RA{}

\LA{}Implication as negated conjunction~{\nwtagstyle{}\subpageref{NW1NjfzY-2qUp6m-1}}\RA{}

\nwused{\\{NW1NjfzY-4dPl0p-1}}\nwendcode{}\nwbegindocs{14}\nwdocspar
%\bigbreak% IFF
\item \checkbox{1} $\vdash A\iff A$
\item \checkbox{1} $\vdash(A\iff B)\implies(B\iff A)$
\item \checkbox{1} Transitivity: $\vdash(A\iff B)\land(B\iff C)\implies(A\iff C)$
\item \checkbox{0} $\vdash((A\iff B)\land A)\implies B$ and
  $\vdash((A\iff B)\land B)\implies A$
\nwenddocs{}\nwbegincode{15}\sublabel{NW1NjfzY-4CUt2L-7}\nwmargintag{{\nwtagstyle{}\subpageref{NW1NjfzY-4CUt2L-7}}}\moddef{Derived results for Intuitionistic propositional logic~{\nwtagstyle{}\subpageref{NW1NjfzY-4CUt2L-1}}}\plusendmoddef\nwstartdeflinemarkup\nwusesondefline{\\{NW1NjfzY-4dPl0p-1}}\nwprevnextdefs{NW1NjfzY-4CUt2L-6}{\relax}\nwenddeflinemarkup
\LA{}Iff is reflexive~{\nwtagstyle{}\subpageref{NW1NjfzY-YtLDN-1}}\RA{}

\LA{}Iff is symmetric~{\nwtagstyle{}\subpageref{NW1NjfzY-3hb2X5-1}}\RA{}

\LA{}Iff is transitive~{\nwtagstyle{}\subpageref{NW1NjfzY-2isj9x-1}}\RA{}
\nwused{\\{NW1NjfzY-4dPl0p-1}}\nwendcode{}\nwbegindocs{16}\nwdocspar
\end{enumerate}
\end{node}

\begin{node}
Given some list of formulas $\Gamma_{1}$ and a theorem
$\Gamma_{2}\vdash A$, infer $\Gamma_{1}\cup\Gamma_{2}\vdash A$
\end{node}

\begin{proof}
\begin{enumerate}
\item Given $\Gamma_{2}\vdash A$
\item $A,\Gamma_{1}\vdash A$ by assuming
\item $\Gamma_{1}\vdash A\implies A$ by discharging $A$ from (2)
\item $\Gamma_{2}\cup\Gamma_{1}\vdash A$ by modus ponens(3, 1)\qedhere
\end{enumerate}
\end{proof}

\nwenddocs{}\nwbegincode{17}\sublabel{NW1NjfzY-2n5HPL-1}\nwmargintag{{\nwtagstyle{}\subpageref{NW1NjfzY-2n5HPL-1}}}\moddef{Weakening inference rule~{\nwtagstyle{}\subpageref{NW1NjfzY-2n5HPL-1}}}\endmoddef\nwstartdeflinemarkup\nwusesondefline{\\{NW1NjfzY-4CUt2L-1}}\nwenddeflinemarkup
fun weaken Gamma th1 =
  let
    val th2 = \nwlinkedidentc{Thm.assume}{NW4ICsJQ-2zKkPU-2} (\nwlinkedidentc{Thm.concl}{NW4ICsJQ-YIPBT-1} th1) Gamma;
    val th3 = \nwlinkedidentc{Thm.disch}{NW4ICsJQ-2zKkPU-2} (\nwlinkedidentc{Thm.concl}{NW4ICsJQ-YIPBT-1} th1) th2;
  in
    \nwlinkedidentc{Thm.modus_ponens}{NW4ICsJQ-2zKkPU-2} th2 th1
  end;
\nwindexdefn{\nwixident{Derived.weaken}}{Derived.weaken}{NW1NjfzY-2n5HPL-1}\eatline
\nwused{\\{NW1NjfzY-4CUt2L-1}}\nwidentdefs{\\{{\nwixident{Derived.weaken}}{Derived.weaken}}}\nwidentuses{\\{{\nwixident{Thm.assume}}{Thm.assume}}\\{{\nwixident{Thm.concl}}{Thm.concl}}\\{{\nwixident{Thm.disch}}{Thm.disch}}\\{{\nwixident{Thm.modus{\_}ponens}}{Thm.modus:unponens}}}\nwindexuse{\nwixident{Thm.assume}}{Thm.assume}{NW1NjfzY-2n5HPL-1}\nwindexuse{\nwixident{Thm.concl}}{Thm.concl}{NW1NjfzY-2n5HPL-1}\nwindexuse{\nwixident{Thm.disch}}{Thm.disch}{NW1NjfzY-2n5HPL-1}\nwindexuse{\nwixident{Thm.modus{\_}ponens}}{Thm.modus:unponens}{NW1NjfzY-2n5HPL-1}\nwendcode{}\nwbegindocs{18}\nwdocspar
\begin{node}
Given $\Gamma_{1}\vdash A\implies B$ and $\Gamma_{2}\vdash B\implies C$,
infer $(\Gamma_{1}\cup\Gamma_{2})\setminus\{A\}\vdash A\implies C$.
\end{node}

\begin{proof}
\begin{enumerate}
\item Given $\Gamma_{1}\vdash A\implies B$
\item Given $\Gamma_{2}\vdash B\implies C$
\item $A,\Gamma_{1}\vdash B$ by undischarging $A$ from (1)
\item $A,\Gamma_{1}\cup\Gamma_{2}\vdash C$ by modus ponens(2, 3)
\item $(\Gamma_{1}\cup\Gamma_{2})\setminus\{A\}\vdash A\implies C$ by discharging $A$\qedhere
\end{enumerate}
\end{proof}

\nwenddocs{}\nwbegincode{19}\sublabel{NW1NjfzY-2CxI45-1}\nwmargintag{{\nwtagstyle{}\subpageref{NW1NjfzY-2CxI45-1}}}\moddef{Hypothetical syllogism derived inference rule~{\nwtagstyle{}\subpageref{NW1NjfzY-2CxI45-1}}}\endmoddef\nwstartdeflinemarkup\nwusesondefline{\\{NW1NjfzY-4CUt2L-1}}\nwenddeflinemarkup
fun hypo_syll thm1 thm2 =
  let
    val (A,B) = \nwlinkedidentc{Formula.dest_imp}{NW3w0iRO-1Yt1Jl-3}(\nwlinkedidentc{Thm.concl}{NW4ICsJQ-YIPBT-1} thm1)
                handle (\nwlinkedidentc{Formula.Dest}{NW3w0iRO-2jnj9W-1} _) =>
                       raise Fail "hypo_syll: first theorem is not an implication";
    val th2 = \nwlinkedidentc{Thm.modus_ponens}{NW4ICsJQ-2zKkPU-2} thm2 (\nwlinkedidentc{Thm.undisch}{NW4ICsJQ-2zKkPU-2} thm1)
              handle (Thm.Fail _) =>
                     raise Fail "hypo_syll: conclusion of first theorem mismatches premise of second theorem";
  in
    \nwlinkedidentc{Thm.disch}{NW4ICsJQ-2zKkPU-2} A th2
  end;
\nwindexdefn{\nwixident{Derived.hypo{\_}syll}}{Derived.hypo:unsyll}{NW1NjfzY-2CxI45-1}\eatline
\nwused{\\{NW1NjfzY-4CUt2L-1}}\nwidentdefs{\\{{\nwixident{Derived.hypo{\_}syll}}{Derived.hypo:unsyll}}}\nwidentuses{\\{{\nwixident{Formula.Dest}}{Formula.Dest}}\\{{\nwixident{Formula.dest{\_}imp}}{Formula.dest:unimp}}\\{{\nwixident{Thm.concl}}{Thm.concl}}\\{{\nwixident{Thm.disch}}{Thm.disch}}\\{{\nwixident{Thm.modus{\_}ponens}}{Thm.modus:unponens}}\\{{\nwixident{Thm.undisch}}{Thm.undisch}}}\nwindexuse{\nwixident{Formula.Dest}}{Formula.Dest}{NW1NjfzY-2CxI45-1}\nwindexuse{\nwixident{Formula.dest{\_}imp}}{Formula.dest:unimp}{NW1NjfzY-2CxI45-1}\nwindexuse{\nwixident{Thm.concl}}{Thm.concl}{NW1NjfzY-2CxI45-1}\nwindexuse{\nwixident{Thm.disch}}{Thm.disch}{NW1NjfzY-2CxI45-1}\nwindexuse{\nwixident{Thm.modus{\_}ponens}}{Thm.modus:unponens}{NW1NjfzY-2CxI45-1}\nwindexuse{\nwixident{Thm.undisch}}{Thm.undisch}{NW1NjfzY-2CxI45-1}\nwendcode{}\nwbegindocs{20}\nwdocspar

\begin{node}
One possible Hilbert proof calculus for Intuitionistic logic has the
axioms
\begin{enumerate}
\item $A\implies(\neg A\implies B)$
\item $(A\implies B)\implies((A\implies\neg B)\implies\neg A)$
\end{enumerate}
We will obtain these as derived inference rules.

Note: Kleene referred to $\vdash A\implies(\neg A\implies B)$ as
``weak negation elimination'', so we will use it here.
\end{node}

\nwenddocs{}\nwbegincode{21}\sublabel{NW1NjfzY-16XbWD-1}\nwmargintag{{\nwtagstyle{}\subpageref{NW1NjfzY-16XbWD-1}}}\moddef{Weak negation elimination~{\nwtagstyle{}\subpageref{NW1NjfzY-16XbWD-1}}}\endmoddef\nwstartdeflinemarkup\nwusesondefline{\\{NW1NjfzY-4CUt2L-1}}\nwenddeflinemarkup
fun weak_not_elim A B =
  let
    val not_A = \nwlinkedidentc{Formula.mk_not}{NW3w0iRO-3SlqND-2} A;
    val th1 = \nwlinkedidentc{Thm.assume}{NW4ICsJQ-2zKkPU-2} A [not_A];     (* A, not A |- A *)
    val th2 = \nwlinkedidentc{Thm.assume}{NW4ICsJQ-2zKkPU-2} not_A [A];     (* A, not A |- not A *)
    val th3 = \nwlinkedidentc{Thm.not_elim}{NW4ICsJQ-2zKkPU-5} th2;         (* A, not A |- A implies false *)
    val th4 = \nwlinkedidentc{Thm.modus_ponens}{NW4ICsJQ-2zKkPU-2} th3 th1; (* A, not A |- false *)
    val th5 = \nwlinkedidentc{Thm.contr}{NW4ICsJQ-2zKkPU-6} B th4;          (* A, not A |- B *)
    val th6 = \nwlinkedidentc{Thm.disch}{NW4ICsJQ-2zKkPU-2} not_A th5;      (* A |- not A implies B *)
  in
    \nwlinkedidentc{Thm.disch}{NW4ICsJQ-2zKkPU-2} A th6
  end;
\nwindexdefn{\nwixident{Derived.weak{\_}not{\_}elim}}{Derived.weak:unnot:unelim}{NW1NjfzY-16XbWD-1}\eatline
\nwused{\\{NW1NjfzY-4CUt2L-1}}\nwidentdefs{\\{{\nwixident{Derived.weak{\_}not{\_}elim}}{Derived.weak:unnot:unelim}}}\nwidentuses{\\{{\nwixident{Formula.mk{\_}not}}{Formula.mk:unnot}}\\{{\nwixident{Thm.assume}}{Thm.assume}}\\{{\nwixident{Thm.contr}}{Thm.contr}}\\{{\nwixident{Thm.disch}}{Thm.disch}}\\{{\nwixident{Thm.modus{\_}ponens}}{Thm.modus:unponens}}\\{{\nwixident{Thm.not{\_}elim}}{Thm.not:unelim}}}\nwindexuse{\nwixident{Formula.mk{\_}not}}{Formula.mk:unnot}{NW1NjfzY-16XbWD-1}\nwindexuse{\nwixident{Thm.assume}}{Thm.assume}{NW1NjfzY-16XbWD-1}\nwindexuse{\nwixident{Thm.contr}}{Thm.contr}{NW1NjfzY-16XbWD-1}\nwindexuse{\nwixident{Thm.disch}}{Thm.disch}{NW1NjfzY-16XbWD-1}\nwindexuse{\nwixident{Thm.modus{\_}ponens}}{Thm.modus:unponens}{NW1NjfzY-16XbWD-1}\nwindexuse{\nwixident{Thm.not{\_}elim}}{Thm.not:unelim}{NW1NjfzY-16XbWD-1}\nwendcode{}\nwbegindocs{22}\nwdocspar
\begin{node}[Absurdum]\label{node:derived.nw:absurdum}
$\vdash(A\implies B)\implies((A\implies\neg B)\implies\neg A)$
\end{node}

\nwenddocs{}\nwbegincode{23}\sublabel{NW1NjfzY-4ZKGAh-1}\nwmargintag{{\nwtagstyle{}\subpageref{NW1NjfzY-4ZKGAh-1}}}\moddef{Prove $\vdash(A\implies B)\implies((A\implies\neg B)\implies\neg A)$~{\nwtagstyle{}\subpageref{NW1NjfzY-4ZKGAh-1}}}\endmoddef\nwstartdeflinemarkup\nwusesondefline{\\{NW1NjfzY-4CUt2L-1}}\nwenddeflinemarkup
fun absurdum A B =
  let
    val not_B = \nwlinkedidentc{Formula.mk_not}{NW3w0iRO-3SlqND-2} B;
    val A_imp_B = \nwlinkedidentc{Formula.mk_imp}{NW3w0iRO-3SlqND-3}(A, B);
    val A_imp_not_B = \nwlinkedidentc{Formula.mk_imp}{NW3w0iRO-3SlqND-3}(A, not_B);
    val th1 = \nwlinkedidentc{Thm.assume}{NW4ICsJQ-2zKkPU-2} A_imp_B [A_imp_not_B]; (* A => B, A => not B |- A => B *)
    val th2 = \nwlinkedidentc{Thm.undisch}{NW4ICsJQ-2zKkPU-2} th1;                  (* A, A => B, A => not B |- B *)
    val th3 = \nwlinkedidentc{Thm.assume}{NW4ICsJQ-2zKkPU-2} A_imp_not_B [A_imp_B]; (* A => B, A => not B |- A => not B *)
    val th4 = \nwlinkedidentc{Thm.undisch}{NW4ICsJQ-2zKkPU-2} th3;                  (* A, A => B, A => not B |- not B *)
    val th5 = \nwlinkedidentc{Thm.not_elim}{NW4ICsJQ-2zKkPU-5} th4;                 (* A, A => B, A => not B |- B => false *)
    val th6 = \nwlinkedidentc{Thm.disch}{NW4ICsJQ-2zKkPU-2} A (\nwlinkedidentc{Thm.modus_ponens}{NW4ICsJQ-2zKkPU-2} th5 th2); (* A => B, A => not B |- A => false *)
    val th7 = \nwlinkedidentc{Thm.not_intro}{NW4ICsJQ-2zKkPU-5} th6;                (* A => B, A => not B |- not A *)
    val th8 = \nwlinkedidentc{Thm.disch}{NW4ICsJQ-2zKkPU-2} A_imp_not_B th7;        (* A => B |- (A => not B) => not A *)
  in
    \nwlinkedidentc{Thm.disch}{NW4ICsJQ-2zKkPU-2} A_imp_B th8
  end;
\nwindexdefn{\nwixident{Derived.absurdum}}{Derived.absurdum}{NW1NjfzY-4ZKGAh-1}\eatline
\nwused{\\{NW1NjfzY-4CUt2L-1}}\nwidentdefs{\\{{\nwixident{Derived.absurdum}}{Derived.absurdum}}}\nwidentuses{\\{{\nwixident{Formula.mk{\_}imp}}{Formula.mk:unimp}}\\{{\nwixident{Formula.mk{\_}not}}{Formula.mk:unnot}}\\{{\nwixident{Thm.assume}}{Thm.assume}}\\{{\nwixident{Thm.disch}}{Thm.disch}}\\{{\nwixident{Thm.modus{\_}ponens}}{Thm.modus:unponens}}\\{{\nwixident{Thm.not{\_}elim}}{Thm.not:unelim}}\\{{\nwixident{Thm.not{\_}intro}}{Thm.not:unintro}}\\{{\nwixident{Thm.undisch}}{Thm.undisch}}}\nwindexuse{\nwixident{Formula.mk{\_}imp}}{Formula.mk:unimp}{NW1NjfzY-4ZKGAh-1}\nwindexuse{\nwixident{Formula.mk{\_}not}}{Formula.mk:unnot}{NW1NjfzY-4ZKGAh-1}\nwindexuse{\nwixident{Thm.assume}}{Thm.assume}{NW1NjfzY-4ZKGAh-1}\nwindexuse{\nwixident{Thm.disch}}{Thm.disch}{NW1NjfzY-4ZKGAh-1}\nwindexuse{\nwixident{Thm.modus{\_}ponens}}{Thm.modus:unponens}{NW1NjfzY-4ZKGAh-1}\nwindexuse{\nwixident{Thm.not{\_}elim}}{Thm.not:unelim}{NW1NjfzY-4ZKGAh-1}\nwindexuse{\nwixident{Thm.not{\_}intro}}{Thm.not:unintro}{NW1NjfzY-4ZKGAh-1}\nwindexuse{\nwixident{Thm.undisch}}{Thm.undisch}{NW1NjfzY-4ZKGAh-1}\nwendcode{}\nwbegindocs{24}\nwdocspar
\begin{node}
$\vdash A\lor A\implies A$
\end{node}

\begin{proof}
\begin{enumerate}
\item $A\vdash A$ by assuming $A$
\item $\vdash A\implies A$ by discharging $A$
\item $\vdash A\lor A\implies A$ by $\lor$-elim(2, 2)\qedhere
\end{enumerate}
\end{proof}

\nwenddocs{}\nwbegincode{25}\sublabel{NW1NjfzY-3wJoEd-1}\nwmargintag{{\nwtagstyle{}\subpageref{NW1NjfzY-3wJoEd-1}}}\moddef{$\vdash A\lor A\implies A$~{\nwtagstyle{}\subpageref{NW1NjfzY-3wJoEd-1}}}\endmoddef\nwstartdeflinemarkup\nwusesondefline{\\{NW1NjfzY-4CUt2L-2}}\nwenddeflinemarkup
fun or_idem A =
  let
    val th1 = \nwlinkedidentc{Thm.assume}{NW4ICsJQ-2zKkPU-2} A [];
    val th2 = \nwlinkedidentc{Thm.disch}{NW4ICsJQ-2zKkPU-2} A th1;
  in
    \nwlinkedidentc{Thm.or_elim}{NW4ICsJQ-2zKkPU-4} th2 th2
  end;
\nwused{\\{NW1NjfzY-4CUt2L-2}}\nwidentuses{\\{{\nwixident{Thm.assume}}{Thm.assume}}\\{{\nwixident{Thm.disch}}{Thm.disch}}\\{{\nwixident{Thm.or{\_}elim}}{Thm.or:unelim}}}\nwindexuse{\nwixident{Thm.assume}}{Thm.assume}{NW1NjfzY-3wJoEd-1}\nwindexuse{\nwixident{Thm.disch}}{Thm.disch}{NW1NjfzY-3wJoEd-1}\nwindexuse{\nwixident{Thm.or{\_}elim}}{Thm.or:unelim}{NW1NjfzY-3wJoEd-1}\nwendcode{}\nwbegindocs{26}\nwdocspar

\begin{node}
$\vdash (A\lor B)\implies(B\lor A)$
\end{node}

\begin{proof}
\begin{enumerate}
\item $A\vdash A$ by assuming $A$
\item $A\vdash B\lor A$ by $\lor$-intro-l($B$)
\item $\vdash A\implies(B\lor A)$ by discharging $A$ from (2)
\item $B\vdash B$ by assuming $B$
\item $B\vdash B\lor A$ by $\lor$-intro-r($A$)
\item $\vdash B\implies(B\lor A)$ by discharging $B$
\item $\vdash(A\lor B)\implies(B\lor A)$ by $\lor$-elim(3, 6)\qedhere
\end{enumerate}
\end{proof}

\nwenddocs{}\nwbegincode{27}\sublabel{NW1NjfzY-3WoMbH-1}\nwmargintag{{\nwtagstyle{}\subpageref{NW1NjfzY-3WoMbH-1}}}\moddef{Disjunction is commutative~{\nwtagstyle{}\subpageref{NW1NjfzY-3WoMbH-1}}}\endmoddef\nwstartdeflinemarkup\nwusesondefline{\\{NW1NjfzY-4CUt2L-2}}\nwenddeflinemarkup
fun or_comm A B =
  let
    val th1 = \nwlinkedidentc{Thm.assume}{NW4ICsJQ-2zKkPU-2} A [];
    val th2 = \nwlinkedidentc{Thm.or_intro_l}{NW4ICsJQ-2zKkPU-4} B th1;
    val th3 = \nwlinkedidentc{Thm.disch}{NW4ICsJQ-2zKkPU-2} A th2;
    val th4 = \nwlinkedidentc{Thm.assume}{NW4ICsJQ-2zKkPU-2} B [];
    val th5 = \nwlinkedidentc{Thm.or_intro_r}{NW4ICsJQ-2zKkPU-4} A th4;
    val th6 = \nwlinkedidentc{Thm.disch}{NW4ICsJQ-2zKkPU-2} B th5;
  in
    \nwlinkedidentc{Thm.or_elim}{NW4ICsJQ-2zKkPU-4} th3 th6
  end;
\nwindexdefn{\nwixident{Derived.or{\_}comm}}{Derived.or:uncomm}{NW1NjfzY-3WoMbH-1}\eatline
\nwused{\\{NW1NjfzY-4CUt2L-2}}\nwidentdefs{\\{{\nwixident{Derived.or{\_}comm}}{Derived.or:uncomm}}}\nwidentuses{\\{{\nwixident{Thm.assume}}{Thm.assume}}\\{{\nwixident{Thm.disch}}{Thm.disch}}\\{{\nwixident{Thm.or{\_}elim}}{Thm.or:unelim}}\\{{\nwixident{Thm.or{\_}intro{\_}l}}{Thm.or:unintro:unl}}\\{{\nwixident{Thm.or{\_}intro{\_}r}}{Thm.or:unintro:unr}}}\nwindexuse{\nwixident{Thm.assume}}{Thm.assume}{NW1NjfzY-3WoMbH-1}\nwindexuse{\nwixident{Thm.disch}}{Thm.disch}{NW1NjfzY-3WoMbH-1}\nwindexuse{\nwixident{Thm.or{\_}elim}}{Thm.or:unelim}{NW1NjfzY-3WoMbH-1}\nwindexuse{\nwixident{Thm.or{\_}intro{\_}l}}{Thm.or:unintro:unl}{NW1NjfzY-3WoMbH-1}\nwindexuse{\nwixident{Thm.or{\_}intro{\_}r}}{Thm.or:unintro:unr}{NW1NjfzY-3WoMbH-1}\nwendcode{}\nwbegindocs{28}\nwdocspar
\begin{node}
$\vdash((A\lor B)\lor C)\iff(A\lor(B\lor C))$
\end{node}

\begin{proof}
Claim 1: $\vdash((A\lor B)\lor C)\implies(A\lor(B\lor C))$
\begin{enumerate}
\item $A\vdash A$ by assuming
\item $A\vdash A\lor(B\lor C)$ by $\lor$-intro-r($B\lor C$)
\item\label{step:or-assoc:forward:coup-1} $\vdash A\implies(A\lor(B\lor C))$ by discharging
\item $B\vdash B$ by assuming
\item $B\vdash B\lor C$ by $\lor$-intro-r($C$)
\item $B\vdash A\lor(B\lor C)$ by $\lor$-intro-l($A$)
\item\label{step:or-assoc:forward:coup-2} $\vdash B\implies(A\lor(B\lor C))$ by discharging
\item $C\vdash C$ by assuming
\item $C\vdash B\lor C$ by $\lor$-intro-l($B$)
\item $C\vdash A\lor(B\lor C)$ by $\lor$-intro-l($A$)
\item\label{step:or-assoc:forward:coup-3} $\vdash C\implies(A\lor(B\lor C))$ by discharging
\item\label{step:or-assoc:forward:penultimate-step} $\vdash(A\lor B)\implies(A\lor(B\lor C))$ by $\lor$-elim(\ref{step:or-assoc:forward:coup-1}, \ref{step:or-assoc:forward:coup-2})
\item $\vdash((A\lor B)\lor C)\implies(A\lor(B\lor C))$ by $\lor$-elim(\ref{step:or-assoc:forward:penultimate-step}, \ref{step:or-assoc:forward:coup-3})
\end{enumerate}
Claim 2: $\vdash(A\lor(B\lor C))\implies((A\lor B)\lor C)$ by
virtually the same reasoning:
\begin{enumerate}
\item $A\vdash A$ by assuming $A$
\item $A\vdash A\lor B$ by $\lor$-intro-r($B$)
\item $A\vdash (A\lor B)\lor C$ by $\lor$-intro-r($C$)
\item\label{step:or-assoc:backward:coup-1} $\vdash A\implies((A\lor B)\lor C)$ by discharging $A$
\item $B\vdash B$ by assuming $B$
\item $B\vdash A\lor B$ by $\lor$-intro-l($A$)
\item $B\vdash(A\lor B)\lor C$ by $\lor$-intro-r($C$)
\item\label{step:or-assoc:backward:coup-2} $\vdash B\implies((A\lor B)\lor C)$ by discharging $B$
\item $C\vdash C$ by assuming $C$
\item $C\vdash(A\lor B)\lor C$ by $\lor$-intro-l($A\lor B$)
\item\label{step:or-assoc:backward:coup-3} $\vdash C\implies((A\lor B)\lor C)$ by discharging $C$
\item\label{step:or-assoc:backward:penultimate} $\vdash(B\lor C)\implies((A\lor B)\lor C)$ by $\lor$-elim(\ref{step:or-assoc:backward:coup-3}, \ref{step:or-assoc:backward:coup-2})
\item $\vdash(A\lor(B\lor C))\implies((A\lor B)\lor C)$ by $\lor$-elim(\ref{step:or-assoc:backward:penultimate},
\ref{step:or-assoc:backward:coup-1})
\end{enumerate}
Hence the result by $\Leftrightarrow$-introduction applied to claim 1
and claim 2.
\end{proof}

\nwenddocs{}\nwbegincode{29}\sublabel{NW1NjfzY-J2rA5-1}\nwmargintag{{\nwtagstyle{}\subpageref{NW1NjfzY-J2rA5-1}}}\moddef{Disjunction is associative~{\nwtagstyle{}\subpageref{NW1NjfzY-J2rA5-1}}}\endmoddef\nwstartdeflinemarkup\nwusesondefline{\\{NW1NjfzY-4CUt2L-2}}\nwenddeflinemarkup
local
  fun forward A B C =
    let
      val th1 = \nwlinkedidentc{Thm.assume}{NW4ICsJQ-2zKkPU-2} A [];
      val th2 = \nwlinkedidentc{Thm.or_intro_r}{NW4ICsJQ-2zKkPU-4} (\nwlinkedidentc{Formula.mk_or}{NW3w0iRO-3SlqND-3}(B, C)) th1;
      val th3 = \nwlinkedidentc{Thm.disch}{NW4ICsJQ-2zKkPU-2} A th2;
      val th4 = \nwlinkedidentc{Thm.assume}{NW4ICsJQ-2zKkPU-2} B [];
      val th5 = \nwlinkedidentc{Thm.or_intro_r}{NW4ICsJQ-2zKkPU-4} C th4;
      val th6 = \nwlinkedidentc{Thm.or_intro_l}{NW4ICsJQ-2zKkPU-4} A th5;
      val th7 = \nwlinkedidentc{Thm.disch}{NW4ICsJQ-2zKkPU-2} B th6;
      val th8 = \nwlinkedidentc{Thm.assume}{NW4ICsJQ-2zKkPU-2} C [];
      val th9 = \nwlinkedidentc{Thm.or_intro_l}{NW4ICsJQ-2zKkPU-4} B th8;
      val th10 = \nwlinkedidentc{Thm.or_intro_l}{NW4ICsJQ-2zKkPU-4} A th9;
      val th11 = \nwlinkedidentc{Thm.disch}{NW4ICsJQ-2zKkPU-2} C th10;
      val th12 = \nwlinkedidentc{Thm.or_elim}{NW4ICsJQ-2zKkPU-4} th3 th7;
    in
      \nwlinkedidentc{Thm.or_elim}{NW4ICsJQ-2zKkPU-4} th12 th11
    end;
  fun backward A B C =
    let
      val th1 = \nwlinkedidentc{Thm.assume}{NW4ICsJQ-2zKkPU-2} A [];
      val th2 = \nwlinkedidentc{Thm.or_intro_r}{NW4ICsJQ-2zKkPU-4} B th1;
      val th3 = \nwlinkedidentc{Thm.or_intro_r}{NW4ICsJQ-2zKkPU-4} C th2;
      val th4 = \nwlinkedidentc{Thm.disch}{NW4ICsJQ-2zKkPU-2} A th3;
      val th5 = \nwlinkedidentc{Thm.assume}{NW4ICsJQ-2zKkPU-2} B [];
      val th6 = \nwlinkedidentc{Thm.or_intro_l}{NW4ICsJQ-2zKkPU-4} A th5;
      val th7 = \nwlinkedidentc{Thm.or_intro_r}{NW4ICsJQ-2zKkPU-4} C th6;
      val th8 = \nwlinkedidentc{Thm.disch}{NW4ICsJQ-2zKkPU-2} B th7;
      val th9 = \nwlinkedidentc{Thm.assume}{NW4ICsJQ-2zKkPU-2} C [];
      val th10 = \nwlinkedidentc{Thm.or_intro_l}{NW4ICsJQ-2zKkPU-4} (\nwlinkedidentc{Formula.mk_or}{NW3w0iRO-3SlqND-3}(A, B)) th9;
      val th11 = \nwlinkedidentc{Thm.disch}{NW4ICsJQ-2zKkPU-2} C th10;
      val th12 = \nwlinkedidentc{Thm.or_elim}{NW4ICsJQ-2zKkPU-4} th11 th8;
    in
      \nwlinkedidentc{Thm.or_elim}{NW4ICsJQ-2zKkPU-4} th12 th4
    end;
in
fun or_assoc A B C =
  \nwlinkedidentc{Thm.iff_intro}{NW4ICsJQ-2zKkPU-7} (forward A B C) (backward A B C)
end;
\nwindexdefn{\nwixident{Derived.or{\_}assoc}}{Derived.or:unassoc}{NW1NjfzY-J2rA5-1}\eatline
\nwused{\\{NW1NjfzY-4CUt2L-2}}\nwidentdefs{\\{{\nwixident{Derived.or{\_}assoc}}{Derived.or:unassoc}}}\nwidentuses{\\{{\nwixident{Formula.mk{\_}or}}{Formula.mk:unor}}\\{{\nwixident{Thm.assume}}{Thm.assume}}\\{{\nwixident{Thm.disch}}{Thm.disch}}\\{{\nwixident{Thm.iff{\_}intro}}{Thm.iff:unintro}}\\{{\nwixident{Thm.or{\_}elim}}{Thm.or:unelim}}\\{{\nwixident{Thm.or{\_}intro{\_}l}}{Thm.or:unintro:unl}}\\{{\nwixident{Thm.or{\_}intro{\_}r}}{Thm.or:unintro:unr}}}\nwindexuse{\nwixident{Formula.mk{\_}or}}{Formula.mk:unor}{NW1NjfzY-J2rA5-1}\nwindexuse{\nwixident{Thm.assume}}{Thm.assume}{NW1NjfzY-J2rA5-1}\nwindexuse{\nwixident{Thm.disch}}{Thm.disch}{NW1NjfzY-J2rA5-1}\nwindexuse{\nwixident{Thm.iff{\_}intro}}{Thm.iff:unintro}{NW1NjfzY-J2rA5-1}\nwindexuse{\nwixident{Thm.or{\_}elim}}{Thm.or:unelim}{NW1NjfzY-J2rA5-1}\nwindexuse{\nwixident{Thm.or{\_}intro{\_}l}}{Thm.or:unintro:unl}{NW1NjfzY-J2rA5-1}\nwindexuse{\nwixident{Thm.or{\_}intro{\_}r}}{Thm.or:unintro:unr}{NW1NjfzY-J2rA5-1}\nwendcode{}\nwbegindocs{30}\nwdocspar
\begin{node}
$\vdash(A\implies A')\implies(A\lor B\implies A'\lor B)$
\end{node}

\begin{proof}
\begin{enumerate}
\item $A\implies A'\vdash A\implies A'$ by asssuming $A\implies A'$%1
\item $A,A\implies A'\vdash A'$ by undischarging%2
\item $A,A\implies A'\vdash A'\lor B$ by $\lor$-intro-r($B$)%3
\item $A\implies A'\vdash A\implies A'\lor B$ by discharging $A$%4
\item $A\implies A', B\vdash B$ by assuming%5
\item $A\implies A', B\vdash A'\lor B$ by $\lor$-intro-l($A'$)%6
\item $A\implies A'\vdash B\implies A'\lor B$ by discharging $B$%7
\item $A\implies A'\vdash A\lor B\implies A'\lor B$ by $\lor$-elim(4, 7)%8
\item $\vdash(A\implies A')\implies(A\lor B\implies A'\lor B)$ by discharging
\qedhere
\end{enumerate}
\end{proof}

\nwenddocs{}\nwbegincode{31}\sublabel{NW1NjfzY-4VDMqe-1}\nwmargintag{{\nwtagstyle{}\subpageref{NW1NjfzY-4VDMqe-1}}}\moddef{Congruence of left disjunct~{\nwtagstyle{}\subpageref{NW1NjfzY-4VDMqe-1}}}\endmoddef\nwstartdeflinemarkup\nwusesondefline{\\{NW1NjfzY-4CUt2L-2}}\nwenddeflinemarkup
fun or_cong_l A A' B =
  let
    val A_imp_A' = \nwlinkedidentc{Formula.mk_imp}{NW3w0iRO-3SlqND-3}(A, A');
    val th1 = \nwlinkedidentc{Thm.assume}{NW4ICsJQ-2zKkPU-2} A_imp_A' [];
    val th2 = \nwlinkedidentc{Thm.undisch}{NW4ICsJQ-2zKkPU-2} th1;
    val th3 = \nwlinkedidentc{Thm.or_intro_r}{NW4ICsJQ-2zKkPU-4} B th2;
    val th4 = \nwlinkedidentc{Thm.disch}{NW4ICsJQ-2zKkPU-2} A th3;
    val th5 = \nwlinkedidentc{Thm.assume}{NW4ICsJQ-2zKkPU-2} B [A_imp_A'];
    val th6 = \nwlinkedidentc{Thm.or_intro_l}{NW4ICsJQ-2zKkPU-4} A' th5;
    val th7 = \nwlinkedidentc{Thm.disch}{NW4ICsJQ-2zKkPU-2} B th6;
    val th8 = \nwlinkedidentc{Thm.or_elim}{NW4ICsJQ-2zKkPU-4} th4 th7;
  in
    \nwlinkedidentc{Thm.disch}{NW4ICsJQ-2zKkPU-2} A_imp_A' th8
  end;
\nwindexdefn{\nwixident{Derived.or{\_}cong{\_}l}}{Derived.or:uncong:unl}{NW1NjfzY-4VDMqe-1}\eatline
\nwused{\\{NW1NjfzY-4CUt2L-2}}\nwidentdefs{\\{{\nwixident{Derived.or{\_}cong{\_}l}}{Derived.or:uncong:unl}}}\nwidentuses{\\{{\nwixident{Formula.mk{\_}imp}}{Formula.mk:unimp}}\\{{\nwixident{Thm.assume}}{Thm.assume}}\\{{\nwixident{Thm.disch}}{Thm.disch}}\\{{\nwixident{Thm.or{\_}elim}}{Thm.or:unelim}}\\{{\nwixident{Thm.or{\_}intro{\_}l}}{Thm.or:unintro:unl}}\\{{\nwixident{Thm.or{\_}intro{\_}r}}{Thm.or:unintro:unr}}\\{{\nwixident{Thm.undisch}}{Thm.undisch}}}\nwindexuse{\nwixident{Formula.mk{\_}imp}}{Formula.mk:unimp}{NW1NjfzY-4VDMqe-1}\nwindexuse{\nwixident{Thm.assume}}{Thm.assume}{NW1NjfzY-4VDMqe-1}\nwindexuse{\nwixident{Thm.disch}}{Thm.disch}{NW1NjfzY-4VDMqe-1}\nwindexuse{\nwixident{Thm.or{\_}elim}}{Thm.or:unelim}{NW1NjfzY-4VDMqe-1}\nwindexuse{\nwixident{Thm.or{\_}intro{\_}l}}{Thm.or:unintro:unl}{NW1NjfzY-4VDMqe-1}\nwindexuse{\nwixident{Thm.or{\_}intro{\_}r}}{Thm.or:unintro:unr}{NW1NjfzY-4VDMqe-1}\nwindexuse{\nwixident{Thm.undisch}}{Thm.undisch}{NW1NjfzY-4VDMqe-1}\nwendcode{}\nwbegindocs{32}\nwdocspar
\begin{node}
$\vdash(B\implies B')\implies(A\lor B\implies A\lor B')$
\end{node}

\begin{proof}
\begin{enumerate}
\item $A,B\implies B'\vdash A$ by assuming%1
\item $A,B\implies B'\vdash A\lor B'$ by $\lor$-intro-r($B'$)%2
\item $B\implies B'\vdash A\implies A\lor B'$ by discharging $A$%3
\item $B\implies B'\vdash B\implies B'$ by assuming %4
\item $B,B\implies B'\vdash B'$ by undischarging%5
\item $B,B\implies B'\vdash A\lor B'$ by $\lor$-intro-l($A$)%6
\item $B\implies B'\vdash B\implies A\lor B'$ by discharging $B$%7
\item $B\implies B'\vdash A\lor B\implies A\lor B'$ by $\lor$-elim(3, 7)%8
\item $\vdash(B\implies B')\implies(A\lor B\implies A\lor B')$
  by undischarging 7\qedhere
\end{enumerate}
\end{proof}

\nwenddocs{}\nwbegincode{33}\sublabel{NW1NjfzY-lrWvY-1}\nwmargintag{{\nwtagstyle{}\subpageref{NW1NjfzY-lrWvY-1}}}\moddef{Congruence of right disjunct~{\nwtagstyle{}\subpageref{NW1NjfzY-lrWvY-1}}}\endmoddef\nwstartdeflinemarkup\nwusesondefline{\\{NW1NjfzY-4CUt2L-2}}\nwenddeflinemarkup
fun or_cong_r A B B' =
  let
    val B_imp_B' = \nwlinkedidentc{Formula.mk_imp}{NW3w0iRO-3SlqND-3}(B, B');
    val th1 = \nwlinkedidentc{Thm.assume}{NW4ICsJQ-2zKkPU-2} A [B_imp_B'];
    val th2 = \nwlinkedidentc{Thm.or_intro_r}{NW4ICsJQ-2zKkPU-4} B' th1;
    val th3 = \nwlinkedidentc{Thm.disch}{NW4ICsJQ-2zKkPU-2} A th2;
    val th4 = \nwlinkedidentc{Thm.assume}{NW4ICsJQ-2zKkPU-2} B_imp_B' [];
    val th5 = \nwlinkedidentc{Thm.undisch}{NW4ICsJQ-2zKkPU-2} th4;
    val th6 = \nwlinkedidentc{Thm.or_intro_l}{NW4ICsJQ-2zKkPU-4} A th5;
    val th7 = \nwlinkedidentc{Thm.disch}{NW4ICsJQ-2zKkPU-2} B th6;
    val th8 = \nwlinkedidentc{Thm.or_elim}{NW4ICsJQ-2zKkPU-4} th3 th7;
  in
    \nwlinkedidentc{Thm.disch}{NW4ICsJQ-2zKkPU-2} B_imp_B' th8
  end;
\nwindexdefn{\nwixident{Derived.or{\_}cong{\_}r}}{Derived.or:uncong:unr}{NW1NjfzY-lrWvY-1}\eatline
\nwused{\\{NW1NjfzY-4CUt2L-2}}\nwidentdefs{\\{{\nwixident{Derived.or{\_}cong{\_}r}}{Derived.or:uncong:unr}}}\nwidentuses{\\{{\nwixident{Formula.mk{\_}imp}}{Formula.mk:unimp}}\\{{\nwixident{Thm.assume}}{Thm.assume}}\\{{\nwixident{Thm.disch}}{Thm.disch}}\\{{\nwixident{Thm.or{\_}elim}}{Thm.or:unelim}}\\{{\nwixident{Thm.or{\_}intro{\_}l}}{Thm.or:unintro:unl}}\\{{\nwixident{Thm.or{\_}intro{\_}r}}{Thm.or:unintro:unr}}\\{{\nwixident{Thm.undisch}}{Thm.undisch}}}\nwindexuse{\nwixident{Formula.mk{\_}imp}}{Formula.mk:unimp}{NW1NjfzY-lrWvY-1}\nwindexuse{\nwixident{Thm.assume}}{Thm.assume}{NW1NjfzY-lrWvY-1}\nwindexuse{\nwixident{Thm.disch}}{Thm.disch}{NW1NjfzY-lrWvY-1}\nwindexuse{\nwixident{Thm.or{\_}elim}}{Thm.or:unelim}{NW1NjfzY-lrWvY-1}\nwindexuse{\nwixident{Thm.or{\_}intro{\_}l}}{Thm.or:unintro:unl}{NW1NjfzY-lrWvY-1}\nwindexuse{\nwixident{Thm.or{\_}intro{\_}r}}{Thm.or:unintro:unr}{NW1NjfzY-lrWvY-1}\nwindexuse{\nwixident{Thm.undisch}}{Thm.undisch}{NW1NjfzY-lrWvY-1}\nwendcode{}\nwbegindocs{34}\begin{node}
We have $\vdash A\implies(B\implies A\land B)$.
\end{node}

\begin{proof}
\begin{enumerate}
\item $A\vdash A$ by assuming
\item $B\vdash B$ by assuming
\item $A,B\vdash A\land B$ by $\land$-introduction(1, 2)
\item $A\vdash B\implies(A\land B)$ by discharging $B$ from 3
\item $\vdash A\implies(B\implies(A\land B))$ by discharging $A$ from 4\qedhere
\end{enumerate}
\end{proof}

\nwenddocs{}\nwbegincode{35}\sublabel{NW1NjfzY-1I5AY8-1}\nwmargintag{{\nwtagstyle{}\subpageref{NW1NjfzY-1I5AY8-1}}}\moddef{Conjunction tautology $\vdash A\implies(B\implies A\land B)$~{\nwtagstyle{}\subpageref{NW1NjfzY-1I5AY8-1}}}\endmoddef\nwstartdeflinemarkup\nwusesondefline{\\{NW1NjfzY-4CUt2L-3}}\nwenddeflinemarkup
fun and_thm A B =
  let
    val AB = \nwlinkedidentc{Formula.mk_and}{NW3w0iRO-3SlqND-3}(A, B);
    val th1 = \nwlinkedidentc{Thm.assume}{NW4ICsJQ-2zKkPU-2} A [B];
    val th2 = \nwlinkedidentc{Thm.assume}{NW4ICsJQ-2zKkPU-2} B [A];
    val th3 = Thm.and_intro th1 th2;
    val th4 = \nwlinkedidentc{Thm.disch}{NW4ICsJQ-2zKkPU-2} B th3;
  in
    \nwlinkedidentc{Thm.disch}{NW4ICsJQ-2zKkPU-2} A th4
  end;
\nwindexdefn{\nwixident{Derived.and{\_}thm}}{Derived.and:unthm}{NW1NjfzY-1I5AY8-1}\eatline
\nwused{\\{NW1NjfzY-4CUt2L-3}}\nwidentdefs{\\{{\nwixident{Derived.and{\_}thm}}{Derived.and:unthm}}}\nwidentuses{\\{{\nwixident{Formula.mk{\_}and}}{Formula.mk:unand}}\\{{\nwixident{Thm.assume}}{Thm.assume}}\\{{\nwixident{Thm.disch}}{Thm.disch}}}\nwindexuse{\nwixident{Formula.mk{\_}and}}{Formula.mk:unand}{NW1NjfzY-1I5AY8-1}\nwindexuse{\nwixident{Thm.assume}}{Thm.assume}{NW1NjfzY-1I5AY8-1}\nwindexuse{\nwixident{Thm.disch}}{Thm.disch}{NW1NjfzY-1I5AY8-1}\nwendcode{}\nwbegindocs{36}\nwdocspar
\begin{node}
$\vdash A\land B\implies B\land A$
\end{node}

\begin{proof}
\begin{enumerate}
\item $A\land B\vdash A\land B$ by assuming $A\land B$
\item $A\land B\vdash B$ by $\land$-elim-r(1)
\item $A\land B\vdash A$ by $\land$-elim-l(1)
\item $A\land B\vdash B\land A$ by $\land$-intro(2, 3)
\item $\vdash(A\land B)\implies(B\land A)$ by discharging $A\land B$\qedhere
\end{enumerate}
\end{proof}

\nwenddocs{}\nwbegincode{37}\sublabel{NW1NjfzY-29ie6N-1}\nwmargintag{{\nwtagstyle{}\subpageref{NW1NjfzY-29ie6N-1}}}\moddef{Conjunction is commutative~{\nwtagstyle{}\subpageref{NW1NjfzY-29ie6N-1}}}\endmoddef\nwstartdeflinemarkup\nwusesondefline{\\{NW1NjfzY-4CUt2L-3}}\nwenddeflinemarkup
fun and_comm A B =
  let
    val A_and_B = \nwlinkedidentc{Formula.mk_and}{NW3w0iRO-3SlqND-3}(A, B);
    val th1 = \nwlinkedidentc{Thm.assume}{NW4ICsJQ-2zKkPU-2} A_and_B [];
    val th2 = Thm.and_elim_r th1;
    val th3 = Thm.and_elim_l th1;
    val th4 = Thm.and_intro th2 th3;
  in
    \nwlinkedidentc{Thm.disch}{NW4ICsJQ-2zKkPU-2} A_and_B th4
  end;
\nwindexdefn{\nwixident{Derived.and{\_}comm}}{Derived.and:uncomm}{NW1NjfzY-29ie6N-1}\eatline
\nwused{\\{NW1NjfzY-4CUt2L-3}}\nwidentdefs{\\{{\nwixident{Derived.and{\_}comm}}{Derived.and:uncomm}}}\nwidentuses{\\{{\nwixident{Formula.mk{\_}and}}{Formula.mk:unand}}\\{{\nwixident{Thm.assume}}{Thm.assume}}\\{{\nwixident{Thm.disch}}{Thm.disch}}}\nwindexuse{\nwixident{Formula.mk{\_}and}}{Formula.mk:unand}{NW1NjfzY-29ie6N-1}\nwindexuse{\nwixident{Thm.assume}}{Thm.assume}{NW1NjfzY-29ie6N-1}\nwindexuse{\nwixident{Thm.disch}}{Thm.disch}{NW1NjfzY-29ie6N-1}\nwendcode{}\nwbegindocs{38}\nwdocspar
\begin{node}
$\vdash(A\land(B\land C))\iff((A\land B)\land C)$
\end{node}

\begin{proof}
\begin{enumerate}
\item $(A\land(B\land C))\vdash(A\land(B\land C))$ by assuming $(A\land(B\land C))$
\item $(A\land(B\land C))\vdash A$ by $\land$-elim-l(1)%2
\item $(A\land(B\land C))\vdash(B\land C)$ by $\land$-elim-r(1)
\item $(A\land(B\land C))\vdash B$ by $\land$-elim-l(3)%3
\item $(A\land(B\land C))\vdash C$ by $\land$-elim-r(3)%4
\item $(A\land(B\land C))\vdash A\land B$ by $\land$-intro(2, 4)%5
\item $(A\land(B\land C))\vdash(A\land B)\land C$ by $\land$-intro(5, 4)%6
\item $\vdash(A\land(B\land C))\implies((A\land B)\land C$ by discharging%7
\item $((A\land B)\land C)\vdash((A\land B)\land C)$ by assuming $((A\land B)\land C)$%8
\item $((A\land B)\land C)\vdash(A\land B)$ by $\land$-elim-l(8)%9
\item $((A\land B)\land C)\vdash A$ by $\land$-elim-l(9)%10
\item $((A\land B)\land C)\vdash B$ by $\land$-elim-r(9)%11
\item $((A\land B)\land C)\vdash C$ by $\land$-elim-r(8)%12
\item $((A\land B)\land C)\vdash B\land C$ by $\land$-intro(11, 12)%13
\item $((A\land B)\land C)\vdash A\land(B\land C)$ by $\land$-intro(10, 13)%14
\item $\vdash((A\land B)\land C)\implies(A\land(B\land C))$ by
  discharging 14%15
\item $\vdash(A\land(B\land C))\iff((A\land B)\land C)$ by $\Leftrightarrow$-intro(7, 15)\qedhere
\end{enumerate}
\end{proof}

\nwenddocs{}\nwbegincode{39}\sublabel{NW1NjfzY-328uZt-1}\nwmargintag{{\nwtagstyle{}\subpageref{NW1NjfzY-328uZt-1}}}\moddef{Conjunction is associative~{\nwtagstyle{}\subpageref{NW1NjfzY-328uZt-1}}}\endmoddef\nwstartdeflinemarkup\nwusesondefline{\\{NW1NjfzY-4CUt2L-3}}\nwenddeflinemarkup
local
  fun forward A B C =
    let
      val A_BC = \nwlinkedidentc{Formula.mk_and}{NW3w0iRO-3SlqND-3}(A, \nwlinkedidentc{Formula.mk_and}{NW3w0iRO-3SlqND-3}(B, C));
      val th1 = \nwlinkedidentc{Thm.assume}{NW4ICsJQ-2zKkPU-2} A_BC [];
      val th2 = Thm.and_elim_l th1; (* ... |- A *)
      val th3 = Thm.and_elim_r th1;
      val th4 = Thm.and_elim_l th3; (* ... |- B *)
      val th5 = Thm.and_elim_r th3; (* ... |- C *)
      val th6 = Thm.and_intro th2 th4;
      val th7 = Thm.and_intro th6 th5;
    in
      \nwlinkedidentc{Thm.disch}{NW4ICsJQ-2zKkPU-2} A_BC th7
    end;
  fun backward A B C =
    let
      val AB_C = \nwlinkedidentc{Formula.mk_and}{NW3w0iRO-3SlqND-3}(\nwlinkedidentc{Formula.mk_and}{NW3w0iRO-3SlqND-3}(A, B), C);
      val th1 = \nwlinkedidentc{Thm.assume}{NW4ICsJQ-2zKkPU-2} AB_C [];
      val th2 = Thm.and_elim_l th1; (* ... |- A & B *)
      val th3 = Thm.and_elim_l th2; (* ... |- A *)
      val th4 = Thm.and_elim_r th2; (* ... |- B *)
      val th5 = Thm.and_elim_r th1; (* ... |- C *)
      val th6 = Thm.and_intro th4 th5; (* ... |- B & C *)
      val th7 = Thm.and_intro th3 th6; (* ... |- A & (B & C) *)
    in
      \nwlinkedidentc{Thm.disch}{NW4ICsJQ-2zKkPU-2} AB_C th7
    end;
in
fun and_assoc A B C =
  \nwlinkedidentc{Thm.iff_intro}{NW4ICsJQ-2zKkPU-7} (forward A B C) (backward A B C)
end;
\nwindexdefn{\nwixident{Derived.and{\_}assoc}}{Derived.and:unassoc}{NW1NjfzY-328uZt-1}\eatline
\nwused{\\{NW1NjfzY-4CUt2L-3}}\nwidentdefs{\\{{\nwixident{Derived.and{\_}assoc}}{Derived.and:unassoc}}}\nwidentuses{\\{{\nwixident{Formula.mk{\_}and}}{Formula.mk:unand}}\\{{\nwixident{Thm.assume}}{Thm.assume}}\\{{\nwixident{Thm.disch}}{Thm.disch}}\\{{\nwixident{Thm.iff{\_}intro}}{Thm.iff:unintro}}}\nwindexuse{\nwixident{Formula.mk{\_}and}}{Formula.mk:unand}{NW1NjfzY-328uZt-1}\nwindexuse{\nwixident{Thm.assume}}{Thm.assume}{NW1NjfzY-328uZt-1}\nwindexuse{\nwixident{Thm.disch}}{Thm.disch}{NW1NjfzY-328uZt-1}\nwindexuse{\nwixident{Thm.iff{\_}intro}}{Thm.iff:unintro}{NW1NjfzY-328uZt-1}\nwendcode{}\nwbegindocs{40}\nwdocspar
\begin{node}
$\vdash(A\implies A')\implies(A\land B\implies A'\land B)$
\end{node}

\begin{proof}
\begin{enumerate}
\item $A\implies A', A\land B\vdash A\implies A'$ by assuming%1
\item $A\implies A', A\land B\vdash A\land B$ by assuming%2
\item $A\implies A', A\land B\vdash A$ by $\land$-elim-l(2)%3
\item $A\implies A', A\land B\vdash A'$ by modus ponens(1, 3)%4
\item $A\implies A', A\land B\vdash B$ by $\land$-elim-r(2)%5
\item $A\implies A', A\land B\vdash A'\land B$ by $\land$-intro(4,5)%6
\item $A\implies A'\vdash A\land B\implies A'\land B$ by discharging
\item $\vdash(A\implies A')\implies (A\land B\implies A'\land B)$ by discharging
\qedhere
\end{enumerate}
\end{proof}

\nwenddocs{}\nwbegincode{41}\sublabel{NW1NjfzY-4YXcw4-1}\nwmargintag{{\nwtagstyle{}\subpageref{NW1NjfzY-4YXcw4-1}}}\moddef{Conjunction left congruence~{\nwtagstyle{}\subpageref{NW1NjfzY-4YXcw4-1}}}\endmoddef\nwstartdeflinemarkup\nwusesondefline{\\{NW1NjfzY-4CUt2L-3}}\nwenddeflinemarkup
fun and_cong_l A A' B =
  let
    val A_imp_A' = \nwlinkedidentc{Formula.mk_imp}{NW3w0iRO-3SlqND-3}(A, A');
    val A_and_B = \nwlinkedidentc{Formula.mk_and}{NW3w0iRO-3SlqND-3}(A, B);
    val th1 = \nwlinkedidentc{Thm.assume}{NW4ICsJQ-2zKkPU-2} A_imp_A' [A_and_B];
    val th2 = \nwlinkedidentc{Thm.assume}{NW4ICsJQ-2zKkPU-2} A_and_B [A_imp_A'];
    val th3 = Thm.and_elim_l th2;
    val th4 = \nwlinkedidentc{Thm.modus_ponens}{NW4ICsJQ-2zKkPU-2} th1 th3;
    val th5 = Thm.and_elim_r th2;
    val th6 = Thm.and_intro th4 th5;
    val th7 = \nwlinkedidentc{Thm.disch}{NW4ICsJQ-2zKkPU-2} A_and_B th6;
  in
    \nwlinkedidentc{Thm.disch}{NW4ICsJQ-2zKkPU-2} A_imp_A' th7
  end;
\nwindexdefn{\nwixident{Derived.and{\_}cong{\_}l}}{Derived.and:uncong:unl}{NW1NjfzY-4YXcw4-1}\eatline
\nwused{\\{NW1NjfzY-4CUt2L-3}}\nwidentdefs{\\{{\nwixident{Derived.and{\_}cong{\_}l}}{Derived.and:uncong:unl}}}\nwidentuses{\\{{\nwixident{Formula.mk{\_}and}}{Formula.mk:unand}}\\{{\nwixident{Formula.mk{\_}imp}}{Formula.mk:unimp}}\\{{\nwixident{Thm.assume}}{Thm.assume}}\\{{\nwixident{Thm.disch}}{Thm.disch}}\\{{\nwixident{Thm.modus{\_}ponens}}{Thm.modus:unponens}}}\nwindexuse{\nwixident{Formula.mk{\_}and}}{Formula.mk:unand}{NW1NjfzY-4YXcw4-1}\nwindexuse{\nwixident{Formula.mk{\_}imp}}{Formula.mk:unimp}{NW1NjfzY-4YXcw4-1}\nwindexuse{\nwixident{Thm.assume}}{Thm.assume}{NW1NjfzY-4YXcw4-1}\nwindexuse{\nwixident{Thm.disch}}{Thm.disch}{NW1NjfzY-4YXcw4-1}\nwindexuse{\nwixident{Thm.modus{\_}ponens}}{Thm.modus:unponens}{NW1NjfzY-4YXcw4-1}\nwendcode{}\nwbegindocs{42}\nwdocspar
\begin{node}
$\vdash(B\implies B')\implies(A\land B\implies A\land B')$
\end{node}

\begin{proof}
\begin{enumerate}
\item $A\land B,B\implies B'\vdash B\implies B'$ by assuming%1
\item $A\land B,B\implies B'\vdash A\land B$ by assuming%2
\item $A\land B,B\implies B'\vdash A$ by $\land$-elim-l(2)%3
\item $A\land B,B\implies B'\vdash B$ by $\land$-elim-r(2)%4
\item $A\land B,B\implies B'\vdash B'$ by modus ponens(1, 4)%5
\item $A\land B,B\implies B'\vdash A\land B'$ by $\land$-intro(3, 5)%6
\item $B\implies B'\vdash A\land B \implies A\land B'$ by discharging%7
\item $\vdash(B\implies B')\implies(A\land B \implies A\land B')$ by discharging
\qedhere
\end{enumerate}
\end{proof}

\nwenddocs{}\nwbegincode{43}\sublabel{NW1NjfzY-TrI6l-1}\nwmargintag{{\nwtagstyle{}\subpageref{NW1NjfzY-TrI6l-1}}}\moddef{Conjunction right congruence~{\nwtagstyle{}\subpageref{NW1NjfzY-TrI6l-1}}}\endmoddef\nwstartdeflinemarkup\nwusesondefline{\\{NW1NjfzY-4CUt2L-3}}\nwenddeflinemarkup
fun and_cong_r A B B' =
  let
    val A_and_B = \nwlinkedidentc{Formula.mk_and}{NW3w0iRO-3SlqND-3}(A, B);
    val B_imp_B' = \nwlinkedidentc{Formula.mk_imp}{NW3w0iRO-3SlqND-3}(B, B');
    val th1 = \nwlinkedidentc{Thm.assume}{NW4ICsJQ-2zKkPU-2} B_imp_B' [];
    val th2 = \nwlinkedidentc{Thm.assume}{NW4ICsJQ-2zKkPU-2} A_and_B [];
    val th3 = Thm.and_elim_l th2;
    val th4 = Thm.and_elim_r th2;
    val th5 = \nwlinkedidentc{Thm.modus_ponens}{NW4ICsJQ-2zKkPU-2} th1 th4;
    val th6 = Thm.and_intro th3 th5;
    val th7 = \nwlinkedidentc{Thm.disch}{NW4ICsJQ-2zKkPU-2} A_and_B th6;
  in
    \nwlinkedidentc{Thm.disch}{NW4ICsJQ-2zKkPU-2} B_imp_B' th7
  end;
\nwindexdefn{\nwixident{Derived.and{\_}cong{\_}r}}{Derived.and:uncong:unr}{NW1NjfzY-TrI6l-1}\eatline
\nwused{\\{NW1NjfzY-4CUt2L-3}}\nwidentdefs{\\{{\nwixident{Derived.and{\_}cong{\_}r}}{Derived.and:uncong:unr}}}\nwidentuses{\\{{\nwixident{Formula.mk{\_}and}}{Formula.mk:unand}}\\{{\nwixident{Formula.mk{\_}imp}}{Formula.mk:unimp}}\\{{\nwixident{Thm.assume}}{Thm.assume}}\\{{\nwixident{Thm.disch}}{Thm.disch}}\\{{\nwixident{Thm.modus{\_}ponens}}{Thm.modus:unponens}}}\nwindexuse{\nwixident{Formula.mk{\_}and}}{Formula.mk:unand}{NW1NjfzY-TrI6l-1}\nwindexuse{\nwixident{Formula.mk{\_}imp}}{Formula.mk:unimp}{NW1NjfzY-TrI6l-1}\nwindexuse{\nwixident{Thm.assume}}{Thm.assume}{NW1NjfzY-TrI6l-1}\nwindexuse{\nwixident{Thm.disch}}{Thm.disch}{NW1NjfzY-TrI6l-1}\nwindexuse{\nwixident{Thm.modus{\_}ponens}}{Thm.modus:unponens}{NW1NjfzY-TrI6l-1}\nwendcode{}\nwbegindocs{44}\nwdocspar
\begin{node}[Krzysztof Grabczewski]
$\vdash A\lor B\implies((A\implies C)\implies((B\implies D)\implies(C\lor D)))$
\end{node}

\begin{proof}
\begin{enumerate}
\item $\vdash(A\implies C)\implies((A\lor B)\implies(C\lor B))$ by
  disjunction congruence left%1
\item $A\lor B,A\implies C\vdash C\lor B$ by undischarging (1) twice%2
\item $\vdash(B\implies D)\implies((C\lor B)\implies(C\lor D))$ by
  disjunction congruence right%3
\item $B\implies D\vdash(C\lor B)\implies(C\lor D)$ by undischarging(3)%4
\item $A\lor B,A\implies C,B\implies D\vdash C\lor D$ by modus
  ponens(4, 2)
\item $A\lor B,A\implies C\vdash(B\implies D)\implies(C\lor D)$
\item $A\lor B\vdash(A\implies C)\implies((B\implies D)\implies(C\lor D))$
\item $\vdash(A\lor B)\implies((A\implies C)\implies((B\implies D)\implies(C\lor D)))$\qedhere
\end{enumerate}
\end{proof}

\nwenddocs{}\nwbegincode{45}\sublabel{NW1NjfzY-h4O6T-1}\nwmargintag{{\nwtagstyle{}\subpageref{NW1NjfzY-h4O6T-1}}}\moddef{Disjunction two-sided congruence~{\nwtagstyle{}\subpageref{NW1NjfzY-h4O6T-1}}}\endmoddef\nwstartdeflinemarkup\nwusesondefline{\\{NW1NjfzY-4CUt2L-4}}\nwenddeflinemarkup
fun disj_imp_disj A B C D =
  let
    val B_imp_D = \nwlinkedidentc{Formula.mk_imp}{NW3w0iRO-3SlqND-3}(B, D);
    val A_imp_C = \nwlinkedidentc{Formula.mk_imp}{NW3w0iRO-3SlqND-3}(A, C);
    val A_or_B = \nwlinkedidentc{Formula.mk_or}{NW3w0iRO-3SlqND-3}(A, B);
    val th1 = or_cong_l A C B;
    val th2 = \nwlinkedidentc{Thm.undisch}{NW4ICsJQ-2zKkPU-2} (\nwlinkedidentc{Thm.undisch}{NW4ICsJQ-2zKkPU-2} th1);
    val th3 = or_cong_r C B D;
    val th4 = \nwlinkedidentc{Thm.undisch}{NW4ICsJQ-2zKkPU-2} th3;
    val th5 = \nwlinkedidentc{Thm.modus_ponens}{NW4ICsJQ-2zKkPU-2} th4 th2;
  in
    \nwlinkedidentc{Thm.disch}{NW4ICsJQ-2zKkPU-2} A_or_B (\nwlinkedidentc{Thm.disch}{NW4ICsJQ-2zKkPU-2} A_imp_C (\nwlinkedidentc{Thm.disch}{NW4ICsJQ-2zKkPU-2} B_imp_D th5))
  end;
\nwindexdefn{\nwixident{Derived.disj{\_}imp{\_}disj}}{Derived.disj:unimp:undisj}{NW1NjfzY-h4O6T-1}\eatline
\nwused{\\{NW1NjfzY-4CUt2L-4}}\nwidentdefs{\\{{\nwixident{Derived.disj{\_}imp{\_}disj}}{Derived.disj:unimp:undisj}}}\nwidentuses{\\{{\nwixident{Formula.mk{\_}imp}}{Formula.mk:unimp}}\\{{\nwixident{Formula.mk{\_}or}}{Formula.mk:unor}}\\{{\nwixident{Thm.disch}}{Thm.disch}}\\{{\nwixident{Thm.modus{\_}ponens}}{Thm.modus:unponens}}\\{{\nwixident{Thm.undisch}}{Thm.undisch}}}\nwindexuse{\nwixident{Formula.mk{\_}imp}}{Formula.mk:unimp}{NW1NjfzY-h4O6T-1}\nwindexuse{\nwixident{Formula.mk{\_}or}}{Formula.mk:unor}{NW1NjfzY-h4O6T-1}\nwindexuse{\nwixident{Thm.disch}}{Thm.disch}{NW1NjfzY-h4O6T-1}\nwindexuse{\nwixident{Thm.modus{\_}ponens}}{Thm.modus:unponens}{NW1NjfzY-h4O6T-1}\nwindexuse{\nwixident{Thm.undisch}}{Thm.undisch}{NW1NjfzY-h4O6T-1}\nwendcode{}\nwbegindocs{46}\nwdocspar
\begin{node}% conj_disj_distribL
$\vdash A\land(B\lor C)\iff((A\land B)\lor(A\land C))$
\end{node}

\begin{proof}
The forward proof:
\begin{enumerate}
\item $B\lor C\vdash B\lor C$ by assuming
\item $\vdash (B\lor C)\implies[(B\implies A\land B)\implies((C\implies A\land C)\implies(A\land B)\lor(A\land C))]$
  by {\Tt{}disj{\_}imp{\_}disj\nwendquote} $B$ $C$ $A\land B$ $A\land C$
\item $\vdash B\lor C\vdash(B\implies A\land B)\implies((C\implies A\land C)\implies(A\land B)\lor(A\land C))$
  by modus ponens(2, 1)%3
\item $A\vdash B\implies A\land B$%4
\item $A, B\lor C\vdash(C\implies A\land C)\implies(A\land B)\lor(A\land C)$
  by modus ponens(3, 4)%5
\item $A\vdash C\implies A\land C$%6
\item $A, B\lor C\vdash(A\land B)\lor(A\land C)$
  by modus ponens(5, 6)%7
\item $\vdash A\implies(B\lor C\implies(A\land B)\lor(A\land C))$
  by discharging the hypotheses%8
\item $\vdash A\land(B\lor C)\implies(A\land B)\lor(A\land C)$
  by uncurrying(8)%
\end{enumerate}
The backward direction:
\begin{enumerate}
\item $A\land B\vdash A\land B$ by assuming $A\land B$%1
\item $A\land B\vdash B$ by $\land$-elim-r(1)%2
\item $A\land B\vdash B\lor C$ by $\lor$-intro-r($C$, 2)%3
\item $A\land B\vdash A$ by $\land$-elim-l(1)%4
\item $A\land B\vdash A\land(B\lor C)$ by $\land$-intro(4, 3)%5
\item $A\land C\vdash A\land C$ by assuming $A\land C$%6
\item $A\land C\vdash C$ by $\land$-elim-r(6)%7
\item $A\land C\vdash B\lor C$ by $\lor$-intro-l($B$, 7)%8
\item $A\land C\vdash A$ by $\land$-elim-l(6)%9
\item $A\land C\vdash A\land(B\lor C)$ by $\land$-intro(9, 8)%10
\item $\vdash(A\land C)\implies A\land(B\lor C)$ by discharging 10
\item $\vdash(A\land B)\implies A\land(B\lor C)$ by discharging 5
\item $\vdash((A\land B)\lor(A\land C))\implies A\land(B\lor C)$ by
  $\lor$-elim(12, 11)
\qedhere
\end{enumerate}
\end{proof}

\nwenddocs{}\nwbegincode{47}\sublabel{NW1NjfzY-efA9w-1}\nwmargintag{{\nwtagstyle{}\subpageref{NW1NjfzY-efA9w-1}}}\moddef{$\vdash A\land(B\lor C)\iff((A\land B)\lor(A\land C))$~{\nwtagstyle{}\subpageref{NW1NjfzY-efA9w-1}}}\endmoddef\nwstartdeflinemarkup\nwusesondefline{\\{NW1NjfzY-4CUt2L-4}}\nwenddeflinemarkup
local
  fun uncurry_imp A B C =
  let
    val AB = \nwlinkedidentc{Formula.mk_and}{NW3w0iRO-3SlqND-3}(A, B);
    val A_imp_B_imp_C = \nwlinkedidentc{Formula.mk_imp}{NW3w0iRO-3SlqND-3}(A, \nwlinkedidentc{Formula.mk_imp}{NW3w0iRO-3SlqND-3}(B, C));
    val th1 = \nwlinkedidentc{Thm.assume}{NW4ICsJQ-2zKkPU-2} A_imp_B_imp_C [AB];
    val th2 = \nwlinkedidentc{Thm.assume}{NW4ICsJQ-2zKkPU-2} AB [A_imp_B_imp_C];
    val th3 = Thm.and_elim_l th2; (* ... |- A *)
    val th4 = \nwlinkedidentc{Thm.modus_ponens}{NW4ICsJQ-2zKkPU-2} th1 th3; (* ... |- B => C *)
    val th5 = Thm.and_elim_r th2; (* ... |- B *)
    val th6 = \nwlinkedidentc{Thm.modus_ponens}{NW4ICsJQ-2zKkPU-2} th4 th5; (* ... |- C *);
    val th7 = \nwlinkedidentc{Thm.disch}{NW4ICsJQ-2zKkPU-2} AB th6;
  in
    \nwlinkedidentc{Thm.disch}{NW4ICsJQ-2zKkPU-2} A_imp_B_imp_C th7
  end;
  fun forward A B C =
    let
      (* taut X produces "A |- X => (A & X)" *)
      fun taut X =
        let
          val lm1 = \nwlinkedidentc{Thm.assume}{NW4ICsJQ-2zKkPU-2} X [A];
          val lm2 = \nwlinkedidentc{Thm.assume}{NW4ICsJQ-2zKkPU-2} A [X];
          val lm3 = Thm.and_intro lm2 lm1;
        in
          \nwlinkedidentc{Thm.disch}{NW4ICsJQ-2zKkPU-2} X lm3
        end;
      val B_or_C = \nwlinkedidentc{Formula.mk_or}{NW3w0iRO-3SlqND-3}(B, C);
      val AB_or_AC = \nwlinkedidentc{Formula.mk_or}{NW3w0iRO-3SlqND-3}(\nwlinkedidentc{Formula.mk_and}{NW3w0iRO-3SlqND-3}(A, B),
                                   \nwlinkedidentc{Formula.mk_and}{NW3w0iRO-3SlqND-3}(A, C));
      val th1 = disj_imp_disj B C (\nwlinkedidentc{Formula.mk_and}{NW3w0iRO-3SlqND-3}(A,B)) (\nwlinkedidentc{Formula.mk_and}{NW3w0iRO-3SlqND-3}(A,C));
      val th2 = \nwlinkedidentc{Thm.assume}{NW4ICsJQ-2zKkPU-2} B_or_C [];
      val th3 = \nwlinkedidentc{Thm.modus_ponens}{NW4ICsJQ-2zKkPU-2} th1 th2;
      val th4 = \nwlinkedidentc{Thm.modus_ponens}{NW4ICsJQ-2zKkPU-2} th3 (taut B);
      val th5 = \nwlinkedidentc{Thm.modus_ponens}{NW4ICsJQ-2zKkPU-2} th4 (taut C);
      val th6 = \nwlinkedidentc{Thm.disch}{NW4ICsJQ-2zKkPU-2} B_or_C th5;
      val th7 = \nwlinkedidentc{Thm.disch}{NW4ICsJQ-2zKkPU-2} A th6; (* |- A => ((B or C) => (A & B) or (A & C)) *)
    in
      \nwlinkedidentc{Thm.modus_ponens}{NW4ICsJQ-2zKkPU-2} (uncurry_imp A B_or_C AB_or_AC) th7
    end;
  fun backward A B C =
    let
      val AB = \nwlinkedidentc{Formula.mk_and}{NW3w0iRO-3SlqND-3}(A, B);
      val AC = \nwlinkedidentc{Formula.mk_and}{NW3w0iRO-3SlqND-3}(A, C);
      val th1 = \nwlinkedidentc{Thm.assume}{NW4ICsJQ-2zKkPU-2} AB [];
      val th2 = Thm.and_elim_r th1;
      val th3 = \nwlinkedidentc{Thm.or_intro_r}{NW4ICsJQ-2zKkPU-4} C th2;
      val th4 = Thm.and_elim_l th1;
      val th5 = \nwlinkedidentc{Thm.disch}{NW4ICsJQ-2zKkPU-2} AB (Thm.and_intro th4 th3);
      val th6 = \nwlinkedidentc{Thm.assume}{NW4ICsJQ-2zKkPU-2} AC [];
      val th7 = Thm.and_elim_r th6;
      val th8 = \nwlinkedidentc{Thm.or_intro_l}{NW4ICsJQ-2zKkPU-4} B th7;
      val th9 = Thm.and_elim_l th6;
      val th10 = \nwlinkedidentc{Thm.disch}{NW4ICsJQ-2zKkPU-2} AC (Thm.and_intro th9 th8);
    in
      \nwlinkedidentc{Thm.or_elim}{NW4ICsJQ-2zKkPU-4} th5 th10
    end;
in
fun conj_disj_distribL A B C =
  \nwlinkedidentc{Thm.iff_intro}{NW4ICsJQ-2zKkPU-7} (forward A B C) (backward A B C)
end;
\nwindexdefn{\nwixident{Derived.conj{\_}disj{\_}distribL}}{Derived.conj:undisj:undistribL}{NW1NjfzY-efA9w-1}\eatline
\nwused{\\{NW1NjfzY-4CUt2L-4}}\nwidentdefs{\\{{\nwixident{Derived.conj{\_}disj{\_}distribL}}{Derived.conj:undisj:undistribL}}}\nwidentuses{\\{{\nwixident{Formula.mk{\_}and}}{Formula.mk:unand}}\\{{\nwixident{Formula.mk{\_}imp}}{Formula.mk:unimp}}\\{{\nwixident{Formula.mk{\_}or}}{Formula.mk:unor}}\\{{\nwixident{Thm.assume}}{Thm.assume}}\\{{\nwixident{Thm.disch}}{Thm.disch}}\\{{\nwixident{Thm.iff{\_}intro}}{Thm.iff:unintro}}\\{{\nwixident{Thm.modus{\_}ponens}}{Thm.modus:unponens}}\\{{\nwixident{Thm.or{\_}elim}}{Thm.or:unelim}}\\{{\nwixident{Thm.or{\_}intro{\_}l}}{Thm.or:unintro:unl}}\\{{\nwixident{Thm.or{\_}intro{\_}r}}{Thm.or:unintro:unr}}}\nwindexuse{\nwixident{Formula.mk{\_}and}}{Formula.mk:unand}{NW1NjfzY-efA9w-1}\nwindexuse{\nwixident{Formula.mk{\_}imp}}{Formula.mk:unimp}{NW1NjfzY-efA9w-1}\nwindexuse{\nwixident{Formula.mk{\_}or}}{Formula.mk:unor}{NW1NjfzY-efA9w-1}\nwindexuse{\nwixident{Thm.assume}}{Thm.assume}{NW1NjfzY-efA9w-1}\nwindexuse{\nwixident{Thm.disch}}{Thm.disch}{NW1NjfzY-efA9w-1}\nwindexuse{\nwixident{Thm.iff{\_}intro}}{Thm.iff:unintro}{NW1NjfzY-efA9w-1}\nwindexuse{\nwixident{Thm.modus{\_}ponens}}{Thm.modus:unponens}{NW1NjfzY-efA9w-1}\nwindexuse{\nwixident{Thm.or{\_}elim}}{Thm.or:unelim}{NW1NjfzY-efA9w-1}\nwindexuse{\nwixident{Thm.or{\_}intro{\_}l}}{Thm.or:unintro:unl}{NW1NjfzY-efA9w-1}\nwindexuse{\nwixident{Thm.or{\_}intro{\_}r}}{Thm.or:unintro:unr}{NW1NjfzY-efA9w-1}\nwendcode{}\nwbegindocs{48}\nwdocspar
\begin{node}%conj_disj_distribR
$\vdash (A\lor B)\land C\iff((A\land C)\lor(B\land C))$
\end{node}

\begin{proof}
The forward direction:
\begin{enumerate}
\item $\vdash C\land(A\lor B)\iff(C\land A)\lor(C\land B)$ by
  {\Tt{}conj{\_}disj{\_}distribL\nwendquote}($C$, $A$, $B$)%1
\item $\vdash C\land(A\lor B)\implies(C\land A)\lor(C\land B)$ by $\Leftrightarrow$-elim-l(1)%2
\item $\vdash (A\lor B)\land C\implies C\land(A\lor B)$ by
  {\Tt{}and{\_}comm\nwendquote}($A\lor B$, $C$)%3
\item $\vdash (A\lor B)\land C\implies(C\land A)\lor(C\land B)$ by
  hypothetical syllogism(3, 2)%4
\item $(A\lor B)\land C\vdash(C\land A)\lor(C\land B)$ by undischarging(4)%5
\item $\vdash(C\land A)\implies(A\land C)$ by {\Tt{}and{\_}comm\nwendquote}($C$, $A$)%6
\item $\vdash((C\land A)\implies(A\land C))\implies(((C\land A)\lor(C\land B))\implies((A\land C)\lor(C\land B)))$
  by {\Tt{}or{\_}cong{\_}l\nwendquote}($C\land A$, $A\land C$, $C\land B$)%7
\item $\vdash((C\land A)\lor(C\land B))\implies((A\land C)\lor(C\land B))$
  by modus ponens(7, 6)% 8
\item $(A\lor B)\land C\vdash(A\land C)\lor(C\land B)$ by modus
  ponens(8, 5)%9
\item $\vdash(C\land B)\implies(B\land C)$ by {\Tt{}and{\_}comm\nwendquote}($C$, $B$)%10
\item $\vdash((C\land B)\implies(B\land C))\implies(((A\land
  C)\lor(C\land B))\implies((A\land C)\lor(B\land C)))$ by
  {\Tt{}or{\_}cong{\_}r\nwendquote}($A\land C$, $C\land B$, $B\land C$)%11
\item $\vdash((A\land C)\lor(C\land B))\implies((A\land C)\lor(B\land C))$ by modus
  ponens(11, 10)%12
\item $(A\lor B)\land C\vdash(A\land C)\lor(B\land C)$ by modus
  ponens(9, 12)%13
\item $\vdash((A\lor B)\land C)\implies(A\land C)\lor(B\land C)$ by
  discharging 13%14
\end{enumerate}
The backward direction:
\begin{enumerate}\setcounter{enumi}{14}
\item $A\land C\vdash A\land C$ by assuming%15
\item $A\land C\vdash A$ by $\land$-elim-l(15)%16
\item $A\land C\vdash A\lor B$ by $\lor$-intro-r(16, $B$)%17
\item $A\land C\vdash C$ by $\land$-elim-r(15)%18
\item $(A\land C)\vdash(A\lor B)\land C$ by $\land$-intro(17,18)%19
\item $\vdash(A\land C)\implies(A\lor B)\land C$ by discharging(19)%20
\item $B\land C\vdash B\land C$ by assuming%21
\item $B\land C\vdash B$ by $\land$-elim-l(21)%22
\item $B\land C\vdash A\lor B$ by $\lor$-intro-l($A$, 22)%23
\item $B\land C\vdash C$ by $\land$-elim-r(21)%24
\item $B\land C\vdash(A\lor B)\land C$ by $\land$-intro(23, 24)%25
\item $\vdash(B\land C)\implies((A\lor B)\land C)$ by discarging(25)%26
\item $\vdash(A\land C)\lor(B\land C)\implies((A\lor B)\land C)$ by $\lor$-intro(20, 26)\qedhere
\end{enumerate}
\end{proof}

\nwenddocs{}\nwbegincode{49}\sublabel{NW1NjfzY-8du9Q-1}\nwmargintag{{\nwtagstyle{}\subpageref{NW1NjfzY-8du9Q-1}}}\moddef{Right distributivity of conjunction over disjunction~{\nwtagstyle{}\subpageref{NW1NjfzY-8du9Q-1}}}\endmoddef\nwstartdeflinemarkup\nwusesondefline{\\{NW1NjfzY-4CUt2L-4}}\nwenddeflinemarkup
local
  infixr &;
  fun (A & B) = \nwlinkedidentc{Formula.mk_and}{NW3w0iRO-3SlqND-3}(A, B);
  fun forward A B C =
    let
      val A_or_B = \nwlinkedidentc{Formula.mk_or}{NW3w0iRO-3SlqND-3}(A, B);
      val th1 = conj_disj_distribL C A B;
      val th2 = \nwlinkedidentc{Thm.iff_elim_l}{NW4ICsJQ-2zKkPU-7} th1;
      val th3 = and_comm A_or_B C;
      val th4 = hypo_syll th3 th2;
      val th5 = \nwlinkedidentc{Thm.undisch}{NW4ICsJQ-2zKkPU-2} th4;
      val th6 = and_comm C A;
      val th7 = or_cong_l (C & A) (A & C) (C & B);
      val th8 = \nwlinkedidentc{Thm.modus_ponens}{NW4ICsJQ-2zKkPU-2} th7 th6;
      val th9 = \nwlinkedidentc{Thm.modus_ponens}{NW4ICsJQ-2zKkPU-2} th8 th5;
      val th10 = and_comm C B;
      val th11 = or_cong_r (A & C) (C & B) (B & C);
      val th12 = \nwlinkedidentc{Thm.modus_ponens}{NW4ICsJQ-2zKkPU-2} th11 th10;
      val th13 = \nwlinkedidentc{Thm.modus_ponens}{NW4ICsJQ-2zKkPU-2} th9 th12;
    in
      \nwlinkedidentc{Thm.disch}{NW4ICsJQ-2zKkPU-2} (\nwlinkedidentc{Formula.mk_or}{NW3w0iRO-3SlqND-3}(A,B) & C) th13
    end;
  fun backward A B C =
    let
      fun iter X Y =
        let
          val lm1 = \nwlinkedidentc{Thm.assume}{NW4ICsJQ-2zKkPU-2} (X & C) [];
          val lm2 = Thm.and_elim_l lm1;
          val lm3 = \nwlinkedidentc{Thm.or_intro_r}{NW4ICsJQ-2zKkPU-4} Y lm2;
          val lm4 = Thm.and_elim_r lm1;
          val lm5 = Thm.and_intro lm3 lm4;
        in
          \nwlinkedidentc{Thm.disch}{NW4ICsJQ-2zKkPU-2} (X & C) lm5
        end;
      val th20 = iter A B;
      val th26 = iter B A;
    in
      Thm.and_intro th20 th26
    end;
in
fun conj_disj_distribR A B C =
  \nwlinkedidentc{Thm.iff_intro}{NW4ICsJQ-2zKkPU-7} (forward A B C) (backward A B C)
end;
\nwindexdefn{\nwixident{Derived.conj{\_}disj{\_}distribR}}{Derived.conj:undisj:undistribR}{NW1NjfzY-8du9Q-1}\eatline
\nwused{\\{NW1NjfzY-4CUt2L-4}}\nwidentdefs{\\{{\nwixident{Derived.conj{\_}disj{\_}distribR}}{Derived.conj:undisj:undistribR}}}\nwidentuses{\\{{\nwixident{Formula.mk{\_}and}}{Formula.mk:unand}}\\{{\nwixident{Formula.mk{\_}or}}{Formula.mk:unor}}\\{{\nwixident{Thm.assume}}{Thm.assume}}\\{{\nwixident{Thm.disch}}{Thm.disch}}\\{{\nwixident{Thm.iff{\_}elim{\_}l}}{Thm.iff:unelim:unl}}\\{{\nwixident{Thm.iff{\_}intro}}{Thm.iff:unintro}}\\{{\nwixident{Thm.modus{\_}ponens}}{Thm.modus:unponens}}\\{{\nwixident{Thm.or{\_}intro{\_}r}}{Thm.or:unintro:unr}}\\{{\nwixident{Thm.undisch}}{Thm.undisch}}}\nwindexuse{\nwixident{Formula.mk{\_}and}}{Formula.mk:unand}{NW1NjfzY-8du9Q-1}\nwindexuse{\nwixident{Formula.mk{\_}or}}{Formula.mk:unor}{NW1NjfzY-8du9Q-1}\nwindexuse{\nwixident{Thm.assume}}{Thm.assume}{NW1NjfzY-8du9Q-1}\nwindexuse{\nwixident{Thm.disch}}{Thm.disch}{NW1NjfzY-8du9Q-1}\nwindexuse{\nwixident{Thm.iff{\_}elim{\_}l}}{Thm.iff:unelim:unl}{NW1NjfzY-8du9Q-1}\nwindexuse{\nwixident{Thm.iff{\_}intro}}{Thm.iff:unintro}{NW1NjfzY-8du9Q-1}\nwindexuse{\nwixident{Thm.modus{\_}ponens}}{Thm.modus:unponens}{NW1NjfzY-8du9Q-1}\nwindexuse{\nwixident{Thm.or{\_}intro{\_}r}}{Thm.or:unintro:unr}{NW1NjfzY-8du9Q-1}\nwindexuse{\nwixident{Thm.undisch}}{Thm.undisch}{NW1NjfzY-8du9Q-1}\nwendcode{}\nwbegindocs{50}\nwdocspar
\begin{node}%disj_conj_distribL
$\vdash A\lor(B\land C)\iff((A\lor B)\land(A\lor C))$
\end{node}

\begin{proof}
The forward direction:
\begin{enumerate}
\item $A\vdash A$ by assuming%1
\item $A\vdash A\lor B$ by $\lor$-intro(1, $B$)%2
\item $A\vdash A\lor C$ by $\lor$-intro(1, $C$)%3
\item $A\vdash(A\lor B)\land(A\lor C)$ by $\land$-intro(2, 3)%4
\item $\vdash A\implies(A\lor B)\land(A\lor C)$ by discharging(4)%5
\item $B\land C\vdash B\land C$ by assuming%6
\item $B\land C\vdash B$ by $\land$-elim-l(6)%7
\item $B\land C\vdash A\lor B$ by $\lor$-intro-l(7, $A$)%8
\item $B\land C\vdash C$ by $\land$-elim-r(6)%9
\item $B\land C\vdash A\lor C$ by $\lor$-intr-l(9, $A$)%10
\item $B\land C\vdash (A\lor B)\land(A\lor C)$ by $\land$-intro(8, 10)%11
\item $\vdash B\land C\implies(A\lor B)\land(A\lor C)$ by discharging(11)%12
\item $\vdash A\lor(B\land C)\implies(A\lor B)\land(A\lor C)$ by $\lor$-elim(5, 12)%13
\end{enumerate}
The other direction ($\vdash(A\lor B)\land(A\lor C)\implies A\lor(B\land C)$):
\begin{enumerate}
\item $A\vdash A$ by assuming%1
\item $A\vdash A\lor(B\land C)$ by $\lor$-intro-r($B\land C$, 1)%2
\item $\vdash A\implies A\lor(B\land C)$ by discharging $A$ in (2)%3
\item $B,C\vdash B$ by assuming%4
\item $B,C\vdash C$ by assuming%5
\item $B,C\vdash B\land C$ by $\land$-intro(4, 5)%6
\item $B,C\vdash A\lor(B\land C)$ by $\lor$-intro-l($A$, 6)%7
\item $B\vdash C\implies A\lor(B\land C)$ by discharging $C$ from 7%8
\item $B\vdash A\lor C\implies A\lor(B\land C)$ by $\lor$-elim(3, 8)%9
\item $\vdash B\implies(A\lor C\implies A\lor(B\land C))$ by
  discharging $B$ from 9%10
\item $A,A\lor C\vdash A\lor(B\land C)$ by weakening 3 with $A\lor C$%11
\item $A\vdash A\lor C\implies A\lor(B\land C)$ by discharging $A\lor C$ from 11%12
\item $\vdash A\implies(A\lor C\implies A\lor(B\land C)$ by
  discharging $A$ from 12%13
\item $\vdash (A\lor B)\implies((A\lor C)\implies A\lor(B\land C))$ by
  $\lor$-elim(13, 10)%14
\item $\vdash ((A\lor B)\implies((A\lor C)\implies A\lor(B\land C)))\implies((A\lor B)\land(A\lor C)\implies A\lor(B\land C))$
  by uncurrying $A\lor B$, $A\lor C$, $A\lor(B\land C)$%15
\item $\vdash(A\lor B)\land(A\lor C)\implies A\lor(B\land C)$ by modus
  ponens(15, 11)\qedhere
\end{enumerate}
\end{proof}

\nwenddocs{}\nwbegincode{51}\sublabel{NW1NjfzY-2uP77E-1}\nwmargintag{{\nwtagstyle{}\subpageref{NW1NjfzY-2uP77E-1}}}\moddef{Left distributivity of disjunction over conjunction~{\nwtagstyle{}\subpageref{NW1NjfzY-2uP77E-1}}}\endmoddef\nwstartdeflinemarkup\nwenddeflinemarkup
local
  infixr & or;
  fun (A & B) = \nwlinkedidentc{Formula.mk_and}{NW3w0iRO-3SlqND-3}(A, B);
  fun (A or B) = \nwlinkedidentc{Formula.mk_or}{NW3w0iRO-3SlqND-3}(A, B);
  fun uncurry A B C =
    let
      val AB = A & B;
      val A_imp_B_imp_C = \nwlinkedidentc{Formula.mk_imp}{NW3w0iRO-3SlqND-3}(A, \nwlinkedidentc{Formula.mk_imp}{NW3w0iRO-3SlqND-3}(B, C));
      val th1 = \nwlinkedidentc{Thm.assume}{NW4ICsJQ-2zKkPU-2} A_imp_B_imp_C [AB];
      val th2 = \nwlinkedidentc{Thm.assume}{NW4ICsJQ-2zKkPU-2} AB [A_imp_B_imp_C];
      val th3 = Thm.and_elim_l th2; (* ... |- A *)
      val th4 = \nwlinkedidentc{Thm.modus_ponens}{NW4ICsJQ-2zKkPU-2} th1 th3; (* ... |- B => C *)
      val th5 = Thm.and_elim_r th2; (* ... |- B *)
      val th6 = \nwlinkedidentc{Thm.modus_ponens}{NW4ICsJQ-2zKkPU-2} th4 th5; (* ... |- C *);
      val th7 = \nwlinkedidentc{Thm.disch}{NW4ICsJQ-2zKkPU-2} AB th6;
    in
      \nwlinkedidentc{Thm.disch}{NW4ICsJQ-2zKkPU-2} A_imp_B_imp_C th7
    end;
  fun forward A B C =
    let
      val th1 = \nwlinkedidentc{Thm.assume}{NW4ICsJQ-2zKkPU-2} A [];
      val th2 = \nwlinkedidentc{Thm.or_intro_r}{NW4ICsJQ-2zKkPU-4} B th1;
      val th3 = \nwlinkedidentc{Thm.or_intro_r}{NW4ICsJQ-2zKkPU-4} C th1;
      val th4 = Thm.and_intro th2 th3;
      val th5 = \nwlinkedidentc{Thm.disch}{NW4ICsJQ-2zKkPU-2} A th4;
      val th6 = \nwlinkedidentc{Thm.assume}{NW4ICsJQ-2zKkPU-2} (B & C) [];
      val th7 = Thm.and_elim_l th6;
      val th8 = \nwlinkedidentc{Thm.or_intro_l}{NW4ICsJQ-2zKkPU-4} A th7;
      val th9 = Thm.and_elim_r th6;
      val th10 = \nwlinkedidentc{Thm.or_intro_l}{NW4ICsJQ-2zKkPU-4} A th9;
      val th11 = Thm.and_intro th8 th10;
      val th12 = \nwlinkedidentc{Thm.disch}{NW4ICsJQ-2zKkPU-2} (B & C) th11;
    in
      \nwlinkedidentc{Thm.or_elim}{NW4ICsJQ-2zKkPU-4} th5 th12
    end;
  fun backward A B C =
    let
      val th1 = Thm.assuming A [];
      val th2 = \nwlinkedidentc{Thm.or_intro_r}{NW4ICsJQ-2zKkPU-4} (B & C) th1;
      val th3 = \nwlinkedidentc{Thm.disch}{NW4ICsJQ-2zKkPU-2} A th2;
      val th4 = \nwlinkedidentc{Thm.assume}{NW4ICsJQ-2zKkPU-2} B [C];
      val th5 = \nwlinkedidentc{Thm.assume}{NW4ICsJQ-2zKkPU-2} C [B];
      val th6 = Thm.and_intro th4 th5;
      val th7 = \nwlinkedidentc{Thm.or_intro_l}{NW4ICsJQ-2zKkPU-4} A th6;
      val th8 = \nwlinkedidentc{Thm.disch}{NW4ICsJQ-2zKkPU-2} C th7;
      val th9 = \nwlinkedidentc{Thm.or_elim}{NW4ICsJQ-2zKkPU-4} th3 th8;
      val th10 = \nwlinkedidentc{Thm.disch}{NW4ICsJQ-2zKkPU-2} B th9;
      val th11 = weaken [(A or C)] th3;
      val th12 = \nwlinkedidentc{Thm.disch}{NW4ICsJQ-2zKkPU-2} (A or C) th11;
      val th13 = \nwlinkedidentc{Thm.disch}{NW4ICsJQ-2zKkPU-2} A th11;
      val th14 = \nwlinkedidentc{Thm.or_elim}{NW4ICsJQ-2zKkPU-4} th13 th10;
      val th15 = uncurry (A or B) (A or C) (A & (B or C));
    in
      \nwlinkedidentc{Thm.modus_ponens}{NW4ICsJQ-2zKkPU-2} th15 th14
    end;
in
fun disj_conj_distribL A B C =
  \nwlinkedidentc{Thm.iff_intro}{NW4ICsJQ-2zKkPU-7} (forward A B C) (backward A B C)
end;
\nwnotused{Left distributivity of disjunction over conjunction}\nwidentuses{\\{{\nwixident{Formula.mk{\_}and}}{Formula.mk:unand}}\\{{\nwixident{Formula.mk{\_}imp}}{Formula.mk:unimp}}\\{{\nwixident{Formula.mk{\_}or}}{Formula.mk:unor}}\\{{\nwixident{Thm.assume}}{Thm.assume}}\\{{\nwixident{Thm.disch}}{Thm.disch}}\\{{\nwixident{Thm.iff{\_}intro}}{Thm.iff:unintro}}\\{{\nwixident{Thm.modus{\_}ponens}}{Thm.modus:unponens}}\\{{\nwixident{Thm.or{\_}elim}}{Thm.or:unelim}}\\{{\nwixident{Thm.or{\_}intro{\_}l}}{Thm.or:unintro:unl}}\\{{\nwixident{Thm.or{\_}intro{\_}r}}{Thm.or:unintro:unr}}}\nwindexuse{\nwixident{Formula.mk{\_}and}}{Formula.mk:unand}{NW1NjfzY-2uP77E-1}\nwindexuse{\nwixident{Formula.mk{\_}imp}}{Formula.mk:unimp}{NW1NjfzY-2uP77E-1}\nwindexuse{\nwixident{Formula.mk{\_}or}}{Formula.mk:unor}{NW1NjfzY-2uP77E-1}\nwindexuse{\nwixident{Thm.assume}}{Thm.assume}{NW1NjfzY-2uP77E-1}\nwindexuse{\nwixident{Thm.disch}}{Thm.disch}{NW1NjfzY-2uP77E-1}\nwindexuse{\nwixident{Thm.iff{\_}intro}}{Thm.iff:unintro}{NW1NjfzY-2uP77E-1}\nwindexuse{\nwixident{Thm.modus{\_}ponens}}{Thm.modus:unponens}{NW1NjfzY-2uP77E-1}\nwindexuse{\nwixident{Thm.or{\_}elim}}{Thm.or:unelim}{NW1NjfzY-2uP77E-1}\nwindexuse{\nwixident{Thm.or{\_}intro{\_}l}}{Thm.or:unintro:unl}{NW1NjfzY-2uP77E-1}\nwindexuse{\nwixident{Thm.or{\_}intro{\_}r}}{Thm.or:unintro:unr}{NW1NjfzY-2uP77E-1}\nwendcode{}\nwbegindocs{52}\nwdocspar

\begin{node}%disj_conj_distribR
$\vdash (A\land B)\lor C\iff((A\lor C)\land(B\lor C))$
\end{node}

There are two ways to prove this: use {\Tt{}disj{\_}conj{\_}distribL\nwendquote} and
congruences and commutativity to get the result (requiring 23 steps); or in a way analogous
to the proof of {\Tt{}disj{\_}conj{\_}distribL\nwendquote} (requiring 27 steps). 

\begin{proof}
\begin{enumerate}
\item $\vdash (A\land B)\lor C\iff C\lor(A\land B)$ by $\lor$-commutativity
\item $\vdash C\lor(A\land B)\iff (C\lor A)\land(C\lor B)$ by {\Tt{}disj{\_}conj{\_}distribL\nwendquote} $C$ $A$ $B$
\item $\vdash[(A\land B)\lor C\iff C\lor(A\land B)]\land[C\lor(A\land B)\iff (C\lor A)\land(C\lor B)]$ by $\land$-intro(1, 2)%3
\item $\vdash(A\land B)\lor C\iff (C\lor A)\land(C\lor B)$ by modus ponens ({\Tt{}iff{\_}trans\nwendquote}
$((A\land B)\lor C)$ $(C\lor(A\land B))$ $[C\lor(A\land B)\iff (C\lor A)\land(C\lor B)]$),
3%4
\end{enumerate}
\bigbreak
Then we have the forward proof:
\begin{enumerate}\setcounter{enumi}{4}
\item $\vdash(A\land B)\lor C\implies (C\lor A)\land(C\lor B)$ by $\Leftrightarrow$-elim-r(4)%5
\item $\vdash(C\lor A)\implies(A\lor C)$ by $\lor$-commutativity%6
\item $\vdash((C\lor A)\implies(A\lor C))\implies((C\lor A)\land(C\lor B)\implies(A\lor C)\land(C\lor B))$
  by $\land$-congruence%7
\item $\vdash(C\lor A)\land(C\lor B)\implies(A\lor C)\land(C\lor B)$
  by modus ponens(7, 6)%8
\item $\vdash(A\land B)\lor C\implies (A\lor C)\land(C\lor B)$
  by hypothetical syllogism(5, 8)%9
\item $\vdash(C\lor B)\implies(B\lor C)$ by $\lor$-commutativity%10
\item $\vdash((C\lor B)\implies(B\lor C))\implies((A\lor
  C))\land(C\lor B)\implies(A\lor C)\land(B\lor C))$
  by $\lor$-congruence-r%11
\item $\vdash (A\lor C))\land(C\lor B)\implies(A\lor C)\land(B\lor C)$
  by modus ponens(11, 10)%12
\item $\vdash(A\land B)\lor C\implies (A\lor C)\land(B\lor C)$
  by hypothetical syllogism(9, 12)%13
\end{enumerate}
\bigbreak
The backward proof:
\begin{enumerate}\setcounter{enumi}{13}
\item $\vdash(C\lor A)\land(C\lor B)\implies(A\land B)\lor C$ by $\Leftrightarrow$-elim-r(4)%14
\item $\vdash(A\lor C)\implies(C\lor A)$ by $\lor$-commutativity%15
\item $\vdash((A\lor C)\implies(C\lor A))\implies((A\lor C)\land(C\lor B)\implies(C\lor A)\land(C\lor B))$
  by $\land$-congruence-r%16
\item $\vdash(A\lor C)\land(C\lor B)\implies(C\lor A)\land(C\lor B)$
  by modus ponens(16, 15)%17
\item $\vdash(A\lor C)\land(C\lor B)\implies(A\land B)\lor C$ by
  hypothetical syllogism(17, 14)%18
\item $\vdash(B\lor C)\implies(C\lor B)$ by $\lor$-commutativity%19
\item $\vdash((B\lor C)\implies(C\lor B))\implies((A\lor C)\land(B\lor
  C)\implies(A\lor C)\land(C\lor B))$
  by $\land$-congruence-r%20
\item $\vdash (A\lor C)\land(B\lor C)\implies(A\lor C)\land(C\lor B)$
  by modus ponens(20, 19)%21
\item $\vdash(A\lor C)\land(B\lor C)\implies(A\land B)\lor C$ by
  hypothetical syllogism(21, 18)%22
\item $\vdash(A\land B)\lor C\iff(A\lor C)\land(B\lor C)$ by
  $\Leftrightarrow$-intro(13, 22)\qedhere
\end{enumerate}
\end{proof}

\begin{node}
$\vdash A\implies A$.
\end{node}

\nwenddocs{}\nwbegincode{53}\sublabel{NW1NjfzY-2PXiqP-1}\nwmargintag{{\nwtagstyle{}\subpageref{NW1NjfzY-2PXiqP-1}}}\moddef{Theorem: implications is reflexive~{\nwtagstyle{}\subpageref{NW1NjfzY-2PXiqP-1}}}\endmoddef\nwstartdeflinemarkup\nwusesondefline{\\{NW1NjfzY-4CUt2L-5}}\nwenddeflinemarkup
fun imp_reflexive A =
  \nwlinkedidentc{Thm.disch}{NW4ICsJQ-2zKkPU-2} A (\nwlinkedidentc{Thm.assume}{NW4ICsJQ-2zKkPU-2} A []);
\nwused{\\{NW1NjfzY-4CUt2L-5}}\nwidentuses{\\{{\nwixident{Thm.assume}}{Thm.assume}}\\{{\nwixident{Thm.disch}}{Thm.disch}}}\nwindexuse{\nwixident{Thm.assume}}{Thm.assume}{NW1NjfzY-2PXiqP-1}\nwindexuse{\nwixident{Thm.disch}}{Thm.disch}{NW1NjfzY-2PXiqP-1}\nwendcode{}\nwbegindocs{54}\nwdocspar

\begin{theorem}[Frege]
$\vdash(A\implies(B\implies C))\implies((A\implies B)\implies(A\implies C))$
\end{theorem}

\begin{proof}
\begin{enumerate}
\item $(A\implies(B\implies C)), (A\implies B), A\vdash A$
\item $(A\implies(B\implies C)), (A\implies B), A\vdash A\implies B$
\item $(A\implies(B\implies C)), (A\implies B), A\vdash B$ by modus ponens(2,1)
\item $(A\implies(B\implies C)), (A\implies B), A\vdash A\implies(B\implies C)$
\item $(A\implies(B\implies C)), (A\implies B), A\vdash B\implies C$
  by modus ponens(4, 1)
\item $(A\implies(B\implies C)), (A\implies B), A\vdash C$
  by modus ponens(5, 3)
\item $(A\implies(B\implies C)), (A\implies B)\vdash A\implies C$
  by discharging $A$
\item $(A\implies(B\implies C))\vdash (A\implies B)\implies(A\implies C)$
  by discharging $A\implies B$
\item $\vdash(A\implies(B\implies C))\implies\bigl((A\implies B)\implies(A\implies C)\bigr)$
  by discharging $A\implies(B\implies C)$\qedhere
\end{enumerate}
\end{proof}

\nwenddocs{}\nwbegincode{55}\sublabel{NW1NjfzY-bcGD8-1}\nwmargintag{{\nwtagstyle{}\subpageref{NW1NjfzY-bcGD8-1}}}\moddef{Frege's tautology (propositional logic)~{\nwtagstyle{}\subpageref{NW1NjfzY-bcGD8-1}}}\endmoddef\nwstartdeflinemarkup\nwusesondefline{\\{NW1NjfzY-4CUt2L-5}}\nwenddeflinemarkup
fun frege_id A B C =
  let
    val A_imp_B_imp_C = \nwlinkedidentc{Formula.mk_imp}{NW3w0iRO-3SlqND-3}(A, \nwlinkedidentc{Formula.mk_imp}{NW3w0iRO-3SlqND-3}(B, C));
    val A_imp_B = \nwlinkedidentc{Formula.mk_imp}{NW3w0iRO-3SlqND-3}(A, B);

    val th1 = \nwlinkedidentc{Thm.assume}{NW4ICsJQ-2zKkPU-2} A [A_imp_B_imp_C, A_imp_B]; (* ... |- A *)
    val th2 = \nwlinkedidentc{Thm.assume}{NW4ICsJQ-2zKkPU-2} A_imp_B [A_imp_B_imp_C, A];
    val th3 = \nwlinkedidentc{Thm.modus_ponens}{NW4ICsJQ-2zKkPU-2} th2 th1; (* ... |- B *)
    
    val th4 = \nwlinkedidentc{Thm.assume}{NW4ICsJQ-2zKkPU-2} A_imp_B_imp_C [A_imp_B, A];
    val th5 = \nwlinkedidentc{Thm.modus_ponens}{NW4ICsJQ-2zKkPU-2} th4 th1; (* ... |- B => C *)
    val th6 = \nwlinkedidentc{Thm.modus_ponens}{NW4ICsJQ-2zKkPU-2} th5 th2; (* ... |- C *)
    val th7 = \nwlinkedidentc{Thm.disch}{NW4ICsJQ-2zKkPU-2} A th6; (* ... |- A => C *)
    val th8 = \nwlinkedidentc{Thm.disch}{NW4ICsJQ-2zKkPU-2} A_imp_B th7;
  in
    \nwlinkedidentc{Thm.disch}{NW4ICsJQ-2zKkPU-2} A_imp_B_imp_C th8
  end;
\nwindexdefn{\nwixident{Derived.frege{\_}id}}{Derived.frege:unid}{NW1NjfzY-bcGD8-1}\eatline
\nwused{\\{NW1NjfzY-4CUt2L-5}}\nwidentdefs{\\{{\nwixident{Derived.frege{\_}id}}{Derived.frege:unid}}}\nwidentuses{\\{{\nwixident{Formula.mk{\_}imp}}{Formula.mk:unimp}}\\{{\nwixident{Thm.assume}}{Thm.assume}}\\{{\nwixident{Thm.disch}}{Thm.disch}}\\{{\nwixident{Thm.modus{\_}ponens}}{Thm.modus:unponens}}}\nwindexuse{\nwixident{Formula.mk{\_}imp}}{Formula.mk:unimp}{NW1NjfzY-bcGD8-1}\nwindexuse{\nwixident{Thm.assume}}{Thm.assume}{NW1NjfzY-bcGD8-1}\nwindexuse{\nwixident{Thm.disch}}{Thm.disch}{NW1NjfzY-bcGD8-1}\nwindexuse{\nwixident{Thm.modus{\_}ponens}}{Thm.modus:unponens}{NW1NjfzY-bcGD8-1}\nwendcode{}\nwbegindocs{56}\nwdocspar
\begin{node}
We can ``curry'' implications to conjunction, and vice-versa. That is
to say, we can replace $A\land B\implies C$ with $A\implies(B\implies C)$
and vice-versa.
\end{node}

\begin{proof}
$\vdash(A\land B\implies C)\implies(A\implies(B\implies C))$
\begin{enumerate}
\item $A,B,A\land B\implies C\vdash A$ by assuming
\item $A,B,A\land B\implies C\vdash B$ by assuming
\item $A,B,A\land B\implies C\vdash A\land B$ by $\land$-intro(1, 2)
\item $A,B,A\land B\implies C\vdash A\land B\implies C$ by assuming
\item $A,B,A\land B\implies C\vdash C$ by modus ponens(4, 3)
\item $A,A\land B\implies C\vdash B\implies C$ by discharging $B$
\item $A\land B\implies C\vdash A\implies(B\implies C)$ by discharging $A$
\item $\vdash(A\land B\implies C)\implies(A\implies(B\implies C))$ by discharging \qedhere
\end{enumerate}
\end{proof}

\nwenddocs{}\nwbegincode{57}\sublabel{NW1NjfzY-1iptma-1}\nwmargintag{{\nwtagstyle{}\subpageref{NW1NjfzY-1iptma-1}}}\moddef{Curry implication~{\nwtagstyle{}\subpageref{NW1NjfzY-1iptma-1}}}\endmoddef\nwstartdeflinemarkup\nwusesondefline{\\{NW1NjfzY-4CUt2L-5}}\nwenddeflinemarkup
fun curry_imp A B C =
  let
    val AB = \nwlinkedidentc{Formula.mk_and}{NW3w0iRO-3SlqND-3}(A, B);
    val AB_imp_C = \nwlinkedidentc{Formula.mk_imp}{NW3w0iRO-3SlqND-3}(AB, C);
    val th1 = \nwlinkedidentc{Thm.assume}{NW4ICsJQ-2zKkPU-2} A [B, AB_imp_C];
    val th2 = \nwlinkedidentc{Thm.assume}{NW4ICsJQ-2zKkPU-2} B [A, AB_imp_C];
    val th3 = Thm.and_intro th1 th2; (* ... |- A & B *)
    val th4 = \nwlinkedidentc{Thm.assume}{NW4ICsJQ-2zKkPU-2} AB_imp_C [A, B];
    val th5 = \nwlinkedidentc{Thm.modus_ponens}{NW4ICsJQ-2zKkPU-2} th4 th3; (* |- C *)
    val th6 = \nwlinkedidentc{Thm.disch}{NW4ICsJQ-2zKkPU-2} B th5;
    val th7 = \nwlinkedidentc{Thm.disch}{NW4ICsJQ-2zKkPU-2} A th6;
  in
    \nwlinkedidentc{Thm.disch}{NW4ICsJQ-2zKkPU-2} AB_imp_C th7
  end;
\nwindexdefn{\nwixident{Derived.curry{\_}imp}}{Derived.curry:unimp}{NW1NjfzY-1iptma-1}\eatline
\nwused{\\{NW1NjfzY-4CUt2L-5}}\nwidentdefs{\\{{\nwixident{Derived.curry{\_}imp}}{Derived.curry:unimp}}}\nwidentuses{\\{{\nwixident{Formula.mk{\_}and}}{Formula.mk:unand}}\\{{\nwixident{Formula.mk{\_}imp}}{Formula.mk:unimp}}\\{{\nwixident{Thm.assume}}{Thm.assume}}\\{{\nwixident{Thm.disch}}{Thm.disch}}\\{{\nwixident{Thm.modus{\_}ponens}}{Thm.modus:unponens}}}\nwindexuse{\nwixident{Formula.mk{\_}and}}{Formula.mk:unand}{NW1NjfzY-1iptma-1}\nwindexuse{\nwixident{Formula.mk{\_}imp}}{Formula.mk:unimp}{NW1NjfzY-1iptma-1}\nwindexuse{\nwixident{Thm.assume}}{Thm.assume}{NW1NjfzY-1iptma-1}\nwindexuse{\nwixident{Thm.disch}}{Thm.disch}{NW1NjfzY-1iptma-1}\nwindexuse{\nwixident{Thm.modus{\_}ponens}}{Thm.modus:unponens}{NW1NjfzY-1iptma-1}\nwendcode{}\nwbegindocs{58}\nwdocspar
\begin{node}
Uncurrying $\vdash(A\implies(B\implies C))\implies(A\land B\implies C)$
\end{node}

\begin{proof}
$\vdash(A\implies(B\implies C))\implies(A\land B\implies C)$
\begin{enumerate}
\item\label{step:derived.nw:uncurrying:step-1} $A\land B,A\implies(B\implies C)\vdash A\implies(B\implies C)$ by assuming
\item $A\land B,A\implies(B\implies C)\vdash A\land B$ by assuming
\item\label{step:derived.nw:uncurrying:obtain-a} $A\land B,A\implies(B\implies C)\vdash A$ by $\land$-elim-l~(2)
\item\label{step:derived.nw:uncurrying:obtain-b-implies-c} $A\land B,A\implies(B\implies C)\vdash B\implies C$ by modus ponens(\ref{step:derived.nw:uncurrying:step-1}, \ref{step:derived.nw:uncurrying:obtain-a})
\item\label{step:derived.nw:uncurrying:obtain-b} $A\land B,A\implies(B\implies C)\vdash B$ by $\land$-elim-r~(2)
\item $A\land B,A\implies(B\implies C)\vdash C$ by modus ponens(\ref{step:derived.nw:uncurrying:obtain-b-implies-c}, \ref{step:derived.nw:uncurrying:obtain-b})
\item $A\implies(B\implies C)\vdash A\land B\implies C$ by discharging
\item $\vdash(A\implies(B\implies C))\implies(A\land B\implies C)$ by discharging\qedhere
\end{enumerate}
\end{proof}

\nwenddocs{}\nwbegincode{59}\sublabel{NW1NjfzY-1JZnQr-1}\nwmargintag{{\nwtagstyle{}\subpageref{NW1NjfzY-1JZnQr-1}}}\moddef{Uncurrying implication~{\nwtagstyle{}\subpageref{NW1NjfzY-1JZnQr-1}}}\endmoddef\nwstartdeflinemarkup\nwusesondefline{\\{NW1NjfzY-4CUt2L-5}}\nwenddeflinemarkup
fun uncurry_imp A B C =
  let
    val AB = \nwlinkedidentc{Formula.mk_and}{NW3w0iRO-3SlqND-3}(A, B);
    val A_imp_B_imp_C = \nwlinkedidentc{Formula.mk_imp}{NW3w0iRO-3SlqND-3}(A, \nwlinkedidentc{Formula.mk_imp}{NW3w0iRO-3SlqND-3}(B, C));
    val th1 = \nwlinkedidentc{Thm.assume}{NW4ICsJQ-2zKkPU-2} A_imp_B_imp_C [AB];
    val th2 = \nwlinkedidentc{Thm.assume}{NW4ICsJQ-2zKkPU-2} AB [A_imp_B_imp_C];
    val th3 = Thm.and_elim_l th2; (* ... |- A *)
    val th4 = \nwlinkedidentc{Thm.modus_ponens}{NW4ICsJQ-2zKkPU-2} th1 th3; (* ... |- B => C *)
    val th5 = Thm.and_elim_r th2; (* ... |- B *)
    val th6 = \nwlinkedidentc{Thm.modus_ponens}{NW4ICsJQ-2zKkPU-2} th4 th5; (* ... |- C *);
    val th7 = \nwlinkedidentc{Thm.disch}{NW4ICsJQ-2zKkPU-2} AB th6;
  in
    \nwlinkedidentc{Thm.disch}{NW4ICsJQ-2zKkPU-2} A_imp_B_imp_C th7
  end;
\nwindexdefn{\nwixident{Derived.uncurry{\_}imp}}{Derived.uncurry:unimp}{NW1NjfzY-1JZnQr-1}\eatline
\nwused{\\{NW1NjfzY-4CUt2L-5}}\nwidentdefs{\\{{\nwixident{Derived.uncurry{\_}imp}}{Derived.uncurry:unimp}}}\nwidentuses{\\{{\nwixident{Formula.mk{\_}and}}{Formula.mk:unand}}\\{{\nwixident{Formula.mk{\_}imp}}{Formula.mk:unimp}}\\{{\nwixident{Thm.assume}}{Thm.assume}}\\{{\nwixident{Thm.disch}}{Thm.disch}}\\{{\nwixident{Thm.modus{\_}ponens}}{Thm.modus:unponens}}}\nwindexuse{\nwixident{Formula.mk{\_}and}}{Formula.mk:unand}{NW1NjfzY-1JZnQr-1}\nwindexuse{\nwixident{Formula.mk{\_}imp}}{Formula.mk:unimp}{NW1NjfzY-1JZnQr-1}\nwindexuse{\nwixident{Thm.assume}}{Thm.assume}{NW1NjfzY-1JZnQr-1}\nwindexuse{\nwixident{Thm.disch}}{Thm.disch}{NW1NjfzY-1JZnQr-1}\nwindexuse{\nwixident{Thm.modus{\_}ponens}}{Thm.modus:unponens}{NW1NjfzY-1JZnQr-1}\nwendcode{}\nwbegindocs{60}\nwdocspar
\begin{node}
The hypothetical syllogism $\vdash(A\implies B)\implies((B\implies C)\implies(A\implies C))$.
\end{node}

\begin{proof}
\begin{enumerate}
\item $A,A\implies B, B\implies C\vdash A$ by assume
\item $A,A\implies B, B\implies C\vdash A\implies B$ by assume
\item $A,A\implies B, B\implies C\vdash B$ by modus ponens(2, 1)
\item $A,A\implies B, B\implies C\vdash B\implies C$ by assume
\item $A,A\implies B, B\implies C\vdash C$ by modus ponens(4, 3)
\item $A\implies B, B\implies C\vdash A\implies C$ by discharging $A$
\item $A\implies B\vdash(B\implies C)\implies(A\implies C)$ by discharging $B\implies C$
\item $\vdash(A\implies B)\implies\bigl((B\implies C)\implies(A\implies C)\bigr)$ by discharging $A\implies B$\qedhere
\end{enumerate}
\end{proof}

\nwenddocs{}\nwbegincode{61}\sublabel{NW1NjfzY-1KT69U-1}\nwmargintag{{\nwtagstyle{}\subpageref{NW1NjfzY-1KT69U-1}}}\moddef{Hypothetical syllogism~{\nwtagstyle{}\subpageref{NW1NjfzY-1KT69U-1}}}\endmoddef\nwstartdeflinemarkup\nwusesondefline{\\{NW1NjfzY-4CUt2L-5}}\nwenddeflinemarkup
fun hypo_syll_thm A B C =
  let
    val A_imp_B = \nwlinkedidentc{Formula.mk_imp}{NW3w0iRO-3SlqND-3}(A, B);
    val B_imp_C = \nwlinkedidentc{Formula.mk_imp}{NW3w0iRO-3SlqND-3}(B, C);
    val th1 = \nwlinkedidentc{Thm.assume}{NW4ICsJQ-2zKkPU-2} A [A_imp_B, B_imp_C];
    val th2 = \nwlinkedidentc{Thm.assume}{NW4ICsJQ-2zKkPU-2} A_imp_B [A, B_imp_C];
    val th3 = \nwlinkedidentc{Thm.modus_ponens}{NW4ICsJQ-2zKkPU-2} th2 th1;
    val th4 = \nwlinkedidentc{Thm.assume}{NW4ICsJQ-2zKkPU-2} B_imp_C [A, B_imp_C];
    val th5 = \nwlinkedidentc{Thm.modus_ponens}{NW4ICsJQ-2zKkPU-2} th4 th3;
    val th6 = \nwlinkedidentc{Thm.disch}{NW4ICsJQ-2zKkPU-2} A th5;
    val th7 = \nwlinkedidentc{Thm.disch}{NW4ICsJQ-2zKkPU-2} B_imp_C th6;
  in
    \nwlinkedidentc{Thm.disch}{NW4ICsJQ-2zKkPU-2} A_imp_B th7
  end;
\nwindexdefn{\nwixident{Derived.hypo{\_}syll{\_}thm}}{Derived.hypo:unsyll:unthm}{NW1NjfzY-1KT69U-1}\eatline
\nwused{\\{NW1NjfzY-4CUt2L-5}}\nwidentdefs{\\{{\nwixident{Derived.hypo{\_}syll{\_}thm}}{Derived.hypo:unsyll:unthm}}}\nwidentuses{\\{{\nwixident{Formula.mk{\_}imp}}{Formula.mk:unimp}}\\{{\nwixident{Thm.assume}}{Thm.assume}}\\{{\nwixident{Thm.disch}}{Thm.disch}}\\{{\nwixident{Thm.modus{\_}ponens}}{Thm.modus:unponens}}}\nwindexuse{\nwixident{Formula.mk{\_}imp}}{Formula.mk:unimp}{NW1NjfzY-1KT69U-1}\nwindexuse{\nwixident{Thm.assume}}{Thm.assume}{NW1NjfzY-1KT69U-1}\nwindexuse{\nwixident{Thm.disch}}{Thm.disch}{NW1NjfzY-1KT69U-1}\nwindexuse{\nwixident{Thm.modus{\_}ponens}}{Thm.modus:unponens}{NW1NjfzY-1KT69U-1}\nwendcode{}\nwbegindocs{62}\nwdocspar
\begin{node}
Transitivity follows from uncurrying and hypothetical syllogism.
\end{node}

\begin{proof}
\begin{enumerate}
\item $\vdash((A\implies B)\implies((B\implies C)\implies(A\implies C)))\implies((A\implies B)\land (B\implies C)\implies (A\implies C))$
  by uncurrying $A\implies B$, $B\implies C$, $A\implies C$
\item $\vdash(A\implies B)\implies((B\implies C)\implies(A\implies C))$
  by hypothetical syllogism $A$, $B$, $C$
\item $\vdash(A\implies B)\land (B\implies C)\implies (A\implies C)$
  by modus ponens(1, 2)\qedhere
\end{enumerate}
\end{proof}

\nwenddocs{}\nwbegincode{63}\sublabel{NW1NjfzY-1sm2GN-1}\nwmargintag{{\nwtagstyle{}\subpageref{NW1NjfzY-1sm2GN-1}}}\moddef{Transitivity of implies~{\nwtagstyle{}\subpageref{NW1NjfzY-1sm2GN-1}}}\endmoddef\nwstartdeflinemarkup\nwusesondefline{\\{NW1NjfzY-4CUt2L-5}}\nwenddeflinemarkup
fun imp_trans A B C =
  \nwlinkedidentc{Thm.modus_ponens}{NW4ICsJQ-2zKkPU-2} (uncurry_imp (\nwlinkedidentc{Formula.mk_imp}{NW3w0iRO-3SlqND-3}(A, B))
                                (\nwlinkedidentc{Formula.mk_imp}{NW3w0iRO-3SlqND-3}(B, C))
                                (\nwlinkedidentc{Formula.mk_imp}{NW3w0iRO-3SlqND-3}(A, C)))
                   (hypo_syll_thm A B C);
\nwindexdefn{\nwixident{Derived.imp{\_}trans}}{Derived.imp:untrans}{NW1NjfzY-1sm2GN-1}\eatline
\nwused{\\{NW1NjfzY-4CUt2L-5}}\nwidentdefs{\\{{\nwixident{Derived.imp{\_}trans}}{Derived.imp:untrans}}}\nwidentuses{\\{{\nwixident{Formula.mk{\_}imp}}{Formula.mk:unimp}}\\{{\nwixident{Thm.modus{\_}ponens}}{Thm.modus:unponens}}}\nwindexuse{\nwixident{Formula.mk{\_}imp}}{Formula.mk:unimp}{NW1NjfzY-1sm2GN-1}\nwindexuse{\nwixident{Thm.modus{\_}ponens}}{Thm.modus:unponens}{NW1NjfzY-1sm2GN-1}\nwendcode{}\nwbegindocs{64}\nwdocspar


\begin{node}
The contrapositive comes in two flavors for Intuitionistic logic. The
first form we will prove is $\vdash(A\implies B)\implies(\neg B\implies\neg A)$.
\end{node}

\nwenddocs{}\nwbegincode{65}\sublabel{NW1NjfzY-4KFbFg-1}\nwmargintag{{\nwtagstyle{}\subpageref{NW1NjfzY-4KFbFg-1}}}\moddef{Contrapositive derived rules~{\nwtagstyle{}\subpageref{NW1NjfzY-4KFbFg-1}}}\endmoddef\nwstartdeflinemarkup\nwusesondefline{\\{NW1NjfzY-4CUt2L-5}}\nwprevnextdefs{\relax}{NW1NjfzY-4KFbFg-2}\nwenddeflinemarkup
fun contrapositive A B =
  let
    val A_imp_B = \nwlinkedidentc{Formula.mk_imp}{NW3w0iRO-3SlqND-3}(A, B);
    val not_B = \nwlinkedidentc{Formula.mk_not}{NW3w0iRO-3SlqND-2}(B);
    val th1 = \nwlinkedidentc{Thm.assume}{NW4ICsJQ-2zKkPU-2} A_imp_B [A, not_B];
    val th2 = \nwlinkedidentc{Thm.assume}{NW4ICsJQ-2zKkPU-2} A [A_imp_B, not_B];
    val th3 = \nwlinkedidentc{Thm.modus_ponens}{NW4ICsJQ-2zKkPU-2} th1 th2;      (* A, not B, A => B |- B *)
    val th4 = \nwlinkedidentc{Thm.assume}{NW4ICsJQ-2zKkPU-2} not_B [A, A_imp_B]; (* A, not B, A => B |- not B *)
    val th5 = \nwlinkedidentc{Thm.not_elim}{NW4ICsJQ-2zKkPU-5} th4;
    val th6 = \nwlinkedidentc{Thm.modus_ponens}{NW4ICsJQ-2zKkPU-2} th5 th3;      (* A, not B, A => B |- false *)
    val th7 = \nwlinkedidentc{Thm.disch}{NW4ICsJQ-2zKkPU-2} A th6;           (* not B, A => B |- A => false *)
    val th8 = \nwlinkedidentc{Thm.not_intro}{NW4ICsJQ-2zKkPU-5} th7;             (* not B, A => B |- not A *)
    val th9 = \nwlinkedidentc{Thm.disch}{NW4ICsJQ-2zKkPU-2} not_B th8;
  in
    \nwlinkedidentc{Thm.disch}{NW4ICsJQ-2zKkPU-2} A_imp_B th9
  end;
\nwindexdefn{\nwixident{Derived.contrapositive}}{Derived.contrapositive}{NW1NjfzY-4KFbFg-1}\eatline
\nwalsodefined{\\{NW1NjfzY-4KFbFg-2}}\nwused{\\{NW1NjfzY-4CUt2L-5}}\nwidentdefs{\\{{\nwixident{Derived.contrapositive}}{Derived.contrapositive}}}\nwidentuses{\\{{\nwixident{Formula.mk{\_}imp}}{Formula.mk:unimp}}\\{{\nwixident{Formula.mk{\_}not}}{Formula.mk:unnot}}\\{{\nwixident{Thm.assume}}{Thm.assume}}\\{{\nwixident{Thm.disch}}{Thm.disch}}\\{{\nwixident{Thm.modus{\_}ponens}}{Thm.modus:unponens}}\\{{\nwixident{Thm.not{\_}elim}}{Thm.not:unelim}}\\{{\nwixident{Thm.not{\_}intro}}{Thm.not:unintro}}}\nwindexuse{\nwixident{Formula.mk{\_}imp}}{Formula.mk:unimp}{NW1NjfzY-4KFbFg-1}\nwindexuse{\nwixident{Formula.mk{\_}not}}{Formula.mk:unnot}{NW1NjfzY-4KFbFg-1}\nwindexuse{\nwixident{Thm.assume}}{Thm.assume}{NW1NjfzY-4KFbFg-1}\nwindexuse{\nwixident{Thm.disch}}{Thm.disch}{NW1NjfzY-4KFbFg-1}\nwindexuse{\nwixident{Thm.modus{\_}ponens}}{Thm.modus:unponens}{NW1NjfzY-4KFbFg-1}\nwindexuse{\nwixident{Thm.not{\_}elim}}{Thm.not:unelim}{NW1NjfzY-4KFbFg-1}\nwindexuse{\nwixident{Thm.not{\_}intro}}{Thm.not:unintro}{NW1NjfzY-4KFbFg-1}\nwendcode{}\nwbegindocs{66}\nwdocspar
\begin{node}
Since Intuitionistic logic lacks the law of the excluded middle, we
also want to prove $\vdash(A\implies\neg B)\implies(B\implies\neg A)$.
\end{node}

\begin{proof}
\begin{enumerate}
\item $(A\implies B),A,\neg B\vdash A\implies B$ by assume
\item $(A\implies B),A,\neg B\vdash A$ by assume
\item $(A\implies B),A,\neg B\vdash B$ by MP(1,2)
\item $(A\implies B),A,\neg B\vdash \neg B$ by assume
\item $(A\implies B),A,\neg B\vdash B\implies\bot$ by $\neg$-elim(4)
\item $(A\implies B),A,\neg B\vdash \bot$ by MP(5,3)
\item $(A\implies B),\neg B\vdash A\implies\bot$ by discharging $A$
\item $(A\implies B),\neg B\vdash \neg A$ by $\neg$-intro(7)
\item $(A\implies B)\vdash\neg B\implies \neg A$ by discharging $\neg B$
\item $\vdash(A\implies B)\implies(\neg B\implies \neg A)$ by discharging $\neg B$\qedhere
\end{enumerate}
\end{proof}

\nwenddocs{}\nwbegincode{67}\sublabel{NW1NjfzY-4KFbFg-2}\nwmargintag{{\nwtagstyle{}\subpageref{NW1NjfzY-4KFbFg-2}}}\moddef{Contrapositive derived rules~{\nwtagstyle{}\subpageref{NW1NjfzY-4KFbFg-1}}}\plusendmoddef\nwstartdeflinemarkup\nwusesondefline{\\{NW1NjfzY-4CUt2L-5}}\nwprevnextdefs{NW1NjfzY-4KFbFg-1}{\relax}\nwenddeflinemarkup
fun contrapositive_not A B =
  let
    val not_B = \nwlinkedidentc{Formula.mk_not}{NW3w0iRO-3SlqND-2} B;
    val not_A = \nwlinkedidentc{Formula.mk_not}{NW3w0iRO-3SlqND-2} A;
    val A_imp_not_B = \nwlinkedidentc{Formula.mk_imp}{NW3w0iRO-3SlqND-3}(A, not_B);
    val not_A_imp_B = \nwlinkedidentc{Formula.mk_imp}{NW3w0iRO-3SlqND-3}(not_A, B);
    val th1 = \nwlinkedidentc{Thm.assume}{NW4ICsJQ-2zKkPU-2} A_imp_not_B [A, B];
    val th2 = \nwlinkedidentc{Thm.assume}{NW4ICsJQ-2zKkPU-2} A [A_imp_not_B, B];
    val th3 = \nwlinkedidentc{Thm.modus_ponens}{NW4ICsJQ-2zKkPU-2} th1 th2; (* A, B, A => not B |- not B *)
    val th4 = \nwlinkedidentc{Thm.not_elim}{NW4ICsJQ-2zKkPU-5} th3;         (* A, B, A => not B |- B => false *)
    val th5 = \nwlinkedidentc{Thm.assume}{NW4ICsJQ-2zKkPU-2} B [A_imp_not_B, A];
    val th6 = \nwlinkedidentc{Thm.modus_ponens}{NW4ICsJQ-2zKkPU-2} th4 th5; (* A, B, A => not B |- false *)
    val th7 = \nwlinkedidentc{Thm.disch}{NW4ICsJQ-2zKkPU-2} A th6;      (* B, A => not B |- A => false *)
    val th8 = \nwlinkedidentc{Thm.not_intro}{NW4ICsJQ-2zKkPU-5} th7;        (* B, A => not B |- not A *)
    val th9 = \nwlinkedidentc{Thm.disch}{NW4ICsJQ-2zKkPU-2} B th8;      (* A => not B |- B => not A *)
  in
    \nwlinkedidentc{Thm.disch}{NW4ICsJQ-2zKkPU-2} A_imp_not_B th9
  end;
\nwindexdefn{\nwixident{Derived.contrapositive{\_}not}}{Derived.contrapositive:unnot}{NW1NjfzY-4KFbFg-2}\eatline
\nwused{\\{NW1NjfzY-4CUt2L-5}}\nwidentdefs{\\{{\nwixident{Derived.contrapositive{\_}not}}{Derived.contrapositive:unnot}}}\nwidentuses{\\{{\nwixident{Formula.mk{\_}imp}}{Formula.mk:unimp}}\\{{\nwixident{Formula.mk{\_}not}}{Formula.mk:unnot}}\\{{\nwixident{Thm.assume}}{Thm.assume}}\\{{\nwixident{Thm.disch}}{Thm.disch}}\\{{\nwixident{Thm.modus{\_}ponens}}{Thm.modus:unponens}}\\{{\nwixident{Thm.not{\_}elim}}{Thm.not:unelim}}\\{{\nwixident{Thm.not{\_}intro}}{Thm.not:unintro}}}\nwindexuse{\nwixident{Formula.mk{\_}imp}}{Formula.mk:unimp}{NW1NjfzY-4KFbFg-2}\nwindexuse{\nwixident{Formula.mk{\_}not}}{Formula.mk:unnot}{NW1NjfzY-4KFbFg-2}\nwindexuse{\nwixident{Thm.assume}}{Thm.assume}{NW1NjfzY-4KFbFg-2}\nwindexuse{\nwixident{Thm.disch}}{Thm.disch}{NW1NjfzY-4KFbFg-2}\nwindexuse{\nwixident{Thm.modus{\_}ponens}}{Thm.modus:unponens}{NW1NjfzY-4KFbFg-2}\nwindexuse{\nwixident{Thm.not{\_}elim}}{Thm.not:unelim}{NW1NjfzY-4KFbFg-2}\nwindexuse{\nwixident{Thm.not{\_}intro}}{Thm.not:unintro}{NW1NjfzY-4KFbFg-2}\nwendcode{}\nwbegindocs{68}\nwdocspar
\begin{node}
$\vdash\neg(A\lor B)\iff((\neg A)\land(\neg B))$
\end{node}

\begin{proof}
Backwards direction:
\begin{enumerate}
\item $\neg A\land\neg B\vdash \neg A\land\neg B$ by assuming $\neg A\land\neg B$%1
\item $\neg A\land\neg B\vdash \neg A$ by $\land$-elim-l(1)%2
\item $\neg A\land\neg B\vdash A\implies\bot$ by $\neg$-elim(2)%3
\item $\neg A\land\neg B\vdash\neg B$ by $\land$-elim-r(1)%4
\item $\neg A\land\neg B\vdash B\implies\bot$ by $\neg$-elim(4)%5
\item $\neg A\land\neg B\vdash A\lor B\implies\bot$ by $\lor$-elim(3, 5)%6
\item $\neg A\land\neg B\vdash\neg(A\lor B)$ by $\neg$-intro(6)%7
\item $\vdash(\neg A\land\neg B)\implies\neg(A\lor B)$ by discharging 7
\end{enumerate}
Forwards direction:
\begin{enumerate}
\item $A\lor B,\bot\vdash\bot$ by assuming $\bot$%1
\item $A\lor B,\bot\vdash\neg A\land\neg B$ by contr($\neg A\land\neg B$)%2
\item $A\lor B\vdash\bot\implies(\neg A\land\neg B)$ by discharging $\bot$%3
\item $\vdash(A\lor B)\implies(\bot\implies(\neg A\land\neg B))$ by discharging%4
\item $\vdash((A\lor B)\implies(\bot\implies(\neg A\land\neg B)))\implies(((A\lor B)\implies\bot)\implies(\neg A\land\neg B))$
  by {\Tt{}frege{\_}id\nwendquote}($A\lor B$, $\bot$, $\neg A\land\neg B$)%5
\item $\vdash(((A\lor B)\implies\bot)\implies(\neg A\land\neg B))$ by
  modus ponens(5, 4)%6
\item $\vdash(\neg(A\lor B))\implies((A\lor B)\implies\bot)$%7
\item $\vdash(\neg(A\lor B))\implies(\neg A\land\neg B)$ by
  hypothetical syllogism(7, 6)\qedhere
\end{enumerate}
\end{proof}

\nwenddocs{}\nwbegincode{69}\sublabel{NW1NjfzY-4UPWei-1}\nwmargintag{{\nwtagstyle{}\subpageref{NW1NjfzY-4UPWei-1}}}\moddef{Negation of disjunction~{\nwtagstyle{}\subpageref{NW1NjfzY-4UPWei-1}}}\endmoddef\nwstartdeflinemarkup\nwusesondefline{\\{NW1NjfzY-4CUt2L-6}}\nwenddeflinemarkup
local
  fun backward A B =
    let
      val not_A_and_not_B = \nwlinkedidentc{Formula.mk_and}{NW3w0iRO-3SlqND-3}(\nwlinkedidentc{Formula.mk_not}{NW3w0iRO-3SlqND-2}(A),
                                           \nwlinkedidentc{Formula.mk_not}{NW3w0iRO-3SlqND-2}(B));
      val th1 = \nwlinkedidentc{Thm.assume}{NW4ICsJQ-2zKkPU-2} not_A_and_not_B [];
      val th2 = Thm.and_elim_l th1;
      val th3 = \nwlinkedidentc{Thm.not_elim}{NW4ICsJQ-2zKkPU-5} th2;
      val th4 = Thm.and_elim_r th1;
      val th5 = \nwlinkedidentc{Thm.not_elim}{NW4ICsJQ-2zKkPU-5} th4;
      val th6 = \nwlinkedidentc{Thm.or_elim}{NW4ICsJQ-2zKkPU-4} th3 th5;
      val th7 = \nwlinkedidentc{Thm.not_intro}{NW4ICsJQ-2zKkPU-5} th6;
    in
      \nwlinkedidentc{Thm.disch}{NW4ICsJQ-2zKkPU-2} not_A_and_not_B th7
    end;
  fun forward A B =
    let
      val A_or_B = \nwlinkedidentc{Formula.mk_or}{NW3w0iRO-3SlqND-3}(A, B);
      val not_A_and_not_B = \nwlinkedidentc{Formula.mk_and}{NW3w0iRO-3SlqND-3}(\nwlinkedidentc{Formula.mk_not}{NW3w0iRO-3SlqND-2}(A),
                                           \nwlinkedidentc{Formula.mk_not}{NW3w0iRO-3SlqND-2}(B));
      val F = \nwlinkedidentc{Formula.mk_falsum}{NW3w0iRO-3SlqND-1} ();
      val th1 = \nwlinkedidentc{Thm.assume}{NW4ICsJQ-2zKkPU-2} F [A_or_B];
      val th2 = \nwlinkedidentc{Thm.contr}{NW4ICsJQ-2zKkPU-6} not_A_and_not_B th1;
      val th3 = \nwlinkedidentc{Thm.disch}{NW4ICsJQ-2zKkPU-2} F th2;
      val th4 = \nwlinkedidentc{Thm.disch}{NW4ICsJQ-2zKkPU-2} A_or_B th3;
      val th5 = frege_id A_or_B F not_A_and_not_B;
      val th6 = \nwlinkedidentc{Thm.modus_ponens}{NW4ICsJQ-2zKkPU-2} th5 th4;
      val th7 = \nwlinkedidentc{Thm.disch}{NW4ICsJQ-2zKkPU-2} not_A_and_not_B
                              (\nwlinkedidentc{Thm.not_elim}{NW4ICsJQ-2zKkPU-5} (\nwlinkedidentc{Thm.assume}{NW4ICsJQ-2zKkPU-2} not_A_and_not_B []));
    in
      hypo_syll th7 th6
    end;
in
fun not_or A B =
  \nwlinkedidentc{Thm.iff_intro}{NW4ICsJQ-2zKkPU-7} (forward A B) (backward A B)
end;
\nwindexdefn{\nwixident{Derived.not{\_}or}}{Derived.not:unor}{NW1NjfzY-4UPWei-1}\eatline
\nwused{\\{NW1NjfzY-4CUt2L-6}}\nwidentdefs{\\{{\nwixident{Derived.not{\_}or}}{Derived.not:unor}}}\nwidentuses{\\{{\nwixident{Formula.mk{\_}and}}{Formula.mk:unand}}\\{{\nwixident{Formula.mk{\_}falsum}}{Formula.mk:unfalsum}}\\{{\nwixident{Formula.mk{\_}not}}{Formula.mk:unnot}}\\{{\nwixident{Formula.mk{\_}or}}{Formula.mk:unor}}\\{{\nwixident{Thm.assume}}{Thm.assume}}\\{{\nwixident{Thm.contr}}{Thm.contr}}\\{{\nwixident{Thm.disch}}{Thm.disch}}\\{{\nwixident{Thm.iff{\_}intro}}{Thm.iff:unintro}}\\{{\nwixident{Thm.modus{\_}ponens}}{Thm.modus:unponens}}\\{{\nwixident{Thm.not{\_}elim}}{Thm.not:unelim}}\\{{\nwixident{Thm.not{\_}intro}}{Thm.not:unintro}}\\{{\nwixident{Thm.or{\_}elim}}{Thm.or:unelim}}}\nwindexuse{\nwixident{Formula.mk{\_}and}}{Formula.mk:unand}{NW1NjfzY-4UPWei-1}\nwindexuse{\nwixident{Formula.mk{\_}falsum}}{Formula.mk:unfalsum}{NW1NjfzY-4UPWei-1}\nwindexuse{\nwixident{Formula.mk{\_}not}}{Formula.mk:unnot}{NW1NjfzY-4UPWei-1}\nwindexuse{\nwixident{Formula.mk{\_}or}}{Formula.mk:unor}{NW1NjfzY-4UPWei-1}\nwindexuse{\nwixident{Thm.assume}}{Thm.assume}{NW1NjfzY-4UPWei-1}\nwindexuse{\nwixident{Thm.contr}}{Thm.contr}{NW1NjfzY-4UPWei-1}\nwindexuse{\nwixident{Thm.disch}}{Thm.disch}{NW1NjfzY-4UPWei-1}\nwindexuse{\nwixident{Thm.iff{\_}intro}}{Thm.iff:unintro}{NW1NjfzY-4UPWei-1}\nwindexuse{\nwixident{Thm.modus{\_}ponens}}{Thm.modus:unponens}{NW1NjfzY-4UPWei-1}\nwindexuse{\nwixident{Thm.not{\_}elim}}{Thm.not:unelim}{NW1NjfzY-4UPWei-1}\nwindexuse{\nwixident{Thm.not{\_}intro}}{Thm.not:unintro}{NW1NjfzY-4UPWei-1}\nwindexuse{\nwixident{Thm.or{\_}elim}}{Thm.or:unelim}{NW1NjfzY-4UPWei-1}\nwendcode{}\nwbegindocs{70}\nwdocspar
\begin{node}
$\vdash((\neg A)\lor(\neg B))\implies\neg(A\land B)$
\end{node}

\begin{proof}
\begin{enumerate}
\item $A\land B\vdash A\land B$ by assuming $A\land B$%1
\item $A\land B\vdash A$ by $\land$-elim-l(1)%2
\item $\vdash(A\land B)\implies A$ by discharging 2%3
\item $\vdash((A\land B)\implies A)\implies(\neg A\implies\neg(A\land B))$
  by {\Tt{}contrapositive\nwendquote}($A\land B$, $A$)%4
\item\label{step:disjunction-of-nots:coup-1} $\vdash\neg A\implies\neg(A\land B)$
by modus ponens(4, 3)%5
\item $A\land B\vdash B$ by $\land$-elim-r(1)%6
\item $\vdash(A\land B)\implies B$ by discharging 6%7
\item $\vdash((A\land B)\implies B)\implies(\neg B\implies\neg(A\land B))$
  by {\Tt{}contrapositive\nwendquote}($A\land B$, $B$)%8
\item\label{step:disjunction-of-nots:coup-2} $\vdash\neg B\implies\neg(A\land B)$
  by modus ponens(8, 7)% 9
\item $\vdash((\neg A)\lor(\neg B))\implies\neg(A\land B)$
  by $\lor$-elim(\ref{step:disjunction-of-nots:coup-1}, \ref{step:disjunction-of-nots:coup-2})\qedhere
\end{enumerate}
\end{proof}

\nwenddocs{}\nwbegincode{71}\sublabel{NW1NjfzY-wYirn-1}\nwmargintag{{\nwtagstyle{}\subpageref{NW1NjfzY-wYirn-1}}}\moddef{Negation of conjunction~{\nwtagstyle{}\subpageref{NW1NjfzY-wYirn-1}}}\endmoddef\nwstartdeflinemarkup\nwusesondefline{\\{NW1NjfzY-4CUt2L-6}}\nwenddeflinemarkup
fun not_and A B =
  let
    val A_and_B = \nwlinkedidentc{Formula.mk_and}{NW3w0iRO-3SlqND-3}(A, B);
    val th1 = \nwlinkedidentc{Thm.assume}{NW4ICsJQ-2zKkPU-2} A_and_B [];
    fun iter X and_elim =
      let
        val lm2 = and_elim th1;
        val lm3 = \nwlinkedidentc{Thm.disch}{NW4ICsJQ-2zKkPU-2} A_and_B lm2;
        val lm4 = contrapositive A_and_B X;
      in
        \nwlinkedidentc{Thm.modus_ponens}{NW4ICsJQ-2zKkPU-2} lm4 lm3
      end;
    val th6 = iter A (Thm.and_elim_l);
    val th9 = iter B (Thm.and_elim_r);
  in
    \nwlinkedidentc{Thm.or_elim}{NW4ICsJQ-2zKkPU-4} th6 th9
  end;
\nwindexdefn{\nwixident{Derived.not{\_}and}}{Derived.not:unand}{NW1NjfzY-wYirn-1}\eatline
\nwused{\\{NW1NjfzY-4CUt2L-6}}\nwidentdefs{\\{{\nwixident{Derived.not{\_}and}}{Derived.not:unand}}}\nwidentuses{\\{{\nwixident{Formula.mk{\_}and}}{Formula.mk:unand}}\\{{\nwixident{Thm.assume}}{Thm.assume}}\\{{\nwixident{Thm.disch}}{Thm.disch}}\\{{\nwixident{Thm.modus{\_}ponens}}{Thm.modus:unponens}}\\{{\nwixident{Thm.or{\_}elim}}{Thm.or:unelim}}}\nwindexuse{\nwixident{Formula.mk{\_}and}}{Formula.mk:unand}{NW1NjfzY-wYirn-1}\nwindexuse{\nwixident{Thm.assume}}{Thm.assume}{NW1NjfzY-wYirn-1}\nwindexuse{\nwixident{Thm.disch}}{Thm.disch}{NW1NjfzY-wYirn-1}\nwindexuse{\nwixident{Thm.modus{\_}ponens}}{Thm.modus:unponens}{NW1NjfzY-wYirn-1}\nwindexuse{\nwixident{Thm.or{\_}elim}}{Thm.or:unelim}{NW1NjfzY-wYirn-1}\nwendcode{}\nwbegindocs{72}\nwdocspar
\begin{node}
$\vdash((\neg A)\lor B)\implies(A\implies B)$
\end{node}

We want to prove $\vdash\neg A\implies(A\implies B)$ and $\vdash B\implies(A\implies B)$

\begin{proof}
\begin{enumerate}
\item $\vdash A\implies(\neg A\implies B)$ by {\Tt{}weak{\_}not{\_}elim\ A\ B\nwendquote}
\item $A,\neg A\vdash B$ by undischarging $A$, $\neg A$
\item $\vdash(\neg A)\implies(A\implies B)$ by discharging $A$ and
  then discharging $\neg A$
\item $A,B\vdash B$ by assume
\item $B\vdash A\implies B$ by discharging $A$
\item $\vdash B\implies(A\implies B)$ by discharging $B$
\item $\vdash((\neg A)\lor B)\implies(A\implies B)$ by $\lor$-elimination(3,6)
\qedhere
\end{enumerate}
\end{proof}

\nwenddocs{}\nwbegincode{73}\sublabel{NW1NjfzY-3NUP7w-1}\nwmargintag{{\nwtagstyle{}\subpageref{NW1NjfzY-3NUP7w-1}}}\moddef{$(\neg A)\lor B$ sufficient to prove $A\implies B$~{\nwtagstyle{}\subpageref{NW1NjfzY-3NUP7w-1}}}\endmoddef\nwstartdeflinemarkup\nwusesondefline{\\{NW1NjfzY-4CUt2L-6}}\nwenddeflinemarkup
fun not_disj_to_imp A B =
  let
    val not_A = \nwlinkedidentc{Formula.mk_not}{NW3w0iRO-3SlqND-2}(A);
    val th1 = weak_not_elim A B;
    val th2 = \nwlinkedidentc{Thm.undisch}{NW4ICsJQ-2zKkPU-2} (\nwlinkedidentc{Thm.undisch}{NW4ICsJQ-2zKkPU-2} th1);
    val th3 = \nwlinkedidentc{Thm.disch}{NW4ICsJQ-2zKkPU-2} not_A (\nwlinkedidentc{Thm.disch}{NW4ICsJQ-2zKkPU-2} A th2);

    val th4 = \nwlinkedidentc{Thm.assume}{NW4ICsJQ-2zKkPU-2} B [A];
    val th5 = \nwlinkedidentc{Thm.disch}{NW4ICsJQ-2zKkPU-2} A th4;
    val th6 = \nwlinkedidentc{Thm.disch}{NW4ICsJQ-2zKkPU-2} B th5;
  in
    \nwlinkedidentc{Thm.or_elim}{NW4ICsJQ-2zKkPU-4} th3 th6
  end;
\nwindexdefn{\nwixident{Derived.not{\_}disj{\_}to{\_}imp}}{Derived.not:undisj:unto:unimp}{NW1NjfzY-3NUP7w-1}\eatline
\nwused{\\{NW1NjfzY-4CUt2L-6}}\nwidentdefs{\\{{\nwixident{Derived.not{\_}disj{\_}to{\_}imp}}{Derived.not:undisj:unto:unimp}}}\nwidentuses{\\{{\nwixident{Formula.mk{\_}not}}{Formula.mk:unnot}}\\{{\nwixident{Thm.assume}}{Thm.assume}}\\{{\nwixident{Thm.disch}}{Thm.disch}}\\{{\nwixident{Thm.or{\_}elim}}{Thm.or:unelim}}\\{{\nwixident{Thm.undisch}}{Thm.undisch}}}\nwindexuse{\nwixident{Formula.mk{\_}not}}{Formula.mk:unnot}{NW1NjfzY-3NUP7w-1}\nwindexuse{\nwixident{Thm.assume}}{Thm.assume}{NW1NjfzY-3NUP7w-1}\nwindexuse{\nwixident{Thm.disch}}{Thm.disch}{NW1NjfzY-3NUP7w-1}\nwindexuse{\nwixident{Thm.or{\_}elim}}{Thm.or:unelim}{NW1NjfzY-3NUP7w-1}\nwindexuse{\nwixident{Thm.undisch}}{Thm.undisch}{NW1NjfzY-3NUP7w-1}\nwendcode{}\nwbegindocs{74}\nwdocspar
\begin{node}
$\vdash(A\implies B)\implies\neg(A\land\neg B)$
\end{node}

\begin{proof}
\begin{enumerate}
\item $A\land\neg B, A\implies B\vdash A\land\neg B$ by assume
\item $A\land\neg B, A\implies B\vdash A$ by $\land$-elim-l(1)
\item $A\land\neg B, A\implies B\vdash\neg B$ by $\land$-elim-r(1)
\item $A\land\neg B, A\implies B\vdash B\implies\bot$ by $\neg$-elim(3)
\item $A\land\neg B, A\implies B\vdash A\implies B$ by assume
\item $A\land\neg B, A\implies B\vdash B$ by modus ponens(5, 2)
\item $A\land\neg B, A\implies B\vdash \bot$ by modus ponens(4, 6)
\item $A\implies B\vdash(A\land\neg B)\implies \bot$ by discharging($A\land\neg B$)
\item $A\implies B\vdash\neg(A\land\neg B)$ by $\neg$-introduction
\item $\vdash(A\implies B)\implies\neg(A\land\neg B)$ by discharging\qedhere
\end{enumerate}
\end{proof}

\nwenddocs{}\nwbegincode{75}\sublabel{NW1NjfzY-2qUp6m-1}\nwmargintag{{\nwtagstyle{}\subpageref{NW1NjfzY-2qUp6m-1}}}\moddef{Implication as negated conjunction~{\nwtagstyle{}\subpageref{NW1NjfzY-2qUp6m-1}}}\endmoddef\nwstartdeflinemarkup\nwusesondefline{\\{NW1NjfzY-4CUt2L-6}}\nwenddeflinemarkup
fun imp_is_not_and A B =
  let
    val not_B = \nwlinkedidentc{Formula.mk_not}{NW3w0iRO-3SlqND-2}(B);
    val A_and_not_B = \nwlinkedidentc{Formula.mk_and}{NW3w0iRO-3SlqND-3}(A, not_B);
    val A_imp_B = \nwlinkedidentc{Formula.mk_imp}{NW3w0iRO-3SlqND-3}(A, B);
    val th1 = \nwlinkedidentc{Thm.assume}{NW4ICsJQ-2zKkPU-2} A_and_not_B [A_imp_B];
    val th2 = Thm.and_elim_l th1;            (* ... |- A *)
    val th3 = Thm.and_elim_r th1;            (* ... |- not B *)
    val th4 = \nwlinkedidentc{Thm.not_elim}{NW4ICsJQ-2zKkPU-5} th3;              (* ... |- B => false *)
    val th5 = \nwlinkedidentc{Thm.assume}{NW4ICsJQ-2zKkPU-2} A_imp_B [A_and_not_B]; (* ... |- A => B *)
    val th6 = \nwlinkedidentc{Thm.modus_ponens}{NW4ICsJQ-2zKkPU-2} th5 th2;      (* ... |- B *)
    val th7 = \nwlinkedidentc{Thm.modus_ponens}{NW4ICsJQ-2zKkPU-2} th4 th6;      (* ... |- false *)
    val th8 = \nwlinkedidentc{Thm.disch}{NW4ICsJQ-2zKkPU-2} A_and_not_B th7; (* ... |- (A & not B) => false *)
    val th9 = \nwlinkedidentc{Thm.not_intro}{NW4ICsJQ-2zKkPU-5} th8;             (* ... |- not (A & not B) *)
  in
    \nwlinkedidentc{Thm.disch}{NW4ICsJQ-2zKkPU-2} A_imp_B th9
  end;
\nwindexdefn{\nwixident{Derived.imp{\_}is{\_}not{\_}and}}{Derived.imp:unis:unnot:unand}{NW1NjfzY-2qUp6m-1}\eatline
\nwused{\\{NW1NjfzY-4CUt2L-6}}\nwidentdefs{\\{{\nwixident{Derived.imp{\_}is{\_}not{\_}and}}{Derived.imp:unis:unnot:unand}}}\nwidentuses{\\{{\nwixident{Formula.mk{\_}and}}{Formula.mk:unand}}\\{{\nwixident{Formula.mk{\_}imp}}{Formula.mk:unimp}}\\{{\nwixident{Formula.mk{\_}not}}{Formula.mk:unnot}}\\{{\nwixident{Thm.assume}}{Thm.assume}}\\{{\nwixident{Thm.disch}}{Thm.disch}}\\{{\nwixident{Thm.modus{\_}ponens}}{Thm.modus:unponens}}\\{{\nwixident{Thm.not{\_}elim}}{Thm.not:unelim}}\\{{\nwixident{Thm.not{\_}intro}}{Thm.not:unintro}}}\nwindexuse{\nwixident{Formula.mk{\_}and}}{Formula.mk:unand}{NW1NjfzY-2qUp6m-1}\nwindexuse{\nwixident{Formula.mk{\_}imp}}{Formula.mk:unimp}{NW1NjfzY-2qUp6m-1}\nwindexuse{\nwixident{Formula.mk{\_}not}}{Formula.mk:unnot}{NW1NjfzY-2qUp6m-1}\nwindexuse{\nwixident{Thm.assume}}{Thm.assume}{NW1NjfzY-2qUp6m-1}\nwindexuse{\nwixident{Thm.disch}}{Thm.disch}{NW1NjfzY-2qUp6m-1}\nwindexuse{\nwixident{Thm.modus{\_}ponens}}{Thm.modus:unponens}{NW1NjfzY-2qUp6m-1}\nwindexuse{\nwixident{Thm.not{\_}elim}}{Thm.not:unelim}{NW1NjfzY-2qUp6m-1}\nwindexuse{\nwixident{Thm.not{\_}intro}}{Thm.not:unintro}{NW1NjfzY-2qUp6m-1}\nwendcode{}\nwbegindocs{76}\nwdocspar
\begin{node}
$\vdash A\iff A$
\end{node}

\nwenddocs{}\nwbegincode{77}\sublabel{NW1NjfzY-YtLDN-1}\nwmargintag{{\nwtagstyle{}\subpageref{NW1NjfzY-YtLDN-1}}}\moddef{Iff is reflexive~{\nwtagstyle{}\subpageref{NW1NjfzY-YtLDN-1}}}\endmoddef\nwstartdeflinemarkup\nwusesondefline{\\{NW1NjfzY-4CUt2L-7}}\nwenddeflinemarkup
fun iff_refl A =
  let val th1 = imp_reflexive A
  in \nwlinkedidentc{Thm.iff_intro}{NW4ICsJQ-2zKkPU-7} th1 th1 end;
\nwindexdefn{\nwixident{Derived.iff{\_}refl}}{Derived.iff:unrefl}{NW1NjfzY-YtLDN-1}\eatline
\nwused{\\{NW1NjfzY-4CUt2L-7}}\nwidentdefs{\\{{\nwixident{Derived.iff{\_}refl}}{Derived.iff:unrefl}}}\nwidentuses{\\{{\nwixident{Thm.iff{\_}intro}}{Thm.iff:unintro}}}\nwindexuse{\nwixident{Thm.iff{\_}intro}}{Thm.iff:unintro}{NW1NjfzY-YtLDN-1}\nwendcode{}\nwbegindocs{78}\nwdocspar
\begin{node}
$\vdash(A\iff B)\implies(B\iff A)$
\end{node}

\nwenddocs{}\nwbegincode{79}\sublabel{NW1NjfzY-3hb2X5-1}\nwmargintag{{\nwtagstyle{}\subpageref{NW1NjfzY-3hb2X5-1}}}\moddef{Iff is symmetric~{\nwtagstyle{}\subpageref{NW1NjfzY-3hb2X5-1}}}\endmoddef\nwstartdeflinemarkup\nwusesondefline{\\{NW1NjfzY-4CUt2L-7}}\nwenddeflinemarkup
fun iff_sym A B =
  let
    val A_iff_B = \nwlinkedidentc{Formula.mk_iff}{NW3w0iRO-3SlqND-3}(A, B);
    val th1 = \nwlinkedidentc{Thm.assume}{NW4ICsJQ-2zKkPU-2} A_iff_B [];
    val th2 = \nwlinkedidentc{Thm.iff_elim_l}{NW4ICsJQ-2zKkPU-7} th1; (* ... |- A => B *)
    val th3 = \nwlinkedidentc{Thm.iff_elim_r}{NW4ICsJQ-2zKkPU-7} th1; (* ... |- B => A *)
  in
    \nwlinkedidentc{Thm.disch}{NW4ICsJQ-2zKkPU-2} A_iff_B (\nwlinkedidentc{Thm.iff_intro}{NW4ICsJQ-2zKkPU-7} th3 th2)
  end;
\nwindexdefn{\nwixident{Derived.iff{\_}sym}}{Derived.iff:unsym}{NW1NjfzY-3hb2X5-1}\eatline
\nwused{\\{NW1NjfzY-4CUt2L-7}}\nwidentdefs{\\{{\nwixident{Derived.iff{\_}sym}}{Derived.iff:unsym}}}\nwidentuses{\\{{\nwixident{Formula.mk{\_}iff}}{Formula.mk:uniff}}\\{{\nwixident{Thm.assume}}{Thm.assume}}\\{{\nwixident{Thm.disch}}{Thm.disch}}\\{{\nwixident{Thm.iff{\_}elim{\_}l}}{Thm.iff:unelim:unl}}\\{{\nwixident{Thm.iff{\_}elim{\_}r}}{Thm.iff:unelim:unr}}\\{{\nwixident{Thm.iff{\_}intro}}{Thm.iff:unintro}}}\nwindexuse{\nwixident{Formula.mk{\_}iff}}{Formula.mk:uniff}{NW1NjfzY-3hb2X5-1}\nwindexuse{\nwixident{Thm.assume}}{Thm.assume}{NW1NjfzY-3hb2X5-1}\nwindexuse{\nwixident{Thm.disch}}{Thm.disch}{NW1NjfzY-3hb2X5-1}\nwindexuse{\nwixident{Thm.iff{\_}elim{\_}l}}{Thm.iff:unelim:unl}{NW1NjfzY-3hb2X5-1}\nwindexuse{\nwixident{Thm.iff{\_}elim{\_}r}}{Thm.iff:unelim:unr}{NW1NjfzY-3hb2X5-1}\nwindexuse{\nwixident{Thm.iff{\_}intro}}{Thm.iff:unintro}{NW1NjfzY-3hb2X5-1}\nwendcode{}\nwbegindocs{80}\nwdocspar
\begin{node}
$\vdash (A\iff B)\land(B\iff C)\implies(A\iff C)$
\end{node}

\begin{proof}
\begin{enumerate}
\item $(A\iff B)\land(B\iff C)\vdash(A\iff B)\land(B\iff C)$ by assuming%1
\item $(A\iff B)\land(B\iff C)\vdash A\iff B$ by $\land$-elim-l(1)%2
\item $(A\iff B)\land(B\iff C)\vdash A\implies B$ by $\iff$-elim-l(2)%3
\item $(A\iff B)\land(B\iff C)\vdash B\iff C$ by $\land$-elim-r(1)%4
\item $(A\iff B)\land(B\iff C)\vdash B\implies C$ by $\iff$-elim-r(4)%5
\item $(A\iff B)\land(B\iff C)\vdash A\implies C$ by {\Tt{}imp{\_}trans\nwendquote}(3, 5)%6
\item $(A\iff B)\land(B\iff C)\vdash C\implies B$ by $\iff$-elim-l(4)%7
\item $(A\iff B)\land(B\iff C)\vdash B\implies A$ by $\iff$-elim-r(2)%8
\item $(A\iff B)\land(B\iff C)\vdash C\implies A$ by {\Tt{}imp{\_}trans\nwendquote}(8, 9)%9
\item $(A\iff B)\land(B\iff C)\vdash A\iff C$ by $\Leftrightarrow$-intro(6, 9)
\item $\vdash((A\iff B)\land(B\iff C))\implies(A\iff C)$ by discharging hypothesis from step 10\qedhere
\end{enumerate}
\end{proof}

\nwenddocs{}\nwbegincode{81}\sublabel{NW1NjfzY-2isj9x-1}\nwmargintag{{\nwtagstyle{}\subpageref{NW1NjfzY-2isj9x-1}}}\moddef{Iff is transitive~{\nwtagstyle{}\subpageref{NW1NjfzY-2isj9x-1}}}\endmoddef\nwstartdeflinemarkup\nwusesondefline{\\{NW1NjfzY-4CUt2L-7}}\nwenddeflinemarkup
fun iff_trans A B C =
  let
    val A_iff_B = \nwlinkedidentc{Formula.mk_iff}{NW3w0iRO-3SlqND-3}(A, B);
    val B_iff_C = \nwlinkedidentc{Formula.mk_iff}{NW3w0iRO-3SlqND-3}(B, C);
    val AB_and_BC = \nwlinkedidentc{Formula.mk_and}{NW3w0iRO-3SlqND-3}(A_iff_B, B_iff_C);
    val th1 = \nwlinkedidentc{Thm.assume}{NW4ICsJQ-2zKkPU-2} AB_and_BC [];
    val th2 = Thm.and_elim_l th1; (* ... |- A iff B *)
    val th3 = \nwlinkedidentc{Thm.iff_elim_l}{NW4ICsJQ-2zKkPU-7} th2; (* ... |- A => B *)
    val th4 = Thm.and_elim_r th1; (* ... |- B iff C *)
    val th5 = \nwlinkedidentc{Thm.iff_elim_r}{NW4ICsJQ-2zKkPU-7} th4; (* ... |- B => C *)
    val th6 = Thm.and_intro th3 th5; (* ... |- (A => B) & (B => C *)
    val th7 = \nwlinkedidentc{Thm.modus_ponens}{NW4ICsJQ-2zKkPU-2} (imp_trans A B C) th6; (* ... |- A => C *)
    val th8 = \nwlinkedidentc{Thm.iff_elim_r}{NW4ICsJQ-2zKkPU-7} th2; (* ... |- B => A *)
    val th9 = \nwlinkedidentc{Thm.iff_elim_l}{NW4ICsJQ-2zKkPU-7} th4; (* ... |- C => B *)
    val th10 = Thm.and_intro th4 th8; (* ... |- (C => B) & (B => A *)
    val th11 = \nwlinkedidentc{Thm.modus_ponens}{NW4ICsJQ-2zKkPU-2} (imp_trans C B A) th10; (* ... |- C => A *)
    val th12 = \nwlinkedidentc{Thm.iff_intro}{NW4ICsJQ-2zKkPU-7} th7 th11; (* ... |- A iff C *)
  in
    \nwlinkedidentc{Thm.disch}{NW4ICsJQ-2zKkPU-2} AB_and_BC th12
  end;
\nwindexdefn{\nwixident{Derived.iff{\_}trans}}{Derived.iff:untrans}{NW1NjfzY-2isj9x-1}\eatline
\nwused{\\{NW1NjfzY-4CUt2L-7}}\nwidentdefs{\\{{\nwixident{Derived.iff{\_}trans}}{Derived.iff:untrans}}}\nwidentuses{\\{{\nwixident{Formula.mk{\_}and}}{Formula.mk:unand}}\\{{\nwixident{Formula.mk{\_}iff}}{Formula.mk:uniff}}\\{{\nwixident{Thm.assume}}{Thm.assume}}\\{{\nwixident{Thm.disch}}{Thm.disch}}\\{{\nwixident{Thm.iff{\_}elim{\_}l}}{Thm.iff:unelim:unl}}\\{{\nwixident{Thm.iff{\_}elim{\_}r}}{Thm.iff:unelim:unr}}\\{{\nwixident{Thm.iff{\_}intro}}{Thm.iff:unintro}}\\{{\nwixident{Thm.modus{\_}ponens}}{Thm.modus:unponens}}}\nwindexuse{\nwixident{Formula.mk{\_}and}}{Formula.mk:unand}{NW1NjfzY-2isj9x-1}\nwindexuse{\nwixident{Formula.mk{\_}iff}}{Formula.mk:uniff}{NW1NjfzY-2isj9x-1}\nwindexuse{\nwixident{Thm.assume}}{Thm.assume}{NW1NjfzY-2isj9x-1}\nwindexuse{\nwixident{Thm.disch}}{Thm.disch}{NW1NjfzY-2isj9x-1}\nwindexuse{\nwixident{Thm.iff{\_}elim{\_}l}}{Thm.iff:unelim:unl}{NW1NjfzY-2isj9x-1}\nwindexuse{\nwixident{Thm.iff{\_}elim{\_}r}}{Thm.iff:unelim:unr}{NW1NjfzY-2isj9x-1}\nwindexuse{\nwixident{Thm.iff{\_}intro}}{Thm.iff:unintro}{NW1NjfzY-2isj9x-1}\nwindexuse{\nwixident{Thm.modus{\_}ponens}}{Thm.modus:unponens}{NW1NjfzY-2isj9x-1}\nwendcode{}\nwbegindocs{82}\nwdocspar
\vfill\eject
\begin{node}[Quantifier theorems]\label{node:fs0:derived.nw:toc:predicate-logic}
Let $C$ be such that $x$ is not free in $C$ (i.e., $x\notin FV(C)$).
We also have derived rules for quantifiers:
\begin{enumerate}\setcounter{enumi}{\value{node}}
\item \checkbox{1} $\vdash((\forall x\ldotp A[x])\lor C)\implies\forall x\ldotp
(A[x]\lor C)$%a(done)
\item \checkbox{1} $\vdash((\exists x\ldotp A[x])\lor(\exists x\ldotp B[x]))\iff\exists x\ldotp(A[x]\lor B[x])$%b(done)
\item \checkbox{1} $\vdash(\forall x\ldotp A[x]\land B[x])\iff((\forall x\ldotp A[x])\land(\forall x\ldotp B[x]))$%c(done)
\item \checkbox{1} $\vdash(\exists x\ldotp A[x]\land B[x])\implies((\exists x\ldotp A[x])\land(\exists x\ldotp B[x]))$%d(done)
\item \checkbox{1} $\vdash(\forall x\ldotp(A[x]\implies C))\iff((\exists x\ldotp A[x])\implies C)$%e(done)
\item \checkbox{1} $\vdash(\forall x\ldotp\neg A[x])\iff\neg(\exists x\ldotp A[x])$%e'(done)
\bigbreak
\item \checkbox{1} $\vdash(\forall x\ldotp(C\implies A[x]))\iff(C\implies\forall x\ldotp A[x])$%f
\item \checkbox{1} $\vdash(\exists x\ldotp(A[x]\implies C))\implies((\forall x\ldotp A[x])\implies C)$%g
\item \checkbox{1} $\vdash(\exists x\ldotp\neg A[x])\implies\neg(\forall x\ldotp A[x])$%g'
\item \checkbox{1} $\vdash(\exists x\ldotp(C\implies A[x]))\implies(C\implies\exists x\ldotp A[x])$%h
\item \checkbox{0} $\vdash(\exists x\ldotp C)\iff C$ and $\vdash(\forall x\ldotp C)\iff C$
\end{enumerate}
\end{node}

\nwenddocs{}\nwbegincode{83}\sublabel{NW1NjfzY-uZfts-1}\nwmargintag{{\nwtagstyle{}\subpageref{NW1NjfzY-uZfts-1}}}\moddef{Derived results for Intuitionistic predicate logic~{\nwtagstyle{}\subpageref{NW1NjfzY-uZfts-1}}}\endmoddef\nwstartdeflinemarkup\nwusesondefline{\\{NW1NjfzY-4dPl0p-1}}\nwprevnextdefs{\relax}{NW1NjfzY-uZfts-2}\nwenddeflinemarkup

\nwalsodefined{\\{NW1NjfzY-uZfts-2}\\{NW1NjfzY-uZfts-3}\\{NW1NjfzY-uZfts-4}\\{NW1NjfzY-uZfts-5}\\{NW1NjfzY-uZfts-6}\\{NW1NjfzY-uZfts-7}\\{NW1NjfzY-uZfts-8}\\{NW1NjfzY-uZfts-9}\\{NW1NjfzY-uZfts-A}\\{NW1NjfzY-uZfts-B}\\{NW1NjfzY-uZfts-C}\\{NW1NjfzY-uZfts-D}\\{NW1NjfzY-uZfts-E}}\nwused{\\{NW1NjfzY-4dPl0p-1}}\nwendcode{}\nwbegindocs{84}\nwdocspar

\begin{node}% all_disj?
$\vdash(\forall x\ldotp A[x]\lor C)\implies\forall x\ldotp(A[x]\lor C)$
\end{node}

\begin{proof}
\begin{enumerate}
\item $C\vdash C$ by assuming%1
\item $C\vdash A[x]\lor C$ by $\lor$-intro-l($A[x]$)%2
\item $C\vdash\forall x\ldotp(A[x]\lor C)$ by $\forall$-intro(2)%3
\item $\vdash C\implies\forall x\ldotp(A[x]\lor C)$ by discharging%4
\item $\forall x\ldotp A[x]\vdash\forall x\ldotp A[x]$ by assuming $\forall x\ldotp A[x]$%5
\item $\forall x\ldotp A[x]\vdash A[x]$ by $\forall$-elimination(5, $x$)%6
\item $\forall x\ldotp A[x]\vdash A[x]\lor C$ by $\lor$-intro-r($C$)%7
\item $\forall x\ldotp A[x]\vdash\forall x\ldotp(A[x]\lor C)$ by
  $\forall$-intro(7, $x$)%8
\item $\vdash(\forall x\ldotp A[x])\implies\forall x\ldotp(A[x]\lor C)$ by
  discharging 8%9
\item $\vdash((\forall x\ldotp A[x])\lor C)\implies\forall x\ldotp(A[x]\lor C)$
  by $\lor$-elimination(9, 4)\qedhere
\end{enumerate}
\end{proof}

\nwenddocs{}\nwbegincode{85}\sublabel{NW1NjfzY-uZfts-2}\nwmargintag{{\nwtagstyle{}\subpageref{NW1NjfzY-uZfts-2}}}\moddef{Derived results for Intuitionistic predicate logic~{\nwtagstyle{}\subpageref{NW1NjfzY-uZfts-1}}}\plusendmoddef\nwstartdeflinemarkup\nwusesondefline{\\{NW1NjfzY-4dPl0p-1}}\nwprevnextdefs{NW1NjfzY-uZfts-1}{NW1NjfzY-uZfts-3}\nwenddeflinemarkup
(* all_disjL : \nwlinkedidentc{Term.t}{NW1rTmOJ-4LsRWC-1}. -> \nwlinkedidentc{Formula.t}{NW3w0iRO-2jnj9W-1} -> \nwlinkedidentc{Formula.t}{NW3w0iRO-2jnj9W-1} -> \nwlinkedidentc{Thm.t}{NW4ICsJQ-YIPBT-1} *)
fun all_disjL x A C =
  let
    val th1 = \nwlinkedidentc{Thm.assume}{NW4ICsJQ-2zKkPU-2} C [];
    val th2 = \nwlinkedidentc{Thm.or_intro_l}{NW4ICsJQ-2zKkPU-4} A th1;
    val th3 = \nwlinkedidentc{Thm.forall_intro}{NW4ICsJQ-2zKkPU-8} x th2;
    val th4 = \nwlinkedidentc{Thm.disch}{NW4ICsJQ-2zKkPU-2} C th3;
    val th5 = \nwlinkedidentc{Thm.assume}{NW4ICsJQ-2zKkPU-2} (\nwlinkedidentc{Formula.mk_forall}{NW3w0iRO-3SlqND-4}(x, A)) [];
    val th6 = \nwlinkedidentc{Thm.forall_elim}{NW4ICsJQ-2zKkPU-9} x th5;
    val th7 = \nwlinkedidentc{Thm.or_intro_r}{NW4ICsJQ-2zKkPU-4} C th6;
    val th8 = \nwlinkedidentc{Thm.forall_intro}{NW4ICsJQ-2zKkPU-8} x th7;
    val th9 = \nwlinkedidentc{Thm.disch}{NW4ICsJQ-2zKkPU-2} (\nwlinkedidentc{Formula.mk_forall}{NW3w0iRO-3SlqND-4}(x, A)) th8;
  in
    \nwlinkedidentc{Thm.or_elim}{NW4ICsJQ-2zKkPU-4} th9 th4
  end;
\nwindexdefn{\nwixident{Derived.all{\_}disjL}}{Derived.all:undisjL}{NW1NjfzY-uZfts-2}\eatline
\nwused{\\{NW1NjfzY-4dPl0p-1}}\nwidentdefs{\\{{\nwixident{Derived.all{\_}disjL}}{Derived.all:undisjL}}}\nwidentuses{\\{{\nwixident{Formula.mk{\_}forall}}{Formula.mk:unforall}}\\{{\nwixident{Formula.t}}{Formula.t}}\\{{\nwixident{Term.t}}{Term.t}}\\{{\nwixident{Thm.assume}}{Thm.assume}}\\{{\nwixident{Thm.disch}}{Thm.disch}}\\{{\nwixident{Thm.forall{\_}elim}}{Thm.forall:unelim}}\\{{\nwixident{Thm.forall{\_}intro}}{Thm.forall:unintro}}\\{{\nwixident{Thm.or{\_}elim}}{Thm.or:unelim}}\\{{\nwixident{Thm.or{\_}intro{\_}l}}{Thm.or:unintro:unl}}\\{{\nwixident{Thm.or{\_}intro{\_}r}}{Thm.or:unintro:unr}}\\{{\nwixident{Thm.t}}{Thm.t}}}\nwindexuse{\nwixident{Formula.mk{\_}forall}}{Formula.mk:unforall}{NW1NjfzY-uZfts-2}\nwindexuse{\nwixident{Formula.t}}{Formula.t}{NW1NjfzY-uZfts-2}\nwindexuse{\nwixident{Term.t}}{Term.t}{NW1NjfzY-uZfts-2}\nwindexuse{\nwixident{Thm.assume}}{Thm.assume}{NW1NjfzY-uZfts-2}\nwindexuse{\nwixident{Thm.disch}}{Thm.disch}{NW1NjfzY-uZfts-2}\nwindexuse{\nwixident{Thm.forall{\_}elim}}{Thm.forall:unelim}{NW1NjfzY-uZfts-2}\nwindexuse{\nwixident{Thm.forall{\_}intro}}{Thm.forall:unintro}{NW1NjfzY-uZfts-2}\nwindexuse{\nwixident{Thm.or{\_}elim}}{Thm.or:unelim}{NW1NjfzY-uZfts-2}\nwindexuse{\nwixident{Thm.or{\_}intro{\_}l}}{Thm.or:unintro:unl}{NW1NjfzY-uZfts-2}\nwindexuse{\nwixident{Thm.or{\_}intro{\_}r}}{Thm.or:unintro:unr}{NW1NjfzY-uZfts-2}\nwindexuse{\nwixident{Thm.t}}{Thm.t}{NW1NjfzY-uZfts-2}\nwendcode{}\nwbegindocs{86}\nwdocspar
We also add a test to check that the result works.

\nwenddocs{}\nwbegincode{87}\sublabel{NW1NjfzY-3rvgpx-1}\nwmargintag{{\nwtagstyle{}\subpageref{NW1NjfzY-3rvgpx-1}}}\moddef{Unit tests for \code{}Derived.sml\edoc{}~{\nwtagstyle{}\subpageref{NW1NjfzY-3rvgpx-1}}}\endmoddef\nwstartdeflinemarkup\nwprevnextdefs{\relax}{NW1NjfzY-3rvgpx-2}\nwenddeflinemarkup
fun all_disjL_test () =
  let
    val x = Term.Var("x", \nwlinkedidentc{Sort.IND}{NW1rTmOJ-3ZqggC-1});
    val A = \nwlinkedidentc{Formula.mk_pred}{NW3w0iRO-3SlqND-1}("A", [x]);
    val C = \nwlinkedidentc{Formula.mk_pred}{NW3w0iRO-3SlqND-1}("C", []);
    val expected = \nwlinkedidentc{Formula.mk_imp}{NW3w0iRO-3SlqND-3}(\nwlinkedidentc{Formula.mk_or}{NW3w0iRO-3SlqND-3}(\nwlinkedidentc{Formula.mk_forall}{NW3w0iRO-3SlqND-4}(x,A),C),
                                  \nwlinkedidentc{Formula.mk_forall}{NW3w0iRO-3SlqND-4}(x, \nwlinkedidentc{Formula.mk_or}{NW3w0iRO-3SlqND-3}(A,C)));
    val actual = \nwlinkedidentc{Thm.concl}{NW4ICsJQ-YIPBT-1} (\nwlinkedidentc{Derived.all_disjL}{NW1NjfzY-uZfts-2} x A C);
  in
    \nwlinkedidentc{Formula.eq}{NW3w0iRO-5gQzE-1} expected actual
  end;
\nwalsodefined{\\{NW1NjfzY-3rvgpx-2}\\{NW1NjfzY-3rvgpx-3}\\{NW1NjfzY-3rvgpx-4}\\{NW1NjfzY-3rvgpx-5}\\{NW1NjfzY-3rvgpx-6}\\{NW1NjfzY-3rvgpx-7}\\{NW1NjfzY-3rvgpx-8}\\{NW1NjfzY-3rvgpx-9}\\{NW1NjfzY-3rvgpx-A}}\nwnotused{Unit tests for [[Derived.sml]]}\nwidentuses{\\{{\nwixident{Derived.all{\_}disjL}}{Derived.all:undisjL}}\\{{\nwixident{Formula.eq}}{Formula.eq}}\\{{\nwixident{Formula.mk{\_}forall}}{Formula.mk:unforall}}\\{{\nwixident{Formula.mk{\_}imp}}{Formula.mk:unimp}}\\{{\nwixident{Formula.mk{\_}or}}{Formula.mk:unor}}\\{{\nwixident{Formula.mk{\_}pred}}{Formula.mk:unpred}}\\{{\nwixident{Sort.IND}}{Sort.IND}}\\{{\nwixident{Thm.concl}}{Thm.concl}}}\nwindexuse{\nwixident{Derived.all{\_}disjL}}{Derived.all:undisjL}{NW1NjfzY-3rvgpx-1}\nwindexuse{\nwixident{Formula.eq}}{Formula.eq}{NW1NjfzY-3rvgpx-1}\nwindexuse{\nwixident{Formula.mk{\_}forall}}{Formula.mk:unforall}{NW1NjfzY-3rvgpx-1}\nwindexuse{\nwixident{Formula.mk{\_}imp}}{Formula.mk:unimp}{NW1NjfzY-3rvgpx-1}\nwindexuse{\nwixident{Formula.mk{\_}or}}{Formula.mk:unor}{NW1NjfzY-3rvgpx-1}\nwindexuse{\nwixident{Formula.mk{\_}pred}}{Formula.mk:unpred}{NW1NjfzY-3rvgpx-1}\nwindexuse{\nwixident{Sort.IND}}{Sort.IND}{NW1NjfzY-3rvgpx-1}\nwindexuse{\nwixident{Thm.concl}}{Thm.concl}{NW1NjfzY-3rvgpx-1}\nwendcode{}\nwbegindocs{88}\nwdocspar

\begin{node}%ex_disj_distrib
$\vdash((\exists x\ldotp A[x])\lor(\exists x\ldotp B[x]))\iff\exists x\ldotp(A[x]\lor B[x])$
\end{node}

\begin{proof}
The forward proof:
\begin{enumerate}
\item $A[x]\vdash A[x]$ by assuming $A[x]$%1
\item $A[x]\vdash A[x]\lor B[x]$ by $\lor$-intro-r($B[x]$)%2
\item $A[x]\vdash\exists x\ldotp(A[x]\lor B[x])$ by $\exists$-intro(2, $x$)%3
\item $\vdash A[x]\implies\exists x\ldotp(A[x]\lor B[x])$ by
  discharging $A[x]$%4
\item $\vdash(\exists x\ldotp A[x])\implies\exists x\ldotp(A[x]\lor B[x])$ by
  $\exists$-elimination(4)%5
\item $B[x]\vdash B[x]$ by assuming $B[x]$%6
\item $B[x]\vdash A[x]\lor B[x]$ by $\lor$-intro-l(6, $A[x]$)%7
\item $B[x]\vdash\exists x\ldotp(A[x]\lor B[x])$ by $\exists$-intro(7, $x$)%8
\item $\vdash B[x]\implies\exists x\ldotp(A[x]\lor B[x])$ by
  discharging $B[x]$%9
\item $\vdash(\exists x\ldotp B[x])\implies\exists x\ldotp(A[x]\lor B[x])$ by
  $\exists$-elimination(9)%10
\item $\vdash(\exists x\ldotp A[x])\lor(\exists x\ldotp B[x])\implies\exists x\ldotp(A[x]\lor B[x])$
  by $\lor$-intro(5, 10)%11
\end{enumerate}
The backward proof:
\begin{enumerate}\setcounter{enumi}{11}
\item $A[x]\vdash \exists x\ldotp A[x]$ by $\exists$-intro(1)%12
\item $A[x]\vdash(\exists x\ldotp A[x])\lor(\exists x\ldotp B[x])$ by $\lor$-intro-r($\exists x\ldotp B[x]$)%13
\item $\vdash A[x]\implies((\exists x\ldotp A[x])\lor(\exists x\ldotp B[x]))$
  by discharging $A[x]$%14
\item $B[x]\vdash\exists x\ldotp B[x]$ by $\exists$-intro(6)%15
\item $B[x]\vdash(\exists x\ldotp A[x])\lor(\exists x\ldotp B[x])$ by $\lor$-intro-l($\exists x\ldotp A[x]$)%16
\item $\vdash B[x]\implies((\exists x\ldotp A[x])\lor(\exists x\ldotp B[x]))$
  by discharging $B[x]$ from 16%17
\item $\vdash (A[x]\lor B[x])\implies((\exists x\ldotp A[x])\lor(\exists x\ldotp B[x]))$
  by $\lor$-elimination(14, 17)%18
\item $\vdash\exists x\ldotp(A[x]\lor B[x])\implies((\exists x\ldotp A[x])\lor(\exists x\ldotp B[x]))$
  by $\exists$-elimination(18)\qedhere
\end{enumerate}
\end{proof}

\nwenddocs{}\nwbegincode{89}\sublabel{NW1NjfzY-uZfts-3}\nwmargintag{{\nwtagstyle{}\subpageref{NW1NjfzY-uZfts-3}}}\moddef{Derived results for Intuitionistic predicate logic~{\nwtagstyle{}\subpageref{NW1NjfzY-uZfts-1}}}\plusendmoddef\nwstartdeflinemarkup\nwusesondefline{\\{NW1NjfzY-4dPl0p-1}}\nwprevnextdefs{NW1NjfzY-uZfts-2}{NW1NjfzY-uZfts-4}\nwenddeflinemarkup
local
  fun forward x A B =
    let
      val th1 = \nwlinkedidentc{Thm.assume}{NW4ICsJQ-2zKkPU-2} A [];
      val th2 = \nwlinkedidentc{Thm.or_intro_r}{NW4ICsJQ-2zKkPU-4} B th1;
      val th3 = \nwlinkedidentc{Thm.exists_intro}{NW4ICsJQ-2zKkPU-A} x x (\nwlinkedidentc{Formula.mk_or}{NW3w0iRO-3SlqND-3}(A, B))
                                 th2;
      val th4 = \nwlinkedidentc{Thm.disch}{NW4ICsJQ-2zKkPU-2} A th3;
      val th5 = \nwlinkedidentc{Thm.exists_elim}{NW4ICsJQ-2zKkPU-B} x th4;
      val th6 = \nwlinkedidentc{Thm.assume}{NW4ICsJQ-2zKkPU-2} B [];
      val th7 = \nwlinkedidentc{Thm.or_intro_l}{NW4ICsJQ-2zKkPU-4} A th6;
      val th8 = \nwlinkedidentc{Thm.exists_intro}{NW4ICsJQ-2zKkPU-A} x x (\nwlinkedidentc{Formula.mk_or}{NW3w0iRO-3SlqND-3}(A, B)) th7;
      val th9 = \nwlinkedidentc{Thm.disch}{NW4ICsJQ-2zKkPU-2} B th8;
      val th10 = \nwlinkedidentc{Thm.exists_elim}{NW4ICsJQ-2zKkPU-B} x th9;
    in
      \nwlinkedidentc{Thm.or_elim}{NW4ICsJQ-2zKkPU-4} th5 th10
    end;
  fun backward x A B =
    let
      val th1 = \nwlinkedidentc{Thm.assume}{NW4ICsJQ-2zKkPU-2} A [];
      val th2 = \nwlinkedidentc{Thm.exists_intro}{NW4ICsJQ-2zKkPU-A} x x A th1;
      val th3 = \nwlinkedidentc{Thm.or_intro_r}{NW4ICsJQ-2zKkPU-4} (\nwlinkedidentc{Formula.mk_exists}{NW3w0iRO-3SlqND-4}(x, B)) th2;
      val th4 = \nwlinkedidentc{Thm.disch}{NW4ICsJQ-2zKkPU-2} A th3;
      val th5 = \nwlinkedidentc{Thm.assume}{NW4ICsJQ-2zKkPU-2} B [];
      val th6 = \nwlinkedidentc{Thm.exists_intro}{NW4ICsJQ-2zKkPU-A} x x B th5;
      val th7 = \nwlinkedidentc{Thm.or_intro_l}{NW4ICsJQ-2zKkPU-4} (\nwlinkedidentc{Formula.mk_exists}{NW3w0iRO-3SlqND-4}(x, A)) th6;
      val th8 = \nwlinkedidentc{Thm.disch}{NW4ICsJQ-2zKkPU-2} B th7;
      val th9 = \nwlinkedidentc{Thm.or_elim}{NW4ICsJQ-2zKkPU-4} th4 th8;
    in
      \nwlinkedidentc{Thm.exists_elim}{NW4ICsJQ-2zKkPU-B} x th9
    end;
in
fun ex_disj_distrib x A B =
  \nwlinkedidentc{Thm.iff_intro}{NW4ICsJQ-2zKkPU-7} (forward x A B) (backward x A B)
end;
\nwindexdefn{\nwixident{Derived.ex{\_}disj{\_}distrib}}{Derived.ex:undisj:undistrib}{NW1NjfzY-uZfts-3}\eatline
\nwused{\\{NW1NjfzY-4dPl0p-1}}\nwidentdefs{\\{{\nwixident{Derived.ex{\_}disj{\_}distrib}}{Derived.ex:undisj:undistrib}}}\nwidentuses{\\{{\nwixident{Formula.mk{\_}exists}}{Formula.mk:unexists}}\\{{\nwixident{Formula.mk{\_}or}}{Formula.mk:unor}}\\{{\nwixident{Thm.assume}}{Thm.assume}}\\{{\nwixident{Thm.disch}}{Thm.disch}}\\{{\nwixident{Thm.exists{\_}elim}}{Thm.exists:unelim}}\\{{\nwixident{Thm.exists{\_}intro}}{Thm.exists:unintro}}\\{{\nwixident{Thm.iff{\_}intro}}{Thm.iff:unintro}}\\{{\nwixident{Thm.or{\_}elim}}{Thm.or:unelim}}\\{{\nwixident{Thm.or{\_}intro{\_}l}}{Thm.or:unintro:unl}}\\{{\nwixident{Thm.or{\_}intro{\_}r}}{Thm.or:unintro:unr}}}\nwindexuse{\nwixident{Formula.mk{\_}exists}}{Formula.mk:unexists}{NW1NjfzY-uZfts-3}\nwindexuse{\nwixident{Formula.mk{\_}or}}{Formula.mk:unor}{NW1NjfzY-uZfts-3}\nwindexuse{\nwixident{Thm.assume}}{Thm.assume}{NW1NjfzY-uZfts-3}\nwindexuse{\nwixident{Thm.disch}}{Thm.disch}{NW1NjfzY-uZfts-3}\nwindexuse{\nwixident{Thm.exists{\_}elim}}{Thm.exists:unelim}{NW1NjfzY-uZfts-3}\nwindexuse{\nwixident{Thm.exists{\_}intro}}{Thm.exists:unintro}{NW1NjfzY-uZfts-3}\nwindexuse{\nwixident{Thm.iff{\_}intro}}{Thm.iff:unintro}{NW1NjfzY-uZfts-3}\nwindexuse{\nwixident{Thm.or{\_}elim}}{Thm.or:unelim}{NW1NjfzY-uZfts-3}\nwindexuse{\nwixident{Thm.or{\_}intro{\_}l}}{Thm.or:unintro:unl}{NW1NjfzY-uZfts-3}\nwindexuse{\nwixident{Thm.or{\_}intro{\_}r}}{Thm.or:unintro:unr}{NW1NjfzY-uZfts-3}\nwendcode{}\nwbegindocs{90}\nwdocspar
\nwenddocs{}\nwbegincode{91}\sublabel{NW1NjfzY-3rvgpx-2}\nwmargintag{{\nwtagstyle{}\subpageref{NW1NjfzY-3rvgpx-2}}}\moddef{Unit tests for \code{}Derived.sml\edoc{}~{\nwtagstyle{}\subpageref{NW1NjfzY-3rvgpx-1}}}\plusendmoddef\nwstartdeflinemarkup\nwprevnextdefs{NW1NjfzY-3rvgpx-1}{NW1NjfzY-3rvgpx-3}\nwenddeflinemarkup
fun ex_disj_distrib_test () =
  let
    val x = Term.Var("x", \nwlinkedidentc{Sort.IND}{NW1rTmOJ-3ZqggC-1});
    val A = \nwlinkedidentc{Formula.mk_pred}{NW3w0iRO-3SlqND-1}("A", [x]);
    val B = \nwlinkedidentc{Formula.mk_pred}{NW3w0iRO-3SlqND-1}("B", [x]);
    val expected =
      \nwlinkedidentc{Formula.mk_iff}{NW3w0iRO-3SlqND-3}(\nwlinkedidentc{Formula.mk_or}{NW3w0iRO-3SlqND-3}(\nwlinkedidentc{Formula.mk_exists}{NW3w0iRO-3SlqND-4}(x,A),
                                   \nwlinkedidentc{Formula.mk_exists}{NW3w0iRO-3SlqND-4}(x,B)),
                     \nwlinkedidentc{Formula.mk_exists}{NW3w0iRO-3SlqND-4}(x, \nwlinkedidentc{Formula.mk_or}{NW3w0iRO-3SlqND-3}(A,B)));
    val actual = \nwlinkedidentc{Thm.concl}{NW4ICsJQ-YIPBT-1} (\nwlinkedidentc{Derived.ex_disj_distrib}{NW1NjfzY-uZfts-3} x A B);
  in
    \nwlinkedidentc{Formula.eq}{NW3w0iRO-5gQzE-1} expected actual
  end;
\nwidentuses{\\{{\nwixident{Derived.ex{\_}disj{\_}distrib}}{Derived.ex:undisj:undistrib}}\\{{\nwixident{Formula.eq}}{Formula.eq}}\\{{\nwixident{Formula.mk{\_}exists}}{Formula.mk:unexists}}\\{{\nwixident{Formula.mk{\_}iff}}{Formula.mk:uniff}}\\{{\nwixident{Formula.mk{\_}or}}{Formula.mk:unor}}\\{{\nwixident{Formula.mk{\_}pred}}{Formula.mk:unpred}}\\{{\nwixident{Sort.IND}}{Sort.IND}}\\{{\nwixident{Thm.concl}}{Thm.concl}}}\nwindexuse{\nwixident{Derived.ex{\_}disj{\_}distrib}}{Derived.ex:undisj:undistrib}{NW1NjfzY-3rvgpx-2}\nwindexuse{\nwixident{Formula.eq}}{Formula.eq}{NW1NjfzY-3rvgpx-2}\nwindexuse{\nwixident{Formula.mk{\_}exists}}{Formula.mk:unexists}{NW1NjfzY-3rvgpx-2}\nwindexuse{\nwixident{Formula.mk{\_}iff}}{Formula.mk:uniff}{NW1NjfzY-3rvgpx-2}\nwindexuse{\nwixident{Formula.mk{\_}or}}{Formula.mk:unor}{NW1NjfzY-3rvgpx-2}\nwindexuse{\nwixident{Formula.mk{\_}pred}}{Formula.mk:unpred}{NW1NjfzY-3rvgpx-2}\nwindexuse{\nwixident{Sort.IND}}{Sort.IND}{NW1NjfzY-3rvgpx-2}\nwindexuse{\nwixident{Thm.concl}}{Thm.concl}{NW1NjfzY-3rvgpx-2}\nwendcode{}\nwbegindocs{92}\nwdocspar

\begin{node}%all_conj_distrib
$\vdash(\forall x\ldotp A[x]\land B[x])\iff((\forall x\ldotp A[x])\land(\forall x\ldotp B[x]))$
\end{node}

\begin{proof}
Forward proof:
\begin{enumerate}
\item $\forall x\ldotp(A[x]\land B[x])\vdash\forall x\ldotp(A[x]\land B[x])$
  by assuming $\forall x\ldotp(A[x]\land B[x])$%1
\item $\forall x\ldotp(A[x]\land B[x])\vdash A[x]\land B[x]$ by $\forall$-elim%2
\item $\forall x\ldotp(A[x]\land B[x])\vdash A[x]$ by $\land$-elim-l(2)%3
\item $\forall x\ldotp(A[x]\land B[x])\vdash\forall x\ldotp A[x]$ by $\forall$-intro(3)%4
\item $\forall x\ldotp(A[x]\land B[x])\vdash B[x]$ by $\land$-elim-r(2)%5
\item $\forall x\ldotp(A[x]\land B[x])\vdash\forall x\ldotp B[x]$ by $\forall$-intro(5)%6
\item $\forall x\ldotp(A[x]\land B[x])\vdash(\forall x\ldotp A[x])\land(\forall x\ldotp B[x])$
  by $\land$-intro(4, 6)%7
\item $\vdash(\forall x\ldotp(A[x]\land B[x]))\implies((\forall x\ldotp A[x])\land(\forall x\ldotp B[x]))$
  by discharging (7)%8
\end{enumerate}
Backward proof
\begin{enumerate}\setcounter{enumi}{8}
\item $(\forall x\ldotp A[x])\land(\forall x\ldotp B[x])\vdash(\forall x\ldotp A[x])\land(\forall x\ldotp B[x])$
  by assuming $(\forall x\ldotp A[x])\land(\forall x\ldotp B[x])$%9
\item $(\forall x\ldotp A[x])\land(\forall x\ldotp B[x])\vdash(\forall x\ldotp A[x])$
  by $\land$-elim-l(9)%10
\item $(\forall x\ldotp A[x])\land(\forall x\ldotp B[x])\vdash A[x]$
  by $\forall$-elim(10)%11
\item $(\forall x\ldotp A[x])\land(\forall x\ldotp B[x])\vdash\forall x\ldotp B[x]$
  by $\land$-elim-r(9)%12
\item $(\forall x\ldotp A[x])\land(\forall x\ldotp B[x])\vdash B[x]$
  by $\forall$-elim(12)%13
\item $(\forall x\ldotp A[x])\land(\forall x\ldotp B[x])\vdash
  A[x]\land B[x]$ by $\land$-intro(11, 13)% 14
\item $(\forall x\ldotp A[x])\land(\forall x\ldotp B[x])\vdash\forall
  x\ldotp(A[x]\land B[x])$ by $\forall$-intro(14)% 15
\item $\vdash((\forall x\ldotp A[x])\land(\forall x\ldotp B[x]))\implies\forall
  x\ldotp(A[x]\land B[x])$ by discharging 15\qedhere
\end{enumerate}
\end{proof}

\nwenddocs{}\nwbegincode{93}\sublabel{NW1NjfzY-uZfts-4}\nwmargintag{{\nwtagstyle{}\subpageref{NW1NjfzY-uZfts-4}}}\moddef{Derived results for Intuitionistic predicate logic~{\nwtagstyle{}\subpageref{NW1NjfzY-uZfts-1}}}\plusendmoddef\nwstartdeflinemarkup\nwusesondefline{\\{NW1NjfzY-4dPl0p-1}}\nwprevnextdefs{NW1NjfzY-uZfts-3}{NW1NjfzY-uZfts-5}\nwenddeflinemarkup
fun all_conj_distrib x A B =
  let
    (* backward proof *)
    val th1 = \nwlinkedidentc{Thm.assume}{NW4ICsJQ-2zKkPU-2} (\nwlinkedidentc{Formula.mk_forall}{NW3w0iRO-3SlqND-4}(x, \nwlinkedidentc{Formula.mk_and}{NW3w0iRO-3SlqND-3}(A, B))) [];
    val th2 = \nwlinkedidentc{Thm.forall_elim}{NW4ICsJQ-2zKkPU-9} x th1;
    val th3 = Thm.and_elim_l th2; (* ... |- A[x] *)
    val th4 = \nwlinkedidentc{Thm.forall_intro}{NW4ICsJQ-2zKkPU-8} x th3;
    val th5 = Thm.and_elim_r th2; (* ... |- B[x] *)
    val th6 = \nwlinkedidentc{Thm.forall_intro}{NW4ICsJQ-2zKkPU-8} x th5;
    val th7 = Thm.and_intro th4 th6;
    val th8 = \nwlinkedidentc{Thm.disch}{NW4ICsJQ-2zKkPU-2} (\nwlinkedidentc{Formula.mk_forall}{NW3w0iRO-3SlqND-4}(x, \nwlinkedidentc{Formula.mk_and}{NW3w0iRO-3SlqND-3}(A, B))) th7;
    (* forward proof *)
    val th9 = \nwlinkedidentc{Thm.assume}{NW4ICsJQ-2zKkPU-2} (\nwlinkedidentc{Formula.mk_and}{NW3w0iRO-3SlqND-3}(\nwlinkedidentc{Formula.mk_forall}{NW3w0iRO-3SlqND-4}(x, A),
                                         \nwlinkedidentc{Formula.mk_forall}{NW3w0iRO-3SlqND-4}(x, B)))
                         [];
    val th10 = Thm.and_elim_l th9;
    val th11 = \nwlinkedidentc{Thm.forall_elim}{NW4ICsJQ-2zKkPU-9} x th10;
    val th12 = Thm.and_elim_r th9;
    val th13 = \nwlinkedidentc{Thm.forall_elim}{NW4ICsJQ-2zKkPU-9} x th12;
    val th14 = Thm.and_intro th11 th13;
    val th15 = \nwlinkedidentc{Thm.forall_intro}{NW4ICsJQ-2zKkPU-8} x th14;
    val th16 = \nwlinkedidentc{Thm.disch}{NW4ICsJQ-2zKkPU-2} (\nwlinkedidentc{Formula.mk_and}{NW3w0iRO-3SlqND-3}(\nwlinkedidentc{Formula.mk_forall}{NW3w0iRO-3SlqND-4}(x, A),
                                         \nwlinkedidentc{Formula.mk_forall}{NW3w0iRO-3SlqND-4}(x, B)))
                         th15;
  in
    \nwlinkedidentc{Thm.iff_intro}{NW4ICsJQ-2zKkPU-7} th16 th8
  end;
\nwindexdefn{\nwixident{Derived.all{\_}conj{\_}distrib}}{Derived.all:unconj:undistrib}{NW1NjfzY-uZfts-4}\eatline
\nwused{\\{NW1NjfzY-4dPl0p-1}}\nwidentdefs{\\{{\nwixident{Derived.all{\_}conj{\_}distrib}}{Derived.all:unconj:undistrib}}}\nwidentuses{\\{{\nwixident{Formula.mk{\_}and}}{Formula.mk:unand}}\\{{\nwixident{Formula.mk{\_}forall}}{Formula.mk:unforall}}\\{{\nwixident{Thm.assume}}{Thm.assume}}\\{{\nwixident{Thm.disch}}{Thm.disch}}\\{{\nwixident{Thm.forall{\_}elim}}{Thm.forall:unelim}}\\{{\nwixident{Thm.forall{\_}intro}}{Thm.forall:unintro}}\\{{\nwixident{Thm.iff{\_}intro}}{Thm.iff:unintro}}}\nwindexuse{\nwixident{Formula.mk{\_}and}}{Formula.mk:unand}{NW1NjfzY-uZfts-4}\nwindexuse{\nwixident{Formula.mk{\_}forall}}{Formula.mk:unforall}{NW1NjfzY-uZfts-4}\nwindexuse{\nwixident{Thm.assume}}{Thm.assume}{NW1NjfzY-uZfts-4}\nwindexuse{\nwixident{Thm.disch}}{Thm.disch}{NW1NjfzY-uZfts-4}\nwindexuse{\nwixident{Thm.forall{\_}elim}}{Thm.forall:unelim}{NW1NjfzY-uZfts-4}\nwindexuse{\nwixident{Thm.forall{\_}intro}}{Thm.forall:unintro}{NW1NjfzY-uZfts-4}\nwindexuse{\nwixident{Thm.iff{\_}intro}}{Thm.iff:unintro}{NW1NjfzY-uZfts-4}\nwendcode{}\nwbegindocs{94}\nwdocspar
\nwenddocs{}\nwbegincode{95}\sublabel{NW1NjfzY-3rvgpx-3}\nwmargintag{{\nwtagstyle{}\subpageref{NW1NjfzY-3rvgpx-3}}}\moddef{Unit tests for \code{}Derived.sml\edoc{}~{\nwtagstyle{}\subpageref{NW1NjfzY-3rvgpx-1}}}\plusendmoddef\nwstartdeflinemarkup\nwprevnextdefs{NW1NjfzY-3rvgpx-2}{NW1NjfzY-3rvgpx-4}\nwenddeflinemarkup
fun all_conj_distrib_test () =
  let
    val x = Term.Var("x", \nwlinkedidentc{Sort.IND}{NW1rTmOJ-3ZqggC-1});
    val A = \nwlinkedidentc{Formula.mk_pred}{NW3w0iRO-3SlqND-1}("A", [x]);
    val B = \nwlinkedidentc{Formula.mk_pred}{NW3w0iRO-3SlqND-1}("B", [x]);
    val expected =
      \nwlinkedidentc{Formula.mk_iff}{NW3w0iRO-3SlqND-3}(\nwlinkedidentc{Formula.mk_and}{NW3w0iRO-3SlqND-3}(\nwlinkedidentc{Formula.mk_forall}{NW3w0iRO-3SlqND-4}(x,A),
                                    \nwlinkedidentc{Formula.mk_forall}{NW3w0iRO-3SlqND-4}(x,B)),
                     \nwlinkedidentc{Formula.mk_forall}{NW3w0iRO-3SlqND-4}(x, \nwlinkedidentc{Formula.mk_and}{NW3w0iRO-3SlqND-3}(A,B)));
    val actual = \nwlinkedidentc{Thm.concl}{NW4ICsJQ-YIPBT-1} (\nwlinkedidentc{Derived.all_conj_distrib}{NW1NjfzY-uZfts-4} x A B);
  in
    if not(\nwlinkedidentc{Formula.eq}{NW3w0iRO-5gQzE-1} expected actual)
    then "Expecting: "^(\nwlinkedidentc{Formula.serialize}{NW3w0iRO-3H9tUA-1} expected)^"\\nActual: "^
         (\nwlinkedidentc{Formula.serialize}{NW3w0iRO-3H9tUA-1} actual)^"\\n"
    else ""
  end;
\nwidentuses{\\{{\nwixident{Derived.all{\_}conj{\_}distrib}}{Derived.all:unconj:undistrib}}\\{{\nwixident{Formula.eq}}{Formula.eq}}\\{{\nwixident{Formula.mk{\_}and}}{Formula.mk:unand}}\\{{\nwixident{Formula.mk{\_}forall}}{Formula.mk:unforall}}\\{{\nwixident{Formula.mk{\_}iff}}{Formula.mk:uniff}}\\{{\nwixident{Formula.mk{\_}pred}}{Formula.mk:unpred}}\\{{\nwixident{Formula.serialize}}{Formula.serialize}}\\{{\nwixident{Sort.IND}}{Sort.IND}}\\{{\nwixident{Thm.concl}}{Thm.concl}}}\nwindexuse{\nwixident{Derived.all{\_}conj{\_}distrib}}{Derived.all:unconj:undistrib}{NW1NjfzY-3rvgpx-3}\nwindexuse{\nwixident{Formula.eq}}{Formula.eq}{NW1NjfzY-3rvgpx-3}\nwindexuse{\nwixident{Formula.mk{\_}and}}{Formula.mk:unand}{NW1NjfzY-3rvgpx-3}\nwindexuse{\nwixident{Formula.mk{\_}forall}}{Formula.mk:unforall}{NW1NjfzY-3rvgpx-3}\nwindexuse{\nwixident{Formula.mk{\_}iff}}{Formula.mk:uniff}{NW1NjfzY-3rvgpx-3}\nwindexuse{\nwixident{Formula.mk{\_}pred}}{Formula.mk:unpred}{NW1NjfzY-3rvgpx-3}\nwindexuse{\nwixident{Formula.serialize}}{Formula.serialize}{NW1NjfzY-3rvgpx-3}\nwindexuse{\nwixident{Sort.IND}}{Sort.IND}{NW1NjfzY-3rvgpx-3}\nwindexuse{\nwixident{Thm.concl}}{Thm.concl}{NW1NjfzY-3rvgpx-3}\nwendcode{}\nwbegindocs{96}\nwdocspar

\begin{node}%ex_conj_distrib
$\vdash(\exists x\ldotp A[x]\land B[x])\implies((\exists x\ldotp A[x])\land(\exists x\ldotp B[x]))$
\end{node}

\begin{proof}
\begin{enumerate}
\item $A[x]\land B[x]\vdash A[x]\land B[x]$ 
  by assuming $A[x]\land B[x]$%1
\item $A[x]\land B[x]\vdash A[x]$ by $\land$-elim-l(1)%2
\item $A[x]\land B[x]\vdash \exists x\ldotp A[x]$ by $\exists$-intro(2, $x$)%3
\item $\vdash A[x]\land B[x]\implies\exists x\ldotp A[x]$ by
  discharging 3%4
\item $A[x]\land B[x]\vdash B[x]$ by $\land$-elim-r(1)%5
\item $A[x]\land B[x]\vdash \exists x\ldotp  B[x]$ by $\exists$-intro(5)%6
\item $\vdash A[x]\land B[x]\implies \exists x\ldotp  B[x]$ by
  discharging 6%7
\item $\vdash A[x]\land B[x]\implies(\exists x\ldotp A[x])\land(\exists x\ldotp B[x])$
  by $\land$-intro(4, 7)%8
\item $\vdash(\exists x\ldotp A[x]\land B[x])\implies(\exists x\ldotp A[x])\land(\exists x\ldotp B[x])$
  by $\exists$-elim(8)\qedhere
\end{enumerate}
\end{proof}

\nwenddocs{}\nwbegincode{97}\sublabel{NW1NjfzY-uZfts-5}\nwmargintag{{\nwtagstyle{}\subpageref{NW1NjfzY-uZfts-5}}}\moddef{Derived results for Intuitionistic predicate logic~{\nwtagstyle{}\subpageref{NW1NjfzY-uZfts-1}}}\plusendmoddef\nwstartdeflinemarkup\nwusesondefline{\\{NW1NjfzY-4dPl0p-1}}\nwprevnextdefs{NW1NjfzY-uZfts-4}{NW1NjfzY-uZfts-6}\nwenddeflinemarkup
fun ex_conj_distrib x A B =
  let
    val AB = \nwlinkedidentc{Formula.mk_and}{NW3w0iRO-3SlqND-3}(A, B);
    val th1 = \nwlinkedidentc{Thm.assume}{NW4ICsJQ-2zKkPU-2} AB [];
    val th2 = Thm.and_elim_l th1;
    val th3 = \nwlinkedidentc{Thm.exists_intro}{NW4ICsJQ-2zKkPU-A} x x A th2;
    val th4 = Thm.and_elim_r th1;
    val th5 = \nwlinkedidentc{Thm.exists_intro}{NW4ICsJQ-2zKkPU-A} x x B th4;
    val th6 = Thm.and_intro th3 th5;
    val th7 = \nwlinkedidentc{Thm.disch}{NW4ICsJQ-2zKkPU-2} AB th6;
  in
    \nwlinkedidentc{Thm.exists_elim}{NW4ICsJQ-2zKkPU-B} x th7
  end;
\nwindexdefn{\nwixident{Derived.ex{\_}conj{\_}distrib}}{Derived.ex:unconj:undistrib}{NW1NjfzY-uZfts-5}\eatline
\nwused{\\{NW1NjfzY-4dPl0p-1}}\nwidentdefs{\\{{\nwixident{Derived.ex{\_}conj{\_}distrib}}{Derived.ex:unconj:undistrib}}}\nwidentuses{\\{{\nwixident{Formula.mk{\_}and}}{Formula.mk:unand}}\\{{\nwixident{Thm.assume}}{Thm.assume}}\\{{\nwixident{Thm.disch}}{Thm.disch}}\\{{\nwixident{Thm.exists{\_}elim}}{Thm.exists:unelim}}\\{{\nwixident{Thm.exists{\_}intro}}{Thm.exists:unintro}}}\nwindexuse{\nwixident{Formula.mk{\_}and}}{Formula.mk:unand}{NW1NjfzY-uZfts-5}\nwindexuse{\nwixident{Thm.assume}}{Thm.assume}{NW1NjfzY-uZfts-5}\nwindexuse{\nwixident{Thm.disch}}{Thm.disch}{NW1NjfzY-uZfts-5}\nwindexuse{\nwixident{Thm.exists{\_}elim}}{Thm.exists:unelim}{NW1NjfzY-uZfts-5}\nwindexuse{\nwixident{Thm.exists{\_}intro}}{Thm.exists:unintro}{NW1NjfzY-uZfts-5}\nwendcode{}\nwbegindocs{98}\nwdocspar
\nwenddocs{}\nwbegincode{99}\sublabel{NW1NjfzY-3rvgpx-4}\nwmargintag{{\nwtagstyle{}\subpageref{NW1NjfzY-3rvgpx-4}}}\moddef{Unit tests for \code{}Derived.sml\edoc{}~{\nwtagstyle{}\subpageref{NW1NjfzY-3rvgpx-1}}}\plusendmoddef\nwstartdeflinemarkup\nwprevnextdefs{NW1NjfzY-3rvgpx-3}{NW1NjfzY-3rvgpx-5}\nwenddeflinemarkup
fun ex_conj_distrib_test () =
  let
    val x = Term.Var("x", \nwlinkedidentc{Sort.IND}{NW1rTmOJ-3ZqggC-1});
    val A = \nwlinkedidentc{Formula.mk_pred}{NW3w0iRO-3SlqND-1}("A", [x]);
    val B = \nwlinkedidentc{Formula.mk_pred}{NW3w0iRO-3SlqND-1}("B", [x]);
    val expected =
      \nwlinkedidentc{Formula.mk_imp}{NW3w0iRO-3SlqND-3}(\nwlinkedidentc{Formula.mk_exists}{NW3w0iRO-3SlqND-4}(x, \nwlinkedidentc{Formula.mk_and}{NW3w0iRO-3SlqND-3}(A,B)),
                     \nwlinkedidentc{Formula.mk_and}{NW3w0iRO-3SlqND-3}(\nwlinkedidentc{Formula.mk_exists}{NW3w0iRO-3SlqND-4}(x, A),
                                    \nwlinkedidentc{Formula.mk_exists}{NW3w0iRO-3SlqND-4}(x, B)));
    val actual = \nwlinkedidentc{Thm.concl}{NW4ICsJQ-YIPBT-1} (\nwlinkedidentc{Derived.ex_conj_distrib}{NW1NjfzY-uZfts-5} x A B);
  in
    if not(\nwlinkedidentc{Formula.eq}{NW3w0iRO-5gQzE-1} expected actual)
    then "Expecting: "^(\nwlinkedidentc{Formula.serialize}{NW3w0iRO-3H9tUA-1} expected)^"\\nActual: "^
         (\nwlinkedidentc{Formula.serialize}{NW3w0iRO-3H9tUA-1} actual)^"\\n"
    else ""
  end;
\nwidentuses{\\{{\nwixident{Derived.ex{\_}conj{\_}distrib}}{Derived.ex:unconj:undistrib}}\\{{\nwixident{Formula.eq}}{Formula.eq}}\\{{\nwixident{Formula.mk{\_}and}}{Formula.mk:unand}}\\{{\nwixident{Formula.mk{\_}exists}}{Formula.mk:unexists}}\\{{\nwixident{Formula.mk{\_}imp}}{Formula.mk:unimp}}\\{{\nwixident{Formula.mk{\_}pred}}{Formula.mk:unpred}}\\{{\nwixident{Formula.serialize}}{Formula.serialize}}\\{{\nwixident{Sort.IND}}{Sort.IND}}\\{{\nwixident{Thm.concl}}{Thm.concl}}}\nwindexuse{\nwixident{Derived.ex{\_}conj{\_}distrib}}{Derived.ex:unconj:undistrib}{NW1NjfzY-3rvgpx-4}\nwindexuse{\nwixident{Formula.eq}}{Formula.eq}{NW1NjfzY-3rvgpx-4}\nwindexuse{\nwixident{Formula.mk{\_}and}}{Formula.mk:unand}{NW1NjfzY-3rvgpx-4}\nwindexuse{\nwixident{Formula.mk{\_}exists}}{Formula.mk:unexists}{NW1NjfzY-3rvgpx-4}\nwindexuse{\nwixident{Formula.mk{\_}imp}}{Formula.mk:unimp}{NW1NjfzY-3rvgpx-4}\nwindexuse{\nwixident{Formula.mk{\_}pred}}{Formula.mk:unpred}{NW1NjfzY-3rvgpx-4}\nwindexuse{\nwixident{Formula.serialize}}{Formula.serialize}{NW1NjfzY-3rvgpx-4}\nwindexuse{\nwixident{Sort.IND}}{Sort.IND}{NW1NjfzY-3rvgpx-4}\nwindexuse{\nwixident{Thm.concl}}{Thm.concl}{NW1NjfzY-3rvgpx-4}\nwendcode{}\nwbegindocs{100}\nwdocspar

\begin{node}%imp_ex
$\vdash(\forall x\ldotp(A[x]\implies C))\iff((\exists x\ldotp A[x])\implies C)$

This is more useful as two theorems:
\begin{equation}
\vdash(\forall x\ldotp(A[x]\implies C))\implies((\exists x\ldotp A[x])\implies C)
\end{equation}  
\begin{equation}
\vdash((\exists x\ldotp A[x])\implies C)\implies(\forall x\ldotp(A[x]\implies C))
\end{equation}  
\end{node}

\begin{proof}
$\vdash\bigl((\exists x\ldotp A[x])\implies C\bigr)\implies\bigl(\forall x\ldotp(A[x]\implies C)\bigr)$
\begin{enumerate}
\item $A[x],(\exists x\ldotp A[x])\implies C\vdash A[x]$ by assume
\item $A[x],(\exists x\ldotp A[x])\implies C\vdash\exists x\ldotp A[x]$ by $\exists$-intro
\item $A[x],(\exists x\ldotp A[x])\implies C\vdash(\exists x\ldotp A[x])\implies C$ by assume
\item $A[x],(\exists x\ldotp A[x])\implies C\vdash C$ by modus ponens(3, 2)
\item $(\exists x\ldotp A[x])\implies C\vdash A[x]\implies C$ by
  discharging $A[x]$
\item $(\exists x\ldotp A[x])\implies C\vdash\forall x\ldotp(A[x]\implies C)$ by
  $\forall$-introduction
\item $\vdash\bigl((\exists x\ldotp A[x])\implies C\bigr)\implies\bigl(\forall x\ldotp(A[x]\implies C)\bigr)$ by
  discharging the remaining hypothesis.
\end{enumerate}
$\vdash\bigl(\forall x\ldotp(A[x]\implies C)\bigr)\implies\bigl((\exists x\ldotp A[x])\implies C\bigr)$
\begin{enumerate}
\item $\bigl(\forall x\ldotp(A[x]\implies C)\bigr)\vdash\bigl(\forall x\ldotp(A[x]\implies C)\bigr)$
  by assuming
\item $\bigl(\forall x\ldotp(A[x]\implies C)\bigr)\vdash A[x]\implies C$
  by $\forall$-elimination
\item $\bigl(\forall x\ldotp(A[x]\implies C)\bigr)\vdash(\exists x\ldotp A[x])\implies C$
  by $\exists$-elimination
\item $\vdash\bigl(\forall x\ldotp(A[x]\implies C)\bigr)\implies\bigl((\exists x\ldotp A[x])\implies C\bigr)$
  by discharging the only hypothesis.\qedhere
\end{enumerate}
\end{proof}

\nwenddocs{}\nwbegincode{101}\sublabel{NW1NjfzY-uZfts-6}\nwmargintag{{\nwtagstyle{}\subpageref{NW1NjfzY-uZfts-6}}}\moddef{Derived results for Intuitionistic predicate logic~{\nwtagstyle{}\subpageref{NW1NjfzY-uZfts-1}}}\plusendmoddef\nwstartdeflinemarkup\nwusesondefline{\\{NW1NjfzY-4dPl0p-1}}\nwprevnextdefs{NW1NjfzY-uZfts-5}{NW1NjfzY-uZfts-7}\nwenddeflinemarkup
local
fun backward x A C =
  let
    val ex_A = \nwlinkedidentc{Formula.mk_exists}{NW3w0iRO-3SlqND-4}(x, A);
    val ex_A_imp_C = \nwlinkedidentc{Formula.mk_imp}{NW3w0iRO-3SlqND-3}(ex_A, C);
    (* val A_imp_C = \nwlinkedidentc{Formula.mk_imp}{NW3w0iRO-3SlqND-3}(A, C); *)
    val th1 = \nwlinkedidentc{Thm.assume}{NW4ICsJQ-2zKkPU-2} A [ex_A_imp_C];
    val th2 = \nwlinkedidentc{Thm.exists_intro}{NW4ICsJQ-2zKkPU-A} x x A th1;
    val th3 = \nwlinkedidentc{Thm.assume}{NW4ICsJQ-2zKkPU-2} ex_A_imp_C [A];
    val th4 = \nwlinkedidentc{Thm.modus_ponens}{NW4ICsJQ-2zKkPU-2} th3 th2;
    val th5 = \nwlinkedidentc{Thm.disch}{NW4ICsJQ-2zKkPU-2} A th4;
    val th6 = \nwlinkedidentc{Thm.forall_intro}{NW4ICsJQ-2zKkPU-8} x th5;
  in
    \nwlinkedidentc{Thm.disch}{NW4ICsJQ-2zKkPU-2} ex_A_imp_C th6
  end;

fun forward x A C =
  let
    val A_imp_C = \nwlinkedidentc{Formula.mk_imp}{NW3w0iRO-3SlqND-3}(A, C);
    val for_x_st_A_holds_C = \nwlinkedidentc{Formula.mk_forall}{NW3w0iRO-3SlqND-4}(x, A_imp_C);
    val th1 = \nwlinkedidentc{Thm.assume}{NW4ICsJQ-2zKkPU-2} for_x_st_A_holds_C [];
    val th2 = \nwlinkedidentc{Thm.forall_elim}{NW4ICsJQ-2zKkPU-9} x th1;
    val th3 = \nwlinkedidentc{Thm.exists_elim}{NW4ICsJQ-2zKkPU-B} x th2;
  in
    \nwlinkedidentc{Thm.disch}{NW4ICsJQ-2zKkPU-2} for_x_st_A_holds_C th3
  end;
in
fun imp_ex x A C =
  \nwlinkedidentc{Thm.iff_intro}{NW4ICsJQ-2zKkPU-7} (forward x A C) (backward x A C)
end;
\nwindexdefn{\nwixident{Derived.imp{\_}ex}}{Derived.imp:unex}{NW1NjfzY-uZfts-6}\eatline
\nwused{\\{NW1NjfzY-4dPl0p-1}}\nwidentdefs{\\{{\nwixident{Derived.imp{\_}ex}}{Derived.imp:unex}}}\nwidentuses{\\{{\nwixident{Formula.mk{\_}exists}}{Formula.mk:unexists}}\\{{\nwixident{Formula.mk{\_}forall}}{Formula.mk:unforall}}\\{{\nwixident{Formula.mk{\_}imp}}{Formula.mk:unimp}}\\{{\nwixident{Thm.assume}}{Thm.assume}}\\{{\nwixident{Thm.disch}}{Thm.disch}}\\{{\nwixident{Thm.exists{\_}elim}}{Thm.exists:unelim}}\\{{\nwixident{Thm.exists{\_}intro}}{Thm.exists:unintro}}\\{{\nwixident{Thm.forall{\_}elim}}{Thm.forall:unelim}}\\{{\nwixident{Thm.forall{\_}intro}}{Thm.forall:unintro}}\\{{\nwixident{Thm.iff{\_}intro}}{Thm.iff:unintro}}\\{{\nwixident{Thm.modus{\_}ponens}}{Thm.modus:unponens}}}\nwindexuse{\nwixident{Formula.mk{\_}exists}}{Formula.mk:unexists}{NW1NjfzY-uZfts-6}\nwindexuse{\nwixident{Formula.mk{\_}forall}}{Formula.mk:unforall}{NW1NjfzY-uZfts-6}\nwindexuse{\nwixident{Formula.mk{\_}imp}}{Formula.mk:unimp}{NW1NjfzY-uZfts-6}\nwindexuse{\nwixident{Thm.assume}}{Thm.assume}{NW1NjfzY-uZfts-6}\nwindexuse{\nwixident{Thm.disch}}{Thm.disch}{NW1NjfzY-uZfts-6}\nwindexuse{\nwixident{Thm.exists{\_}elim}}{Thm.exists:unelim}{NW1NjfzY-uZfts-6}\nwindexuse{\nwixident{Thm.exists{\_}intro}}{Thm.exists:unintro}{NW1NjfzY-uZfts-6}\nwindexuse{\nwixident{Thm.forall{\_}elim}}{Thm.forall:unelim}{NW1NjfzY-uZfts-6}\nwindexuse{\nwixident{Thm.forall{\_}intro}}{Thm.forall:unintro}{NW1NjfzY-uZfts-6}\nwindexuse{\nwixident{Thm.iff{\_}intro}}{Thm.iff:unintro}{NW1NjfzY-uZfts-6}\nwindexuse{\nwixident{Thm.modus{\_}ponens}}{Thm.modus:unponens}{NW1NjfzY-uZfts-6}\nwendcode{}\nwbegindocs{102}\nwdocspar
\nwenddocs{}\nwbegincode{103}\sublabel{NW1NjfzY-3rvgpx-5}\nwmargintag{{\nwtagstyle{}\subpageref{NW1NjfzY-3rvgpx-5}}}\moddef{Unit tests for \code{}Derived.sml\edoc{}~{\nwtagstyle{}\subpageref{NW1NjfzY-3rvgpx-1}}}\plusendmoddef\nwstartdeflinemarkup\nwprevnextdefs{NW1NjfzY-3rvgpx-4}{NW1NjfzY-3rvgpx-6}\nwenddeflinemarkup
fun imp_ex_test () =
  let
    val x = Term.Var("x", \nwlinkedidentc{Sort.IND}{NW1rTmOJ-3ZqggC-1});
    val A = \nwlinkedidentc{Formula.mk_pred}{NW3w0iRO-3SlqND-1}("A", [x]);
    val C = \nwlinkedidentc{Formula.mk_pred}{NW3w0iRO-3SlqND-1}("C", []);
    val expected =
      \nwlinkedidentc{Formula.mk_iff}{NW3w0iRO-3SlqND-3}(\nwlinkedidentc{Formula.mk_forall}{NW3w0iRO-3SlqND-4}(x, \nwlinkedidentc{Formula.mk_imp}{NW3w0iRO-3SlqND-3}(A,C)),
                     \nwlinkedidentc{Formula.mk_imp}{NW3w0iRO-3SlqND-3}(\nwlinkedidentc{Formula.mk_exists}{NW3w0iRO-3SlqND-4}(x, A),
                                    C));
    val actual = \nwlinkedidentc{Thm.concl}{NW4ICsJQ-YIPBT-1} (\nwlinkedidentc{Derived.imp_ex}{NW1NjfzY-uZfts-6} x A C);
  in
    if not(\nwlinkedidentc{Formula.eq}{NW3w0iRO-5gQzE-1} expected actual)
    then "Expecting: "^(\nwlinkedidentc{Formula.serialize}{NW3w0iRO-3H9tUA-1} expected)^"\\nActual: "^
         (\nwlinkedidentc{Formula.serialize}{NW3w0iRO-3H9tUA-1} actual)^"\\n"
    else ""
  end;
\nwidentuses{\\{{\nwixident{Derived.imp{\_}ex}}{Derived.imp:unex}}\\{{\nwixident{Formula.eq}}{Formula.eq}}\\{{\nwixident{Formula.mk{\_}exists}}{Formula.mk:unexists}}\\{{\nwixident{Formula.mk{\_}forall}}{Formula.mk:unforall}}\\{{\nwixident{Formula.mk{\_}iff}}{Formula.mk:uniff}}\\{{\nwixident{Formula.mk{\_}imp}}{Formula.mk:unimp}}\\{{\nwixident{Formula.mk{\_}pred}}{Formula.mk:unpred}}\\{{\nwixident{Formula.serialize}}{Formula.serialize}}\\{{\nwixident{Sort.IND}}{Sort.IND}}\\{{\nwixident{Thm.concl}}{Thm.concl}}}\nwindexuse{\nwixident{Derived.imp{\_}ex}}{Derived.imp:unex}{NW1NjfzY-3rvgpx-5}\nwindexuse{\nwixident{Formula.eq}}{Formula.eq}{NW1NjfzY-3rvgpx-5}\nwindexuse{\nwixident{Formula.mk{\_}exists}}{Formula.mk:unexists}{NW1NjfzY-3rvgpx-5}\nwindexuse{\nwixident{Formula.mk{\_}forall}}{Formula.mk:unforall}{NW1NjfzY-3rvgpx-5}\nwindexuse{\nwixident{Formula.mk{\_}iff}}{Formula.mk:uniff}{NW1NjfzY-3rvgpx-5}\nwindexuse{\nwixident{Formula.mk{\_}imp}}{Formula.mk:unimp}{NW1NjfzY-3rvgpx-5}\nwindexuse{\nwixident{Formula.mk{\_}pred}}{Formula.mk:unpred}{NW1NjfzY-3rvgpx-5}\nwindexuse{\nwixident{Formula.serialize}}{Formula.serialize}{NW1NjfzY-3rvgpx-5}\nwindexuse{\nwixident{Sort.IND}}{Sort.IND}{NW1NjfzY-3rvgpx-5}\nwindexuse{\nwixident{Thm.concl}}{Thm.concl}{NW1NjfzY-3rvgpx-5}\nwendcode{}\nwbegindocs{104}\nwdocspar

\begin{node}
$\vdash(\forall x\ldotp\neg A[x])\iff(\neg(\exists x\ldotp A[x]))$

This is a variant of the previous result, except that it's faster.
\end{node}

\nwenddocs{}\nwbegincode{105}\sublabel{NW1NjfzY-uZfts-7}\nwmargintag{{\nwtagstyle{}\subpageref{NW1NjfzY-uZfts-7}}}\moddef{Derived results for Intuitionistic predicate logic~{\nwtagstyle{}\subpageref{NW1NjfzY-uZfts-1}}}\plusendmoddef\nwstartdeflinemarkup\nwusesondefline{\\{NW1NjfzY-4dPl0p-1}}\nwprevnextdefs{NW1NjfzY-uZfts-6}{NW1NjfzY-uZfts-8}\nwenddeflinemarkup
local
fun backward x A =
  let
    val ex_A = \nwlinkedidentc{Formula.mk_exists}{NW3w0iRO-3SlqND-4}(x, A);
    val F = \nwlinkedidentc{Formula.mk_falsum}{NW3w0iRO-3SlqND-1} ();
    val ex_A_imp_F = \nwlinkedidentc{Formula.mk_imp}{NW3w0iRO-3SlqND-3}(ex_A, F);
    (* val A_imp_F = \nwlinkedidentc{Formula.mk_imp}{NW3w0iRO-3SlqND-3}(A, F); *)
    val th1 = \nwlinkedidentc{Thm.assume}{NW4ICsJQ-2zKkPU-2} A [ex_A_imp_F];
    val th2 = \nwlinkedidentc{Thm.exists_intro}{NW4ICsJQ-2zKkPU-A} x x A th1;
    val th3 = \nwlinkedidentc{Thm.assume}{NW4ICsJQ-2zKkPU-2} ex_A_imp_F [A];
    val th4 = \nwlinkedidentc{Thm.modus_ponens}{NW4ICsJQ-2zKkPU-2} th3 th2; (* ... |- false *)
    val th5 = \nwlinkedidentc{Thm.disch}{NW4ICsJQ-2zKkPU-2} A th4; (* ... |- A => false *)
    val th6 = \nwlinkedidentc{Thm.not_intro}{NW4ICsJQ-2zKkPU-5} th5; (* ... |- not A[x] *)
    val th7 = \nwlinkedidentc{Thm.forall_intro}{NW4ICsJQ-2zKkPU-8} x th6; (* ... |- for x holds not A[x] *)
    val th8 = \nwlinkedidentc{Thm.assume}{NW4ICsJQ-2zKkPU-2} (\nwlinkedidentc{Formula.mk_not}{NW3w0iRO-3SlqND-2} ex_A) [];
    val th9 = \nwlinkedidentc{Thm.not_elim}{NW4ICsJQ-2zKkPU-5} th8;
  in
    hypo_syll (\nwlinkedidentc{Thm.disch}{NW4ICsJQ-2zKkPU-2} (\nwlinkedidentc{Formula.mk_not}{NW3w0iRO-3SlqND-2} ex_A) th9)
              (\nwlinkedidentc{Thm.disch}{NW4ICsJQ-2zKkPU-2} ex_A_imp_F th7)
  end;

fun forward x A =
  let
    val F = \nwlinkedidentc{Formula.mk_falsum}{NW3w0iRO-3SlqND-1} ();
    val A_imp_F = \nwlinkedidentc{Formula.mk_imp}{NW3w0iRO-3SlqND-3}(A, F);
    val for_x_holds_not_A = \nwlinkedidentc{Formula.mk_forall}{NW3w0iRO-3SlqND-4}(x, \nwlinkedidentc{Formula.mk_not}{NW3w0iRO-3SlqND-2}(A));
    val th1 = \nwlinkedidentc{Thm.assume}{NW4ICsJQ-2zKkPU-2} for_x_holds_not_A []; 
    val th2 = \nwlinkedidentc{Thm.forall_elim}{NW4ICsJQ-2zKkPU-9} x th1; (* ... |- not A[x] *)
    val th3 = \nwlinkedidentc{Thm.not_elim}{NW4ICsJQ-2zKkPU-5} th2; (* ... |- A[x] => false *)
    val th4 = \nwlinkedidentc{Thm.exists_elim}{NW4ICsJQ-2zKkPU-B} x th3; (* ... |- (exists x. A[x]) => false *)
    val th5 = \nwlinkedidentc{Thm.not_intro}{NW4ICsJQ-2zKkPU-5} th4; (* ... |- not(exists x. A[x]) *)
  in
    \nwlinkedidentc{Thm.disch}{NW4ICsJQ-2zKkPU-2} for_x_holds_not_A th5
  end;
in
fun not_all_iff_ex_not x A =
  \nwlinkedidentc{Thm.iff_intro}{NW4ICsJQ-2zKkPU-7} (forward x A) (backward x A)
end;
\nwindexdefn{\nwixident{Derived.not{\_}all{\_}iff{\_}ex{\_}not}}{Derived.not:unall:uniff:unex:unnot}{NW1NjfzY-uZfts-7}\eatline
\nwused{\\{NW1NjfzY-4dPl0p-1}}\nwidentdefs{\\{{\nwixident{Derived.not{\_}all{\_}iff{\_}ex{\_}not}}{Derived.not:unall:uniff:unex:unnot}}}\nwidentuses{\\{{\nwixident{Formula.mk{\_}exists}}{Formula.mk:unexists}}\\{{\nwixident{Formula.mk{\_}falsum}}{Formula.mk:unfalsum}}\\{{\nwixident{Formula.mk{\_}forall}}{Formula.mk:unforall}}\\{{\nwixident{Formula.mk{\_}imp}}{Formula.mk:unimp}}\\{{\nwixident{Formula.mk{\_}not}}{Formula.mk:unnot}}\\{{\nwixident{Thm.assume}}{Thm.assume}}\\{{\nwixident{Thm.disch}}{Thm.disch}}\\{{\nwixident{Thm.exists{\_}elim}}{Thm.exists:unelim}}\\{{\nwixident{Thm.exists{\_}intro}}{Thm.exists:unintro}}\\{{\nwixident{Thm.forall{\_}elim}}{Thm.forall:unelim}}\\{{\nwixident{Thm.forall{\_}intro}}{Thm.forall:unintro}}\\{{\nwixident{Thm.iff{\_}intro}}{Thm.iff:unintro}}\\{{\nwixident{Thm.modus{\_}ponens}}{Thm.modus:unponens}}\\{{\nwixident{Thm.not{\_}elim}}{Thm.not:unelim}}\\{{\nwixident{Thm.not{\_}intro}}{Thm.not:unintro}}}\nwindexuse{\nwixident{Formula.mk{\_}exists}}{Formula.mk:unexists}{NW1NjfzY-uZfts-7}\nwindexuse{\nwixident{Formula.mk{\_}falsum}}{Formula.mk:unfalsum}{NW1NjfzY-uZfts-7}\nwindexuse{\nwixident{Formula.mk{\_}forall}}{Formula.mk:unforall}{NW1NjfzY-uZfts-7}\nwindexuse{\nwixident{Formula.mk{\_}imp}}{Formula.mk:unimp}{NW1NjfzY-uZfts-7}\nwindexuse{\nwixident{Formula.mk{\_}not}}{Formula.mk:unnot}{NW1NjfzY-uZfts-7}\nwindexuse{\nwixident{Thm.assume}}{Thm.assume}{NW1NjfzY-uZfts-7}\nwindexuse{\nwixident{Thm.disch}}{Thm.disch}{NW1NjfzY-uZfts-7}\nwindexuse{\nwixident{Thm.exists{\_}elim}}{Thm.exists:unelim}{NW1NjfzY-uZfts-7}\nwindexuse{\nwixident{Thm.exists{\_}intro}}{Thm.exists:unintro}{NW1NjfzY-uZfts-7}\nwindexuse{\nwixident{Thm.forall{\_}elim}}{Thm.forall:unelim}{NW1NjfzY-uZfts-7}\nwindexuse{\nwixident{Thm.forall{\_}intro}}{Thm.forall:unintro}{NW1NjfzY-uZfts-7}\nwindexuse{\nwixident{Thm.iff{\_}intro}}{Thm.iff:unintro}{NW1NjfzY-uZfts-7}\nwindexuse{\nwixident{Thm.modus{\_}ponens}}{Thm.modus:unponens}{NW1NjfzY-uZfts-7}\nwindexuse{\nwixident{Thm.not{\_}elim}}{Thm.not:unelim}{NW1NjfzY-uZfts-7}\nwindexuse{\nwixident{Thm.not{\_}intro}}{Thm.not:unintro}{NW1NjfzY-uZfts-7}\nwendcode{}\nwbegindocs{106}\nwdocspar
\nwenddocs{}\nwbegincode{107}\sublabel{NW1NjfzY-3rvgpx-6}\nwmargintag{{\nwtagstyle{}\subpageref{NW1NjfzY-3rvgpx-6}}}\moddef{Unit tests for \code{}Derived.sml\edoc{}~{\nwtagstyle{}\subpageref{NW1NjfzY-3rvgpx-1}}}\plusendmoddef\nwstartdeflinemarkup\nwprevnextdefs{NW1NjfzY-3rvgpx-5}{NW1NjfzY-3rvgpx-7}\nwenddeflinemarkup
fun not_all_iff_ex_not_test () =
  let
    val x = Term.Var("x", \nwlinkedidentc{Sort.IND}{NW1rTmOJ-3ZqggC-1});
    val A = \nwlinkedidentc{Formula.mk_pred}{NW3w0iRO-3SlqND-1}("A", [x]);
    val expected =
      \nwlinkedidentc{Formula.mk_iff}{NW3w0iRO-3SlqND-3}(\nwlinkedidentc{Formula.mk_forall}{NW3w0iRO-3SlqND-4}(x, \nwlinkedidentc{Formula.mk_not}{NW3w0iRO-3SlqND-2}(A)),
                     \nwlinkedidentc{Formula.mk_not}{NW3w0iRO-3SlqND-2}(\nwlinkedidentc{Formula.mk_exists}{NW3w0iRO-3SlqND-4}(x, A)));
    val actual = \nwlinkedidentc{Thm.concl}{NW4ICsJQ-YIPBT-1} (\nwlinkedidentc{Derived.not_all_iff_ex_not}{NW1NjfzY-uZfts-7} x A);
  in
    if not(\nwlinkedidentc{Formula.eq}{NW3w0iRO-5gQzE-1} expected actual)
    then "Expecting: "^(\nwlinkedidentc{Formula.serialize}{NW3w0iRO-3H9tUA-1} expected)^"\\nActual: "^
         (\nwlinkedidentc{Formula.serialize}{NW3w0iRO-3H9tUA-1} actual)^"\\n"
    else ""
  end;
\nwidentuses{\\{{\nwixident{Derived.not{\_}all{\_}iff{\_}ex{\_}not}}{Derived.not:unall:uniff:unex:unnot}}\\{{\nwixident{Formula.eq}}{Formula.eq}}\\{{\nwixident{Formula.mk{\_}exists}}{Formula.mk:unexists}}\\{{\nwixident{Formula.mk{\_}forall}}{Formula.mk:unforall}}\\{{\nwixident{Formula.mk{\_}iff}}{Formula.mk:uniff}}\\{{\nwixident{Formula.mk{\_}not}}{Formula.mk:unnot}}\\{{\nwixident{Formula.mk{\_}pred}}{Formula.mk:unpred}}\\{{\nwixident{Formula.serialize}}{Formula.serialize}}\\{{\nwixident{Sort.IND}}{Sort.IND}}\\{{\nwixident{Thm.concl}}{Thm.concl}}}\nwindexuse{\nwixident{Derived.not{\_}all{\_}iff{\_}ex{\_}not}}{Derived.not:unall:uniff:unex:unnot}{NW1NjfzY-3rvgpx-6}\nwindexuse{\nwixident{Formula.eq}}{Formula.eq}{NW1NjfzY-3rvgpx-6}\nwindexuse{\nwixident{Formula.mk{\_}exists}}{Formula.mk:unexists}{NW1NjfzY-3rvgpx-6}\nwindexuse{\nwixident{Formula.mk{\_}forall}}{Formula.mk:unforall}{NW1NjfzY-3rvgpx-6}\nwindexuse{\nwixident{Formula.mk{\_}iff}}{Formula.mk:uniff}{NW1NjfzY-3rvgpx-6}\nwindexuse{\nwixident{Formula.mk{\_}not}}{Formula.mk:unnot}{NW1NjfzY-3rvgpx-6}\nwindexuse{\nwixident{Formula.mk{\_}pred}}{Formula.mk:unpred}{NW1NjfzY-3rvgpx-6}\nwindexuse{\nwixident{Formula.serialize}}{Formula.serialize}{NW1NjfzY-3rvgpx-6}\nwindexuse{\nwixident{Sort.IND}}{Sort.IND}{NW1NjfzY-3rvgpx-6}\nwindexuse{\nwixident{Thm.concl}}{Thm.concl}{NW1NjfzY-3rvgpx-6}\nwendcode{}\nwbegindocs{108}\nwdocspar

\begin{node}
$\vdash(\forall x\ldotp(C\implies A[x]))\iff(C\implies\forall x\ldotp A[x])$
\end{node}

\begin{proof}
\begin{enumerate}
\item $\forall x\ldotp(C\implies A[x])\vdash\forall x\ldotp(C\implies A[x])$
by assuming%1
\item $\forall x\ldotp(C\implies A[x])\vdash C\implies A[x]$
  by $\forall$-elim($x$, 1)%2
\item $C,\forall x\ldotp(C\implies A[x])\vdash A[x]$
  by undischarging%3
\item $C,\forall x\ldotp(C\implies A[x])\vdash\forall x\ldotp A[x]$
  by $\forall$-intro($x$, 3)%4
\item $\forall x\ldotp(C\implies A[x])\vdash C\implies\forall x\ldotp A[x]$
  by discharging $C$%5
\item $\vdash(\forall x\ldotp(C\implies A[x]))\implies(C\implies\forall x\ldotp A[x])$
  by discharging%6
\item $C\implies\forall x\ldotp A[x]\vdash C\implies\forall x\ldotp A[x]$
  by assuming%7
\item $C,C\implies\forall x\ldotp A[x]\vdash\forall x\ldotp A[x]$
  by undischarging%8
\item $C,C\implies\forall x\ldotp A[x]\vdash A[x]$
  by $\forall$-elim($x$, 8)%9
\item $C\implies\forall x\ldotp A[x]\vdash C\implies A[x]$
  by discharging $C$%10
\item $C\implies\forall x\ldotp A[x]\vdash\forall x\ldotp(C\implies A[x])$
  by $\forall$-intro($x$, 10)%11
\item $\vdash(C\implies\forall x\ldotp A[x])\implies(\forall x\ldotp(C\implies A[x]))$
  by discharging%13
\item $\vdash(\forall x\ldotp(C\implies A[x]))\iff(C\implies\forall x\ldotp A[x])$
  by $\Leftrightarrow$-intro(6, 13)\qedhere%14
\end{enumerate}
\end{proof}

\nwenddocs{}\nwbegincode{109}\sublabel{NW1NjfzY-uZfts-8}\nwmargintag{{\nwtagstyle{}\subpageref{NW1NjfzY-uZfts-8}}}\moddef{Derived results for Intuitionistic predicate logic~{\nwtagstyle{}\subpageref{NW1NjfzY-uZfts-1}}}\plusendmoddef\nwstartdeflinemarkup\nwusesondefline{\\{NW1NjfzY-4dPl0p-1}}\nwprevnextdefs{NW1NjfzY-uZfts-7}{NW1NjfzY-uZfts-9}\nwenddeflinemarkup
(* all_imp : \nwlinkedidentc{Term.t}{NW1rTmOJ-4LsRWC-1} -> \nwlinkedidentc{Formula.t}{NW3w0iRO-2jnj9W-1} -> \nwlinkedidentc{Formula.t}{NW3w0iRO-2jnj9W-1} -> \nwlinkedidentc{Thm.t}{NW4ICsJQ-YIPBT-1}

all_imp x C A =
|- (for x st C holds A[x]) iff (C implies for x holds A[x]) *)
fun all_imp x C A =
  let
    val all_C_imp_A = \nwlinkedidentc{Formula.mk_forall}{NW3w0iRO-3SlqND-4}(x, \nwlinkedidentc{Formula.mk_imp}{NW3w0iRO-3SlqND-3}(C, A));
    (* forward direction *)
    val th1 = \nwlinkedidentc{Thm.assume}{NW4ICsJQ-2zKkPU-2} all_C_imp_A [];
    val th2 = \nwlinkedidentc{Thm.forall_elim}{NW4ICsJQ-2zKkPU-9} x th1;
    val th3 = \nwlinkedidentc{Thm.undisch}{NW4ICsJQ-2zKkPU-2} th2;
    val th4 = \nwlinkedidentc{Thm.forall_intro}{NW4ICsJQ-2zKkPU-8} x th3;
    val th5 = \nwlinkedidentc{Thm.disch}{NW4ICsJQ-2zKkPU-2} C th4;
    val th6 = \nwlinkedidentc{Thm.disch}{NW4ICsJQ-2zKkPU-2} all_C_imp_A th5;
    (* backward direction *)
    val C_imp_all_A = \nwlinkedidentc{Formula.mk_imp}{NW3w0iRO-3SlqND-3}(C, \nwlinkedidentc{Formula.mk_forall}{NW3w0iRO-3SlqND-4}(x, A));
    val th7 = \nwlinkedidentc{Thm.assume}{NW4ICsJQ-2zKkPU-2} C_imp_all_A [];
    val th8 = \nwlinkedidentc{Thm.undisch}{NW4ICsJQ-2zKkPU-2} th7;
    val th9 = \nwlinkedidentc{Thm.forall_elim}{NW4ICsJQ-2zKkPU-9} x th8;
    val th10 = \nwlinkedidentc{Thm.disch}{NW4ICsJQ-2zKkPU-2} C th9;
    val th11 = \nwlinkedidentc{Thm.forall_intro}{NW4ICsJQ-2zKkPU-8} x th10;
    val th12 = \nwlinkedidentc{Thm.disch}{NW4ICsJQ-2zKkPU-2} C_imp_all_A th11;
  in
    \nwlinkedidentc{Thm.iff_intro}{NW4ICsJQ-2zKkPU-7} th6 th12
  end;
\nwindexdefn{\nwixident{Derived.all{\_}imp}}{Derived.all:unimp}{NW1NjfzY-uZfts-8}\eatline
\nwused{\\{NW1NjfzY-4dPl0p-1}}\nwidentdefs{\\{{\nwixident{Derived.all{\_}imp}}{Derived.all:unimp}}}\nwidentuses{\\{{\nwixident{Formula.mk{\_}forall}}{Formula.mk:unforall}}\\{{\nwixident{Formula.mk{\_}imp}}{Formula.mk:unimp}}\\{{\nwixident{Formula.t}}{Formula.t}}\\{{\nwixident{Term.t}}{Term.t}}\\{{\nwixident{Thm.assume}}{Thm.assume}}\\{{\nwixident{Thm.disch}}{Thm.disch}}\\{{\nwixident{Thm.forall{\_}elim}}{Thm.forall:unelim}}\\{{\nwixident{Thm.forall{\_}intro}}{Thm.forall:unintro}}\\{{\nwixident{Thm.iff{\_}intro}}{Thm.iff:unintro}}\\{{\nwixident{Thm.t}}{Thm.t}}\\{{\nwixident{Thm.undisch}}{Thm.undisch}}}\nwindexuse{\nwixident{Formula.mk{\_}forall}}{Formula.mk:unforall}{NW1NjfzY-uZfts-8}\nwindexuse{\nwixident{Formula.mk{\_}imp}}{Formula.mk:unimp}{NW1NjfzY-uZfts-8}\nwindexuse{\nwixident{Formula.t}}{Formula.t}{NW1NjfzY-uZfts-8}\nwindexuse{\nwixident{Term.t}}{Term.t}{NW1NjfzY-uZfts-8}\nwindexuse{\nwixident{Thm.assume}}{Thm.assume}{NW1NjfzY-uZfts-8}\nwindexuse{\nwixident{Thm.disch}}{Thm.disch}{NW1NjfzY-uZfts-8}\nwindexuse{\nwixident{Thm.forall{\_}elim}}{Thm.forall:unelim}{NW1NjfzY-uZfts-8}\nwindexuse{\nwixident{Thm.forall{\_}intro}}{Thm.forall:unintro}{NW1NjfzY-uZfts-8}\nwindexuse{\nwixident{Thm.iff{\_}intro}}{Thm.iff:unintro}{NW1NjfzY-uZfts-8}\nwindexuse{\nwixident{Thm.t}}{Thm.t}{NW1NjfzY-uZfts-8}\nwindexuse{\nwixident{Thm.undisch}}{Thm.undisch}{NW1NjfzY-uZfts-8}\nwendcode{}\nwbegincode{110}\sublabel{NW1NjfzY-3rvgpx-7}\nwmargintag{{\nwtagstyle{}\subpageref{NW1NjfzY-3rvgpx-7}}}\moddef{Unit tests for \code{}Derived.sml\edoc{}~{\nwtagstyle{}\subpageref{NW1NjfzY-3rvgpx-1}}}\plusendmoddef\nwstartdeflinemarkup\nwprevnextdefs{NW1NjfzY-3rvgpx-6}{NW1NjfzY-3rvgpx-8}\nwenddeflinemarkup
fun all_imp_test () =
  let
    val x = Term.Var("x", \nwlinkedidentc{Sort.IND}{NW1rTmOJ-3ZqggC-1});
    val A = \nwlinkedidentc{Formula.mk_pred}{NW3w0iRO-3SlqND-1}("A", [x]);
    val C = \nwlinkedidentc{Formula.mk_pred}{NW3w0iRO-3SlqND-1}("C", []);
    val expected =
      \nwlinkedidentc{Formula.mk_iff}{NW3w0iRO-3SlqND-3}(\nwlinkedidentc{Formula.mk_forall}{NW3w0iRO-3SlqND-4}(x, \nwlinkedidentc{Formula.mk_imp}{NW3w0iRO-3SlqND-3}(C, A)),
                     \nwlinkedidentc{Formula.mk_imp}{NW3w0iRO-3SlqND-3}(C, \nwlinkedidentc{Formula.mk_forall}{NW3w0iRO-3SlqND-4}(x, A)));
    val actual = \nwlinkedidentc{Thm.concl}{NW4ICsJQ-YIPBT-1} (\nwlinkedidentc{Derived.all_imp}{NW1NjfzY-uZfts-8} x C A);
  in
    if not(\nwlinkedidentc{Formula.eq}{NW3w0iRO-5gQzE-1} expected actual)
    then "Expecting: "^(\nwlinkedidentc{Formula.serialize}{NW3w0iRO-3H9tUA-1} expected)^"\\nActual: "^
         (\nwlinkedidentc{Formula.serialize}{NW3w0iRO-3H9tUA-1} actual)^"\\n"
    else ""
  end;
\nwidentuses{\\{{\nwixident{Derived.all{\_}imp}}{Derived.all:unimp}}\\{{\nwixident{Formula.eq}}{Formula.eq}}\\{{\nwixident{Formula.mk{\_}forall}}{Formula.mk:unforall}}\\{{\nwixident{Formula.mk{\_}iff}}{Formula.mk:uniff}}\\{{\nwixident{Formula.mk{\_}imp}}{Formula.mk:unimp}}\\{{\nwixident{Formula.mk{\_}pred}}{Formula.mk:unpred}}\\{{\nwixident{Formula.serialize}}{Formula.serialize}}\\{{\nwixident{Sort.IND}}{Sort.IND}}\\{{\nwixident{Thm.concl}}{Thm.concl}}}\nwindexuse{\nwixident{Derived.all{\_}imp}}{Derived.all:unimp}{NW1NjfzY-3rvgpx-7}\nwindexuse{\nwixident{Formula.eq}}{Formula.eq}{NW1NjfzY-3rvgpx-7}\nwindexuse{\nwixident{Formula.mk{\_}forall}}{Formula.mk:unforall}{NW1NjfzY-3rvgpx-7}\nwindexuse{\nwixident{Formula.mk{\_}iff}}{Formula.mk:uniff}{NW1NjfzY-3rvgpx-7}\nwindexuse{\nwixident{Formula.mk{\_}imp}}{Formula.mk:unimp}{NW1NjfzY-3rvgpx-7}\nwindexuse{\nwixident{Formula.mk{\_}pred}}{Formula.mk:unpred}{NW1NjfzY-3rvgpx-7}\nwindexuse{\nwixident{Formula.serialize}}{Formula.serialize}{NW1NjfzY-3rvgpx-7}\nwindexuse{\nwixident{Sort.IND}}{Sort.IND}{NW1NjfzY-3rvgpx-7}\nwindexuse{\nwixident{Thm.concl}}{Thm.concl}{NW1NjfzY-3rvgpx-7}\nwendcode{}\nwbegindocs{111}\nwdocspar

\begin{node}
$\vdash(\exists x\ldotp(A[x]\implies C))\implies((\forall x\ldotp A[x])\implies C)$
\end{node}

\begin{proof}
\begin{enumerate}
\item $A[x]\implies C\vdash A[x]\implies C$
  by assuming%1
\item $\forall x\ldotp A[x]\vdash\forall x\ldotp A[x]$ by assuming%2
\item $\forall x\ldotp A[x]\vdash A[x]$ by $\forall$-elim($x$, 2)%3
\item $(\forall x\ldotp A[x]),(A[x]\implies C)\vdash C$
  by modus ponens(1, 3)%4
\item $(A[x]\implies C)\vdash(\forall x\ldotp A[x])\implies C$
  by discharging $\forall x\ldotp A[x]$%5
\item $\vdash(A[x]\implies C)\implies((\forall x\ldotp A[x])\implies C)$
  by discharging%6
\item $\vdash(\exists x\ldotp(A[x]\implies C))\implies((\forall x\ldotp A[x])\implies C)$
  by $\exists$-elim($x$, 6)\qedhere
\end{enumerate}
\end{proof}

\nwenddocs{}\nwbegincode{112}\sublabel{NW1NjfzY-uZfts-9}\nwmargintag{{\nwtagstyle{}\subpageref{NW1NjfzY-uZfts-9}}}\moddef{Derived results for Intuitionistic predicate logic~{\nwtagstyle{}\subpageref{NW1NjfzY-uZfts-1}}}\plusendmoddef\nwstartdeflinemarkup\nwusesondefline{\\{NW1NjfzY-4dPl0p-1}}\nwprevnextdefs{NW1NjfzY-uZfts-8}{NW1NjfzY-uZfts-A}\nwenddeflinemarkup
fun imp_all x A C =
  let
    val A_imp_C = \nwlinkedidentc{Formula.mk_imp}{NW3w0iRO-3SlqND-3}(A, C);
    val all_A = \nwlinkedidentc{Formula.mk_forall}{NW3w0iRO-3SlqND-4}(x, A);
    val th1 = \nwlinkedidentc{Thm.assume}{NW4ICsJQ-2zKkPU-2} A_imp_C [];
    val th2 = \nwlinkedidentc{Thm.assume}{NW4ICsJQ-2zKkPU-2} all_A [];
    val th3 = \nwlinkedidentc{Thm.forall_elim}{NW4ICsJQ-2zKkPU-9} x th2;
    val th4 = \nwlinkedidentc{Thm.modus_ponens}{NW4ICsJQ-2zKkPU-2} th1 th3;
    val th5 = \nwlinkedidentc{Thm.disch}{NW4ICsJQ-2zKkPU-2} all_A th4;
    val th6 = \nwlinkedidentc{Thm.disch}{NW4ICsJQ-2zKkPU-2} A_imp_C th5;
  in
    \nwlinkedidentc{Thm.exists_elim}{NW4ICsJQ-2zKkPU-B} x th6
  end;
\nwindexdefn{\nwixident{Derived.imp{\_}all}}{Derived.imp:unall}{NW1NjfzY-uZfts-9}\eatline
\nwused{\\{NW1NjfzY-4dPl0p-1}}\nwidentdefs{\\{{\nwixident{Derived.imp{\_}all}}{Derived.imp:unall}}}\nwidentuses{\\{{\nwixident{Formula.mk{\_}forall}}{Formula.mk:unforall}}\\{{\nwixident{Formula.mk{\_}imp}}{Formula.mk:unimp}}\\{{\nwixident{Thm.assume}}{Thm.assume}}\\{{\nwixident{Thm.disch}}{Thm.disch}}\\{{\nwixident{Thm.exists{\_}elim}}{Thm.exists:unelim}}\\{{\nwixident{Thm.forall{\_}elim}}{Thm.forall:unelim}}\\{{\nwixident{Thm.modus{\_}ponens}}{Thm.modus:unponens}}}\nwindexuse{\nwixident{Formula.mk{\_}forall}}{Formula.mk:unforall}{NW1NjfzY-uZfts-9}\nwindexuse{\nwixident{Formula.mk{\_}imp}}{Formula.mk:unimp}{NW1NjfzY-uZfts-9}\nwindexuse{\nwixident{Thm.assume}}{Thm.assume}{NW1NjfzY-uZfts-9}\nwindexuse{\nwixident{Thm.disch}}{Thm.disch}{NW1NjfzY-uZfts-9}\nwindexuse{\nwixident{Thm.exists{\_}elim}}{Thm.exists:unelim}{NW1NjfzY-uZfts-9}\nwindexuse{\nwixident{Thm.forall{\_}elim}}{Thm.forall:unelim}{NW1NjfzY-uZfts-9}\nwindexuse{\nwixident{Thm.modus{\_}ponens}}{Thm.modus:unponens}{NW1NjfzY-uZfts-9}\nwendcode{}\nwbegindocs{113}\nwdocspar
\nwenddocs{}\nwbegincode{114}\sublabel{NW1NjfzY-3rvgpx-8}\nwmargintag{{\nwtagstyle{}\subpageref{NW1NjfzY-3rvgpx-8}}}\moddef{Unit tests for \code{}Derived.sml\edoc{}~{\nwtagstyle{}\subpageref{NW1NjfzY-3rvgpx-1}}}\plusendmoddef\nwstartdeflinemarkup\nwprevnextdefs{NW1NjfzY-3rvgpx-7}{NW1NjfzY-3rvgpx-9}\nwenddeflinemarkup
fun imp_all_test () =
  let
    val x = Term.Var("x", \nwlinkedidentc{Sort.IND}{NW1rTmOJ-3ZqggC-1});
    val A = \nwlinkedidentc{Formula.mk_pred}{NW3w0iRO-3SlqND-1}("A", [x]);
    val C = \nwlinkedidentc{Formula.mk_pred}{NW3w0iRO-3SlqND-1}("C", []);
    val expected =
      \nwlinkedidentc{Formula.mk_imp}{NW3w0iRO-3SlqND-3}(\nwlinkedidentc{Formula.mk_exists}{NW3w0iRO-3SlqND-4}(x, \nwlinkedidentc{Formula.mk_imp}{NW3w0iRO-3SlqND-3}(A,C)),
                     \nwlinkedidentc{Formula.mk_imp}{NW3w0iRO-3SlqND-3}(\nwlinkedidentc{Formula.mk_forall}{NW3w0iRO-3SlqND-4}(x, A),
                                    C));
    val actual = \nwlinkedidentc{Thm.concl}{NW4ICsJQ-YIPBT-1} (\nwlinkedidentc{Derived.imp_all}{NW1NjfzY-uZfts-9} x A C);
  in
    if not(\nwlinkedidentc{Formula.eq}{NW3w0iRO-5gQzE-1} expected actual)
    then "Expecting: "^(\nwlinkedidentc{Formula.serialize}{NW3w0iRO-3H9tUA-1} expected)^"\\nActual: "^
         (\nwlinkedidentc{Formula.serialize}{NW3w0iRO-3H9tUA-1} actual)^"\\n"
    else ""
  end;
\nwidentuses{\\{{\nwixident{Derived.imp{\_}all}}{Derived.imp:unall}}\\{{\nwixident{Formula.eq}}{Formula.eq}}\\{{\nwixident{Formula.mk{\_}exists}}{Formula.mk:unexists}}\\{{\nwixident{Formula.mk{\_}forall}}{Formula.mk:unforall}}\\{{\nwixident{Formula.mk{\_}imp}}{Formula.mk:unimp}}\\{{\nwixident{Formula.mk{\_}pred}}{Formula.mk:unpred}}\\{{\nwixident{Formula.serialize}}{Formula.serialize}}\\{{\nwixident{Sort.IND}}{Sort.IND}}\\{{\nwixident{Thm.concl}}{Thm.concl}}}\nwindexuse{\nwixident{Derived.imp{\_}all}}{Derived.imp:unall}{NW1NjfzY-3rvgpx-8}\nwindexuse{\nwixident{Formula.eq}}{Formula.eq}{NW1NjfzY-3rvgpx-8}\nwindexuse{\nwixident{Formula.mk{\_}exists}}{Formula.mk:unexists}{NW1NjfzY-3rvgpx-8}\nwindexuse{\nwixident{Formula.mk{\_}forall}}{Formula.mk:unforall}{NW1NjfzY-3rvgpx-8}\nwindexuse{\nwixident{Formula.mk{\_}imp}}{Formula.mk:unimp}{NW1NjfzY-3rvgpx-8}\nwindexuse{\nwixident{Formula.mk{\_}pred}}{Formula.mk:unpred}{NW1NjfzY-3rvgpx-8}\nwindexuse{\nwixident{Formula.serialize}}{Formula.serialize}{NW1NjfzY-3rvgpx-8}\nwindexuse{\nwixident{Sort.IND}}{Sort.IND}{NW1NjfzY-3rvgpx-8}\nwindexuse{\nwixident{Thm.concl}}{Thm.concl}{NW1NjfzY-3rvgpx-8}\nwendcode{}\nwbegindocs{115}\nwdocspar

\begin{node}
$\vdash(\exists x\ldotp\neg A[x])\implies\neg(\forall x\ldotp A[x])$
\end{node}

\begin{proof}
\begin{enumerate}
\item $\neg A[x]\vdash\neg A[x]$
  by assuming%1
\item $\neg A[x]\vdash A[x]\implies\bot$ by $\neg$-elim(1)%2
\item $\forall x\ldotp A[x]\vdash\forall x\ldotp A[x]$ by assuming%3
\item $\forall x\ldotp A[x]\vdash A[x]$ by $\forall$-elim($x$, 3)%4
\item $(\forall x\ldotp A[x]),(\neg A[x])\vdash\bot$
  by modus ponens(2, 4)%5
\item $(\neg A[x])\vdash(\forall x\ldotp A[x]) \implies \bot$
  by discharging $\forall x\ldotp A[x]$%6
\item $(\neg A[x])\vdash\neg(\forall x\ldotp A[x])$
  by $\neg$-intro(6)%7
\item $\vdash(\neg A[x])\implies\neg(\forall x\ldotp A[x])$
  by discharging $\neg A[x]$%8
\item $\vdash(\exists x\ldotp\neg A[x])\implies\neg(\forall x\ldotp A[x])$
  by $\exists$-elim($x$, 8)\qedhere%9
\end{enumerate}
\end{proof}

\nwenddocs{}\nwbegincode{116}\sublabel{NW1NjfzY-uZfts-A}\nwmargintag{{\nwtagstyle{}\subpageref{NW1NjfzY-uZfts-A}}}\moddef{Derived results for Intuitionistic predicate logic~{\nwtagstyle{}\subpageref{NW1NjfzY-uZfts-1}}}\plusendmoddef\nwstartdeflinemarkup\nwusesondefline{\\{NW1NjfzY-4dPl0p-1}}\nwprevnextdefs{NW1NjfzY-uZfts-9}{NW1NjfzY-uZfts-B}\nwenddeflinemarkup
fun ex_not x A =
  let
    val not_A = \nwlinkedidentc{Formula.mk_not}{NW3w0iRO-3SlqND-2} A;
    val all_A = \nwlinkedidentc{Formula.mk_forall}{NW3w0iRO-3SlqND-4}(x, A);
    val F = \nwlinkedidentc{Formula.mk_falsum}{NW3w0iRO-3SlqND-1} ();
    val th1 = \nwlinkedidentc{Thm.assume}{NW4ICsJQ-2zKkPU-2} not_A [];
    val th2 = \nwlinkedidentc{Thm.not_elim}{NW4ICsJQ-2zKkPU-5} th1;
    val th3 = \nwlinkedidentc{Thm.assume}{NW4ICsJQ-2zKkPU-2} all_A [];
    val th4 = \nwlinkedidentc{Thm.forall_elim}{NW4ICsJQ-2zKkPU-9} x th3;
    val th5 = \nwlinkedidentc{Thm.modus_ponens}{NW4ICsJQ-2zKkPU-2} th2 th4;
    val th6 = \nwlinkedidentc{Thm.disch}{NW4ICsJQ-2zKkPU-2} all_A th5;
    val th7 = \nwlinkedidentc{Thm.not_intro}{NW4ICsJQ-2zKkPU-5} th6;
    val th8 = \nwlinkedidentc{Thm.disch}{NW4ICsJQ-2zKkPU-2} not_A th7;
  in
    \nwlinkedidentc{Thm.exists_elim}{NW4ICsJQ-2zKkPU-B} x th8
  end;
\nwindexdefn{\nwixident{Derived.ex{\_}not}}{Derived.ex:unnot}{NW1NjfzY-uZfts-A}\eatline
\nwused{\\{NW1NjfzY-4dPl0p-1}}\nwidentdefs{\\{{\nwixident{Derived.ex{\_}not}}{Derived.ex:unnot}}}\nwidentuses{\\{{\nwixident{Formula.mk{\_}falsum}}{Formula.mk:unfalsum}}\\{{\nwixident{Formula.mk{\_}forall}}{Formula.mk:unforall}}\\{{\nwixident{Formula.mk{\_}not}}{Formula.mk:unnot}}\\{{\nwixident{Thm.assume}}{Thm.assume}}\\{{\nwixident{Thm.disch}}{Thm.disch}}\\{{\nwixident{Thm.exists{\_}elim}}{Thm.exists:unelim}}\\{{\nwixident{Thm.forall{\_}elim}}{Thm.forall:unelim}}\\{{\nwixident{Thm.modus{\_}ponens}}{Thm.modus:unponens}}\\{{\nwixident{Thm.not{\_}elim}}{Thm.not:unelim}}\\{{\nwixident{Thm.not{\_}intro}}{Thm.not:unintro}}}\nwindexuse{\nwixident{Formula.mk{\_}falsum}}{Formula.mk:unfalsum}{NW1NjfzY-uZfts-A}\nwindexuse{\nwixident{Formula.mk{\_}forall}}{Formula.mk:unforall}{NW1NjfzY-uZfts-A}\nwindexuse{\nwixident{Formula.mk{\_}not}}{Formula.mk:unnot}{NW1NjfzY-uZfts-A}\nwindexuse{\nwixident{Thm.assume}}{Thm.assume}{NW1NjfzY-uZfts-A}\nwindexuse{\nwixident{Thm.disch}}{Thm.disch}{NW1NjfzY-uZfts-A}\nwindexuse{\nwixident{Thm.exists{\_}elim}}{Thm.exists:unelim}{NW1NjfzY-uZfts-A}\nwindexuse{\nwixident{Thm.forall{\_}elim}}{Thm.forall:unelim}{NW1NjfzY-uZfts-A}\nwindexuse{\nwixident{Thm.modus{\_}ponens}}{Thm.modus:unponens}{NW1NjfzY-uZfts-A}\nwindexuse{\nwixident{Thm.not{\_}elim}}{Thm.not:unelim}{NW1NjfzY-uZfts-A}\nwindexuse{\nwixident{Thm.not{\_}intro}}{Thm.not:unintro}{NW1NjfzY-uZfts-A}\nwendcode{}\nwbegindocs{117}\nwdocspar
\nwenddocs{}\nwbegincode{118}\sublabel{NW1NjfzY-3rvgpx-9}\nwmargintag{{\nwtagstyle{}\subpageref{NW1NjfzY-3rvgpx-9}}}\moddef{Unit tests for \code{}Derived.sml\edoc{}~{\nwtagstyle{}\subpageref{NW1NjfzY-3rvgpx-1}}}\plusendmoddef\nwstartdeflinemarkup\nwprevnextdefs{NW1NjfzY-3rvgpx-8}{NW1NjfzY-3rvgpx-A}\nwenddeflinemarkup
fun ex_not_test () =
  let
    val x = Term.Var("x", \nwlinkedidentc{Sort.IND}{NW1rTmOJ-3ZqggC-1});
    val A = \nwlinkedidentc{Formula.mk_pred}{NW3w0iRO-3SlqND-1}("A", [x]);
    val expected =
      \nwlinkedidentc{Formula.mk_imp}{NW3w0iRO-3SlqND-3}(\nwlinkedidentc{Formula.mk_exists}{NW3w0iRO-3SlqND-4}(x, \nwlinkedidentc{Formula.mk_not}{NW3w0iRO-3SlqND-2}(A)),
                     \nwlinkedidentc{Formula.mk_not}{NW3w0iRO-3SlqND-2}(\nwlinkedidentc{Formula.mk_forall}{NW3w0iRO-3SlqND-4}(x, A)));
    val actual = \nwlinkedidentc{Thm.concl}{NW4ICsJQ-YIPBT-1} (\nwlinkedidentc{Derived.ex_not}{NW1NjfzY-uZfts-A} x A);
  in
    if not(\nwlinkedidentc{Formula.eq}{NW3w0iRO-5gQzE-1} expected actual)
    then "Expecting: "^(\nwlinkedidentc{Formula.serialize}{NW3w0iRO-3H9tUA-1} expected)^"\\nActual: "^
         (\nwlinkedidentc{Formula.serialize}{NW3w0iRO-3H9tUA-1} actual)^"\\n"
    else ""
  end;
\nwidentuses{\\{{\nwixident{Derived.ex{\_}not}}{Derived.ex:unnot}}\\{{\nwixident{Formula.eq}}{Formula.eq}}\\{{\nwixident{Formula.mk{\_}exists}}{Formula.mk:unexists}}\\{{\nwixident{Formula.mk{\_}forall}}{Formula.mk:unforall}}\\{{\nwixident{Formula.mk{\_}imp}}{Formula.mk:unimp}}\\{{\nwixident{Formula.mk{\_}not}}{Formula.mk:unnot}}\\{{\nwixident{Formula.mk{\_}pred}}{Formula.mk:unpred}}\\{{\nwixident{Formula.serialize}}{Formula.serialize}}\\{{\nwixident{Sort.IND}}{Sort.IND}}\\{{\nwixident{Thm.concl}}{Thm.concl}}}\nwindexuse{\nwixident{Derived.ex{\_}not}}{Derived.ex:unnot}{NW1NjfzY-3rvgpx-9}\nwindexuse{\nwixident{Formula.eq}}{Formula.eq}{NW1NjfzY-3rvgpx-9}\nwindexuse{\nwixident{Formula.mk{\_}exists}}{Formula.mk:unexists}{NW1NjfzY-3rvgpx-9}\nwindexuse{\nwixident{Formula.mk{\_}forall}}{Formula.mk:unforall}{NW1NjfzY-3rvgpx-9}\nwindexuse{\nwixident{Formula.mk{\_}imp}}{Formula.mk:unimp}{NW1NjfzY-3rvgpx-9}\nwindexuse{\nwixident{Formula.mk{\_}not}}{Formula.mk:unnot}{NW1NjfzY-3rvgpx-9}\nwindexuse{\nwixident{Formula.mk{\_}pred}}{Formula.mk:unpred}{NW1NjfzY-3rvgpx-9}\nwindexuse{\nwixident{Formula.serialize}}{Formula.serialize}{NW1NjfzY-3rvgpx-9}\nwindexuse{\nwixident{Sort.IND}}{Sort.IND}{NW1NjfzY-3rvgpx-9}\nwindexuse{\nwixident{Thm.concl}}{Thm.concl}{NW1NjfzY-3rvgpx-9}\nwendcode{}\nwbegindocs{119}\nwdocspar

\begin{node}
$\vdash(\exists x\ldotp(C\implies A[x]))\implies(C\implies\exists x\ldotp A[x])$
\end{node}

\begin{proof}
\begin{enumerate}
\item $C\implies A[x]\vdash C\implies A[x]$ by assuming
\item $C\vdash C$ by assuming
\item $C, C\implies A[x]\vdash A[x]$ by modus ponens(1, 2)%3
\item $C, C\implies A[x]\vdash\exists x\ldotp A[x]$ by
  $\exists$-intro($x$, 3)%4
\item $C\implies A[x]\vdash C\implies\exists x\ldotp A[x]$ by
  discharging $C$%5
\item $\vdash(C\implies A[x])\implies(C\implies\exists x\ldotp A[x])$ by
  discharging $C\implies A[x]$%6
\item $\vdash(\exists x\ldotp(C\implies A[x]))\implies(C\implies\exists x\ldotp A[x])$ by
  $\exists$-elim($x$, 6)\qedhere%7
\end{enumerate}
\end{proof}

\nwenddocs{}\nwbegincode{120}\sublabel{NW1NjfzY-uZfts-B}\nwmargintag{{\nwtagstyle{}\subpageref{NW1NjfzY-uZfts-B}}}\moddef{Derived results for Intuitionistic predicate logic~{\nwtagstyle{}\subpageref{NW1NjfzY-uZfts-1}}}\plusendmoddef\nwstartdeflinemarkup\nwusesondefline{\\{NW1NjfzY-4dPl0p-1}}\nwprevnextdefs{NW1NjfzY-uZfts-A}{NW1NjfzY-uZfts-C}\nwenddeflinemarkup
(* ex_imp : \nwlinkedidentc{Term.t}{NW1rTmOJ-4LsRWC-1} -> \nwlinkedidentc{Formula.t}{NW3w0iRO-2jnj9W-1} -> \nwlinkedidentc{Formula.t}{NW3w0iRO-2jnj9W-1} -> \nwlinkedidentc{Thm.t}{NW4ICsJQ-YIPBT-1}

|- (ex x st C implies A[x]) implies (C implies ex x st A[x]) *)
fun ex_imp x C A =
  let
    val C_imp_A = \nwlinkedidentc{Formula.mk_imp}{NW3w0iRO-3SlqND-3}(C, A);
    val th1 = \nwlinkedidentc{Thm.assume}{NW4ICsJQ-2zKkPU-2} C_imp_A [];
    val th2 = \nwlinkedidentc{Thm.assume}{NW4ICsJQ-2zKkPU-2} C [];
    val th3 = \nwlinkedidentc{Thm.modus_ponens}{NW4ICsJQ-2zKkPU-2} th1 th2;
    val th4 = \nwlinkedidentc{Thm.exists_intro}{NW4ICsJQ-2zKkPU-A} x x A th3;
    val th5 = \nwlinkedidentc{Thm.disch}{NW4ICsJQ-2zKkPU-2} C th4;
    val th6 = \nwlinkedidentc{Thm.disch}{NW4ICsJQ-2zKkPU-2} C_imp_A th5;
  in
    \nwlinkedidentc{Thm.exists_elim}{NW4ICsJQ-2zKkPU-B} x th6
  end;
\nwindexdefn{\nwixident{Derived.ex{\_}imp}}{Derived.ex:unimp}{NW1NjfzY-uZfts-B}\eatline
\nwused{\\{NW1NjfzY-4dPl0p-1}}\nwidentdefs{\\{{\nwixident{Derived.ex{\_}imp}}{Derived.ex:unimp}}}\nwidentuses{\\{{\nwixident{Formula.mk{\_}imp}}{Formula.mk:unimp}}\\{{\nwixident{Formula.t}}{Formula.t}}\\{{\nwixident{Term.t}}{Term.t}}\\{{\nwixident{Thm.assume}}{Thm.assume}}\\{{\nwixident{Thm.disch}}{Thm.disch}}\\{{\nwixident{Thm.exists{\_}elim}}{Thm.exists:unelim}}\\{{\nwixident{Thm.exists{\_}intro}}{Thm.exists:unintro}}\\{{\nwixident{Thm.modus{\_}ponens}}{Thm.modus:unponens}}\\{{\nwixident{Thm.t}}{Thm.t}}}\nwindexuse{\nwixident{Formula.mk{\_}imp}}{Formula.mk:unimp}{NW1NjfzY-uZfts-B}\nwindexuse{\nwixident{Formula.t}}{Formula.t}{NW1NjfzY-uZfts-B}\nwindexuse{\nwixident{Term.t}}{Term.t}{NW1NjfzY-uZfts-B}\nwindexuse{\nwixident{Thm.assume}}{Thm.assume}{NW1NjfzY-uZfts-B}\nwindexuse{\nwixident{Thm.disch}}{Thm.disch}{NW1NjfzY-uZfts-B}\nwindexuse{\nwixident{Thm.exists{\_}elim}}{Thm.exists:unelim}{NW1NjfzY-uZfts-B}\nwindexuse{\nwixident{Thm.exists{\_}intro}}{Thm.exists:unintro}{NW1NjfzY-uZfts-B}\nwindexuse{\nwixident{Thm.modus{\_}ponens}}{Thm.modus:unponens}{NW1NjfzY-uZfts-B}\nwindexuse{\nwixident{Thm.t}}{Thm.t}{NW1NjfzY-uZfts-B}\nwendcode{}\nwbegindocs{121}\nwdocspar
\nwenddocs{}\nwbegincode{122}\sublabel{NW1NjfzY-3rvgpx-A}\nwmargintag{{\nwtagstyle{}\subpageref{NW1NjfzY-3rvgpx-A}}}\moddef{Unit tests for \code{}Derived.sml\edoc{}~{\nwtagstyle{}\subpageref{NW1NjfzY-3rvgpx-1}}}\plusendmoddef\nwstartdeflinemarkup\nwprevnextdefs{NW1NjfzY-3rvgpx-9}{\relax}\nwenddeflinemarkup
fun ex_imp_test () =
  let
    val x = Term.Var("x", \nwlinkedidentc{Sort.IND}{NW1rTmOJ-3ZqggC-1});
    val A = \nwlinkedidentc{Formula.mk_pred}{NW3w0iRO-3SlqND-1}("A", [x]);
    val C = \nwlinkedidentc{Formula.mk_pred}{NW3w0iRO-3SlqND-1}("C", []);
    val expected =
      \nwlinkedidentc{Formula.mk_imp}{NW3w0iRO-3SlqND-3}(\nwlinkedidentc{Formula.mk_exists}{NW3w0iRO-3SlqND-4}(x, \nwlinkedidentc{Formula.mk_imp}{NW3w0iRO-3SlqND-3}(C,A)),
                     \nwlinkedidentc{Formula.mk_imp}{NW3w0iRO-3SlqND-3}(C,
                                    \nwlinkedidentc{Formula.mk_exists}{NW3w0iRO-3SlqND-4}(x, A)));
    val actual = \nwlinkedidentc{Thm.concl}{NW4ICsJQ-YIPBT-1} (ex_imp x C A);
  in
    if not(\nwlinkedidentc{Formula.eq}{NW3w0iRO-5gQzE-1} expected actual)
    then "Expecting: "^(\nwlinkedidentc{Formula.serialize}{NW3w0iRO-3H9tUA-1} expected)^"\\nActual: "^
         (\nwlinkedidentc{Formula.serialize}{NW3w0iRO-3H9tUA-1} actual)^"\\n"
    else ""
  end;
\nwidentuses{\\{{\nwixident{Formula.eq}}{Formula.eq}}\\{{\nwixident{Formula.mk{\_}exists}}{Formula.mk:unexists}}\\{{\nwixident{Formula.mk{\_}imp}}{Formula.mk:unimp}}\\{{\nwixident{Formula.mk{\_}pred}}{Formula.mk:unpred}}\\{{\nwixident{Formula.serialize}}{Formula.serialize}}\\{{\nwixident{Sort.IND}}{Sort.IND}}\\{{\nwixident{Thm.concl}}{Thm.concl}}}\nwindexuse{\nwixident{Formula.eq}}{Formula.eq}{NW1NjfzY-3rvgpx-A}\nwindexuse{\nwixident{Formula.mk{\_}exists}}{Formula.mk:unexists}{NW1NjfzY-3rvgpx-A}\nwindexuse{\nwixident{Formula.mk{\_}imp}}{Formula.mk:unimp}{NW1NjfzY-3rvgpx-A}\nwindexuse{\nwixident{Formula.mk{\_}pred}}{Formula.mk:unpred}{NW1NjfzY-3rvgpx-A}\nwindexuse{\nwixident{Formula.serialize}}{Formula.serialize}{NW1NjfzY-3rvgpx-A}\nwindexuse{\nwixident{Sort.IND}}{Sort.IND}{NW1NjfzY-3rvgpx-A}\nwindexuse{\nwixident{Thm.concl}}{Thm.concl}{NW1NjfzY-3rvgpx-A}\nwendcode{}\nwbegindocs{123}\nwdocspar

\begin{node}
$\vdash(\exists x\ldotp C)\iff C$ and $\vdash(\forall x\ldotp C)\iff C$
\end{node}

\begin{node}\label{node:derived:existential-elimination}
Lastly, we will need to use a modified version of existential
elimination~(\S\ref{node:kernel:existential-elimination}) for
LCF-style tactics. We want something of the form:
\begin{equation*}
\vcenter{\infer[\exists\mbox{-elim}(v,\Gamma_{1}\vdash\exists x\ldotp A)]%
  {\Gamma_{1}\cup(\Gamma_{2}\setminus(A[v/x]))\vdash B}%
  {\Gamma_{1}\vdash\exists x\ldotp A[x], & \Gamma_{2}\vdash B}}
\end{equation*}
with the side condition that $v$ is not free in $\Gamma_{2}\setminus\{A[v/x]\}$,
$A$, or $B$.
\end{node}

\begin{proof}
\begin{enumerate}
\item\textsc{Given:} $\Gamma_{1}\vdash\exists x\ldotp A$%1
\item\textsc{Given:} $\Gamma_{2}\vdash B$%2
\item $\{A[v/x]\}\cup\Gamma_{2}\vdash B$ by weakening(2)%3
\item $\Gamma_{2}\setminus\{A[v/x]\}\vdash A[v]\implies B$ by discharging($A[v]$,3)%4
\item $\Gamma_{2}\setminus\{A[v/x]\}\vdash \exists v\ldotp A[v]\implies B$
  by existential elimination(4)%5
\item $\vdash\exists x\ldotp A[x]\implies\exists v\ldotp A[v]$
  by renaming bound variable%6
\item $\Gamma_{2}\setminus\{A[v/x]\}\vdash \exists x\ldotp A[x]\implies B$
  by hypothetical syllogism(6, 5)%7
\item $\Gamma_{1}\cup(\Gamma_{2}\setminus\{A[v/x]\})\vdash B$
  by syllogism(7, 1)\qedhere%8
\end{enumerate}
\end{proof}

\nwenddocs{}\nwbegincode{124}\sublabel{NW1NjfzY-uZfts-C}\nwmargintag{{\nwtagstyle{}\subpageref{NW1NjfzY-uZfts-C}}}\moddef{Derived results for Intuitionistic predicate logic~{\nwtagstyle{}\subpageref{NW1NjfzY-uZfts-1}}}\plusendmoddef\nwstartdeflinemarkup\nwusesondefline{\\{NW1NjfzY-4dPl0p-1}}\nwprevnextdefs{NW1NjfzY-uZfts-B}{NW1NjfzY-uZfts-D}\nwenddeflinemarkup
(* exists_elim : \nwlinkedidentc{Term.t}{NW1rTmOJ-4LsRWC-1} -> \nwlinkedidentc{Thm.t}{NW4ICsJQ-YIPBT-1} -> \nwlinkedidentc{Thm.t}{NW4ICsJQ-YIPBT-1} -> \nwlinkedidentc{Thm.t}{NW4ICsJQ-YIPBT-1} *)
fun exists_elim (v : \nwlinkedidentc{Term.t}{NW1rTmOJ-4LsRWC-1}) (th1 : \nwlinkedidentc{Thm.t}{NW4ICsJQ-YIPBT-1}) (th2 : \nwlinkedidentc{Thm.t}{NW4ICsJQ-YIPBT-1}) : \nwlinkedidentc{Thm.t}{NW4ICsJQ-YIPBT-1} =
  let
    val (x,A) = Formula.dest_exists(\nwlinkedidentc{Thm.concl}{NW4ICsJQ-YIPBT-1} th1);
    val Av = \nwlinkedidentc{Formula.subst}{NW3w0iRO-1jtQRV-1} x v A;
    val th3 = weaken [Av] th2;
    val th4 = \nwlinkedidentc{Thm.disch}{NW4ICsJQ-2zKkPU-2} Av th3;
    val th5 = \nwlinkedidentc{Thm.exists_elim}{NW4ICsJQ-2zKkPU-B} v th4;
    (* aside on renaming bound variables *)
    val lm1 = \nwlinkedidentc{Thm.assume}{NW4ICsJQ-2zKkPU-2} (\nwlinkedidentc{Formula.mk_exists}{NW3w0iRO-3SlqND-4}(x,A)) [Av];
    val lm2 = \nwlinkedidentc{Thm.disch}{NW4ICsJQ-2zKkPU-2} Av lm1;
    val lm3 = \nwlinkedidentc{Thm.disch}{NW4ICsJQ-2zKkPU-2} (\nwlinkedidentc{Formula.mk_exists}{NW3w0iRO-3SlqND-4}(x,A)) lm1;
    val th6 = \nwlinkedidentc{Thm.modus_ponens}{NW4ICsJQ-2zKkPU-2} lm3 th1;
    val th7 = hypo_syll th6 th5;
  in
    \nwlinkedidentc{Thm.modus_ponens}{NW4ICsJQ-2zKkPU-2} th7 th1
  end;
\nwindexdefn{\nwixident{Derived.exists{\_}elim}}{Derived.exists:unelim}{NW1NjfzY-uZfts-C}\eatline
\nwused{\\{NW1NjfzY-4dPl0p-1}}\nwidentdefs{\\{{\nwixident{Derived.exists{\_}elim}}{Derived.exists:unelim}}}\nwidentuses{\\{{\nwixident{Formula.mk{\_}exists}}{Formula.mk:unexists}}\\{{\nwixident{Formula.subst}}{Formula.subst}}\\{{\nwixident{Term.t}}{Term.t}}\\{{\nwixident{Thm.assume}}{Thm.assume}}\\{{\nwixident{Thm.concl}}{Thm.concl}}\\{{\nwixident{Thm.disch}}{Thm.disch}}\\{{\nwixident{Thm.exists{\_}elim}}{Thm.exists:unelim}}\\{{\nwixident{Thm.modus{\_}ponens}}{Thm.modus:unponens}}\\{{\nwixident{Thm.t}}{Thm.t}}}\nwindexuse{\nwixident{Formula.mk{\_}exists}}{Formula.mk:unexists}{NW1NjfzY-uZfts-C}\nwindexuse{\nwixident{Formula.subst}}{Formula.subst}{NW1NjfzY-uZfts-C}\nwindexuse{\nwixident{Term.t}}{Term.t}{NW1NjfzY-uZfts-C}\nwindexuse{\nwixident{Thm.assume}}{Thm.assume}{NW1NjfzY-uZfts-C}\nwindexuse{\nwixident{Thm.concl}}{Thm.concl}{NW1NjfzY-uZfts-C}\nwindexuse{\nwixident{Thm.disch}}{Thm.disch}{NW1NjfzY-uZfts-C}\nwindexuse{\nwixident{Thm.exists{\_}elim}}{Thm.exists:unelim}{NW1NjfzY-uZfts-C}\nwindexuse{\nwixident{Thm.modus{\_}ponens}}{Thm.modus:unponens}{NW1NjfzY-uZfts-C}\nwindexuse{\nwixident{Thm.t}}{Thm.t}{NW1NjfzY-uZfts-C}\nwendcode{}\nwbegindocs{125}\nwdocspar
\begin{node}
We have the following \emph{congruence of universal quantifier}
theorem, $\vdash (\forall x\ldotp A[x])\implies(\forall y\ldotp A[y])$.
\end{node}

\begin{proof}
\begin{enumerate}
\item $\forall x\ldotp A[x]\vdash\forall x\ldotp A[x]$ by assuming%1
\item $\forall x\ldotp A[x]\vdash A[y]$ by $\forall$-elim($y$, 1)%2
\item $\forall x\ldotp A[x]\vdash\forall y\ldotp A[y]$ by $\forall$-intro(2)%3
\item $\vdash(\forall x\ldotp A[x])\implies(\forall y\ldotp A[y])$
  by discharging(3)\qedhere
\end{enumerate}
\end{proof}

\nwenddocs{}\nwbegincode{126}\sublabel{NW1NjfzY-uZfts-D}\nwmargintag{{\nwtagstyle{}\subpageref{NW1NjfzY-uZfts-D}}}\moddef{Derived results for Intuitionistic predicate logic~{\nwtagstyle{}\subpageref{NW1NjfzY-uZfts-1}}}\plusendmoddef\nwstartdeflinemarkup\nwusesondefline{\\{NW1NjfzY-4dPl0p-1}}\nwprevnextdefs{NW1NjfzY-uZfts-C}{NW1NjfzY-uZfts-E}\nwenddeflinemarkup
(* forall_cong : \nwlinkedidentc{Term.t}{NW1rTmOJ-4LsRWC-1} -> \nwlinkedidentc{Term.t}{NW1rTmOJ-4LsRWC-1} -> (\nwlinkedidentc{Term.t}{NW1rTmOJ-4LsRWC-1} -> \nwlinkedidentc{Formula.t}{NW3w0iRO-2jnj9W-1} -> \nwlinkedidentc{Thm.t}{NW4ICsJQ-YIPBT-1} *)
fun forall_cong x y (A : \nwlinkedidentc{Term.t}{NW1rTmOJ-4LsRWC-1} -> \nwlinkedidentc{Formula.t}{NW3w0iRO-2jnj9W-1}) =
  if not(\nwlinkedidentc{Term.have_same_sort}{NW1rTmOJ-1WQNps-1} x y)
  then raise Fail "forall_cong: variables have different sorts"
  else
    let
      val Ax = \nwlinkedidentc{Formula.mk_forall}{NW3w0iRO-3SlqND-4}(x, A x);
      val th1 = \nwlinkedidentc{Thm.assume}{NW4ICsJQ-2zKkPU-2} Ax [];
      val th2 = \nwlinkedidentc{Thm.forall_elim}{NW4ICsJQ-2zKkPU-9} y th1;
      val th3 = \nwlinkedidentc{Thm.forall_intro}{NW4ICsJQ-2zKkPU-8} y th2;
    in
      \nwlinkedidentc{Thm.disch}{NW4ICsJQ-2zKkPU-2} Ax th3
    end;
\nwindexdefn{\nwixident{Derived.forall{\_}cong}}{Derived.forall:uncong}{NW1NjfzY-uZfts-D}\eatline
\nwused{\\{NW1NjfzY-4dPl0p-1}}\nwidentdefs{\\{{\nwixident{Derived.forall{\_}cong}}{Derived.forall:uncong}}}\nwidentuses{\\{{\nwixident{Formula.mk{\_}forall}}{Formula.mk:unforall}}\\{{\nwixident{Formula.t}}{Formula.t}}\\{{\nwixident{Term.have{\_}same{\_}sort}}{Term.have:unsame:unsort}}\\{{\nwixident{Term.t}}{Term.t}}\\{{\nwixident{Thm.assume}}{Thm.assume}}\\{{\nwixident{Thm.disch}}{Thm.disch}}\\{{\nwixident{Thm.forall{\_}elim}}{Thm.forall:unelim}}\\{{\nwixident{Thm.forall{\_}intro}}{Thm.forall:unintro}}\\{{\nwixident{Thm.t}}{Thm.t}}}\nwindexuse{\nwixident{Formula.mk{\_}forall}}{Formula.mk:unforall}{NW1NjfzY-uZfts-D}\nwindexuse{\nwixident{Formula.t}}{Formula.t}{NW1NjfzY-uZfts-D}\nwindexuse{\nwixident{Term.have{\_}same{\_}sort}}{Term.have:unsame:unsort}{NW1NjfzY-uZfts-D}\nwindexuse{\nwixident{Term.t}}{Term.t}{NW1NjfzY-uZfts-D}\nwindexuse{\nwixident{Thm.assume}}{Thm.assume}{NW1NjfzY-uZfts-D}\nwindexuse{\nwixident{Thm.disch}}{Thm.disch}{NW1NjfzY-uZfts-D}\nwindexuse{\nwixident{Thm.forall{\_}elim}}{Thm.forall:unelim}{NW1NjfzY-uZfts-D}\nwindexuse{\nwixident{Thm.forall{\_}intro}}{Thm.forall:unintro}{NW1NjfzY-uZfts-D}\nwindexuse{\nwixident{Thm.t}}{Thm.t}{NW1NjfzY-uZfts-D}\nwendcode{}\nwbegindocs{127}\nwdocspar
\begin{node}
Similarly, we have congruence for existential quantifiers,
\begin{equation*}
\vdash(\exists x\ldotp A[x])\implies(\exists y\ldotp A[y]).
\end{equation*}
\end{node}

\begin{proof}
\begin{enumerate}
\item $A[x]\vdash A[x]$ by assuming%1
\item $A[x]\vdash\exists y\ldotp A[y]$ by $\exists$-intro($y$, 1)%2
\item $\vdash A[x]\implies(\exists y\ldotp A[y])$ by discharging(2)%3
\item $\vdash(\exists x\ldotp A[x])\implies(\exists y\ldotp A[y])$
  by $\exists$-elim(4)\qedhere
\end{enumerate}
\end{proof}

\nwenddocs{}\nwbegincode{128}\sublabel{NW1NjfzY-uZfts-E}\nwmargintag{{\nwtagstyle{}\subpageref{NW1NjfzY-uZfts-E}}}\moddef{Derived results for Intuitionistic predicate logic~{\nwtagstyle{}\subpageref{NW1NjfzY-uZfts-1}}}\plusendmoddef\nwstartdeflinemarkup\nwusesondefline{\\{NW1NjfzY-4dPl0p-1}}\nwprevnextdefs{NW1NjfzY-uZfts-D}{\relax}\nwenddeflinemarkup
(* exists_cong : \nwlinkedidentc{Term.t}{NW1rTmOJ-4LsRWC-1} -> \nwlinkedidentc{Term.t}{NW1rTmOJ-4LsRWC-1} -> (\nwlinkedidentc{Term.t}{NW1rTmOJ-4LsRWC-1} -> \nwlinkedidentc{Formula.t}{NW3w0iRO-2jnj9W-1} -> \nwlinkedidentc{Thm.t}{NW4ICsJQ-YIPBT-1} *)
fun exists_cong x y (A : \nwlinkedidentc{Term.t}{NW1rTmOJ-4LsRWC-1} -> \nwlinkedidentc{Formula.t}{NW3w0iRO-2jnj9W-1}) =
  if not(\nwlinkedidentc{Term.have_same_sort}{NW1rTmOJ-1WQNps-1} x y)
  then raise Fail "forall_cong: variables have different sorts"
  else
    let
      val th1 = \nwlinkedidentc{Thm.assume}{NW4ICsJQ-2zKkPU-2} (A x) [];
      val th2 = \nwlinkedidentc{Thm.exists_intro}{NW4ICsJQ-2zKkPU-A} y x (A y) th1;
      val th3 = \nwlinkedidentc{Thm.disch}{NW4ICsJQ-2zKkPU-2} (A x) th2;
    in
      \nwlinkedidentc{Thm.exists_elim}{NW4ICsJQ-2zKkPU-B} x th3
    end;
\nwindexdefn{\nwixident{Derived.exists{\_}cong}}{Derived.exists:uncong}{NW1NjfzY-uZfts-E}\eatline
\nwused{\\{NW1NjfzY-4dPl0p-1}}\nwidentdefs{\\{{\nwixident{Derived.exists{\_}cong}}{Derived.exists:uncong}}}\nwidentuses{\\{{\nwixident{Formula.t}}{Formula.t}}\\{{\nwixident{Term.have{\_}same{\_}sort}}{Term.have:unsame:unsort}}\\{{\nwixident{Term.t}}{Term.t}}\\{{\nwixident{Thm.assume}}{Thm.assume}}\\{{\nwixident{Thm.disch}}{Thm.disch}}\\{{\nwixident{Thm.exists{\_}elim}}{Thm.exists:unelim}}\\{{\nwixident{Thm.exists{\_}intro}}{Thm.exists:unintro}}\\{{\nwixident{Thm.t}}{Thm.t}}}\nwindexuse{\nwixident{Formula.t}}{Formula.t}{NW1NjfzY-uZfts-E}\nwindexuse{\nwixident{Term.have{\_}same{\_}sort}}{Term.have:unsame:unsort}{NW1NjfzY-uZfts-E}\nwindexuse{\nwixident{Term.t}}{Term.t}{NW1NjfzY-uZfts-E}\nwindexuse{\nwixident{Thm.assume}}{Thm.assume}{NW1NjfzY-uZfts-E}\nwindexuse{\nwixident{Thm.disch}}{Thm.disch}{NW1NjfzY-uZfts-E}\nwindexuse{\nwixident{Thm.exists{\_}elim}}{Thm.exists:unelim}{NW1NjfzY-uZfts-E}\nwindexuse{\nwixident{Thm.exists{\_}intro}}{Thm.exists:unintro}{NW1NjfzY-uZfts-E}\nwindexuse{\nwixident{Thm.t}}{Thm.t}{NW1NjfzY-uZfts-E}\nwendcode{}\nwbegindocs{129}\nwdocspar
\nwenddocs{}\nwbegincode{130}\sublabel{NW1NjfzY-3DeU7m-1}\nwmargintag{{\nwtagstyle{}\subpageref{NW1NjfzY-3DeU7m-1}}}\moddef{Unit tests for derived rules~{\nwtagstyle{}\subpageref{NW1NjfzY-3DeU7m-1}}}\endmoddef\nwstartdeflinemarkup\nwenddeflinemarkup
fun forall_cong_test () =
  let
    fun A x = Fun.mk_pred("A",[x]);
    val x = Term.Var("x",\nwlinkedidentc{Sort.IND}{NW1rTmOJ-3ZqggC-1});
    val y = Term.Var("y",\nwlinkedidentc{Sort.IND}{NW1rTmOJ-3ZqggC-1});
    val expected = \nwlinkedidentc{Formula.mk_imp}{NW3w0iRO-3SlqND-3}(\nwlinkedidentc{Formula.mk_forall}{NW3w0iRO-3SlqND-4}(x, A x),
                                  \nwlinkedidentc{Formula.mk_forall}{NW3w0iRO-3SlqND-4}(y, A y));
    val actual = \nwlinkedidentc{Thm.concl}{NW4ICsJQ-YIPBT-1}(\nwlinkedidentc{Derived.forall_cong}{NW1NjfzY-uZfts-D} x y A);
  in
    \nwlinkedidentc{Formula.eq}{NW3w0iRO-5gQzE-1} expected actual
  end;

fun exists_cong_test () =
  let
    fun A x = Fun.mk_pred("A",[x]);
    val x = Term.Var("x",\nwlinkedidentc{Sort.IND}{NW1rTmOJ-3ZqggC-1});
    val y = Term.Var("y",\nwlinkedidentc{Sort.IND}{NW1rTmOJ-3ZqggC-1});
    val expected = \nwlinkedidentc{Formula.mk_imp}{NW3w0iRO-3SlqND-3}(\nwlinkedidentc{Formula.mk_exists}{NW3w0iRO-3SlqND-4}(x, A x),
                                  \nwlinkedidentc{Formula.mk_exists}{NW3w0iRO-3SlqND-4}(y, A y));
    val actual = \nwlinkedidentc{Thm.concl}{NW4ICsJQ-YIPBT-1}(\nwlinkedidentc{Derived.exists_cong}{NW1NjfzY-uZfts-E} x y A);
  in
    \nwlinkedidentc{Formula.eq}{NW3w0iRO-5gQzE-1} expected actual
  end;
\nwnotused{Unit tests for derived rules}\nwidentuses{\\{{\nwixident{Derived.exists{\_}cong}}{Derived.exists:uncong}}\\{{\nwixident{Derived.forall{\_}cong}}{Derived.forall:uncong}}\\{{\nwixident{Formula.eq}}{Formula.eq}}\\{{\nwixident{Formula.mk{\_}exists}}{Formula.mk:unexists}}\\{{\nwixident{Formula.mk{\_}forall}}{Formula.mk:unforall}}\\{{\nwixident{Formula.mk{\_}imp}}{Formula.mk:unimp}}\\{{\nwixident{Sort.IND}}{Sort.IND}}\\{{\nwixident{Thm.concl}}{Thm.concl}}}\nwindexuse{\nwixident{Derived.exists{\_}cong}}{Derived.exists:uncong}{NW1NjfzY-3DeU7m-1}\nwindexuse{\nwixident{Derived.forall{\_}cong}}{Derived.forall:uncong}{NW1NjfzY-3DeU7m-1}\nwindexuse{\nwixident{Formula.eq}}{Formula.eq}{NW1NjfzY-3DeU7m-1}\nwindexuse{\nwixident{Formula.mk{\_}exists}}{Formula.mk:unexists}{NW1NjfzY-3DeU7m-1}\nwindexuse{\nwixident{Formula.mk{\_}forall}}{Formula.mk:unforall}{NW1NjfzY-3DeU7m-1}\nwindexuse{\nwixident{Formula.mk{\_}imp}}{Formula.mk:unimp}{NW1NjfzY-3DeU7m-1}\nwindexuse{\nwixident{Sort.IND}}{Sort.IND}{NW1NjfzY-3DeU7m-1}\nwindexuse{\nwixident{Thm.concl}}{Thm.concl}{NW1NjfzY-3DeU7m-1}\nwendcode{}\nwbegindocs{131}\nwdocspar
\nwenddocs{}\nwfilename{nw/goal.nw}\nwbegindocs{0}% -*- mode: poly-noweb; noweb-code-mode: sml-mode; -*-
%%
%% goal.nw
%% 
%% Made by Alex Nelson <pqnelson@gmail.com>
%% Login   <alex@lisp>
%% 
%% Started on  2026-04-09T15:40:57-0700
%% Last update 2026-04-09T15:40:57-0700
%% 

\section{Tactics and goals}

\begin{node}
The idea behind LCF-style provers is that we can use
inference rules as the ``smart constructors'' for abstract theorem
objects. This is great, it ensures ``correctness by construction''.
But Robin Milner~\cite{milner1984use} realized we can also program ``tactics'' to prove theorems for us.

Remember, for natural deduction proofs, it is more ``natural'' to read
the proof from the bottom up. It is also more ``natural'' to
\emph{construct} natural deduction proofs this way.

We abstract the notion of a ``proof obligation'' or ``proof goal'' as
a \define{Goal}. This is the ``bottom'' of the natural deduction
proof. 

We then describe the idea of applying inference rules ``backwards''
(i.e., from the bottom of the proof upwards) as \define{Tactics}. When
an inference rule has multiple premises, this amounts to requiring the
user prove several subgoals. We will keep a stack (in the computer
science sense of the term) of subgoals.

Goals will also be given the ``justification'' or ``validation'' function, which
essentially says ``If you give me theorems for the subgoals, I will
return to you a theorem establishing the proof obligation at this point.''
We should view this as ``going back down'' the natural deduction proof
using inference rules (or some derived inference rules).
\end{node}

\nwenddocs{}\nwbegincode{1}\sublabel{NW1F32vJ-KAFG0-1}\nwmargintag{{\nwtagstyle{}\subpageref{NW1F32vJ-KAFG0-1}}}\moddef{Goal.sig~{\nwtagstyle{}\subpageref{NW1F32vJ-KAFG0-1}}}\endmoddef\nwstartdeflinemarkup\nwenddeflinemarkup
signature Goal = sig
  exception Unsolved of string;
  exception Mismatch of string;
  
  type t = \{ goals : ((string * \nwlinkedidentc{Formula.t}{NW3w0iRO-2jnj9W-1}) list * \nwlinkedidentc{Formula.t}{NW3w0iRO-2jnj9W-1}) list
           , justify : \nwlinkedidentc{Thm.t}{NW4ICsJQ-YIPBT-1} list -> \nwlinkedidentc{Thm.t}{NW4ICsJQ-YIPBT-1} \};

  val subgoals : t -> ((string * \nwlinkedidentc{Formula.t}{NW3w0iRO-2jnj9W-1}) list * \nwlinkedidentc{Formula.t}{NW3w0iRO-2jnj9W-1}) list;

  val justify : t -> \nwlinkedidentc{Thm.t}{NW4ICsJQ-YIPBT-1} list -> \nwlinkedidentc{Thm.t}{NW4ICsJQ-YIPBT-1};
  
  val to_string : t -> string;

  val setup : \nwlinkedidentc{Formula.t}{NW3w0iRO-2jnj9W-1} -> t;

  val extract_thm : t -> \nwlinkedidentc{Thm.t}{NW4ICsJQ-YIPBT-1};
end;
\nwnotused{Goal.sig}\nwidentuses{\\{{\nwixident{Formula.t}}{Formula.t}}\\{{\nwixident{Thm.t}}{Thm.t}}}\nwindexuse{\nwixident{Formula.t}}{Formula.t}{NW1F32vJ-KAFG0-1}\nwindexuse{\nwixident{Thm.t}}{Thm.t}{NW1F32vJ-KAFG0-1}\nwendcode{}\nwbegindocs{2}\nwdocspar

\nwenddocs{}\nwbegincode{3}\sublabel{NW1F32vJ-2MKBgq-1}\nwmargintag{{\nwtagstyle{}\subpageref{NW1F32vJ-2MKBgq-1}}}\moddef{Goal.sml~{\nwtagstyle{}\subpageref{NW1F32vJ-2MKBgq-1}}}\endmoddef\nwstartdeflinemarkup\nwenddeflinemarkup
structure Goal : Goal = struct
exception Unsolved of string;
exception Mismatch of string;

type t = \{ goals : ((string * \nwlinkedidentc{Formula.t}{NW3w0iRO-2jnj9W-1}) list * \nwlinkedidentc{Formula.t}{NW3w0iRO-2jnj9W-1}) list,
           justify : \nwlinkedidentc{Thm.t}{NW4ICsJQ-YIPBT-1} list -> \nwlinkedidentc{Thm.t}{NW4ICsJQ-YIPBT-1} \};

fun subgoals (\{goals, ...\} : t) = goals;

fun justify (\{justify = jfn, ...\} : t) = jfn;

local
  fun single (asl, fm) =
    (String.concatWith
       ", "
       (map (fn (label, p) =>
                if "" = label then \nwlinkedidentc{Formula.serialize}{NW3w0iRO-3H9tUA-1} p
                else "\\""^label^"\\": "^(\nwlinkedidentc{Formula.serialize}{NW3w0iRO-3H9tUA-1} p))
            asl)) ^
    (if (List.null asl) then "|- "
     else " |- ") ^ 
    (\nwlinkedidentc{Formula.serialize}{NW3w0iRO-3H9tUA-1} fm);
  fun iter n acc [] = acc
    | iter n acc ((asl, fm)::gls) =
        iter (n + 1)
             (acc ^
              "Subgoal "^
              (Int.toString n)^
              ": " ^
              (single (asl, fm)) ^"\\n")
             gls;
in
fun to_string (\{goals, justify\}) =
 (case goals of
   [] => "No more goals"
 | [g] => (single g)
 | (gls) => (iter 1 "" gls))
end;

(* setup : \nwlinkedidentc{Formula.t}{NW3w0iRO-2jnj9W-1} -> Goal.t *)
fun setup fm =
  let
    fun check [thm] =
      if \nwlinkedidentc{Formula.eq}{NW3w0iRO-5gQzE-1} fm (\nwlinkedidentc{Thm.concl}{NW4ICsJQ-YIPBT-1} thm)
      then thm
      else raise Mismatch "Goal.setup: wrong theorem"
  in
    \{ goals = [([], fm)]
    , justify = check \}
  end;

(* extract_thm : Goal.t -> \nwlinkedidentc{Thm.t}{NW4ICsJQ-YIPBT-1} *)
fun extract_thm (\{goals, justify\} : t) = 
  if List.null goals
  then justify []
  else raise Unsolved "Goal.extract_thm: unsolved goals";
end;
\nwnotused{Goal.sml}\nwidentuses{\\{{\nwixident{Formula.eq}}{Formula.eq}}\\{{\nwixident{Formula.serialize}}{Formula.serialize}}\\{{\nwixident{Formula.t}}{Formula.t}}\\{{\nwixident{Thm.concl}}{Thm.concl}}\\{{\nwixident{Thm.t}}{Thm.t}}}\nwindexuse{\nwixident{Formula.eq}}{Formula.eq}{NW1F32vJ-2MKBgq-1}\nwindexuse{\nwixident{Formula.serialize}}{Formula.serialize}{NW1F32vJ-2MKBgq-1}\nwindexuse{\nwixident{Formula.t}}{Formula.t}{NW1F32vJ-2MKBgq-1}\nwindexuse{\nwixident{Thm.concl}}{Thm.concl}{NW1F32vJ-2MKBgq-1}\nwindexuse{\nwixident{Thm.t}}{Thm.t}{NW1F32vJ-2MKBgq-1}\nwendcode{}\nwbegindocs{4}\nwdocspar

\begin{node}
A tactic is literally an inference rule ``upside down''. For example,
conjunction-introduction is the rule
\begin{equation}
\vcenter{\infer[\land\mbox{-intro}]{\Gamma\vdash A\land B}{\Gamma\vdash A, & \Gamma\vdash B}}
\end{equation}
Then the tactic is written with a double inference line
\begin{equation}
\vcenter{\infer=[\land\mbox{-intro-tac}]{\Gamma\vdash A, \qquad \Gamma\vdash B}{\Gamma\vdash A\land B}}
\end{equation}
It begins with a goal trying to prove $A\land B$ under hypotheses
$\Gamma$, and breaks it down into two subgoals: try to prove $A$ under
those hypotheses, and try to prove $B$ under those hypotheses.
\end{node}

\nwenddocs{}\nwbegincode{5}\sublabel{NW1F32vJ-VKxLl-1}\nwmargintag{{\nwtagstyle{}\subpageref{NW1F32vJ-VKxLl-1}}}\moddef{Tactic.sig~{\nwtagstyle{}\subpageref{NW1F32vJ-VKxLl-1}}}\endmoddef\nwstartdeflinemarkup\nwenddeflinemarkup
signature Tactic = sig
  type t = Goal.t -> Goal.t;

  val prove : \nwlinkedidentc{Formula.t}{NW3w0iRO-2jnj9W-1} -> t list -> \nwlinkedidentc{Thm.t}{NW4ICsJQ-YIPBT-1};
  
  \LA{}Specification for tactics~{\nwtagstyle{}\subpageref{NW1F32vJ-2kzjE-1}}\RA{}
end;
\nwnotused{Tactic.sig}\nwidentuses{\\{{\nwixident{Formula.t}}{Formula.t}}\\{{\nwixident{Thm.t}}{Thm.t}}}\nwindexuse{\nwixident{Formula.t}}{Formula.t}{NW1F32vJ-VKxLl-1}\nwindexuse{\nwixident{Thm.t}}{Thm.t}{NW1F32vJ-VKxLl-1}\nwendcode{}\nwbegindocs{6}\nwdocspar

\nwenddocs{}\nwbegincode{7}\sublabel{NW1F32vJ-2XH0oT-1}\nwmargintag{{\nwtagstyle{}\subpageref{NW1F32vJ-2XH0oT-1}}}\moddef{Tactic.sml~{\nwtagstyle{}\subpageref{NW1F32vJ-2XH0oT-1}}}\endmoddef\nwstartdeflinemarkup\nwenddeflinemarkup
structure Tactic : Tactic = struct
type t = Goal.t -> Goal.t;

\LA{}Proving a theorem from tactics~{\nwtagstyle{}\subpageref{NW1F32vJ-2zEBzF-1}}\RA{}

\LA{}Check theorem hypotheses form subset of assumptions list~{\nwtagstyle{}\subpageref{NW1F32vJ-1bk6hJ-1}}\RA{}

\LA{}Implementation of tactics~{\nwtagstyle{}\subpageref{NW1F32vJ-3Cb3oS-1}}\RA{}

end;
\nwnotused{Tactic.sml}\nwendcode{}\nwbegindocs{8}\nwdocspar

\begin{node}
We can then prove a formula using a ``proof'' {\Tt{}pf\nwendquote}, which is just a
list of tactics. The only subtlety here is that tactics are literally
the inference rules backwards, so we need to apply the tactics in
reverse order.
\end{node}

\nwenddocs{}\nwbegincode{9}\sublabel{NW1F32vJ-2zEBzF-1}\nwmargintag{{\nwtagstyle{}\subpageref{NW1F32vJ-2zEBzF-1}}}\moddef{Proving a theorem from tactics~{\nwtagstyle{}\subpageref{NW1F32vJ-2zEBzF-1}}}\endmoddef\nwstartdeflinemarkup\nwusesondefline{\\{NW1F32vJ-2XH0oT-1}}\nwenddeflinemarkup
local
  fun tactic_proof (g : Goal.t) pf =
    Goal.extract_thm(foldr (fn (tac, g') => tac g')
                           g
                           pf);
in
fun prove (fm : \nwlinkedidentc{Formula.t}{NW3w0iRO-2jnj9W-1}) pf =
  tactic_proof (Goal.setup fm) pf
end;
\nwindexdefn{\nwixident{Tactic.prove}}{Tactic.prove}{NW1F32vJ-2zEBzF-1}\eatline
\nwused{\\{NW1F32vJ-2XH0oT-1}}\nwidentdefs{\\{{\nwixident{Tactic.prove}}{Tactic.prove}}}\nwidentuses{\\{{\nwixident{Formula.t}}{Formula.t}}}\nwindexuse{\nwixident{Formula.t}}{Formula.t}{NW1F32vJ-2zEBzF-1}\nwendcode{}\nwbegindocs{10}\nwdocspar
\begin{node}
When we the current goal requires proving $\Gamma\vdash A\land B$,
we can transform it into two subgoals: $\Gamma\vdash A$ and
$\Gamma\vdash B$.

We also need to modify the justification during this step. We are
stipulating the user will be able to produce these theorems
$\Gamma\vdash A$ and $\Gamma\vdash B$. Once obtained, we pop them from
the stack of theorems, then just use the
{\Tt{}Thm.and{\_}intro\nwendquote} inference rule to obtain $\Gamma\vdash A\land B$,
and push it back onto the stack of theorems. Then we use the
accumulated justification to prove the original claim.
\end{node}

\nwenddocs{}\nwbegincode{11}\sublabel{NW1F32vJ-2kzjE-1}\nwmargintag{{\nwtagstyle{}\subpageref{NW1F32vJ-2kzjE-1}}}\moddef{Specification for tactics~{\nwtagstyle{}\subpageref{NW1F32vJ-2kzjE-1}}}\endmoddef\nwstartdeflinemarkup\nwusesondefline{\\{NW1F32vJ-VKxLl-1}}\nwprevnextdefs{\relax}{NW1F32vJ-2kzjE-2}\nwenddeflinemarkup
  val and_intro : t;
\nwalsodefined{\\{NW1F32vJ-2kzjE-2}\\{NW1F32vJ-2kzjE-3}\\{NW1F32vJ-2kzjE-4}\\{NW1F32vJ-2kzjE-5}\\{NW1F32vJ-2kzjE-6}\\{NW1F32vJ-2kzjE-7}\\{NW1F32vJ-2kzjE-8}\\{NW1F32vJ-2kzjE-9}\\{NW1F32vJ-2kzjE-A}\\{NW1F32vJ-2kzjE-B}}\nwused{\\{NW1F32vJ-VKxLl-1}}\nwendcode{}\nwbegindocs{12}\nwdocspar

\nwenddocs{}\nwbegincode{13}\sublabel{NW1F32vJ-3Cb3oS-1}\nwmargintag{{\nwtagstyle{}\subpageref{NW1F32vJ-3Cb3oS-1}}}\moddef{Implementation of tactics~{\nwtagstyle{}\subpageref{NW1F32vJ-3Cb3oS-1}}}\endmoddef\nwstartdeflinemarkup\nwusesondefline{\\{NW1F32vJ-2XH0oT-1}}\nwprevnextdefs{\relax}{NW1F32vJ-3Cb3oS-2}\nwenddeflinemarkup
(* and_intro : Goal.t -> Goal.t *)
fun and_intro \{goals = ((asl, and_p_q)::gls), justify\} =
  let
    val (p,q) = \nwlinkedidentc{Formula.dest_and}{NW3w0iRO-1Yt1Jl-3} and_p_q;
    fun jfn' (thp::thq::ths) = justify ((Thm.and_intro thp thq)::ths);
  in
    \{ goals = (asl,p)::(asl,q)::gls
    , justify = jfn' \}
  end;
\nwindexdefn{\nwixident{Tactic.and{\_}intro}}{Tactic.and:unintro}{NW1F32vJ-3Cb3oS-1}\eatline
\nwalsodefined{\\{NW1F32vJ-3Cb3oS-2}\\{NW1F32vJ-3Cb3oS-3}\\{NW1F32vJ-3Cb3oS-4}\\{NW1F32vJ-3Cb3oS-5}\\{NW1F32vJ-3Cb3oS-6}\\{NW1F32vJ-3Cb3oS-7}\\{NW1F32vJ-3Cb3oS-8}\\{NW1F32vJ-3Cb3oS-9}\\{NW1F32vJ-3Cb3oS-A}\\{NW1F32vJ-3Cb3oS-B}}\nwused{\\{NW1F32vJ-2XH0oT-1}}\nwidentdefs{\\{{\nwixident{Tactic.and{\_}intro}}{Tactic.and:unintro}}}\nwidentuses{\\{{\nwixident{Formula.dest{\_}and}}{Formula.dest:unand}}}\nwindexuse{\nwixident{Formula.dest{\_}and}}{Formula.dest:unand}{NW1F32vJ-3Cb3oS-1}\nwendcode{}\nwbegindocs{14}\nwdocspar
\begin{node}
The $\land$-elimination inference rules, recall, look like:
\begin{equation}
\vcenter{\infer[\land\mbox{-elim-l}]{\Gamma\vdash A}{\Gamma\vdash A\land B}}
\end{equation}
\begin{equation}
\vcenter{\infer[\land\mbox{-elim-r}]{\Gamma\vdash B}{\Gamma\vdash A\land B}}
\end{equation}
The corresponding tactics would look like:
\begin{equation}
\vcenter{\infer=[\land\mbox{-elim-l-tac}(B)]{\Gamma\vdash A\land B}{\Gamma\vdash A}}
\end{equation}
\begin{equation}
\vcenter{\infer=[\land\mbox{-elim-r-tac}(A)]{\Gamma\vdash A\land B}{\Gamma\vdash B}}
\end{equation}
These would change the current goal from proving just $\Gamma\vdash A$
to some conjunction $\Gamma\vdash A\land B$, the user just needs to
specify what formula $B$ should be included in the conjunction.

Here we can start to collect dividends on placing all our tactics in
the {\Tt{}Tactic\nwendquote} structure, because we don't have to add a {\Tt{}{\_}tac\nwendquote} (or
{\Tt{}{\_}tactic\nwendquote}) suffix to these function names. We can mirror the names
used in the kernel for the inference rules.
\end{node}

\nwenddocs{}\nwbegincode{15}\sublabel{NW1F32vJ-2kzjE-2}\nwmargintag{{\nwtagstyle{}\subpageref{NW1F32vJ-2kzjE-2}}}\moddef{Specification for tactics~{\nwtagstyle{}\subpageref{NW1F32vJ-2kzjE-1}}}\plusendmoddef\nwstartdeflinemarkup\nwusesondefline{\\{NW1F32vJ-VKxLl-1}}\nwprevnextdefs{NW1F32vJ-2kzjE-1}{NW1F32vJ-2kzjE-3}\nwenddeflinemarkup
  val and_elim_l : \nwlinkedidentc{Formula.t}{NW3w0iRO-2jnj9W-1} -> t;
  val and_elim_r : \nwlinkedidentc{Formula.t}{NW3w0iRO-2jnj9W-1} -> t;
\nwused{\\{NW1F32vJ-VKxLl-1}}\nwidentuses{\\{{\nwixident{Formula.t}}{Formula.t}}}\nwindexuse{\nwixident{Formula.t}}{Formula.t}{NW1F32vJ-2kzjE-2}\nwendcode{}\nwbegindocs{16}\nwdocspar

\nwenddocs{}\nwbegincode{17}\sublabel{NW1F32vJ-3Cb3oS-2}\nwmargintag{{\nwtagstyle{}\subpageref{NW1F32vJ-3Cb3oS-2}}}\moddef{Implementation of tactics~{\nwtagstyle{}\subpageref{NW1F32vJ-3Cb3oS-1}}}\plusendmoddef\nwstartdeflinemarkup\nwusesondefline{\\{NW1F32vJ-2XH0oT-1}}\nwprevnextdefs{NW1F32vJ-3Cb3oS-1}{NW1F32vJ-3Cb3oS-3}\nwenddeflinemarkup
(* and_elim_l : \nwlinkedidentc{Formula.t}{NW3w0iRO-2jnj9W-1} -> Goal.t -> Goal.t *)
fun and_elim_l B \{goals = ((asl, A)::gls), justify\} =
  let
    val AB = \nwlinkedidentc{Formula.mk_and}{NW3w0iRO-3SlqND-3}(A, B);
    fun jfn' (thAB::ths) = justify ((Thm.and_elim_l thAB)::ths);
  in
    \{ goals = (asl,AB)::gls
    , justify = jfn' \}
  end;
(* and_elim_r : \nwlinkedidentc{Formula.t}{NW3w0iRO-2jnj9W-1} -> Goal.t -> Goal.t *)
fun and_elim_r A \{goals = ((asl, B)::gls), justify\} =
  let
    val AB = \nwlinkedidentc{Formula.mk_and}{NW3w0iRO-3SlqND-3}(A, B);
    fun jfn' (thAB::ths) = justify ((Thm.and_elim_r thAB)::ths);
  in
    \{ goals = (asl,AB)::gls
    , justify = jfn' \}
  end;
\nwindexdefn{\nwixident{Tactic.and{\_}elim{\_}l}}{Tactic.and:unelim:unl}{NW1F32vJ-3Cb3oS-2}\nwindexdefn{\nwixident{Tactic.and{\_}elim{\_}r}}{Tactic.and:unelim:unr}{NW1F32vJ-3Cb3oS-2}\eatline
\nwused{\\{NW1F32vJ-2XH0oT-1}}\nwidentdefs{\\{{\nwixident{Tactic.and{\_}elim{\_}l}}{Tactic.and:unelim:unl}}\\{{\nwixident{Tactic.and{\_}elim{\_}r}}{Tactic.and:unelim:unr}}}\nwidentuses{\\{{\nwixident{Formula.mk{\_}and}}{Formula.mk:unand}}\\{{\nwixident{Formula.t}}{Formula.t}}}\nwindexuse{\nwixident{Formula.mk{\_}and}}{Formula.mk:unand}{NW1F32vJ-3Cb3oS-2}\nwindexuse{\nwixident{Formula.t}}{Formula.t}{NW1F32vJ-3Cb3oS-2}\nwendcode{}\nwbegindocs{18}\nwdocspar
\begin{node}
We had some rules for $\Longrightarrow$-introduction ({\Tt{}disch\nwendquote}) and
elimination ({\Tt{}modus{\_}ponens\nwendquote}). For the tactic associated with
{\Tt{}disch\nwendquote}, we see it should look like
\begin{equation}
\vcenter{\infer=[\texttt{disch}(A)]{A,\Gamma\vdash B}{\Gamma\setminus\{A\}\vdash A\implies B}}
\end{equation}
We are adding $A$ as an assumption (where $A$ is determined from the
current goal $A\implies B$), so we need a label for it.

Note that negation $\neg A$ is just an abbreviation for
$A\implies\bot$ in Intuitionistic logic. So this tactic handles the
negation situation like:
\begin{equation}
\vcenter{\infer=[\mbox{disch}]{A,\Gamma\vdash\bot}{\Gamma\vdash\neg A}}
\end{equation}
This is consistent with how HOL handles negated
formulas.\footnote{Compare to HOL Light \url{https://hol-light.github.io/references/HTML/DISCH_TAC.html}
and HOL4 \url{https://hol-theorem-prover.org/trindemossen-2-helpdocs/help/Docfiles/HTML/Tactic.DISCH_TAC.html}}
\end{node}

\nwenddocs{}\nwbegincode{19}\sublabel{NW1F32vJ-2kzjE-3}\nwmargintag{{\nwtagstyle{}\subpageref{NW1F32vJ-2kzjE-3}}}\moddef{Specification for tactics~{\nwtagstyle{}\subpageref{NW1F32vJ-2kzjE-1}}}\plusendmoddef\nwstartdeflinemarkup\nwusesondefline{\\{NW1F32vJ-VKxLl-1}}\nwprevnextdefs{NW1F32vJ-2kzjE-2}{NW1F32vJ-2kzjE-4}\nwenddeflinemarkup
  val disch : string -> t;
\nwused{\\{NW1F32vJ-VKxLl-1}}\nwendcode{}\nwbegindocs{20}\nwdocspar

\nwenddocs{}\nwbegincode{21}\sublabel{NW1F32vJ-3Cb3oS-3}\nwmargintag{{\nwtagstyle{}\subpageref{NW1F32vJ-3Cb3oS-3}}}\moddef{Implementation of tactics~{\nwtagstyle{}\subpageref{NW1F32vJ-3Cb3oS-1}}}\plusendmoddef\nwstartdeflinemarkup\nwusesondefline{\\{NW1F32vJ-2XH0oT-1}}\nwprevnextdefs{NW1F32vJ-3Cb3oS-2}{NW1F32vJ-3Cb3oS-4}\nwenddeflinemarkup
fun disch (label : string) \{ goals = (asl, A_imp_B)::gls
                           , justify = jfn \} =
  if \nwlinkedidentc{Formula.is_imp}{NW3w0iRO-3NyJLv-3} A_imp_B
  then \LA{}Tactic to discharge an implication~{\nwtagstyle{}\subpageref{NW1F32vJ-IHUmU-1}}\RA{}
  else if \nwlinkedidentc{Formula.is_not}{NW3w0iRO-3NyJLv-2} A_imp_B
  then \LA{}Tactic for negation~{\nwtagstyle{}\subpageref{NW1F32vJ-1G1Cmu-1}}\RA{}
  else raise Fail "\nwlinkedidentc{Tactic.disch}{NW1F32vJ-3Cb3oS-3}"
\nwindexdefn{\nwixident{Tactic.disch}}{Tactic.disch}{NW1F32vJ-3Cb3oS-3}\eatline
\nwused{\\{NW1F32vJ-2XH0oT-1}}\nwidentdefs{\\{{\nwixident{Tactic.disch}}{Tactic.disch}}}\nwidentuses{\\{{\nwixident{Formula.is{\_}imp}}{Formula.is:unimp}}\\{{\nwixident{Formula.is{\_}not}}{Formula.is:unnot}}}\nwindexuse{\nwixident{Formula.is{\_}imp}}{Formula.is:unimp}{NW1F32vJ-3Cb3oS-3}\nwindexuse{\nwixident{Formula.is{\_}not}}{Formula.is:unnot}{NW1F32vJ-3Cb3oS-3}\nwendcode{}\nwbegindocs{22}\nwdocspar
\nwenddocs{}\nwbegincode{23}\sublabel{NW1F32vJ-1G1Cmu-1}\nwmargintag{{\nwtagstyle{}\subpageref{NW1F32vJ-1G1Cmu-1}}}\moddef{Tactic for negation~{\nwtagstyle{}\subpageref{NW1F32vJ-1G1Cmu-1}}}\endmoddef\nwstartdeflinemarkup\nwusesondefline{\\{NW1F32vJ-3Cb3oS-3}}\nwenddeflinemarkup
let
  val A = \nwlinkedidentc{Formula.dest_not}{NW3w0iRO-1Yt1Jl-2} A_imp_B;
  fun jfn' (th::thms) = jfn ((\nwlinkedidentc{Thm.disch}{NW4ICsJQ-2zKkPU-2} A (\nwlinkedidentc{Thm.not_elim}{NW4ICsJQ-2zKkPU-5} th))::thms);
in
  \{ goals = (((label, A)::asl, \nwlinkedidentc{Formula.mk_falsum}{NW3w0iRO-3SlqND-1}())::gls)
  , justify = jfn' \}
end
\nwused{\\{NW1F32vJ-3Cb3oS-3}}\nwidentuses{\\{{\nwixident{Formula.dest{\_}not}}{Formula.dest:unnot}}\\{{\nwixident{Formula.mk{\_}falsum}}{Formula.mk:unfalsum}}\\{{\nwixident{Thm.disch}}{Thm.disch}}\\{{\nwixident{Thm.not{\_}elim}}{Thm.not:unelim}}}\nwindexuse{\nwixident{Formula.dest{\_}not}}{Formula.dest:unnot}{NW1F32vJ-1G1Cmu-1}\nwindexuse{\nwixident{Formula.mk{\_}falsum}}{Formula.mk:unfalsum}{NW1F32vJ-1G1Cmu-1}\nwindexuse{\nwixident{Thm.disch}}{Thm.disch}{NW1F32vJ-1G1Cmu-1}\nwindexuse{\nwixident{Thm.not{\_}elim}}{Thm.not:unelim}{NW1F32vJ-1G1Cmu-1}\nwendcode{}\nwbegindocs{24}\nwdocspar

\nwenddocs{}\nwbegincode{25}\sublabel{NW1F32vJ-IHUmU-1}\nwmargintag{{\nwtagstyle{}\subpageref{NW1F32vJ-IHUmU-1}}}\moddef{Tactic to discharge an implication~{\nwtagstyle{}\subpageref{NW1F32vJ-IHUmU-1}}}\endmoddef\nwstartdeflinemarkup\nwusesondefline{\\{NW1F32vJ-3Cb3oS-3}}\nwenddeflinemarkup
let
  val (A,B) = \nwlinkedidentc{Formula.dest_imp}{NW3w0iRO-1Yt1Jl-3}(A_imp_B);
  fun jfn' (th::thms) = jfn ((\nwlinkedidentc{Thm.disch}{NW4ICsJQ-2zKkPU-2} A th)::thms);
in
  \{ goals = (((label, A)::asl, B)::gls)
  , justify = jfn' \}
end
\nwused{\\{NW1F32vJ-3Cb3oS-3}}\nwidentuses{\\{{\nwixident{Formula.dest{\_}imp}}{Formula.dest:unimp}}\\{{\nwixident{Thm.disch}}{Thm.disch}}}\nwindexuse{\nwixident{Formula.dest{\_}imp}}{Formula.dest:unimp}{NW1F32vJ-IHUmU-1}\nwindexuse{\nwixident{Thm.disch}}{Thm.disch}{NW1F32vJ-IHUmU-1}\nwendcode{}\nwbegindocs{26}\nwdocspar

\begin{node}
The {\Tt{}modus{\_}ponens\nwendquote} inference rule,
\begin{equation}
\vcenter{\infer[\texttt{modus\_ponens}]{\Gamma_{1}\cup\Gamma_{2}\vdash B}{\Gamma_{1}\vdash A\implies B,&\Gamma_{2}\vdash A}}
\end{equation}
has a corresponding tactic
\begin{equation}
\vcenter{\infer=[\texttt{modus\_ponens}(A)]{\Gamma\vdash A\implies B,\qquad\Gamma\vdash A}{\Gamma\vdash B}}
\end{equation}
This changes the goal from trying to prove $\Gamma\vdash B$ to two
subgoals $\Gamma\vdash A\implies B$ and $\Gamma\vdash A$.
\end{node}

\nwenddocs{}\nwbegincode{27}\sublabel{NW1F32vJ-2kzjE-4}\nwmargintag{{\nwtagstyle{}\subpageref{NW1F32vJ-2kzjE-4}}}\moddef{Specification for tactics~{\nwtagstyle{}\subpageref{NW1F32vJ-2kzjE-1}}}\plusendmoddef\nwstartdeflinemarkup\nwusesondefline{\\{NW1F32vJ-VKxLl-1}}\nwprevnextdefs{NW1F32vJ-2kzjE-3}{NW1F32vJ-2kzjE-5}\nwenddeflinemarkup
  val modus_ponens : \nwlinkedidentc{Formula.t}{NW3w0iRO-2jnj9W-1} -> t;
\nwused{\\{NW1F32vJ-VKxLl-1}}\nwidentuses{\\{{\nwixident{Formula.t}}{Formula.t}}}\nwindexuse{\nwixident{Formula.t}}{Formula.t}{NW1F32vJ-2kzjE-4}\nwendcode{}\nwbegindocs{28}\nwdocspar

\nwenddocs{}\nwbegincode{29}\sublabel{NW1F32vJ-3Cb3oS-4}\nwmargintag{{\nwtagstyle{}\subpageref{NW1F32vJ-3Cb3oS-4}}}\moddef{Implementation of tactics~{\nwtagstyle{}\subpageref{NW1F32vJ-3Cb3oS-1}}}\plusendmoddef\nwstartdeflinemarkup\nwusesondefline{\\{NW1F32vJ-2XH0oT-1}}\nwprevnextdefs{NW1F32vJ-3Cb3oS-3}{NW1F32vJ-3Cb3oS-5}\nwenddeflinemarkup
fun modus_ponens A \{goals = (asl, B)::gls, justify = jfn\} =
  let
    val A_imp_B = \nwlinkedidentc{Formula.mk_imp}{NW3w0iRO-3SlqND-3}(A, B);
    fun jfn' (th_A_imp_B::th_A::thms) =
      jfn ((\nwlinkedidentc{Thm.modus_ponens}{NW4ICsJQ-2zKkPU-2} th_A_imp_B th_A)::thms);
  in
    \{ goals = (asl, A_imp_B)::(asl, A)::gls
    , justify = jfn' \}
  end;
\nwindexdefn{\nwixident{Tactic.modus{\_}ponens}}{Tactic.modus:unponens}{NW1F32vJ-3Cb3oS-4}\eatline
\nwused{\\{NW1F32vJ-2XH0oT-1}}\nwidentdefs{\\{{\nwixident{Tactic.modus{\_}ponens}}{Tactic.modus:unponens}}}\nwidentuses{\\{{\nwixident{Formula.mk{\_}imp}}{Formula.mk:unimp}}\\{{\nwixident{Thm.modus{\_}ponens}}{Thm.modus:unponens}}}\nwindexuse{\nwixident{Formula.mk{\_}imp}}{Formula.mk:unimp}{NW1F32vJ-3Cb3oS-4}\nwindexuse{\nwixident{Thm.modus{\_}ponens}}{Thm.modus:unponens}{NW1F32vJ-3Cb3oS-4}\nwendcode{}\nwbegindocs{30}\nwdocspar
\begin{node}\label{node:tactic:MP}
Well, \textit{modus ponens} can be a little annoying to use,
especially if you have a theorem handy which proves $\Gamma'\vdash A$,
so HOL provers typically use a modified variant of the tactic:
\begin{equation*}
\vcenter{\infer=[\mbox{MP}(\Gamma'\vdash A)]{\Gamma\vdash A\implies B}{\Gamma\vdash B}}
\end{equation*}
where $\Gamma'\subset\Gamma$ (and the tactic fails otherwise).
\end{node}

\nwenddocs{}\nwbegincode{31}\sublabel{NW1F32vJ-2kzjE-5}\nwmargintag{{\nwtagstyle{}\subpageref{NW1F32vJ-2kzjE-5}}}\moddef{Specification for tactics~{\nwtagstyle{}\subpageref{NW1F32vJ-2kzjE-1}}}\plusendmoddef\nwstartdeflinemarkup\nwusesondefline{\\{NW1F32vJ-VKxLl-1}}\nwprevnextdefs{NW1F32vJ-2kzjE-4}{NW1F32vJ-2kzjE-6}\nwenddeflinemarkup
val mp : \nwlinkedidentc{Thm.t}{NW4ICsJQ-YIPBT-1} -> t;
\nwused{\\{NW1F32vJ-VKxLl-1}}\nwidentuses{\\{{\nwixident{Thm.t}}{Thm.t}}}\nwindexuse{\nwixident{Thm.t}}{Thm.t}{NW1F32vJ-2kzjE-5}\nwendcode{}\nwbegindocs{32}\nwdocspar

\nwenddocs{}\nwbegincode{33}\sublabel{NW1F32vJ-3Cb3oS-5}\nwmargintag{{\nwtagstyle{}\subpageref{NW1F32vJ-3Cb3oS-5}}}\moddef{Implementation of tactics~{\nwtagstyle{}\subpageref{NW1F32vJ-3Cb3oS-1}}}\plusendmoddef\nwstartdeflinemarkup\nwusesondefline{\\{NW1F32vJ-2XH0oT-1}}\nwprevnextdefs{NW1F32vJ-3Cb3oS-4}{NW1F32vJ-3Cb3oS-6}\nwenddeflinemarkup
fun mp thm \{goals = (asl, B)::gls, justify = jfn\} =
  if not(\nwlinkedidentc{assumptions_contain_hypotheses}{NW1F32vJ-1bk6hJ-1} thm asl)
  then raise Fail("mp: assumptions do not contain all hypotheses of theorem")
  else
    let
      val A = \nwlinkedidentc{Thm.concl}{NW4ICsJQ-YIPBT-1} thm;
      val A_imp_B = \nwlinkedidentc{Formula.mk_imp}{NW3w0iRO-3SlqND-3}(A, B);
      fun jfn' (th_A_imp_B::thms) =
        jfn ((\nwlinkedidentc{Thm.modus_ponens}{NW4ICsJQ-2zKkPU-2} th_A_imp_B thm)::thms);
    in
      \{ goals = (asl, A_imp_B)::gls
      , justify = jfn' \}
    end;
\nwindexdefn{\nwixident{Tactic.mp}}{Tactic.mp}{NW1F32vJ-3Cb3oS-5}\eatline
\nwused{\\{NW1F32vJ-2XH0oT-1}}\nwidentdefs{\\{{\nwixident{Tactic.mp}}{Tactic.mp}}}\nwidentuses{\\{{\nwixident{assumptions{\_}contain{\_}hypotheses}}{assumptions:uncontain:unhypotheses}}\\{{\nwixident{Formula.mk{\_}imp}}{Formula.mk:unimp}}\\{{\nwixident{Thm.concl}}{Thm.concl}}\\{{\nwixident{Thm.modus{\_}ponens}}{Thm.modus:unponens}}}\nwindexuse{\nwixident{assumptions{\_}contain{\_}hypotheses}}{assumptions:uncontain:unhypotheses}{NW1F32vJ-3Cb3oS-5}\nwindexuse{\nwixident{Formula.mk{\_}imp}}{Formula.mk:unimp}{NW1F32vJ-3Cb3oS-5}\nwindexuse{\nwixident{Thm.concl}}{Thm.concl}{NW1F32vJ-3Cb3oS-5}\nwindexuse{\nwixident{Thm.modus{\_}ponens}}{Thm.modus:unponens}{NW1F32vJ-3Cb3oS-5}\nwendcode{}\nwbegindocs{34}\nwdocspar
\begin{node}
Disjunction elimination is a bit tricky as an inference rule,
\begin{equation}
\vcenter{\infer[\lor\mbox{-elim}]{\Gamma_{1}\cup\Gamma_{2}\vdash(A\lor B)\implies C}{\Gamma_{1}\vdash A\implies C,&\Gamma_{2}\vdash B\implies C}}
\end{equation}
But as a tactic, we can transform it to:
\begin{equation}
\vcenter{\infer=[\lor\mbox{-elim}?]{\Gamma_{1}\vdash A\implies C,\quad\Gamma_{2}\vdash B\implies C}{\Gamma_{1}\cup\Gamma_{2}\vdash(A\lor B)\implies C}}
\end{equation}
In practice, we often use the modified inference rule
\begin{equation}
\vcenter{\infer[\lor\mbox{-elim}]{\Gamma\vdash C}{\Gamma\vdash A\lor B,&A,\Gamma\vdash C,&B,\Gamma\vdash C}}
\end{equation}
This just uses {\Tt{}\nwlinkedidentq{Thm.undisch}{NW4ICsJQ-2zKkPU-2}\nwendquote} on the premises of the usual
$\lor$-elimination, and adds a premise for the disjunction. Then using
\textit{modus ponens} and the usual $\lor$-elimination, we obtain our
new version of $\lor$-elimination. The tactic is then
\begin{equation}
\vcenter{\infer=[\lor\mbox{-elim}(\Gamma\vdash A\lor B)]{\Gamma\vdash A\implies C,\quad\Gamma\vdash B\implies C}{\Gamma\vdash C}}
\end{equation}
which matches what can be found in, e.g., HOL.\footnote{See, e.g., HOL
Light \url{https://hol-light.github.io/references/HTML/DISJ_CASES_TAC.html};
HOL4 \url{https://hol-theorem-prover.org/trindemossen-2-helpdocs/help/Docfiles/HTML/Tactic.DISJ_CASES_TAC.html}}

Why do this? Well, it matches how Mathematicians use
$\lor$-elimination in practice (as a {\Tt{}per\ cases\nwendquote} argument).
\end{node}

\nwenddocs{}\nwbegincode{35}\sublabel{NW1F32vJ-2kzjE-6}\nwmargintag{{\nwtagstyle{}\subpageref{NW1F32vJ-2kzjE-6}}}\moddef{Specification for tactics~{\nwtagstyle{}\subpageref{NW1F32vJ-2kzjE-1}}}\plusendmoddef\nwstartdeflinemarkup\nwusesondefline{\\{NW1F32vJ-VKxLl-1}}\nwprevnextdefs{NW1F32vJ-2kzjE-5}{NW1F32vJ-2kzjE-7}\nwenddeflinemarkup
  val disj_cases : \nwlinkedidentc{Thm.t}{NW4ICsJQ-YIPBT-1} -> t;
\nwused{\\{NW1F32vJ-VKxLl-1}}\nwidentuses{\\{{\nwixident{Thm.t}}{Thm.t}}}\nwindexuse{\nwixident{Thm.t}}{Thm.t}{NW1F32vJ-2kzjE-6}\nwendcode{}\nwbegindocs{36}\nwdocspar

\nwenddocs{}\nwbegincode{37}\sublabel{NW1F32vJ-3Cb3oS-6}\nwmargintag{{\nwtagstyle{}\subpageref{NW1F32vJ-3Cb3oS-6}}}\moddef{Implementation of tactics~{\nwtagstyle{}\subpageref{NW1F32vJ-3Cb3oS-1}}}\plusendmoddef\nwstartdeflinemarkup\nwusesondefline{\\{NW1F32vJ-2XH0oT-1}}\nwprevnextdefs{NW1F32vJ-3Cb3oS-5}{NW1F32vJ-3Cb3oS-7}\nwenddeflinemarkup
fun disj_cases thm \{goals = (asl, C)::gls, justify = jfn\} =
  let
    val A_or_B = \nwlinkedidentc{Thm.concl}{NW4ICsJQ-YIPBT-1} thm;
    val (A,B) = \nwlinkedidentc{Formula.dest_or}{NW3w0iRO-1Yt1Jl-3} A_or_B
                handle _ => raise Fail "\nwlinkedidentc{Tactic.disj_cases}{NW1F32vJ-3Cb3oS-6}: theorem is not a disjunction";
    val A_imp_C = \nwlinkedidentc{Formula.mk_imp}{NW3w0iRO-3SlqND-3}(A,C);
    val B_imp_C = \nwlinkedidentc{Formula.mk_imp}{NW3w0iRO-3SlqND-3}(B,C);
    fun jfn' (th_A_imp_C::th_B_imp_C::thms) =
      jfn ((\nwlinkedidentc{Thm.modus_ponens}{NW4ICsJQ-2zKkPU-2} (\nwlinkedidentc{Thm.or_elim}{NW4ICsJQ-2zKkPU-4} th_A_imp_C th_B_imp_C) thm)::thms);
  in
    \{ goals = (asl, A_imp_C)::(asl, B_imp_C)::gls
    , justify = jfn' \}
  end;
\nwindexdefn{\nwixident{Tactic.disj{\_}cases}}{Tactic.disj:uncases}{NW1F32vJ-3Cb3oS-6}\eatline
\nwused{\\{NW1F32vJ-2XH0oT-1}}\nwidentdefs{\\{{\nwixident{Tactic.disj{\_}cases}}{Tactic.disj:uncases}}}\nwidentuses{\\{{\nwixident{Formula.dest{\_}or}}{Formula.dest:unor}}\\{{\nwixident{Formula.mk{\_}imp}}{Formula.mk:unimp}}\\{{\nwixident{Thm.concl}}{Thm.concl}}\\{{\nwixident{Thm.modus{\_}ponens}}{Thm.modus:unponens}}\\{{\nwixident{Thm.or{\_}elim}}{Thm.or:unelim}}}\nwindexuse{\nwixident{Formula.dest{\_}or}}{Formula.dest:unor}{NW1F32vJ-3Cb3oS-6}\nwindexuse{\nwixident{Formula.mk{\_}imp}}{Formula.mk:unimp}{NW1F32vJ-3Cb3oS-6}\nwindexuse{\nwixident{Thm.concl}}{Thm.concl}{NW1F32vJ-3Cb3oS-6}\nwindexuse{\nwixident{Thm.modus{\_}ponens}}{Thm.modus:unponens}{NW1F32vJ-3Cb3oS-6}\nwindexuse{\nwixident{Thm.or{\_}elim}}{Thm.or:unelim}{NW1F32vJ-3Cb3oS-6}\nwendcode{}\nwbegindocs{38}\nwdocspar
\begin{node}
We also have two tactics corresponding to the disjunction introduction
rules
\begin{equation}
\vcenter{\infer=[\lor\mbox{-intro-l}]{\Gamma\vdash A}{\Gamma\vdash A\lor B}}
\end{equation}
\begin{equation}
\vcenter{\infer=[\lor\mbox{-intro-r}]{\Gamma\vdash B}{\Gamma\vdash A\lor B}}
\end{equation}
\end{node}
% https://hol-light.github.io/references/HTML/DISJ1_TAC.html
% https://hol-theorem-prover.org/trindemossen-2-helpdocs/help/Docfiles/HTML/Tactic.DISJ1_TAC.html

\nwenddocs{}\nwbegincode{39}\sublabel{NW1F32vJ-2kzjE-7}\nwmargintag{{\nwtagstyle{}\subpageref{NW1F32vJ-2kzjE-7}}}\moddef{Specification for tactics~{\nwtagstyle{}\subpageref{NW1F32vJ-2kzjE-1}}}\plusendmoddef\nwstartdeflinemarkup\nwusesondefline{\\{NW1F32vJ-VKxLl-1}}\nwprevnextdefs{NW1F32vJ-2kzjE-6}{NW1F32vJ-2kzjE-8}\nwenddeflinemarkup
val or_l : t;
val or_r : t;
\nwused{\\{NW1F32vJ-VKxLl-1}}\nwendcode{}\nwbegindocs{40}\nwdocspar

\nwenddocs{}\nwbegincode{41}\sublabel{NW1F32vJ-3Cb3oS-7}\nwmargintag{{\nwtagstyle{}\subpageref{NW1F32vJ-3Cb3oS-7}}}\moddef{Implementation of tactics~{\nwtagstyle{}\subpageref{NW1F32vJ-3Cb3oS-1}}}\plusendmoddef\nwstartdeflinemarkup\nwusesondefline{\\{NW1F32vJ-2XH0oT-1}}\nwprevnextdefs{NW1F32vJ-3Cb3oS-6}{NW1F32vJ-3Cb3oS-8}\nwenddeflinemarkup
fun or_l \{goals = (asl, A_or_B)::gls, justify = jfn\} =
  let
    val (A,B) = \nwlinkedidentc{Formula.dest_or}{NW3w0iRO-1Yt1Jl-3} A_or_B;
    fun jfn' (th_A::thms) =
      jfn ((\nwlinkedidentc{Thm.or_intro_r}{NW4ICsJQ-2zKkPU-4} B th_A)::thms);
  in
    \{ goals = (asl, A)::gls
    , justify = jfn' \}
  end;

fun or_r \{goals = (asl, A_or_B)::gls, justify = jfn\} =
  let
    val (A,B) = \nwlinkedidentc{Formula.dest_or}{NW3w0iRO-1Yt1Jl-3} A_or_B;
    fun jfn' (th_B::thms) =
      jfn ((\nwlinkedidentc{Thm.or_intro_l}{NW4ICsJQ-2zKkPU-4} A th_B)::thms);
  in
    \{ goals = (asl, B)::gls
    , justify = jfn' \}
  end;
\nwindexdefn{\nwixident{Tactic.or{\_}l}}{Tactic.or:unl}{NW1F32vJ-3Cb3oS-7}\nwindexdefn{\nwixident{Tactic.or{\_}r}}{Tactic.or:unr}{NW1F32vJ-3Cb3oS-7}\eatline
\nwused{\\{NW1F32vJ-2XH0oT-1}}\nwidentdefs{\\{{\nwixident{Tactic.or{\_}l}}{Tactic.or:unl}}\\{{\nwixident{Tactic.or{\_}r}}{Tactic.or:unr}}}\nwidentuses{\\{{\nwixident{Formula.dest{\_}or}}{Formula.dest:unor}}\\{{\nwixident{Thm.or{\_}intro{\_}l}}{Thm.or:unintro:unl}}\\{{\nwixident{Thm.or{\_}intro{\_}r}}{Thm.or:unintro:unr}}}\nwindexuse{\nwixident{Formula.dest{\_}or}}{Formula.dest:unor}{NW1F32vJ-3Cb3oS-7}\nwindexuse{\nwixident{Thm.or{\_}intro{\_}l}}{Thm.or:unintro:unl}{NW1F32vJ-3Cb3oS-7}\nwindexuse{\nwixident{Thm.or{\_}intro{\_}r}}{Thm.or:unintro:unr}{NW1F32vJ-3Cb3oS-7}\nwendcode{}\nwbegindocs{42}\nwdocspar
% contr
\begin{node}
The explosion principle as a tactic.
\begin{equation}
\vcenter{\infer=[\mbox{contr}(\Gamma_{0}\vdash\bot)]{}{\Gamma\vdash A}}
\end{equation}
where $\Gamma_{0}\subset\Gamma$. If $\Gamma_{0}\not\subset\Gamma$,
then this tactic will fail. Also, if the theorem supplied does not
prove $\Gamma\vdash\bot$, then this tactic will fail
immediately---this is strictly speaking not necessary, because
{\Tt{}\nwlinkedidentq{Thm.contr}{NW4ICsJQ-2zKkPU-6}\nwendquote} will fail, but that will happen when trying to rebuild
the proof (and it will seem random to the user).

This will solve the current subgoal, popping it off the stack of
subgoals.
% CONTR_TAC
% https://hol-theorem-prover.org/trindemossen-2-helpdocs/help/Docfiles/HTML/Tactic.CONTR_TAC.html
\end{node}

\nwenddocs{}\nwbegincode{43}\sublabel{NW1F32vJ-2kzjE-8}\nwmargintag{{\nwtagstyle{}\subpageref{NW1F32vJ-2kzjE-8}}}\moddef{Specification for tactics~{\nwtagstyle{}\subpageref{NW1F32vJ-2kzjE-1}}}\plusendmoddef\nwstartdeflinemarkup\nwusesondefline{\\{NW1F32vJ-VKxLl-1}}\nwprevnextdefs{NW1F32vJ-2kzjE-7}{NW1F32vJ-2kzjE-9}\nwenddeflinemarkup
val contr : \nwlinkedidentc{Thm.t}{NW4ICsJQ-YIPBT-1} -> t;
\nwused{\\{NW1F32vJ-VKxLl-1}}\nwidentuses{\\{{\nwixident{Thm.t}}{Thm.t}}}\nwindexuse{\nwixident{Thm.t}}{Thm.t}{NW1F32vJ-2kzjE-8}\nwendcode{}\nwbegindocs{44}\nwdocspar

\nwenddocs{}\nwbegincode{45}\sublabel{NW1F32vJ-3Cb3oS-8}\nwmargintag{{\nwtagstyle{}\subpageref{NW1F32vJ-3Cb3oS-8}}}\moddef{Implementation of tactics~{\nwtagstyle{}\subpageref{NW1F32vJ-3Cb3oS-1}}}\plusendmoddef\nwstartdeflinemarkup\nwusesondefline{\\{NW1F32vJ-2XH0oT-1}}\nwprevnextdefs{NW1F32vJ-3Cb3oS-7}{NW1F32vJ-3Cb3oS-9}\nwenddeflinemarkup
fun contr thm \{goals = (asl, A)::gls, justify = jfn\} =
  if not(\nwlinkedidentc{assumptions_contain_hypotheses}{NW1F32vJ-1bk6hJ-1} thm asl)
  then raise Fail("contr: assumptions do not contain all hypotheses of theorem")
  else if not(\nwlinkedidentc{Formula.is_falsum}{NW3w0iRO-3NyJLv-1}(\nwlinkedidentc{Thm.concl}{NW4ICsJQ-YIPBT-1} thm))
  then raise Fail("contr: theorem is not a contradiction")
  else
    let
      fun jfn' thms =
        jfn ((\nwlinkedidentc{Thm.contr}{NW4ICsJQ-2zKkPU-6} A thm)::thms);
    in
      \{ goals = gls
      , justify = jfn' \}
    end;
\nwindexdefn{\nwixident{Tactic.contr}}{Tactic.contr}{NW1F32vJ-3Cb3oS-8}\eatline
\nwused{\\{NW1F32vJ-2XH0oT-1}}\nwidentdefs{\\{{\nwixident{Tactic.contr}}{Tactic.contr}}}\nwidentuses{\\{{\nwixident{assumptions{\_}contain{\_}hypotheses}}{assumptions:uncontain:unhypotheses}}\\{{\nwixident{Formula.is{\_}falsum}}{Formula.is:unfalsum}}\\{{\nwixident{Thm.concl}}{Thm.concl}}\\{{\nwixident{Thm.contr}}{Thm.contr}}}\nwindexuse{\nwixident{assumptions{\_}contain{\_}hypotheses}}{assumptions:uncontain:unhypotheses}{NW1F32vJ-3Cb3oS-8}\nwindexuse{\nwixident{Formula.is{\_}falsum}}{Formula.is:unfalsum}{NW1F32vJ-3Cb3oS-8}\nwindexuse{\nwixident{Thm.concl}}{Thm.concl}{NW1F32vJ-3Cb3oS-8}\nwindexuse{\nwixident{Thm.contr}}{Thm.contr}{NW1F32vJ-3Cb3oS-8}\nwendcode{}\nwbegindocs{46}\nwdocspar
\begin{node}
Checking if the hypotheses of a theorem form a subset of the
assumptions list is a straightforward (if rather boring) task.
You just check each hypothesis is syntactically equal to some
assumption, and iterate over the hypotheses of the theorem.
\end{node}

\nwenddocs{}\nwbegincode{47}\sublabel{NW1F32vJ-1bk6hJ-1}\nwmargintag{{\nwtagstyle{}\subpageref{NW1F32vJ-1bk6hJ-1}}}\moddef{Check theorem hypotheses form subset of assumptions list~{\nwtagstyle{}\subpageref{NW1F32vJ-1bk6hJ-1}}}\endmoddef\nwstartdeflinemarkup\nwusesondefline{\\{NW1F32vJ-2XH0oT-1}}\nwenddeflinemarkup
local
  fun iter [] _ = true
    | iter (h::hs) asl =
        (List.exists (\nwlinkedidentc{Formula.eq}{NW3w0iRO-5gQzE-1} h) asl) andalso
        iter hs asl;
in
fun \nwlinkedidentc{assumptions_contain_hypotheses}{NW1F32vJ-1bk6hJ-1} (thm : \nwlinkedidentc{Thm.t}{NW4ICsJQ-YIPBT-1}) (asl : (string * \nwlinkedidentc{Formula.t}{NW3w0iRO-2jnj9W-1}) list) =
  iter (\nwlinkedidentc{Thm.hyps}{NW4ICsJQ-YIPBT-1} thm) (map (fn (_,fm) => fm) asl)
end;
\nwindexdefn{\nwixident{assumptions{\_}contain{\_}hypotheses}}{assumptions:uncontain:unhypotheses}{NW1F32vJ-1bk6hJ-1}\eatline
\nwused{\\{NW1F32vJ-2XH0oT-1}}\nwidentdefs{\\{{\nwixident{assumptions{\_}contain{\_}hypotheses}}{assumptions:uncontain:unhypotheses}}}\nwidentuses{\\{{\nwixident{Formula.eq}}{Formula.eq}}\\{{\nwixident{Formula.t}}{Formula.t}}\\{{\nwixident{Thm.hyps}}{Thm.hyps}}\\{{\nwixident{Thm.t}}{Thm.t}}}\nwindexuse{\nwixident{Formula.eq}}{Formula.eq}{NW1F32vJ-1bk6hJ-1}\nwindexuse{\nwixident{Formula.t}}{Formula.t}{NW1F32vJ-1bk6hJ-1}\nwindexuse{\nwixident{Thm.hyps}}{Thm.hyps}{NW1F32vJ-1bk6hJ-1}\nwindexuse{\nwixident{Thm.t}}{Thm.t}{NW1F32vJ-1bk6hJ-1}\nwendcode{}\nwbegindocs{48}\nwdocspar
\begin{node}
We can transform the subgoal $\Gamma\vdash A\iff B$ into two subgoals
$A\implies B$ and $B\implies A$,
\begin{equation}
\vcenter{\infer=[\mbox{iff}]{\Gamma\vdash A\implies B,\qquad\Gamma\vdash B\implies A}{\Gamma\vdash A\iff B}}
\end{equation}
The reverse tactic has a subgoal be $\Gamma\vdash B$ and then for any
formula $A$ transforms it to $\Gamma\vdash A\iff B$,
\begin{equation}
\vcenter{\infer=[\mbox{iff-mp}(A)]{\Gamma\vdash A\iff B}{\Gamma\vdash B}}
\end{equation}
I'm not certain how useful this second tactic would be.
% https://hol-light.github.io/references/HTML/EQ_TAC.html
% https://hol-theorem-prover.org/trindemossen-2-helpdocs/help/Docfiles/HTML/Tactic.EQ_TAC.html
\end{node}

\nwenddocs{}\nwbegincode{49}\sublabel{NW1F32vJ-2kzjE-9}\nwmargintag{{\nwtagstyle{}\subpageref{NW1F32vJ-2kzjE-9}}}\moddef{Specification for tactics~{\nwtagstyle{}\subpageref{NW1F32vJ-2kzjE-1}}}\plusendmoddef\nwstartdeflinemarkup\nwusesondefline{\\{NW1F32vJ-VKxLl-1}}\nwprevnextdefs{NW1F32vJ-2kzjE-8}{NW1F32vJ-2kzjE-A}\nwenddeflinemarkup
val iff : t;
\nwused{\\{NW1F32vJ-VKxLl-1}}\nwendcode{}\nwbegindocs{50}\nwdocspar

\nwenddocs{}\nwbegincode{51}\sublabel{NW1F32vJ-3Cb3oS-9}\nwmargintag{{\nwtagstyle{}\subpageref{NW1F32vJ-3Cb3oS-9}}}\moddef{Implementation of tactics~{\nwtagstyle{}\subpageref{NW1F32vJ-3Cb3oS-1}}}\plusendmoddef\nwstartdeflinemarkup\nwusesondefline{\\{NW1F32vJ-2XH0oT-1}}\nwprevnextdefs{NW1F32vJ-3Cb3oS-8}{NW1F32vJ-3Cb3oS-A}\nwenddeflinemarkup
fun iff \{goals = (asl, A_iff_B)::gls, justify = jfn\} =
  let
    val (A,B) = \nwlinkedidentc{Formula.dest_iff}{NW3w0iRO-1Yt1Jl-3} A_iff_B;
    val A_imp_B = \nwlinkedidentc{Formula.mk_imp}{NW3w0iRO-3SlqND-3}(A, B);
    val B_imp_A = \nwlinkedidentc{Formula.mk_imp}{NW3w0iRO-3SlqND-3}(B, A);
    fun jfn' (th_A_imp_B::th_B_imp_A::thms) =
      jfn((\nwlinkedidentc{Thm.iff_intro}{NW4ICsJQ-2zKkPU-7} th_A_imp_B th_B_imp_A)::thms);
  in
    \{ goals = (asl, A_imp_B)::(asl, B_imp_A)::gls
    , justify = jfn' \}
  end;
\nwindexdefn{\nwixident{Thm.iff}}{Thm.iff}{NW1F32vJ-3Cb3oS-9}\eatline
\nwused{\\{NW1F32vJ-2XH0oT-1}}\nwidentdefs{\\{{\nwixident{Thm.iff}}{Thm.iff}}}\nwidentuses{\\{{\nwixident{Formula.dest{\_}iff}}{Formula.dest:uniff}}\\{{\nwixident{Formula.mk{\_}imp}}{Formula.mk:unimp}}\\{{\nwixident{Thm.iff{\_}intro}}{Thm.iff:unintro}}}\nwindexuse{\nwixident{Formula.dest{\_}iff}}{Formula.dest:uniff}{NW1F32vJ-3Cb3oS-9}\nwindexuse{\nwixident{Formula.mk{\_}imp}}{Formula.mk:unimp}{NW1F32vJ-3Cb3oS-9}\nwindexuse{\nwixident{Thm.iff{\_}intro}}{Thm.iff:unintro}{NW1F32vJ-3Cb3oS-9}\nwendcode{}\nwbegindocs{52}\nwdocspar
\begin{node}
The last tactics which need discussing concern themselves with quantifiers.
We will first focus on universal quantifiers.

When the current goal is $\Gamma\vdash\forall x\ldotp A$, we find a
fresh variable $x'$ (usually $x$) for $x$ and transform the goal to
$\Gamma\vdash A[x'/x]$. This is the reverse of $\forall$-introduction,
but since it proves the ``generalization'' of a claim, we'll refer to
it as the \define{Generalization Tactic}:
\begin{equation}
\vcenter{\infer=[\mbox{gen}]{\Gamma\vdash A[x'/x]}{\Gamma\vdash\forall x\ldotp A}}
\end{equation}
Conversely, if we want to transform the current goal $\Gamma\vdash A$
into $\Gamma\vdash\forall x\ldotp A[x/t]$ where $t$ is a term and $x$
is a variable, then we're really applying the specialization inference
rule ($\forall$-elimination).
\begin{equation}
\vcenter{\infer=[\mbox{spec}(t,x)]{\Gamma\vdash\forall x\ldotp A[x/t]}{\Gamma\vdash A}}
\end{equation}
Note that ``spec'' fails when $t$ and $x$ are of different sorts.
\end{node}
% https://hol-light.github.io/references/HTML/SPEC_TAC.html

\nwenddocs{}\nwbegincode{53}\sublabel{NW1F32vJ-2kzjE-A}\nwmargintag{{\nwtagstyle{}\subpageref{NW1F32vJ-2kzjE-A}}}\moddef{Specification for tactics~{\nwtagstyle{}\subpageref{NW1F32vJ-2kzjE-1}}}\plusendmoddef\nwstartdeflinemarkup\nwusesondefline{\\{NW1F32vJ-VKxLl-1}}\nwprevnextdefs{NW1F32vJ-2kzjE-9}{NW1F32vJ-2kzjE-B}\nwenddeflinemarkup
val gen : t;
val spec : (\nwlinkedidentc{Term.t}{NW1rTmOJ-4LsRWC-1} * \nwlinkedidentc{Term.t}{NW1rTmOJ-4LsRWC-1}) -> t;
\nwused{\\{NW1F32vJ-VKxLl-1}}\nwidentuses{\\{{\nwixident{Term.t}}{Term.t}}}\nwindexuse{\nwixident{Term.t}}{Term.t}{NW1F32vJ-2kzjE-A}\nwendcode{}\nwbegindocs{54}\nwdocspar

\nwenddocs{}\nwbegincode{55}\sublabel{NW1F32vJ-3Cb3oS-A}\nwmargintag{{\nwtagstyle{}\subpageref{NW1F32vJ-3Cb3oS-A}}}\moddef{Implementation of tactics~{\nwtagstyle{}\subpageref{NW1F32vJ-3Cb3oS-1}}}\plusendmoddef\nwstartdeflinemarkup\nwusesondefline{\\{NW1F32vJ-2XH0oT-1}}\nwprevnextdefs{NW1F32vJ-3Cb3oS-9}{NW1F32vJ-3Cb3oS-B}\nwenddeflinemarkup
fun gen \{goals = (asl, forall_A)::gls, justify = jfn\} =
  let
    val (x as (Term.Var(x0,s0)),A) = Formula.dest_forall forall_A;
    val fvs = List.concat (map (fn (_,fm) => \nwlinkedidentc{Formula.fv}{NW3w0iRO-jRgYZ-2} fm) asl);
    val x' = \nwlinkedidentc{Term.fresh_var}{NW1rTmOJ-4Ykq9M-1} x0 s0 fvs;
    fun jfn' (th_A::thms) = jfn((\nwlinkedidentc{Thm.forall_intro}{NW4ICsJQ-2zKkPU-8} x' th_A)::thms);
  in
    \{ goals = (asl, \nwlinkedidentc{Formula.subst}{NW3w0iRO-1jtQRV-1} x' x A)::gls
    , justify = jfn' \}
  end;

fun spec (t, x) \{goals = (asl, A)::gls, justify = jfn\} =
  let
    val A' = \nwlinkedidentc{Formula.mk_forall}{NW3w0iRO-3SlqND-4}(x, \nwlinkedidentc{Formula.subst}{NW3w0iRO-1jtQRV-1} x t A);
    fun jfn' (thm::thms) =
      jfn((\nwlinkedidentc{Thm.forall_elim}{NW4ICsJQ-2zKkPU-9} t thm)::thms);
  in
    \{ goals = (asl, A')::gls
    , justify = jfn' \}
  end;
\nwindexdefn{\nwixident{Tactic.gen}}{Tactic.gen}{NW1F32vJ-3Cb3oS-A}\nwindexdefn{\nwixident{Tactic.spec}}{Tactic.spec}{NW1F32vJ-3Cb3oS-A}\eatline
\nwused{\\{NW1F32vJ-2XH0oT-1}}\nwidentdefs{\\{{\nwixident{Tactic.gen}}{Tactic.gen}}\\{{\nwixident{Tactic.spec}}{Tactic.spec}}}\nwidentuses{\\{{\nwixident{Formula.fv}}{Formula.fv}}\\{{\nwixident{Formula.mk{\_}forall}}{Formula.mk:unforall}}\\{{\nwixident{Formula.subst}}{Formula.subst}}\\{{\nwixident{Term.fresh{\_}var}}{Term.fresh:unvar}}\\{{\nwixident{Thm.forall{\_}elim}}{Thm.forall:unelim}}\\{{\nwixident{Thm.forall{\_}intro}}{Thm.forall:unintro}}}\nwindexuse{\nwixident{Formula.fv}}{Formula.fv}{NW1F32vJ-3Cb3oS-A}\nwindexuse{\nwixident{Formula.mk{\_}forall}}{Formula.mk:unforall}{NW1F32vJ-3Cb3oS-A}\nwindexuse{\nwixident{Formula.subst}}{Formula.subst}{NW1F32vJ-3Cb3oS-A}\nwindexuse{\nwixident{Term.fresh{\_}var}}{Term.fresh:unvar}{NW1F32vJ-3Cb3oS-A}\nwindexuse{\nwixident{Thm.forall{\_}elim}}{Thm.forall:unelim}{NW1F32vJ-3Cb3oS-A}\nwindexuse{\nwixident{Thm.forall{\_}intro}}{Thm.forall:unintro}{NW1F32vJ-3Cb3oS-A}\nwendcode{}\nwbegindocs{56}\nwdocspar
\begin{node}
Existential quantifier tactics. The tactic corresponding to the
\emph{derived} $\exists$-elimination rule:
\begin{equation}
\vcenter{\infer=[\mbox{choose}(\Gamma'\vdash\exists x\ldotp A)]{(A[x'/x]),\Gamma\vdash B}{\Gamma\vdash B}}
\end{equation}
Care must be taken: although this corresponds to
$\exists$-elimination, the form of $\exists$-elimination in the kernel
does not adequately match this tactic.

The tactic corresponding to $\exists$-introduction:
\begin{equation}
\vcenter{\infer=[\mbox{exists}(t)]{\Gamma\vdash A[t/x]}{\Gamma\vdash\exists x\ldotp A}}
\end{equation}
\end{node}
% https://hol-light.github.io/references/HTML/CHOOSE_TAC.html
% https://hol-light.github.io/references/HTML/EXISTS_TAC.html

\nwenddocs{}\nwbegincode{57}\sublabel{NW1F32vJ-2kzjE-B}\nwmargintag{{\nwtagstyle{}\subpageref{NW1F32vJ-2kzjE-B}}}\moddef{Specification for tactics~{\nwtagstyle{}\subpageref{NW1F32vJ-2kzjE-1}}}\plusendmoddef\nwstartdeflinemarkup\nwusesondefline{\\{NW1F32vJ-VKxLl-1}}\nwprevnextdefs{NW1F32vJ-2kzjE-A}{\relax}\nwenddeflinemarkup
val choose : \nwlinkedidentc{Thm.t}{NW4ICsJQ-YIPBT-1} -> t;
val exists : \nwlinkedidentc{Term.t}{NW1rTmOJ-4LsRWC-1} -> t;
\nwused{\\{NW1F32vJ-VKxLl-1}}\nwidentuses{\\{{\nwixident{Term.t}}{Term.t}}\\{{\nwixident{Thm.t}}{Thm.t}}}\nwindexuse{\nwixident{Term.t}}{Term.t}{NW1F32vJ-2kzjE-B}\nwindexuse{\nwixident{Thm.t}}{Thm.t}{NW1F32vJ-2kzjE-B}\nwendcode{}\nwbegindocs{58}\nwdocspar

\nwenddocs{}\nwbegincode{59}\sublabel{NW1F32vJ-3Cb3oS-B}\nwmargintag{{\nwtagstyle{}\subpageref{NW1F32vJ-3Cb3oS-B}}}\moddef{Implementation of tactics~{\nwtagstyle{}\subpageref{NW1F32vJ-3Cb3oS-1}}}\plusendmoddef\nwstartdeflinemarkup\nwusesondefline{\\{NW1F32vJ-2XH0oT-1}}\nwprevnextdefs{NW1F32vJ-3Cb3oS-A}{\relax}\nwenddeflinemarkup
(* choose : \nwlinkedidentc{Thm.t}{NW4ICsJQ-YIPBT-1} -> Tactic.t *)
fun choose th \{goals = (asl, B)::gls, justify = jfn\} =
  if not(\nwlinkedidentc{assumptions_contain_hypotheses}{NW1F32vJ-1bk6hJ-1} th asl)
  then raise Fail "choose: assumptions do not contain all hypotheses of theorem"
  else
    let
      val (x as (Term.Var(x0,s)),A) = Formula.dest_exists(\nwlinkedidentc{Thm.concl}{NW4ICsJQ-YIPBT-1} th);
      val fvs = List.concat (map (fn (_,fm) => \nwlinkedidentc{Formula.fv}{NW3w0iRO-jRgYZ-2} fm) asl);
      val x' = \nwlinkedidentc{Term.fresh_var}{NW1rTmOJ-4Ykq9M-1} x0 s fvs;
      fun jfn' (thm::thms) = jfn((\nwlinkedidentc{Derived.exists_elim}{NW1NjfzY-uZfts-C} x' th thm)::thms);
    in
      \{ goals = (("",\nwlinkedidentc{Formula.subst}{NW3w0iRO-1jtQRV-1} x x' A)::asl, B)::gls
      , justify = jfn' \}
    end;

(* exists : \nwlinkedidentc{Term.t}{NW1rTmOJ-4LsRWC-1} -> Tactic.t *)
fun exists tm \{goals = (asl, ex_A)::gls, justify = jfn\} =
  let
    val (x,A) = Formula.dest_exists(ex_A);
    fun jfn' (thm::thms) = jfn((\nwlinkedidentc{Thm.exists_intro}{NW4ICsJQ-2zKkPU-A} x tm A thm)::thms);
  in
    \{ goals = (asl, \nwlinkedidentc{Formula.subst}{NW3w0iRO-1jtQRV-1} tm x A)::gls
    , justify = jfn' \}
  end;
\nwused{\\{NW1F32vJ-2XH0oT-1}}\nwidentuses{\\{{\nwixident{assumptions{\_}contain{\_}hypotheses}}{assumptions:uncontain:unhypotheses}}\\{{\nwixident{Derived.exists{\_}elim}}{Derived.exists:unelim}}\\{{\nwixident{Formula.fv}}{Formula.fv}}\\{{\nwixident{Formula.subst}}{Formula.subst}}\\{{\nwixident{Term.fresh{\_}var}}{Term.fresh:unvar}}\\{{\nwixident{Term.t}}{Term.t}}\\{{\nwixident{Thm.concl}}{Thm.concl}}\\{{\nwixident{Thm.exists{\_}intro}}{Thm.exists:unintro}}\\{{\nwixident{Thm.t}}{Thm.t}}}\nwindexuse{\nwixident{assumptions{\_}contain{\_}hypotheses}}{assumptions:uncontain:unhypotheses}{NW1F32vJ-3Cb3oS-B}\nwindexuse{\nwixident{Derived.exists{\_}elim}}{Derived.exists:unelim}{NW1F32vJ-3Cb3oS-B}\nwindexuse{\nwixident{Formula.fv}}{Formula.fv}{NW1F32vJ-3Cb3oS-B}\nwindexuse{\nwixident{Formula.subst}}{Formula.subst}{NW1F32vJ-3Cb3oS-B}\nwindexuse{\nwixident{Term.fresh{\_}var}}{Term.fresh:unvar}{NW1F32vJ-3Cb3oS-B}\nwindexuse{\nwixident{Term.t}}{Term.t}{NW1F32vJ-3Cb3oS-B}\nwindexuse{\nwixident{Thm.concl}}{Thm.concl}{NW1F32vJ-3Cb3oS-B}\nwindexuse{\nwixident{Thm.exists{\_}intro}}{Thm.exists:unintro}{NW1F32vJ-3Cb3oS-B}\nwindexuse{\nwixident{Thm.t}}{Thm.t}{NW1F32vJ-3Cb3oS-B}\nwendcode{}\nwbegindocs{60}\nwdocspar

\begin{xca}
Prove these tactics are ``complete'', in the sense that any proof we
can do with the kernel's inference rules can also be done using these
tactics. [Hint: it's probably easiest to first prove we can derive the kernel's $\exists$-elimination rule~(\S\ref{node:kernel:existential-elimination}) from the derived $\exists$-elimination rule (\S\ref{node:derived:existential-elimination}).]
\end{xca}

\begin{problem}
How exactly do we ``use'' assumptions in a tactics proof?
\end{problem}

\begin{node}
After all, we have the assumptions list for a subgoal accompanying us
at every turn, how do we use them and assert something like, ``This subgoal is precisely assumption 3''?
One way to do it is use \textit{modus ponens} tactic (\S\ref{node:tactic:MP})
with the theorem $A_{3}\vdash A_{3}$ by ``{\Tt{}\nwlinkedidentq{Thm.assume}{NW4ICsJQ-2zKkPU-2}\nwendquote} $A_{3}$ \texttt{[]}''.
\end{node}

\begin{node}
More broadly, there is a related problem: suppose the current subgoal
is already proven as a theorem elsewhere? How do we ``solve'' the
subgoal \emph{using} that theorem? HOL family provers use the
{\Tt{}ACCEPT{\_}TAC\nwendquote} to accomplish this.

However, it is not hard to imagine situations where the variables
appearing in the theorem differ from what the goal expects. If we are
trying to prove $\vdash\forall x\ldotp P[x]$ and we have a theorem
$\vdash\forall y\ldotp P[y]$ where $x\neq y$ are different variables,
then it would be nice to convert the variables to match. Borrowing the
language of $\lambda$-calculus, it would be nice for the tactic to
perform $\alpha$-conversion. (Note: this would also have to transform
free variables affected by the $\alpha$-conversion which appear in the
assumptions for the subgoal.)
\end{node}

\begin{node}
Suppose our goal is of the form $A_{1},\dots,A_{n}\vdash^{\mkern-8.75mu\scriptstyle\rule[-.9ex]{0pt}{0pt}?} A$ and this
is really the same as a theorem $A'_{1},\dots,A'_{m}\vdash A'$ where
$m\leq n$ and there is some injective map $\pi\colon\{1,\dots,m\}\to\{1,\dots,n\}$
such that $A'_{i}$ looks like $A_{\pi(i)}$. We want to essentially
construct a derived inference rule taking
\begin{enumerate}
\item The ordered pair $([A_{1},\dots,A_{n}], A)$ corresponding to the goal;
\item The theorem $A'_{1},\dots,A'_{m}\vdash A'$
\end{enumerate}
and produce a theorem of the form
\begin{equation*}
\vdash(A'_{1}\land\cdots\land A'_{m}\implies A')\implies(A_{1}\land\cdots\land A_{n}\implies A).
\end{equation*}
This is where all the hard work lies, making this derived inference
rule. Then we can perform \textit{modus ponens}, repeatedly uncurry
the result, and undischarge until we obtain the theorem $A_{1},\dots,A_{n}\vdash A$.
This would give us some {\Tt{}match\nwendquote} derived rule of inference, which we
can then turn into a tactic.
\end{node}



\begin{node}[References]
For more about tactics, it is useful to see how other LCF-style
provers employ them. See, e.g., Chapter 4 of HOL's description~\cite{HOLdescription}.
\end{node}
\nwenddocs{}\nwfilename{nw/unification.nw}\nwbegindocs{0}% -*- mode: poly-noweb; noweb-code-mode: sml-mode; -*-
%%
%% unification.nw
%% 
%% Made by Alex Nelson <pqnelson@gmail.com>
%% Login   <alex@lisp>
%% 
%% Started on  2026-04-25T15:53:23-0700
%% Last update 2026-04-25T15:53:23-0700
%% 

\section{Unification}


\begin{node}
Unification, for first-order logic, given two ``expressions'' $e_{1}$
and $e_{2}$ (a pair of terms or a pair of formulas) tries to find a
``substitution'' $\sigma$ which can be applied to obtain
\emph{syntactically} identical expressions
$\sigma(e_{1})=\sigma(e_{2})$ (where $\sigma(-)$ essentially
corresponds to a ``parallel'' version of {\Tt{}\nwlinkedidentq{Formula.subst}{NW3w0iRO-1jtQRV-1}\nwendquote} or {\Tt{}\nwlinkedidentq{Term.subst}{NW1rTmOJ-3OIyM5-1}\nwendquote}). The {\Tt{}\nwlinkedidentq{Formula.eq}{NW3w0iRO-5gQzE-1}\nwendquote} and {\Tt{}\nwlinkedidentq{Term.eq}{NW1rTmOJ-3810sY-1}\nwendquote}
predicates test if two formulas (resp., terms) are syntactically
identical. If two expressions are unifiable, then it returns such a
substitution. 

We will need only ``first-order unification'', i.e., we work with only
``first-order variables'' and not ``higher-order
variables''. Higher-order logic allows us to quantify over predicates
and functions using higher-order variables, but unification for
higher-order variables is rather complicated. First-order unification
merely allows replacing variables with terms, which is far more tame
by comparison.

We will only need to unify terms and literals (predicates or their
negation). 
\end{node}

\nwenddocs{}\nwbegincode{1}\sublabel{NW1BOcHB-3z172B-1}\nwmargintag{{\nwtagstyle{}\subpageref{NW1BOcHB-3z172B-1}}}\moddef{Unif.sig~{\nwtagstyle{}\subpageref{NW1BOcHB-3z172B-1}}}\endmoddef\nwstartdeflinemarkup\nwenddeflinemarkup
signature Unif = sig
  structure Env :
            sig
              type t;
              val undefined : unit -> t;
              val apply : t -> \nwlinkedidentc{Term.t}{NW1rTmOJ-4LsRWC-1} -> \nwlinkedidentc{Term.t}{NW1rTmOJ-4LsRWC-1};
              val defined : t -> \nwlinkedidentc{Term.t}{NW1rTmOJ-4LsRWC-1} -> bool;
              val insert : t -> \nwlinkedidentc{Term.t}{NW1rTmOJ-4LsRWC-1} -> \nwlinkedidentc{Term.t}{NW1rTmOJ-4LsRWC-1} -> t;
            end;
  val y : \nwlinkedidentc{Env.t}{NW1BOcHB-25AKca-1} -> (\nwlinkedidentc{Term.t}{NW1rTmOJ-4LsRWC-1} * \nwlinkedidentc{Term.t}{NW1rTmOJ-4LsRWC-1}) list -> \nwlinkedidentc{Env.t}{NW1BOcHB-25AKca-1};
  val y_literals : \nwlinkedidentc{Env.t}{NW1BOcHB-25AKca-1} -> \nwlinkedidentc{Formula.t}{NW3w0iRO-2jnj9W-1} * \nwlinkedidentc{Formula.t}{NW3w0iRO-2jnj9W-1} -> \nwlinkedidentc{Env.t}{NW1BOcHB-25AKca-1};
  val y_complements : \nwlinkedidentc{Env.t}{NW1BOcHB-25AKca-1} -> \nwlinkedidentc{Formula.t}{NW3w0iRO-2jnj9W-1} * \nwlinkedidentc{Formula.t}{NW3w0iRO-2jnj9W-1} -> \nwlinkedidentc{Env.t}{NW1BOcHB-25AKca-1};
end;
\nwnotused{Unif.sig}\nwidentuses{\\{{\nwixident{Env.t}}{Env.t}}\\{{\nwixident{Formula.t}}{Formula.t}}\\{{\nwixident{Term.t}}{Term.t}}}\nwindexuse{\nwixident{Env.t}}{Env.t}{NW1BOcHB-3z172B-1}\nwindexuse{\nwixident{Formula.t}}{Formula.t}{NW1BOcHB-3z172B-1}\nwindexuse{\nwixident{Term.t}}{Term.t}{NW1BOcHB-3z172B-1}\nwendcode{}\nwbegindocs{2}\nwdocspar

\nwenddocs{}\nwbegincode{3}\sublabel{NW1BOcHB-1KP4Zf-1}\nwmargintag{{\nwtagstyle{}\subpageref{NW1BOcHB-1KP4Zf-1}}}\moddef{Unif.sml~{\nwtagstyle{}\subpageref{NW1BOcHB-1KP4Zf-1}}}\endmoddef\nwstartdeflinemarkup\nwenddeflinemarkup
structure Unif :> Unif = struct
\LA{}Finite map for unification~{\nwtagstyle{}\subpageref{NW1BOcHB-25AKca-1}}\RA{}

\LA{}Poor man's hashmap~{\nwtagstyle{}\subpageref{nw@notdef}}\RA{}

\LA{}Unify two terms~{\nwtagstyle{}\subpageref{NW1BOcHB-1kcz6k-1}}\RA{}

\LA{}Unify two literals~{\nwtagstyle{}\subpageref{NW1BOcHB-3w0JS3-1}}\RA{}
end;
\nwnotused{Unif.sml}\nwendcode{}\nwbegindocs{4}\nwdocspar

\begin{node}
Unification tracks the substitution constructed using a finite partial
map. We recursively build the map from equations given for unification.
We will write $t\sim t'$ for the ``equation'' submitted for
unification involving the two terms $t$ and $t'$. This will produce a
finite partial map, and we will denote a specific entry in it as
$x\mapsto t$ to indicate $x$ would be substituted for the term $t$.
\end{node}

\begin{definition}
Suppose $M$ is a finite partial map from terms to terms.
\begin{enumerate}
\item Let $x\mapsto t$ be an entry in the map $M$. We write $x\to y$ if there exists a $y\in FV(t)$.
\item A \define{Cycle} in $M$ is a nonempty finite sequence leading
  back to the starting point $x_{0}\to x_{1}\to\dots\to x_{n}\to x_{0}$.
\item If there are no cycles in $M$, we say it is \define{Cycle-free}. Otherwise,
  we call $M$ \define{Cyclic}.
\end{enumerate}
\end{definition}

\begin{example}[{Duffy~\cite[p.143]{duffy1991principles}}]
Consider unifying $P[x,f(x)]$ with $P[g(y),y]$ where $P[-,-]$ is a
predicate of two terms, $f(-)$ and $g(-)$ are functions.

The first step will be to break this up into two equations
\begin{equation}
x\sim g(y)
\end{equation}
and
\begin{equation}
f(x)\sim y.
\end{equation}
We can solve the first equation with the replacement $x\mapsto g(y)$.

The second equation can be ``rewritten'' as
\begin{equation}
y\sim f(x),
\end{equation}
by symmetry. However, if we try adjoining the replacement $y\mapsto f(x)$
to our current partial map, then we end up with a cycle $x\to y\to x$.

We want to avoid this. Why? Because if we consider the formula,
\begin{equation}\label{eq:unification:example-of-why-we-want-to-avoid-cycles}
(\forall x\exists y\ldotp P[x,y])\implies(\exists y\forall x\ldotp P[x,y]),
\end{equation}
the negation of this formula has conjunctive normal form
\begin{equation}
\forall x,y\ldotp P[x,f(x)]\land\neg P[g(y),y].
\end{equation}
A resolution proof can be found for this, if we allowed cycles during
unification. But Equation~\eqref{eq:unification:example-of-why-we-want-to-avoid-cycles}
is not a theorem in classical logic. Had we allowed cycles, soundness
would be compromised. 
\end{example}

\begin{node}
The finite map used for tracking variable assignment could be
anything. For simplicity, at the moment, we use association lists.

A ``Poor Man's Table'' is just an association list.
If we implement a more optimized hash table, like a Red-Black tree or
a 2-3 tree, then we can jettison this implementation and use something
superior. The contract:
\begin{enumerate}
\item {\Tt{}undefined\nwendquote} is an empty table
\item {\Tt{}apply\ table\ key\nwendquote} gets the value associated with the {\Tt{}key\nwendquote}
  in the {\Tt{}table\nwendquote}
\item {\Tt{}insert\ table\ key\ value\nwendquote} inserts the association of {\Tt{}value\nwendquote}
  with {\Tt{}key\nwendquote} into the {\Tt{}table\nwendquote}, raising a {\Tt{}Fail\nwendquote} exception if
  this would produce a cycle in the substitution;
\item {\Tt{}defined\ table\ key\nwendquote} tests if the {\Tt{}table\nwendquote} contains the
  {\Tt{}key\nwendquote} at all
\end{enumerate}
These are the table operations.

Right now, this is a huge bottleneck for performance. The low-hanging
fruit would be to make this replacement.
\end{node}

\begin{invariant}[Partial maps remain cycle free]\label{invariant:unification:cycle-free}
Every [public] function in the {\Tt{}Unif\nwendquote} structure requires any
{\Tt{}\nwlinkedidentq{Env.t}{NW1BOcHB-25AKca-1}\nwendquote} partial map given as an argument to be cycle free, and
ensures the {\Tt{}\nwlinkedidentq{Env.t}{NW1BOcHB-25AKca-1}\nwendquote} partial map produced (if any) to be cycle-free.
\end{invariant}

\nwenddocs{}\nwbegincode{5}\sublabel{NW1BOcHB-25AKca-1}\nwmargintag{{\nwtagstyle{}\subpageref{NW1BOcHB-25AKca-1}}}\moddef{Finite map for unification~{\nwtagstyle{}\subpageref{NW1BOcHB-25AKca-1}}}\endmoddef\nwstartdeflinemarkup\nwusesondefline{\\{NW1BOcHB-1KP4Zf-1}}\nwenddeflinemarkup
structure Env :>
            sig
              type t;
              val undefined : unit -> t;
              val apply : t -> \nwlinkedidentc{Term.t}{NW1rTmOJ-4LsRWC-1} -> \nwlinkedidentc{Term.t}{NW1rTmOJ-4LsRWC-1};
              val defined : t -> \nwlinkedidentc{Term.t}{NW1rTmOJ-4LsRWC-1} -> bool;
              val insert : t -> \nwlinkedidentc{Term.t}{NW1rTmOJ-4LsRWC-1} -> \nwlinkedidentc{Term.t}{NW1rTmOJ-4LsRWC-1} -> t;
            end =
struct
type t = (\nwlinkedidentc{Term.t}{NW1rTmOJ-4LsRWC-1} * \nwlinkedidentc{Term.t}{NW1rTmOJ-4LsRWC-1}) list;

fun undefined () = [];

fun apply [] _ = raise Fail "apply"
  | apply ((y,tm)::ys) x = if \nwlinkedidentc{Term.eq}{NW1rTmOJ-3810sY-1} x y
                           then tm
                           else apply ys x;

fun defined [] _ = false
  | defined ((y,_)::env) x = (\nwlinkedidentc{Term.eq}{NW1rTmOJ-3810sY-1} x y) orelse
                             defined env x;
local
\LA{}Test if a new equations is trivial for unification~{\nwtagstyle{}\subpageref{NW1BOcHB-4edp1Q-1}}\RA{}
in
fun insert env (x as (Term.Var _)) y = 
      if is_cyclic env x y then env 
      else (x,y)::env
  | insert _ _ _ = raise Fail "\nwlinkedidentc{Env.insert}{NW1BOcHB-25AKca-1}: trying to insert function as key"
end;
end;
\nwindexdefn{\nwixident{Env.t}}{Env.t}{NW1BOcHB-25AKca-1}\nwindexdefn{\nwixident{Env.undefined}}{Env.undefined}{NW1BOcHB-25AKca-1}\nwindexdefn{\nwixident{Env.apply}}{Env.apply}{NW1BOcHB-25AKca-1}\nwindexdefn{\nwixident{Env.insert}}{Env.insert}{NW1BOcHB-25AKca-1}\nwindexdefn{\nwixident{Env.defined}}{Env.defined}{NW1BOcHB-25AKca-1}\eatline
\nwused{\\{NW1BOcHB-1KP4Zf-1}}\nwidentdefs{\\{{\nwixident{Env.apply}}{Env.apply}}\\{{\nwixident{Env.defined}}{Env.defined}}\\{{\nwixident{Env.insert}}{Env.insert}}\\{{\nwixident{Env.t}}{Env.t}}\\{{\nwixident{Env.undefined}}{Env.undefined}}}\nwidentuses{\\{{\nwixident{Term.eq}}{Term.eq}}\\{{\nwixident{Term.t}}{Term.t}}}\nwindexuse{\nwixident{Term.eq}}{Term.eq}{NW1BOcHB-25AKca-1}\nwindexuse{\nwixident{Term.t}}{Term.t}{NW1BOcHB-25AKca-1}\nwendcode{}\nwbegindocs{6}\nwdocspar
\begin{node}
We also have the really interesting part of solving the unification
problem: testing if $x=t$ is a trivial equation or not. This would
happen for literally trivial equations $x=x$ which shed no light on anything,
but also for $x=y$ when we already have $y=t'$.

However, there is the possibility that we might produce a cycle. This
is handled by raising an error.

Specifically: checking there are no cycles when we try to associate
variable $x$ with term $t$ depends on what $t$ is. When
\begin{enumerate}
\item $t$ is a variable $t=y$: This boils down to two subcases.
  \begin{enumerate}
  \item $x=y$ 
  \item $x\neq y$, but $y$ is bound to something in the
    environment $y=t_{y}$. Then we just check recursively that associating $x$ to
    $t_{y}$ will be cycle-free.
  \end{enumerate}
\item $t$ is a function $f(\vec{t})$: There is no variable $y\in
  FV(t)$ such that $y\to^{*}x$ (i.e., there is no sequence of zero or more
  steps leading from $y$ to $x$)
\end{enumerate}
We can check the first condition elsewhere (namely, in the {\Tt{}unify\nwendquote}
function). The second condition will be tested in {\Tt{}is{\_}cyclic\nwendquote},
which will return \textbf{false} if the condition is satisfied.
\end{node}

\begin{node}
If we extend {\Tt{}env\nwendquote} with a new entry $x\sim t$, then we should
consider the following situations:
\begin{enumerate}
\item $t$ is a variable $t=y$.
  \begin{enumerate}
  \item If $y=x$ are the same variable, then this is a trivial
    extension, and we should return \textbf{true} to indicate \emph{do
    no extend the environment}.
  \item If $y\neq x$ is a different variable, then we should check if
    $y\sim t'$ is in the environment and if so whether we can add
    $x\sim t'$ to the environment.

    If there is no entry $y\sim t'$ in the environment, then return
    \textbf{false} to indicate \emph{we can extend the environment
    while maintaining the cycle-free invariant}.

    If there is an entry $y\sim t'$ in the environment, then we just
    recursively test if $x\sim t'$ violates the cycle-free invariant
    of the environment.
  \end{enumerate}
\item If $t$ is a function $f(t_{1},\dots,t_{n})$, then we want to test 
  that none of the possible extensions $x\sim t_{1}$, \dots, $x\sim t_{n}$
  would lead to a cyclic environment. If there is a possible extension
  $x\sim t_{k}$, then return \textbf{true}. Otherwise,
  return \textbf{false} to indicate \emph{we can extend the
  environment while maintaining the cycle-free invariant}.
\end{enumerate}
\end{node}

\nwenddocs{}\nwbegincode{7}\sublabel{NW1BOcHB-4edp1Q-1}\nwmargintag{{\nwtagstyle{}\subpageref{NW1BOcHB-4edp1Q-1}}}\moddef{Test if a new equations is trivial for unification~{\nwtagstyle{}\subpageref{NW1BOcHB-4edp1Q-1}}}\endmoddef\nwstartdeflinemarkup\nwusesondefline{\\{NW1BOcHB-25AKca-1}}\nwenddeflinemarkup
(* REQUIRES: `x` is not defined in the `env` *)
fun is_cyclic env x tm =
 (case tm of
   (y as (Term.Var _)) => \nwlinkedidentc{Term.eq}{NW1rTmOJ-3810sY-1} x y orelse (* trivial case *)
                          \nwlinkedidentc{Term.have_same_sort}{NW1rTmOJ-1WQNps-1} x y andalso
                          defined env y andalso
                          is_cyclic env x (apply env y)
 | (Term.Fun(f,s,args)) => List.exists (is_cyclic env x) args andalso
                           raise Fail "cyclic");
\nwused{\\{NW1BOcHB-25AKca-1}}\nwidentuses{\\{{\nwixident{Term.eq}}{Term.eq}}\\{{\nwixident{Term.have{\_}same{\_}sort}}{Term.have:unsame:unsort}}}\nwindexuse{\nwixident{Term.eq}}{Term.eq}{NW1BOcHB-4edp1Q-1}\nwindexuse{\nwixident{Term.have{\_}same{\_}sort}}{Term.have:unsame:unsort}{NW1BOcHB-4edp1Q-1}\nwendcode{}\nwbegindocs{8}\nwdocspar

\begin{node}
Unification tries to ``solve'' a bunch of equations. We have the
following situations:
\begin{enumerate}
\item $f(t_{1},\dots,t_{m})=g(t'_{1},\dots,t'_{n})$ requires $f$ and
  $g$ to be the same function symbol with the same sort, and $m=n$,
  then we just need to solve the equations $t_{1}=t'_{1}$, \dots, $t_{n}=t'_{n}$.
\item $x=t$ first requires checking if we have already ``solved'' for
  $x$, and if so try to match it against $t$.

  Otherwise, if $x=t$ is a trivial equation, then we just move along
  to try solving the remaining equations.

  If $x=t$ is not a trivial equation, then we extend the current
  environment (``solutions'') with this entry, and then try solving
  the remaining equations.
\item $t=x$ is just the same situation.
\end{enumerate}
There cannot be any other possible situation.
\end{node}

\nwenddocs{}\nwbegincode{9}\sublabel{NW1BOcHB-1kcz6k-1}\nwmargintag{{\nwtagstyle{}\subpageref{NW1BOcHB-1kcz6k-1}}}\moddef{Unify two terms~{\nwtagstyle{}\subpageref{NW1BOcHB-1kcz6k-1}}}\endmoddef\nwstartdeflinemarkup\nwusesondefline{\\{NW1BOcHB-1KP4Zf-1}}\nwenddeflinemarkup
(* \nwlinkedidentc{Unif.y}{NW1BOcHB-1kcz6k-1} : \nwlinkedidentc{Env.t}{NW1BOcHB-25AKca-1} -> (\nwlinkedidentc{Term.t}{NW1rTmOJ-4LsRWC-1} * \nwlinkedidentc{Term.t}{NW1rTmOJ-4LsRWC-1}) list -> \nwlinkedidentc{Env.t}{NW1BOcHB-25AKca-1}

REQUIRES: `env` is cycle-free
REQUIRES: the set `env @ eqns` has exactly the same set of
          unifiers as the original problem. *)
fun y env ([] : (\nwlinkedidentc{Term.t}{NW1rTmOJ-4LsRWC-1} * \nwlinkedidentc{Term.t}{NW1rTmOJ-4LsRWC-1}) list) = env
  | y env ((Term.Fun(f,s1,args1),Term.Fun(g,s2,args2))::eqns) = 
      if f = g andalso s1 = s2 andalso length args1 = length args2
      then y env ((ListPair.zip(args1, args2)) @ eqns)
      else raise Fail "impossible unification"
  | y env (((x as (Term.Var _)), t)::eqns) =
      if not(\nwlinkedidentc{Term.have_same_sort}{NW1rTmOJ-1WQNps-1} x t)
      then raise Fail "impossible unification: variable has different sort than term"
      else if \nwlinkedidentc{Env.defined}{NW1BOcHB-25AKca-1} env x (* x already in env *)
      then y env ((\nwlinkedidentc{Env.apply}{NW1BOcHB-25AKca-1} env x, t)::eqns)
      else y (\nwlinkedidentc{Env.insert}{NW1BOcHB-25AKca-1} env x t) eqns
  | y env ((t, x as (Term.Var _))::eqns) = y env ((x,t)::eqns);
\nwindexdefn{\nwixident{Unif.y}}{Unif.y}{NW1BOcHB-1kcz6k-1}\eatline
\nwused{\\{NW1BOcHB-1KP4Zf-1}}\nwidentdefs{\\{{\nwixident{Unif.y}}{Unif.y}}}\nwidentuses{\\{{\nwixident{Env.apply}}{Env.apply}}\\{{\nwixident{Env.defined}}{Env.defined}}\\{{\nwixident{Env.insert}}{Env.insert}}\\{{\nwixident{Env.t}}{Env.t}}\\{{\nwixident{Term.have{\_}same{\_}sort}}{Term.have:unsame:unsort}}\\{{\nwixident{Term.t}}{Term.t}}}\nwindexuse{\nwixident{Env.apply}}{Env.apply}{NW1BOcHB-1kcz6k-1}\nwindexuse{\nwixident{Env.defined}}{Env.defined}{NW1BOcHB-1kcz6k-1}\nwindexuse{\nwixident{Env.insert}}{Env.insert}{NW1BOcHB-1kcz6k-1}\nwindexuse{\nwixident{Env.t}}{Env.t}{NW1BOcHB-1kcz6k-1}\nwindexuse{\nwixident{Term.have{\_}same{\_}sort}}{Term.have:unsame:unsort}{NW1BOcHB-1kcz6k-1}\nwindexuse{\nwixident{Term.t}}{Term.t}{NW1BOcHB-1kcz6k-1}\nwendcode{}\nwbegindocs{10}\nwdocspar
\begin{node}
Unifying two literals $P[t_{1},\dots,t_{m}]$ with $Q[u_{1},\dots,u_{n}]$
requires checking $P=Q$ (they have the same identifier/name) and $m=n$ (they have the same number of arguments), then invoking
{\Tt{}unify\nwendquote} to solve the system $t_{1}=u_{1}$, \dots, $t_{n}=u_{n}$.
If $P\neq Q$ are distinct identifiers, or $m\neq n$, then a failure is raised.
If $Q$ is the negation of $P$, then a failure is raised (there is no
way to unify a predicate with its negation).
\end{node}

\nwenddocs{}\nwbegincode{11}\sublabel{NW1BOcHB-3w0JS3-1}\nwmargintag{{\nwtagstyle{}\subpageref{NW1BOcHB-3w0JS3-1}}}\moddef{Unify two literals~{\nwtagstyle{}\subpageref{NW1BOcHB-3w0JS3-1}}}\endmoddef\nwstartdeflinemarkup\nwusesondefline{\\{NW1BOcHB-1KP4Zf-1}}\nwenddeflinemarkup
(* \nwlinkedidentc{Unif.y_literals}{NW1BOcHB-3w0JS3-1} : \nwlinkedidentc{Env.t}{NW1BOcHB-25AKca-1} -> \nwlinkedidentc{Term.t}{NW1rTmOJ-4LsRWC-1} * \nwlinkedidentc{Term.t}{NW1rTmOJ-4LsRWC-1} -> \nwlinkedidentc{Env.t}{NW1BOcHB-25AKca-1} *)
fun y_literals env (lhs,rhs) =
  if \nwlinkedidentc{Formula.is_not}{NW3w0iRO-3NyJLv-2} lhs andalso \nwlinkedidentc{Formula.is_not}{NW3w0iRO-3NyJLv-2} rhs
  then y_literals env (\nwlinkedidentc{Formula.dest_not}{NW3w0iRO-1Yt1Jl-2} lhs, \nwlinkedidentc{Formula.dest_not}{NW3w0iRO-1Yt1Jl-2} rhs)
  else if \nwlinkedidentc{Formula.is_falsum}{NW3w0iRO-3NyJLv-1} lhs andalso \nwlinkedidentc{Formula.is_falsum}{NW3w0iRO-3NyJLv-1} rhs
  then env
  else if \nwlinkedidentc{Formula.is_pred}{NW3w0iRO-3NyJLv-1} lhs andalso \nwlinkedidentc{Formula.is_pred}{NW3w0iRO-3NyJLv-1} rhs
  then let val (p1,args1) = \nwlinkedidentc{Formula.dest_pred}{NW3w0iRO-1Yt1Jl-1} lhs;
           val (p2,args2) = \nwlinkedidentc{Formula.dest_pred}{NW3w0iRO-1Yt1Jl-1} rhs;
       in if p1 = p2 andalso length args1 = length args2
          then y env (ListPair.zip(args1, args2))
          else raise Fail "\nwlinkedidentc{Unif.y_literals}{NW1BOcHB-3w0JS3-1}: impossible unification"
       end
  else raise Fail "\nwlinkedidentc{Unif.y_literals}{NW1BOcHB-3w0JS3-1}: Cannot unify literals";

fun negate A = if \nwlinkedidentc{Formula.is_not}{NW3w0iRO-3NyJLv-2} A
               then \nwlinkedidentc{Formula.dest_not}{NW3w0iRO-1Yt1Jl-2} A
               else \nwlinkedidentc{Formula.mk_not}{NW3w0iRO-3SlqND-2} A;

fun y_complements env (A,B) =
  y_literals env (A, negate B);
\nwindexdefn{\nwixident{Unif.y{\_}complements}}{Unif.y:uncomplements}{NW1BOcHB-3w0JS3-1}\nwindexdefn{\nwixident{Unif.y{\_}literals}}{Unif.y:unliterals}{NW1BOcHB-3w0JS3-1}\eatline
\nwused{\\{NW1BOcHB-1KP4Zf-1}}\nwidentdefs{\\{{\nwixident{Unif.y{\_}complements}}{Unif.y:uncomplements}}\\{{\nwixident{Unif.y{\_}literals}}{Unif.y:unliterals}}}\nwidentuses{\\{{\nwixident{Env.t}}{Env.t}}\\{{\nwixident{Formula.dest{\_}not}}{Formula.dest:unnot}}\\{{\nwixident{Formula.dest{\_}pred}}{Formula.dest:unpred}}\\{{\nwixident{Formula.is{\_}falsum}}{Formula.is:unfalsum}}\\{{\nwixident{Formula.is{\_}not}}{Formula.is:unnot}}\\{{\nwixident{Formula.is{\_}pred}}{Formula.is:unpred}}\\{{\nwixident{Formula.mk{\_}not}}{Formula.mk:unnot}}\\{{\nwixident{Term.t}}{Term.t}}}\nwindexuse{\nwixident{Env.t}}{Env.t}{NW1BOcHB-3w0JS3-1}\nwindexuse{\nwixident{Formula.dest{\_}not}}{Formula.dest:unnot}{NW1BOcHB-3w0JS3-1}\nwindexuse{\nwixident{Formula.dest{\_}pred}}{Formula.dest:unpred}{NW1BOcHB-3w0JS3-1}\nwindexuse{\nwixident{Formula.is{\_}falsum}}{Formula.is:unfalsum}{NW1BOcHB-3w0JS3-1}\nwindexuse{\nwixident{Formula.is{\_}not}}{Formula.is:unnot}{NW1BOcHB-3w0JS3-1}\nwindexuse{\nwixident{Formula.is{\_}pred}}{Formula.is:unpred}{NW1BOcHB-3w0JS3-1}\nwindexuse{\nwixident{Formula.mk{\_}not}}{Formula.mk:unnot}{NW1BOcHB-3w0JS3-1}\nwindexuse{\nwixident{Term.t}}{Term.t}{NW1BOcHB-3w0JS3-1}\nwendcode{}\nwfilename{nw/classical.nw}\nwbegindocs{0}% -*- mode: poly-noweb; noweb-code-mode: sml-mode; -*-
%%
%% classical.nw
%% 
%% Made by Alex Nelson <pqnelson@gmail.com>
%% Login   <alex@lisp>
%% 
%% Started on  2026-04-21T12:45:57-0700
%% Last update 2026-04-21T12:45:57-0700
%% 

\section{Classical tableaux}

\begin{node}
We want to investigate if certain $FS_{0}$ formulas are theorems, had
we used Classical logic instead of Intuitionistic logic. After all,
every Intuitionistically valid formula is also valid Classically (but
not all Classically valid formulas are Intuitionistically acceptable).
Tableaux is one way to investigate this question.
\end{node}

\nwenddocs{}\nwbegincode{1}\sublabel{NW1mR7zJ-3NMRe1-1}\nwmargintag{{\nwtagstyle{}\subpageref{NW1mR7zJ-3NMRe1-1}}}\moddef{Classical.sig~{\nwtagstyle{}\subpageref{NW1mR7zJ-3NMRe1-1}}}\endmoddef\nwstartdeflinemarkup\nwenddeflinemarkup
signature Classical = sig
  val nnf : \nwlinkedidentc{Formula.t}{NW3w0iRO-2jnj9W-1} -> \nwlinkedidentc{Formula.t}{NW3w0iRO-2jnj9W-1};
  val askolemize : \nwlinkedidentc{Formula.t}{NW3w0iRO-2jnj9W-1} -> \nwlinkedidentc{Formula.t}{NW3w0iRO-2jnj9W-1};
  val tab : \nwlinkedidentc{Formula.t}{NW3w0iRO-2jnj9W-1} -> int -> int;
end;
\nwnotused{Classical.sig}\nwidentuses{\\{{\nwixident{Formula.t}}{Formula.t}}}\nwindexuse{\nwixident{Formula.t}}{Formula.t}{NW1mR7zJ-3NMRe1-1}\nwendcode{}\nwbegindocs{2}\nwdocspar

\nwenddocs{}\nwbegincode{3}\sublabel{NW1mR7zJ-l5jX3-1}\nwmargintag{{\nwtagstyle{}\subpageref{NW1mR7zJ-l5jX3-1}}}\moddef{Classical.sml~{\nwtagstyle{}\subpageref{NW1mR7zJ-l5jX3-1}}}\endmoddef\nwstartdeflinemarkup\nwenddeflinemarkup
structure Classical :> Classical = struct
\LA{}Transform formula to negation normal form~{\nwtagstyle{}\subpageref{NW1mR7zJ-4TkWUc-1}}\RA{}

\LA{}Simplify formulas using classical tautologies~{\nwtagstyle{}\subpageref{NW1mR7zJ-2RorhN-1}}\RA{}

\LA{}Skolemization~{\nwtagstyle{}\subpageref{NW1mR7zJ-4QrsTb-1}}\RA{}

\LA{}Tableaux refuter~{\nwtagstyle{}\subpageref{NW1mR7zJ-yCIr5-1}}\RA{}
end;
\nwnotused{Classical.sml}\nwendcode{}\nwbegindocs{4}\nwdocspar

\nwenddocs{}\nwbegincode{5}\sublabel{NW1mR7zJ-10sN88-1}\nwmargintag{{\nwtagstyle{}\subpageref{NW1mR7zJ-10sN88-1}}}\moddef{tests/ClassicalSuite.sml~{\nwtagstyle{}\subpageref{NW1mR7zJ-10sN88-1}}}\endmoddef\nwstartdeflinemarkup\nwenddeflinemarkup
structure ClassicalSuite : TestSuite = struct
(* helper functions to make unit tests prettier *)
val A = \nwlinkedidentc{Formula.mk_pred}{NW3w0iRO-3SlqND-1}("A", []);
val B = \nwlinkedidentc{Formula.mk_pred}{NW3w0iRO-3SlqND-1}("B", []);
infixr ==> <=>;
fun (p ==> q) = \nwlinkedidentc{Formula.mk_imp}{NW3w0iRO-3SlqND-3}(p, q);
fun (p <=> q) = \nwlinkedidentc{Formula.mk_iff}{NW3w0iRO-3SlqND-3}(p, q);
infixr &;
fun (p & q) = \nwlinkedidentc{Formula.mk_and}{NW3w0iRO-3SlqND-3}(p, q);

fun ~ p = \nwlinkedidentc{Formula.mk_not}{NW3w0iRO-3SlqND-2}(p);

fun exists(x, p) = \nwlinkedidentc{Formula.mk_exists}{NW3w0iRO-3SlqND-4}(x, p);
fun all(x, p) = \nwlinkedidentc{Formula.mk_forall}{NW3w0iRO-3SlqND-4}(x, p);

(* the tests themselves *)
val suite = \nwlinkedidentc{Test.suite}{NWsR2q7-3fh2dK-1} "ClassicalSuite" [
  \LA{}Unit tests for \code{}Classical\edoc{}~{\nwtagstyle{}\subpageref{NW1mR7zJ-4QAFgj-1}}\RA{}
];
end;
\nwnotused{tests/ClassicalSuite.sml}\nwidentuses{\\{{\nwixident{Formula.mk{\_}and}}{Formula.mk:unand}}\\{{\nwixident{Formula.mk{\_}exists}}{Formula.mk:unexists}}\\{{\nwixident{Formula.mk{\_}forall}}{Formula.mk:unforall}}\\{{\nwixident{Formula.mk{\_}iff}}{Formula.mk:uniff}}\\{{\nwixident{Formula.mk{\_}imp}}{Formula.mk:unimp}}\\{{\nwixident{Formula.mk{\_}not}}{Formula.mk:unnot}}\\{{\nwixident{Formula.mk{\_}pred}}{Formula.mk:unpred}}\\{{\nwixident{Test.suite}}{Test.suite}}}\nwindexuse{\nwixident{Formula.mk{\_}and}}{Formula.mk:unand}{NW1mR7zJ-10sN88-1}\nwindexuse{\nwixident{Formula.mk{\_}exists}}{Formula.mk:unexists}{NW1mR7zJ-10sN88-1}\nwindexuse{\nwixident{Formula.mk{\_}forall}}{Formula.mk:unforall}{NW1mR7zJ-10sN88-1}\nwindexuse{\nwixident{Formula.mk{\_}iff}}{Formula.mk:uniff}{NW1mR7zJ-10sN88-1}\nwindexuse{\nwixident{Formula.mk{\_}imp}}{Formula.mk:unimp}{NW1mR7zJ-10sN88-1}\nwindexuse{\nwixident{Formula.mk{\_}not}}{Formula.mk:unnot}{NW1mR7zJ-10sN88-1}\nwindexuse{\nwixident{Formula.mk{\_}pred}}{Formula.mk:unpred}{NW1mR7zJ-10sN88-1}\nwindexuse{\nwixident{Test.suite}}{Test.suite}{NW1mR7zJ-10sN88-1}\nwendcode{}\nwbegindocs{6}\nwdocspar

\nwenddocs{}\nwbegincode{7}\sublabel{NW1mR7zJ-2ohDA9-1}\nwmargintag{{\nwtagstyle{}\subpageref{NW1mR7zJ-2ohDA9-1}}}\moddef{Register unit tests~{\nwtagstyle{}\subpageref{NW1mR7zJ-2ohDA9-1}}}\endmoddef\nwstartdeflinemarkup\nwusesondefline{\\{NWsR2q7-3MeQDa-1}}\nwenddeflinemarkup
, ClassicalSuite.suite
\nwused{\\{NWsR2q7-3MeQDa-1}}\nwendcode{}\nwbegindocs{8}\nwdocspar

\nwenddocs{}\nwbegincode{9}\sublabel{NW1mR7zJ-3lyaJm-1}\nwmargintag{{\nwtagstyle{}\subpageref{NW1mR7zJ-3lyaJm-1}}}\moddef{Unit tests to run~{\nwtagstyle{}\subpageref{NW1mR7zJ-3lyaJm-1}}}\endmoddef\nwstartdeflinemarkup\nwusesondefline{\\{NWsR2q7-2v996w-1}}\nwenddeflinemarkup
tests/ClassicalSuite.sml
\nwused{\\{NWsR2q7-2v996w-1}}\nwendcode{}\nwbegindocs{10}\nwdocspar

\begin{node}
Signed formulas $T\phi$ and $F\phi$ encode whether we have to prove
$\phi$ is true or false, respectively. Tableaux constructs a tree from
signed formulas. For example $T(A\land B)$ can extend the current
branch of the Tableaux with two new signed formulas $TA$ and $TB$. But
$F(A\land B)$ creates two new branches in the tableaux, one starts
with $FA$ and the other starts with $FB$.

Continuing in this manner, we continue writing a branch until a
contradiction has been found. More precisely, we have $TA$ and $FA$
occur within the same branch. In this situation, the branch ``closes''
and we write $\times$ to indicate this. What this means is that we
can stop working on this branch, and we can go to the next branch. If
there are no more branches to work on, then the Tableaux refutes the formula.
\end{node}


\begin{definition}
A \define{Literal} is a formula such that it is either an atomic formula
(a predicate, or equality between terms [if working in first-order logic
  with equality])
or the negation of an atomic formula.
\end{definition}

\begin{definition}
A formula $\Phi$ is said to be in \define{Negation Normal Form} if
\begin{enumerate}
\item $\Phi$ is a literal, or
\item $\Phi=A\land B$ where $A$ and $B$ are in negation normal form,
  or
\item $\Phi=A\lor B$ where $A$ and $B$ are in negation normal form, or
\item $\Phi=\forall x\ldotp A$ or $\Phi=\exists x\ldotp A$, where $A$ is again in negation normal form.
\end{enumerate}
In Classical logic, we can always transform \emph{any} formula to be
in negation normal form.
\end{definition}

\begin{node}
The basic gambit for classical tableaux is:
\begin{enumerate}
\item Given a formula $\phi$, transform it to negation normal form $\psi$;
\item Skolemize $\psi$ to obtain $\chi$;
\item Perform the tableaux method on $\chi$.
\end{enumerate}
\end{node}

\begin{node}%163
We transform a formula into negation normal form by replacing
$A\implies B$ by $(\neg A)\lor B$ and so on. We begin by transforming
the basic connectives and quantifiers by applying the transformation
to the immediate subformulas, using the rules:
\begin{enumerate}
\item $NNF(A\land B)=NNF(A)\land NNF(B)$
\item $NNF(A\lor B)=NNF(A)\lor NNF(B)$
\item $NNF(A\implies B)=NNF(\neg A)\lor NNF(B)$
\item $NNF(A\iff B)=(NNF(A)\land NNF(B))\lor(NNF(\neg A)\land NNF(\neg B))$
\item $NNF(\exists x\ldotp A)=\exists x\ldotp NNF(A)$
\item $NNF(\forall x\ldotp A)=\forall x\ldotp NNF(A)$
\end{enumerate}
\end{node}

\nwenddocs{}\nwbegincode{11}\sublabel{NW1mR7zJ-4TkWUc-1}\nwmargintag{{\nwtagstyle{}\subpageref{NW1mR7zJ-4TkWUc-1}}}\moddef{Transform formula to negation normal form~{\nwtagstyle{}\subpageref{NW1mR7zJ-4TkWUc-1}}}\endmoddef\nwstartdeflinemarkup\nwusesondefline{\\{NW1mR7zJ-l5jX3-1}}\nwenddeflinemarkup
(* nnf : \nwlinkedidentc{Formula.t}{NW3w0iRO-2jnj9W-1} -> \nwlinkedidentc{Formula.t}{NW3w0iRO-2jnj9W-1} *)
fun nnf fm =
  if \nwlinkedidentc{Formula.is_and}{NW3w0iRO-3NyJLv-3} fm
  then let val (A,B) = \nwlinkedidentc{Formula.dest_and}{NW3w0iRO-1Yt1Jl-3} fm
       in \nwlinkedidentc{Formula.mk_and}{NW3w0iRO-3SlqND-3}(nnf A, nnf B) end
  else if \nwlinkedidentc{Formula.is_or}{NW3w0iRO-3NyJLv-3} fm
  then let val (A,B) = \nwlinkedidentc{Formula.dest_or}{NW3w0iRO-1Yt1Jl-3} fm
       in \nwlinkedidentc{Formula.mk_or}{NW3w0iRO-3SlqND-3}(nnf A, nnf B) end
  else if \nwlinkedidentc{Formula.is_imp}{NW3w0iRO-3NyJLv-3} fm
  then let val (A,B) = \nwlinkedidentc{Formula.dest_imp}{NW3w0iRO-1Yt1Jl-3} fm
       in \nwlinkedidentc{Formula.mk_or}{NW3w0iRO-3SlqND-3}(nnf (\nwlinkedidentc{Formula.mk_not}{NW3w0iRO-3SlqND-2} A),
                        nnf B) end
  else if \nwlinkedidentc{Formula.is_iff}{NW3w0iRO-3NyJLv-3} fm
  then let val (A,B) = \nwlinkedidentc{Formula.dest_iff}{NW3w0iRO-1Yt1Jl-3} fm
       in \nwlinkedidentc{Formula.mk_or}{NW3w0iRO-3SlqND-3}(\nwlinkedidentc{Formula.mk_and}{NW3w0iRO-3SlqND-3}(nnf A, nnf B),
                        \nwlinkedidentc{Formula.mk_and}{NW3w0iRO-3SlqND-3}(nnf(\nwlinkedidentc{Formula.mk_not}{NW3w0iRO-3SlqND-2}(A)), nnf(\nwlinkedidentc{Formula.mk_not}{NW3w0iRO-3SlqND-2}(B)))) end
  else if Formula.is_forall fm
  then let val(x,A) = Formula.dest_forall fm
       in \nwlinkedidentc{Formula.mk_forall}{NW3w0iRO-3SlqND-4}(x, nnf A) end
  else if Formula.is_exists fm
  then let val(x,A) = Formula.dest_exists fm
       in \nwlinkedidentc{Formula.mk_exists}{NW3w0iRO-3SlqND-4}(x, nnf A) end
  \LA{}Negation normal form for negated formulas~{\nwtagstyle{}\subpageref{NW1mR7zJ-4HwyeF-1}}\RA{}

\nwused{\\{NW1mR7zJ-l5jX3-1}}\nwidentuses{\\{{\nwixident{Formula.dest{\_}and}}{Formula.dest:unand}}\\{{\nwixident{Formula.dest{\_}iff}}{Formula.dest:uniff}}\\{{\nwixident{Formula.dest{\_}imp}}{Formula.dest:unimp}}\\{{\nwixident{Formula.dest{\_}or}}{Formula.dest:unor}}\\{{\nwixident{Formula.is{\_}and}}{Formula.is:unand}}\\{{\nwixident{Formula.is{\_}iff}}{Formula.is:uniff}}\\{{\nwixident{Formula.is{\_}imp}}{Formula.is:unimp}}\\{{\nwixident{Formula.is{\_}or}}{Formula.is:unor}}\\{{\nwixident{Formula.mk{\_}and}}{Formula.mk:unand}}\\{{\nwixident{Formula.mk{\_}exists}}{Formula.mk:unexists}}\\{{\nwixident{Formula.mk{\_}forall}}{Formula.mk:unforall}}\\{{\nwixident{Formula.mk{\_}not}}{Formula.mk:unnot}}\\{{\nwixident{Formula.mk{\_}or}}{Formula.mk:unor}}\\{{\nwixident{Formula.t}}{Formula.t}}}\nwindexuse{\nwixident{Formula.dest{\_}and}}{Formula.dest:unand}{NW1mR7zJ-4TkWUc-1}\nwindexuse{\nwixident{Formula.dest{\_}iff}}{Formula.dest:uniff}{NW1mR7zJ-4TkWUc-1}\nwindexuse{\nwixident{Formula.dest{\_}imp}}{Formula.dest:unimp}{NW1mR7zJ-4TkWUc-1}\nwindexuse{\nwixident{Formula.dest{\_}or}}{Formula.dest:unor}{NW1mR7zJ-4TkWUc-1}\nwindexuse{\nwixident{Formula.is{\_}and}}{Formula.is:unand}{NW1mR7zJ-4TkWUc-1}\nwindexuse{\nwixident{Formula.is{\_}iff}}{Formula.is:uniff}{NW1mR7zJ-4TkWUc-1}\nwindexuse{\nwixident{Formula.is{\_}imp}}{Formula.is:unimp}{NW1mR7zJ-4TkWUc-1}\nwindexuse{\nwixident{Formula.is{\_}or}}{Formula.is:unor}{NW1mR7zJ-4TkWUc-1}\nwindexuse{\nwixident{Formula.mk{\_}and}}{Formula.mk:unand}{NW1mR7zJ-4TkWUc-1}\nwindexuse{\nwixident{Formula.mk{\_}exists}}{Formula.mk:unexists}{NW1mR7zJ-4TkWUc-1}\nwindexuse{\nwixident{Formula.mk{\_}forall}}{Formula.mk:unforall}{NW1mR7zJ-4TkWUc-1}\nwindexuse{\nwixident{Formula.mk{\_}not}}{Formula.mk:unnot}{NW1mR7zJ-4TkWUc-1}\nwindexuse{\nwixident{Formula.mk{\_}or}}{Formula.mk:unor}{NW1mR7zJ-4TkWUc-1}\nwindexuse{\nwixident{Formula.t}}{Formula.t}{NW1mR7zJ-4TkWUc-1}\nwendcode{}\nwbegindocs{12}\nwdocspar

\begin{node}
The rules for transforming a negated formula into negation normal form
are:
\begin{enumerate}
\item $NNF(\neg\neg A)=NNF(A)$
\item $NNF(\neg(A\land B))=NNF(\neg A)\lor NNF(\neg B)$
\item $NNF(\neg(A\lor B))=NNF(\neg A)\land NNF(\neg B)$
\item $NNF(\neg(A\implies B))=NNF(A)\land NNF(\neg B)$
\item $NNF(\neg(A\iff B))=(NNF(A)\land NNF(\neg B))\lor(NNF(\neg A)\land NNF(B))$
\item $NNF(\neg\exists x\ldotp A)=\forall x\ldotp NNF(\neg A)$
\item $NNF(\neg\forall x\ldotp A)=\exists x\ldotp NNF(\neg A)$
\end{enumerate}
\end{node}

\nwenddocs{}\nwbegincode{13}\sublabel{NW1mR7zJ-4HwyeF-1}\nwmargintag{{\nwtagstyle{}\subpageref{NW1mR7zJ-4HwyeF-1}}}\moddef{Negation normal form for negated formulas~{\nwtagstyle{}\subpageref{NW1mR7zJ-4HwyeF-1}}}\endmoddef\nwstartdeflinemarkup\nwusesondefline{\\{NW1mR7zJ-4TkWUc-1}}\nwenddeflinemarkup
else if \nwlinkedidentc{Formula.is_not}{NW3w0iRO-3NyJLv-2} fm
then let val fm' = \nwlinkedidentc{Formula.dest_not}{NW3w0iRO-1Yt1Jl-2} fm
     in if \nwlinkedidentc{Formula.is_not}{NW3w0iRO-3NyJLv-2} fm'
        then let val A = \nwlinkedidentc{Formula.dest_not}{NW3w0iRO-1Yt1Jl-2} fm'
             in nnf A end
        else if \nwlinkedidentc{Formula.is_and}{NW3w0iRO-3NyJLv-3} fm'
        then let val (A,B) = \nwlinkedidentc{Formula.dest_and}{NW3w0iRO-1Yt1Jl-3} fm'
             in \nwlinkedidentc{Formula.mk_or}{NW3w0iRO-3SlqND-3}(nnf(\nwlinkedidentc{Formula.mk_not}{NW3w0iRO-3SlqND-2}(A)),
                              nnf(\nwlinkedidentc{Formula.mk_not}{NW3w0iRO-3SlqND-2}(B))) end
        else if \nwlinkedidentc{Formula.is_or}{NW3w0iRO-3NyJLv-3} fm'
        then let val (A,B) = \nwlinkedidentc{Formula.dest_or}{NW3w0iRO-1Yt1Jl-3} fm'
             in \nwlinkedidentc{Formula.mk_and}{NW3w0iRO-3SlqND-3}(nnf(\nwlinkedidentc{Formula.mk_not}{NW3w0iRO-3SlqND-2}(A)),
                              nnf(\nwlinkedidentc{Formula.mk_not}{NW3w0iRO-3SlqND-2}(B))) end
        else if \nwlinkedidentc{Formula.is_imp}{NW3w0iRO-3NyJLv-3} fm'
        then let val (A,B) = \nwlinkedidentc{Formula.dest_imp}{NW3w0iRO-1Yt1Jl-3} fm'
             in \nwlinkedidentc{Formula.mk_and}{NW3w0iRO-3SlqND-3}(nnf A,
                               nnf(\nwlinkedidentc{Formula.mk_not}{NW3w0iRO-3SlqND-2}(B))) end
        else if \nwlinkedidentc{Formula.is_iff}{NW3w0iRO-3NyJLv-3} fm'
        then let val (A,B) = \nwlinkedidentc{Formula.dest_iff}{NW3w0iRO-1Yt1Jl-3} fm'
             in \nwlinkedidentc{Formula.mk_or}{NW3w0iRO-3SlqND-3}(\nwlinkedidentc{Formula.mk_and}{NW3w0iRO-3SlqND-3}(nnf A,
                                             nnf(\nwlinkedidentc{Formula.mk_not}{NW3w0iRO-3SlqND-2}(B))),
                              \nwlinkedidentc{Formula.mk_and}{NW3w0iRO-3SlqND-3}(nnf(\nwlinkedidentc{Formula.mk_not}{NW3w0iRO-3SlqND-2}(A)),
                                             nnf B))
             end
        else if Formula.is_exists fm'
        then let val (x,A) = Formula.dest_exists fm'
             in \nwlinkedidentc{Formula.mk_forall}{NW3w0iRO-3SlqND-4}(x, nnf(\nwlinkedidentc{Formula.mk_not}{NW3w0iRO-3SlqND-2}(A))) end
        else if Formula.is_forall fm'
        then let val (x,A) = Formula.dest_forall fm'
             in \nwlinkedidentc{Formula.mk_exists}{NW3w0iRO-3SlqND-4}(x, nnf(\nwlinkedidentc{Formula.mk_not}{NW3w0iRO-3SlqND-2}(A))) end
        else fm
     end
else fm;
\nwused{\\{NW1mR7zJ-4TkWUc-1}}\nwidentuses{\\{{\nwixident{Formula.dest{\_}and}}{Formula.dest:unand}}\\{{\nwixident{Formula.dest{\_}iff}}{Formula.dest:uniff}}\\{{\nwixident{Formula.dest{\_}imp}}{Formula.dest:unimp}}\\{{\nwixident{Formula.dest{\_}not}}{Formula.dest:unnot}}\\{{\nwixident{Formula.dest{\_}or}}{Formula.dest:unor}}\\{{\nwixident{Formula.is{\_}and}}{Formula.is:unand}}\\{{\nwixident{Formula.is{\_}iff}}{Formula.is:uniff}}\\{{\nwixident{Formula.is{\_}imp}}{Formula.is:unimp}}\\{{\nwixident{Formula.is{\_}not}}{Formula.is:unnot}}\\{{\nwixident{Formula.is{\_}or}}{Formula.is:unor}}\\{{\nwixident{Formula.mk{\_}and}}{Formula.mk:unand}}\\{{\nwixident{Formula.mk{\_}exists}}{Formula.mk:unexists}}\\{{\nwixident{Formula.mk{\_}forall}}{Formula.mk:unforall}}\\{{\nwixident{Formula.mk{\_}not}}{Formula.mk:unnot}}\\{{\nwixident{Formula.mk{\_}or}}{Formula.mk:unor}}}\nwindexuse{\nwixident{Formula.dest{\_}and}}{Formula.dest:unand}{NW1mR7zJ-4HwyeF-1}\nwindexuse{\nwixident{Formula.dest{\_}iff}}{Formula.dest:uniff}{NW1mR7zJ-4HwyeF-1}\nwindexuse{\nwixident{Formula.dest{\_}imp}}{Formula.dest:unimp}{NW1mR7zJ-4HwyeF-1}\nwindexuse{\nwixident{Formula.dest{\_}not}}{Formula.dest:unnot}{NW1mR7zJ-4HwyeF-1}\nwindexuse{\nwixident{Formula.dest{\_}or}}{Formula.dest:unor}{NW1mR7zJ-4HwyeF-1}\nwindexuse{\nwixident{Formula.is{\_}and}}{Formula.is:unand}{NW1mR7zJ-4HwyeF-1}\nwindexuse{\nwixident{Formula.is{\_}iff}}{Formula.is:uniff}{NW1mR7zJ-4HwyeF-1}\nwindexuse{\nwixident{Formula.is{\_}imp}}{Formula.is:unimp}{NW1mR7zJ-4HwyeF-1}\nwindexuse{\nwixident{Formula.is{\_}not}}{Formula.is:unnot}{NW1mR7zJ-4HwyeF-1}\nwindexuse{\nwixident{Formula.is{\_}or}}{Formula.is:unor}{NW1mR7zJ-4HwyeF-1}\nwindexuse{\nwixident{Formula.mk{\_}and}}{Formula.mk:unand}{NW1mR7zJ-4HwyeF-1}\nwindexuse{\nwixident{Formula.mk{\_}exists}}{Formula.mk:unexists}{NW1mR7zJ-4HwyeF-1}\nwindexuse{\nwixident{Formula.mk{\_}forall}}{Formula.mk:unforall}{NW1mR7zJ-4HwyeF-1}\nwindexuse{\nwixident{Formula.mk{\_}not}}{Formula.mk:unnot}{NW1mR7zJ-4HwyeF-1}\nwindexuse{\nwixident{Formula.mk{\_}or}}{Formula.mk:unor}{NW1mR7zJ-4HwyeF-1}\nwendcode{}\nwbegindocs{14}\nwdocspar

\begin{node}
Skolemization eliminates existentially quantified variables by making
them functions of universally quantified formulas. The idea is that
the following are equivalent:
\begin{enumerate}
\item $\forall x\ldotp\exists y\ldotp P[x,y]$
\item $\forall x\ldotp P[x, f(x)]$ where $f$ is a ``skolemization function''.
\end{enumerate}
This is typically understood as requiring the axiom of choice (since
$f$ is choosing ``some'' constant, in some unprescribed way).
\end{node}

\begin{node}
We need to name Skolem functions to be something new and fresh. We
gather the name of all the functions currently used.
\end{node}

\nwenddocs{}\nwbegincode{15}\sublabel{NW1mR7zJ-43eePE-1}\nwmargintag{{\nwtagstyle{}\subpageref{NW1mR7zJ-43eePE-1}}}\moddef{Gather all the functions used~{\nwtagstyle{}\subpageref{NW1mR7zJ-43eePE-1}}}\endmoddef\nwstartdeflinemarkup\nwusesondefline{\\{NW1mR7zJ-4B0sXE-1}}\nwenddeflinemarkup
(* funcs : \nwlinkedidentc{Term.t}{NW1rTmOJ-4LsRWC-1} -> (string * \nwlinkedidentc{Sort.t}{NW1rTmOJ-3ZqggC-1} * int) list *)
fun funcs (Term.Var _) = []
  | funcs (Term.Fun(f,s,args)) =
      List.foldr (fn (arg, acc) =>
                   (funcs arg)@acc)
                 [(f,s,length args)]
                 args;

(* functions : \nwlinkedidentc{Formula.t}{NW3w0iRO-2jnj9W-1} -> (string * \nwlinkedidentc{Sort.t}{NW1rTmOJ-3ZqggC-1} * int) list *)
fun functions fm =
  \nwlinkedidentc{Formula.atom_union}{NW3w0iRO-28e6MC-1}
    (fn p => let val (_,args) = \nwlinkedidentc{Formula.dest_pred}{NW3w0iRO-1Yt1Jl-1} p
             in List.foldr (fn (arg, acc) => (funcs arg)@acc)
                           []
                           args
             end)
    fm;
\nwused{\\{NW1mR7zJ-4B0sXE-1}}\nwidentuses{\\{{\nwixident{Formula.atom{\_}union}}{Formula.atom:ununion}}\\{{\nwixident{Formula.dest{\_}pred}}{Formula.dest:unpred}}\\{{\nwixident{Formula.t}}{Formula.t}}\\{{\nwixident{Sort.t}}{Sort.t}}\\{{\nwixident{Term.t}}{Term.t}}}\nwindexuse{\nwixident{Formula.atom{\_}union}}{Formula.atom:ununion}{NW1mR7zJ-43eePE-1}\nwindexuse{\nwixident{Formula.dest{\_}pred}}{Formula.dest:unpred}{NW1mR7zJ-43eePE-1}\nwindexuse{\nwixident{Formula.t}}{Formula.t}{NW1mR7zJ-43eePE-1}\nwindexuse{\nwixident{Sort.t}}{Sort.t}{NW1mR7zJ-43eePE-1}\nwindexuse{\nwixident{Term.t}}{Term.t}{NW1mR7zJ-43eePE-1}\nwendcode{}\nwbegindocs{16}\nwdocspar

\begin{node}
We will want to take the variant of an identifier to make sure it is
not currently used for a function. We ignore the arity of the function
name, we just need a ``fresh'' function name.
\end{node}

\nwenddocs{}\nwbegincode{17}\sublabel{NW1mR7zJ-2QPys7-1}\nwmargintag{{\nwtagstyle{}\subpageref{NW1mR7zJ-2QPys7-1}}}\moddef{Variant of an identifier~{\nwtagstyle{}\subpageref{NW1mR7zJ-2QPys7-1}}}\endmoddef\nwstartdeflinemarkup\nwusesondefline{\\{NW1mR7zJ-4B0sXE-1}}\nwenddeflinemarkup
(* variant : string -> \nwlinkedidentc{Sort.t}{NW1rTmOJ-3ZqggC-1} -> (string * \nwlinkedidentc{Sort.t}{NW1rTmOJ-3ZqggC-1} * int) list -> string *)
local
  fun iter n name sort funcs =
    if List.exists (fn (x,s,_) =>
                       s = sort andalso
                       (name ^ (Int.toString n)) = x)
                   funcs
    then iter (n + 1) name sort funcs
    else (name ^ (Int.toString n), sort);
in
fun variant name sort funcs : string * \nwlinkedidentc{Sort.t}{NW1rTmOJ-3ZqggC-1} =
  if List.exists (fn (x,s,_) => s = sort andalso name = x)
                 funcs
  then iter 0 name sort funcs
  else (name, sort);
end;
\nwused{\\{NW1mR7zJ-4B0sXE-1}}\nwidentuses{\\{{\nwixident{Sort.t}}{Sort.t}}}\nwindexuse{\nwixident{Sort.t}}{Sort.t}{NW1mR7zJ-2QPys7-1}\nwendcode{}\nwbegindocs{18}\nwdocspar

\begin{node}
Skolemization acts on closed formulas in negation normal form.
So negation has been pulled down to the literals. The basic steps in the procedure is:
\begin{enumerate}
\item Given $\exists x\ldotp A[x]$, collect the distinct free
  variables in a list $(x_{1},\dots,x_{n})$ (which, since we assumed
  the given formula is closed, means the free variables are
  universally quantified and this is being recursively called on a
  subformula), then introduce the Skolem
  function $f(x_{1},\dots,x_{n})$ and return $A[f(x_{1},\dots,x_{n})]$
\item Given $A\land B$, simply return the conjunction of the
  {\Tt{}skolem\nwendquote} function applied to $A$ and $B$
\item Given $A\lor B$, simply return the disjunction of the
  {\Tt{}skolem\nwendquote} function applied to $A$ and $B$
\item Given $\forall x\ldotp A[x]$, recursively apply {\Tt{}skolem\nwendquote} to
  $A$ to obtain $A'[x]$, then return $\forall x\ldotp A'[x]$.
\end{enumerate}
\end{node}

\nwenddocs{}\nwbegincode{19}\sublabel{NW1mR7zJ-4B0sXE-1}\nwmargintag{{\nwtagstyle{}\subpageref{NW1mR7zJ-4B0sXE-1}}}\moddef{Skolemize a formula~{\nwtagstyle{}\subpageref{NW1mR7zJ-4B0sXE-1}}}\endmoddef\nwstartdeflinemarkup\nwusesondefline{\\{NW1mR7zJ-4QrsTb-1}}\nwenddeflinemarkup
\LA{}Gather all the functions used~{\nwtagstyle{}\subpageref{NW1mR7zJ-43eePE-1}}\RA{}

\LA{}Variant of an identifier~{\nwtagstyle{}\subpageref{NW1mR7zJ-2QPys7-1}}\RA{}

local
  fun iter [] acc = acc
    | iter (x::xs) acc = if List.exists (\nwlinkedidentc{Term.eq}{NW1rTmOJ-3810sY-1} x) xs
                         then iter xs acc
                         else iter xs (x::acc);
in
fun dedupe vars = iter vars []
end;

(* skolem : \nwlinkedidentc{Formula.t}{NW3w0iRO-2jnj9W-1} ->  (string * \nwlinkedidentc{Sort.t}{NW1rTmOJ-3ZqggC-1} * int) list ->
     \nwlinkedidentc{Formula.t}{NW3w0iRO-2jnj9W-1} * (string * \nwlinkedidentc{Sort.t}{NW1rTmOJ-3ZqggC-1} * int) list

REQUIRES: given formula is in NNF
ENSURES: resulting formula is Skolemized, and produces the
         list of Skolem functions/constants. *)
fun skolem (fm : \nwlinkedidentc{Formula.t}{NW3w0iRO-2jnj9W-1}) (fns : (string * \nwlinkedidentc{Sort.t}{NW1rTmOJ-3ZqggC-1} * int) list)
  : \nwlinkedidentc{Formula.t}{NW3w0iRO-2jnj9W-1} * ((string * \nwlinkedidentc{Sort.t}{NW1rTmOJ-3ZqggC-1} * int) list) =
  if Formula.is_exists fm
  then let val((var as (Term.Var(y,s))), A) = Formula.dest_exists fm;
           val xs = dedupe(\nwlinkedidentc{Formula.fv}{NW3w0iRO-jRgYZ-2} fm);
           val (const as (f,s')) = variant (if List.null xs
                                 then ("c_"^y)
                                 else ("f_"^y))
                                s
                                fns;
           val fx = Term.Fun(f,s,xs);
       in skolem (\nwlinkedidentc{Formula.subst}{NW3w0iRO-1jtQRV-1} var fx A) ((f,s',length xs)::fns) end
  else if Formula.is_forall fm
  then let val (x, A) = Formula.dest_forall fm;
           val (A', fns') = skolem A fns;
       in (\nwlinkedidentc{Formula.mk_forall}{NW3w0iRO-3SlqND-4}(x, A'), fns') end
  else if \nwlinkedidentc{Formula.is_and}{NW3w0iRO-3NyJLv-3} fm
  then let val (A,B) = \nwlinkedidentc{Formula.dest_and}{NW3w0iRO-1Yt1Jl-3} fm
       in skolem2 \nwlinkedidentc{Formula.mk_and}{NW3w0iRO-3SlqND-3} (A,B) fns end
  else if \nwlinkedidentc{Formula.is_or}{NW3w0iRO-3NyJLv-3} fm
  then let val (A,B) = \nwlinkedidentc{Formula.dest_or}{NW3w0iRO-1Yt1Jl-3} fm
       in skolem2 \nwlinkedidentc{Formula.mk_or}{NW3w0iRO-3SlqND-3} (A,B) fns end
  else (fm,fns)
and skolem2 conn (A,B) fns =
    let val (A',fns') = skolem A fns;
        val (B',fns'') = skolem B fns';
    in (conn(A', B'), fns'') end;
\nwused{\\{NW1mR7zJ-4QrsTb-1}}\nwidentuses{\\{{\nwixident{Formula.dest{\_}and}}{Formula.dest:unand}}\\{{\nwixident{Formula.dest{\_}or}}{Formula.dest:unor}}\\{{\nwixident{Formula.fv}}{Formula.fv}}\\{{\nwixident{Formula.is{\_}and}}{Formula.is:unand}}\\{{\nwixident{Formula.is{\_}or}}{Formula.is:unor}}\\{{\nwixident{Formula.mk{\_}and}}{Formula.mk:unand}}\\{{\nwixident{Formula.mk{\_}forall}}{Formula.mk:unforall}}\\{{\nwixident{Formula.mk{\_}or}}{Formula.mk:unor}}\\{{\nwixident{Formula.subst}}{Formula.subst}}\\{{\nwixident{Formula.t}}{Formula.t}}\\{{\nwixident{Sort.t}}{Sort.t}}\\{{\nwixident{Term.eq}}{Term.eq}}}\nwindexuse{\nwixident{Formula.dest{\_}and}}{Formula.dest:unand}{NW1mR7zJ-4B0sXE-1}\nwindexuse{\nwixident{Formula.dest{\_}or}}{Formula.dest:unor}{NW1mR7zJ-4B0sXE-1}\nwindexuse{\nwixident{Formula.fv}}{Formula.fv}{NW1mR7zJ-4B0sXE-1}\nwindexuse{\nwixident{Formula.is{\_}and}}{Formula.is:unand}{NW1mR7zJ-4B0sXE-1}\nwindexuse{\nwixident{Formula.is{\_}or}}{Formula.is:unor}{NW1mR7zJ-4B0sXE-1}\nwindexuse{\nwixident{Formula.mk{\_}and}}{Formula.mk:unand}{NW1mR7zJ-4B0sXE-1}\nwindexuse{\nwixident{Formula.mk{\_}forall}}{Formula.mk:unforall}{NW1mR7zJ-4B0sXE-1}\nwindexuse{\nwixident{Formula.mk{\_}or}}{Formula.mk:unor}{NW1mR7zJ-4B0sXE-1}\nwindexuse{\nwixident{Formula.subst}}{Formula.subst}{NW1mR7zJ-4B0sXE-1}\nwindexuse{\nwixident{Formula.t}}{Formula.t}{NW1mR7zJ-4B0sXE-1}\nwindexuse{\nwixident{Sort.t}}{Sort.t}{NW1mR7zJ-4B0sXE-1}\nwindexuse{\nwixident{Term.eq}}{Term.eq}{NW1mR7zJ-4B0sXE-1}\nwendcode{}\nwbegindocs{20}\nwdocspar

\begin{node}
We also want to simplify formulas using various identities. The basic
structure of simplification amounts to performing these
transformations once using the {\Tt{}simplify1\nwendquote} function, and propagate
them down subformulas using the {\Tt{}simplify\nwendquote} function.
\end{node}

\nwenddocs{}\nwbegincode{21}\sublabel{NW1mR7zJ-2RorhN-1}\nwmargintag{{\nwtagstyle{}\subpageref{NW1mR7zJ-2RorhN-1}}}\moddef{Simplify formulas using classical tautologies~{\nwtagstyle{}\subpageref{NW1mR7zJ-2RorhN-1}}}\endmoddef\nwstartdeflinemarkup\nwusesondefline{\\{NW1mR7zJ-l5jX3-1}}\nwenddeflinemarkup
local
  fun simplify1 fm =
    \LA{}Simplify the body of a quantified formula~{\nwtagstyle{}\subpageref{NW1mR7zJ-3pvELX-1}}\RA{}
    \LA{}Simplify double negation~{\nwtagstyle{}\subpageref{NW1mR7zJ-tCoNW-1}}\RA{}
    \LA{}Simplify conjunction~{\nwtagstyle{}\subpageref{NW1mR7zJ-19fhtN-1}}\RA{}
    \LA{}Simplify disjunction~{\nwtagstyle{}\subpageref{NW1mR7zJ-4ZXXl0-1}}\RA{}
    \LA{}Simplify implication~{\nwtagstyle{}\subpageref{NW1mR7zJ-3HyT53-1}}\RA{}
    \LA{}Simplify logical equivalence~{\nwtagstyle{}\subpageref{NW1mR7zJ-48args-1}}\RA{}
    else fm;
in
fun simplify fm =
  if \nwlinkedidentc{Formula.is_not}{NW3w0iRO-3NyJLv-2} fm
  then let val A = \nwlinkedidentc{Formula.dest_not}{NW3w0iRO-1Yt1Jl-2} fm
       in simplify1(\nwlinkedidentc{Formula.mk_not}{NW3w0iRO-3SlqND-2}(simplify A)) end
  else if \nwlinkedidentc{Formula.is_and}{NW3w0iRO-3NyJLv-3} fm
  then let val (A,B) = \nwlinkedidentc{Formula.dest_and}{NW3w0iRO-1Yt1Jl-3} fm
       in simplify1(\nwlinkedidentc{Formula.mk_and}{NW3w0iRO-3SlqND-3}(simplify A,simplify B)) end
  else if \nwlinkedidentc{Formula.is_or}{NW3w0iRO-3NyJLv-3} fm
  then let val (A,B) = \nwlinkedidentc{Formula.dest_or}{NW3w0iRO-1Yt1Jl-3} fm
       in simplify1(\nwlinkedidentc{Formula.mk_or}{NW3w0iRO-3SlqND-3}(simplify A,simplify B)) end
  else if \nwlinkedidentc{Formula.is_imp}{NW3w0iRO-3NyJLv-3} fm
  then let val (A,B) = \nwlinkedidentc{Formula.dest_imp}{NW3w0iRO-1Yt1Jl-3} fm
       in simplify1(\nwlinkedidentc{Formula.mk_imp}{NW3w0iRO-3SlqND-3}(simplify A,simplify B)) end
  else if \nwlinkedidentc{Formula.is_iff}{NW3w0iRO-3NyJLv-3} fm
  then let val (A,B) = \nwlinkedidentc{Formula.dest_iff}{NW3w0iRO-1Yt1Jl-3} fm
       in simplify1(\nwlinkedidentc{Formula.mk_iff}{NW3w0iRO-3SlqND-3}(simplify A,simplify B)) end
  else if Formula.is_exists fm
  then let val (x,A) = Formula.dest_exists fm
       in simplify1(\nwlinkedidentc{Formula.mk_exists}{NW3w0iRO-3SlqND-4}(x, simplify A)) end
  else if Formula.is_forall fm
  then let val (x,A) = Formula.dest_forall fm
       in simplify1(\nwlinkedidentc{Formula.mk_forall}{NW3w0iRO-3SlqND-4}(x, simplify A)) end
  else fm (* falsum or predicates are already simplified *)
end;
\nwused{\\{NW1mR7zJ-l5jX3-1}}\nwidentuses{\\{{\nwixident{Formula.dest{\_}and}}{Formula.dest:unand}}\\{{\nwixident{Formula.dest{\_}iff}}{Formula.dest:uniff}}\\{{\nwixident{Formula.dest{\_}imp}}{Formula.dest:unimp}}\\{{\nwixident{Formula.dest{\_}not}}{Formula.dest:unnot}}\\{{\nwixident{Formula.dest{\_}or}}{Formula.dest:unor}}\\{{\nwixident{Formula.is{\_}and}}{Formula.is:unand}}\\{{\nwixident{Formula.is{\_}iff}}{Formula.is:uniff}}\\{{\nwixident{Formula.is{\_}imp}}{Formula.is:unimp}}\\{{\nwixident{Formula.is{\_}not}}{Formula.is:unnot}}\\{{\nwixident{Formula.is{\_}or}}{Formula.is:unor}}\\{{\nwixident{Formula.mk{\_}and}}{Formula.mk:unand}}\\{{\nwixident{Formula.mk{\_}exists}}{Formula.mk:unexists}}\\{{\nwixident{Formula.mk{\_}forall}}{Formula.mk:unforall}}\\{{\nwixident{Formula.mk{\_}iff}}{Formula.mk:uniff}}\\{{\nwixident{Formula.mk{\_}imp}}{Formula.mk:unimp}}\\{{\nwixident{Formula.mk{\_}not}}{Formula.mk:unnot}}\\{{\nwixident{Formula.mk{\_}or}}{Formula.mk:unor}}}\nwindexuse{\nwixident{Formula.dest{\_}and}}{Formula.dest:unand}{NW1mR7zJ-2RorhN-1}\nwindexuse{\nwixident{Formula.dest{\_}iff}}{Formula.dest:uniff}{NW1mR7zJ-2RorhN-1}\nwindexuse{\nwixident{Formula.dest{\_}imp}}{Formula.dest:unimp}{NW1mR7zJ-2RorhN-1}\nwindexuse{\nwixident{Formula.dest{\_}not}}{Formula.dest:unnot}{NW1mR7zJ-2RorhN-1}\nwindexuse{\nwixident{Formula.dest{\_}or}}{Formula.dest:unor}{NW1mR7zJ-2RorhN-1}\nwindexuse{\nwixident{Formula.is{\_}and}}{Formula.is:unand}{NW1mR7zJ-2RorhN-1}\nwindexuse{\nwixident{Formula.is{\_}iff}}{Formula.is:uniff}{NW1mR7zJ-2RorhN-1}\nwindexuse{\nwixident{Formula.is{\_}imp}}{Formula.is:unimp}{NW1mR7zJ-2RorhN-1}\nwindexuse{\nwixident{Formula.is{\_}not}}{Formula.is:unnot}{NW1mR7zJ-2RorhN-1}\nwindexuse{\nwixident{Formula.is{\_}or}}{Formula.is:unor}{NW1mR7zJ-2RorhN-1}\nwindexuse{\nwixident{Formula.mk{\_}and}}{Formula.mk:unand}{NW1mR7zJ-2RorhN-1}\nwindexuse{\nwixident{Formula.mk{\_}exists}}{Formula.mk:unexists}{NW1mR7zJ-2RorhN-1}\nwindexuse{\nwixident{Formula.mk{\_}forall}}{Formula.mk:unforall}{NW1mR7zJ-2RorhN-1}\nwindexuse{\nwixident{Formula.mk{\_}iff}}{Formula.mk:uniff}{NW1mR7zJ-2RorhN-1}\nwindexuse{\nwixident{Formula.mk{\_}imp}}{Formula.mk:unimp}{NW1mR7zJ-2RorhN-1}\nwindexuse{\nwixident{Formula.mk{\_}not}}{Formula.mk:unnot}{NW1mR7zJ-2RorhN-1}\nwindexuse{\nwixident{Formula.mk{\_}or}}{Formula.mk:unor}{NW1mR7zJ-2RorhN-1}\nwendcode{}\nwbegindocs{22}\nwdocspar

\begin{node}
We simplify $\forall x\ldotp A[x]$ by simplifying $A[x]$ to $A'[x]$
(whatever that might be), then returning $\forall x\ldotp A'[x]$.
Similar simplification applies to existentially quantified formulas.
\end{node}

\nwenddocs{}\nwbegincode{23}\sublabel{NW1mR7zJ-3pvELX-1}\nwmargintag{{\nwtagstyle{}\subpageref{NW1mR7zJ-3pvELX-1}}}\moddef{Simplify the body of a quantified formula~{\nwtagstyle{}\subpageref{NW1mR7zJ-3pvELX-1}}}\endmoddef\nwstartdeflinemarkup\nwusesondefline{\\{NW1mR7zJ-2RorhN-1}}\nwenddeflinemarkup
if Formula.is_forall fm
then let val (x,A) = Formula.dest_forall fm
     in if List.exists (\nwlinkedidentc{Term.eq}{NW1rTmOJ-3810sY-1} x) (\nwlinkedidentc{Formula.fv}{NW3w0iRO-jRgYZ-2} A)
        then fm
        else A
     end
else if Formula.is_exists fm
then let val (x,A) = Formula.dest_exists fm
     in if List.exists (\nwlinkedidentc{Term.eq}{NW1rTmOJ-3810sY-1} x) (\nwlinkedidentc{Formula.fv}{NW3w0iRO-jRgYZ-2} A)
        then fm
        else A
     end
\nwused{\\{NW1mR7zJ-2RorhN-1}}\nwidentuses{\\{{\nwixident{Formula.fv}}{Formula.fv}}\\{{\nwixident{Term.eq}}{Term.eq}}}\nwindexuse{\nwixident{Formula.fv}}{Formula.fv}{NW1mR7zJ-3pvELX-1}\nwindexuse{\nwixident{Term.eq}}{Term.eq}{NW1mR7zJ-3pvELX-1}\nwendcode{}\nwbegindocs{24}\nwdocspar

\begin{node}
We simplify $\neg\neg A$ to be just $A$.
\end{node}

\nwenddocs{}\nwbegincode{25}\sublabel{NW1mR7zJ-tCoNW-1}\nwmargintag{{\nwtagstyle{}\subpageref{NW1mR7zJ-tCoNW-1}}}\moddef{Simplify double negation~{\nwtagstyle{}\subpageref{NW1mR7zJ-tCoNW-1}}}\endmoddef\nwstartdeflinemarkup\nwusesondefline{\\{NW1mR7zJ-2RorhN-1}}\nwenddeflinemarkup
else if \nwlinkedidentc{Formula.is_not}{NW3w0iRO-3NyJLv-2} fm
then let val A = \nwlinkedidentc{Formula.dest_not}{NW3w0iRO-1Yt1Jl-2} fm
     in if \nwlinkedidentc{Formula.is_not}{NW3w0iRO-3NyJLv-2} A
        then \nwlinkedidentc{Formula.dest_not}{NW3w0iRO-1Yt1Jl-2} A
        else fm
     end
\nwused{\\{NW1mR7zJ-2RorhN-1}}\nwidentuses{\\{{\nwixident{Formula.dest{\_}not}}{Formula.dest:unnot}}\\{{\nwixident{Formula.is{\_}not}}{Formula.is:unnot}}}\nwindexuse{\nwixident{Formula.dest{\_}not}}{Formula.dest:unnot}{NW1mR7zJ-tCoNW-1}\nwindexuse{\nwixident{Formula.is{\_}not}}{Formula.is:unnot}{NW1mR7zJ-tCoNW-1}\nwendcode{}\nwbegindocs{26}\nwdocspar

\begin{node}
We simplify $\bot\land B$ and $A\land\bot$ to $\bot$. In the cases
where neither conjunct is falsum: 
$\top\land B$ (i.e., $A=\top$) simplifies to $B$, and $A\land\top$
(i.e., $B=\top$) simplifies to $A$.
\end{node}

\nwenddocs{}\nwbegincode{27}\sublabel{NW1mR7zJ-19fhtN-1}\nwmargintag{{\nwtagstyle{}\subpageref{NW1mR7zJ-19fhtN-1}}}\moddef{Simplify conjunction~{\nwtagstyle{}\subpageref{NW1mR7zJ-19fhtN-1}}}\endmoddef\nwstartdeflinemarkup\nwusesondefline{\\{NW1mR7zJ-2RorhN-1}}\nwenddeflinemarkup
else if \nwlinkedidentc{Formula.is_and}{NW3w0iRO-3NyJLv-3} fm
then let val (A,B) = \nwlinkedidentc{Formula.dest_and}{NW3w0iRO-1Yt1Jl-3} fm
     in if \nwlinkedidentc{Formula.is_falsum}{NW3w0iRO-3NyJLv-1} A orelse \nwlinkedidentc{Formula.is_falsum}{NW3w0iRO-3NyJLv-1} B
        then \nwlinkedidentc{Formula.mk_falsum}{NW3w0iRO-3SlqND-1} ()
        else if \nwlinkedidentc{Formula.is_verum}{NW3w0iRO-3NyJLv-1} A
        then B
        else if \nwlinkedidentc{Formula.is_verum}{NW3w0iRO-3NyJLv-1} B
        then A
        else fm
     end
\nwused{\\{NW1mR7zJ-2RorhN-1}}\nwidentuses{\\{{\nwixident{Formula.dest{\_}and}}{Formula.dest:unand}}\\{{\nwixident{Formula.is{\_}and}}{Formula.is:unand}}\\{{\nwixident{Formula.is{\_}falsum}}{Formula.is:unfalsum}}\\{{\nwixident{Formula.is{\_}verum}}{Formula.is:unverum}}\\{{\nwixident{Formula.mk{\_}falsum}}{Formula.mk:unfalsum}}}\nwindexuse{\nwixident{Formula.dest{\_}and}}{Formula.dest:unand}{NW1mR7zJ-19fhtN-1}\nwindexuse{\nwixident{Formula.is{\_}and}}{Formula.is:unand}{NW1mR7zJ-19fhtN-1}\nwindexuse{\nwixident{Formula.is{\_}falsum}}{Formula.is:unfalsum}{NW1mR7zJ-19fhtN-1}\nwindexuse{\nwixident{Formula.is{\_}verum}}{Formula.is:unverum}{NW1mR7zJ-19fhtN-1}\nwindexuse{\nwixident{Formula.mk{\_}falsum}}{Formula.mk:unfalsum}{NW1mR7zJ-19fhtN-1}\nwendcode{}\nwbegindocs{28}\nwdocspar

\begin{node}
We simplify $\bot\lor A$ and $A\lor\bot$ to $A$, and $\top\lor A$
and $A\lor\top$ to $\top$. Otherwise we don't touch the formula.
\end{node}

\nwenddocs{}\nwbegincode{29}\sublabel{NW1mR7zJ-4ZXXl0-1}\nwmargintag{{\nwtagstyle{}\subpageref{NW1mR7zJ-4ZXXl0-1}}}\moddef{Simplify disjunction~{\nwtagstyle{}\subpageref{NW1mR7zJ-4ZXXl0-1}}}\endmoddef\nwstartdeflinemarkup\nwusesondefline{\\{NW1mR7zJ-2RorhN-1}}\nwenddeflinemarkup
else if \nwlinkedidentc{Formula.is_or}{NW3w0iRO-3NyJLv-3} fm
then let val (A,B) = \nwlinkedidentc{Formula.dest_or}{NW3w0iRO-1Yt1Jl-3} fm
     in if \nwlinkedidentc{Formula.is_falsum}{NW3w0iRO-3NyJLv-1} A orelse \nwlinkedidentc{Formula.is_verum}{NW3w0iRO-3NyJLv-1} B
        then B
        else if \nwlinkedidentc{Formula.is_falsum}{NW3w0iRO-3NyJLv-1} B orelse \nwlinkedidentc{Formula.is_verum}{NW3w0iRO-3NyJLv-1} A
        then A
        else fm
     end
\nwused{\\{NW1mR7zJ-2RorhN-1}}\nwidentuses{\\{{\nwixident{Formula.dest{\_}or}}{Formula.dest:unor}}\\{{\nwixident{Formula.is{\_}falsum}}{Formula.is:unfalsum}}\\{{\nwixident{Formula.is{\_}or}}{Formula.is:unor}}\\{{\nwixident{Formula.is{\_}verum}}{Formula.is:unverum}}}\nwindexuse{\nwixident{Formula.dest{\_}or}}{Formula.dest:unor}{NW1mR7zJ-4ZXXl0-1}\nwindexuse{\nwixident{Formula.is{\_}falsum}}{Formula.is:unfalsum}{NW1mR7zJ-4ZXXl0-1}\nwindexuse{\nwixident{Formula.is{\_}or}}{Formula.is:unor}{NW1mR7zJ-4ZXXl0-1}\nwindexuse{\nwixident{Formula.is{\_}verum}}{Formula.is:unverum}{NW1mR7zJ-4ZXXl0-1}\nwendcode{}\nwbegindocs{30}\nwdocspar

\begin{node}
We simplify $\bot\implies A$ and $A\implies\top$ both to just $\top$.
We similarly simplify $\top\implies B$ to just $B$. Otherwise we leave
the formula alone.
\end{node}

\nwenddocs{}\nwbegincode{31}\sublabel{NW1mR7zJ-3HyT53-1}\nwmargintag{{\nwtagstyle{}\subpageref{NW1mR7zJ-3HyT53-1}}}\moddef{Simplify implication~{\nwtagstyle{}\subpageref{NW1mR7zJ-3HyT53-1}}}\endmoddef\nwstartdeflinemarkup\nwusesondefline{\\{NW1mR7zJ-2RorhN-1}}\nwenddeflinemarkup
else if \nwlinkedidentc{Formula.is_imp}{NW3w0iRO-3NyJLv-3} fm
then let val (A,B) = \nwlinkedidentc{Formula.dest_imp}{NW3w0iRO-1Yt1Jl-3} fm
     in if \nwlinkedidentc{Formula.is_falsum}{NW3w0iRO-3NyJLv-1} A orelse \nwlinkedidentc{Formula.is_verum}{NW3w0iRO-3NyJLv-1} B
        then \nwlinkedidentc{Formula.mk_verum}{NW3w0iRO-3SlqND-1} ()
        else if \nwlinkedidentc{Formula.is_verum}{NW3w0iRO-3NyJLv-1} A
        then B
        else fm
     end
\nwused{\\{NW1mR7zJ-2RorhN-1}}\nwidentuses{\\{{\nwixident{Formula.dest{\_}imp}}{Formula.dest:unimp}}\\{{\nwixident{Formula.is{\_}falsum}}{Formula.is:unfalsum}}\\{{\nwixident{Formula.is{\_}imp}}{Formula.is:unimp}}\\{{\nwixident{Formula.is{\_}verum}}{Formula.is:unverum}}\\{{\nwixident{Formula.mk{\_}verum}}{Formula.mk:unverum}}}\nwindexuse{\nwixident{Formula.dest{\_}imp}}{Formula.dest:unimp}{NW1mR7zJ-3HyT53-1}\nwindexuse{\nwixident{Formula.is{\_}falsum}}{Formula.is:unfalsum}{NW1mR7zJ-3HyT53-1}\nwindexuse{\nwixident{Formula.is{\_}imp}}{Formula.is:unimp}{NW1mR7zJ-3HyT53-1}\nwindexuse{\nwixident{Formula.is{\_}verum}}{Formula.is:unverum}{NW1mR7zJ-3HyT53-1}\nwindexuse{\nwixident{Formula.mk{\_}verum}}{Formula.mk:unverum}{NW1mR7zJ-3HyT53-1}\nwendcode{}\nwbegindocs{32}\nwdocspar

\begin{node}
We simplify $A\iff B$ according to the following rules:
\begin{enumerate}
\item $\top\iff B$ becomes $B$
\item $A\iff\top$ becomes $A$
\item $\bot\iff B$ becomes $\neg B$
\item $A\iff\bot$ becomes $\neg A$
\end{enumerate}
Otherwise, we just leave the formula alone.
\end{node}

\nwenddocs{}\nwbegincode{33}\sublabel{NW1mR7zJ-48args-1}\nwmargintag{{\nwtagstyle{}\subpageref{NW1mR7zJ-48args-1}}}\moddef{Simplify logical equivalence~{\nwtagstyle{}\subpageref{NW1mR7zJ-48args-1}}}\endmoddef\nwstartdeflinemarkup\nwusesondefline{\\{NW1mR7zJ-2RorhN-1}}\nwenddeflinemarkup
else if \nwlinkedidentc{Formula.is_iff}{NW3w0iRO-3NyJLv-3} fm
then let val (A,B) = \nwlinkedidentc{Formula.dest_iff}{NW3w0iRO-1Yt1Jl-3} fm
     in if \nwlinkedidentc{Formula.is_verum}{NW3w0iRO-3NyJLv-1} A
        then B
        else if \nwlinkedidentc{Formula.is_verum}{NW3w0iRO-3NyJLv-1} B
        then A
        else if \nwlinkedidentc{Formula.is_falsum}{NW3w0iRO-3NyJLv-1} A
        then \nwlinkedidentc{Formula.mk_not}{NW3w0iRO-3SlqND-2} B
        else if \nwlinkedidentc{Formula.is_falsum}{NW3w0iRO-3NyJLv-1} B
        then \nwlinkedidentc{Formula.mk_not}{NW3w0iRO-3SlqND-2} A
        else fm
     end
\nwused{\\{NW1mR7zJ-2RorhN-1}}\nwidentuses{\\{{\nwixident{Formula.dest{\_}iff}}{Formula.dest:uniff}}\\{{\nwixident{Formula.is{\_}falsum}}{Formula.is:unfalsum}}\\{{\nwixident{Formula.is{\_}iff}}{Formula.is:uniff}}\\{{\nwixident{Formula.is{\_}verum}}{Formula.is:unverum}}\\{{\nwixident{Formula.mk{\_}not}}{Formula.mk:unnot}}}\nwindexuse{\nwixident{Formula.dest{\_}iff}}{Formula.dest:uniff}{NW1mR7zJ-48args-1}\nwindexuse{\nwixident{Formula.is{\_}falsum}}{Formula.is:unfalsum}{NW1mR7zJ-48args-1}\nwindexuse{\nwixident{Formula.is{\_}iff}}{Formula.is:uniff}{NW1mR7zJ-48args-1}\nwindexuse{\nwixident{Formula.is{\_}verum}}{Formula.is:unverum}{NW1mR7zJ-48args-1}\nwindexuse{\nwixident{Formula.mk{\_}not}}{Formula.mk:unnot}{NW1mR7zJ-48args-1}\nwendcode{}\nwbegindocs{34}\nwdocspar

\begin{node}
Skolemization then applies the algorithm to replace existentially
quantified variables with skolem constants, to the negation normal
form of the given formula (after simplifying the formula).
\end{node}

\nwenddocs{}\nwbegincode{35}\sublabel{NW1mR7zJ-4QrsTb-1}\nwmargintag{{\nwtagstyle{}\subpageref{NW1mR7zJ-4QrsTb-1}}}\moddef{Skolemization~{\nwtagstyle{}\subpageref{NW1mR7zJ-4QrsTb-1}}}\endmoddef\nwstartdeflinemarkup\nwusesondefline{\\{NW1mR7zJ-l5jX3-1}}\nwenddeflinemarkup
\LA{}Skolemize a formula~{\nwtagstyle{}\subpageref{NW1mR7zJ-4B0sXE-1}}\RA{}

(* askolemize : \nwlinkedidentc{Formula.t}{NW3w0iRO-2jnj9W-1} -> \nwlinkedidentc{Formula.t}{NW3w0iRO-2jnj9W-1}

REQUIRES: The given formula is a closed formula.
ENSURES: The resulting formula is in Negation Normal Form with
         all existential quantifiers replaced by Skolem
         functions/constants. *)
fun askolemize fm =
  let val (result, _) = skolem (nnf(simplify(fm))) (functions fm)
  in result end;
\nwused{\\{NW1mR7zJ-l5jX3-1}}\nwidentuses{\\{{\nwixident{Formula.t}}{Formula.t}}}\nwindexuse{\nwixident{Formula.t}}{Formula.t}{NW1mR7zJ-4QrsTb-1}\nwendcode{}\nwbegindocs{36}\nwdocspar

\section{Tableaux}

\begin{node}
Most tableaux methods follows \leanTAP's algorithm~\cite{beckert1995leantap}.
The basic tableaux method works as follows:
\begin{enumerate}
\item for $A\land B$, separately assume $A$ and $B$; this amounts to
  working through $A$, then working through $B$, then working through
  the remaining unexplored subformulas.
\item for $A\lor B$, we ``branch'' --- that is to say, we perform two
  separate refutations, one for $A$, and another for $B$;
\item for $\forall x\ldotp A[x]$, introduce a fresh variable $y$ and
  assume $A[y]$, but we keep the original $\forall x\ldotp A[x]$
  in case multiple instances are needed.
\item for a literal $L$, try to refute it based on the literals
  already encountered (which is the only way to close the branch);
  otherwise, we add $L$ to the ``database'' of
  literals, and look at the next unexpanded formula.
\end{enumerate}
There is no ordering to the subformulas being explored. Smullyan noted
that it's probably more optimal to explore the conjunctions first,
then the disjunctions. This would probably be low hanging fruit.
\end{node}

\nwenddocs{}\nwbegincode{37}\sublabel{NW1mR7zJ-1k5nOE-1}\nwmargintag{{\nwtagstyle{}\subpageref{NW1mR7zJ-1k5nOE-1}}}\moddef{Tableau method~{\nwtagstyle{}\subpageref{NW1mR7zJ-1k5nOE-1}}}\endmoddef\nwstartdeflinemarkup\nwusesondefline{\\{NW1mR7zJ-yCIr5-1}}\nwenddeflinemarkup
fun tryfind f [] = raise Fail "tryfind"
  | tryfind f (x::xs) = (f x
                         handle _ => tryfind f xs);

(* tableau : (\nwlinkedidentc{Formula.t}{NW3w0iRO-2jnj9W-1} list * \nwlinkedidentc{Formula.t}{NW3w0iRO-2jnj9W-1} list * int) ->
                (Unif.\nwlinkedidentc{Env.t}{NW1BOcHB-25AKca-1} * int -> 'a) ->
                Unif.\nwlinkedidentc{Env.t}{NW1BOcHB-25AKca-1} * int -> 'a *)
fun tableau (fms, lits, n) (cont : Unif.\nwlinkedidentc{Env.t}{NW1BOcHB-25AKca-1} * int -> 'a) (env, k) =
  if n < 0 then raise Fail "no proof at this level"
  else
    (case fms of
      [] => raise Fail "tableau: no proof"
    | (phi::unexp) =>
        if \nwlinkedidentc{Formula.is_and}{NW3w0iRO-3NyJLv-3} phi
        then let val (A,B) = \nwlinkedidentc{Formula.dest_and}{NW3w0iRO-1Yt1Jl-3} phi
             in tableau (A::B::unexp,lits,n) cont (env, k) end
        else if \nwlinkedidentc{Formula.is_or}{NW3w0iRO-3NyJLv-3} phi
        then let val (A,B) = \nwlinkedidentc{Formula.dest_or}{NW3w0iRO-1Yt1Jl-3} phi
             in tableau (A::unexp,lits,n)
                        (tableau (B::unexp,lits,n) cont)
                        (env, k)
             end
        else if Formula.is_forall phi
        then let val(x,A) = Formula.dest_forall phi;
                 val y = Term.Var("_"^(Int.toString k), \nwlinkedidentc{Term.sort}{NW1rTmOJ-1WQNps-1}(x));
                 val A' = \nwlinkedidentc{Formula.subst}{NW3w0iRO-1jtQRV-1} x y A;
             in tableau (A'::unexp@[\nwlinkedidentc{Formula.mk_forall}{NW3w0iRO-3SlqND-4}(x,A)],lits,n-1) cont (env, k+1) end
        else
         (tryfind (fn l => cont(\nwlinkedidentc{Unif.y_complements}{NW1BOcHB-3w0JS3-1} env (phi,l),k))
                  lits
         handle _ => tableau(unexp,phi::lits,n) cont (env,k)));
\nwused{\\{NW1mR7zJ-yCIr5-1}}\nwidentuses{\\{{\nwixident{Env.t}}{Env.t}}\\{{\nwixident{Formula.dest{\_}and}}{Formula.dest:unand}}\\{{\nwixident{Formula.dest{\_}or}}{Formula.dest:unor}}\\{{\nwixident{Formula.is{\_}and}}{Formula.is:unand}}\\{{\nwixident{Formula.is{\_}or}}{Formula.is:unor}}\\{{\nwixident{Formula.mk{\_}forall}}{Formula.mk:unforall}}\\{{\nwixident{Formula.subst}}{Formula.subst}}\\{{\nwixident{Formula.t}}{Formula.t}}\\{{\nwixident{Term.sort}}{Term.sort}}\\{{\nwixident{Unif.y{\_}complements}}{Unif.y:uncomplements}}}\nwindexuse{\nwixident{Env.t}}{Env.t}{NW1mR7zJ-1k5nOE-1}\nwindexuse{\nwixident{Formula.dest{\_}and}}{Formula.dest:unand}{NW1mR7zJ-1k5nOE-1}\nwindexuse{\nwixident{Formula.dest{\_}or}}{Formula.dest:unor}{NW1mR7zJ-1k5nOE-1}\nwindexuse{\nwixident{Formula.is{\_}and}}{Formula.is:unand}{NW1mR7zJ-1k5nOE-1}\nwindexuse{\nwixident{Formula.is{\_}or}}{Formula.is:unor}{NW1mR7zJ-1k5nOE-1}\nwindexuse{\nwixident{Formula.mk{\_}forall}}{Formula.mk:unforall}{NW1mR7zJ-1k5nOE-1}\nwindexuse{\nwixident{Formula.subst}}{Formula.subst}{NW1mR7zJ-1k5nOE-1}\nwindexuse{\nwixident{Formula.t}}{Formula.t}{NW1mR7zJ-1k5nOE-1}\nwindexuse{\nwixident{Term.sort}}{Term.sort}{NW1mR7zJ-1k5nOE-1}\nwindexuse{\nwixident{Unif.y{\_}complements}}{Unif.y:uncomplements}{NW1mR7zJ-1k5nOE-1}\nwendcode{}\nwbegindocs{38}\nwdocspar

\begin{node}
The basic idea is to proving a formula $\phi$ classically with method
of tableaux amounts to refuting $\neg\phi$ with tableaux. We need to
universally quantify free variables, then negate the result, then
Skolemize the result, before applying the tableaux method.
\end{node}

\nwenddocs{}\nwbegincode{39}\sublabel{NW1mR7zJ-3SEP5w-1}\nwmargintag{{\nwtagstyle{}\subpageref{NW1mR7zJ-3SEP5w-1}}}\moddef{Tableaux refutation~{\nwtagstyle{}\subpageref{NW1mR7zJ-3SEP5w-1}}}\endmoddef\nwstartdeflinemarkup\nwusesondefline{\\{NW1mR7zJ-yCIr5-1}}\nwenddeflinemarkup
fun deepen f n maxiter =
  (print("Searching with depth limit "^
         (Int.toString n)^"\\n");
   f n)
  handle (Fail _) => if n < maxiter then deepen f (n + 1) maxiter
                     else raise Fail "maximum iterations encountered";

local
  fun id x = x;
in
fun tabrefute fms maxiter =
  deepen (fn n =>
             (tableau (fms,[],n)
                      id
                      (Unif.\nwlinkedidentc{Env.undefined}{NW1BOcHB-25AKca-1} (),0);
              n))
         0
         maxiter
end;
\nwused{\\{NW1mR7zJ-yCIr5-1}}\nwidentuses{\\{{\nwixident{Env.undefined}}{Env.undefined}}}\nwindexuse{\nwixident{Env.undefined}}{Env.undefined}{NW1mR7zJ-3SEP5w-1}\nwendcode{}\nwbegindocs{40}\nwdocspar

\begin{node}
We can combine everything to give the user {\Tt{}\nwlinkedidentq{Classical.tab}{NW1mR7zJ-yCIr5-1}\nwendquote}, which
attempts to prove the validity of the given formula by refuting its
negation. The integer parameter supplied limits the number of
iterations performed, to guarantee it will terminate.

One possibly optimization is to transform the Skolemized formula into
Disjunctive normal form, then applying the tableaux method to each
disjunct.
\end{node}

\nwenddocs{}\nwbegincode{41}\sublabel{NW1mR7zJ-yCIr5-1}\nwmargintag{{\nwtagstyle{}\subpageref{NW1mR7zJ-yCIr5-1}}}\moddef{Tableaux refuter~{\nwtagstyle{}\subpageref{NW1mR7zJ-yCIr5-1}}}\endmoddef\nwstartdeflinemarkup\nwusesondefline{\\{NW1mR7zJ-l5jX3-1}}\nwenddeflinemarkup
\LA{}Tableau method~{\nwtagstyle{}\subpageref{NW1mR7zJ-1k5nOE-1}}\RA{}

\LA{}Tableaux refutation~{\nwtagstyle{}\subpageref{NW1mR7zJ-3SEP5w-1}}\RA{}

(* generalize : \nwlinkedidentc{Formula.t}{NW3w0iRO-2jnj9W-1} -> \nwlinkedidentc{Formula.t}{NW3w0iRO-2jnj9W-1} *)
fun generalize fm =
  List.foldr \nwlinkedidentc{Formula.mk_forall}{NW3w0iRO-3SlqND-4} fm (dedupe (\nwlinkedidentc{Formula.fv}{NW3w0iRO-jRgYZ-2} fm));

fun tab fm maxiter =
  let val fm' = askolemize(\nwlinkedidentc{Formula.mk_not}{NW3w0iRO-3SlqND-2}(generalize fm))
  in if \nwlinkedidentc{Formula.is_falsum}{NW3w0iRO-3NyJLv-1} fm' then 0
     else tabrefute [fm'] maxiter
  end;
\nwindexdefn{\nwixident{Classical.tab}}{Classical.tab}{NW1mR7zJ-yCIr5-1}\eatline
\nwused{\\{NW1mR7zJ-l5jX3-1}}\nwidentdefs{\\{{\nwixident{Classical.tab}}{Classical.tab}}}\nwidentuses{\\{{\nwixident{Formula.fv}}{Formula.fv}}\\{{\nwixident{Formula.is{\_}falsum}}{Formula.is:unfalsum}}\\{{\nwixident{Formula.mk{\_}forall}}{Formula.mk:unforall}}\\{{\nwixident{Formula.mk{\_}not}}{Formula.mk:unnot}}\\{{\nwixident{Formula.t}}{Formula.t}}}\nwindexuse{\nwixident{Formula.fv}}{Formula.fv}{NW1mR7zJ-yCIr5-1}\nwindexuse{\nwixident{Formula.is{\_}falsum}}{Formula.is:unfalsum}{NW1mR7zJ-yCIr5-1}\nwindexuse{\nwixident{Formula.mk{\_}forall}}{Formula.mk:unforall}{NW1mR7zJ-yCIr5-1}\nwindexuse{\nwixident{Formula.mk{\_}not}}{Formula.mk:unnot}{NW1mR7zJ-yCIr5-1}\nwindexuse{\nwixident{Formula.t}}{Formula.t}{NW1mR7zJ-yCIr5-1}\nwendcode{}\nwbegindocs{42}\nwdocspar
\begin{node}
As a couple of unit tests, we should check that the simpler problems
found in Pelletier's problem set ``work''.
\end{node}

\nwenddocs{}\nwbegincode{43}\sublabel{NW1mR7zJ-4QAFgj-1}\nwmargintag{{\nwtagstyle{}\subpageref{NW1mR7zJ-4QAFgj-1}}}\moddef{Unit tests for \code{}Classical\edoc{}~{\nwtagstyle{}\subpageref{NW1mR7zJ-4QAFgj-1}}}\endmoddef\nwstartdeflinemarkup\nwusesondefline{\\{NW1mR7zJ-10sN88-1}}\nwenddeflinemarkup
\nwlinkedidentc{Test.mk}{NWsR2q7-3fh2dK-1} "Pelletier #1" (fn () =>
  let
    val fm = (A ==> B) <=> ((~ B) ==> (~ A));
    val expected = 0;
    val actual = \nwlinkedidentc{Classical.tab}{NW1mR7zJ-yCIr5-1} fm 3;
  in
    \nwlinkedidentc{Assert.eq}{NWsR2q7-4gQMdZ-1} expected
              actual
              (concat ["EXPECTED: ",
                       Int.toString expected,
                       "\\nACTUAL: ",
                       Int.toString actual])
  end)
, \nwlinkedidentc{Test.mk}{NWsR2q7-3fh2dK-1} "Pelletier #2" (fn () =>
    let
      val fm = (~ (~ A)) <=> A;
      val expected = 0;
      val actual = \nwlinkedidentc{Classical.tab}{NW1mR7zJ-yCIr5-1} fm 3;
    in
      \nwlinkedidentc{Assert.eq}{NWsR2q7-4gQMdZ-1} expected
                actual
                (concat ["EXPECTED: ",
                         Int.toString expected,
                         "\\nACTUAL: ",
                         Int.toString actual])
    end)
, \nwlinkedidentc{Test.mk}{NWsR2q7-3fh2dK-1} "Pelletier #3" (fn () =>
    let
      val fm = (~(A ==> B)) ==> (B ==> A);
      val expected = 0;
      val actual = \nwlinkedidentc{Classical.tab}{NW1mR7zJ-yCIr5-1} fm 3;
    in
      \nwlinkedidentc{Assert.eq}{NWsR2q7-4gQMdZ-1} expected
                actual
                (concat ["EXPECTED: ",
                         Int.toString expected,
                         "\\nACTUAL: ",
                         Int.toString actual])
    end)
, \nwlinkedidentc{Test.mk}{NWsR2q7-3fh2dK-1} "Pelletier #4" (fn () =>
    let
      val fm = (((~ A) ==> B) <=> ((~ B) ==> A));
      val expected = 0;
      val actual = \nwlinkedidentc{Classical.tab}{NW1mR7zJ-yCIr5-1} fm 3;
    in
      \nwlinkedidentc{Assert.eq}{NWsR2q7-4gQMdZ-1} expected
                actual
                (concat ["EXPECTED: ",
                         Int.toString expected,
                         "\\nACTUAL: ",
                         Int.toString actual])
    end)
, \nwlinkedidentc{Test.mk}{NWsR2q7-3fh2dK-1} "Pelletier #5" (fn () =>
    let
      val fm = \nwlinkedidentc{Formula.mk_or}{NW3w0iRO-3SlqND-3}(A, ~ A);
      val expected = 0;
      val actual = \nwlinkedidentc{Classical.tab}{NW1mR7zJ-yCIr5-1} fm 3;
    in
      \nwlinkedidentc{Assert.eq}{NWsR2q7-4gQMdZ-1} expected
                actual
                (concat ["EXPECTED: ",
                         Int.toString expected,
                         "\\nACTUAL: ",
                         Int.toString actual])
    end)
, \nwlinkedidentc{Test.mk}{NWsR2q7-3fh2dK-1} "Pelletier #18" (fn () =>
    let
      val x = Term.Var("x",\nwlinkedidentc{Sort.IND}{NW1rTmOJ-3ZqggC-1});
      val y = Term.Var("y",\nwlinkedidentc{Sort.IND}{NW1rTmOJ-3ZqggC-1});
      fun P t = \nwlinkedidentc{Formula.mk_pred}{NW3w0iRO-3SlqND-1}("P", [t]);
      val fm = (exists(y, all(x, (P y) ==> (P x))));
      val expected = 2;
      val actual = \nwlinkedidentc{Classical.tab}{NW1mR7zJ-yCIr5-1} fm 5;
    in
      \nwlinkedidentc{Assert.eq}{NWsR2q7-4gQMdZ-1} expected
                actual
                (concat ["EXPECTED: ",
                         Int.toString expected,
                         "\\nACTUAL: ",
                         Int.toString actual])
    end)
, \nwlinkedidentc{Test.mk}{NWsR2q7-3fh2dK-1} "Pelletier #19" (fn () =>
    let
      val x = Term.Var("x",\nwlinkedidentc{Sort.IND}{NW1rTmOJ-3ZqggC-1});
      val y = Term.Var("y",\nwlinkedidentc{Sort.IND}{NW1rTmOJ-3ZqggC-1});
      val z = Term.Var("z",\nwlinkedidentc{Sort.IND}{NW1rTmOJ-3ZqggC-1});
      fun P t = \nwlinkedidentc{Formula.mk_pred}{NW3w0iRO-3SlqND-1}("P", [t]);
      fun Q t = \nwlinkedidentc{Formula.mk_pred}{NW3w0iRO-3SlqND-1}("Q", [t]);
      val fm = exists(x,
                      all(y,
                          all(z, ((P y) ==> (Q z))
                                   ==>
                                 (P x) ==> (Q x))));
      val expected = 2;
      val actual = \nwlinkedidentc{Classical.tab}{NW1mR7zJ-yCIr5-1} fm 5;
    in
      \nwlinkedidentc{Assert.eq}{NWsR2q7-4gQMdZ-1} expected
                actual
                (concat ["EXPECTED: ",
                         Int.toString expected,
                         "\\nACTUAL: ",
                         Int.toString actual])
    end)
\nwused{\\{NW1mR7zJ-10sN88-1}}\nwidentuses{\\{{\nwixident{Assert.eq}}{Assert.eq}}\\{{\nwixident{Classical.tab}}{Classical.tab}}\\{{\nwixident{Formula.mk{\_}or}}{Formula.mk:unor}}\\{{\nwixident{Formula.mk{\_}pred}}{Formula.mk:unpred}}\\{{\nwixident{Sort.IND}}{Sort.IND}}\\{{\nwixident{Test.mk}}{Test.mk}}}\nwindexuse{\nwixident{Assert.eq}}{Assert.eq}{NW1mR7zJ-4QAFgj-1}\nwindexuse{\nwixident{Classical.tab}}{Classical.tab}{NW1mR7zJ-4QAFgj-1}\nwindexuse{\nwixident{Formula.mk{\_}or}}{Formula.mk:unor}{NW1mR7zJ-4QAFgj-1}\nwindexuse{\nwixident{Formula.mk{\_}pred}}{Formula.mk:unpred}{NW1mR7zJ-4QAFgj-1}\nwindexuse{\nwixident{Sort.IND}}{Sort.IND}{NW1mR7zJ-4QAFgj-1}\nwindexuse{\nwixident{Test.mk}}{Test.mk}{NW1mR7zJ-4QAFgj-1}\nwendcode{}\nwbegindocs{44}\nwdocspar

\begin{node}
I am curious whether the claim ``$\forall x\ldotp x=0\lor\exists
y,z\ldotp x=(y,z)$'' is provable in $FS_{0}$. The relevant axioms
conjoined together as the premises for the claim:
\begin{multline*}
(\forall x_{1},x_{2}\ldotp 0\neq(x_{1},x_{2}))\land
(\forall x_{3},x_{4}\ldotp P_{1}(x_{3},x_{4})=x_{3})\land
(\forall x_{4},x_{6}\ldotp P_{2}(x_{5},x_{6})=x_{6})\land\\
(0\in U)\land(\forall x_{7},x_{8}\ldotp x_{7}\in U\land x_{8}\in U\implies (x_{7},x_{8})\in U)\land\\
(\forall X\ldotp 0\in X\land\forall x_{9},x_{10}\ldotp(x_{9}\in X\land x_{10}\in X\implies(x_{9},x_{10})\in X)\implies\forall x_{11}\ldotp x_{11}\in X)\\
(\forall t_{1},t_{2}\ldotp t_{1}=t_{2}\implies t_{2}=t_{1})\land(\forall t_{3},t_{4},t_{5}\ldotp t_{3}=t_{4}\land t_{4}=t_{5}\implies t_{3}=t_{5})\\
(\forall t_{6},t_{7},f\ldotp t_{6}=t_{7}\implies ft_{6}=ft_{7})\implies\\
  \implies(\forall x_{12}\ldotp 0=x_{12}\lor\exists x_{13},x_{14}\ldotp x_{12}=(x_{13},x_{14}))
\end{multline*}
The Skolemized version of this formula is
\begin{align*}
(\forall x_{1},x_{2}\ldotp 0\neq(x_{1},x_{2}))\land\\
(\forall x_{3},x_{4}\ldotp P_{1}(x_{3},x_{4})=x_{3})\land\\
(\forall x_{4},x_{6}\ldotp P_{2}(x_{5},x_{6})=x_{6})\land\\
(0\in U)\land\\
  (\forall x_{7},x_{8}\ldotp x_{7}\notin U\lor x_{8}\notin U\lor(x_{7},x_{8})\in U)\land\\
(0\notin X\lor\bigl(f_{x}(X)\in X\land f_{y}(X)\in X\land(f_{x}(X),f_{y}(X))\notin X\bigr)\lor(\forall x\ldotp x\in X))\land\\
  (\forall t_{1},t_{2}\ldotp t_{1}\neq t_{2}\lor t_{2}=t_{1})\land\\
  (\forall t_{3},t_{4},t_{5}\ldotp t_{3}\neq t_{4}\lor t_{4}\neq t_{5}\lor t_{3}=t_{5})\land\\
(\forall t_{6},t_{7},f\ldotp t_{6}\neq t_{7}\lor ft_{6}=ft_{7})\land\\
  \land(0\neq c_{x_{12}}\land\forall x_{13},x_{14}\ldotp c_{x_{12}}\neq(x_{13},x_{14}))
\end{align*}



%% And(
%%   And(Forall(Var(x1, IND), Forall(Var(x2, IND), Not(Pred(=[Fun(cons, IND, [Var(x1, IND), Var(x2, IND)]), Fun(0, IND, [])])))),
%%   And(Forall(Var(x3, IND), Forall(Var(x4, IND), Pred(=[Fun(apply, IND, [Fun(P1, FUN, []), Fun(cons, IND, [Var(x3, IND), Var(x4, IND)])]), Var(x3, IND)]))),
%%   And(Forall(Var(x5, IND), Forall(Var(x6, IND), Pred(=[Fun(apply, IND, [Fun(P2, FUN, []), Fun(cons, IND, [Var(x5, IND), Var(x6, IND)])]), Var(x6, IND)]))),
%%   And(Pred(in[Fun(0, IND, []), Fun(U, CLASS, [])]),
%%   And(Forall(Var(x7, IND), Forall(Var(x8, IND), Or(Or(Not(Pred(in[Var(x7, IND), Fun(U, CLASS, [])])), Not(Pred(in[Var(x8, IND), Fun(U, CLASS, [])]))), Pred(in[Fun(cons, IND, [Var(x7, IND), Var(x8, IND)]), Fun(U, CLASS, [])])))),
%%   And(Forall(Var(X, CLASS), Or(Or(Not(Pred(in[Fun(0, IND, []), Var(X, CLASS)])), And(And(Pred(in[Fun(f_x, IND, [Var(X, CLASS)]), Var(X, CLASS)]), Pred(in[Fun(f_y, IND, [Var(X, CLASS)]), Var(X, CLASS)])), Not(Pred(in[Fun(cons, IND, [Fun(f_x, IND, [Var(X, CLASS)]), Fun(f_y, IND, [Var(X, CLASS)])]), Var(X, CLASS)])))), Forall(Var(x, IND), Pred(in[Var(x, IND), Var(X, CLASS)])))),
%%   And(Forall(Var(x9, IND), Forall(Var(x10, IND), Or(Not(Pred(=[Var(x9, IND), Var(x10, IND)])), Pred(=[Var(x10, IND), Var(x9, IND)])))),
%%   And(Forall(Var(x12, IND), Forall(Var(x13, IND), Forall(Var(x11, IND), Or(Or(Not(Pred(=[Var(x11, IND), Var(x12, IND)])), Not(Pred(=[Var(x12, IND), Var(x13, IND)]))), Pred(=[Var(x11, IND), Var(x13, IND)]))))),
%%       Forall(Var(x15, IND), Forall(Var(f, FUN), Forall(Var(x14, IND), Or(Not(Pred(=[Var(x14, IND), Var(x15, IND)])), Pred(=[Fun(apply, IND, [Var(f, FUN), Var(x14, IND)]), Fun(apply, IND, [Var(f, FUN), Var(x15, IND)])]))))))))))))), 
%% And(Not(Pred(=[Fun(0, IND, []), Fun(c_x16, IND, [])])),
%%     Forall(Var(x17, IND), Forall(Var(x18, IND), Not(Pred(=[Fun(c_x16, IND, []), Fun(cons, IND, [Var(x17, IND), Var(x18, IND)])]))))))
\end{node}
\nwenddocs{}\nwfilename{nw/printer.nw}\nwbegindocs{0}% -*- mode: poly-noweb; noweb-code-mode: sml-mode; -*-
%%
%% printer.nw
%% 
%% Made by Alex Nelson <pqnelson@gmail.com>
%% Login   <alex@lisp>
%% 
%% Started on  2026-04-18T12:17:53-0700
%% Last update 2026-04-18T12:17:53-0700
%% 

\section{Pretty printer}

\begin{node}
Eventually, it will get rather tedious to look at the raw serialization of
terms and formulas (and theorems and goals and\dots). We want to
support the user's ability to \emph{prettyprint} these things, which
is to say: we want to create a string representation of these things
closer to what might be found in a textbook. A pretty printer is a
collection of four functions to produce ``pretty'' string
representations for terms, formulas, theorems, and goals (given as the
collection of subgoals).
\end{node}

\nwenddocs{}\nwbegincode{1}\sublabel{NW24q7Pv-3x49D5-1}\nwmargintag{{\nwtagstyle{}\subpageref{NW24q7Pv-3x49D5-1}}}\moddef{Printer.sig~{\nwtagstyle{}\subpageref{NW24q7Pv-3x49D5-1}}}\endmoddef\nwstartdeflinemarkup\nwenddeflinemarkup
signature Printer = sig
  val term : \nwlinkedidentc{Term.t}{NW1rTmOJ-4LsRWC-1} -> string;
  val formula : \nwlinkedidentc{Formula.t}{NW3w0iRO-2jnj9W-1} -> string;
  val thm : \nwlinkedidentc{Thm.t}{NW4ICsJQ-YIPBT-1} -> string;
  val goals : ((string * \nwlinkedidentc{Formula.t}{NW3w0iRO-2jnj9W-1}) list * \nwlinkedidentc{Formula.t}{NW3w0iRO-2jnj9W-1}) list -> string;
end;
\nwnotused{Printer.sig}\nwidentuses{\\{{\nwixident{Formula.t}}{Formula.t}}\\{{\nwixident{Term.t}}{Term.t}}\\{{\nwixident{Thm.t}}{Thm.t}}}\nwindexuse{\nwixident{Formula.t}}{Formula.t}{NW24q7Pv-3x49D5-1}\nwindexuse{\nwixident{Term.t}}{Term.t}{NW24q7Pv-3x49D5-1}\nwindexuse{\nwixident{Thm.t}}{Thm.t}{NW24q7Pv-3x49D5-1}\nwendcode{}\nwbegindocs{2}\nwdocspar

\begin{node}
We will have a {\Tt{}TBPrinter\nwendquote} (``\underline{t}ext\underline{b}ook
pretty printer'').
\end{node}

\nwenddocs{}\nwbegincode{3}\sublabel{NW24q7Pv-3JbCWS-1}\nwmargintag{{\nwtagstyle{}\subpageref{NW24q7Pv-3JbCWS-1}}}\moddef{TBPrinter.sml~{\nwtagstyle{}\subpageref{NW24q7Pv-3JbCWS-1}}}\endmoddef\nwstartdeflinemarkup\nwenddeflinemarkup
structure TBPrinter : Printer = struct
\LA{}TBPrinter for a term~{\nwtagstyle{}\subpageref{NW24q7Pv-3IZAOd-1}}\RA{}

\LA{}TBPrinter for a formula~{\nwtagstyle{}\subpageref{NW24q7Pv-14xJt9-1}}\RA{}

\LA{}TBPrinter for a theorem~{\nwtagstyle{}\subpageref{NW24q7Pv-2pSwbS-1}}\RA{}

\LA{}TBPrinter for goals~{\nwtagstyle{}\subpageref{NW24q7Pv-3yfJXJ-1}}\RA{}
end;
\nwnotused{TBPrinter.sml}\nwendcode{}\nwbegindocs{4}\nwdocspar

\begin{node}
Printing terms is fairly straightforward: variables are just
themselves, and functions are of the form $f(t_{1},\dots,t_{n})$.
Constants (``nullary functions'') are written as just $f$ and not $f()$.
\end{node}

\nwenddocs{}\nwbegincode{5}\sublabel{NW24q7Pv-3IZAOd-1}\nwmargintag{{\nwtagstyle{}\subpageref{NW24q7Pv-3IZAOd-1}}}\moddef{TBPrinter for a term~{\nwtagstyle{}\subpageref{NW24q7Pv-3IZAOd-1}}}\endmoddef\nwstartdeflinemarkup\nwusesondefline{\\{NW24q7Pv-3JbCWS-1}}\nwenddeflinemarkup
fun term (Term.Var(x,_)) = x
  | term (Term.Fun(f, _, args)) =
      if List.null args
      then f
      else f^"("^(String.concatWith ", "
                                    (map term args))^
           ")";
\nwindexdefn{\nwixident{TBPrinter.term}}{TBPrinter.term}{NW24q7Pv-3IZAOd-1}\eatline
\nwused{\\{NW24q7Pv-3JbCWS-1}}\nwidentdefs{\\{{\nwixident{TBPrinter.term}}{TBPrinter.term}}}\nwendcode{}\nwbegindocs{6}\nwdocspar
\begin{node}
Formulas are a bit tricky. We need to print parentheses when
precedence causes ambiguities. For example $A\implies B\implies C$
is read as $A\implies(B\implies C)$ and not $(A\implies B)\implies C$.
Similarly $A\lor B\land C$ means $A\lor(B\land C)$.

We wish to be careful: for implications and logical equivalence, if
either subformula is another implication (or logical equivalence),
then we include the parentheses for clarity.
\end{node}

\nwenddocs{}\nwbegincode{7}\sublabel{NW24q7Pv-14xJt9-1}\nwmargintag{{\nwtagstyle{}\subpageref{NW24q7Pv-14xJt9-1}}}\moddef{TBPrinter for a formula~{\nwtagstyle{}\subpageref{NW24q7Pv-14xJt9-1}}}\endmoddef\nwstartdeflinemarkup\nwusesondefline{\\{NW24q7Pv-3JbCWS-1}}\nwenddeflinemarkup
(* "∀" (* U+2200 *)
   "∃" (* U+2203 *)
"∧" U+2227
"∨"; (* U+2228 *)
"⇔" U+21D4
 *)
\LA{}Gathering bound variables for pretty printing~{\nwtagstyle{}\subpageref{NW24q7Pv-2b9skq-1}}\RA{}
fun formula fm =
  if \nwlinkedidentc{Formula.is_falsum}{NW3w0iRO-3NyJLv-1} fm
  then "\\u00E2\\u008A\\u00A5" (* U+22A5 *)
  else if \nwlinkedidentc{Formula.is_pred}{NW3w0iRO-3NyJLv-1} fm
  then let val (P,args) = \nwlinkedidentc{Formula.dest_pred}{NW3w0iRO-1Yt1Jl-1} fm
       in P^"["^(String.concatWith ", " (map term args))^"]" end
  else if \nwlinkedidentc{Formula.is_not}{NW3w0iRO-3NyJLv-2} fm
  then let val A = \nwlinkedidentc{Formula.dest_not}{NW3w0iRO-1Yt1Jl-2} fm;
           val conn = "\\u00C2\\u00AC"; (* U+00AC *)
       in conn ^ (paren_form A fm) end
  else if \nwlinkedidentc{Formula.is_and}{NW3w0iRO-3NyJLv-3} fm
  then let val (A, B) = \nwlinkedidentc{Formula.dest_and}{NW3w0iRO-1Yt1Jl-3} fm;
           val conn = "\\u00e2\\u0088\\u00a7"; (* U+2227 *)
       in (paren_form A fm) ^ conn ^ (paren_form B fm) end
  else if \nwlinkedidentc{Formula.is_or}{NW3w0iRO-3NyJLv-3} fm
  then let val (A, B) = \nwlinkedidentc{Formula.dest_or}{NW3w0iRO-1Yt1Jl-3} fm;
           val conn = "\\u00e2\\u0088\\u00a8"; (* U+2228 *)
       in (paren_form A fm) ^ conn ^ (paren_form B fm) end
  else if \nwlinkedidentc{Formula.is_imp}{NW3w0iRO-3NyJLv-3} fm
  then let val (A, B) = \nwlinkedidentc{Formula.dest_imp}{NW3w0iRO-1Yt1Jl-3} fm;
           val conn = "\\u00e2\\u009f\\u00b9"; (* U+27F9 for long version, U+21D2 *)
       in (paren_form A fm) ^ conn ^ (paren_form B fm) end
  else if \nwlinkedidentc{Formula.is_iff}{NW3w0iRO-3NyJLv-3} fm
  then let val (A, B) = \nwlinkedidentc{Formula.dest_iff}{NW3w0iRO-1Yt1Jl-3} fm;
           val conn = "\\u00e2\\u009f\\u00ba"; (* U+27FA for long version, U+21D4 *)
       in (paren_form A fm) ^ conn ^ (paren_form B fm) end
  else if Formula.is_forall fm
  then let
    val (var, body) = Formula.dest_forall fm
    val (prefix, A) = gather_bvars Formula.is_forall
                                   Formula.dest_forall
                                   ("\\u00e2\\u0088\\u0080" ^ (* U+2200 *)
                                    (term var))
                                   body
  in prefix ^ "." ^ (formula body) end
  else if Formula.is_exists fm
  then let
    val (var, body) = Formula.dest_exists fm
    val (prefix, A) = gather_bvars Formula.is_exists
                                   Formula.dest_exists
                                   ("\\u00E2\\u0088\\u0083" ^ (* U+2203 *)
                                    (term var))
                                   body
  in prefix ^ "." ^ (formula body) end
  else raise Fail("unknown connective: "^(\nwlinkedidentc{Formula.serialize}{NW3w0iRO-3H9tUA-1} fm))
and paren_form sub fm =
    (case \nwlinkedidentc{Formula.precedence}{NW3w0iRO-3yI6te-1} sub fm of
      LESS => "("^(formula sub)^")"
    | EQUAL => if (\nwlinkedidentc{Formula.is_imp}{NW3w0iRO-3NyJLv-3} fm orelse \nwlinkedidentc{Formula.is_iff}{NW3w0iRO-3NyJLv-3} fm) andalso
                  (\nwlinkedidentc{Formula.is_imp}{NW3w0iRO-3NyJLv-3} sub orelse \nwlinkedidentc{Formula.is_iff}{NW3w0iRO-3NyJLv-3} sub)
               then "("^(formula sub)^")"
               else formula sub
    | _ => formula sub);
\nwindexdefn{\nwixident{TBPrinter.formula}}{TBPrinter.formula}{NW24q7Pv-14xJt9-1}\eatline
\nwused{\\{NW24q7Pv-3JbCWS-1}}\nwidentdefs{\\{{\nwixident{TBPrinter.formula}}{TBPrinter.formula}}}\nwidentuses{\\{{\nwixident{Formula.dest{\_}and}}{Formula.dest:unand}}\\{{\nwixident{Formula.dest{\_}iff}}{Formula.dest:uniff}}\\{{\nwixident{Formula.dest{\_}imp}}{Formula.dest:unimp}}\\{{\nwixident{Formula.dest{\_}not}}{Formula.dest:unnot}}\\{{\nwixident{Formula.dest{\_}or}}{Formula.dest:unor}}\\{{\nwixident{Formula.dest{\_}pred}}{Formula.dest:unpred}}\\{{\nwixident{Formula.is{\_}and}}{Formula.is:unand}}\\{{\nwixident{Formula.is{\_}falsum}}{Formula.is:unfalsum}}\\{{\nwixident{Formula.is{\_}iff}}{Formula.is:uniff}}\\{{\nwixident{Formula.is{\_}imp}}{Formula.is:unimp}}\\{{\nwixident{Formula.is{\_}not}}{Formula.is:unnot}}\\{{\nwixident{Formula.is{\_}or}}{Formula.is:unor}}\\{{\nwixident{Formula.is{\_}pred}}{Formula.is:unpred}}\\{{\nwixident{Formula.precedence}}{Formula.precedence}}\\{{\nwixident{Formula.serialize}}{Formula.serialize}}}\nwindexuse{\nwixident{Formula.dest{\_}and}}{Formula.dest:unand}{NW24q7Pv-14xJt9-1}\nwindexuse{\nwixident{Formula.dest{\_}iff}}{Formula.dest:uniff}{NW24q7Pv-14xJt9-1}\nwindexuse{\nwixident{Formula.dest{\_}imp}}{Formula.dest:unimp}{NW24q7Pv-14xJt9-1}\nwindexuse{\nwixident{Formula.dest{\_}not}}{Formula.dest:unnot}{NW24q7Pv-14xJt9-1}\nwindexuse{\nwixident{Formula.dest{\_}or}}{Formula.dest:unor}{NW24q7Pv-14xJt9-1}\nwindexuse{\nwixident{Formula.dest{\_}pred}}{Formula.dest:unpred}{NW24q7Pv-14xJt9-1}\nwindexuse{\nwixident{Formula.is{\_}and}}{Formula.is:unand}{NW24q7Pv-14xJt9-1}\nwindexuse{\nwixident{Formula.is{\_}falsum}}{Formula.is:unfalsum}{NW24q7Pv-14xJt9-1}\nwindexuse{\nwixident{Formula.is{\_}iff}}{Formula.is:uniff}{NW24q7Pv-14xJt9-1}\nwindexuse{\nwixident{Formula.is{\_}imp}}{Formula.is:unimp}{NW24q7Pv-14xJt9-1}\nwindexuse{\nwixident{Formula.is{\_}not}}{Formula.is:unnot}{NW24q7Pv-14xJt9-1}\nwindexuse{\nwixident{Formula.is{\_}or}}{Formula.is:unor}{NW24q7Pv-14xJt9-1}\nwindexuse{\nwixident{Formula.is{\_}pred}}{Formula.is:unpred}{NW24q7Pv-14xJt9-1}\nwindexuse{\nwixident{Formula.precedence}}{Formula.precedence}{NW24q7Pv-14xJt9-1}\nwindexuse{\nwixident{Formula.serialize}}{Formula.serialize}{NW24q7Pv-14xJt9-1}\nwendcode{}\nwbegindocs{8}\nwdocspar
\begin{node}
We have our pretty printer gather comma-separated variables when we
have the same binder nested multiple times. For example, $\forall
x\ldotp\forall y\ldotp\forall z\ldotp P[x,y,z]$ would be pretty
printed as $\forall x,y,z\ldotp P[x,y,z]$. To support this, we have a
helper function to accumulate the bound variables, separating them by
commas. We assume the printer has encountered the first quantifier,
and has made its bound variable the initial {\Tt{}acc\nwendquote} value. The body
for the quantifier is the initial {\Tt{}fm\nwendquote}.
\end{node}

\nwenddocs{}\nwbegincode{9}\sublabel{NW24q7Pv-2b9skq-1}\nwmargintag{{\nwtagstyle{}\subpageref{NW24q7Pv-2b9skq-1}}}\moddef{Gathering bound variables for pretty printing~{\nwtagstyle{}\subpageref{NW24q7Pv-2b9skq-1}}}\endmoddef\nwstartdeflinemarkup\nwusesondefline{\\{NW24q7Pv-14xJt9-1}}\nwenddeflinemarkup
fun gather_bvars is_q dest_q acc fm =
  if not(is_q fm) then (acc, fm)
  else let val (var, body) = dest_q fm
       in gather_bvars is_q dest_q (acc^", "^(term var)) body
       end;
\nwused{\\{NW24q7Pv-14xJt9-1}}\nwendcode{}\nwbegindocs{10}\nwdocspar

\begin{node}
Theorems are just a bunch of formulas separated by a turnstile.
\end{node}

\nwenddocs{}\nwbegincode{11}\sublabel{NW24q7Pv-2pSwbS-1}\nwmargintag{{\nwtagstyle{}\subpageref{NW24q7Pv-2pSwbS-1}}}\moddef{TBPrinter for a theorem~{\nwtagstyle{}\subpageref{NW24q7Pv-2pSwbS-1}}}\endmoddef\nwstartdeflinemarkup\nwusesondefline{\\{NW24q7Pv-3JbCWS-1}}\nwenddeflinemarkup
(* NB: "⊦" U+22A6 is assertion
 "⊢" U+22A2 is just right tack *)
fun thm th =
  let
    val hyps = \nwlinkedidentc{Thm.hyps}{NW4ICsJQ-YIPBT-1} th;
    val concl = \nwlinkedidentc{Thm.concl}{NW4ICsJQ-YIPBT-1} th;
    val turnstile = if List.null hyps
                    then "\\u00E2\\u008A\\u00A2 "
                    else " \\u00E2\\u008A\\u00A2 ";
  in
    (String.concatWith ", " (map formula hyps)) ^
    turnstile ^
    (formula concl)
  end;
\nwindexdefn{\nwixident{TBPrinter.thm}}{TBPrinter.thm}{NW24q7Pv-2pSwbS-1}\eatline
\nwused{\\{NW24q7Pv-3JbCWS-1}}\nwidentdefs{\\{{\nwixident{TBPrinter.thm}}{TBPrinter.thm}}}\nwidentuses{\\{{\nwixident{Thm.concl}}{Thm.concl}}\\{{\nwixident{Thm.hyps}}{Thm.hyps}}}\nwindexuse{\nwixident{Thm.concl}}{Thm.concl}{NW24q7Pv-2pSwbS-1}\nwindexuse{\nwixident{Thm.hyps}}{Thm.hyps}{NW24q7Pv-2pSwbS-1}\nwendcode{}\nwbegindocs{12}\nwdocspar
\begin{node}
A goal is similar to a theorem, but the hypotheses may be labeled.
\end{node}

\nwenddocs{}\nwbegincode{13}\sublabel{NW24q7Pv-3yfJXJ-1}\nwmargintag{{\nwtagstyle{}\subpageref{NW24q7Pv-3yfJXJ-1}}}\moddef{TBPrinter for goals~{\nwtagstyle{}\subpageref{NW24q7Pv-3yfJXJ-1}}}\endmoddef\nwstartdeflinemarkup\nwusesondefline{\\{NW24q7Pv-3JbCWS-1}}\nwenddeflinemarkup
local
  fun turnstile [] = "\\u00E2\\u008A\\u00A2 "
    | turnstile _ = " \\u00E2\\u008A\\u00A2 ";
  fun single (asl, fm) =
    (String.concatWith
       ", "
       (map (fn (label, p) =>
                if "" = label then formula p
                else "\\""^label^"\\": "^(formula p))
            asl)) ^
    (turnstile asl) ^ 
    (formula fm);
  fun iter n acc [] = acc
    | iter n acc ((asl, fm)::gls) =
        iter (n + 1)
             (acc ^
              "Subgoal "^
              (Int.toString n)^
              ": " ^
              (single (asl, fm)) ^
              (if List.null gls then "" else "\\n"))
             gls;
in
fun goals [] = "No more goals"
  | goals [g] = single g
  | goals gls = iter 1 "" gls
end;
\nwindexdefn{\nwixident{TBPrinter.goals}}{TBPrinter.goals}{NW24q7Pv-3yfJXJ-1}\eatline

\nwused{\\{NW24q7Pv-3JbCWS-1}}\nwidentdefs{\\{{\nwixident{TBPrinter.goals}}{TBPrinter.goals}}}\nwendcode{}\nwfilename{nw/xunit.nw}\nwbegindocs{0}% -*- mode: poly-noweb; noweb-code-mode: sml-mode; -*-
%%
%% xunit.nw
%% 
%% Made by Alex Nelson <pqnelson@gmail.com>
%% Login   <alex@lisp>
%% 
%% Started on  2026-05-04T16:56:51-0700
%% Last update 2026-05-04T16:56:51-0700
%% 

\section{Unit testing framework}

\begin{node}
How do we know our code works as expected? One solution is to
\emph{prove} the code works according to some specification and some
formal semantics of the programming language. This is hard,
time-consuming, and brittle.

Another solution is to test code against a handful of
examples. This is easier, more robust, and acts as a form of
documentation (want to know how to use the code? Look at some
examples, in the form of tests). Further, testing gives us immediate
feedback about whether our expected use case works or
not. (Admittedly, in the static mode of presentation, this dynamic
feedback loop is missing.)
\end{node}

\begin{node}[Overview of design]
We define each \define{Test Case} as a function which checks one
scenario for a ``system under test'' (usually a specific
function). The typical test cases are either (1) simple ``check the
actual output matches the expected output'' tests, or (2) check
specific exceptions are raised in illegal situations.

A generic test cases has four ``phases'' or ``steps'':
\begin{enumerate}
\item \textsc{Setup} the test fixture [all the data needed to describe
  the ``before executing the system under test''];
\item \textsc{Exercise} the system under test (i.e., run the code you
  want to test against the test fixture);
\item \textsc{Verify} the ``actual output'' matches [in some
  appropriate sense] the ``expected output''
\item \textsc{Teardown} the test fixtures.
\end{enumerate}
\SML/ written in a functional style handles the teardown for us
automatically, sparing us one step of montony.

The xUnit framework uses a \define{Assertion Method} to verify the
actual output matches the expected output. 

A collection of tests form a \define{Test Suite}, which is itself a
test. In particular, a test suite may contain a nested test
suite. Every test is either a test case or a test suite.

The \define{Test Runner} takes tests, executes them, constructing
\define{Test Results} containing the outcome of each test along with
other metadata (like the amount of time it took to perform). We keep
things as simple as possible, and manually register tests with the
test runner.

We then have a \define{Test Reporter} take test results produced by
the runner, then prints a summary to the screen (or otherwise produces
some kind of artifact, like an XML file). Some people advocate saving
an XML file for each time you run unit tests, other people do not. We
are agnostic, but want to recognize this is a source of change (so we
abstract it away into its own module).
\end{node}

\begin{node}[Assertions]
We have a simple module to handle assertions for equality, or
asserting that a certain condition is satisfied. User-specific
messages are supplied for failure reporting.

When the assertion is not satisfied, a custom {\Tt{}\nwlinkedidentq{Assert.Fail}{NWsR2q7-4gQMdZ-1}\nwendquote}
exception is propagated. Later, the test runner will handle these
exceptions differently than any other: assertion failures are recorded
as ``failed tests'' whereas all other exceptional situations are
recorded as ``unexpected exception raised'' as its result.
\end{node}

\nwenddocs{}\nwbegincode{1}\sublabel{NWsR2q7-4gQMdZ-1}\nwmargintag{{\nwtagstyle{}\subpageref{NWsR2q7-4gQMdZ-1}}}\moddef{Assert.sml~{\nwtagstyle{}\subpageref{NWsR2q7-4gQMdZ-1}}}\endmoddef\nwstartdeflinemarkup\nwenddeflinemarkup
structure Assert = struct
exception Fail of string;

(* that : bool -> string -> unit *)
fun that cond msg =
  if cond then ()
  else raise Fail msg;

(* eq : ''a -> ''a -> string -> unit *)
fun eq expected actual msg =
  that (expected = actual) msg;
end;
\nwindexdefn{\nwixident{Assert.Fail}}{Assert.Fail}{NWsR2q7-4gQMdZ-1}\nwindexdefn{\nwixident{Assert.that}}{Assert.that}{NWsR2q7-4gQMdZ-1}\nwindexdefn{\nwixident{Assert.eq}}{Assert.eq}{NWsR2q7-4gQMdZ-1}\eatline
\nwnotused{Assert.sml}\nwidentdefs{\\{{\nwixident{Assert.eq}}{Assert.eq}}\\{{\nwixident{Assert.Fail}}{Assert.Fail}}\\{{\nwixident{Assert.that}}{Assert.that}}}\nwendcode{}\nwbegindocs{2}\nwdocspar
\begin{node}[Test]
We begin with the data structure for a test case and a test suite.

A test case has a name and a ``test method'' which the runner
executes. We construct a test case with {\Tt{}\nwlinkedidentq{Test.mk}{NWsR2q7-3fh2dK-1}\nwendquote} (since {\Tt{}mk\nwendquote} is
the standard constructor prefix/``name component'').

A test suite has a name and a collection of tests. We construct a test
suite with the {\Tt{}\nwlinkedidentq{Test.suite}{NWsR2q7-3fh2dK-1}\nwendquote} method.

Executing a test suite or test case is handled with the {\Tt{}\nwlinkedidentq{Test.run}{NWsR2q7-1bpV8D-1}\nwendquote}
method. It will produce a {\Tt{}\nwlinkedidentq{Test.Result.t}{NWsR2q7-oQiQW-1}\nwendquote} object storing the outcome
and metadata.
\end{node}

\nwenddocs{}\nwbegincode{3}\sublabel{NWsR2q7-1Fzsdc-1}\nwmargintag{{\nwtagstyle{}\subpageref{NWsR2q7-1Fzsdc-1}}}\moddef{Test.sig~{\nwtagstyle{}\subpageref{NWsR2q7-1Fzsdc-1}}}\endmoddef\nwstartdeflinemarkup\nwenddeflinemarkup
signature Test = sig
  type t;
  structure Result : sig
              \LA{}Specification for test result~{\nwtagstyle{}\subpageref{NWsR2q7-1ETFk0-1}}\RA{}
            end;
  val suite : string -> t list -> t;
  val mk : string -> (unit -> unit) -> t;
  val run : t -> Result.t;
end;
\nwnotused{Test.sig}\nwendcode{}\nwbegindocs{4}\nwdocspar

\nwenddocs{}\nwbegincode{5}\sublabel{NWsR2q7-3ucTCC-1}\nwmargintag{{\nwtagstyle{}\subpageref{NWsR2q7-3ucTCC-1}}}\moddef{Test.sml~{\nwtagstyle{}\subpageref{NWsR2q7-3ucTCC-1}}}\endmoddef\nwstartdeflinemarkup\nwenddeflinemarkup
structure Test :> Test = struct

structure Result = struct
\LA{}Implementation of test result~{\nwtagstyle{}\subpageref{NWsR2q7-oQiQW-1}}\RA{}
end;

\LA{}Constructors for tests~{\nwtagstyle{}\subpageref{NWsR2q7-3fh2dK-1}}\RA{}

\LA{}Run a test case or suite~{\nwtagstyle{}\subpageref{NWsR2q7-1bpV8D-1}}\RA{}
end;
\nwnotused{Test.sml}\nwendcode{}\nwbegindocs{6}\nwdocspar

\begin{node}
The constructors for a test case and test suite are rather straightforward.
\end{node}

\nwenddocs{}\nwbegincode{7}\sublabel{NWsR2q7-3fh2dK-1}\nwmargintag{{\nwtagstyle{}\subpageref{NWsR2q7-3fh2dK-1}}}\moddef{Constructors for tests~{\nwtagstyle{}\subpageref{NWsR2q7-3fh2dK-1}}}\endmoddef\nwstartdeflinemarkup\nwusesondefline{\\{NWsR2q7-3ucTCC-1}}\nwenddeflinemarkup
datatype t = Case of string * (unit -> unit)
           | Suite of string * (t list);

fun mk name method = Case(name, method);

fun suite name tests = Suite(name, tests);
\nwindexdefn{\nwixident{Test.t}}{Test.t}{NWsR2q7-3fh2dK-1}\nwindexdefn{\nwixident{Test.mk}}{Test.mk}{NWsR2q7-3fh2dK-1}\nwindexdefn{\nwixident{Test.suite}}{Test.suite}{NWsR2q7-3fh2dK-1}\eatline
\nwused{\\{NWsR2q7-3ucTCC-1}}\nwidentdefs{\\{{\nwixident{Test.mk}}{Test.mk}}\\{{\nwixident{Test.suite}}{Test.suite}}\\{{\nwixident{Test.t}}{Test.t}}}\nwendcode{}\nwbegindocs{8}\nwdocspar
\begin{node}
We abstract away the outcome of the test result, since more may be
added later. Some people want to allow ``skipping'' tests and wish to
add ``skipped'' as an outcome. The outcomes we will report are:
Success (the test's outcome was as expected), ``failure'' (the test's
outcome was not as expected, so we need the failure message), and
``error'' (an unexpected exception was raised, which we store in the
outcome).

The test result for a test case includes the test name, the amount of
time elapsed while running the test as measured against a ``clock'',
and the test outcome.

The results for a test suite are similar, but instead of a single test
outcome we have a collection of test results (for all the tests
included in the suite).
\end{node}

\nwenddocs{}\nwbegincode{9}\sublabel{NWsR2q7-1ETFk0-1}\nwmargintag{{\nwtagstyle{}\subpageref{NWsR2q7-1ETFk0-1}}}\moddef{Specification for test result~{\nwtagstyle{}\subpageref{NWsR2q7-1ETFk0-1}}}\endmoddef\nwstartdeflinemarkup\nwusesondefline{\\{NWsR2q7-1Fzsdc-1}}\nwprevnextdefs{\relax}{NWsR2q7-1ETFk0-2}\nwenddeflinemarkup
type t;
\nwalsodefined{\\{NWsR2q7-1ETFk0-2}\\{NWsR2q7-1ETFk0-3}\\{NWsR2q7-1ETFk0-4}\\{NWsR2q7-1ETFk0-5}\\{NWsR2q7-1ETFk0-6}\\{NWsR2q7-1ETFk0-7}\\{NWsR2q7-1ETFk0-8}\\{NWsR2q7-1ETFk0-9}}\nwused{\\{NWsR2q7-1Fzsdc-1}}\nwendcode{}\nwbegindocs{10}\nwdocspar

\nwenddocs{}\nwbegincode{11}\sublabel{NWsR2q7-oQiQW-1}\nwmargintag{{\nwtagstyle{}\subpageref{NWsR2q7-oQiQW-1}}}\moddef{Implementation of test result~{\nwtagstyle{}\subpageref{NWsR2q7-oQiQW-1}}}\endmoddef\nwstartdeflinemarkup\nwusesondefline{\\{NWsR2q7-3ucTCC-1}}\nwprevnextdefs{\relax}{NWsR2q7-oQiQW-2}\nwenddeflinemarkup
datatype Outcome = Success
                 | Failure of string
                 | Error of exn;

datatype t = Case of string * Time.time * Outcome
           | Suite of string * Time.time * (t list);
\nwindexdefn{\nwixident{Test.Result.t}}{Test.Result.t}{NWsR2q7-oQiQW-1}\eatline
\nwalsodefined{\\{NWsR2q7-oQiQW-2}\\{NWsR2q7-oQiQW-3}\\{NWsR2q7-oQiQW-4}\\{NWsR2q7-oQiQW-5}\\{NWsR2q7-oQiQW-6}\\{NWsR2q7-oQiQW-7}\\{NWsR2q7-oQiQW-8}\\{NWsR2q7-oQiQW-9}}\nwused{\\{NWsR2q7-3ucTCC-1}}\nwidentdefs{\\{{\nwixident{Test.Result.t}}{Test.Result.t}}}\nwendcode{}\nwbegindocs{12}\nwdocspar
\begin{node}
How do we actually ``construct'' a test result from a test? Well, we
run it. For a test case, this amounts to ``clocking'' the test and
recording the outcome. We'll delegate this work to {\Tt{}\nwlinkedidentq{Test.Result.for{\_}case}{NWsR2q7-oQiQW-2}\nwendquote}.

For a test suite, we have a bunch of tests to run. We'll just lazily
run all these tests.
\end{node}

\nwenddocs{}\nwbegincode{13}\sublabel{NWsR2q7-1bpV8D-1}\nwmargintag{{\nwtagstyle{}\subpageref{NWsR2q7-1bpV8D-1}}}\moddef{Run a test case or suite~{\nwtagstyle{}\subpageref{NWsR2q7-1bpV8D-1}}}\endmoddef\nwstartdeflinemarkup\nwusesondefline{\\{NWsR2q7-3ucTCC-1}}\nwenddeflinemarkup
(* run : \nwlinkedidentc{Test.t}{NWsR2q7-3fh2dK-1} -> \nwlinkedidentc{Test.Result.t}{NWsR2q7-oQiQW-1} *)
fun run (Case (name, method)) = Result.for_case name method
  | run (Suite (name, tests)) =
      Result.for_suite name (fn () => map run tests);
\nwindexdefn{\nwixident{Test.run}}{Test.run}{NWsR2q7-1bpV8D-1}\eatline
\nwused{\\{NWsR2q7-3ucTCC-1}}\nwidentdefs{\\{{\nwixident{Test.run}}{Test.run}}}\nwidentuses{\\{{\nwixident{Test.Result.t}}{Test.Result.t}}\\{{\nwixident{Test.t}}{Test.t}}}\nwindexuse{\nwixident{Test.Result.t}}{Test.Result.t}{NWsR2q7-1bpV8D-1}\nwindexuse{\nwixident{Test.t}}{Test.t}{NWsR2q7-1bpV8D-1}\nwendcode{}\nwbegindocs{14}\nwdocspar
\begin{node}
Given the name (as a string) and the four-phase test procedure as a
function, we:
\begin{enumerate}
\item Record the start time
\item Execute the test
\item In the default case: record the time interval since we started
  the test, along with the name, and that the test outcome was a {\Tt{}Success\nwendquote}.
\item If an {\Tt{}\nwlinkedidentq{Assert.Fail}{NWsR2q7-4gQMdZ-1}\ msg\nwendquote} was raised, then we record the time
  interval since we started, along with the name, and that the outcome
  was a {\Tt{}Failure\nwendquote} with the given {\Tt{}msg\nwendquote}.
\item If any other exception {\Tt{}e\nwendquote} was raised, then we record the time
  interval since we started, along with the test case's name, and that
  the outcome was an {\Tt{}Error\ e\nwendquote}.
\end{enumerate}
\end{node}

\nwenddocs{}\nwbegincode{15}\sublabel{NWsR2q7-1ETFk0-2}\nwmargintag{{\nwtagstyle{}\subpageref{NWsR2q7-1ETFk0-2}}}\moddef{Specification for test result~{\nwtagstyle{}\subpageref{NWsR2q7-1ETFk0-1}}}\plusendmoddef\nwstartdeflinemarkup\nwusesondefline{\\{NWsR2q7-1Fzsdc-1}}\nwprevnextdefs{NWsR2q7-1ETFk0-1}{NWsR2q7-1ETFk0-3}\nwenddeflinemarkup
val for_case : string -> (unit -> unit) -> t;
\nwused{\\{NWsR2q7-1Fzsdc-1}}\nwendcode{}\nwbegindocs{16}\nwdocspar

\nwenddocs{}\nwbegincode{17}\sublabel{NWsR2q7-oQiQW-2}\nwmargintag{{\nwtagstyle{}\subpageref{NWsR2q7-oQiQW-2}}}\moddef{Implementation of test result~{\nwtagstyle{}\subpageref{NWsR2q7-oQiQW-1}}}\plusendmoddef\nwstartdeflinemarkup\nwusesondefline{\\{NWsR2q7-3ucTCC-1}}\nwprevnextdefs{NWsR2q7-oQiQW-1}{NWsR2q7-oQiQW-3}\nwenddeflinemarkup
fun for_case name assertion =
  let
    val start = Time.now();
  in
    (assertion();
     Case(name,
          (Time.-)(Time.now(), start),
          Success))
    handle \nwlinkedidentc{Assert.Fail}{NWsR2q7-4gQMdZ-1} msg =>
           Case (name,
                 (Time.-)(Time.now(), start),
                 Failure msg)
         | e => Case(name,
                     (Time.-)(Time.now(), start),
                     Error e)
  end;
\nwindexdefn{\nwixident{Test.Result.for{\_}case}}{Test.Result.for:uncase}{NWsR2q7-oQiQW-2}\eatline
\nwused{\\{NWsR2q7-3ucTCC-1}}\nwidentdefs{\\{{\nwixident{Test.Result.for{\_}case}}{Test.Result.for:uncase}}}\nwidentuses{\\{{\nwixident{Assert.Fail}}{Assert.Fail}}}\nwindexuse{\nwixident{Assert.Fail}}{Assert.Fail}{NWsR2q7-oQiQW-2}\nwendcode{}\nwbegindocs{18}\nwdocspar
\begin{node}
The test result for a test suite is determined by:
\begin{enumerate}
\item Clock the start time; then
\item Calculate the result of each test by recursively invoking
  {\Tt{}\nwlinkedidentq{Test.run}{NWsR2q7-1bpV8D-1}\nwendquote} on each test in the suite; then
\item Clock the end time, recording the duration of running the tests
  in the suite; then
\item Finally construct the test result object for the test suite.
\end{enumerate}
\end{node}

\nwenddocs{}\nwbegincode{19}\sublabel{NWsR2q7-1ETFk0-3}\nwmargintag{{\nwtagstyle{}\subpageref{NWsR2q7-1ETFk0-3}}}\moddef{Specification for test result~{\nwtagstyle{}\subpageref{NWsR2q7-1ETFk0-1}}}\plusendmoddef\nwstartdeflinemarkup\nwusesondefline{\\{NWsR2q7-1Fzsdc-1}}\nwprevnextdefs{NWsR2q7-1ETFk0-2}{NWsR2q7-1ETFk0-4}\nwenddeflinemarkup
val for_suite : string -> (unit -> (t list)) -> t;
\nwused{\\{NWsR2q7-1Fzsdc-1}}\nwendcode{}\nwbegindocs{20}\nwdocspar

\nwenddocs{}\nwbegincode{21}\sublabel{NWsR2q7-oQiQW-3}\nwmargintag{{\nwtagstyle{}\subpageref{NWsR2q7-oQiQW-3}}}\moddef{Implementation of test result~{\nwtagstyle{}\subpageref{NWsR2q7-oQiQW-1}}}\plusendmoddef\nwstartdeflinemarkup\nwusesondefline{\\{NWsR2q7-3ucTCC-1}}\nwprevnextdefs{NWsR2q7-oQiQW-2}{NWsR2q7-oQiQW-4}\nwenddeflinemarkup
fun for_suite name mk_results =
  let
    val start = Time.now();
    val results = mk_results ();
    val dt = (Time.-)(Time.now(), start);
  in
    Suite(name, dt, results)
  end;
\nwindexdefn{\nwixident{Test.Result.for{\_}suite}}{Test.Result.for:unsuite}{NWsR2q7-oQiQW-3}\eatline
\nwused{\\{NWsR2q7-3ucTCC-1}}\nwidentdefs{\\{{\nwixident{Test.Result.for{\_}suite}}{Test.Result.for:unsuite}}}\nwendcode{}\nwbegindocs{22}\nwdocspar
\begin{node}[Accessor methods for test result objects]
We have abstracted the test result internals away, so we won't have to
worry about them. We have standard ways to get the name of the test,
the failure message (if any---otherwise return the empty string), and
the exception message (if any---otherwise return the empty string).
\end{node}

\nwenddocs{}\nwbegincode{23}\sublabel{NWsR2q7-1ETFk0-4}\nwmargintag{{\nwtagstyle{}\subpageref{NWsR2q7-1ETFk0-4}}}\moddef{Specification for test result~{\nwtagstyle{}\subpageref{NWsR2q7-1ETFk0-1}}}\plusendmoddef\nwstartdeflinemarkup\nwusesondefline{\\{NWsR2q7-1Fzsdc-1}}\nwprevnextdefs{NWsR2q7-1ETFk0-3}{NWsR2q7-1ETFk0-5}\nwenddeflinemarkup
val name : t -> string;
val msg : t -> string;
val exn_message : t -> string;
\nwused{\\{NWsR2q7-1Fzsdc-1}}\nwendcode{}\nwbegindocs{24}\nwdocspar

\nwenddocs{}\nwbegincode{25}\sublabel{NWsR2q7-oQiQW-4}\nwmargintag{{\nwtagstyle{}\subpageref{NWsR2q7-oQiQW-4}}}\moddef{Implementation of test result~{\nwtagstyle{}\subpageref{NWsR2q7-oQiQW-1}}}\plusendmoddef\nwstartdeflinemarkup\nwusesondefline{\\{NWsR2q7-3ucTCC-1}}\nwprevnextdefs{NWsR2q7-oQiQW-3}{NWsR2q7-oQiQW-5}\nwenddeflinemarkup
(* name : \nwlinkedidentc{Test.Result.t}{NWsR2q7-oQiQW-1} -> string *)
fun name (Case (n,_,_)) = n
  | name (Suite (n,_,_)) = n;

(* msg : \nwlinkedidentc{Test.Result.t}{NWsR2q7-oQiQW-1} -> string *)
fun msg (Case (_,_,Failure msg)) = msg
  | msg _ = "";

(* exn_msg : \nwlinkedidentc{Test.Result.t}{NWsR2q7-oQiQW-1} -> string *)
fun exn_message (Case (_,_,Error e)) = exnMessage e
  | exn_message _ = ""; 
\nwindexdefn{\nwixident{Test.Result.name}}{Test.Result.name}{NWsR2q7-oQiQW-4}\nwindexdefn{\nwixident{Test.Result.msg}}{Test.Result.msg}{NWsR2q7-oQiQW-4}\nwindexdefn{\nwixident{Test.Result.exn{\_}message}}{Test.Result.exn:unmessage}{NWsR2q7-oQiQW-4}\eatline
\nwused{\\{NWsR2q7-3ucTCC-1}}\nwidentdefs{\\{{\nwixident{Test.Result.exn{\_}message}}{Test.Result.exn:unmessage}}\\{{\nwixident{Test.Result.msg}}{Test.Result.msg}}\\{{\nwixident{Test.Result.name}}{Test.Result.name}}}\nwidentuses{\\{{\nwixident{Test.Result.t}}{Test.Result.t}}}\nwindexuse{\nwixident{Test.Result.t}}{Test.Result.t}{NWsR2q7-oQiQW-4}\nwendcode{}\nwbegindocs{26}\nwdocspar
\begin{node}
For test cases, we can access the time it took to execute.

For test suites, there is some overhead by recursively invoking
{\Tt{}\nwlinkedidentq{Test.run}{NWsR2q7-1bpV8D-1}\nwendquote} lazily. If the user wants to omit this, they can
calculate the sum of time intervals it took to run each test case.
\end{node}

\nwenddocs{}\nwbegincode{27}\sublabel{NWsR2q7-1ETFk0-5}\nwmargintag{{\nwtagstyle{}\subpageref{NWsR2q7-1ETFk0-5}}}\moddef{Specification for test result~{\nwtagstyle{}\subpageref{NWsR2q7-1ETFk0-1}}}\plusendmoddef\nwstartdeflinemarkup\nwusesondefline{\\{NWsR2q7-1Fzsdc-1}}\nwprevnextdefs{NWsR2q7-1ETFk0-4}{NWsR2q7-1ETFk0-6}\nwenddeflinemarkup
(* realtime excludes the time it took a TestSuite to
   allocate infrastructure in memory. *)
val realtime : t -> Time.time;
val runtime : t -> Time.time;
\nwused{\\{NWsR2q7-1Fzsdc-1}}\nwendcode{}\nwbegindocs{28}\nwdocspar

\nwenddocs{}\nwbegincode{29}\sublabel{NWsR2q7-oQiQW-5}\nwmargintag{{\nwtagstyle{}\subpageref{NWsR2q7-oQiQW-5}}}\moddef{Implementation of test result~{\nwtagstyle{}\subpageref{NWsR2q7-oQiQW-1}}}\plusendmoddef\nwstartdeflinemarkup\nwusesondefline{\\{NWsR2q7-3ucTCC-1}}\nwprevnextdefs{NWsR2q7-oQiQW-4}{NWsR2q7-oQiQW-6}\nwenddeflinemarkup
(* runtime : t -> Time.time

How long did it take to run the test(s) and construct the
result(s)? *)
  fun runtime (Case (_,dt,_)) = dt
    | runtime (Suite (_,dt,_)) = dt;

  (* realtime : t -> Time.time

How long did it take just to run the test(s)? *)
  fun realtime (Case (_,dt,_)) = dt
    | realtime (Suite (_,_,[])) = Time.zeroTime
    | realtime (Suite (_,_,[x])) = realtime x
    | realtime (Suite (_,_,r::rs)) =
      foldl (fn (result,dt) =>
                (Time.+)(dt, realtime result))
            (realtime r)
            rs;
\nwindexdefn{\nwixident{Test.Result.runtime}}{Test.Result.runtime}{NWsR2q7-oQiQW-5}\nwindexdefn{\nwixident{Test.Result.realtime}}{Test.Result.realtime}{NWsR2q7-oQiQW-5}\eatline
\nwused{\\{NWsR2q7-3ucTCC-1}}\nwidentdefs{\\{{\nwixident{Test.Result.realtime}}{Test.Result.realtime}}\\{{\nwixident{Test.Result.runtime}}{Test.Result.runtime}}}\nwendcode{}\nwbegindocs{30}\nwdocspar
\begin{node}
When reporting test results for a test suite, we will need to access
the results for each test in the suite.
\end{node}

\nwenddocs{}\nwbegincode{31}\sublabel{NWsR2q7-1ETFk0-6}\nwmargintag{{\nwtagstyle{}\subpageref{NWsR2q7-1ETFk0-6}}}\moddef{Specification for test result~{\nwtagstyle{}\subpageref{NWsR2q7-1ETFk0-1}}}\plusendmoddef\nwstartdeflinemarkup\nwusesondefline{\\{NWsR2q7-1Fzsdc-1}}\nwprevnextdefs{NWsR2q7-1ETFk0-5}{NWsR2q7-1ETFk0-7}\nwenddeflinemarkup
val results : t -> t list;
\nwused{\\{NWsR2q7-1Fzsdc-1}}\nwendcode{}\nwbegindocs{32}\nwdocspar

\nwenddocs{}\nwbegincode{33}\sublabel{NWsR2q7-oQiQW-6}\nwmargintag{{\nwtagstyle{}\subpageref{NWsR2q7-oQiQW-6}}}\moddef{Implementation of test result~{\nwtagstyle{}\subpageref{NWsR2q7-oQiQW-1}}}\plusendmoddef\nwstartdeflinemarkup\nwusesondefline{\\{NWsR2q7-3ucTCC-1}}\nwprevnextdefs{NWsR2q7-oQiQW-5}{NWsR2q7-oQiQW-7}\nwenddeflinemarkup
fun results (Suite (_,_,rs)) = rs;
\nwindexdefn{\nwixident{Test.Result.results}}{Test.Result.results}{NWsR2q7-oQiQW-6}\eatline
\nwused{\\{NWsR2q7-3ucTCC-1}}\nwidentdefs{\\{{\nwixident{Test.Result.results}}{Test.Result.results}}}\nwendcode{}\nwbegindocs{34}\nwdocspar
\begin{node}
We can tabulate the number of successes, failures, errors, and total
number of tests. The logic for these functions are all rather similar.
\begin{itemize}
\item For a test case, if the outcome is a {\Tt{}Success\nwendquote} then return
  $1$, otherwise return $0$.
\item For a test suite, just count all the successes in each test in
  the suite.
\end{itemize}
Counting the total number of test cases just sums the total number of
successes, failures, and errors.
\end{node}

\nwenddocs{}\nwbegincode{35}\sublabel{NWsR2q7-1ETFk0-7}\nwmargintag{{\nwtagstyle{}\subpageref{NWsR2q7-1ETFk0-7}}}\moddef{Specification for test result~{\nwtagstyle{}\subpageref{NWsR2q7-1ETFk0-1}}}\plusendmoddef\nwstartdeflinemarkup\nwusesondefline{\\{NWsR2q7-1Fzsdc-1}}\nwprevnextdefs{NWsR2q7-1ETFk0-6}{NWsR2q7-1ETFk0-8}\nwenddeflinemarkup
val count_successes : t -> int;
val count_failures : t -> int;
val count_errors : t -> int;
val count_total : t -> int;
\nwused{\\{NWsR2q7-1Fzsdc-1}}\nwendcode{}\nwbegindocs{36}\nwdocspar

\nwenddocs{}\nwbegincode{37}\sublabel{NWsR2q7-oQiQW-7}\nwmargintag{{\nwtagstyle{}\subpageref{NWsR2q7-oQiQW-7}}}\moddef{Implementation of test result~{\nwtagstyle{}\subpageref{NWsR2q7-oQiQW-1}}}\plusendmoddef\nwstartdeflinemarkup\nwusesondefline{\\{NWsR2q7-3ucTCC-1}}\nwprevnextdefs{NWsR2q7-oQiQW-6}{NWsR2q7-oQiQW-8}\nwenddeflinemarkup
fun count_successes (Case (_,_,Success)) = 1
  | count_successes (Case _) = 0 
  | count_successes (Suite (_,_,outcomes)) =
      foldl (op +) 0 (map count_successes outcomes);

fun count_failures (Case (_,_,Failure _)) = 1
  | count_failures (Case _) = 0
  | count_failures (Suite (_,_,outcomes)) =
      foldl (op +) 0 (map count_failures outcomes);

fun count_errors (Case (_,_,Error _)) = 1
  | count_errors (Case _) = 0
  | count_errors (Suite (_,_,outcomes)) =
      foldl (op +) 0 (map count_errors outcomes);

fun count_total x =
  count_successes x + count_failures x + count_errors x;
\nwindexdefn{\nwixident{Test.Result.count{\_}total}}{Test.Result.count:untotal}{NWsR2q7-oQiQW-7}\nwindexdefn{\nwixident{Test.Result.count{\_}successes}}{Test.Result.count:unsuccesses}{NWsR2q7-oQiQW-7}\nwindexdefn{\nwixident{Test.Result.count{\_}failures}}{Test.Result.count:unfailures}{NWsR2q7-oQiQW-7}\nwindexdefn{\nwixident{Test.Result.count{\_}errors}}{Test.Result.count:unerrors}{NWsR2q7-oQiQW-7}\eatline
\nwused{\\{NWsR2q7-3ucTCC-1}}\nwidentdefs{\\{{\nwixident{Test.Result.count{\_}errors}}{Test.Result.count:unerrors}}\\{{\nwixident{Test.Result.count{\_}failures}}{Test.Result.count:unfailures}}\\{{\nwixident{Test.Result.count{\_}successes}}{Test.Result.count:unsuccesses}}\\{{\nwixident{Test.Result.count{\_}total}}{Test.Result.count:untotal}}}\nwendcode{}\nwbegindocs{38}\nwdocspar
\begin{node}[Predicate for test outcome]
We can say a test suite or a test case is a success if the {\Tt{}count{\_}successes\nwendquote} equals the {\Tt{}count{\_}total\nwendquote}.

Failure for a test case is rather obvious (the outcome was a
{\Tt{}Failure\nwendquote}). But what qualifies as a ``failure'' for a test suite? We
will treat \emph{any} reported failure as indicating the test suite
should qualify as a ``failure''.

Similarly, any error result in a test contained in a test suite
qualifies that test suite as a ``error''.
\end{node}

\nwenddocs{}\nwbegincode{39}\sublabel{NWsR2q7-1ETFk0-8}\nwmargintag{{\nwtagstyle{}\subpageref{NWsR2q7-1ETFk0-8}}}\moddef{Specification for test result~{\nwtagstyle{}\subpageref{NWsR2q7-1ETFk0-1}}}\plusendmoddef\nwstartdeflinemarkup\nwusesondefline{\\{NWsR2q7-1Fzsdc-1}}\nwprevnextdefs{NWsR2q7-1ETFk0-7}{NWsR2q7-1ETFk0-9}\nwenddeflinemarkup
val is_success : t -> bool;
val is_failure : t -> bool;
val is_error : t -> bool;
\nwused{\\{NWsR2q7-1Fzsdc-1}}\nwendcode{}\nwbegindocs{40}\nwdocspar

\nwenddocs{}\nwbegincode{41}\sublabel{NWsR2q7-oQiQW-8}\nwmargintag{{\nwtagstyle{}\subpageref{NWsR2q7-oQiQW-8}}}\moddef{Implementation of test result~{\nwtagstyle{}\subpageref{NWsR2q7-oQiQW-1}}}\plusendmoddef\nwstartdeflinemarkup\nwusesondefline{\\{NWsR2q7-3ucTCC-1}}\nwprevnextdefs{NWsR2q7-oQiQW-7}{NWsR2q7-oQiQW-9}\nwenddeflinemarkup
fun is_success x = (count_total x) = (count_successes x);

fun is_failure (Case (_,_, Failure _)) = true
  | is_failure (Case _) = false
  | is_failure (s as Suite _) = (count_failures s) > 0;

fun is_error (Case (_,_, Error _)) = true
  | is_error (Case _) = false
  | is_error (s as Suite _) = (count_errors s) > 0;
\nwindexdefn{\nwixident{Test.Result.is{\_}success}}{Test.Result.is:unsuccess}{NWsR2q7-oQiQW-8}\nwindexdefn{\nwixident{Test.Result.is{\_}failure}}{Test.Result.is:unfailure}{NWsR2q7-oQiQW-8}\nwindexdefn{\nwixident{Test.Result.is{\_}error}}{Test.Result.is:unerror}{NWsR2q7-oQiQW-8}\eatline
\nwused{\\{NWsR2q7-3ucTCC-1}}\nwidentdefs{\\{{\nwixident{Test.Result.is{\_}error}}{Test.Result.is:unerror}}\\{{\nwixident{Test.Result.is{\_}failure}}{Test.Result.is:unfailure}}\\{{\nwixident{Test.Result.is{\_}success}}{Test.Result.is:unsuccess}}}\nwendcode{}\nwbegindocs{42}\nwdocspar
\begin{node}[Predicates for test results]
Since the test result is an opaque data structure, we need to provide
predicates to indicate whether a test result is for a suite or a test
case. This will be necessary for reporting results.
\end{node}

\nwenddocs{}\nwbegincode{43}\sublabel{NWsR2q7-1ETFk0-9}\nwmargintag{{\nwtagstyle{}\subpageref{NWsR2q7-1ETFk0-9}}}\moddef{Specification for test result~{\nwtagstyle{}\subpageref{NWsR2q7-1ETFk0-1}}}\plusendmoddef\nwstartdeflinemarkup\nwusesondefline{\\{NWsR2q7-1Fzsdc-1}}\nwprevnextdefs{NWsR2q7-1ETFk0-8}{\relax}\nwenddeflinemarkup
val is_case : t -> bool;
val is_suite : t -> bool;
\nwused{\\{NWsR2q7-1Fzsdc-1}}\nwendcode{}\nwbegindocs{44}\nwdocspar

\nwenddocs{}\nwbegincode{45}\sublabel{NWsR2q7-oQiQW-9}\nwmargintag{{\nwtagstyle{}\subpageref{NWsR2q7-oQiQW-9}}}\moddef{Implementation of test result~{\nwtagstyle{}\subpageref{NWsR2q7-oQiQW-1}}}\plusendmoddef\nwstartdeflinemarkup\nwusesondefline{\\{NWsR2q7-3ucTCC-1}}\nwprevnextdefs{NWsR2q7-oQiQW-8}{\relax}\nwenddeflinemarkup
fun is_case (Case _) = true
  | is_case _ = false;

fun is_suite (Suite _) = true
  | is_suite _ = false;
\nwindexdefn{\nwixident{Test.Result.is{\_}case}}{Test.Result.is:uncase}{NWsR2q7-oQiQW-9}\nwindexdefn{\nwixident{Test.Result.is{\_}suite}}{Test.Result.is:unsuite}{NWsR2q7-oQiQW-9}\eatline
\nwused{\\{NWsR2q7-3ucTCC-1}}\nwidentdefs{\\{{\nwixident{Test.Result.is{\_}case}}{Test.Result.is:uncase}}\\{{\nwixident{Test.Result.is{\_}suite}}{Test.Result.is:unsuite}}}\nwendcode{}\nwbegindocs{46}\nwdocspar
\begin{node}[Test reporter specification]
So far, we have implemented the assertion methods, the test case and
suite structures, and the test result data structure. We just need a
test runner and test reporter modules. Let us now turn to the test reporter.

We will abstract away ``reporting'' as any \SML/ function producing a
unit type (which is what happens when we write to a file or print to
the screen). Ostensibly, the user may want to send an email (or post a
TikTok video, or page someone, or\dots), so we just abstract away
whatever ``reporting'' mechanism this could be.
\end{node}

\nwenddocs{}\nwbegincode{47}\sublabel{NWsR2q7-jZez2-1}\nwmargintag{{\nwtagstyle{}\subpageref{NWsR2q7-jZez2-1}}}\moddef{Reporter.sig~{\nwtagstyle{}\subpageref{NWsR2q7-jZez2-1}}}\endmoddef\nwstartdeflinemarkup\nwenddeflinemarkup
signature Reporter = sig
  val report : \nwlinkedidentc{Test.Result.t}{NWsR2q7-oQiQW-1} -> unit;
  val report_all : \nwlinkedidentc{Test.Result.t}{NWsR2q7-oQiQW-1} list -> unit;
end;
\nwnotused{Reporter.sig}\nwidentuses{\\{{\nwixident{Test.Result.t}}{Test.Result.t}}}\nwindexuse{\nwixident{Test.Result.t}}{Test.Result.t}{NWsR2q7-jZez2-1}\nwendcode{}\nwbegindocs{48}\nwdocspar

\begin{node}
I'm used the JUnit's report to the terminal, so I'm going to stick to
that. We will report the time interval in microseconds.
\end{node}

\nwenddocs{}\nwbegincode{49}\sublabel{NWsR2q7-bqErk-1}\nwmargintag{{\nwtagstyle{}\subpageref{NWsR2q7-bqErk-1}}}\moddef{JUnitTt.sml~{\nwtagstyle{}\subpageref{NWsR2q7-bqErk-1}}}\endmoddef\nwstartdeflinemarkup\nwenddeflinemarkup
(* 
Print to the terminal (hence the "-Tt" suffix) a summary of test
results in the style of JUnit.

If any failures or errors occur, print those out with a bit more
detail.
*)
structure JUnitTt :> Reporter = struct
structure Result = Test.Result;

fun interval_to_string (dt : Time.time) =
  (LargeInt.toString (Time.toMicroseconds dt))^"ms";

\LA{}Iteratively report JUnit-style report for a single outcome~{\nwtagstyle{}\subpageref{NWsR2q7-4wVoS-1}}\RA{}

val report = print o (report_iter "");

val report_all = app report;
end;
\nwnotused{JUnitTt.sml}\nwendcode{}\nwbegindocs{50}\nwdocspar

\begin{node}
How we report results depends on whether we are reporting for a test
case or a test suite:
\begin{itemize}
\item Reporting a test case:
\begin{enumerate}
\item If the test failed, we should print the path to the test, announce it failed, and give the failure message.
\item If the test err'd, we should do everything similarly to the failure case, except announce it was an ERROR and produce the exception message.
\item If the test succeeded, then don't do anything: we do not wish to pollute the terminal.
\end{enumerate}
\item For a test suite:
\begin{enumerate}
\item Write the test suite's full path to the screen.
\item Write the results for each test in the suite to the screen.
\item Write to the screen a summary line counting the total number of tests run, the number of failures, the number of errors, and the name of the test suite.
\end{enumerate}
\end{itemize}
\end{node}

\nwenddocs{}\nwbegincode{51}\sublabel{NWsR2q7-4wVoS-1}\nwmargintag{{\nwtagstyle{}\subpageref{NWsR2q7-4wVoS-1}}}\moddef{Iteratively report JUnit-style report for a single outcome~{\nwtagstyle{}\subpageref{NWsR2q7-4wVoS-1}}}\endmoddef\nwstartdeflinemarkup\nwusesondefline{\\{NWsR2q7-bqErk-1}}\nwenddeflinemarkup
fun report_iter p =
  let
    fun path "" s = s
      | path p s = p^"."^s; 
    fun for_case c =
      if Result.is_failure c
      then concat [path p (Result.name c),
                   " FAIL: ",
                   Result.msg c,
                   "\\n"]
      else if Result.is_error c
      then concat [path p (Result.name c),
                   " ERROR: ",
                   (Result.exn_message c),
                   "\\n"]
      else "";
    fun for_suite result =
      let
        val p2 = path p (Result.name result);
        val rs = \nwlinkedidentc{Test.Result.results}{NWsR2q7-oQiQW-6} result;
      in
        concat ["Running ", p2, "\\n",
                concat (map (report_iter p2) rs),
                "Tests run: ",
                Int.toString(Result.count_total result),
                ", Failures: ",
                Int.toString(Result.count_failures result),
                ", Errors: ",
                Int.toString(Result.count_errors result),
                ", Time elapsed: ",
                interval_to_string (Result.realtime result),
                " - in ",
                (Result.name result), "\\n"]
      end;
  in
    fn r => if \nwlinkedidentc{Test.Result.is_case}{NWsR2q7-oQiQW-9} r then for_case r
            else for_suite r
  end;
\nwused{\\{NWsR2q7-bqErk-1}}\nwidentuses{\\{{\nwixident{Test.Result.is{\_}case}}{Test.Result.is:uncase}}\\{{\nwixident{Test.Result.results}}{Test.Result.results}}}\nwindexuse{\nwixident{Test.Result.is{\_}case}}{Test.Result.is:uncase}{NWsR2q7-4wVoS-1}\nwindexuse{\nwixident{Test.Result.results}}{Test.Result.results}{NWsR2q7-4wVoS-1}\nwendcode{}\nwbegindocs{52}\nwdocspar

\begin{node}[Test runner]
Now we have arrived at the final stage of the unit testing framework:
we want to take a list of tests, execute them, and report the results
using some test reporter. We will simply create a ``parametrized module''
(\SML/ functor) which allows for any arbitrary test reporter to be used.

The test runner has a single method, which expects a list of tests to
run and report.
\end{node}

\nwenddocs{}\nwbegincode{53}\sublabel{NWsR2q7-4Peuwc-1}\nwmargintag{{\nwtagstyle{}\subpageref{NWsR2q7-4Peuwc-1}}}\moddef{TestRunner.sig~{\nwtagstyle{}\subpageref{NWsR2q7-4Peuwc-1}}}\endmoddef\nwstartdeflinemarkup\nwenddeflinemarkup
signature TestRunner = sig
  val run : \nwlinkedidentc{Test.t}{NWsR2q7-3fh2dK-1} list -> unit;
end;
\nwnotused{TestRunner.sig}\nwidentuses{\\{{\nwixident{Test.t}}{Test.t}}}\nwindexuse{\nwixident{Test.t}}{Test.t}{NWsR2q7-4Peuwc-1}\nwendcode{}\nwbegindocs{54}\nwdocspar

\begin{node}
The test runner:
\begin{enumerate}
\item Runs the tests to obtain the test results,
\item Reports the results,
\end{enumerate}
and that's it.
\end{node}

\nwenddocs{}\nwbegincode{55}\sublabel{NWsR2q7-2dKvEw-1}\nwmargintag{{\nwtagstyle{}\subpageref{NWsR2q7-2dKvEw-1}}}\moddef{tests/MkRunner.sml~{\nwtagstyle{}\subpageref{NWsR2q7-2dKvEw-1}}}\endmoddef\nwstartdeflinemarkup\nwenddeflinemarkup
(*
Using any test reporter, we can run the tests, then
report the results.
*)
functor MkRunner(Reporter : Reporter) :> TestRunner = struct
  fun run tests =
    let
      val results = map \nwlinkedidentc{Test.run}{NWsR2q7-1bpV8D-1} tests;
    in
      Reporter.report_all results;
      ()
    end;
end;
\nwnotused{tests/MkRunner.sml}\nwidentuses{\\{{\nwixident{Test.run}}{Test.run}}}\nwindexuse{\nwixident{Test.run}}{Test.run}{NWsR2q7-2dKvEw-1}\nwendcode{}\nwbegindocs{56}\nwdocspar

\begin{node}
The ``main'' function is rather straightforward for unit tests.

Working with MLton and Poly/ML, we also define the {\Tt{}main\nwendquote} function here
for future usage.

We also initialize the unit tests to be run with an empty test
suite. This way, when we register test suites, we will do it in a
uniform manner (to aid readability).
\end{node}

\nwenddocs{}\nwbegincode{57}\sublabel{NWsR2q7-3MeQDa-1}\nwmargintag{{\nwtagstyle{}\subpageref{NWsR2q7-3MeQDa-1}}}\moddef{tests/main.sml~{\nwtagstyle{}\subpageref{NWsR2q7-3MeQDa-1}}}\endmoddef\nwstartdeflinemarkup\nwenddeflinemarkup
(*
Using the different reporters, we can run the tests, then
report the results.

Defaults to the JUnitTt reporter.

For inspecting per test result, VerboseTt.report may be useful
*)
structure Runner :> TestRunner = MkRunner(JUnitTt);

val tests : \nwlinkedidentc{Test.t}{NWsR2q7-3fh2dK-1} list = [
  \nwlinkedidentc{Test.suite}{NWsR2q7-3fh2dK-1} "" []
\LA{}Register unit tests~{\nwtagstyle{}\subpageref{NW1mR7zJ-2ohDA9-1}}\RA{}
];

fun main () = Runner.run tests;
\nwnotused{tests/main.sml}\nwidentuses{\\{{\nwixident{Test.suite}}{Test.suite}}\\{{\nwixident{Test.t}}{Test.t}}}\nwindexuse{\nwixident{Test.suite}}{Test.suite}{NWsR2q7-3MeQDa-1}\nwindexuse{\nwixident{Test.t}}{Test.t}{NWsR2q7-3MeQDa-1}\nwendcode{}\nwbegindocs{58}\nwdocspar

\begin{node}[SML/NJ adapter]
If we want to support SML/NJ, then we should organize test suites as
structures implementing the {\Tt{}TestSuite\nwendquote} signature, then add to
the {\Tt{}\LA{}Unit tests~{\nwtagstyle{}\subpageref{nw@notdef}}\RA{}\nwendquote} chunk the {\Tt{}MyTestSuite.suite\nwendquote}. We informally
stipulate that the {\Tt{}suite\nwendquote} value should be a test suite, but there's
no way to enforce this.
\end{node}

\nwenddocs{}\nwbegincode{59}\sublabel{NWsR2q7-1mlFrS-1}\nwmargintag{{\nwtagstyle{}\subpageref{NWsR2q7-1mlFrS-1}}}\moddef{tests/TestSuite.sig~{\nwtagstyle{}\subpageref{NWsR2q7-1mlFrS-1}}}\endmoddef\nwstartdeflinemarkup\nwenddeflinemarkup
signature TestSuite = sig
  val suite : \nwlinkedidentc{Test.t}{NWsR2q7-3fh2dK-1}
end;
\nwnotused{tests/TestSuite.sig}\nwidentuses{\\{{\nwixident{Test.t}}{Test.t}}}\nwindexuse{\nwixident{Test.t}}{Test.t}{NWsR2q7-1mlFrS-1}\nwendcode{}\nwbegindocs{60}\nwdocspar

\begin{node}[Build script for MLton]
Lastly, it will be convenient to have a build script for MLton (and
Poly/ML if the user employs \texttt{polymlb}).
\end{node}

\nwenddocs{}\nwbegincode{61}\sublabel{NWsR2q7-2v996w-1}\nwmargintag{{\nwtagstyle{}\subpageref{NWsR2q7-2v996w-1}}}\moddef{test.mlb~{\nwtagstyle{}\subpageref{NWsR2q7-2v996w-1}}}\endmoddef\nwstartdeflinemarkup\nwenddeflinemarkup
$(SML_LIB)/basis/basis.mlb
tests/Assert.sml
tests/Test.sig
tests/Test.sml
tests/Reporter.sig
tests/JUnitTt.sml
tests/TestRunner.sig
tests/MkRunner.sml
tests/TestSuite.sig
common.mlb
\LA{}Unit tests to run~{\nwtagstyle{}\subpageref{NW1mR7zJ-3lyaJm-1}}\RA{}
tests/main.sml
src/main.sml
\nwnotused{test.mlb}\nwendcode{}\nwbegindocs{62}\nwdocspar

\begin{node}[References]
Beck~\cite{beck1994simple} coded the progenitor SUnit framework for
treating specification documents as Smalltalk code.
Following Meszaros~\cite{meszaros2007xunit}, we describe the basic
design of a simplified testing framework for our needs.
\end{node}

\begin{xca}[Saving artifacts]
Extreme programming practices advocate test reporters saving an
``artifact'' of some kind when running the tests. Write a test
reporter producing some file recording the results.
\begin{enumerate}
\item Determine the markup for the file (XML? JSON? YAML?)
\item Write a specification for the artifact produced. 
\item Write a reporter producing this sort of artifact. Where will the artifact be placed? 
\item Write a tool which reads in the artifact, and prints out a summary to the screen. 
\end{enumerate}
\end{xca}

\begin{xca}[Combine test reporters]
Write a \SML/ functor \texttt{MkJointReporter(R1 : Reporter, R2 : Reporter) : Reporter}
which will perform all the reporting in {\Tt{}R1\nwendquote} and {\Tt{}R2\nwendquote}.
\end{xca}

\begin{xca}[Capture output]
Write code to capture \texttt{stdout} (and \texttt{stderr}), and store
anything written to these output streams as metadata in a test result.
\end{xca}

\begin{xca}[Exit status]
Modify the test runner code so it will exit with a failure status if
there were failed unit tests (or errors encountered).
\end{xca}

\begin{xca}[Testing the tester]
We mentioned that unit testing originated as a way to transform a
specification document into executable code.
\begin{enumerate}
\item Write specifications for the code in this section (i.e.,
  specifications for the unit testing framework itself)
\item Transform these specifications into unit tests, so we can check
  the unit testing code is doing what we expect.
\end{enumerate}
\end{xca}

\begin{xca}[BDD testing]
There are as many ways to test code as there are programmers (if not more).
A popular alternative to the minimal design we've implemented is
``behaviour driven development'' where you first write down the tests
in \define{Gherkin} to specify the behaviour expected in different
situations (``Given [setup], when [code is executed], expected
[outcome]'' is usually the template they use), then you run the tests
(to make sure they fail), then you implement the code to satisfy the
specified behaviour.
\begin{enumerate}
\item Research Gherkin and this approach to unit testing. What are its
  merits and drawbacks?
\item What would be needed to modify our code to implement a
  Gherkin-style testing framework?
\item Would this style of coding be well-suited to developing a proof
  assistant? Why or why not? (This could easily be an essay, a term
  paper, or a PhD thesis.)
\end{enumerate}
\end{xca}
\nwenddocs{}

\nwixlogsorted{c}{{$(\neg A)\lor B$ sufficient to prove $A\implies B$}{NW1NjfzY-3NUP7w-1}{\nwixu{NW1NjfzY-4CUt2L-6}\nwixd{NW1NjfzY-3NUP7w-1}}}%
\nwixlogsorted{c}{{$\vdash A\land(B\lor C)\iff((A\land B)\lor(A\land C))$}{NW1NjfzY-efA9w-1}{\nwixu{NW1NjfzY-4CUt2L-4}\nwixd{NW1NjfzY-efA9w-1}}}%
\nwixlogsorted{c}{{$\vdash A\lor A\implies A$}{NW1NjfzY-3wJoEd-1}{\nwixu{NW1NjfzY-4CUt2L-2}\nwixd{NW1NjfzY-3wJoEd-1}}}%
\nwixlogsorted{c}{{Apply two functions componentwise to an $FS_{0}$ pair}{NW2DqMzY-1lMRH2-1}{\nwixu{NW2DqMzY-3aj0iq-1}\nwixd{NW2DqMzY-1lMRH2-1}}}%
\nwixlogsorted{c}{{Assert.sml}{NWsR2q7-4gQMdZ-1}{\nwixd{NWsR2q7-4gQMdZ-1}}}%
\nwixlogsorted{c}{{Axioms for $FS_{0}$}{NW4ICsJQ-2aDdLu-1}{\nwixu{NW4ICsJQ-YIPBT-1}\nwixd{NW4ICsJQ-2aDdLu-1}\nwixd{NW4ICsJQ-2aDdLu-2}\nwixd{NW4ICsJQ-2aDdLu-3}\nwixd{NW4ICsJQ-2aDdLu-4}\nwixd{NW4ICsJQ-2aDdLu-5}\nwixd{NW4ICsJQ-2aDdLu-6}\nwixd{NW4ICsJQ-2aDdLu-7}\nwixd{NW4ICsJQ-2aDdLu-8}\nwixd{NW4ICsJQ-2aDdLu-9}\nwixd{NW4ICsJQ-2aDdLu-A}\nwixd{NW4ICsJQ-2aDdLu-B}\nwixd{NW4ICsJQ-2aDdLu-C}\nwixd{NW4ICsJQ-2aDdLu-D}}}%
\nwixlogsorted{c}{{Check if a formula contains a variable}{NW3w0iRO-4DOLFK-1}{\nwixu{NW3w0iRO-2jnj9W-1}\nwixd{NW3w0iRO-4DOLFK-1}}}%
\nwixlogsorted{c}{{Check theorem hypotheses form subset of assumptions list}{NW1F32vJ-1bk6hJ-1}{\nwixu{NW1F32vJ-2XH0oT-1}\nwixd{NW1F32vJ-1bk6hJ-1}}}%
\nwixlogsorted{c}{{Classical.sig}{NW1mR7zJ-3NMRe1-1}{\nwixd{NW1mR7zJ-3NMRe1-1}}}%
\nwixlogsorted{c}{{Classical.sml}{NW1mR7zJ-l5jX3-1}{\nwixd{NW1mR7zJ-l5jX3-1}}}%
\nwixlogsorted{c}{{Comparing formulas}{NW3w0iRO-35s47m-1}{\nwixu{NW3w0iRO-2jnj9W-1}\nwixd{NW3w0iRO-35s47m-1}}}%
\nwixlogsorted{c}{{Comparing two terms}{NW1rTmOJ-2yJYsC-1}{\nwixu{NW1rTmOJ-4LsRWC-1}\nwixd{NW1rTmOJ-2yJYsC-1}}}%
\nwixlogsorted{c}{{Compose two $FS_{0}$ functions together}{NW2DqMzY-1dWhhL-1}{\nwixu{NW2DqMzY-3aj0iq-1}\nwixd{NW2DqMzY-1dWhhL-1}}}%
\nwixlogsorted{c}{{Congruence of left disjunct}{NW1NjfzY-4VDMqe-1}{\nwixu{NW1NjfzY-4CUt2L-2}\nwixd{NW1NjfzY-4VDMqe-1}}}%
\nwixlogsorted{c}{{Congruence of right disjunct}{NW1NjfzY-lrWvY-1}{\nwixu{NW1NjfzY-4CUt2L-2}\nwixd{NW1NjfzY-lrWvY-1}}}%
\nwixlogsorted{c}{{Conjunction is associative}{NW1NjfzY-328uZt-1}{\nwixu{NW1NjfzY-4CUt2L-3}\nwixd{NW1NjfzY-328uZt-1}}}%
\nwixlogsorted{c}{{Conjunction is commutative}{NW1NjfzY-29ie6N-1}{\nwixu{NW1NjfzY-4CUt2L-3}\nwixd{NW1NjfzY-29ie6N-1}}}%
\nwixlogsorted{c}{{Conjunction left congruence}{NW1NjfzY-4YXcw4-1}{\nwixu{NW1NjfzY-4CUt2L-3}\nwixd{NW1NjfzY-4YXcw4-1}}}%
\nwixlogsorted{c}{{Conjunction right congruence}{NW1NjfzY-TrI6l-1}{\nwixu{NW1NjfzY-4CUt2L-3}\nwixd{NW1NjfzY-TrI6l-1}}}%
\nwixlogsorted{c}{{Conjunction tautology $\vdash A\implies(B\implies A\land B)$}{NW1NjfzY-1I5AY8-1}{\nwixu{NW1NjfzY-4CUt2L-3}\nwixd{NW1NjfzY-1I5AY8-1}}}%
\nwixlogsorted{c}{{Constant function constructor}{NW2DqMzY-22MFug-1}{\nwixu{NW2DqMzY-3aj0iq-1}\nwixd{NW2DqMzY-22MFug-1}}}%
\nwixlogsorted{c}{{Construct an ordered pair from $FS_{0}$ individuals}{NW2DqMzY-4SkDla-1}{\nwixu{NW2DqMzY-1NUSQs-1}\nwixd{NW2DqMzY-4SkDla-1}}}%
\nwixlogsorted{c}{{Constructor for inductively defined $FS_{0}$ class}{NW2DqMzY-1mMvkv-1}{\nwixu{NW2DqMzY-TSvER-1}\nwixd{NW2DqMzY-1mMvkv-1}}}%
\nwixlogsorted{c}{{Constructors for $FS_{0}$ classes}{NW2DqMzY-TSvER-1}{\nwixu{NW2DqMzY-mVmU8-1}\nwixd{NW2DqMzY-TSvER-1}}}%
\nwixlogsorted{c}{{Constructors for $FS_{0}$ functions}{NW2DqMzY-3aj0iq-1}{\nwixu{NW2DqMzY-mVmU8-1}\nwixd{NW2DqMzY-3aj0iq-1}}}%
\nwixlogsorted{c}{{Constructors for $FS_{0}$ individuals}{NW2DqMzY-1NUSQs-1}{\nwixu{NW2DqMzY-mVmU8-1}\nwixd{NW2DqMzY-1NUSQs-1}}}%
\nwixlogsorted{c}{{Constructors for $FS_{0}$ predicates}{NW2DqMzY-4OjF2N-1}{\nwixu{NW2DqMzY-mVmU8-1}\nwixd{NW2DqMzY-4OjF2N-1}}}%
\nwixlogsorted{c}{{Constructors for formulas}{NW3w0iRO-3SlqND-1}{\nwixu{NW3w0iRO-2jnj9W-1}\nwixd{NW3w0iRO-3SlqND-1}\nwixd{NW3w0iRO-3SlqND-2}\nwixd{NW3w0iRO-3SlqND-3}\nwixd{NW3w0iRO-3SlqND-4}}}%
\nwixlogsorted{c}{{Constructors for tests}{NWsR2q7-3fh2dK-1}{\nwixu{NWsR2q7-3ucTCC-1}\nwixd{NWsR2q7-3fh2dK-1}}}%
\nwixlogsorted{c}{{Contrapositive derived rules}{NW1NjfzY-4KFbFg-1}{\nwixu{NW1NjfzY-4CUt2L-5}\nwixd{NW1NjfzY-4KFbFg-1}\nwixd{NW1NjfzY-4KFbFg-2}}}%
\nwixlogsorted{c}{{Curry implication}{NW1NjfzY-1iptma-1}{\nwixu{NW1NjfzY-4CUt2L-5}\nwixd{NW1NjfzY-1iptma-1}}}%
\nwixlogsorted{c}{{Decide if two individuals are equal}{NW2DqMzY-2vOgWm-1}{\nwixu{NW2DqMzY-3aj0iq-1}\nwixd{NW2DqMzY-2vOgWm-1}}}%
\nwixlogsorted{c}{{Derived results for Intuitionistic predicate logic}{NW1NjfzY-uZfts-1}{\nwixu{NW1NjfzY-4dPl0p-1}\nwixd{NW1NjfzY-uZfts-1}\nwixd{NW1NjfzY-uZfts-2}\nwixd{NW1NjfzY-uZfts-3}\nwixd{NW1NjfzY-uZfts-4}\nwixd{NW1NjfzY-uZfts-5}\nwixd{NW1NjfzY-uZfts-6}\nwixd{NW1NjfzY-uZfts-7}\nwixd{NW1NjfzY-uZfts-8}\nwixd{NW1NjfzY-uZfts-9}\nwixd{NW1NjfzY-uZfts-A}\nwixd{NW1NjfzY-uZfts-B}\nwixd{NW1NjfzY-uZfts-C}\nwixd{NW1NjfzY-uZfts-D}\nwixd{NW1NjfzY-uZfts-E}}}%
\nwixlogsorted{c}{{Derived results for Intuitionistic propositional logic}{NW1NjfzY-4CUt2L-1}{\nwixu{NW1NjfzY-4dPl0p-1}\nwixd{NW1NjfzY-4CUt2L-1}\nwixd{NW1NjfzY-4CUt2L-2}\nwixd{NW1NjfzY-4CUt2L-3}\nwixd{NW1NjfzY-4CUt2L-4}\nwixd{NW1NjfzY-4CUt2L-5}\nwixd{NW1NjfzY-4CUt2L-6}\nwixd{NW1NjfzY-4CUt2L-7}}}%
\nwixlogsorted{c}{{Derived.sml}{NW1NjfzY-4dPl0p-1}{\nwixd{NW1NjfzY-4dPl0p-1}}}%
\nwixlogsorted{c}{{Destructors for formulas}{NW3w0iRO-1Yt1Jl-1}{\nwixu{NW3w0iRO-2jnj9W-1}\nwixd{NW3w0iRO-1Yt1Jl-1}\nwixd{NW3w0iRO-1Yt1Jl-2}\nwixd{NW3w0iRO-1Yt1Jl-3}\nwixd{NW3w0iRO-1Yt1Jl-4}}}%
\nwixlogsorted{c}{{Destructure \code{}const\edoc{}}{NW2DqMzY-1qEma8-1}{\nwixu{NW2DqMzY-22MFug-1}\nwixd{NW2DqMzY-1qEma8-1}}}%
\nwixlogsorted{c}{{Destructure \code{}induct\edoc{}}{NW2DqMzY-4ToA2q-1}{\nwixu{NW2DqMzY-1mMvkv-1}\nwixd{NW2DqMzY-4ToA2q-1}}}%
\nwixlogsorted{c}{{Destructure \code{}intersection\edoc{}}{NW2DqMzY-4EzxyM-1}{\nwixu{NW2DqMzY-19bkPi-1}\nwixd{NW2DqMzY-4EzxyM-1}}}%
\nwixlogsorted{c}{{Destructure \code{}juxt\edoc{}}{NW2DqMzY-hvApp-1}{\nwixu{NW2DqMzY-1lMRH2-1}\nwixd{NW2DqMzY-hvApp-1}}}%
\nwixlogsorted{c}{{Destructure \code{}preimage\edoc{}}{NW2DqMzY-LSOi5-1}{\nwixu{NW2DqMzY-31WKtJ-1}\nwixd{NW2DqMzY-LSOi5-1}}}%
\nwixlogsorted{c}{{Destructure \code{}recur\edoc{}}{NW2DqMzY-Ir9cb-1}{\nwixu{NW2DqMzY-2fxlqE-1}\nwixd{NW2DqMzY-Ir9cb-1}}}%
\nwixlogsorted{c}{{Destructure \code{}singleton\edoc{}}{NW2DqMzY-Qd4qS-1}{\nwixu{NW2DqMzY-14kolG-1}\nwixd{NW2DqMzY-Qd4qS-1}}}%
\nwixlogsorted{c}{{Destructure \code{}union\edoc{}}{NW2DqMzY-27dSmE-1}{\nwixu{NW2DqMzY-22W7o7-1}\nwixd{NW2DqMzY-27dSmE-1}}}%
\nwixlogsorted{c}{{Destructuring \code{}apply\edoc{}}{NW2DqMzY-1VpWgD-1}{\nwixu{NW2DqMzY-1NUSQs-1}\nwixd{NW2DqMzY-1VpWgD-1}}}%
\nwixlogsorted{c}{{Destructuring \code{}comp\edoc{}}{NW2DqMzY-22XS9g-1}{\nwixu{NW2DqMzY-1dWhhL-1}\nwixd{NW2DqMzY-22XS9g-1}}}%
\nwixlogsorted{c}{{Destructuring \code{}cons\edoc{}}{NW2DqMzY-21ZNyJ-1}{\nwixu{NW2DqMzY-1NUSQs-1}\nwixd{NW2DqMzY-21ZNyJ-1}}}%
\nwixlogsorted{c}{{Destructuring \code{}equal\edoc{}}{NW2DqMzY-2FxXJs-1}{\nwixu{NW2DqMzY-ocj48-1}\nwixd{NW2DqMzY-2FxXJs-1}}}%
\nwixlogsorted{c}{{Destructuring \code{}in\edoc{}}{NW2DqMzY-4I0O5b-1}{\nwixu{NW2DqMzY-eTij6-1}\nwixd{NW2DqMzY-4I0O5b-1}}}%
\nwixlogsorted{c}{{Determine a fresh variable given some free variables}{NW1rTmOJ-4Ykq9M-1}{\nwixu{NW1rTmOJ-4LsRWC-1}\nwixd{NW1rTmOJ-4Ykq9M-1}}}%
\nwixlogsorted{c}{{Disjunction is associative}{NW1NjfzY-J2rA5-1}{\nwixu{NW1NjfzY-4CUt2L-2}\nwixd{NW1NjfzY-J2rA5-1}}}%
\nwixlogsorted{c}{{Disjunction is commutative}{NW1NjfzY-3WoMbH-1}{\nwixu{NW1NjfzY-4CUt2L-2}\nwixd{NW1NjfzY-3WoMbH-1}}}%
\nwixlogsorted{c}{{Disjunction two-sided congruence}{NW1NjfzY-h4O6T-1}{\nwixu{NW1NjfzY-4CUt2L-4}\nwixd{NW1NjfzY-h4O6T-1}}}%
\nwixlogsorted{c}{{Equality constructor for $FS_{0}$ terms}{NW2DqMzY-ocj48-1}{\nwixu{NW2DqMzY-4OjF2N-1}\nwixd{NW2DqMzY-ocj48-1}}}%
\nwixlogsorted{c}{{Equality of two terms}{NW1rTmOJ-3810sY-1}{\nwixu{NW1rTmOJ-4LsRWC-1}\nwixd{NW1rTmOJ-3810sY-1}}}%
\nwixlogsorted{c}{{Finite map for unification}{NW1BOcHB-25AKca-1}{\nwixu{NW1BOcHB-1KP4Zf-1}\nwixd{NW1BOcHB-25AKca-1}}}%
\nwixlogsorted{c}{{Folding formulas}{NW3w0iRO-1UOPDY-1}{\nwixu{NW3w0iRO-2jnj9W-1}\nwixd{NW3w0iRO-1UOPDY-1}}}%
\nwixlogsorted{c}{{Formula.sig}{NW3w0iRO-5BH7S-1}{\nwixd{NW3w0iRO-5BH7S-1}}}%
\nwixlogsorted{c}{{Formula.sml}{NW3w0iRO-2jnj9W-1}{\nwixd{NW3w0iRO-2jnj9W-1}}}%
\nwixlogsorted{c}{{Free and bound variables in a formula}{NW3w0iRO-jRgYZ-1}{\nwixu{NW3w0iRO-2jnj9W-1}\nwixd{NW3w0iRO-jRgYZ-1}\nwixd{NW3w0iRO-jRgYZ-2}}}%
\nwixlogsorted{c}{{Free variables in a term}{NW1rTmOJ-2nRtky-1}{\nwixu{NW1rTmOJ-4LsRWC-1}\nwixd{NW1rTmOJ-2nRtky-1}}}%
\nwixlogsorted{c}{{Frege's tautology (propositional logic)}{NW1NjfzY-bcGD8-1}{\nwixu{NW1NjfzY-4CUt2L-5}\nwixd{NW1NjfzY-bcGD8-1}}}%
\nwixlogsorted{c}{{Fresh variable for a formula}{NW4ICsJQ-2auZja-1}{\nwixu{NW4ICsJQ-2aDdLu-8}\nwixd{NW4ICsJQ-2auZja-1}}}%
\nwixlogsorted{c}{{FS0.sig}{NW2DqMzY-3R3yGC-1}{\nwixd{NW2DqMzY-3R3yGC-1}}}%
\nwixlogsorted{c}{{FS0.sml}{NW2DqMzY-mVmU8-1}{\nwixd{NW2DqMzY-mVmU8-1}}}%
\nwixlogsorted{c}{{Gather all the functions used}{NW1mR7zJ-43eePE-1}{\nwixd{NW1mR7zJ-43eePE-1}\nwixu{NW1mR7zJ-4B0sXE-1}}}%
\nwixlogsorted{c}{{Gathering bound variables for pretty printing}{NW24q7Pv-2b9skq-1}{\nwixu{NW24q7Pv-14xJt9-1}\nwixd{NW24q7Pv-2b9skq-1}}}%
\nwixlogsorted{c}{{Goal.sig}{NW1F32vJ-KAFG0-1}{\nwixd{NW1F32vJ-KAFG0-1}}}%
\nwixlogsorted{c}{{Goal.sml}{NW1F32vJ-2MKBgq-1}{\nwixd{NW1F32vJ-2MKBgq-1}}}%
\nwixlogsorted{c}{{Hypothetical syllogism}{NW1NjfzY-1KT69U-1}{\nwixu{NW1NjfzY-4CUt2L-5}\nwixd{NW1NjfzY-1KT69U-1}}}%
\nwixlogsorted{c}{{Hypothetical syllogism derived inference rule}{NW1NjfzY-2CxI45-1}{\nwixu{NW1NjfzY-4CUt2L-1}\nwixd{NW1NjfzY-2CxI45-1}}}%
\nwixlogsorted{c}{{Identity function $FS_{0}$ primitive function constant}{NW2DqMzY-D9caK-1}{\nwixu{NW2DqMzY-3aj0iq-1}\nwixd{NW2DqMzY-D9caK-1}}}%
\nwixlogsorted{c}{{Iff is reflexive}{NW1NjfzY-YtLDN-1}{\nwixu{NW1NjfzY-4CUt2L-7}\nwixd{NW1NjfzY-YtLDN-1}}}%
\nwixlogsorted{c}{{Iff is symmetric}{NW1NjfzY-3hb2X5-1}{\nwixu{NW1NjfzY-4CUt2L-7}\nwixd{NW1NjfzY-3hb2X5-1}}}%
\nwixlogsorted{c}{{Iff is transitive}{NW1NjfzY-2isj9x-1}{\nwixu{NW1NjfzY-4CUt2L-7}\nwixd{NW1NjfzY-2isj9x-1}}}%
\nwixlogsorted{c}{{Implementation of tactics}{NW1F32vJ-3Cb3oS-1}{\nwixu{NW1F32vJ-2XH0oT-1}\nwixd{NW1F32vJ-3Cb3oS-1}\nwixd{NW1F32vJ-3Cb3oS-2}\nwixd{NW1F32vJ-3Cb3oS-3}\nwixd{NW1F32vJ-3Cb3oS-4}\nwixd{NW1F32vJ-3Cb3oS-5}\nwixd{NW1F32vJ-3Cb3oS-6}\nwixd{NW1F32vJ-3Cb3oS-7}\nwixd{NW1F32vJ-3Cb3oS-8}\nwixd{NW1F32vJ-3Cb3oS-9}\nwixd{NW1F32vJ-3Cb3oS-A}\nwixd{NW1F32vJ-3Cb3oS-B}}}%
\nwixlogsorted{c}{{Implementation of test result}{NWsR2q7-oQiQW-1}{\nwixu{NWsR2q7-3ucTCC-1}\nwixd{NWsR2q7-oQiQW-1}\nwixd{NWsR2q7-oQiQW-2}\nwixd{NWsR2q7-oQiQW-3}\nwixd{NWsR2q7-oQiQW-4}\nwixd{NWsR2q7-oQiQW-5}\nwixd{NWsR2q7-oQiQW-6}\nwixd{NWsR2q7-oQiQW-7}\nwixd{NWsR2q7-oQiQW-8}\nwixd{NWsR2q7-oQiQW-9}}}%
\nwixlogsorted{c}{{Implication as negated conjunction}{NW1NjfzY-2qUp6m-1}{\nwixu{NW1NjfzY-4CUt2L-6}\nwixd{NW1NjfzY-2qUp6m-1}}}%
\nwixlogsorted{c}{{Inference rules}{NW4ICsJQ-2zKkPU-1}{\nwixu{NW4ICsJQ-YIPBT-1}\nwixd{NW4ICsJQ-2zKkPU-1}\nwixd{NW4ICsJQ-2zKkPU-2}\nwixd{NW4ICsJQ-2zKkPU-3}\nwixd{NW4ICsJQ-2zKkPU-4}\nwixd{NW4ICsJQ-2zKkPU-5}\nwixd{NW4ICsJQ-2zKkPU-6}\nwixd{NW4ICsJQ-2zKkPU-7}\nwixd{NW4ICsJQ-2zKkPU-8}\nwixd{NW4ICsJQ-2zKkPU-9}\nwixd{NW4ICsJQ-2zKkPU-A}\nwixd{NW4ICsJQ-2zKkPU-B}}}%
\nwixlogsorted{c}{{Intersection of two $FS_{0}$ classes}{NW2DqMzY-19bkPi-1}{\nwixu{NW2DqMzY-TSvER-1}\nwixd{NW2DqMzY-19bkPi-1}}}%
\nwixlogsorted{c}{{Iteratively report JUnit-style report for a single outcome}{NWsR2q7-4wVoS-1}{\nwixu{NWsR2q7-bqErk-1}\nwixd{NWsR2q7-4wVoS-1}}}%
\nwixlogsorted{c}{{JUnitTt.sml}{NWsR2q7-bqErk-1}{\nwixd{NWsR2q7-bqErk-1}}}%
\nwixlogsorted{c}{{Left distributivity of disjunction over conjunction}{NW1NjfzY-2uP77E-1}{\nwixd{NW1NjfzY-2uP77E-1}}}%
\nwixlogsorted{c}{{Negation normal form for negated formulas}{NW1mR7zJ-4HwyeF-1}{\nwixu{NW1mR7zJ-4TkWUc-1}\nwixd{NW1mR7zJ-4HwyeF-1}}}%
\nwixlogsorted{c}{{Negation of conjunction}{NW1NjfzY-wYirn-1}{\nwixu{NW1NjfzY-4CUt2L-6}\nwixd{NW1NjfzY-wYirn-1}}}%
\nwixlogsorted{c}{{Negation of disjunction}{NW1NjfzY-4UPWei-1}{\nwixu{NW1NjfzY-4CUt2L-6}\nwixd{NW1NjfzY-4UPWei-1}}}%
\nwixlogsorted{c}{{Poor man's hashmap}{nw@notdef}{\nwixu{NW1BOcHB-1KP4Zf-1}}}%
\nwixlogsorted{c}{{Precedence for connectives}{NW3w0iRO-3yI6te-1}{\nwixu{NW3w0iRO-2jnj9W-1}\nwixd{NW3w0iRO-3yI6te-1}}}%
\nwixlogsorted{c}{{Predicates for formulas}{NW3w0iRO-3NyJLv-1}{\nwixu{NW3w0iRO-2jnj9W-1}\nwixd{NW3w0iRO-3NyJLv-1}\nwixd{NW3w0iRO-3NyJLv-2}\nwixd{NW3w0iRO-3NyJLv-3}\nwixd{NW3w0iRO-3NyJLv-4}}}%
\nwixlogsorted{c}{{Predicates of terms about their sorts}{NW1rTmOJ-3FfXTO-1}{\nwixu{NW1rTmOJ-4LsRWC-1}\nwixd{NW1rTmOJ-3FfXTO-1}}}%
\nwixlogsorted{c}{{Preimage of an $FS_{0}$ class under an $FS_{0}$ function}{NW2DqMzY-31WKtJ-1}{\nwixu{NW2DqMzY-TSvER-1}\nwixd{NW2DqMzY-31WKtJ-1}}}%
\nwixlogsorted{c}{{Primitive \code{}apply\edoc{} constructor for functions to individuals}{NW2DqMzY-2oGoBx-1}{\nwixu{NW2DqMzY-1NUSQs-1}\nwixd{NW2DqMzY-2oGoBx-1}}}%
\nwixlogsorted{c}{{Printer.sig}{NW24q7Pv-3x49D5-1}{\nwixd{NW24q7Pv-3x49D5-1}}}%
\nwixlogsorted{c}{{Project out the first component of an $FS_{0}$ ordered pair}{NW2DqMzY-ilITB-1}{\nwixu{NW2DqMzY-3aj0iq-1}\nwixd{NW2DqMzY-ilITB-1}}}%
\nwixlogsorted{c}{{Project out the second component of an $FS_{0}$ ordered pair}{NW2DqMzY-HQiIh-1}{\nwixu{NW2DqMzY-3aj0iq-1}\nwixd{NW2DqMzY-HQiIh-1}}}%
\nwixlogsorted{c}{{Prove $\vdash(A\implies B)\implies((A\implies\neg B)\implies\neg A)$}{NW1NjfzY-4ZKGAh-1}{\nwixu{NW1NjfzY-4CUt2L-1}\nwixd{NW1NjfzY-4ZKGAh-1}}}%
\nwixlogsorted{c}{{Proving a theorem from tactics}{NW1F32vJ-2zEBzF-1}{\nwixu{NW1F32vJ-2XH0oT-1}\nwixd{NW1F32vJ-2zEBzF-1}}}%
\nwixlogsorted{c}{{Recursor constructor for primitive recursion}{NW2DqMzY-2fxlqE-1}{\nwixu{NW2DqMzY-3aj0iq-1}\nwixd{NW2DqMzY-2fxlqE-1}}}%
\nwixlogsorted{c}{{Register unit tests}{NW1mR7zJ-2ohDA9-1}{\nwixd{NW1mR7zJ-2ohDA9-1}\nwixu{NWsR2q7-3MeQDa-1}}}%
\nwixlogsorted{c}{{Reporter.sig}{NWsR2q7-jZez2-1}{\nwixd{NWsR2q7-jZez2-1}}}%
\nwixlogsorted{c}{{Right distributivity of conjunction over disjunction}{NW1NjfzY-8du9Q-1}{\nwixu{NW1NjfzY-4CUt2L-4}\nwixd{NW1NjfzY-8du9Q-1}}}%
\nwixlogsorted{c}{{Run a test case or suite}{NWsR2q7-1bpV8D-1}{\nwixu{NWsR2q7-3ucTCC-1}\nwixd{NWsR2q7-1bpV8D-1}}}%
\nwixlogsorted{c}{{Serialize a formula}{NW3w0iRO-3H9tUA-1}{\nwixu{NW3w0iRO-2jnj9W-1}\nwixd{NW3w0iRO-3H9tUA-1}}}%
\nwixlogsorted{c}{{Serialize a term}{NW1rTmOJ-40h2U3-1}{\nwixu{NW1rTmOJ-4LsRWC-1}\nwixd{NW1rTmOJ-40h2U3-1}}}%
\nwixlogsorted{c}{{Set membership constructor for $FS_{0}$}{NW2DqMzY-eTij6-1}{\nwixu{NW2DqMzY-4OjF2N-1}\nwixd{NW2DqMzY-eTij6-1}}}%
\nwixlogsorted{c}{{Simplify conjunction}{NW1mR7zJ-19fhtN-1}{\nwixu{NW1mR7zJ-2RorhN-1}\nwixd{NW1mR7zJ-19fhtN-1}}}%
\nwixlogsorted{c}{{Simplify disjunction}{NW1mR7zJ-4ZXXl0-1}{\nwixu{NW1mR7zJ-2RorhN-1}\nwixd{NW1mR7zJ-4ZXXl0-1}}}%
\nwixlogsorted{c}{{Simplify double negation}{NW1mR7zJ-tCoNW-1}{\nwixu{NW1mR7zJ-2RorhN-1}\nwixd{NW1mR7zJ-tCoNW-1}}}%
\nwixlogsorted{c}{{Simplify formulas using classical tautologies}{NW1mR7zJ-2RorhN-1}{\nwixu{NW1mR7zJ-l5jX3-1}\nwixd{NW1mR7zJ-2RorhN-1}}}%
\nwixlogsorted{c}{{Simplify implication}{NW1mR7zJ-3HyT53-1}{\nwixu{NW1mR7zJ-2RorhN-1}\nwixd{NW1mR7zJ-3HyT53-1}}}%
\nwixlogsorted{c}{{Simplify logical equivalence}{NW1mR7zJ-48args-1}{\nwixu{NW1mR7zJ-2RorhN-1}\nwixd{NW1mR7zJ-48args-1}}}%
\nwixlogsorted{c}{{Simplify the body of a quantified formula}{NW1mR7zJ-3pvELX-1}{\nwixu{NW1mR7zJ-2RorhN-1}\nwixd{NW1mR7zJ-3pvELX-1}}}%
\nwixlogsorted{c}{{Singleton $FS_{0}$ class}{NW2DqMzY-14kolG-1}{\nwixu{NW2DqMzY-TSvER-1}\nwixd{NW2DqMzY-14kolG-1}}}%
\nwixlogsorted{c}{{Skolemization}{NW1mR7zJ-4QrsTb-1}{\nwixu{NW1mR7zJ-l5jX3-1}\nwixd{NW1mR7zJ-4QrsTb-1}}}%
\nwixlogsorted{c}{{Skolemize a formula}{NW1mR7zJ-4B0sXE-1}{\nwixd{NW1mR7zJ-4B0sXE-1}\nwixu{NW1mR7zJ-4QrsTb-1}}}%
\nwixlogsorted{c}{{Sort.sml}{NW1rTmOJ-3ZqggC-1}{\nwixd{NW1rTmOJ-3ZqggC-1}}}%
\nwixlogsorted{c}{{Specification for $FS_{0}$ axioms}{NW4ICsJQ-hYCIS-1}{\nwixu{NW4ICsJQ-2cY0B1-1}\nwixd{NW4ICsJQ-hYCIS-1}\nwixd{NW4ICsJQ-hYCIS-2}\nwixd{NW4ICsJQ-hYCIS-3}\nwixd{NW4ICsJQ-hYCIS-4}\nwixd{NW4ICsJQ-hYCIS-5}\nwixd{NW4ICsJQ-hYCIS-6}\nwixd{NW4ICsJQ-hYCIS-7}\nwixd{NW4ICsJQ-hYCIS-8}\nwixd{NW4ICsJQ-hYCIS-9}\nwixd{NW4ICsJQ-hYCIS-A}\nwixd{NW4ICsJQ-hYCIS-B}\nwixd{NW4ICsJQ-hYCIS-C}\nwixd{NW4ICsJQ-hYCIS-D}}}%
\nwixlogsorted{c}{{Specification for $FS_{0}$ classes}{NW2DqMzY-33zuOn-1}{\nwixu{NW2DqMzY-3R3yGC-1}\nwixd{NW2DqMzY-33zuOn-1}\nwixd{NW2DqMzY-33zuOn-2}\nwixd{NW2DqMzY-33zuOn-3}\nwixd{NW2DqMzY-33zuOn-4}\nwixd{NW2DqMzY-33zuOn-5}}}%
\nwixlogsorted{c}{{Specification for $FS_{0}$ functions}{NW2DqMzY-1uF0T3-1}{\nwixu{NW2DqMzY-3R3yGC-1}\nwixd{NW2DqMzY-1uF0T3-1}\nwixd{NW2DqMzY-1uF0T3-2}\nwixd{NW2DqMzY-1uF0T3-3}\nwixd{NW2DqMzY-1uF0T3-4}\nwixd{NW2DqMzY-1uF0T3-5}\nwixd{NW2DqMzY-1uF0T3-6}\nwixd{NW2DqMzY-1uF0T3-7}\nwixd{NW2DqMzY-1uF0T3-8}}}%
\nwixlogsorted{c}{{Specification for $FS_{0}$ individuals}{NW2DqMzY-24mg92-1}{\nwixu{NW2DqMzY-3R3yGC-1}\nwixd{NW2DqMzY-24mg92-1}}}%
\nwixlogsorted{c}{{Specification for $FS_{0}$ predicates}{NW2DqMzY-194t0N-1}{\nwixu{NW2DqMzY-3R3yGC-1}\nwixd{NW2DqMzY-194t0N-1}\nwixd{NW2DqMzY-194t0N-2}\nwixd{NW2DqMzY-194t0N-3}}}%
\nwixlogsorted{c}{{Specification for formulas}{NW3w0iRO-2HPkc8-1}{\nwixu{NW3w0iRO-5BH7S-1}\nwixd{NW3w0iRO-2HPkc8-1}\nwixd{NW3w0iRO-2HPkc8-2}\nwixd{NW3w0iRO-2HPkc8-3}\nwixd{NW3w0iRO-2HPkc8-4}\nwixd{NW3w0iRO-2HPkc8-5}\nwixd{NW3w0iRO-2HPkc8-6}\nwixd{NW3w0iRO-2HPkc8-7}\nwixd{NW3w0iRO-2HPkc8-8}\nwixd{NW3w0iRO-2HPkc8-9}\nwixd{NW3w0iRO-2HPkc8-A}\nwixd{NW3w0iRO-2HPkc8-B}\nwixd{NW3w0iRO-2HPkc8-C}\nwixd{NW3w0iRO-2HPkc8-D}\nwixd{NW3w0iRO-2HPkc8-E}\nwixd{NW3w0iRO-2HPkc8-F}\nwixd{NW3w0iRO-2HPkc8-G}\nwixd{NW3w0iRO-2HPkc8-H}}}%
\nwixlogsorted{c}{{Specification for inference rules}{NW4ICsJQ-1XQlFy-1}{\nwixu{NW4ICsJQ-2cY0B1-1}\nwixd{NW4ICsJQ-1XQlFy-1}\nwixd{NW4ICsJQ-1XQlFy-2}\nwixd{NW4ICsJQ-1XQlFy-3}\nwixd{NW4ICsJQ-1XQlFy-4}\nwixd{NW4ICsJQ-1XQlFy-5}\nwixd{NW4ICsJQ-1XQlFy-6}\nwixd{NW4ICsJQ-1XQlFy-7}\nwixd{NW4ICsJQ-1XQlFy-8}\nwixd{NW4ICsJQ-1XQlFy-9}\nwixd{NW4ICsJQ-1XQlFy-A}}}%
\nwixlogsorted{c}{{Specification for tactics}{NW1F32vJ-2kzjE-1}{\nwixu{NW1F32vJ-VKxLl-1}\nwixd{NW1F32vJ-2kzjE-1}\nwixd{NW1F32vJ-2kzjE-2}\nwixd{NW1F32vJ-2kzjE-3}\nwixd{NW1F32vJ-2kzjE-4}\nwixd{NW1F32vJ-2kzjE-5}\nwixd{NW1F32vJ-2kzjE-6}\nwixd{NW1F32vJ-2kzjE-7}\nwixd{NW1F32vJ-2kzjE-8}\nwixd{NW1F32vJ-2kzjE-9}\nwixd{NW1F32vJ-2kzjE-A}\nwixd{NW1F32vJ-2kzjE-B}}}%
\nwixlogsorted{c}{{Specification for test result}{NWsR2q7-1ETFk0-1}{\nwixu{NWsR2q7-1Fzsdc-1}\nwixd{NWsR2q7-1ETFk0-1}\nwixd{NWsR2q7-1ETFk0-2}\nwixd{NWsR2q7-1ETFk0-3}\nwixd{NWsR2q7-1ETFk0-4}\nwixd{NWsR2q7-1ETFk0-5}\nwixd{NWsR2q7-1ETFk0-6}\nwixd{NWsR2q7-1ETFk0-7}\nwixd{NWsR2q7-1ETFk0-8}\nwixd{NWsR2q7-1ETFk0-9}}}%
\nwixlogsorted{c}{{State.sig}{NWU4iYe-1UXQzr-1}{\nwixd{NWU4iYe-1UXQzr-1}}}%
\nwixlogsorted{c}{{Substitution in a formula}{NW3w0iRO-1jtQRV-1}{\nwixu{NW3w0iRO-2jnj9W-1}\nwixd{NW3w0iRO-1jtQRV-1}}}%
\nwixlogsorted{c}{{Substitution of a term for a variable in a term}{NW1rTmOJ-3OIyM5-1}{\nwixu{NW1rTmOJ-4LsRWC-1}\nwixd{NW1rTmOJ-3OIyM5-1}}}%
\nwixlogsorted{c}{{Tableau method}{NW1mR7zJ-1k5nOE-1}{\nwixd{NW1mR7zJ-1k5nOE-1}\nwixu{NW1mR7zJ-yCIr5-1}}}%
\nwixlogsorted{c}{{Tableaux refutation}{NW1mR7zJ-3SEP5w-1}{\nwixd{NW1mR7zJ-3SEP5w-1}\nwixu{NW1mR7zJ-yCIr5-1}}}%
\nwixlogsorted{c}{{Tableaux refuter}{NW1mR7zJ-yCIr5-1}{\nwixu{NW1mR7zJ-l5jX3-1}\nwixd{NW1mR7zJ-yCIr5-1}}}%
\nwixlogsorted{c}{{Tactic for negation}{NW1F32vJ-1G1Cmu-1}{\nwixu{NW1F32vJ-3Cb3oS-3}\nwixd{NW1F32vJ-1G1Cmu-1}}}%
\nwixlogsorted{c}{{Tactic to discharge an implication}{NW1F32vJ-IHUmU-1}{\nwixu{NW1F32vJ-3Cb3oS-3}\nwixd{NW1F32vJ-IHUmU-1}}}%
\nwixlogsorted{c}{{Tactic.sig}{NW1F32vJ-VKxLl-1}{\nwixd{NW1F32vJ-VKxLl-1}}}%
\nwixlogsorted{c}{{Tactic.sml}{NW1F32vJ-2XH0oT-1}{\nwixd{NW1F32vJ-2XH0oT-1}}}%
\nwixlogsorted{c}{{Take the union of applying a function to atoms}{NW3w0iRO-28e6MC-1}{\nwixu{NW3w0iRO-2jnj9W-1}\nwixd{NW3w0iRO-28e6MC-1}}}%
\nwixlogsorted{c}{{TBPrinter for a formula}{NW24q7Pv-14xJt9-1}{\nwixu{NW24q7Pv-3JbCWS-1}\nwixd{NW24q7Pv-14xJt9-1}}}%
\nwixlogsorted{c}{{TBPrinter for a term}{NW24q7Pv-3IZAOd-1}{\nwixu{NW24q7Pv-3JbCWS-1}\nwixd{NW24q7Pv-3IZAOd-1}}}%
\nwixlogsorted{c}{{TBPrinter for a theorem}{NW24q7Pv-2pSwbS-1}{\nwixu{NW24q7Pv-3JbCWS-1}\nwixd{NW24q7Pv-2pSwbS-1}}}%
\nwixlogsorted{c}{{TBPrinter for goals}{NW24q7Pv-3yfJXJ-1}{\nwixu{NW24q7Pv-3JbCWS-1}\nwixd{NW24q7Pv-3yfJXJ-1}}}%
\nwixlogsorted{c}{{TBPrinter.sml}{NW24q7Pv-3JbCWS-1}{\nwixd{NW24q7Pv-3JbCWS-1}}}%
\nwixlogsorted{c}{{Term.sml}{NW1rTmOJ-4LsRWC-1}{\nwixd{NW1rTmOJ-4LsRWC-1}}}%
\nwixlogsorted{c}{{Test a term is free for a variable in a formula}{NW3w0iRO-ZZefY-1}{\nwixu{NW3w0iRO-2jnj9W-1}\nwixd{NW3w0iRO-ZZefY-1}}}%
\nwixlogsorted{c}{{Test for $E^{+}$-formula}{NW3w0iRO-1O1TnJ-1}{\nwixu{NW3w0iRO-2jnj9W-1}\nwixd{NW3w0iRO-1O1TnJ-1}}}%
\nwixlogsorted{c}{{Test if a new equations is trivial for unification}{NW1BOcHB-4edp1Q-1}{\nwixu{NW1BOcHB-25AKca-1}\nwixd{NW1BOcHB-4edp1Q-1}}}%
\nwixlogsorted{c}{{Test if term contains subterm}{NW1rTmOJ-Re21r-1}{\nwixu{NW1rTmOJ-4LsRWC-1}\nwixd{NW1rTmOJ-Re21r-1}}}%
\nwixlogsorted{c}{{Test if two formulas are equal}{NW3w0iRO-5gQzE-1}{\nwixu{NW3w0iRO-2jnj9W-1}\nwixd{NW3w0iRO-5gQzE-1}}}%
\nwixlogsorted{c}{{Test if two terms have the same sort}{NW1rTmOJ-1WQNps-1}{\nwixu{NW1rTmOJ-4LsRWC-1}\nwixd{NW1rTmOJ-1WQNps-1}}}%
\nwixlogsorted{c}{{test.mlb}{NWsR2q7-2v996w-1}{\nwixd{NWsR2q7-2v996w-1}}}%
\nwixlogsorted{c}{{Test.sig}{NWsR2q7-1Fzsdc-1}{\nwixd{NWsR2q7-1Fzsdc-1}}}%
\nwixlogsorted{c}{{Test.sml}{NWsR2q7-3ucTCC-1}{\nwixd{NWsR2q7-3ucTCC-1}}}%
\nwixlogsorted{c}{{TestRunner.sig}{NWsR2q7-4Peuwc-1}{\nwixd{NWsR2q7-4Peuwc-1}}}%
\nwixlogsorted{c}{{tests/ClassicalSuite.sml}{NW1mR7zJ-10sN88-1}{\nwixd{NW1mR7zJ-10sN88-1}}}%
\nwixlogsorted{c}{{tests/main.sml}{NWsR2q7-3MeQDa-1}{\nwixd{NWsR2q7-3MeQDa-1}}}%
\nwixlogsorted{c}{{tests/MkRunner.sml}{NWsR2q7-2dKvEw-1}{\nwixd{NWsR2q7-2dKvEw-1}}}%
\nwixlogsorted{c}{{tests/TestSuite.sig}{NWsR2q7-1mlFrS-1}{\nwixd{NWsR2q7-1mlFrS-1}}}%
\nwixlogsorted{c}{{Theorem: implications is reflexive}{NW1NjfzY-2PXiqP-1}{\nwixu{NW1NjfzY-4CUt2L-5}\nwixd{NW1NjfzY-2PXiqP-1}}}%
\nwixlogsorted{c}{{Thm.sig}{NW4ICsJQ-2cY0B1-1}{\nwixd{NW4ICsJQ-2cY0B1-1}}}%
\nwixlogsorted{c}{{Thm.sml}{NW4ICsJQ-YIPBT-1}{\nwixd{NW4ICsJQ-YIPBT-1}}}%
\nwixlogsorted{c}{{Transform formula to negation normal form}{NW1mR7zJ-4TkWUc-1}{\nwixu{NW1mR7zJ-l5jX3-1}\nwixd{NW1mR7zJ-4TkWUc-1}}}%
\nwixlogsorted{c}{{Transitivity of implies}{NW1NjfzY-1sm2GN-1}{\nwixu{NW1NjfzY-4CUt2L-5}\nwixd{NW1NjfzY-1sm2GN-1}}}%
\nwixlogsorted{c}{{Uncurrying implication}{NW1NjfzY-1JZnQr-1}{\nwixu{NW1NjfzY-4CUt2L-5}\nwixd{NW1NjfzY-1JZnQr-1}}}%
\nwixlogsorted{c}{{Unif.sig}{NW1BOcHB-3z172B-1}{\nwixd{NW1BOcHB-3z172B-1}}}%
\nwixlogsorted{c}{{Unif.sml}{NW1BOcHB-1KP4Zf-1}{\nwixd{NW1BOcHB-1KP4Zf-1}}}%
\nwixlogsorted{c}{{Unify two literals}{NW1BOcHB-3w0JS3-1}{\nwixu{NW1BOcHB-1KP4Zf-1}\nwixd{NW1BOcHB-3w0JS3-1}}}%
\nwixlogsorted{c}{{Unify two terms}{NW1BOcHB-1kcz6k-1}{\nwixu{NW1BOcHB-1KP4Zf-1}\nwixd{NW1BOcHB-1kcz6k-1}}}%
\nwixlogsorted{c}{{Unions of two $FS_{0}$ classes}{NW2DqMzY-22W7o7-1}{\nwixu{NW2DqMzY-TSvER-1}\nwixd{NW2DqMzY-22W7o7-1}}}%
\nwixlogsorted{c}{{Unit tests for \code{}Classical\edoc{}}{NW1mR7zJ-4QAFgj-1}{\nwixu{NW1mR7zJ-10sN88-1}\nwixd{NW1mR7zJ-4QAFgj-1}}}%
\nwixlogsorted{c}{{Unit tests for \code{}Derived.sml\edoc{}}{NW1NjfzY-3rvgpx-1}{\nwixd{NW1NjfzY-3rvgpx-1}\nwixd{NW1NjfzY-3rvgpx-2}\nwixd{NW1NjfzY-3rvgpx-3}\nwixd{NW1NjfzY-3rvgpx-4}\nwixd{NW1NjfzY-3rvgpx-5}\nwixd{NW1NjfzY-3rvgpx-6}\nwixd{NW1NjfzY-3rvgpx-7}\nwixd{NW1NjfzY-3rvgpx-8}\nwixd{NW1NjfzY-3rvgpx-9}\nwixd{NW1NjfzY-3rvgpx-A}}}%
\nwixlogsorted{c}{{Unit tests for derived rules}{NW1NjfzY-3DeU7m-1}{\nwixd{NW1NjfzY-3DeU7m-1}}}%
\nwixlogsorted{c}{{Unit tests to run}{NW1mR7zJ-3lyaJm-1}{\nwixd{NW1mR7zJ-3lyaJm-1}\nwixu{NWsR2q7-2v996w-1}}}%
\nwixlogsorted{c}{{Variant of an identifier}{NW1mR7zJ-2QPys7-1}{\nwixd{NW1mR7zJ-2QPys7-1}\nwixu{NW1mR7zJ-4B0sXE-1}}}%
\nwixlogsorted{c}{{Weak negation elimination}{NW1NjfzY-16XbWD-1}{\nwixu{NW1NjfzY-4CUt2L-1}\nwixd{NW1NjfzY-16XbWD-1}}}%
\nwixlogsorted{c}{{Weakening inference rule}{NW1NjfzY-2n5HPL-1}{\nwixu{NW1NjfzY-4CUt2L-1}\nwixd{NW1NjfzY-2n5HPL-1}}}%
\nwixlogsorted{i}{{\nwixident{@@}}{@@}}%
\nwixlogsorted{i}{{\nwixident{And}}{And}}%
\nwixlogsorted{i}{{\nwixident{Assert.eq}}{Assert.eq}}%
\nwixlogsorted{i}{{\nwixident{Assert.Fail}}{Assert.Fail}}%
\nwixlogsorted{i}{{\nwixident{Assert.that}}{Assert.that}}%
\nwixlogsorted{i}{{\nwixident{assumptions{\_}contain{\_}hypotheses}}{assumptions:uncontain:unhypotheses}}%
\nwixlogsorted{i}{{\nwixident{Classical.tab}}{Classical.tab}}%
\nwixlogsorted{i}{{\nwixident{Derived.absurdum}}{Derived.absurdum}}%
\nwixlogsorted{i}{{\nwixident{Derived.all{\_}conj{\_}distrib}}{Derived.all:unconj:undistrib}}%
\nwixlogsorted{i}{{\nwixident{Derived.all{\_}disjL}}{Derived.all:undisjL}}%
\nwixlogsorted{i}{{\nwixident{Derived.all{\_}imp}}{Derived.all:unimp}}%
\nwixlogsorted{i}{{\nwixident{Derived.and{\_}assoc}}{Derived.and:unassoc}}%
\nwixlogsorted{i}{{\nwixident{Derived.and{\_}comm}}{Derived.and:uncomm}}%
\nwixlogsorted{i}{{\nwixident{Derived.and{\_}cong{\_}l}}{Derived.and:uncong:unl}}%
\nwixlogsorted{i}{{\nwixident{Derived.and{\_}cong{\_}r}}{Derived.and:uncong:unr}}%
\nwixlogsorted{i}{{\nwixident{Derived.and{\_}thm}}{Derived.and:unthm}}%
\nwixlogsorted{i}{{\nwixident{Derived.conj{\_}disj{\_}distribL}}{Derived.conj:undisj:undistribL}}%
\nwixlogsorted{i}{{\nwixident{Derived.conj{\_}disj{\_}distribR}}{Derived.conj:undisj:undistribR}}%
\nwixlogsorted{i}{{\nwixident{Derived.contrapositive}}{Derived.contrapositive}}%
\nwixlogsorted{i}{{\nwixident{Derived.contrapositive{\_}not}}{Derived.contrapositive:unnot}}%
\nwixlogsorted{i}{{\nwixident{Derived.curry{\_}imp}}{Derived.curry:unimp}}%
\nwixlogsorted{i}{{\nwixident{Derived.disj{\_}imp{\_}disj}}{Derived.disj:unimp:undisj}}%
\nwixlogsorted{i}{{\nwixident{Derived.ex{\_}conj{\_}distrib}}{Derived.ex:unconj:undistrib}}%
\nwixlogsorted{i}{{\nwixident{Derived.ex{\_}disj{\_}distrib}}{Derived.ex:undisj:undistrib}}%
\nwixlogsorted{i}{{\nwixident{Derived.ex{\_}imp}}{Derived.ex:unimp}}%
\nwixlogsorted{i}{{\nwixident{Derived.ex{\_}not}}{Derived.ex:unnot}}%
\nwixlogsorted{i}{{\nwixident{Derived.exists{\_}cong}}{Derived.exists:uncong}}%
\nwixlogsorted{i}{{\nwixident{Derived.exists{\_}elim}}{Derived.exists:unelim}}%
\nwixlogsorted{i}{{\nwixident{Derived.forall{\_}cong}}{Derived.forall:uncong}}%
\nwixlogsorted{i}{{\nwixident{Derived.frege{\_}id}}{Derived.frege:unid}}%
\nwixlogsorted{i}{{\nwixident{Derived.hypo{\_}syll}}{Derived.hypo:unsyll}}%
\nwixlogsorted{i}{{\nwixident{Derived.hypo{\_}syll{\_}thm}}{Derived.hypo:unsyll:unthm}}%
\nwixlogsorted{i}{{\nwixident{Derived.iff{\_}refl}}{Derived.iff:unrefl}}%
\nwixlogsorted{i}{{\nwixident{Derived.iff{\_}sym}}{Derived.iff:unsym}}%
\nwixlogsorted{i}{{\nwixident{Derived.iff{\_}trans}}{Derived.iff:untrans}}%
\nwixlogsorted{i}{{\nwixident{Derived.imp{\_}all}}{Derived.imp:unall}}%
\nwixlogsorted{i}{{\nwixident{Derived.imp{\_}ex}}{Derived.imp:unex}}%
\nwixlogsorted{i}{{\nwixident{Derived.imp{\_}is{\_}not{\_}and}}{Derived.imp:unis:unnot:unand}}%
\nwixlogsorted{i}{{\nwixident{Derived.imp{\_}trans}}{Derived.imp:untrans}}%
\nwixlogsorted{i}{{\nwixident{Derived.not{\_}all{\_}iff{\_}ex{\_}not}}{Derived.not:unall:uniff:unex:unnot}}%
\nwixlogsorted{i}{{\nwixident{Derived.not{\_}and}}{Derived.not:unand}}%
\nwixlogsorted{i}{{\nwixident{Derived.not{\_}disj{\_}to{\_}imp}}{Derived.not:undisj:unto:unimp}}%
\nwixlogsorted{i}{{\nwixident{Derived.not{\_}or}}{Derived.not:unor}}%
\nwixlogsorted{i}{{\nwixident{Derived.or{\_}assoc}}{Derived.or:unassoc}}%
\nwixlogsorted{i}{{\nwixident{Derived.or{\_}comm}}{Derived.or:uncomm}}%
\nwixlogsorted{i}{{\nwixident{Derived.or{\_}cong{\_}l}}{Derived.or:uncong:unl}}%
\nwixlogsorted{i}{{\nwixident{Derived.or{\_}cong{\_}r}}{Derived.or:uncong:unr}}%
\nwixlogsorted{i}{{\nwixident{Derived.uncurry{\_}imp}}{Derived.uncurry:unimp}}%
\nwixlogsorted{i}{{\nwixident{Derived.weak{\_}not{\_}elim}}{Derived.weak:unnot:unelim}}%
\nwixlogsorted{i}{{\nwixident{Derived.weaken}}{Derived.weaken}}%
\nwixlogsorted{i}{{\nwixident{Env.apply}}{Env.apply}}%
\nwixlogsorted{i}{{\nwixident{Env.defined}}{Env.defined}}%
\nwixlogsorted{i}{{\nwixident{Env.insert}}{Env.insert}}%
\nwixlogsorted{i}{{\nwixident{Env.t}}{Env.t}}%
\nwixlogsorted{i}{{\nwixident{Env.undefined}}{Env.undefined}}%
\nwixlogsorted{i}{{\nwixident{Exists}}{Exists}}%
\nwixlogsorted{i}{{\nwixident{False}}{False}}%
\nwixlogsorted{i}{{\nwixident{Forall}}{Forall}}%
\nwixlogsorted{i}{{\nwixident{Formula.and{\_}elim{\_}l}}{Formula.and:unelim:unl}}%
\nwixlogsorted{i}{{\nwixident{Formula.and{\_}elim{\_}r}}{Formula.and:unelim:unr}}%
\nwixlogsorted{i}{{\nwixident{Formula.and{\_}intro}}{Formula.and:unintro}}%
\nwixlogsorted{i}{{\nwixident{Formula.atom{\_}union}}{Formula.atom:ununion}}%
\nwixlogsorted{i}{{\nwixident{Formula.bv}}{Formula.bv}}%
\nwixlogsorted{i}{{\nwixident{Formula.compare}}{Formula.compare}}%
\nwixlogsorted{i}{{\nwixident{Formula.contains{\_}var}}{Formula.contains:unvar}}%
\nwixlogsorted{i}{{\nwixident{Formula.Dest}}{Formula.Dest}}%
\nwixlogsorted{i}{{\nwixident{Formula.dest{\_}and}}{Formula.dest:unand}}%
\nwixlogsorted{i}{{\nwixident{Formula.dest{\_}iff}}{Formula.dest:uniff}}%
\nwixlogsorted{i}{{\nwixident{Formula.dest{\_}imp}}{Formula.dest:unimp}}%
\nwixlogsorted{i}{{\nwixident{Formula.dest{\_}not}}{Formula.dest:unnot}}%
\nwixlogsorted{i}{{\nwixident{Formula.dest{\_}or}}{Formula.dest:unor}}%
\nwixlogsorted{i}{{\nwixident{Formula.dest{\_}pred}}{Formula.dest:unpred}}%
\nwixlogsorted{i}{{\nwixident{Formula.eq}}{Formula.eq}}%
\nwixlogsorted{i}{{\nwixident{Formula.foldl}}{Formula.foldl}}%
\nwixlogsorted{i}{{\nwixident{Formula.foldr}}{Formula.foldr}}%
\nwixlogsorted{i}{{\nwixident{Formula.free{\_}for}}{Formula.free:unfor}}%
\nwixlogsorted{i}{{\nwixident{Formula.fv}}{Formula.fv}}%
\nwixlogsorted{i}{{\nwixident{Formula.is{\_}and}}{Formula.is:unand}}%
\nwixlogsorted{i}{{\nwixident{Formula.is{\_}e{\_}plus}}{Formula.is:une:unplus}}%
\nwixlogsorted{i}{{\nwixident{Formula.is{\_}falsum}}{Formula.is:unfalsum}}%
\nwixlogsorted{i}{{\nwixident{Formula.is{\_}iff}}{Formula.is:uniff}}%
\nwixlogsorted{i}{{\nwixident{Formula.is{\_}imp}}{Formula.is:unimp}}%
\nwixlogsorted{i}{{\nwixident{Formula.is{\_}not}}{Formula.is:unnot}}%
\nwixlogsorted{i}{{\nwixident{Formula.is{\_}or}}{Formula.is:unor}}%
\nwixlogsorted{i}{{\nwixident{Formula.is{\_}pred}}{Formula.is:unpred}}%
\nwixlogsorted{i}{{\nwixident{Formula.is{\_}verum}}{Formula.is:unverum}}%
\nwixlogsorted{i}{{\nwixident{Formula.mk{\_}and}}{Formula.mk:unand}}%
\nwixlogsorted{i}{{\nwixident{Formula.mk{\_}exists}}{Formula.mk:unexists}}%
\nwixlogsorted{i}{{\nwixident{Formula.mk{\_}falsum}}{Formula.mk:unfalsum}}%
\nwixlogsorted{i}{{\nwixident{Formula.mk{\_}forall}}{Formula.mk:unforall}}%
\nwixlogsorted{i}{{\nwixident{Formula.mk{\_}iff}}{Formula.mk:uniff}}%
\nwixlogsorted{i}{{\nwixident{Formula.mk{\_}imp}}{Formula.mk:unimp}}%
\nwixlogsorted{i}{{\nwixident{Formula.mk{\_}not}}{Formula.mk:unnot}}%
\nwixlogsorted{i}{{\nwixident{Formula.mk{\_}or}}{Formula.mk:unor}}%
\nwixlogsorted{i}{{\nwixident{Formula.mk{\_}pred}}{Formula.mk:unpred}}%
\nwixlogsorted{i}{{\nwixident{Formula.mk{\_}verum}}{Formula.mk:unverum}}%
\nwixlogsorted{i}{{\nwixident{Formula.precedence}}{Formula.precedence}}%
\nwixlogsorted{i}{{\nwixident{Formula.serialize}}{Formula.serialize}}%
\nwixlogsorted{i}{{\nwixident{Formula.subst}}{Formula.subst}}%
\nwixlogsorted{i}{{\nwixident{Formula.t}}{Formula.t}}%
\nwixlogsorted{i}{{\nwixident{FS0.app}}{FS0.app}}%
\nwixlogsorted{i}{{\nwixident{FS0.comp}}{FS0.comp}}%
\nwixlogsorted{i}{{\nwixident{FS0.const}}{FS0.const}}%
\nwixlogsorted{i}{{\nwixident{FS0.dest{\_}app}}{FS0.dest:unapp}}%
\nwixlogsorted{i}{{\nwixident{FS0.dest{\_}comp}}{FS0.dest:uncomp}}%
\nwixlogsorted{i}{{\nwixident{FS0.dest{\_}const}}{FS0.dest:unconst}}%
\nwixlogsorted{i}{{\nwixident{FS0.dest{\_}equal}}{FS0.dest:unequal}}%
\nwixlogsorted{i}{{\nwixident{FS0.dest{\_}in}}{FS0.dest:unin}}%
\nwixlogsorted{i}{{\nwixident{FS0.dest{\_}induct}}{FS0.dest:uninduct}}%
\nwixlogsorted{i}{{\nwixident{FS0.dest{\_}intersection}}{FS0.dest:unintersection}}%
\nwixlogsorted{i}{{\nwixident{FS0.dest{\_}juxt}}{FS0.dest:unjuxt}}%
\nwixlogsorted{i}{{\nwixident{FS0.dest{\_}pair}}{FS0.dest:unpair}}%
\nwixlogsorted{i}{{\nwixident{FS0.dest{\_}preimage}}{FS0.dest:unpreimage}}%
\nwixlogsorted{i}{{\nwixident{FS0.dest{\_}recur}}{FS0.dest:unrecur}}%
\nwixlogsorted{i}{{\nwixident{FS0.dest{\_}singleton}}{FS0.dest:unsingleton}}%
\nwixlogsorted{i}{{\nwixident{FS0.dest{\_}union}}{FS0.dest:ununion}}%
\nwixlogsorted{i}{{\nwixident{FS0.equal}}{FS0.equal}}%
\nwixlogsorted{i}{{\nwixident{FS0.id}}{FS0.id}}%
\nwixlogsorted{i}{{\nwixident{FS0.if{\_}eq}}{FS0.if:uneq}}%
\nwixlogsorted{i}{{\nwixident{FS0.In}}{FS0.In}}%
\nwixlogsorted{i}{{\nwixident{FS0.induct}}{FS0.induct}}%
\nwixlogsorted{i}{{\nwixident{FS0.intersection}}{FS0.intersection}}%
\nwixlogsorted{i}{{\nwixident{FS0.juxt}}{FS0.juxt}}%
\nwixlogsorted{i}{{\nwixident{FS0.mk{\_}id}}{FS0.mk:unid}}%
\nwixlogsorted{i}{{\nwixident{FS0.mk{\_}if{\_}eq}}{FS0.mk:unif:uneq}}%
\nwixlogsorted{i}{{\nwixident{FS0.mk{\_}proj1}}{FS0.mk:unproj1}}%
\nwixlogsorted{i}{{\nwixident{FS0.mk{\_}proj2}}{FS0.mk:unproj2}}%
\nwixlogsorted{i}{{\nwixident{FS0.not{\_}equal}}{FS0.not:unequal}}%
\nwixlogsorted{i}{{\nwixident{FS0.pair}}{FS0.pair}}%
\nwixlogsorted{i}{{\nwixident{FS0.preimage}}{FS0.preimage}}%
\nwixlogsorted{i}{{\nwixident{FS0.proj1}}{FS0.proj1}}%
\nwixlogsorted{i}{{\nwixident{FS0.proj2}}{FS0.proj2}}%
\nwixlogsorted{i}{{\nwixident{FS0.recur}}{FS0.recur}}%
\nwixlogsorted{i}{{\nwixident{FS0.singleton}}{FS0.singleton}}%
\nwixlogsorted{i}{{\nwixident{FS0.union}}{FS0.union}}%
\nwixlogsorted{i}{{\nwixident{FS0.zero}}{FS0.zero}}%
\nwixlogsorted{i}{{\nwixident{Iff}}{Iff}}%
\nwixlogsorted{i}{{\nwixident{Implies}}{Implies}}%
\nwixlogsorted{i}{{\nwixident{Not}}{Not}}%
\nwixlogsorted{i}{{\nwixident{Or}}{Or}}%
\nwixlogsorted{i}{{\nwixident{Pred}}{Pred}}%
\nwixlogsorted{i}{{\nwixident{remove}}{remove}}%
\nwixlogsorted{i}{{\nwixident{Sort.CLASS}}{Sort.CLASS}}%
\nwixlogsorted{i}{{\nwixident{Sort.compare}}{Sort.compare}}%
\nwixlogsorted{i}{{\nwixident{Sort.FUN}}{Sort.FUN}}%
\nwixlogsorted{i}{{\nwixident{Sort.IND}}{Sort.IND}}%
\nwixlogsorted{i}{{\nwixident{Sort.t}}{Sort.t}}%
\nwixlogsorted{i}{{\nwixident{Sort.to{\_}string}}{Sort.to:unstring}}%
\nwixlogsorted{i}{{\nwixident{Tactic.and{\_}elim{\_}l}}{Tactic.and:unelim:unl}}%
\nwixlogsorted{i}{{\nwixident{Tactic.and{\_}elim{\_}r}}{Tactic.and:unelim:unr}}%
\nwixlogsorted{i}{{\nwixident{Tactic.and{\_}intro}}{Tactic.and:unintro}}%
\nwixlogsorted{i}{{\nwixident{Tactic.contr}}{Tactic.contr}}%
\nwixlogsorted{i}{{\nwixident{Tactic.disch}}{Tactic.disch}}%
\nwixlogsorted{i}{{\nwixident{Tactic.disj{\_}cases}}{Tactic.disj:uncases}}%
\nwixlogsorted{i}{{\nwixident{Tactic.gen}}{Tactic.gen}}%
\nwixlogsorted{i}{{\nwixident{Tactic.modus{\_}ponens}}{Tactic.modus:unponens}}%
\nwixlogsorted{i}{{\nwixident{Tactic.mp}}{Tactic.mp}}%
\nwixlogsorted{i}{{\nwixident{Tactic.or{\_}l}}{Tactic.or:unl}}%
\nwixlogsorted{i}{{\nwixident{Tactic.or{\_}r}}{Tactic.or:unr}}%
\nwixlogsorted{i}{{\nwixident{Tactic.prove}}{Tactic.prove}}%
\nwixlogsorted{i}{{\nwixident{Tactic.spec}}{Tactic.spec}}%
\nwixlogsorted{i}{{\nwixident{TBPrinter.formula}}{TBPrinter.formula}}%
\nwixlogsorted{i}{{\nwixident{TBPrinter.goals}}{TBPrinter.goals}}%
\nwixlogsorted{i}{{\nwixident{TBPrinter.term}}{TBPrinter.term}}%
\nwixlogsorted{i}{{\nwixident{TBPrinter.thm}}{TBPrinter.thm}}%
\nwixlogsorted{i}{{\nwixident{Term.compare}}{Term.compare}}%
\nwixlogsorted{i}{{\nwixident{Term.compare{\_}lists}}{Term.compare:unlists}}%
\nwixlogsorted{i}{{\nwixident{Term.eq}}{Term.eq}}%
\nwixlogsorted{i}{{\nwixident{Term.eq{\_}lists}}{Term.eq:unlists}}%
\nwixlogsorted{i}{{\nwixident{Term.fresh{\_}var}}{Term.fresh:unvar}}%
\nwixlogsorted{i}{{\nwixident{Term.fv}}{Term.fv}}%
\nwixlogsorted{i}{{\nwixident{Term.have{\_}same{\_}sort}}{Term.have:unsame:unsort}}%
\nwixlogsorted{i}{{\nwixident{Term.is{\_}class}}{Term.is:unclass}}%
\nwixlogsorted{i}{{\nwixident{Term.is{\_}fun}}{Term.is:unfun}}%
\nwixlogsorted{i}{{\nwixident{Term.is{\_}ind}}{Term.is:unind}}%
\nwixlogsorted{i}{{\nwixident{Term.occurs{\_}in}}{Term.occurs:unin}}%
\nwixlogsorted{i}{{\nwixident{Term.serialize}}{Term.serialize}}%
\nwixlogsorted{i}{{\nwixident{Term.sort}}{Term.sort}}%
\nwixlogsorted{i}{{\nwixident{Term.subst}}{Term.subst}}%
\nwixlogsorted{i}{{\nwixident{Term.t}}{Term.t}}%
\nwixlogsorted{i}{{\nwixident{Test.mk}}{Test.mk}}%
\nwixlogsorted{i}{{\nwixident{Test.Result.count{\_}errors}}{Test.Result.count:unerrors}}%
\nwixlogsorted{i}{{\nwixident{Test.Result.count{\_}failures}}{Test.Result.count:unfailures}}%
\nwixlogsorted{i}{{\nwixident{Test.Result.count{\_}successes}}{Test.Result.count:unsuccesses}}%
\nwixlogsorted{i}{{\nwixident{Test.Result.count{\_}total}}{Test.Result.count:untotal}}%
\nwixlogsorted{i}{{\nwixident{Test.Result.exn{\_}message}}{Test.Result.exn:unmessage}}%
\nwixlogsorted{i}{{\nwixident{Test.Result.for{\_}case}}{Test.Result.for:uncase}}%
\nwixlogsorted{i}{{\nwixident{Test.Result.for{\_}suite}}{Test.Result.for:unsuite}}%
\nwixlogsorted{i}{{\nwixident{Test.Result.is{\_}case}}{Test.Result.is:uncase}}%
\nwixlogsorted{i}{{\nwixident{Test.Result.is{\_}error}}{Test.Result.is:unerror}}%
\nwixlogsorted{i}{{\nwixident{Test.Result.is{\_}failure}}{Test.Result.is:unfailure}}%
\nwixlogsorted{i}{{\nwixident{Test.Result.is{\_}success}}{Test.Result.is:unsuccess}}%
\nwixlogsorted{i}{{\nwixident{Test.Result.is{\_}suite}}{Test.Result.is:unsuite}}%
\nwixlogsorted{i}{{\nwixident{Test.Result.msg}}{Test.Result.msg}}%
\nwixlogsorted{i}{{\nwixident{Test.Result.name}}{Test.Result.name}}%
\nwixlogsorted{i}{{\nwixident{Test.Result.realtime}}{Test.Result.realtime}}%
\nwixlogsorted{i}{{\nwixident{Test.Result.results}}{Test.Result.results}}%
\nwixlogsorted{i}{{\nwixident{Test.Result.runtime}}{Test.Result.runtime}}%
\nwixlogsorted{i}{{\nwixident{Test.Result.t}}{Test.Result.t}}%
\nwixlogsorted{i}{{\nwixident{Test.run}}{Test.run}}%
\nwixlogsorted{i}{{\nwixident{Test.suite}}{Test.suite}}%
\nwixlogsorted{i}{{\nwixident{Test.t}}{Test.t}}%
\nwixlogsorted{i}{{\nwixident{Thm.assume}}{Thm.assume}}%
\nwixlogsorted{i}{{\nwixident{Thm.axiom{\_}class{\_}comp}}{Thm.axiom:unclass:uncomp}}%
\nwixlogsorted{i}{{\nwixident{Thm.axiom{\_}class{\_}eq}}{Thm.axiom:unclass:uneq}}%
\nwixlogsorted{i}{{\nwixident{Thm.axiom{\_}comp}}{Thm.axiom:uncomp}}%
\nwixlogsorted{i}{{\nwixident{Thm.axiom{\_}cond{\_}err}}{Thm.axiom:uncond:unerr}}%
\nwixlogsorted{i}{{\nwixident{Thm.axiom{\_}cond{\_}false}}{Thm.axiom:uncond:unfalse}}%
\nwixlogsorted{i}{{\nwixident{Thm.axiom{\_}cond{\_}true}}{Thm.axiom:uncond:untrue}}%
\nwixlogsorted{i}{{\nwixident{Thm.axiom{\_}const}}{Thm.axiom:unconst}}%
\nwixlogsorted{i}{{\nwixident{Thm.axiom{\_}fun{\_}eq}}{Thm.axiom:unfun:uneq}}%
\nwixlogsorted{i}{{\nwixident{Thm.axiom{\_}id}}{Thm.axiom:unid}}%
\nwixlogsorted{i}{{\nwixident{Thm.axiom{\_}induct{\_}base}}{Thm.axiom:uninduct:unbase}}%
\nwixlogsorted{i}{{\nwixident{Thm.axiom{\_}induct{\_}gen}}{Thm.axiom:uninduct:ungen}}%
\nwixlogsorted{i}{{\nwixident{Thm.axiom{\_}induct{\_}min}}{Thm.axiom:uninduct:unmin}}%
\nwixlogsorted{i}{{\nwixident{Thm.axiom{\_}intersection}}{Thm.axiom:unintersection}}%
\nwixlogsorted{i}{{\nwixident{Thm.axiom{\_}juxt}}{Thm.axiom:unjuxt}}%
\nwixlogsorted{i}{{\nwixident{Thm.axiom{\_}pair}}{Thm.axiom:unpair}}%
\nwixlogsorted{i}{{\nwixident{Thm.axiom{\_}preimage}}{Thm.axiom:unpreimage}}%
\nwixlogsorted{i}{{\nwixident{Thm.axiom{\_}proj1}}{Thm.axiom:unproj1}}%
\nwixlogsorted{i}{{\nwixident{Thm.axiom{\_}proj2}}{Thm.axiom:unproj2}}%
\nwixlogsorted{i}{{\nwixident{Thm.axiom{\_}recur{\_}base}}{Thm.axiom:unrecur:unbase}}%
\nwixlogsorted{i}{{\nwixident{Thm.axiom{\_}recur{\_}rec}}{Thm.axiom:unrecur:unrec}}%
\nwixlogsorted{i}{{\nwixident{Thm.axiom{\_}recur{\_}zero}}{Thm.axiom:unrecur:unzero}}%
\nwixlogsorted{i}{{\nwixident{Thm.axiom{\_}singleton}}{Thm.axiom:unsingleton}}%
\nwixlogsorted{i}{{\nwixident{Thm.axiom{\_}union}}{Thm.axiom:ununion}}%
\nwixlogsorted{i}{{\nwixident{Thm.axiom{\_}universe{\_}ind}}{Thm.axiom:ununiverse:unind}}%
\nwixlogsorted{i}{{\nwixident{Thm.Class{\_}ind}}{Thm.Class:unind}}%
\nwixlogsorted{i}{{\nwixident{Thm.concl}}{Thm.concl}}%
\nwixlogsorted{i}{{\nwixident{Thm.contr}}{Thm.contr}}%
\nwixlogsorted{i}{{\nwixident{Thm.disch}}{Thm.disch}}%
\nwixlogsorted{i}{{\nwixident{Thm.exists{\_}elim}}{Thm.exists:unelim}}%
\nwixlogsorted{i}{{\nwixident{Thm.exists{\_}intro}}{Thm.exists:unintro}}%
\nwixlogsorted{i}{{\nwixident{Thm.forall{\_}elim}}{Thm.forall:unelim}}%
\nwixlogsorted{i}{{\nwixident{Thm.forall{\_}intro}}{Thm.forall:unintro}}%
\nwixlogsorted{i}{{\nwixident{Thm.Fun{\_}ind}}{Thm.Fun:unind}}%
\nwixlogsorted{i}{{\nwixident{Thm.hyps}}{Thm.hyps}}%
\nwixlogsorted{i}{{\nwixident{Thm.iff}}{Thm.iff}}%
\nwixlogsorted{i}{{\nwixident{Thm.iff{\_}elim{\_}l}}{Thm.iff:unelim:unl}}%
\nwixlogsorted{i}{{\nwixident{Thm.iff{\_}elim{\_}r}}{Thm.iff:unelim:unr}}%
\nwixlogsorted{i}{{\nwixident{Thm.iff{\_}intro}}{Thm.iff:unintro}}%
\nwixlogsorted{i}{{\nwixident{Thm.Ind{\_}ind}}{Thm.Ind:unind}}%
\nwixlogsorted{i}{{\nwixident{Thm.mk{\_}axiom}}{Thm.mk:unaxiom}}%
\nwixlogsorted{i}{{\nwixident{Thm.modus{\_}ponens}}{Thm.modus:unponens}}%
\nwixlogsorted{i}{{\nwixident{Thm.not{\_}elim}}{Thm.not:unelim}}%
\nwixlogsorted{i}{{\nwixident{Thm.not{\_}intro}}{Thm.not:unintro}}%
\nwixlogsorted{i}{{\nwixident{Thm.or{\_}elim}}{Thm.or:unelim}}%
\nwixlogsorted{i}{{\nwixident{Thm.or{\_}intro{\_}l}}{Thm.or:unintro:unl}}%
\nwixlogsorted{i}{{\nwixident{Thm.or{\_}intro{\_}r}}{Thm.or:unintro:unr}}%
\nwixlogsorted{i}{{\nwixident{Thm.t}}{Thm.t}}%
\nwixlogsorted{i}{{\nwixident{Thm.undisch}}{Thm.undisch}}%
\nwixlogsorted{i}{{\nwixident{Unif.y}}{Unif.y}}%
\nwixlogsorted{i}{{\nwixident{Unif.y{\_}complements}}{Unif.y:uncomplements}}%
\nwixlogsorted{i}{{\nwixident{Unif.y{\_}literals}}{Unif.y:unliterals}}%

